\chapt{युद्धकाण्डः}

\sect{प्रथमः सर्गः}

\twolineshloka
{श्रुत्वा हनुमतो वाक्यं यथावदभिभाषितम्}
{रामः प्रीतिसमायुक्तो वाक्यमुत्तरमब्रवीत्} %||6-1-1||

\twolineshloka
{कृतं हनुमता कार्यं सुमहद्भुवि दुष्करम्}
{मनसापि यदन्येन न शक्यं धरणीतले} %||6-1-2||

\twolineshloka
{न हि तं परिपश्यामि यस्तरेत महार्णवम्}
{अन्यत्र गरुणाद्वायोरन्यत्र च हनूमतः} %||6-1-3||

\twolineshloka
{देवदानवयक्षाणां गन्धर्वोरगरक्षसाम्}
{अप्रधृष्यां पुरीं लङ्कां रावणेन सुरक्षिताम्} %||6-1-4||

\threelineshloka
{प्रविष्टः सत्त्वमाश्रित्य जीवन्को नाम निष्क्रमेत्}
{को विशेत्सुदुराधर्षां राक्षसैश्च सुरक्षिताम्}
{यो वीर्यबलसम्पन्नो न समः स्याद्धनूमतः} %||6-1-5||

\twolineshloka
{भृत्यकार्यं हनुमता सुग्रीवस्य कृतं महत्}
{एवं विधाय स्वबलं सदृशं विक्रमस्य च} %||6-1-6||

\twolineshloka
{यो हि भृत्यो नियुक्तः सन्भर्त्रा कर्मणि दुष्करे}
{कुर्यात्तदनुरागेण तमाहुः पुरुषोत्तमम्} %||6-1-7||

\twolineshloka
{नियुक्तो नृपतेः कार्यं न कुर्याद्यः समाहितः}
{भृत्यो युक्तः समर्थश्च तमाहुः पुरुषाधमम्} %||6-1-8||

\twolineshloka
{तन्नियोगे नियुक्तेन कृतं कृत्यं हनूमता}
{न चात्मा लघुतां नीतः सुग्रीवश्चापि तोषितः} %||6-1-9||

\twolineshloka
{अहं च रघुवंशश्च लक्ष्मणश्च महाबलः}
{वैदेह्या दर्शनेनाद्य धर्मतः परिरक्षिताः} %||6-1-10||

\twolineshloka
{इदं तु मम दीनस्या मनो भूयः प्रकर्षति}
{यदिहास्य प्रियाख्यातुर्न कुर्मि सदृशं प्रियम्} %||6-1-11||

\twolineshloka
{एष सर्वस्वभूतस्तु परिष्वङ्गो हनूमतः}
{मया कालमिमं प्राप्य दत्तस्तस्य महात्मनः} %||6-1-12||

\twolineshloka
{सर्वथा सुकृतं तावत्सीतायाः परिमार्गणम्}
{सागरं तु समासाद्य पुनर्नष्टं मनो मम} %||6-1-13||

\twolineshloka
{कथं नाम समुद्रस्य दुष्पारस्य महाम्भसः}
{हरयो दक्षिणं पारं गमिष्यन्ति समाहिताः} %||6-1-14||

\twolineshloka
{यद्यप्येष तु वृत्तान्तो वैदेह्या गदितो मम}
{समुद्रपारगमने हरीणां किमिवोत्तरम्} %||6-1-15||

\twolineshloka
{इत्युक्त्वा शोकसम्भ्रान्तो रामः शत्रुनिबर्हणः}
{हनूमन्तं महाबाहुस्ततो ध्यानमुपागमत्} %||6-2-16||


{॥इत्यार्षे श्रीमद्रामायणे वाल्मीकीये आदिकाव्ये युद्धकाण्डे प्रथमः सर्गः॥}

\sect{द्वितीयः सर्गः}

\twolineshloka
{तं तु शोकपरिद्यूनं रामं दशरथात्मजम्}
{उवाच वचनं श्रीमान्सुग्रीवः शोकनाशनम्} %||6-2-1||

\twolineshloka
{किं त्वं सन्तप्यसे वीर यथान्यः प्राकृतस्तथा}
{मैवं भूस्त्यज सन्तापं कृतघ्न इव सौहृदम्} %||6-2-2||

\twolineshloka
{सन्तापस्य च ते स्थानं न हि पश्यामि राघव}
{प्रवृत्तावुपलब्धायां ज्ञाते च निलये रिपोः} %||6-2-3||

\twolineshloka
{धृतिमाञ्शास्त्रवित्प्राज्ञः पण्डितश्चासि राघव}
{त्यजेमां पापिकां बुद्धिं कृत्वात्मेवार्थदूषणीम्} %||6-2-4||

\twolineshloka
{समुद्रं लङ्घयित्वा तु महानक्रसमाकुलम्}
{लङ्कामारोहयिष्यामो हनिष्यामश्च ते रिपुम्} %||6-2-5||

\twolineshloka
{निरुत्साहस्य दीनस्य शोकपर्याकुलात्मनः}
{सर्वार्था व्यवसीदन्ति व्यसनं चाधिगच्छति} %||6-2-6||

\twolineshloka
{इमे शूराः समर्थाश्च सर्वे नो हरियूथपाः}
{त्वत्प्रियार्थं कृतोत्साहाः प्रवेष्टुमपि पावकम्} %||6-2-7||

\twolineshloka
{एषां हर्षेण जानामि तर्कश्चास्मिन्दृढो मम}
{विक्रमेण समानेष्ये सीतां हत्वा यथा रिपुम्} %||6-2-8||

\twolineshloka
{सेतुरत्र यथा वध्येद्यथा पश्येम तां पुरीम्}
{तस्य राक्षसराजस्य तथा त्वं कुरु राघव} %||6-2-9||

\twolineshloka
{दृष्ट्वा तां हि पुरीं लङ्कां त्रिकूटशिखरे स्थिताम्}
{हतं च रावणं युद्धे दर्शनादुपधारय} %||6-2-10||

\twolineshloka
{सेतुबद्धः समुद्रे च यावल्लङ्का समीपतः}
{सर्वं तीर्णं च वै सैन्यं जितमित्युपधार्यताम्} %||6-2-11||

\twolineshloka
{इमे हि समरे शूरा हरयः कामरूपिणः}
{तदलं विक्लवा बुद्धी राजन्सर्वार्थनाशनी} %||6-2-12||

\threelineshloka
{पुरुषस्य हि लोकेऽस्मिञ्शोकः शौर्यापकर्षणः}
{यत्तु कार्यं मनुष्येण शौण्डीर्यमवलम्बता}
{शूराणां हि मनुष्याणां त्वद्विधानां महात्मनाम्} %||6-2-13||

\twolineshloka
{विनष्टे वा प्रनष्टे वा शोकः सर्वार्थनाशनः}
{त्वं तु बुद्धिमतां श्रेष्ठः सर्वशास्त्रार्थकोविदः} %||6-2-14||

\twolineshloka
{मद्विधैः सचिवैः सार्थमरिं जेतुमिहार्हसि}
{न हि पश्याम्यहं कञ्चित्त्रिषु लोकेषु राघव} %||6-2-15||

\twolineshloka
{गृहीतधनुषो यस्ते तिष्ठेदभिमुखो रणे}
{वानरेषु समासक्तं न ते कार्यं विपत्स्यते} %||6-2-16||

\twolineshloka
{अचिराद्द्रक्ष्यसे सीतां तीर्त्वा सागरमक्षयम्}
{तदलं शोकमालम्ब्य क्रोधमालम्ब भूपते} %||6-2-17||

\twolineshloka
{निश्चेष्टाः क्षत्रिया मन्दाः सर्वे चण्डस्य बिभ्यति}
{लङ्गनार्थं च घोरस्य समुद्रस्य नदीपतेः} %||6-2-18||

\twolineshloka
{सहास्माभिरिहोपेतः सूक्ष्मबुद्धिर्विचारय}
{इमे हि समरे शूरा हरयः कामरूपिणः} %||6-2-19||

\twolineshloka
{तानरीन्विधमिष्यन्ति शिलापादपवृष्टिभिः}
{कथञ्चित्परिपश्यामस्ते वयं वरुणालयम्} %||6-2-20||

{किमुक्त्वा बहुधा चापि सर्वथा विजयी भवान्॥\devanumber{21}॥} %||6-3-21|| (Check)

\addtocounter{shlokacount}{1}


{॥इत्यार्षे श्रीमद्रामायणे वाल्मीकीये आदिकाव्ये युद्धकाण्डे द्वितीयः सर्गः॥}

\sect{तृतीयः सर्गः}

\twolineshloka
{सुग्रीवस्य वचः श्रुत्वा हेतुमत्परमार्थवित्}
{प्रतिजग्राह काकुत्स्थो हनूमन्तमथाब्रवीत्} %||6-3-1||

\twolineshloka
{तरसा सेतुबन्धेन सागरोच्छोषणेन वा}
{सर्वथा सुसमर्थोऽस्मि सागरस्यास्य लङ्घने} %||6-3-2||

\twolineshloka
{कति दुर्गाणि दुर्गाया लङ्कायास्तद्ब्रवीहि मे}
{ज्ञातुमिच्छामि तत्सर्वं दर्शनादिव वानर} %||6-3-3||

\twolineshloka
{बलस्य परिमाणं च द्वारदुर्गक्रियामपि}
{गुप्ति कर्म च लङ्काया रक्षसां सदनानि च} %||6-3-4||

\twolineshloka
{यथासुखं यथावच्च लङ्कायामसि दृष्टवान्}
{सरमाचक्ष्व तत्त्वेन सर्वथा कुशलो ह्यसि} %||6-3-5||

\twolineshloka
{श्रुत्वा रामस्य वचनं हनूमान्मारुतात्मजः}
{वाक्यं वाक्यविदां श्रेष्ठो रामं पुनरथाब्रवीत्} %||6-3-6||

\twolineshloka
{श्रूयतां सर्वमाख्यास्ये दुर्गकर्मविधानतः}
{गुप्ता पुरी यथा लङ्का रक्षिता च यथा बलैः} %||6-3-7||

\twolineshloka
{परां समृद्धिं लङ्कायाः सागरस्य च भीमताम्}
{विभागं च बलौघस्य निर्देशं वाहनस्य च} %||6-3-8||

\twolineshloka
{प्रहृष्टा मुदिता लङ्का मत्तद्विपसमाकुला}
{महती रथसम्पूर्णा रक्षोगणसमाकुला} %||6-3-9||

\twolineshloka
{दृढबद्धकवाटानि महापरिघवन्ति च}
{द्वाराणि विपुलान्यस्याश्चत्वारि सुमहान्ति च} %||6-3-10||

\twolineshloka
{वप्रेषूपलयन्त्राणि बलवन्ति महान्ति च}
{आगतं परसैन्यं तैस्तत्र प्रतिनिवार्यते} %||6-3-11||

\twolineshloka
{द्वारेषु संस्कृता भीमाः कालायसमयाः शिताः}
{शतशो रोचिता वीरैः शतघ्न्यो रक्षसां गणैः} %||6-3-12||

\twolineshloka
{सौवर्णश्च महांस्तस्याः प्राकारो दुष्प्रधर्षणः}
{मणिविद्रुमवैदूर्यमुक्ताविचरितान्तरः} %||6-3-13||

\twolineshloka
{सर्वतश्च महाभीमाः शीततोया महाशुभाः}
{अगाधा ग्राहवत्यश्च परिखा मीनसेविताः} %||6-3-14||

\twolineshloka
{द्वारेषु तासां चत्वारः सङ्क्रमाः परमायताः}
{यन्त्रैरुपेता बहुभिर्महद्भिर्दृढसन्धिभिः} %||6-3-15||

\twolineshloka
{त्रायन्ते सङ्क्रमास्तत्र परसैन्यागमे सति}
{यन्त्रैस्तैरवकीर्यन्ते परिखासु समन्ततः} %||6-3-16||

\twolineshloka
{एकस्त्वकम्प्यो बलवान्सङ्क्रमः सुमहादृढः}
{काञ्चनैर्बहुभिः स्तम्भैर्वेदिकाभिश्च शोभितः} %||6-3-17||

\twolineshloka
{स्वयं प्रकृतिसम्पन्नो युयुत्सू राम रावणः}
{उत्थितश्चाप्रमत्तश्च बलानामनुदर्शने} %||6-3-18||

\twolineshloka
{लङ्का पुरी निरालम्बा देवदुर्गा भयावहा}
{नादेयं पार्वतं वन्यं कृत्रिमं च चतुर्विधम्} %||6-3-19||

\twolineshloka
{स्थिता पारे समुद्रस्य दूरपारस्य राघव}
{नौपथश्चापि नास्त्यत्र निरादेशश्च सर्वतः} %||6-3-20||

\twolineshloka
{शैलाग्रे रचिता दुर्गा सा पूर्देवपुरोपमा}
{वाजिवारणसम्पूर्णा लङ्का परमदुर्जया} %||6-3-21||

\twolineshloka
{परिघाश्च शतघ्न्यश्च यन्त्राणि विविधानि च}
{शोभयन्ति पुरीं लङ्कां रावणस्य दुरात्मनः} %||6-3-22||

\twolineshloka
{अयुतं रक्षसामत्र पश्चिमद्वारमाश्रितम्}
{शूलहस्ता दुराधर्षाः सर्वे खड्गाग्रयोधिनः} %||6-3-23||

\twolineshloka
{नियुतं रक्षसामत्र दक्षिणद्वारमाश्रितम्}
{चतुरङ्गेण सैन्येन योधास्तत्राप्यनुत्तमाः} %||6-3-24||

\twolineshloka
{प्रयुतं रक्षसामत्र पूर्वद्वारं समाश्रितम्}
{चर्मखड्गधराः सर्वे तथा सर्वास्त्रकोविदाः} %||6-3-25||

\twolineshloka
{अर्बुदं रक्षसामत्र उत्तरद्वारमाश्रितम्}
{रथिनश्चाश्ववाहाश्च कुलपुत्राः सुपूजिताः} %||6-3-26||

\twolineshloka
{शतं शतसहस्राणां मध्यमं गुल्ममाश्रितम्}
{यातुधाना दुराधर्षाः साग्रकोटिश्च रक्षसाम्} %||6-3-27||

\twolineshloka
{ते मया सङ्क्रमा भग्नाः परिखाश्चावपूरिताः}
{दग्धा च नगरी लङ्का प्राकाराश्चावसादिताः} %||6-3-28||

\twolineshloka
{येन केन तु मार्गेण तराम वरुणालयम्}
{हतेति नगरी लङ्कां वानरैरवधार्यताम्} %||6-3-29||

\twolineshloka
{अङ्गदो द्विविदो मैन्दो जाम्बवान्पनसो नलः}
{नीलः सेनापतिश्चैव बलशेषेण किं तव} %||6-3-30||

\twolineshloka
{प्लवमाना हि गत्वा तां रावणस्य महापुरीम्}
{सप्रकारां सभवनामानयिष्यन्ति मैथिलीम्} %||6-3-31||

\twolineshloka
{एवमाज्ञापय क्षिप्रं बलानां सर्वसङ्ग्रहम्}
{मुहूर्तेन तु युक्तेन प्रस्थानमभिरोचय} %||6-4-32||


{॥इत्यार्षे श्रीमद्रामायणे वाल्मीकीये आदिकाव्ये युद्धकाण्डे तृतीयः सर्गः॥}

\sect{चतुर्थः सर्गः}

\twolineshloka
{श्रुत्वा हनूमतो वाक्यं यथावदनुपूर्वशः}
{ततोऽब्रवीन्महातेजा रामः सत्यपराक्रमः} %||6-4-1||

\twolineshloka
{यां निवेदयसे लङ्कां पुरीं भीमस्य रक्षसः}
{क्षिप्रमेनां वधिष्यामि सत्यमेतद्ब्रवीमि ते} %||6-4-2||

\twolineshloka
{अस्मिन्मुहूर्ते सुग्रीव प्रयाणमभिरोचये}
{युक्तो मुहूर्तो विजयः प्राप्तो मध्यं दिवाकरः} %||6-4-3||

\twolineshloka
{उत्तरा फल्गुनी ह्यद्य श्वस्तु हस्तेन योक्ष्यते}
{अभिप्रयाम सुग्रीव सर्वानीकसमावृताः} %||6-4-4||

\twolineshloka
{निमित्तानि च धन्यानि यानि प्रादुर्भवन्ति मे}
{निहत्य रावणं सीतामानयिष्यामि जानकीम्} %||6-4-5||

\twolineshloka
{उपरिष्टाद्धि नयनं स्फुरमाणमिदं मम}
{विजयं समनुप्राप्तं शंसतीव मनोरथम्} %||6-4-6||

\twolineshloka
{अग्रे यातु बलस्यास्य नीलो मार्गमवेक्षितुम्}
{वृतः शतसहस्रेण वानराणां तरस्विनाम्} %||6-4-7||

\twolineshloka
{फलमूलवता नील शीतकाननवारिणा}
{पथा मधुमता चाशु सेनां सेनापते नय} %||6-4-8||

\twolineshloka
{दूषयेयुर्दुरात्मानः पथि मूलफलोदकम्}
{राक्षसाः परिरक्षेथास्तेभ्यस्त्वं नित्यमुद्यतः} %||6-4-9||

\twolineshloka
{निम्नेषु वनदुर्गेषु वनेषु च वनौकसः}
{अभिप्लुत्याभिपश्येयुः परेषां निहतं बलम्} %||6-4-10||

\twolineshloka
{सागरौघनिभं भीममग्रानीकं महाबलाः}
{कपिसिंहा प्रकर्षन्तु शतशोऽथ सहस्रशः} %||6-4-11||

\twolineshloka
{गजश्च गिरिसङ्काशो गवयश्च महाबलः}
{गवाक्षश्चाग्रतो यान्तु गवां दृप्ता इवर्षभाः} %||6-4-12||

\twolineshloka
{यातु वानरवाहिन्या वानरः प्लवतां पतिः}
{पालयन्दक्षिणं पार्श्वमृषभो वानरर्षभः} %||6-4-13||

\twolineshloka
{गन्धहस्तीव दुर्धर्षस्तरस्वी गन्धमादनः}
{यातु वानरवाहिन्याः सव्यं पार्श्वमधिष्ठितः} %||6-4-14||

\twolineshloka
{यास्यामि बलमध्येऽहं बलौघमभिहर्षयन्}
{अधिरुह्य हनूमन्तमैरावतमिवेश्वरः} %||6-4-15||

\twolineshloka
{अङ्गदेनैष संयातु लक्ष्मणश्चान्तकोपमः}
{सार्वभौमेण भूतेशो द्रविणाधिपतिर्यथा} %||6-4-16||

\twolineshloka
{जाम्बवांश्च सुषेणश्च वेगदर्शी च वानरः}
{ऋक्षराजो महासत्त्वः कुक्षिं रक्षन्तु ते त्रयः} %||6-4-17||

\twolineshloka
{राघवस्य वचः श्रुत्वा सुग्रीवो वाहिनीपतिः}
{व्यादिदेश महावीर्यान्वानरान्वानरर्षभः} %||6-4-18||

\twolineshloka
{ते वानरगणाः सर्वे समुत्पत्य युयुत्सवः}
{गुहाभ्यः शिखरेभ्यश्च आशु पुप्लुविरे तदा} %||6-4-19||

\twolineshloka
{ततो वानरराजेन लक्ष्मणेन च पूजितः}
{जगाम रामो धर्मात्मा ससैन्यो दक्षिणां दिशम्} %||6-4-20||

\twolineshloka
{शतैः शतसहस्रैश्च कोटीभिरयुतैरपि}
{वारणाभिश्च हरिभिर्ययौ परिवृतस्तदा} %||6-4-21||

{तं यान्तमनुयाति स्म महती हरिवाहिनी॥\devanumber{22}॥} %||6-4-22|| (Check)

\addtocounter{shlokacount}{1}

\threelineshloka
{हृष्टाः प्रमुदिताः सर्वे सुग्रीवेणाभिपालिताः}
{आप्लवन्तः प्लवन्तश्च गर्जन्तश्च प्लवङ्गमाः}
{क्ष्वेलन्तो निनदन्तश्च जग्मुर्वै दक्षिणां दिशम्} %||6-4-23||

\twolineshloka
{भक्षयन्तः सुगन्धीनि मधूनि च फलानि च}
{उद्वहन्तो महावृक्षान्मञ्जरीपुञ्जधारिणः} %||6-4-24||

\twolineshloka
{अन्योन्यं सहसा दृष्टा निर्वहन्ति क्षिपन्ति च}
{पतन्तश्चोत्पतन्त्यन्ये पातयन्त्यपरे परान्} %||6-4-25||

\twolineshloka
{रावणो नो निहन्तव्यः सर्वे च रजनीचराः}
{इति गर्जन्ति हरयो राघवस्य समीपतः} %||6-4-26||

\twolineshloka
{पुरस्तादृषभो वीरो नीलः कुमुद एव च}
{पथानं शोधयन्ति स्म वानरैर्बहुभिः सह} %||6-4-27||

\twolineshloka
{मध्ये तु राजा सुग्रीवो रामो लक्ष्मण एव च}
{बहुभिर्बलिभिर्भीमैर्वृताः शत्रुनिबर्हणः} %||6-4-28||

\twolineshloka
{हरिः शतबलिर्वीरः कोटीभिर्दशभिर्वृतः}
{सर्वामेको ह्यवष्टभ्य ररक्ष हरिवाहिनीम्} %||6-4-29||

\twolineshloka
{कोटीशतपरीवारः केसरी पनसो गजः}
{अर्कश्चातिबलः पार्श्वमेकं तस्याभिरक्षति} %||6-4-30||

\twolineshloka
{सुषेणो जाम्बवांश्चैव ऋक्षैर्बहुभिरावृतः}
{सुग्रीवं पुरतः कृत्वा जघनं संररक्षतुः} %||6-4-31||

\twolineshloka
{तेषां सेनापतिर्वीरो नीलो वानरपुङ्गवः}
{सम्पतन्पततां श्रेष्ठस्तद्बलं पर्यपालयत्} %||6-4-32||

\twolineshloka
{दरीमिखः प्रजङ्घश्च जम्भोऽथ रभसः कपिः}
{सर्वतश्च ययुर्वीरास्त्वरयन्तः प्लवङ्गमान्} %||6-4-33||

\twolineshloka
{एवं ते हरिशार्दूला गच्छन्तो बलदर्पिताः}
{अपश्यंस्ते गिरिश्रेष्ठं सह्यं द्रुमलतायुतम्} %||6-4-34||

\twolineshloka
{सागरौघनिभं भीमं तद्वानरबलं महत्}
{निःससर्प महाघोषं भीमवेग इवार्णवः} %||6-4-35||

\twolineshloka
{तस्य दाशरथेः पार्श्वे शूरास्ते कपिकुञ्जराः}
{तूर्णमापुप्लुवुः सर्वे सदश्वा इव चोदिताः} %||6-4-36||

\twolineshloka
{कपिभ्यामुह्यमानौ तौ शुशुभते नरर्षभौ}
{महद्भ्यामिव संस्पृष्टौ ग्राहाभ्यां चन्द्रभास्करौ} %||6-4-37||

\twolineshloka
{तमङ्गदगतो रामं लक्ष्मणः शुभया गिरा}
{उवाच प्रतिपूर्णार्थः स्मृतिमान्प्रतिभानवान्} %||6-4-38||

\twolineshloka
{हृतामवाप्य वैदेहीं क्षिप्रं हत्वा च रावणम्}
{समृद्धार्थः समृद्धार्थामयोध्यां प्रतियास्यसि} %||6-4-39||

\twolineshloka
{महान्ति च निमित्तानि दिवि भूमौ च राघव}
{शुभान्ति तव पश्यामि सर्वाण्येवार्थसिद्धये} %||6-4-40||

\twolineshloka
{अनु वाति शुभो वायुः सेनां मृदुहितः सुखः}
{पूर्णवल्गुस्वराश्चेमे प्रवदन्ति मृगद्विजाः} %||6-4-41||

\twolineshloka
{प्रसन्नाश्च दिशः सर्वा विमलश्च दिवाकरः}
{उशना च प्रसन्नार्चिरनु त्वां भार्गवो गतः} %||6-4-42||

\twolineshloka
{ब्रह्मराशिर्विशुद्धश्च शुद्धाश्च परमर्षयः}
{अर्चिष्मन्तः प्रकाशन्ते ध्रुवं सर्वे प्रदक्षिणम्} %||6-4-43||

\twolineshloka
{त्रिशङ्कुर्विमलो भाति राजर्षिः सपुरोहितः}
{पितामहवरोऽस्माकमिष्क्वाकूणां महात्मनाम्} %||6-4-44||

\twolineshloka
{विमले च प्रकाशेते विशाखे निरुपद्रवे}
{नक्षत्रं परमस्माकमिक्ष्वाकूणां महात्मनाम्} %||6-4-45||

\twolineshloka
{नैरृतं नैरृतानां च नक्षत्रमभिपीड्यते}
{मूलं मूलवता स्पृष्टं धूप्यते धूमकेतुना} %||6-4-46||

\twolineshloka
{सरं चैतद्विनाशाय राक्षसानामुपस्थितम्}
{काले कालगृहीतानां नकत्रं ग्रहपीडितम्} %||6-4-47||

\twolineshloka
{प्रसन्नाः सुरसाश्चापो वनानि फलवन्ति च}
{प्रवान्त्यभ्यधिकं गन्धा यथर्तुकुसुमा द्रुमाः} %||6-4-48||

\twolineshloka
{व्यूढानि कपिसैन्यानि प्रकाशन्तेऽधिकं प्रभो}
{देवानामिव सैन्यानि सङ्ग्रामे तारकामये} %||6-4-49||

\twolineshloka
{एवमार्य समीक्ष्यैतान्प्रीतो भवितुमर्हसि}
{इति भ्रातरमाश्वास्य हृष्टः सौमित्रिरब्रवीत्} %||6-4-50||

\twolineshloka
{अथावृत्य महीं कृत्स्नां जगाम महती चमूः}
{ऋक्षवानरशार्दूलैर्नखदंष्ट्रायुधैर्वृता} %||6-4-51||

\twolineshloka
{कराग्रैश्चरणाग्रैश्च वानरैरुद्धतं रजः}
{भौममन्तर्दधे लोकं निवार्य सवितुः प्रभाम्} %||6-4-52||

\twolineshloka
{सा स्म याति दिवारात्रं महती हरिवाहिनी}
{हृष्टप्रमुदिता सेना सुग्रीवेणाभिरक्षिता} %||6-4-53||

\twolineshloka
{वनरास्त्वरितं यान्ति सर्वे युद्धाभिनन्दनः}
{मुमोक्षयिषवः सीतां मुहूर्तं क्वापि नासत} %||6-4-54||

\twolineshloka
{ततः पादपसम्बाधं नानामृगसमाकुलम्}
{सह्यपर्वतमासेदुर्मलयं च मही धरम्} %||6-4-55||

\twolineshloka
{काननानि विचित्राणि नदीप्रस्रवणानि च}
{पश्यन्नपि ययौ रामः सह्यस्य मलयस्य च} %||6-4-56||

\twolineshloka
{चम्पकांस्तिलकांश्चूतानशोकान्सिन्दुवारकान्}
{करवीरांश्च तिमिशान्भञ्जन्ति स्म प्लवङ्गमाः} %||6-4-57||

\twolineshloka
{फलान्यमृतगन्धीनि मूलानि कुसुमानि च}
{बुभुजुर्वानरास्तत्र पादपानां बलोत्कटाः} %||6-4-58||

\twolineshloka
{द्रोणमात्रप्रमाणानि लम्बमानानि वानराः}
{ययुः पिबन्तो हृष्टास्ते मधूनि मधुपिङ्गलाः} %||6-4-59||

\twolineshloka
{पादपानवभञ्जन्तो विकर्षन्तस्तथा लताः}
{विधमन्तो गिरिवरान्प्रययुः प्लवगर्षभाः} %||6-4-60||

\twolineshloka
{वृक्षेभ्योऽन्ये तु कपयो नर्दन्तो मधुदर्पिताः}
{अन्ये वृक्षान्प्रपद्यन्ते प्रपतन्त्यपि चापरे} %||6-4-61||

\twolineshloka
{बभूव वसुधा तैस्तु सम्पूर्णा हरिपुङ्गवैः}
{यथा कमलकेदारैः पक्वैरिव वसुन्धरा} %||6-4-62||

\twolineshloka
{महेन्द्रमथ सम्प्राप्य रामो राजीवलोचनः}
{अध्यारोहन्महाबाहुः शिखरं द्रुमभूषितम्} %||6-4-63||

\twolineshloka
{ततः शिखरमारुह्य रामो दशरथात्मजः}
{कूर्ममीनसमाकीर्णमपश्यत्सलिलाशयम्} %||6-4-64||

\twolineshloka
{ते सह्यं समतिक्रम्य मलयं च महागिरिम्}
{आसेदुरानुपूर्व्येण समुद्रं भीमनिःस्वनम्} %||6-4-65||

\twolineshloka
{अवरुह्य जगामाशु वेलावनमनुत्तमम्}
{रामो रमयतां श्रेष्ठः ससुग्रीवः सलक्ष्मणः} %||6-4-66||

\twolineshloka
{अथ धौतोपलतलां तोयौघैः सहसोत्थितैः}
{वेलामासाद्य विपुलां रामो वचनमब्रवीत्} %||6-4-67||

\twolineshloka
{एते वयमनुप्राप्ताः सुग्रीव वरुणालयम्}
{इहेदानीं विचिन्ता सा या न पूर्वं समुत्थिता} %||6-4-68||

\twolineshloka
{अतः परमतीरोऽयं सागरः सरितां पति}
{न चायमनुपायेन शक्यस्तरितुमर्णवः} %||6-4-69||

\twolineshloka
{तदिहैव निवेशोऽस्तु मन्त्रः प्रस्तूयतामिह}
{यथेदं वानरबलं परं पारमवाप्नुयात्} %||6-4-70||

\twolineshloka
{इतीव स महाबाहुः सीताहरणकर्शितः}
{रामः सागरमासाद्य वासमाज्ञापयत्तदा} %||6-4-71||

\threelineshloka
{सम्प्राप्तो मन्त्रकालो नः सागरस्येह लङ्घने}
{स्वां स्वां सेनां समुत्सृज्य मा च कश्चित्कुतो व्रजेत्}
{गच्छन्तु वानराः शूरा ज्ञेयं छन्नं भयं च नः} %||6-4-72||

\twolineshloka
{रामस्य वचनं श्रुत्वा सुग्रीवः सहलक्ष्मणः}
{सेनां न्यवेशयत्तीरे सागरस्य द्रुमायुते} %||6-4-73||

\twolineshloka
{विरराज समीपस्थं सागरस्य तु तद्बलम्}
{मधुपाण्डुजलः श्रीमान्द्वितीय इव सागरः} %||6-4-74||

\twolineshloka
{वेलावनमुपागम्य ततस्ते हरिपुङ्गवाः}
{विनिविष्टाः परं पारं काङ्क्षमाणा महोदधेः} %||6-4-75||

\twolineshloka
{सा महार्णवमासाद्य हृष्टा वानरवाहिनी}
{वायुवेगसमाधूतं पश्यमाना महार्णवम्} %||6-4-76||

\twolineshloka
{दूरपारमसम्बाधं रक्षोगणनिषेवितम्}
{पश्यन्तो वरुणावासं निषेदुर्हरियूथपाः} %||6-4-77||

\twolineshloka
{चण्डनक्रग्रहं घोरं क्षपादौ दिवसक्षये}
{चन्द्रोदये समाधूतं प्रतिचन्द्रसमाकुलम्} %||6-4-78||

\twolineshloka
{चण्डानिलमहाग्राहैः कीर्णं तिमितिमिङ्गिलैः}
{दीप्तभोगैरिवाक्रीर्णं भुजङ्गैर्वरुणालयम्} %||6-4-79||

\twolineshloka
{अवगाढं महासत्तैर्नानाशैलसमाकुलम्}
{दुर्गं द्रुगममार्गं तमगाधमसुरालयम्} %||6-4-80||

\twolineshloka
{मकरैर्नागभोगैश्च विगाढा वातलोहिताः}
{उत्पेतुश्च निपेतुश्च प्रवृद्धा जलराशयः} %||6-4-81||

\twolineshloka
{अग्निचूर्णमिवाविद्धं भास्कराम्बुमनोरगम्}
{सुरारिविषयं घोरं पातालविषमं सदा} %||6-4-82||

\twolineshloka
{सागरं चाम्बरप्रख्यमम्बरं सागरोपमम्}
{सागरं चाम्बरं चेति निर्विशेषमदृश्यत} %||6-4-83||

\twolineshloka
{सम्पृक्तं नभसा ह्यम्भः सम्पृक्तं च नभोऽम्भसा}
{तादृग्रूपे स्म दृश्येते तारा रत्नसमाकुले} %||6-4-84||

\twolineshloka
{समुत्पतितमेघस्य वीच्चि मालाकुलस्य च}
{विशेषो न द्वयोरासीत्सागरस्याम्बरस्य च} %||6-4-85||

\twolineshloka
{अन्योन्यैराहताः सक्ताः सस्वनुर्भीमनिःस्वनाः}
{ऊर्मयः सिन्धुराजस्य महाभेर्य इवाहवे} %||6-4-86||

\twolineshloka
{रत्नौघजलसंनादं विषक्तमिव वायुना}
{उत्पतन्तमिव क्रुद्धं यादोगणसमाकुलम्} %||6-4-87||

\threelineshloka
{ददृशुस्ते महात्मानो वाताहतजलाशयम्}
{अनिलोद्धूतमाकाशे प्रवल्गतमिवोर्मिभिः}
{भ्रान्तोर्मिजलसंनादं प्रलोलमिव सागरम्} %||6-5-88||


{॥इत्यार्षे श्रीमद्रामायणे वाल्मीकीये आदिकाव्ये युद्धकाण्डे चतुर्थः सर्गः॥}

\sect{पञ्चमः सर्गः}

\twolineshloka
{सा तु नीलेन विधिवत्स्वारक्षा सुसमाहिता}
{सागरस्योत्तरे तीरे साधु सेना निवेशिता} %||6-5-1||

\twolineshloka
{मैन्दश्च द्विविधश्चोभौ तत्र वानरपुङ्गवौ}
{विचेरतुश्च तां सेनां रक्षार्थं सर्वतो दिशम्} %||6-5-2||

\twolineshloka
{निविष्टायां तु सेनायां तीरे नदनदीपतेः}
{पार्श्वस्थं लक्ष्मणं दृष्ट्वा रामो वचनमब्रवीत्} %||6-5-3||

\twolineshloka
{शोकश्च किल कालेन गच्छता ह्यपगच्छति}
{मम चापश्यतः कान्तामहन्यहनि वर्धते} %||6-5-4||

\twolineshloka
{न मे दुःखं प्रिया दूरे न मे दुःखं हृतेति च}
{एतदेवानुशोचामि वयोऽस्या ह्यतिवर्तते} %||6-5-5||

\twolineshloka
{वाहि वात यतः कन्या तां स्पृष्ट्वा मामपि स्पृश}
{त्वयि मे गात्रसंस्पर्शश्चन्द्रे दृष्टिसमागमः} %||6-5-6||

\twolineshloka
{तन्मे दहति गात्राणि विषं पीतमिवाशये}
{हा नाथेति प्रिया सा मां ह्रियमाणा यदब्रवीत्} %||6-5-7||

\twolineshloka
{तद्वियोगेन्धनवता तच्चिन्ताविपुलार्चिषा}
{रात्रिं दिवं शरीरं मे दह्यते मदनाग्निना} %||6-5-8||

\twolineshloka
{अवगाह्यार्णवं स्वप्स्ये सौमित्रे भवता विना}
{कथञ्चित्प्रज्वलन्कामः समासुप्तं जले दहेत्} %||6-5-9||

\twolineshloka
{बह्वेतत्कामयानस्य शक्यमेतेन जीवितुम्}
{यदहं सा च वामोरुरेकां धरणिमाश्रितौ} %||6-5-10||

\twolineshloka
{केदारस्येव केदारः सोदकस्य निरूदकः}
{उपस्नेहेन जीवामि जीवन्तीं यच्छृणोमि ताम्} %||6-5-11||

\twolineshloka
{कदा तु खलु सुस्शोणीं शतपत्रायतेक्षणाम्}
{विजित्य शत्रून्द्रक्ष्यामि सीतां स्फीतामिव श्रियम्} %||6-5-12||

\twolineshloka
{कदा नु चारुबिम्बौष्ठं तस्याः पद्ममिवाननम्}
{ईषदुन्नम्य पास्यामि रसायनमिवातुरः} %||6-5-13||

\twolineshloka
{तौ तस्याः संहतौ पीनौ स्तनौ तालफलोपमौ}
{कदा नु खलु सोत्कम्पौ हसन्त्या मां भजिष्यतः} %||6-5-14||

\twolineshloka
{सा नूनमसितापाङ्गी रक्षोमध्यगता सती}
{मन्नाथा नाथहीनेव त्रातारं नाधिगच्छति} %||6-5-15||

\twolineshloka
{कदा विक्षोभ्य रक्षांसि सा विधूयोत्पतिष्यति}
{विधूय जलदान्नीलाञ्शशिलेखा शरत्स्विव} %||6-5-16||

\twolineshloka
{स्वभावतनुका नूनं शोकेनानशनेन च}
{भूयस्तनुतरा सीता देशकालविपर्ययात्} %||6-5-17||

\twolineshloka
{कदा नु राक्षसेन्द्रस्य निधायोरसि सायकान्}
{सीतां प्रत्याहरिष्यामि शोकमुत्सृज्य मानसम्} %||6-5-18||

\twolineshloka
{कदा नु खलु मां साध्वी सीतामरसुतोपमा}
{सोत्कण्ठा कण्ठमालम्ब्य मोक्ष्यत्यानन्दजं जलम्} %||6-5-19||

\twolineshloka
{कदा शोकमिमं घोरं मैथिली विप्रयोगजम्}
{सहसा विप्रमोक्ष्यामि वासः शुक्लेतरं यथा} %||6-5-20||

\twolineshloka
{एवं विलपतस्तस्य तत्र रामस्य धीमतः}
{दिनक्षयान्मन्दवपुर्भास्करोऽस्तमुपागमत्} %||6-5-21||

\twolineshloka
{आश्वासितो लक्ष्मणेन रामः सन्ध्यामुपासत}
{स्मरन्कमलपत्राक्षीं सीतां शोकाकुलीकृतः} %||6-6-22||


{॥इत्यार्षे श्रीमद्रामायणे वाल्मीकीये आदिकाव्ये युद्धकाण्डे पञ्चमः सर्गः॥}

\sect{षष्ठः सर्गः}

\threelineshloka
{लङ्कायां तु कृतं कर्म घोरं दृष्ट्वा भवावहम्}
{राक्षसेन्द्रो हनुमता शक्रेणेव महात्मना}
{अब्रवीद्राक्षसान्सर्वान्ह्रिया किञ्चिदवाङ्मुखः} %||6-6-1||

\twolineshloka
{धर्षिता च प्रविष्टा च लङ्का दुष्प्रसहा पुरी}
{तेन वानरमात्रेण दृष्टा सीता च जानकी} %||6-6-2||

\twolineshloka
{प्रसादो धर्षितश्चैत्यः प्रवरा राक्षसा हताः}
{आविला च पुरी लङ्का सर्वा हनुमता कृता} %||6-6-3||

\twolineshloka
{किं करिष्यामि भद्रं वः किं वा युक्तमनन्तरम्}
{उच्यतां नः समर्थं यत्कृतं च सुकृतं भवेत्} %||6-6-4||

\twolineshloka
{मन्त्रमूलं हि विजयं प्राहुरार्या मनस्विनः}
{तस्माद्वै रोचये मन्त्रं रामं प्रति महाबलाः} %||6-6-5||

\twolineshloka
{त्रिविधाः पुरुषा लोके उत्तमाधममध्यमाः}
{तेषां तु समवेतानां गुणदोषं वदाम्यहम्} %||6-6-6||

\twolineshloka
{मन्त्रिभिर्हितसंयुक्तैः समर्थैर्मन्त्रनिर्णये}
{मित्रैर्वापि समानार्थैर्बान्धवैरपि वा हितैः} %||6-6-7||

\twolineshloka
{सहितो मन्त्रयित्वा यः कर्मारम्भान्प्रवर्तयेत्}
{दैवे च कुरुते यत्नं तमाहुः पुरुषोत्तमम्} %||6-6-8||

\twolineshloka
{एकोऽर्थं विमृशेदेको धर्मे प्रकुरुते मनः}
{एकः कार्याणि कुरुते तमाहुर्मध्यमं नरम्} %||6-6-9||

\twolineshloka
{गुणदोषावनिश्चित्य त्यक्त्वा दैवव्यपाश्रयम्}
{करिष्यामीति यः कार्यमुपेक्षेत्स नराधमः} %||6-6-10||

\twolineshloka
{यथेमे पुरुषा नित्यमुत्तमाधममध्यमाः}
{एवं मन्त्रोऽपि विज्ञेय उत्तमाधममध्यमः} %||6-6-11||

\twolineshloka
{ऐकमत्यमुपागम्य शास्त्रदृष्टेन चक्षुषा}
{मन्त्रिणो यत्र निरस्तास्तमाहुर्मन्त्रमुत्तमम्} %||6-6-12||

\twolineshloka
{बह्व्योऽपि मतयो गत्वा मन्त्रिणो ह्यर्थनिर्णये}
{पुनर्यत्रैकतां प्राप्तः स मन्त्रो मध्यमः स्मृतः} %||6-6-13||

\twolineshloka
{अन्योन्यमतिमास्थाय यत्र सम्प्रतिभाष्यते}
{न चैकमत्ये श्रेयोऽस्ति मन्त्रः सोऽधम उच्यते} %||6-6-14||

\twolineshloka
{तस्मात्सुमन्त्रितं साधु भवन्तो मन्त्रिसत्तमाः}
{कार्यं सम्प्रतिपद्यन्तामेतत्कृत्यतमं मम} %||6-6-15||

\twolineshloka
{वानराणां हि वीराणां सहस्रैः परिवारितः}
{रामोऽभ्येति पुरीं लङ्कामस्माकमुपरोधकः} %||6-6-16||

\twolineshloka
{तरिष्यति च सुव्यक्तं राघवः सागरं सुखम्}
{तरसा युक्तरूपेण सानुजः सबलानुगः} %||6-6-17||

\twolineshloka
{अस्मिन्नेवङ्गते कार्ये विरुद्धे वानरैः सह}
{हितं पुरे च सैन्ये च सर्वं सम्मन्त्र्यतां मम} %||6-7-18||


{॥इत्यार्षे श्रीमद्रामायणे वाल्मीकीये आदिकाव्ये युद्धकाण्डे षष्ठः सर्गः॥}

\sect{सप्तमः सर्गः}

\twolineshloka
{इत्युक्ता राक्षसेन्द्रेण राक्षसास्ते महाबलाः}
{ऊचुः प्राञ्जलयः सर्वे रावणं राक्षसेश्वरम्} %||6-7-1||

\twolineshloka
{राजन्परिघशक्त्यृष्टिशूलपट्टससङ्कुलम्}
{सुमहन्नो बलं कस्माद्विषादं भजते भवान्} %||6-7-2||

\twolineshloka
{कैलासशिखरावासी यक्षैर्बहुभिरावृतः}
{सुमहत्कदनं कृत्वा वश्यस्ते धनदः कृतः} %||6-7-3||

\twolineshloka
{स महेश्वरसख्येन श्लाघमानस्त्वया विभो}
{निर्जितः समरे रोषाल्लोकपालो महाबलः} %||6-7-4||

\twolineshloka
{विनिहत्य च यक्षौघान्विक्षोभ्य च विगृह्य च}
{त्वया कैलासशिखराद्विमानमिदमाहृतम्} %||6-7-5||

\twolineshloka
{मयेन दानवेन्द्रेण त्वद्भयात्सख्यमिच्छता}
{दुहिता तव भार्यार्थे दत्ता राक्षसपुङ्गव} %||6-7-6||

\twolineshloka
{दानवेन्द्रो मधुर्नाम वीर्योत्सिक्तो दुरासदः}
{विगृह्य वशमानीतः कुम्भीनस्याः सुखावहः} %||6-7-7||

\twolineshloka
{निर्जितास्ते महाबाहो नागा गत्वा रसातलम्}
{वासुकिस्तक्षकः शङ्खो जटी च वशमाहृताः} %||6-7-8||

\twolineshloka
{अक्षया बलवन्तश्च शूरा लब्धवराः पुनः}
{त्वया संवत्सरं युद्ध्वा समरे दानवा विभो} %||6-7-9||

\twolineshloka
{स्वबलं समुपाश्रित्य नीता वशमरिन्दम}
{मायाश्चाधिगतास्तत्र बहवो राक्षसाधिप} %||6-7-10||

\twolineshloka
{शूराश्च बलवन्तश्च वरुणस्य सुता रणे}
{निर्जितास्ते महाबाहो चतुर्विधबलानुगाः} %||6-7-11||

\twolineshloka
{मृत्युदण्डमहाग्राहं शाल्मलिद्वीपमण्डितम्}
{अवगाह्य त्वया राजन्यमस्य बलसागरम्} %||6-7-12||

\twolineshloka
{जयश्च विप्लुलः प्राप्तो मृत्युश्च प्रतिषेधितः}
{सुयुद्धेन च ते सर्वे लोकास्तत्र सुतोषिताः} %||6-7-13||

\twolineshloka
{क्षत्रियैर्बहुभिर्वीरैः शक्रतुल्यपराक्रमैः}
{आसीद्वसुमती पूर्णा महद्भिरिव पादपैः} %||6-7-14||

\twolineshloka
{तेषां वीर्यगुणोत्साहैर्न समो राघवो रणे}
{प्रसह्य ते त्वया राजन्हताः परमदुर्जयाः} %||6-7-15||

\twolineshloka
{राजन्नापदयुक्तेयमागता प्राकृताज्जनात्}
{हृदि नैव त्वया कार्या त्वं वधिष्यसि राघवम्} %||6-8-16||


{॥इत्यार्षे श्रीमद्रामायणे वाल्मीकीये आदिकाव्ये युद्धकाण्डे सप्तमः सर्गः॥}

\sect{अष्टमः सर्गः}

\twolineshloka
{ततो नीलाम्बुदनिभः प्रहस्तो नाम राक्षसः}
{अब्रवीत्प्राञ्जलिर्वाक्यं शूरः सेनापतिस्तदा} %||6-8-1||

\twolineshloka
{देवदानवगन्धर्वाः पिशाचपतगोरगाः}
{न त्वां धर्षयितुं शक्ताः किं पुनर्वानरा रणे} %||6-8-2||

\twolineshloka
{सर्वे प्रमत्ता विश्वस्ता वञ्चिताः स्म हनूमता}
{न हि मे जीवतो गच्छेज्जीवन्स वनगोचरः} %||6-8-3||

\twolineshloka
{सर्वां सागरपर्यन्तां सशैलवनकाननाम्}
{करोम्यवानरां भूमिमाज्ञापयतु मां भवान्} %||6-8-4||

\twolineshloka
{रक्षां चैव विधास्यामि वानराद्रजनीचर}
{नागमिष्यति ते दुःखं किञ्चिदात्मापराधजम्} %||6-8-5||

\twolineshloka
{अब्रवीच्च सुसङ्क्रुद्धो दुर्मुखो नाम राक्षसः}
{इदं न क्षमणीयं हि सर्वेषां नः प्रधर्षणम्} %||6-8-6||

\twolineshloka
{अयं परिभवो भूयः पुरस्यान्तःपुरस्य च}
{श्रीमतो राक्षसेन्द्रस्य वानरेन्द्रप्रधर्षणम्} %||6-8-7||

\twolineshloka
{अस्मिन्मुहूर्ते हत्वैको निवर्तिष्यामि वानरान्}
{प्रविष्टान्सागरं भीममम्बरं वा रसातलम्} %||6-8-8||

\twolineshloka
{ततोऽब्रवीत्सुसङ्क्रुद्धो वज्रदंष्ट्रो महाबलः}
{प्रगृह्य परिघं घोरं मांसशोणितरूपितम्} %||6-8-9||

\twolineshloka
{किं वो हनुमता कार्यं कृपणेन तपस्विना}
{रामे तिष्ठति दुर्धर्षे सुग्रीवे सहलक्ष्मणे} %||6-8-10||

\twolineshloka
{अद्य रामं ससुग्रीवं परिघेण सलक्ष्मणम्}
{आगमिष्यामि हत्वैको विक्षोभ्य हरिवाहिनीम्} %||6-8-11||

\twolineshloka
{कौम्भकर्णिस्ततो वीरो निकुम्भो नाम वीर्यवान्}
{अब्रवीत्परमकुर्द्धो रावणं लोकरावणम्} %||6-8-12||

\twolineshloka
{सर्वे भवन्तस्तिष्ठन्तु महाराजेन सङ्गताः}
{अहमेको हनिष्यामि राघवं सहलक्ष्मणम्} %||6-8-13||

\twolineshloka
{ततो वज्रहनुर्नाम राक्षसः पर्वतोपमः}
{क्रुद्धः परिलिहन्वक्त्रं जिह्वया वाक्यमब्रवीत्} %||6-8-14||

\twolineshloka
{स्वैरं कुर्वन्तु कार्याणि भवन्तो विगतज्वराः}
{एकोऽहं भक्षयिष्यामि तान्सर्वान्हरियूथपान्} %||6-8-15||

\threelineshloka
{स्वस्थाः क्रीडन्तु निश्चिन्ताः पिबन्तु मधुवारुणीम्}
{अहमेको हनिष्यामि सुग्रीवं सहलक्ष्मणम्}
{साङ्गदं च हनूमन्तं रामं च रणकुञ्जरम्} %||6-9-16||


{॥इत्यार्षे श्रीमद्रामायणे वाल्मीकीये आदिकाव्ये युद्धकाण्डे अष्टमः सर्गः॥}

\sect{नवमः सर्गः}

\twolineshloka
{ततो निकुम्भो रभसः सूर्यशत्रुर्महाबलः}
{सुप्तघ्नो यज्ञकोपश्च महापार्श्वो महोअरः} %||6-9-1||

\twolineshloka
{अग्निकेतुश्च दुर्धर्षो रश्मिकेतुश्च राक्षसः}
{इन्द्रजिच्च महातेजा बलवान्रावणात्मजः} %||6-9-2||

\twolineshloka
{प्रहस्तोऽथ विरूपाक्षो वज्रदंष्ट्रो महाबलः}
{धूम्राक्षश्चातिकायश्च दुर्मुखश्चैव राक्षसः} %||6-9-3||

\twolineshloka
{परिघान्पट्टसान्प्रासाञ्शक्तिशूलपरश्वधान्}
{चापानि च सबाणानि खड्गांश्च विपुलाञ्शितान्} %||6-9-4||

\twolineshloka
{प्रगृह्य परमक्रुद्धाः समुत्पत्य च राक्षसाः}
{अब्रुवन्रावणं सर्वे प्रदीप्ता इव तेजसा} %||6-9-5||

\twolineshloka
{अद्य रामं वधिष्यामः सुग्रीवं च सलक्ष्मणम्}
{कृपणं च हनूमन्तं लङ्का येन प्रधर्षिता} %||6-9-6||

\twolineshloka
{तान्गृहीतायुधान्सर्वान्वारयित्वा विभीषणः}
{अब्रवीत्प्राञ्जलिर्वाक्यं पुनः प्रत्युपवेश्य तान्} %||6-9-7||

\twolineshloka
{अप्युपायैस्त्रिभिस्तात योऽर्थः प्राप्तुं न शक्यते}
{तस्य विक्रमकालांस्तान्युक्तानाहुर्मनीषिणः} %||6-9-8||

\twolineshloka
{प्रमत्तेष्वभियुक्तेषु दैवेन प्रहतेषु च}
{विक्रमास्तात सिध्यन्ति परीक्ष्य विधिना कृताः} %||6-9-9||

\twolineshloka
{अप्रमत्तं कथं तं तु विजिगीषुं बले स्थितम्}
{जितरोषं दुराधर्षं प्रधर्षयितुमिच्छथ} %||6-9-10||

\twolineshloka
{समुद्रं लङ्घयित्वा तु घोरं नदनदीपतिम्}
{कृतं हनुमता कर्म दुष्करं तर्कयेत कः} %||6-9-11||

\twolineshloka
{बलान्यपरिमेयानि वीर्याणि च निशाचराः}
{परेषां सहसावज्ञा न कर्तव्या कथञ्चन} %||6-9-12||

\twolineshloka
{किं च राक्षसराजस्य रामेणापकृतं पुरा}
{आजहार जनस्थानाद्यस्य भार्यां यशस्विनः} %||6-9-13||

\twolineshloka
{खरो यद्यतिवृत्तस्तु रामेण निहतो रणे}
{अवश्यं प्राणिनां प्राणा रक्षितव्या यथा बलम्} %||6-9-14||

\twolineshloka
{एतन्निमित्तं वैदेही भयं नः सुमहद्भवेत्}
{आहृता सा परित्याज्या कलहार्थे कृते न किम्} %||6-9-15||

\twolineshloka
{न नः क्षमं वीर्यवता तेन धर्मानुवर्तिना}
{वैरं निरर्थकं कर्तुं दीयतामस्य मैथिली} %||6-9-16||

\twolineshloka
{यावन्न सगजां साश्वां बहुरत्नसमाकुलाम्}
{पुरीं दारयते बाणैर्दीयतामस्य मैथिली} %||6-9-17||

\twolineshloka
{यावत्सुघोरा महती दुर्धर्षा हरिवाहिनी}
{नावस्कन्दति नो लङ्कां तावत्सीता प्रदीयताम्} %||6-9-18||

\twolineshloka
{विनश्येद्धि पुरी लङ्का शूराः सर्वे च राक्षसाः}
{रामस्य दयिता पत्नी न स्वयं यदि दीयते} %||6-9-19||

\twolineshloka
{प्रसादये त्वां बन्धुत्वात्कुरुष्व वचनं मम}
{हितं पथ्यं त्वहं ब्रूमि दीयतामस्य मैथिली} %||6-9-20||

\fourlineindentedshloka
{पुरा शरत्सूर्यमरीचिसंनिभान्}
{नवाग्रपुङ्खान्सुदृढान्नृपात्मजः}
{सृजत्यमोघान्विशिखान्वधाय ते}
{प्रदीयतां दाशरथाय मैथिली} %||6-9-21||

\fourlineindentedshloka
{त्यजस्व कोपं सुखधर्मनाशनम्}
{भजस्व धर्मं रतिकीर्तिवर्धनम्}
{प्रसीद जीवेम सपुत्रबान्धवाः}
{प्रदीयतां दाशरथाय मैथिली} %||6-10-22||


{॥इत्यार्षे श्रीमद्रामायणे वाल्मीकीये आदिकाव्ये युद्धकाण्डे नवमः सर्गः॥}

\sect{दशमः सर्गः}

\twolineshloka
{सुनिविष्टं हितं वाक्यमुक्तवन्तं विभीषणम्}
{अब्रवीत्परुषं वाक्यं रावणः कालचोदितः} %||6-10-1||

\twolineshloka
{वसेत्सह सपत्नेन क्रुद्धेनाशीविषेण वा}
{न तु मित्रप्रवादेन संवसेच्छत्रुसेविना} %||6-10-2||

\twolineshloka
{जानामि शीलं ज्ञातीनां सर्वलोकेषु राक्षस}
{हृष्यन्ति व्यसनेष्वेते ज्ञातीनां ज्ञातयः सदा} %||6-10-3||

\twolineshloka
{प्रधानं साधकं वैद्यं धर्मशीलं च राक्षस}
{ज्ञातयो ह्यवमन्यन्ते शूरं परिभवन्ति च} %||6-10-4||

\twolineshloka
{नित्यमन्योन्यसंहृष्टा व्यसनेष्वाततायिनः}
{प्रच्छन्नहृदया घोरा ज्ञातयस्तु भयावहाः} %||6-10-5||

\twolineshloka
{श्रूयन्ते हस्तिभिर्गीताः श्लोकाः पद्मवने क्वचित्}
{पाशहस्तान्नरान्दृष्ट्वा शृणु तान्गदतो मम} %||6-10-6||

\twolineshloka
{नाग्निर्नान्यानि शस्त्राणि न नः पाशा भयावहाः}
{घोराः स्वार्थप्रयुक्तास्तु ज्ञातयो नो भयावहाः} %||6-10-7||

\twolineshloka
{उपायमेते वक्ष्यन्ति ग्रहणे नात्र संशयः}
{कृत्स्नाद्भयाज्ज्ञातिभयं सुकष्टं विदितं च नः} %||6-10-8||

\twolineshloka
{विद्यते गोषु सम्पन्नं विद्यते ब्राह्मणे दमः}
{विद्यते स्त्रीषु चापल्यं विद्यते ज्ञातितो भयम्} %||6-10-9||

\twolineshloka
{ततो नेष्टमिदं सौम्य यदहं लोकसत्कृतः}
{ऐश्वर्यमभिजातश्च रिपूणां मूर्ध्नि च स्थितः} %||6-10-10||

\twolineshloka
{अन्यस्त्वेवंविधं ब्रूयाद्वाक्यमेतन्निशाचर}
{अस्मिन्मुहूर्ते न भवेत्त्वां तु धिक्कुलपांसनम्} %||6-10-11||

\twolineshloka
{इत्युक्तः परुषं वाक्यं न्यायवादी विभीषणः}
{उत्पपात गदापाणिश्चतुर्भिः सह राक्षसैः} %||6-10-12||

\twolineshloka
{अब्रवीच्च तदा वाक्यं जातक्रोधो विभीषणः}
{अन्तरिक्षगतः श्रीमान्भ्रातरं राक्षसाधिपम्} %||6-10-13||

\twolineshloka
{स त्वं भ्रातासि मे राजन्ब्रूहि मां यद्यदिच्छसि}
{इदं तु परुषं वाक्यं न क्षमाम्यनृतं तव} %||6-10-14||

\twolineshloka
{सुनीतं हितकामेन वाक्यमुक्तं दशानन}
{न गृह्णन्त्यकृतात्मानः कालस्य वशमागताः} %||6-10-15||

\twolineshloka
{सुलभाः पुरुषा राजन्सततं प्रियवादिनः}
{अप्रियस्य तु पथ्यस्य वक्ता श्रोता च दुर्लभः} %||6-10-16||

\twolineshloka
{बद्धं कालस्य पाशेन सर्वभूतापहारिणा}
{न नश्यन्तमुपेक्षेयं प्रदीप्तं शरणं यथा} %||6-10-17||

\twolineshloka
{दीप्तपावकसङ्काशैः शितैः काञ्चनभूषणैः}
{न त्वामिच्छाम्यहं द्रष्टुं रामेण निहतं शरैः} %||6-10-18||

\twolineshloka
{शूराश्च बलवन्तश्च कृतास्त्राश्च रणाजिरे}
{कालाभिपन्ना सीदन्ति यथा वालुकसेतवः} %||6-10-19||

\twolineshloka
{आत्मानं सर्वथा रक्ष पुरीं चेमां सराक्षसाम्}
{स्वस्ति तेऽस्तु गमिष्यामि सुखी भव मया विना} %||6-10-20||

\fourlineindentedshloka
{निवार्यमाणस्य मया हितैषिणा}
{न रोचते ते वचनं निशाचर}
{परीतकाला हि गतायुषो नरा}
{हितं न गृह्णन्ति सुहृद्भिरीरितम्} %||6-11-21||


{॥इत्यार्षे श्रीमद्रामायणे वाल्मीकीये आदिकाव्ये युद्धकाण्डे दशमः सर्गः॥}

\sect{एकादशः सर्गः}

\twolineshloka
{इत्युक्त्वा परुषं वाक्यं रावणं रावणानुजः}
{आजगाम मुहूर्तेन यत्र रामः सलक्ष्मणः} %||6-11-1||

\twolineshloka
{तं मेरुशिखराकारं दीप्तामिव शतह्रदाम्}
{गगनस्थं महीस्थास्ते ददृशुर्वानराधिपाः} %||6-11-2||

\twolineshloka
{तमात्मपञ्चमं दृष्ट्वा सुग्रीवो वानराधिपः}
{वानरैः सह दुर्धर्षश्चिन्तयामास बुद्धिमान्} %||6-11-3||

\twolineshloka
{चिन्तयित्वा मुहूर्तं तु वानरांस्तानुवाच ह}
{हनूमत्प्रमुखान्सर्वानिदं वचनमुत्तमम्} %||6-11-4||

\twolineshloka
{एष सर्वायुधोपेतश्चतुर्भिः सह राक्षसैः}
{राक्षसोऽभ्येति पश्यध्वमस्मान्हन्तुं न संशयः} %||6-11-5||

\twolineshloka
{सुग्रीवस्य वचः श्रुत्वा सर्वे ते वानरोत्तमाः}
{सालानुद्यम्य शैलांश्च इदं वचनमब्रुवन्} %||6-11-6||

\twolineshloka
{शीघ्रं व्यादिश नो राजन्वधायैषां दुरात्मनाम्}
{निपतन्तु हताश्चैते धरण्यामल्पजीविताः} %||6-11-7||

\twolineshloka
{तेषां सम्भाषमाणानामन्योन्यं स विभीषणः}
{उत्तरं तीरमासाद्य खस्थ एव व्यतिष्ठत} %||6-11-8||

\twolineshloka
{उवाच च महाप्राज्ञः स्वरेण महता महान्}
{सुग्रीवं तांश्च सम्प्रेक्ष्य खस्थ एव विभीषणः} %||6-11-9||

\twolineshloka
{रावणो नाम दुर्वृत्तो राक्षसो राक्षसेश्वरः}
{तस्याहमनुजो भ्राता विभीषण इति श्रुतः} %||6-11-10||

\twolineshloka
{तेन सीता जनस्थानाद्धृता हत्वा जटायुषम्}
{रुद्ध्वा च विवशा दीना राक्षसीभिः सुरक्षिता} %||6-11-11||

\twolineshloka
{तमहं हेतुभिर्वाक्यैर्विविधैश्च न्यदर्शयम्}
{साधु निर्यात्यतां सीता रामायेति पुनः पुनः} %||6-11-12||

\twolineshloka
{स च न प्रतिजग्राह रावणः कालचोदितः}
{उच्यमानो हितं वाक्यं विपरीत इवौषधम्} %||6-11-13||

\twolineshloka
{सोऽहं परुषितस्तेन दासवच्चावमानितः}
{त्यक्त्वा पुत्रांश्च दारांश्च राघवं शरणं गतः} %||6-11-14||

\twolineshloka
{सर्वलोकशरण्याय राघवाय महात्मने}
{निवेदयत मां क्षिप्रं विभीषणमुपस्थितम्} %||6-11-15||

\twolineshloka
{एतत्तु वचनं श्रुत्वा सुग्रीवो लघुविक्रमः}
{लक्ष्मणस्याग्रतो रामं संरब्धमिदमब्रवीत्} %||6-11-16||

\twolineshloka
{रावणस्यानुजो भ्राता विभीषण इति श्रुतः}
{चतुर्भिः सह रक्षोभिर्भवन्तं शरणं गतः} %||6-11-17||

\twolineshloka
{रावणेन प्रणिहितं तमवेहि विभीषणम्}
{तस्याहं निग्रहं मन्ये क्षमं क्षमवतां वर} %||6-11-18||

\twolineshloka
{राक्षसो जिह्मया बुद्ध्या सन्दिष्टोऽयमुपस्थितः}
{प्रहर्तुं मायया छन्नो विश्वस्ते त्वयि राघव} %||6-11-19||

\twolineshloka
{बध्यतामेष तीव्रेण दण्डेन सचिवैः सह}
{रावणस्य नृशंसस्य भ्राता ह्येष विभीषणः} %||6-11-20||

\twolineshloka
{एवमुक्त्वा तु तं रामं संरब्धो वाहिनीपतिः}
{वाक्यज्ञो वाक्यकुशलं ततो मौनमुपागमत्} %||6-11-21||

\twolineshloka
{सुग्रीवस्य तु तद्वाक्यं श्रुत्वा रामो महाबलः}
{समीपस्थानुवाचेदं हनूमत्प्रमुखान्हरीन्} %||6-11-22||

\twolineshloka
{यदुक्तं कपिराजेन रावणावरजं प्रति}
{वाक्यं हेतुमदत्यर्थं भवद्भिरपि तच्छ्रुतम्} %||6-11-23||

\twolineshloka
{सुहृदा ह्यर्थकृच्छेषु युक्तं बुद्धिमता सता}
{समर्थेनापि सन्देष्टुं शाश्वतीं भूतिमिच्छता} %||6-11-24||

\twolineshloka
{इत्येवं परिपृष्टास्ते स्वं स्वं मतमतन्द्रिताः}
{सोपचारं तदा राममूचुर्हितचिकीर्षवः} %||6-11-25||

\twolineshloka
{अज्ञातं नास्ति ते किञ्चित्त्रिषु लोकेषु राघव}
{आत्मानं पूजयन्राम पृच्छस्यस्मान्सुहृत्तया} %||6-11-26||

\twolineshloka
{त्वं हि सत्यव्रतः शूरो धार्मिको दृढविक्रमः}
{परीक्ष्य कारा स्मृतिमान्निसृष्टात्मा सुहृत्सु च} %||6-11-27||

\twolineshloka
{तस्मादेकैकशस्तावद्ब्रुवन्तु सचिवास्तव}
{हेतुतो मतिसम्पन्नाः समर्थाश्च पुनः पुनः} %||6-11-28||

\twolineshloka
{इत्युक्ते राघवायाथ मतिमानङ्गदोऽग्रतः}
{विभीषणपरीक्षार्थमुवाच वचनं हरिः} %||6-11-29||

\twolineshloka
{शत्रोः सकाशात्सम्प्राप्तः सर्वथा शङ्क्य एव हि}
{विश्वासयोग्यः सहसा न कर्तव्यो विभीषणः} %||6-11-30||

\twolineshloka
{छादयित्वात्मभावं हि चरन्ति शठबुद्धयः}
{प्रहरन्ति च रन्ध्रेषु सोऽनर्थः सुमहान्भवेत्} %||6-11-31||

\twolineshloka
{अर्थानर्थौ विनिश्चित्य व्यवसायं भजेत ह}
{गुणतः सङ्ग्रहं कुर्याद्दोषतस्तु विसर्जयेत्} %||6-11-32||

\twolineshloka
{यदि दोषो महांस्तस्मिंस्त्यज्यतामविशङ्कितम्}
{गुणान्वापि बहूञ्ज्ञात्वा सङ्ग्रहः क्रियतां नृप} %||6-11-33||

\twolineshloka
{शरभस्त्वथ निश्चित्य सार्थं वचनमब्रवीत्}
{क्षिप्रमस्मिन्नरव्याघ्र चारः प्रतिविधीयताम्} %||6-11-34||

\twolineshloka
{प्रणिधाय हि चारेण यथावत्सूक्ष्मबुद्धिना}
{परीक्ष्य च ततः कार्यो यथान्यायं परिग्रहः} %||6-11-35||

\twolineshloka
{जाम्बवांस्त्वथ सम्प्रेक्ष्य शास्त्रबुद्ध्या विचक्षणः}
{वाक्यं विज्ञापयामास गुणवद्दोषवर्जितम्} %||6-11-36||

\twolineshloka
{बद्धवैराच्च पापाच्च राक्षसेन्द्राद्विभीषणः}
{अदेश काले सम्प्राप्तः सर्वथा शङ्क्यतामयम्} %||6-11-37||

\twolineshloka
{ततो मैन्दस्तु सम्प्रेक्ष्य नयापनयकोविदः}
{वाक्यं वचनसम्पन्नो बभाषे हेतुमत्तरम्} %||6-11-38||

\twolineshloka
{वचनं नाम तस्यैष रावणस्य विभीषणः}
{पृच्छ्यतां मधुरेणायं शनैर्नरवरेश्वर} %||6-11-39||

\twolineshloka
{भावमस्य तु विज्ञाय ततस्तत्त्वं करिष्यसि}
{यदि दृष्टो न दुष्टो वा बुद्धिपूर्वं नरर्षभ} %||6-11-40||

\twolineshloka
{अथ संस्कारसम्पन्नो हनूमान्सचिवोत्तमः}
{उवाच वचनं श्लक्ष्णमर्थवन्मधुरं लघु} %||6-11-41||

\twolineshloka
{न भवन्तं मतिश्रेष्ठं समर्थं वदतां वरम्}
{अतिशाययितुं शक्तो बृहस्पतिरपि ब्रुवन्} %||6-11-42||

\twolineshloka
{न वादान्नापि सङ्घर्षान्नाधिक्यान्न च कामतः}
{वक्ष्यामि वचनं राजन्यथार्थं रामगौरवात्} %||6-11-43||

\twolineshloka
{अर्थानर्थनिमित्तं हि यदुक्तं सचिवैस्तव}
{तत्र दोषं प्रपश्यामि क्रिया न ह्युपपद्यते} %||6-11-44||

\twolineshloka
{ऋते नियोगात्सामर्थ्यमवबोद्धुं न शक्यते}
{सहसा विनियोगो हि दोषवान्प्रतिभाति मे} %||6-11-45||

\twolineshloka
{चारप्रणिहितं युक्तं यदुक्तं सचिवैस्तव}
{अर्थस्यासम्भवात्तत्र कारणं नोपपद्यते} %||6-11-46||

\twolineshloka
{अदेश काले सम्प्राप्त इत्ययं यद्विभीषणः}
{विवक्षा चात्र मेऽस्तीयं तां निबोध यथा मति} %||6-11-47||

\twolineshloka
{स एष देशः कालश्च भवतीह यथा तथा}
{पुरुषात्पुरुषं प्राप्य तथा दोषगुणावपि} %||6-11-48||

\twolineshloka
{दौरात्म्यं रावणे दृष्ट्वा विक्रमं च तथा त्वयि}
{युक्तमागमनं तस्य सदृशं तस्य बुद्धितः} %||6-11-49||

\twolineshloka
{अज्ञातरूपैः पुरुषैः स राजन्पृच्छ्यतामिति}
{यदुक्तमत्र मे प्रेक्षा काचिदस्ति समीक्षिता} %||6-11-50||

\twolineshloka
{पृच्छ्यमानो विशङ्केत सहसा बुद्धिमान्वचः}
{तत्र मित्रं प्रदुष्येत मिथ्यपृष्टं सुखागतम्} %||6-11-51||

\twolineshloka
{अशक्यः सहसा राजन्भावो वेत्तुं परस्य वै}
{अन्तः स्वभावैर्गीतैस्तैर्नैपुण्यं पश्यता भृशम्} %||6-11-52||

\twolineshloka
{न त्वस्य ब्रुवतो जातु लक्ष्यते दुष्टभावता}
{प्रसन्नं वदनं चापि तस्मान्मे नास्ति संशयः} %||6-11-53||

\twolineshloka
{अशङ्कितमतिः स्वस्थो न शठः परिसर्पति}
{न चास्य दुष्टा वाक्चापि तस्मान्नास्तीह संशयः} %||6-11-54||

\twolineshloka
{आकारश्छाद्यमानोऽपि न शक्यो विनिगूहितुम्}
{बलाद्धि विवृणोत्येव भावमन्तर्गतं नृणाम्} %||6-11-55||

\twolineshloka
{देशकालोपपन्नं च कार्यं कार्यविदां वर}
{सफलं कुरुते क्षिप्रं प्रयोगेणाभिसंहितम्} %||6-11-56||

\twolineshloka
{उद्योगं तव सम्प्रेक्ष्य मिथ्यावृत्तं च रावणम्}
{वालिनश्च वधं श्रुत्वा सुग्रीवं चाभिषेचितम्} %||6-11-57||

\twolineshloka
{राज्यं प्रार्थयमानश्च बुद्धिपूर्वमिहागतः}
{एतावत्तु पुरस्कृत्य युज्यते त्वस्य सङ्ग्रहः} %||6-11-58||

\twolineshloka
{यथाशक्ति मयोक्तं तु राक्षसस्यार्जवं प्रति}
{त्वं प्रमाणं तु शेषस्य श्रुत्वा बुद्धिमतां वर} %||6-12-59||


{॥इत्यार्षे श्रीमद्रामायणे वाल्मीकीये आदिकाव्ये युद्धकाण्डे एकादशः सर्गः॥}

\sect{द्वादशः सर्गः}

\twolineshloka
{अथ रामः प्रसन्नात्मा श्रुत्वा वायुसुतस्य ह}
{प्रत्यभाषत दुर्धर्षः श्रुतवानात्मनि स्थितम्} %||6-12-1||

\twolineshloka
{ममापि तु विवक्षास्ति काचित्प्रति विभीषणम्}
{श्रुतमिच्छामि तत्सर्वं भवद्भिः श्रेयसि स्थितैः} %||6-12-2||

\twolineshloka
{मित्रभावेन सम्प्राप्तं न त्यजेयं कथञ्चन}
{दोषो यद्यपि तस्य स्यात्सतामेतदगर्हितम्} %||6-12-3||

\twolineshloka
{रामस्य वचनं श्रुत्वा सुग्रीवः प्लवगेश्वरः}
{प्रत्यभाषत काकुत्स्थं सौहार्देनाभिचोदितः} %||6-12-4||

\twolineshloka
{किमत्र चित्रं धर्मज्ञ लोकनाथशिखामणे}
{यत्त्वमार्यं प्रभाषेथाः सत्त्ववान्सपथे स्थितः} %||6-12-5||

\twolineshloka
{मम चाप्यन्तरात्मायं शुद्धिं वेत्ति विभीषणम्}
{अनुमनाच्च भावाच्च सर्वतः सुपरीक्षितः} %||6-12-6||

\twolineshloka
{तस्मात्क्षिप्रं सहास्माभिस्तुल्यो भवतु राघव}
{विभीषणो महाप्राज्ञः सखित्वं चाभ्युपैतु नः} %||6-12-7||

\twolineshloka
{स सुग्रीवस्य तद्वाक्यं रामः श्रुत्वा विमृश्य च}
{ततः शुभतरं वाक्यमुवाच हरिपुङ्गवम्} %||6-12-8||

\twolineshloka
{सुदुष्टो वाप्यदुष्टो वा किमेष रजनीचरः}
{सूक्ष्ममप्यहितं कर्तुं ममाशक्तः कथञ्चन} %||6-12-9||

\twolineshloka
{पिशाचान्दानवान्यक्षान्पृथिव्यां चैव राक्षसान्}
{अङ्गुल्यग्रेण तान्हन्यामिच्छन्हरिगणेश्वर} %||6-12-10||

\twolineshloka
{श्रूयते हि कपोतेन शत्रुः शरणमागतः}
{अर्चितश्च यथान्यायं स्वैश्च मांसैर्निमन्त्रितः} %||6-12-11||

\twolineshloka
{स हि तं प्रतिजग्राह भार्या हर्तारमागतम्}
{कपोतो वानरश्रेष्ठ किं पुनर्मद्विधो जनः} %||6-12-12||

\twolineshloka
{ऋषेः कण्वस्य पुत्रेण कण्डुना परमर्षिणा}
{शृणु गाथां पुरा गीतां धर्मिष्ठां सत्यवादिना} %||6-12-13||

\twolineshloka
{बद्धाञ्जलिपुटं दीनं याचन्तं शरणागतम्}
{न हन्यादानृशंस्यार्थमपि शत्रुं परं पत} %||6-12-14||

\twolineshloka
{आर्तो वा यदि वा दृप्तः परेषां शरणं गतः}
{अरिः प्राणान्परित्यज्य रक्षितव्यः कृतात्मना} %||6-12-15||

\twolineshloka
{स चेद्भयाद्वा मोहाद्वा कामाद्वापि न रक्षति}
{स्वया शक्त्या यथातत्त्वं तत्पापं लोकगर्हितम्} %||6-12-16||

\twolineshloka
{विनष्टः पश्यतस्तस्य रक्षिणः शरणागतः}
{आदाय सुकृतं तस्य सर्वं गच्छेदरक्षितः} %||6-12-17||

\twolineshloka
{एवं दोषो महानत्र प्रपन्नानामरक्षणे}
{अस्वर्ग्यं चायशस्यं च बलवीर्यविनाशनम्} %||6-12-18||

\twolineshloka
{करिष्यामि यथार्थं तु कण्डोर्वचनमुत्तमम्}
{धर्मिष्ठं च यशस्यं च स्वर्ग्यं स्यात्तु फलोदये} %||6-12-19||

\twolineshloka
{सकृदेव प्रपन्नाय तवास्मीति च याचते}
{अभयं सर्वभूतेभ्यो ददाम्येतद्व्रतं मम} %||6-12-20||

\twolineshloka
{आनयैनं हरिश्रेष्ठ दत्तमस्याभयं मया}
{विभीषणो वा सुग्रीव यदि वा रावणः स्वयम्} %||6-12-21||

\fourlineindentedshloka
{ततस्तु सुग्रीववचो निशम्य तद्-}
{धरीश्वरेणाभिहितं नरेश्वरः}
{विभीषणेनाशु जगाम सङ्गमम्}
{पतत्रिराजेन यथा पुरन्दरः} %||6-13-22||


{॥इत्यार्षे श्रीमद्रामायणे वाल्मीकीये आदिकाव्ये युद्धकाण्डे द्वादशः सर्गः॥}

\sect{त्रयोदशः सर्गः}

\twolineshloka
{राघवेणाभये दत्ते संनतो रावणानुजः}
{खात्पपातावनिं हृष्टो भक्तैरनुचरैः सह} %||6-13-1||

\twolineshloka
{स तु रामस्य धर्मात्मा निपपात विभीषणः}
{पादयोः शरणान्वेषी चतुर्भिः सह राक्षसैः} %||6-13-2||

\twolineshloka
{अब्रवीच्च तदा रामं वाक्यं तत्र विभीषणः}
{धर्मयुक्तं च युक्तं च साम्प्रतं सम्प्रहर्षणम्} %||6-13-3||

\twolineshloka
{अनुजो रावणस्याहं तेन चास्म्यवमानितः}
{भवन्तं सर्वभूतानां शरण्यं शरणं गतः} %||6-13-4||

\twolineshloka
{परित्यक्ता मया लङ्का मित्राणि च धनानि च}
{भवद्गतं मे राज्यं च जीवितं च सुखानि च} %||6-13-5||

\twolineshloka
{राक्षसानां वधे साह्यं लङ्कायाश्च प्रधर्षणे}
{करिष्यामि यथाप्राणं प्रवेक्ष्यामि च वाहिनीम्} %||6-13-6||

\twolineshloka
{इति ब्रुवाणं रामस्तु परिष्वज्य विभीषणम्}
{अब्रवील्लक्ष्मणं प्रीतः समुद्राज्जलमानय} %||6-13-7||

\twolineshloka
{तेन चेमं महाप्राज्ञमभिषिञ्च विभीषणम्}
{राजानं रक्षसां क्षिप्रं प्रसन्ने मयि मानद} %||6-13-8||

\twolineshloka
{एवमुक्तस्तु सौमित्रिरभ्यषिञ्चद्विभीषणम्}
{मध्ये वानरमुख्यानां राजानं रामशासनात्} %||6-13-9||

\twolineshloka
{तं प्रसादं तु रामस्य दृष्ट्वा सद्यः प्लवङ्गमाः}
{प्रचुक्रुशुर्महानादान्साधु साध्विति चाब्रुवन्} %||6-13-10||

\twolineshloka
{अब्रवीच्च हनूमांश्च सुग्रीवश्च विभीषणम्}
{कथं सागरमक्षोभ्यं तराम वरुणालयम्} %||6-13-11||

\twolineshloka
{उपायैरभिगच्छामो यथा नदनदीपतिम्}
{तराम तरसा सर्वे ससैन्या वरुणालयम्} %||6-13-12||

\twolineshloka
{एवमुक्तस्तु धर्मज्ञः प्रत्युवाच विभीषणः}
{समुद्रं राघवो राजा शरणं गन्तुमर्हति} %||6-13-13||

\twolineshloka
{खानितः सगरेणायमप्रमेयो महोदधिः}
{कर्तुमर्हति रामस्य ज्ञातेः कार्यं महोदधिः} %||6-13-14||

\twolineshloka
{एवं विभीषणेनोक्ते राक्षसेन विपश्चिता}
{प्रकृत्या धर्मशीलस्य राघवस्याप्यरोचत} %||6-13-15||

\twolineshloka
{स लक्ष्मणं महातेजाः सुग्रीवं च हरीश्वरम्}
{सत्क्रियार्थं क्रियादक्षः स्मितपूर्वमुवाच ह} %||6-13-16||

\twolineshloka
{विभीषणस्य मन्त्रोऽयं मम लक्ष्मण रोचते}
{ब्रूहि त्वं सहसुग्रीवस्तवापि यदि रोचते} %||6-13-17||

\twolineshloka
{सुग्रीवः पण्डितो नित्यं भवान्मन्त्रविचक्षणः}
{उभाभ्यां सम्प्रधार्यार्यं रोचते यत्तदुच्यताम्} %||6-13-18||

\twolineshloka
{एवमुक्तौ तु तौ वीरावुभौ सुग्रीवलक्ष्मणौ}
{समुदाचार संयुक्तमिदं वचनमूचतुः} %||6-13-19||

\twolineshloka
{किमर्थं नो नरव्याघ्र न रोचिष्यति राघव}
{विभीषणेन यत्तूक्तमस्मिन्काले सुखावहम्} %||6-13-20||

\twolineshloka
{अबद्ध्वा सागरे सेतुं घोरेऽस्मिन्वरुणालये}
{लङ्का नासादितुं शक्या सेन्द्रैरपि सुरासुरैः} %||6-13-21||

\twolineshloka
{विभीषणस्य शूरस्य यथार्थं क्रियतां वचः}
{अलं कालात्ययं कृत्वा समुद्रोऽयं नियुज्यताम्} %||6-13-22||

\twolineshloka
{एवमुक्तः कुशास्तीर्णे तीरे नदनदीपतेः}
{संविवेश तदा रामो वेद्यामिव हुताशनः} %||6-14-23||


{॥इत्यार्षे श्रीमद्रामायणे वाल्मीकीये आदिकाव्ये युद्धकाण्डे त्रयोदशः सर्गः॥}

\sect{चतुर्दशः सर्गः}

\twolineshloka
{तस्य रामस्य सुप्तस्य कुशास्तीर्णे महीतले}
{नियमादप्रमत्तस्य निशास्तिस्रोऽतिचक्रमुः} %||6-14-1||

\twolineshloka
{न च दर्शयते मन्दस्तदा रामस्य सागरः}
{प्रयतेनापि रामेण यथार्हमभिपूजितः} %||6-14-2||

\twolineshloka
{समुद्रस्य ततः क्रुद्धो रामो रक्तान्तलोचनः}
{समीपस्थमुवाचेदं लक्ष्मणं शुभलक्ष्मणम्} %||6-14-3||

\twolineshloka
{पश्य तावदनार्यस्य पूज्यमानस्य लक्ष्मण}
{अवलेपं समुद्रस्य न दर्शयति यत्स्वयम्} %||6-14-4||

\twolineshloka
{प्रशमश्च क्षमा चैव आर्जवं प्रियवादिता}
{असामर्थ्यं फलन्त्येते निर्गुणेषु सतां गुणाः} %||6-14-5||

\twolineshloka
{आत्मप्रशंसिनं दुष्टं धृष्टं विपरिधावकम्}
{सर्वत्रोत्सृष्टदण्डं च लोकः सत्कुरुते नरम्} %||6-14-6||

\twolineshloka
{न साम्ना शक्यते कीर्तिर्न साम्ना शक्यते यशः}
{प्राप्तुं लक्ष्मण लोकेऽस्मिञ्जयो वा रणमूधनि} %||6-14-7||

\twolineshloka
{अद्य मद्बाणनिर्भिन्नैर्मकरैर्मकरालयम्}
{निरुद्धतोयं सौमित्रे प्लवद्भिः पश्य सर्वतः} %||6-14-8||

\twolineshloka
{महाभोगानि मत्स्यानां करिणां च करानिह}
{भोगांश्च पश्य नागानां मया भिन्नानि लक्ष्मण} %||6-14-9||

\twolineshloka
{सशङ्खशुक्तिका जालं समीनमकरं शरैः}
{अद्य युद्धेन महता समुद्रं परिशोषये} %||6-14-10||

\twolineshloka
{क्षमया हि समायुक्तं मामयं मकरालयः}
{असमर्थं विजानाति धिक्क्षमामीदृशे जने} %||6-14-11||

\twolineshloka
{चापमानय सौमित्रे शरांश्चाशीविषोपमान्}
{अद्याक्षोभ्यमपि क्रुद्धः क्षोभयिष्यामि सागरम्} %||6-14-12||

\twolineshloka
{वेलासु कृतमर्यादं सहसोर्मिसमाकुलम्}
{निर्मर्यादं करिष्यामि सायकैर्वरुणालयम्} %||6-14-13||

\twolineshloka
{एवमुक्त्वा धनुष्पाणिः क्रोधविस्फारितेक्षणः}
{बभूव रामो दुर्धर्षो युगान्ताग्निरिव ज्वलन्} %||6-14-14||

\twolineshloka
{सम्पीड्य च धनुर्घोरं कम्पयित्वा शरैर्जगत्}
{मुमोच विशिखानुग्रान्वज्राणीव शतक्रतुः} %||6-14-15||

\twolineshloka
{ते ज्वलन्तो महावेगास्तेजसा सायकोत्तमाः}
{प्रविशन्ति समुद्रस्य सलिलं त्रस्तपन्नगम्} %||6-14-16||

\twolineshloka
{ततो वेगः समुद्रस्य सनक्रमकरो महान्}
{सम्बभूव महाघोरः समारुतरवस्तदा} %||6-14-17||

\twolineshloka
{महोर्मिमालाविततः शङ्खशुक्तिसमाकुलः}
{सधूमपरिवृत्तोर्मिः सहसाभून्महोदधिः} %||6-14-18||

\twolineshloka
{व्यथिताः पन्नगाश्चासन्दीप्तास्या दीप्तलोचनाः}
{दानवाश्च महावीर्याः पातालतलवासिनः} %||6-14-19||

\twolineshloka
{ऊर्मयः सिन्धुराजस्य सनक्रमकरास्तदा}
{विन्ध्यमन्दरसङ्काशाः समुत्पेतुः सहस्रशः} %||6-14-20||

\twolineshloka
{आघूर्णिततरङ्गौघः सम्भ्रान्तोरगराक्षसः}
{उद्वर्तित महाग्राहः संवृत्तः सलिलाशयः} %||6-15-21||


{॥इत्यार्षे श्रीमद्रामायणे वाल्मीकीये आदिकाव्ये युद्धकाण्डे चतुर्दशः सर्गः॥}

\sect{पञ्चदशः सर्गः}

\threelineshloka
{ततो मध्यात्समुद्रस्य सागरः स्वयमुत्थितः}
{उदयन्हि महाशैलान्मेरोरिव दिवाकरः}
{पन्नगैः सह दीप्तास्यैः समुद्रः प्रत्यदृश्यत} %||6-15-1||

\twolineshloka
{स्निग्धवैदूर्यसङ्काशो जाम्बूनदविभूषितः}
{रक्तमाल्याम्बरधरः पद्मपत्रनिभेक्षणः} %||6-15-2||

\twolineshloka
{सागरः समतिक्रम्य पूर्वमामन्त्र्य वीर्यवान्}
{अब्रवीत्प्राञ्जलिर्वाक्यं राघवं शरपाणिनम्} %||6-15-3||

\twolineshloka
{पृथिवी वायुराकाशमापो ज्योतिश्च राघवः}
{स्वभावे सौम्य तिष्ठन्ति शाश्वतं मार्गमाश्रिताः} %||6-15-4||

\twolineshloka
{तत्स्वभावो ममाप्येष यदगाधोऽहमप्लवः}
{विकारस्तु भवेद्राध एतत्ते प्रवदाम्यहम्} %||6-15-5||

\twolineshloka
{न कामान्न च लोभाद्वा न भयात्पार्थिवात्मज}
{ग्राहनक्राकुलजलं स्तम्भयेयं कथञ्चन} %||6-15-6||

\twolineshloka
{विधास्ये राम येनापि विषहिष्ये ह्यहं तथा}
{ग्राहा न प्रहरिष्यन्ति यावत्सेना तरिष्यति} %||6-15-7||

\twolineshloka
{अयं सौम्य नलो नाम तनुजो विश्वकर्मणः}
{पित्रा दत्तवरः श्रीमान्प्रतिमो विश्वकर्मणः} %||6-15-8||

\twolineshloka
{एष सेतुं महोत्साहः करोतु मयि वानरः}
{तमहं धारयिष्यामि तथा ह्येष यथा पिता} %||6-15-9||

\twolineshloka
{एवमुक्त्वोदधिर्नष्टः समुत्थाय नलस्ततः}
{अब्रवीद्वानरश्रेष्ठो वाक्यं रामं महाबलः} %||6-15-10||

\twolineshloka
{अहं सेतुं करिष्यामि विस्तीर्णे वरुणालये}
{पितुः सामर्थ्यमास्थाय तत्त्वमाह महोदधिः} %||6-15-11||

\twolineshloka
{मम मातुर्वरो दत्तो मन्दरे विश्वकर्मणा}
{औरसस्तस्य पुत्रोऽहं सदृशो विश्वकर्मणा} %||6-15-12||

\twolineshloka
{न चाप्यहमनुक्तो वै प्रब्रूयामात्मनो गुणान्}
{काममद्यैव बध्नन्तु सेतुं वानरपुङ्गवाः} %||6-15-13||

\twolineshloka
{ततो निसृष्टरामेण सर्वतो हरियूथपाः}
{अभिपेतुर्महारण्यं हृष्टाः शतसहस्रशः} %||6-15-14||

\twolineshloka
{ते नगान्नगसङ्काशाः शाखामृगगणर्षभाः}
{बभञ्जुर्वानरास्तत्र प्रचकर्षुश्च सागरम्} %||6-15-15||

\twolineshloka
{ते सालैश्चाश्वकर्णैश्च धवैर्वंशैश्च वानराः}
{कुटजैरर्जुनैस्तालैस्तिकलैस्तिमिशैरपि} %||6-15-16||

\twolineshloka
{बिल्वकैः सप्तपर्णैश्च कर्णिकारैश्च पुष्पितैः}
{चूतैश्चाशोकवृक्षैश्च सागरं समपूरयन्} %||6-15-17||

\twolineshloka
{समूलांश्च विमूलांश्च पादपान्हरिसत्तमाः}
{इन्द्रकेतूनिवोद्यम्य प्रजह्रुर्हरयस्तरून्} %||6-15-18||

\twolineshloka
{प्रक्षिप्यमाणैरचलैः सहसा जलमुद्धतम्}
{समुत्पतितमाकाशमपासर्पत्ततस्ततः} %||6-15-19||

\twolineshloka
{दशयोजनविस्तीर्णं शतयोजनमायतम्}
{नलश्चक्रे महासेतुं मध्ये नदनदीपतेः} %||6-15-20||

\twolineshloka
{शिलानां क्षिप्यमाणानां शैलानां तत्र पात्यताम्}
{बभूव तुमुलः शब्दस्तदा तस्मिन्महोदधौ} %||6-15-21||

\twolineshloka
{स नलेन कृतः सेतुः सागरे मकरालये}
{शुशुभे सुभगः श्रीमान्स्वातीपथ इवाम्बरे} %||6-15-22||

{ततो देवाः सगन्धर्वाः सिद्धाश्च परमर्षयः॥\devanumber{23}॥} %||6-15-23|| (Check)

\addtocounter{shlokacount}{1}

\threelineshloka
{आप्लवन्तः प्लवन्तश्च गर्जन्तश्च प्लवङ्गमाः}
{तमचिन्त्यमसह्यं च अद्भुतं लोमहर्षणम्}
{ददृशुः सर्वभूतानि सागरे सेतुबन्धनम्} %||6-15-24||

\twolineshloka
{तानि कोटिसहस्राणि वानराणां महौजसाम्}
{बध्नन्तः सागरे सेतुं जग्मुः पारं महोदधेः} %||6-15-25||

\twolineshloka
{विशालः सुकृतः श्रीमान्सुभूमिः सुसमाहितः}
{अशोभत महासेतुः सीमन्त इव सागरे} %||6-15-26||

\twolineshloka
{ततः परे समुद्रस्य गदापाणिर्विभीषणः}
{परेषामभिघतार्थमतिष्ठत्सचिवैः सह} %||6-15-27||

\twolineshloka
{अग्रतस्तस्य सैन्यस्य श्रीमान्रामः सलक्ष्मणः}
{जगाम धन्वी धर्मात्मा सुग्रीवेण समन्वितः} %||6-15-28||

\threelineshloka
{अन्ये मध्येन गच्छन्ति पार्श्वतोऽन्ये प्लवङ्गमाः}
{सलिले प्रपतन्त्यन्ये मार्गमन्ये न लेभिरे}
{केचिद्वैहायस गताः सुपर्णा इव पुप्लुवुः} %||6-15-29||

\twolineshloka
{घोषेण महता घोषं सागरस्य समुच्छ्रितम्}
{भीममन्तर्दधे भीमा तरन्ती हरिवाहिनी} %||6-15-30||

\twolineshloka
{वानराणां हि सा तीर्णा वाहिनी नल सेतुना}
{तीरे निविविशे राज्ञा बहुमूलफलोदके} %||6-15-31||

\fourlineindentedshloka
{तदद्भुतं राघव कर्म दुष्करम्}
{समीक्ष्य देवाः सह सिद्धचारणैः}
{उपेत्य रामं सहिता महर्षिभिः}
{समभ्यषिञ्चन्सुशुभैर्जलैः पृथक्} %||6-15-32||

\fourlineindentedshloka
{जयस्व शत्रून्नरदेव मेदिनीम्}
{ससागरां पालय शाश्वतीः समाः}
{इतीव रामं नरदेवसत्कृतम्}
{शुभैर्वचोभिर्विविधैरपूजयन्} %||6-16-33||


{॥इत्यार्षे श्रीमद्रामायणे वाल्मीकीये आदिकाव्ये युद्धकाण्डे पञ्चदशः सर्गः॥}

\sect{षोडशः सर्गः}

\twolineshloka
{सबले सागरं तीर्णे रामे दशरथात्मजे}
{अमात्यौ रावणः श्रीमानब्रवीच्छुकसारणौ} %||6-16-1||

\twolineshloka
{समग्रं सागरं तीर्णं दुस्तरं वानरं बलम्}
{अभूतपूर्वं रामेण सागरे सेतुबन्धनम्} %||6-16-2||

\twolineshloka
{सागरे सेतुबन्धं तु न श्रद्दध्यां कथञ्चन}
{अवश्यं चापि सङ्ख्येयं तन्मया वानरं बलम्} %||6-16-3||

\twolineshloka
{भवन्तौ वानरं सैन्यं प्रविश्यानुपलक्षितौ}
{परिमाणं च वीर्यं च ये च मुख्याः प्लवङ्गमाः} %||6-16-4||

\twolineshloka
{मन्त्रिणो ये च रामस्य सुग्रीवस्य च सम्मताः}
{ये पूर्वमभिवर्तन्ते ये च शूराः प्लवङ्गमाः} %||6-16-5||

\twolineshloka
{स च सेतुर्यथा बद्धः सागरे सलिलार्णवे}
{निवेशश्च यथा तेषां वानराणां महात्मनाम्} %||6-16-6||

\twolineshloka
{रामस्य व्यवसायं च वीर्यं प्रहरणानि च}
{लक्ष्मणस्य च वीरस्य तत्त्वतो ज्ञातुमर्हथ} %||6-16-7||

\twolineshloka
{कश्च सेनापतिस्तेषां वानराणां महौजसाम्}
{एतज्ज्ञात्वा यथातत्त्वं शीघ्रमगन्तुमर्हथः} %||6-16-8||

\twolineshloka
{इति प्रतिसमादिष्टौ राक्षसौ शुकसारणौ}
{हरिरूपधरौ वीरौ प्रविष्टौ वानरं बलम्} %||6-16-9||

\twolineshloka
{ततस्तद्वानरं सैन्यमचिन्त्यं लोमहर्षणम्}
{सङ्ख्यातुं नाध्यगच्छेतां तदा तौ शुकसारणौ} %||6-16-10||

\twolineshloka
{तत्स्थितं पर्वताग्रेषु निर्दरेषु गुहासु च}
{समुद्रस्य च तीरेषु वनेषूपवनेषु च} %||6-16-11||

\twolineshloka
{तरमाणं च तीर्णं च तर्तुकामं च सर्वशः}
{निविष्टं निविशच्चैव भीमनादं महाबलम्} %||6-16-12||

\threelineshloka
{तौ ददर्श महातेजाः प्रच्छन्नौ च विभीषणः}
{आचचक्षेऽथ रामाय गृहीत्वा शुकसारणौ}
{लङ्कायाः समनुप्राप्तौ चारौ परपुरंजयौ} %||6-16-13||

\twolineshloka
{तौ दृष्ट्वा व्यथितौ रामं निराशौ जीविते तदा}
{कृताञ्जलिपुटौ भीतौ वचनं चेदमूचतुः} %||6-16-14||

\twolineshloka
{आवामिहागतौ सौम्य रावणप्रहितावुभौ}
{परिज्ञातुं बलं कृत्स्नं तवेदं रघुनन्दन} %||6-16-15||

\twolineshloka
{तयोस्तद्वचनं श्रुत्वा रामो दशरथात्मजः}
{अब्रवीत्प्रहसन्वाक्यं सर्वभूतहिते रतः} %||6-16-16||

\twolineshloka
{यदि दृष्टं बलं कृत्स्नं वयं वा सुसमीक्षिताः}
{यथोक्तं वा कृतं कार्यं छन्दतः प्रतिगम्यताम्} %||6-16-17||

\twolineshloka
{प्रविश्य नगरीं लङ्कां भवद्भ्यां धनदानुजः}
{वक्तव्यो रक्षसां राजा यथोक्तं वचनं मम} %||6-16-18||

\twolineshloka
{यद्बलं च समाश्रित्य सीतां मे हृतवानसि}
{तद्दर्शय यथाकामं ससैन्यः सहबान्धवः} %||6-16-19||

\twolineshloka
{श्वःकाले नगरीं लङ्कां सप्राकारां सतोरणाम्}
{राक्षसं च बलं पश्य शरैर्विध्वंसितं मया} %||6-16-20||

\twolineshloka
{घोरं रोषमहं मोक्ष्ये बलं धारय रावण}
{श्वःकाले वज्रवान्वज्रं दानवेष्विव वासवः} %||6-16-21||

\twolineshloka
{इति प्रतिसमादिष्टौ राक्षसौ शुकसारणौ}
{आगम्य नगरीं लङ्कामब्रूतां राक्षसाधिपम्} %||6-16-22||

\twolineshloka
{विभीषणगृहीतौ तु वधार्हौ राक्षसेश्वर}
{दृष्ट्वा धर्मात्मना मुक्तौ रामेणामिततेजसा} %||6-16-23||

\twolineshloka
{एकस्थानगता यत्र चत्वारः पुरुषर्षभाः}
{लोकपालोपमाः शूराः कृतास्त्रा दृढविक्रमाः} %||6-16-24||

\twolineshloka
{रामो दाशरथिः श्रीमाँल्लक्ष्मणश्च विभीषणः}
{सुग्रीवश्च महातेजा महेन्द्रसमविक्रमः} %||6-16-25||

\twolineshloka
{एते शक्ताः पुरीं लङ्कां सप्राकारां सतोरणाम्}
{उत्पाट्य सङ्क्रामयितुं सर्वे तिष्ठन्तु वानराः} %||6-16-26||

\twolineshloka
{यादृशं तस्य रामस्य रूपं प्रहरणानि च}
{वधिष्यति पुरीं लङ्कामेकस्तिष्ठन्तु ते त्रयः} %||6-16-27||

\twolineshloka
{रामलक्ष्मणगुप्ता सा सुग्रीवेण च वाहिनी}
{बभूव दुर्धर्षतरा सर्वैरपि सुरासुरैः} %||6-16-28||

\fourlineindentedshloka
{प्रहृष्टरूपा ध्वजिनी वनौकसाम्}
{महात्मनां सम्प्रति योद्धुमिच्छताम्}
{अलं विरोधेन शमो विधीयताम्}
{प्रदीयतां दाशरथाय मैथिली} %||6-17-29||


{॥इत्यार्षे श्रीमद्रामायणे वाल्मीकीये आदिकाव्ये युद्धकाण्डे षोडशः सर्गः॥}

\sect{सप्तदशः सर्गः}

\twolineshloka
{तद्वचः पथ्यमक्लीबं सारणेनाभिभाषितम्}
{निशम्य रावणो राजा प्रत्यभाषत सारणम्} %||6-17-1||

\twolineshloka
{यदि मामभियुञ्जीरन्देवगन्धर्वदानवाः}
{नैव सीतां प्रदास्यामि सर्वलोकभयादपि} %||6-17-2||

\threelineshloka
{त्वं तु सौम्य परित्रस्तो हरिभिर्निर्जितो भृशम्}
{प्रतिप्रदानमद्यैव सीतायाः साधु मन्यसे}
{को हि नाम सपत्नो मां समरे जेतुमर्हति} %||6-17-3||

\threelineshloka
{इत्युक्त्वा परुषं वाक्यं रावणो राक्षसाधिपः}
{आरुरोह ततः श्रीमान्प्रासादं हिमपाण्डुरम्}
{बहुतालसमुत्सेधं रावणोऽथ दिदृक्षया} %||6-17-4||

\threelineshloka
{ताभ्यां चराभ्यां सहितो रावणः क्रोधमूर्छितः}
{पश्यमानः समुद्रं च पर्वतांश्च वनानि च}
{ददर्श पृथिवीदेशं सुसम्पूर्णं प्लवङ्गमैः} %||6-17-5||

\twolineshloka
{तदपारमसङ्ख्येयं वानराणां महद्बलम्}
{आलोक्य रावणो राजा परिपप्रच्छ सारणम्} %||6-17-6||

\twolineshloka
{एषां वानरमुख्यानां के शूराः के महाबलाः}
{के पूर्वमभिवर्तन्ते महोत्साहाः समन्ततः} %||6-17-7||

\twolineshloka
{केषां शृणोति सुग्रीवः के वा यूथपयूथपाः}
{सारणाचक्ष्व मे सर्वं के प्रधानाः प्लवङ्गमाः} %||6-17-8||

\twolineshloka
{सारणो राक्षसेन्द्रस्य वचनं परिपृच्छतः}
{आचचक्षेऽथ मुख्यज्ञो मुख्यांस्तांस्तु वनौकसः} %||6-17-9||

\twolineshloka
{एष योऽभिमुखो लङ्कां नर्दंस्तिष्ठति वानरः}
{यूथपानां सहस्राणां शतेन परिवारितः} %||6-17-10||

\twolineshloka
{यस्य घोषेण महता सप्राकारा सतोरणा}
{लङ्का प्रवेपते सर्वा सशैलवनकानना} %||6-17-11||

\twolineshloka
{सर्वशाखामृगेन्द्रस्य सुग्रीवस्य महात्मनः}
{बलाग्रे तिष्ठते वीरो नीलो नामैष यूथपः} %||6-17-12||

\twolineshloka
{बाहू प्रगृह्य यः पद्भ्यां महीं गच्छति वीर्यवान्}
{लङ्कामभिमुखः कोपादभीक्ष्णं च विजृम्भते} %||6-17-13||

\twolineshloka
{गिरिशृङ्गप्रतीकाशः पद्मकिञ्जल्कसंनिभः}
{स्फोटयत्यभिसंरब्धो लाङ्गूलं च पुनः पुनः} %||6-17-14||

\threelineshloka
{यस्य लाङ्गूलशब्देन स्वनन्तीव दिशो दश}
{एष वानरराजेन सुर्ग्रीवेणाभिषेचितः}
{यौवराज्येऽङ्गदो नाम त्वामाह्वयति संयुगे} %||6-17-15||

\twolineshloka
{ये तु विष्टभ्य गात्राणि क्ष्वेडयन्ति नदन्ति च}
{उत्थाय च विजृम्भन्ते क्रोधेन हरिपुङ्गवाः} %||6-17-16||

\twolineshloka
{एते दुष्प्रसहा घोराश्चण्डाश्चण्डपराक्रमाः}
{अष्टौ शतसहस्राणि दशकोटिशतानि च} %||6-17-17||

\twolineshloka
{य एनमनुगच्छन्ति वीराश्चन्दनवासिनः}
{एष आशंसते लङ्कां स्वेनानीकेन मर्दितुम्} %||6-17-18||

\twolineshloka
{श्वेतो रजतसङ्काशः सबलो भीमविक्रमः}
{बुद्धिमान्वानरः शूरस्त्रिषु लोकेषु विश्रुतः} %||6-17-19||

\twolineshloka
{तूर्णं सुग्रीवमागम्य पुनर्गच्छति वानरः}
{विभजन्वानरीं सेनामनीकानि प्रहर्षयन्} %||6-17-20||

\twolineshloka
{यः पुरा गोमतीतीरे रम्यं पर्येति पर्वतम्}
{नाम्ना सङ्कोचनो नाम नानानगयुतो गिरिः} %||6-17-21||

\twolineshloka
{तत्र राज्यं प्रशास्त्येष कुमुदो नाम यूथपः}
{योऽसौ शतसहस्राणां सहस्रं परिकर्षति} %||6-17-22||

\twolineshloka
{यस्य वाला बहुव्यामा दीर्घलाङ्गूलमाश्रिताः}
{ताम्राः पीताः सिताः श्वेताः प्रकीर्णा घोरकर्मणः} %||6-17-23||

\twolineshloka
{अदीनो रोषणश्चण्डः सङ्ग्राममभिकाङ्क्षति}
{एषैवाशंसते लङ्कां स्वेनानीकेन मर्दितुम्} %||6-17-24||

\twolineshloka
{यस्त्वेष सिंहसङ्काशः कपिलो दीर्घकेसरः}
{निभृतः प्रेक्षते लङ्कां दिधक्षन्निव चक्षुषा} %||6-17-25||

\twolineshloka
{विन्ध्यं कृष्णगिरिं सह्यं पर्वतं च सुदर्शनम्}
{राजन्सततमध्यास्ते रम्भो नामैष यूथपः} %||6-17-26||

\twolineshloka
{शतं शतसहस्राणां त्रिंशच्च हरियूथपाः}
{परिवार्यानुगच्छन्ति लङ्कां मर्दितुमोजसा} %||6-17-27||

\twolineshloka
{यस्तु कर्णौ विवृणुते जृम्भते च पुनः पुनः}
{न च संविजते मृत्योर्न च यूथाद्विधावति} %||6-17-28||

\twolineshloka
{महाबलो वीतभयो रम्यं साल्वेय पर्वतम्}
{राजन्सततमध्यास्ते शरभो नाम यूथपः} %||6-17-29||

\twolineshloka
{एतस्य बलिनः सर्वे विहारा नाम यूथपाः}
{राजञ्शतसहस्राणि चत्वारिंशत्तथैव च} %||6-17-30||

\twolineshloka
{यस्तु मेघ इवाकाशं महानावृत्य तिष्ठति}
{मध्ये वानरवीराणां सुराणामिव वासवः} %||6-17-31||

\twolineshloka
{भेरीणामिव संनादो यस्यैष श्रूयते महान्}
{घोरः शाखामृगेन्द्राणां सङ्ग्राममभिकाङ्क्षताम्} %||6-17-32||

\twolineshloka
{एष पर्वतमध्यास्ते पारियात्रमनुत्तमम्}
{युद्धे दुष्प्रसहो नित्यं पनसो नाम यूथपः} %||6-17-33||

\twolineshloka
{एनं शतसहस्राणां शतार्धं पर्युपासते}
{यूथपा यूथपश्रेष्ठं येषां यूथानि भागशः} %||6-17-34||

\twolineshloka
{यस्तु भीमां प्रवल्गन्तीं चमूं तिष्ठति शोभयन्}
{स्थितां तीरे समुद्रस्य द्वितीय इव सागरः} %||6-17-35||

\twolineshloka
{एष दर्दरसङ्काशो विनतो नाम यूथपः}
{पिबंश्चरति पर्णाशां नदीनामुत्तमां नदीम्} %||6-17-36||

\twolineshloka
{षष्टिः शतसहस्राणि बलमस्य प्लवङ्गमाः}
{त्वामाह्वयति युद्धाय क्रथनो नाम यूथपः} %||6-17-37||

\twolineshloka
{यस्तु गैरिकवर्णाभं वपुः पुष्यति वानरः}
{गवयो नाम तेजस्वी त्वां क्रोधादभिवर्तते} %||6-17-38||

\twolineshloka
{एनं शतसहस्राणि सप्ततिः पर्युपासते}
{एष आशंसते लङ्कां स्वेनानीकेन मर्दितुम्} %||6-17-39||

\twolineshloka
{एते दुष्प्रसहा घोरा बलिनः कामरूपिणः}
{यूथपा यूथपश्रेष्ठा येषां सङ्ख्या न विद्यते} %||6-18-40||


{॥इत्यार्षे श्रीमद्रामायणे वाल्मीकीये आदिकाव्ये युद्धकाण्डे सप्तदशः सर्गः॥}

\sect{अष्टादशः सर्गः}

\twolineshloka
{तांस्तु तेऽहं प्रवक्ष्यामि प्रेक्षमाणस्य यूथपान्}
{राघवार्थे पराक्रान्ता ये न रक्षन्ति जीवितम्} %||6-18-1||

\twolineshloka
{स्निग्धा यस्य बहुश्यामा बाला लाङ्गूलमाश्रिताः}
{ताम्राः पीताः सिताः श्वेताः प्रकीर्णा घोरकर्मणः} %||6-18-2||

\twolineshloka
{प्रगृहीताः प्रकाशन्ते सूर्यस्येव मरीचयः}
{पृथिव्यां चानुकृष्यन्ते हरो नामैष यूथपः} %||6-18-3||

\twolineshloka
{यं पृष्ठतोऽनुगच्छन्ति शतशोऽथ सहस्रशः}
{द्रुमानुद्यम्य सहिता लङ्कारोहणतत्पराः} %||6-18-4||

\twolineshloka
{एष कोटीसहस्रेण वानराणां महौजसाम्}
{आकाङ्क्षते त्वां सङ्ग्रामे जेतुं परपुरंजय} %||6-18-5||

\twolineshloka
{नीलानिव महामेघांस्तिष्ठतो यांस्तु पश्यसि}
{असिताञ्जनसङ्काशान्युद्धे सत्यपराक्रमान्} %||6-18-6||

\twolineshloka
{नखदंष्ट्रायुधान्वीरांस्तीक्ष्णकोपान्भयावहान्}
{असङ्ख्येयाननिर्देश्यान्परं पारमिवोदधेः} %||6-18-7||

\twolineshloka
{पर्वतेषु च ये केचिद्विषमेषु नदीषु च}
{एते त्वामभिवर्तन्ते राजन्नृष्काः सुदारुणाः} %||6-18-8||

\twolineshloka
{एषां मध्ये स्थितो राजन्भीमाक्षो भीमदर्शनः}
{पर्जन्य इव जीमूतैः समन्तात्परिवारितः} %||6-18-9||

\twolineshloka
{ऋक्षवन्तं गिरिश्रेष्ठमध्यास्ते नर्मदां पिबन्}
{सर्वर्क्षाणामधिपतिर्धूम्रो नामैष यूथपः} %||6-18-10||

\twolineshloka
{यवीयानस्य तु भ्राता पश्यैनं पर्वतोपमम्}
{भ्रात्रा समानो रूपेण विशिष्टस्तु पराक्रमे} %||6-18-11||

\twolineshloka
{स एष जाम्बवान्नाम महायूथपयूथपः}
{प्रशान्तो गुरुवर्ती च सम्प्रहारेष्वमर्षणः} %||6-18-12||

\twolineshloka
{एतेन साह्यं सुमहत्कृतं शक्रस्य धीमता}
{देवासुरे जाम्बवता लब्धाश्च बहवो वराः} %||6-18-13||

\twolineshloka
{आरुह्य पर्वताग्रेभ्यो महाभ्रविपुलाः शिलाः}
{मुञ्चन्ति विपुलाकारा न मृत्योरुद्विजन्ति च} %||6-18-14||

\twolineshloka
{राक्षसानां च सदृशाः पिशाचानां च रोमशाः}
{एतस्य सैन्ये बहवो विचरन्त्यग्नितेजसः} %||6-18-15||

\twolineshloka
{यं त्वेनमभिसंरब्धं प्लवमानमिव स्थितम्}
{प्रेक्षन्ते वानराः सर्वे स्थितं यूथपयूथपम्} %||6-18-16||

\twolineshloka
{एष राजन्सहस्राक्षं पर्युपास्ते हरीश्वरः}
{बलेन बलसम्पन्नो रम्भो नामैष यूथपः} %||6-18-17||

\twolineshloka
{यः स्थितं योजने शैलं गच्छन्पार्श्वेन सेवते}
{ऊर्ध्वं तथैव कायेन गतः प्राप्नोति योजनम्} %||6-18-18||

\twolineshloka
{यस्मान्न परमं रूपं चतुष्पादेषु विद्यते}
{श्रुतः संनादनो नाम वानराणां पितामहः} %||6-18-19||

\threelineshloka
{येन युद्धं तदा दत्तं रणे शक्रस्य धीमता}
{पराजयश्च न प्राप्तः सोऽयं यूथपयूथपः}
{यस्य विक्रममाणस्य शक्रस्येव पराक्रमः} %||6-18-20||

\twolineshloka
{एष गन्धर्वकन्यायामुत्पन्नः कृष्णवर्त्मना}
{पुरा देवासुरे युद्धे साह्यार्थं त्रिदिवौकसाम्} %||6-18-21||

\twolineshloka
{यस्य वैश्रवणो राजा जम्बूमुपनिषेवते}
{यो राजा पर्वतेन्द्राणां बहुकिंनरसेविनाम्} %||6-18-22||

\threelineshloka
{विहारसुखदो नित्यं भ्रातुस्ते राक्षसाधिप}
{तत्रैष वसति श्रीमान्बलवान्वानरर्षभः}
{युद्धेष्वकत्थनो नित्यं क्रथनो नाम यूथपः} %||6-18-23||

\twolineshloka
{वृतः कोटिसहस्रेण हरीणां समुपस्थितः}
{एषैवाशंसते लङ्कां स्वेनानीकेन मर्दितुम्} %||6-18-24||

\twolineshloka
{यो गङ्गामनु पर्येति त्रासयन्हस्तियूथपान्}
{हस्तिनां वानराणां च पूर्ववैरमनुस्मरन्} %||6-18-25||

\twolineshloka
{एष यूथपतिर्नेता गच्छन्गिरिगुहाशयः}
{हरीणां वाहिनी मुख्यो नदीं हैमवतीमनु} %||6-18-26||

\twolineshloka
{उशीर बीजमाश्रित्य पर्वतं मन्दरोपमम्}
{रमते वानरश्रेष्ठो दिवि शक्र इव स्वयम्} %||6-18-27||

\twolineshloka
{एनं शतसहस्राणां सहस्रमभिवर्तते}
{एष दुर्मर्षणो राजन्प्रमाथी नाम यूथपः} %||6-18-28||

\twolineshloka
{वातेनेवोद्धतं मेघं यमेनमनुपश्यसि}
{विवर्तमानं बहुशो यत्रैतद्बहुलं रजः} %||6-18-29||

\twolineshloka
{एतेऽसितमुखा घोरा गोलाङ्गूला महाबलाः}
{शतं शतसहस्राणि दृष्ट्वा वै सेतुबन्धनम्} %||6-18-30||

\twolineshloka
{गोलाङ्गूलं महावेगं गवाक्षं नाम यूथपम्}
{परिवार्याभिवर्तन्ते लङ्कां मर्दितुमोजसा} %||6-18-31||

\twolineshloka
{भ्रमराचरिता यत्र सर्वकामफलद्रुमाः}
{यं सूर्यतुल्यवर्णाभमनुपर्येति पर्वतम्} %||6-18-32||

\twolineshloka
{यस्य भासा सदा भान्ति तद्वर्णा मृगपक्षिणः}
{यस्य प्रस्थं महात्मानो न त्यजन्ति महर्षयः} %||6-18-33||

\twolineshloka
{तत्रैष रमते राजन्रम्ये काञ्चनपर्वते}
{मुख्यो वानरमुख्यानां केसरी नाम यूथपः} %||6-18-34||

\twolineshloka
{षष्टिर्गिरिसहस्राणां रम्याः काञ्चनपर्वताः}
{तेषां मध्ये गिरिवरस्त्वमिवानघ रक्षसाम्} %||6-18-35||

\twolineshloka
{तत्रैते कपिलाः श्वेतास्ताम्रास्या मधुपिङ्गलाः}
{निवसन्त्युत्तमगिरौ तीक्ष्णदंष्ट्रानखायुधाः} %||6-18-36||

\twolineshloka
{सिंह इव चतुर्दंष्ट्रा व्याघ्रा इव दुरासदाः}
{सर्वे वैश्वनरसमा ज्वलिताशीविषोपमाः} %||6-18-37||

\twolineshloka
{सुदीर्घाञ्चितलाङ्गूला मत्तमातङ्गसंनिभाः}
{महापर्वतसङ्काशा महाजीमूतनिस्वनाः} %||6-18-38||

\threelineshloka
{एष चैषामधिपतिर्मध्ये तिष्ठति वीर्यवान्}
{नाम्ना पृथिव्यां विख्यातो राजञ्शतबलीति यः}
{एषैवाशंसते लङ्कां स्वेनानीकेन मर्दितुम्} %||6-18-39||

\twolineshloka
{गजो गवाक्षो गवयो नलो नीलश्च वानरः}
{एकैक एव यूथानां कोटिभिर्दशभिर्वृतः} %||6-18-40||

\twolineshloka
{तथान्ये वानरश्रेष्ठा विन्ध्यपर्वतवासिनः}
{न शक्यन्ते बहुत्वात्तु सङ्ख्यातुं लघुविक्रमाः} %||6-18-41||

\fourlineindentedshloka
{सर्वे महाराज महाप्रभावाः}
{सर्वे महाशैलनिकाशकायाः}
{सर्वे समर्थाः पृथिवीं क्षणेन}
{कर्तुं प्रविध्वस्तविकीर्णशैलाम्} %||6-19-42||


{॥इत्यार्षे श्रीमद्रामायणे वाल्मीकीये आदिकाव्ये युद्धकाण्डे अष्टादशः सर्गः॥}

\sect{एकोनविंशः सर्गः}

\twolineshloka
{सारणस्य वचः श्रुत्वा रावणं राक्षसाधिपम्}
{बलमालोकयन्सर्वं शुको वाक्यमथाब्रवीत्} %||6-19-1||

\twolineshloka
{स्थितान्पश्यसि यानेतान्मत्तानिव महाद्विपान्}
{न्यग्रोधानिव गाङ्गेयान्सालान्हैमवतीनिव} %||6-19-2||

\twolineshloka
{एते दुष्प्रसहा राजन्बलिनः कामरूपिणः}
{दैत्यदानवसङ्काशा युद्धे देवपराक्रमाः} %||6-19-3||

\twolineshloka
{एषां कोटिसहस्राणि नव पञ्चच सप्त च}
{तथा शङ्खसहस्राणि तथा वृन्दशतानि च} %||6-19-4||

\twolineshloka
{एते सुग्रीवसचिवाः किष्किन्धानिलयाः सदा}
{हरयो देवगन्धर्वैरुत्पन्नाः कामरूपिणः} %||6-19-5||

\twolineshloka
{यौ तौ पश्यसि तिष्ठन्तौ कुमारौ देवरूपिणौ}
{मैन्दश्च द्विविदश्चोभौ ताभ्यां नास्ति समो युधि} %||6-19-6||

\twolineshloka
{ब्रह्मणा समनुज्ञातावमृतप्राशिनावुभौ}
{आशंसेते युधा लङ्कामेतौ मर्दितुमोजसा} %||6-19-7||

\twolineshloka
{यावेतावेतयोः पार्श्वे स्थितौ पर्वतसंनिभौ}
{सुमुखो विमुखश्चैव मृत्युपुत्रौ पितुः समौ} %||6-19-8||

\twolineshloka
{यं तु पश्यसि तिष्ठन्तं प्रभिन्नमिव कुञ्जरम्}
{यो बलात्क्षोभयेत्क्रुद्धः समुद्रमपि वानरः} %||6-19-9||

\twolineshloka
{एषोऽभिगन्ता लङ्काया वैदेह्यास्तव च प्रभो}
{एनं पश्य पुरा दृष्टं वानरं पुनरागतम्} %||6-19-10||

\twolineshloka
{ज्येष्ठः केसरिणः पुत्रो वातात्मज इति श्रुतः}
{हनूमानिति विख्यातो लङ्घितो येन सागरः} %||6-19-11||

\twolineshloka
{कामरूपी हरिश्रेष्ठो बलरूपसमन्वितः}
{अनिवार्यगतिश्चैव यथा सततगः प्रभुः} %||6-19-12||

\twolineshloka
{उद्यन्तं भास्करं दृष्ट्वा बालः किल पिपासितः}
{त्रियोजनसहस्रं तु अध्वानमवतीर्य हि} %||6-19-13||

\twolineshloka
{आदित्यमाहरिष्यामि न मे क्षुत्प्रतियास्यति}
{इति सञ्चिन्त्य मनसा पुरैष बलदर्पितः} %||6-19-14||

\twolineshloka
{अनाधृष्यतमं देवमपि देवर्षिदानवैः}
{अनासाद्यैव पतितो भास्करोदयने गिरौ} %||6-19-15||

\twolineshloka
{पतितस्य कपेरस्य हनुरेका शिलातले}
{किञ्चिद्भिन्ना दृढहनोर्हनूमानेष तेन वै} %||6-19-16||

\twolineshloka
{सत्यमागमयोगेन ममैष विदितो हरिः}
{नास्य शक्यं बलं रूपं प्रभावो वानुभाषितुम्} %||6-19-17||

\twolineshloka
{एष आशंसते लङ्कामेको मर्दितुमोजसा}
{यश्चैषोऽनन्तरः शूरः श्यामः पद्मनिभेक्षणः} %||6-19-18||

\twolineshloka
{इक्ष्वाकूणामतिरथो लोके विख्यात पौरुषः}
{यस्मिन्न चलते धर्मो यो धर्मं नातिवर्तते} %||6-19-19||

\twolineshloka
{यो ब्राह्ममस्त्रं वेदांश्च वेद वेदविदां वरः}
{यो भिन्द्याद्गगनं बाणैः पर्वतांश्चापि दारयेत्} %||6-19-20||

\twolineshloka
{यस्य मृत्योरिव क्रोधः शक्रस्येव पराक्रमः}
{स एष रामस्त्वां योद्धुं राजन्समभिवर्तते} %||6-19-21||

\twolineshloka
{यश्चैष दक्षिणे पार्श्वे शुद्धजाम्बूनदप्रभः}
{विशालवक्षास्ताम्राक्षो नीलकुञ्चितमूर्धजः} %||6-19-22||

\twolineshloka
{एषोऽस्य लक्ष्मणो नाम भ्राता प्राणसमः प्रियः}
{नये युद्धे च कुशलः सर्वशास्त्रविशारदः} %||6-19-23||

\twolineshloka
{अमर्षी दुर्जयो जेता विक्रान्तो बुद्धिमान्बली}
{रामस्य दक्षिणो बाहुर्नित्यं प्राणो बहिश्चरः} %||6-19-24||

\twolineshloka
{न ह्येष राघवस्यार्थे जीवितं परिरक्षति}
{एषैवाशंसते युद्धे निहन्तुं सर्वराक्षसान्} %||6-19-25||

\twolineshloka
{यस्तु सव्यमसौ पक्षं रामस्याश्रित्य तिष्ठति}
{रक्षोगणपरिक्षिप्तो राजा ह्येष विभीषणः} %||6-19-26||

\twolineshloka
{श्रीमता राजराजेन लङ्कायामभिषेचितः}
{त्वामेव प्रतिसंरब्धो युद्धायैषोऽभिवर्तते} %||6-19-27||

\twolineshloka
{यं तु पश्यसि तिष्ठन्तं मध्ये गिरिमिवाचलम्}
{सर्वशाखामृगेन्द्राणां भर्तारमपराजितम्} %||6-19-28||

\twolineshloka
{तेजसा यशसा बुद्ध्या ज्ञानेनाभिजनेन च}
{यः कपीनति बभ्राज हिमवानिव पर्वतान्} %||6-19-29||

\twolineshloka
{किष्किन्धां यः समध्यास्ते गुहां सगहनद्रुमाम्}
{दुर्गां पर्वतदुर्गस्थां प्रधानैः सह यूथपैः} %||6-19-30||

\twolineshloka
{यस्यैषा काञ्चनी माला शोभते शतपुष्करा}
{कान्ता देवमनुष्याणां यस्यां लक्ष्मीः प्रतिष्ठिता} %||6-19-31||

\twolineshloka
{एतां च मालां तारां च कपिराज्यं च शाश्वतम्}
{सुग्रीवो वालिनं हत्वा रामेण प्रतिपादितः} %||6-19-32||

\twolineshloka
{एवं कोटिसहस्रेण शङ्कूनां च शतेन च}
{सुग्रीवो वानरेन्द्रस्त्वां युद्धार्थमभिवर्तते} %||6-19-33||

\fourlineindentedshloka
{इमां महाराजसमीक्ष्य वाहिनीम्}
{उपस्थितां प्रज्वलितग्रहोपमाम्}
{ततः प्रयत्नः परमो विधीयताम्}
{यथा जयः स्यान्न परैः पराजयः} %||6-20-34||


{॥इत्यार्षे श्रीमद्रामायणे वाल्मीकीये आदिकाव्ये युद्धकाण्डे एकोनविंशः सर्गः॥}

\sect{विंशः सर्गः}

\twolineshloka
{शुकेन तु समाख्यातांस्तान्दृष्ट्वा हरियूथपान्}
{समीपस्थं च रामस्य भ्रातरं स्वं विभीषणम्} %||6-20-1||

\twolineshloka
{लक्ष्मणं च महावीर्यं भुजं रामस्य दक्षिणम्}
{सर्ववानरराजं च सुग्रीवं भीमविक्रमम्} %||6-20-2||

\twolineshloka
{किञ्चिदाविग्नहृदयो जातक्रोधश्च रावणः}
{भर्त्सयामास तौ वीरौ कथान्ते शुकसारणौ} %||6-20-3||

\twolineshloka
{अधोमुखौ तौ प्रणतावब्रवीच्छुकसारणौ}
{रोषगद्गदया वाचा संरब्धः परुषं वचः} %||6-20-4||

\twolineshloka
{न तावत्सदृशं नाम सचिवैरुपजीविभिः}
{विप्रियं नृपतेर्वक्तुं निग्रहप्रग्रहे विभोः} %||6-20-5||

\twolineshloka
{रिपूणां प्रतिकूलानां युद्धार्थमभिवर्तताम्}
{उभाभ्यां सदृशं नाम वक्तुमप्रस्तवे स्तवम्} %||6-20-6||

\twolineshloka
{आचार्या गुरवो वृद्धा वृथा वां पर्युपासिताः}
{सारं यद्राजशास्त्राणामनुजीव्यं न गृह्यते} %||6-20-7||

\twolineshloka
{गृहीतो वा न विज्ञातो भारो ज्ञानस्य वोछ्यते}
{ईदृशैः सचिवैर्युक्तो मूर्खैर्दिष्ट्या धराम्यहम्} %||6-20-8||

\twolineshloka
{किं नु मृत्योर्भयं नास्ति मां वक्तुं परुषं वचः}
{यस्य मे शासतो जिह्वा प्रयच्छति शुभाशुभम्} %||6-20-9||

\twolineshloka
{अप्येव दहनं स्पृष्ट्वा वने तिष्ठन्ति पादपाः}
{राजदोषपरामृष्टास्तिष्ठन्ते नापराधिनः} %||6-20-10||

\twolineshloka
{हन्यामहमिमौ पापौ शत्रुपक्षप्रशंसकौ}
{यदि पूर्वोपकारैर्मे न क्रोधो मृदुतां व्रजेत्} %||6-20-11||

\threelineshloka
{अपध्वंसत गच्छध्वं संनिकर्षादितो मम}
{न हि वां हन्तुमिच्छामि स्मरन्नुपकृतानि वाम्}
{हतावेव कृतघ्नौ तौ मयि स्नेहपराङ्मुखौ} %||6-20-12||

\twolineshloka
{एवमुक्तौ तु सव्रीडौ तावुभौ शुकसारणौ}
{रावणं जयशब्देन प्रतिनन्द्याभिनिःसृतौ} %||6-20-13||

\twolineshloka
{अब्रवीत्स दशग्रीवः समीपस्थं महोदरम्}
{उपस्थापय शीघ्रं मे चारान्नीतिविशारदान्} %||6-20-14||

\twolineshloka
{ततश्चराः सन्त्वरिताः प्राप्ताः पार्थिवशासनात्}
{उपस्थिताः प्राञ्जलयो वर्धयित्वा जयाशिषा} %||6-20-15||

\twolineshloka
{तानब्रवीत्ततो वाक्यं रावणो राक्षसाधिपः}
{चारान्प्रत्ययिकाञ्शूरान्भक्तान्विगतसाध्वसान्} %||6-20-16||

\twolineshloka
{इतो गच्छत रामस्य व्यवसायं परीक्षथ}
{मन्त्रेष्वभ्यन्तरा येऽस्य प्रीत्या तेन समागताः} %||6-20-17||

\twolineshloka
{कथं स्वपिति जागर्ति किमन्यच्च करिष्यति}
{विज्ञाय निपुणं सर्वमागन्तव्यमशेषतः} %||6-20-18||

\twolineshloka
{चारेण विदितः शत्रुः पण्डितैर्वसुधाधिपैः}
{युद्धे स्वल्पेन यत्नेन समासाद्य निरस्यते} %||6-20-19||

\twolineshloka
{चारास्तु ते तथेत्युक्त्वा प्रहृष्टा राक्षसेश्वरम्}
{कृत्वा प्रदक्षिणं जग्मुर्यत्र रामः सलक्ष्मणः} %||6-20-20||

\twolineshloka
{ते सुवेलस्य शैलस्य समीपे रामलक्ष्मणौ}
{प्रच्छन्ना ददृशुर्गत्वा ससुग्रीवविभीषणौ} %||6-20-21||

\twolineshloka
{ते तु धर्मात्मना दृष्टा राक्षसेन्द्रेण राक्षसाः}
{विभीषणेन तत्रस्था निगृहीता यदृच्छया} %||6-20-22||

\twolineshloka
{वानरैरर्दितास्ते तु विक्रान्तैर्लघुविक्रमैः}
{पुनर्लङ्कामनुप्राप्ताः श्वसन्तो नष्टचेतसः} %||6-20-23||

\fourlineindentedshloka
{ततो दशग्रीवमुपस्थितास्ते}
{चारा बहिर्नित्यचरा निशाचराः}
{गिरेः सुवेलस्य समीपवासिनम्}
{न्यवेदयन्भीमबलं महाबलाः} %||6-21-24||


{॥इत्यार्षे श्रीमद्रामायणे वाल्मीकीये आदिकाव्ये युद्धकाण्डे विंशः सर्गः॥}

\sect{एकविंशः सर्गः}

\twolineshloka
{ततस्तमक्षोभ्य बलं लङ्काधिपतये चराः}
{सुवेले राघवं शैले निविष्टं प्रत्यवेदयन्} %||6-21-1||

\twolineshloka
{चाराणां रावणः श्रुत्वा प्राप्तं रामं महाबलम्}
{जातोद्वेगोऽभवत्किञ्चिच्छार्दूलं वाक्यमब्रवीत्} %||6-21-2||

\twolineshloka
{अयथावच्च ते वर्णो दीनश्चासि निशाचर}
{नासि कच्चिदमित्राणां क्रुद्धानां वशमागतः} %||6-21-3||

\twolineshloka
{इति तेनानुशिष्टस्तु वाचं मन्दमुदीरयत्}
{तदा राक्षसशार्दूलं शार्दूलो भयविह्वलः} %||6-21-4||

\twolineshloka
{न ते चारयितुं शक्या राजन्वानरपुङ्गवाः}
{विक्रान्ता बलवन्तश्च राघवेण च रक्षिताः} %||6-21-5||

\twolineshloka
{नापि सम्भाषितुं शक्याः सम्प्रश्नोऽत्र न लभ्यते}
{सर्वतो रक्ष्यते पन्था वानरैः पर्वतोपमैः} %||6-21-6||

\twolineshloka
{प्रविष्टमात्रे ज्ञातोऽहं बले तस्मिन्नचारिते}
{बलाद्गृहीतो बहुभिर्बहुधास्मि विदारितः} %||6-21-7||

\twolineshloka
{जानुभिर्मुष्टिभिर्दन्तैस्तलैश्चाभिहतो भृशम्}
{परिणीतोऽस्मि हरिभिर्बलवद्भिरमर्षणैः} %||6-21-8||

\twolineshloka
{परिणीय च सर्वत्र नीतोऽहं रामसंसदम्}
{रुधिरादिग्धसर्वाङ्गो विह्वलश्चलितेन्द्रियः} %||6-21-9||

\twolineshloka
{हरिभिर्वध्यमानश्च याचमानः कृताञ्जलिः}
{राघवेण परित्रातो जीवामि ह यदृच्छया} %||6-21-10||

\twolineshloka
{एष शैलैः शिलाभिश्च पूरयित्वा महार्णवम्}
{द्वारमाश्रित्य लङ्काया रामस्तिष्ठति सायुधः} %||6-21-11||

\twolineshloka
{गरुडव्यूहमास्थाय सर्वतो हरिभिर्वृतः}
{मां विसृज्य महातेजा लङ्कामेवाभिवर्तते} %||6-21-12||

\twolineshloka
{पुरा प्राकारमायाति क्षिप्रमेकतरं कुरु}
{सीतां चास्मै प्रयच्छाशु सुयुद्धं वा प्रदीयताम्} %||6-21-13||

\twolineshloka
{मनसा सन्ततापाथ तच्छ्रुत्वा राक्षसाधिपः}
{शार्दूलस्य महद्वाक्यमथोवाच स रावणः} %||6-21-14||

\twolineshloka
{यदि मां प्रतियुध्येरन्देवगन्धर्वदानवाः}
{नैव सीतां प्रदास्यामि सर्वलोकभयादपि} %||6-21-15||

\twolineshloka
{एवमुक्त्वा महातेजा रावणः पुनरब्रवीत्}
{चारिता भवता सेना केऽत्र शूराः प्लवङ्गमाः} %||6-21-16||

\twolineshloka
{कीदृशाः किम्प्रभावाश्च वानरा ये दुरासदाः}
{कस्य पुत्राश्च पौत्राश्च तत्त्वमाख्याहि राक्षस} %||6-21-17||

\twolineshloka
{तत्रत्र प्रतिपत्स्यामि ज्ञात्वा तेषां बलाबलम्}
{अवश्यं बलसङ्ख्यानं कर्तव्यं युद्धमिच्छता} %||6-21-18||

\twolineshloka
{अथैवमुक्तः शार्दूलो रावणेनोत्तमश्चरः}
{इदं वचनमारेभे वक्तुं रावणसंनिधौ} %||6-21-19||

\twolineshloka
{अथर्क्षरजसः पुत्रो युधि राजन्सुदुर्जयः}
{गद्गदस्याथ पुत्रोऽत्र जाम्बवानिति विश्रुतः} %||6-21-20||

\twolineshloka
{गद्गदस्यैव पुत्रोऽन्यो गुरुपुत्रः शतक्रतोः}
{कदनं यस्य पुत्रेण कृतमेकेन रक्षसाम्} %||6-21-21||

\twolineshloka
{सुषेणश्चापि धर्मात्मा पुत्रो धर्मस्य वीर्यवान्}
{सौम्यः सोमात्मजश्चात्र राजन्दधिमुखः कपिः} %||6-21-22||

\twolineshloka
{सुमुखो दुर्मुखश्चात्र वेगदर्शी च वानरः}
{मृत्युर्वानररूपेण नूनं सृष्टः स्वयम्भुवा} %||6-21-23||

\twolineshloka
{पुत्रो हुतवहस्याथ नीलः सेनापतिः स्वयम्}
{अनिलस्य च पुत्रोऽत्र हनूमानिति विश्रुतः} %||6-21-24||

\twolineshloka
{नप्ता शक्रस्य दुर्धर्षो बलवानङ्गदो युवा}
{मैन्दश्च द्विविदश्चोभौ बलिनावश्विसम्भवौ} %||6-21-25||

\twolineshloka
{पुत्रा वैवस्वतस्यात्र पञ्चकालान्तकोपमाः}
{गजो गवाक्षो गवयः शरभो गन्धमादनः} %||6-21-26||

\twolineshloka
{श्वेतो ज्योतिर्मुखश्चात्र भास्करस्यात्मसम्भवौ}
{वरुणस्य च पुत्रोऽथ हेमकूटः प्लवङ्गमः} %||6-21-27||

\twolineshloka
{विश्वकर्मसुतो वीरो नलः प्लवगसत्तमः}
{विक्रान्तो वेगवानत्र वसुपुत्रः सुदुर्धरः} %||6-21-28||

\twolineshloka
{दशवानरकोट्यश्च शूराणां युद्धकाङ्क्षिणाम्}
{श्रीमतां देवपुत्राणां शेषान्नाख्यातुमुत्सहे} %||6-21-29||

\twolineshloka
{पुत्रो दशरथस्यैष सिंहसंहननो युवा}
{दूषणो निहतो येन खरश्च त्रिशिरास्तथा} %||6-21-30||

\twolineshloka
{नास्ति रामस्य सदृशो विक्रमे भुवि कश्चन}
{विराधो निहतो येन कबन्धश्चान्तकोपमः} %||6-21-31||

\twolineshloka
{वक्तुं न शक्तो रामस्य नरः कश्चिद्गुणान्क्षितौ}
{जनस्थानगता येन तावन्तो राक्षसा हताः} %||6-21-32||

\twolineshloka
{लक्ष्मणश्चात्र धर्मात्मा मातङ्गानामिवर्षभः}
{यस्य बाणपथं प्राप्य न जीवेदपि वासवः} %||6-21-33||

\twolineshloka
{राक्षसानां वरिष्ठश्च तव भ्राता विभीषणः}
{परिगृह्य पुरीं लङ्कां राघवस्य हिते रतः} %||6-21-34||

\twolineshloka
{इति सर्वं समाख्यातं तवेदं वानरं बलम्}
{सुवेलेऽधिष्ठितं शैले शेषकार्ये भवान्गतिः} %||6-22-35||


{॥इत्यार्षे श्रीमद्रामायणे वाल्मीकीये आदिकाव्ये युद्धकाण्डे एकविंशः सर्गः॥}

\sect{द्वाविंशः सर्गः}

\twolineshloka
{ततस्तमक्षोभ्यबलं लङ्कायां नृपतेश्चरः}
{सुवेले राघवं शैले निविष्टं प्रत्यवेदयन्} %||6-22-1||

\twolineshloka
{चाराणां रावणः श्रुत्वा प्राप्तं रामं महाबलम्}
{जातोद्वेगोऽभवत्किञ्चित्सचिवांश्चेदमब्रवीत्} %||6-22-2||

\twolineshloka
{मन्त्रिणः शीघ्रमायान्तु सर्वे वै सुसमाहिताः}
{अयं नो मन्त्रकालो हि सम्प्राप्त इव राक्षसाः} %||6-22-3||

\twolineshloka
{तस्य तच्छासनं श्रुत्वा मन्त्रिणोऽभ्यागमन्द्रुतम्}
{ततः सम्मन्त्रयामास सचिवै राक्षसैः सह} %||6-22-4||

\twolineshloka
{मन्त्रयित्वा स दुर्धर्षः क्षमं यत्समनन्तरम्}
{विसर्जयित्वा सचिवान्प्रविवेश स्वमालयम्} %||6-22-5||

\twolineshloka
{ततो राक्षसमाहूय विद्युज्जिह्वं महाबलम्}
{मायाविदं महामायः प्राविशद्यत्र मैथिली} %||6-22-6||

\twolineshloka
{विद्युज्जिह्वं च मायाज्ञमब्रवीद्राक्षसाधिपः}
{मोहयिष्यामहे सीतां मायया जनकात्मजाम्} %||6-22-7||

\twolineshloka
{शिरो मायामयं गृह्य राघवस्य निशाचर}
{मां त्वं समुपतिष्ठस्व महच्च सशरं धनुः} %||6-22-8||

\twolineshloka
{एवमुक्तस्तथेत्याह विद्युज्जिह्वो निशाचरः}
{तस्य तुष्टोऽभवद्राजा प्रददौ च विभूषणम्} %||6-22-9||

\threelineshloka
{अशोकवनिकायां तु प्रविवेश महाबलः}
{ततो दीनामदैन्यार्हां ददर्श धनदानुजः}
{अधोमुखीं शोकपरामुपविष्टां महीतले} %||6-22-10||

\twolineshloka
{भर्तारमेव ध्यायन्तीमशोकवनिकां गताम्}
{उपास्यमानां घोराभी राक्षसीभिरदूरतः} %||6-22-11||

\twolineshloka
{उपसृत्य ततः सीतां प्रहर्षन्नाम कीर्तयन्}
{इदं च वचनं धृष्टमुवाच जनकात्मजाम्} %||6-22-12||

\twolineshloka
{सान्त्व्यमाना मया भद्रे यमुपाश्रित्य वल्गसे}
{खर हन्ता स ते भर्ता राघवः समरे हतः} %||6-22-13||

\twolineshloka
{छिन्नं ते सर्वतो मूलं दर्पस्ते निहतो मया}
{व्यसनेनात्मनः सीते मम भार्या भविष्यसि} %||6-22-14||

\twolineshloka
{अल्पपुण्ये निवृत्तार्थे मूढे पण्डितमानिनि}
{शृणु भर्तृबधं सीते घोरं वृत्रवधं यथा} %||6-22-15||

\twolineshloka
{समायातः समुद्रान्तं मां हन्तुं किल राघवः}
{वानरेन्द्रप्रणीतेन बलेन महता वृतः} %||6-22-16||

\twolineshloka
{संनिविष्टः समुद्रस्य तीरमासाद्य दक्षिणम्}
{बलेन महता रामो व्रजत्यस्तं दिवाकरे} %||6-22-17||

\twolineshloka
{अथाध्वनि परिश्रान्तमर्धरात्रे स्थितं बलम्}
{सुखसुप्तं समासाद्य चारितं प्रथमं चरैः} %||6-22-18||

\twolineshloka
{तत्प्रहस्तप्रणीतेन बलेन महता मम}
{बलमस्य हतं रात्रौ यत्र रामः सुलक्ष्मणः} %||6-22-19||

\twolineshloka
{पट्टसान्परिघान्खड्गांश्चक्रान्दण्डान्महायसान्}
{बाणजालानि शूलानि भास्वरान्कूटमुद्गरान्} %||6-22-20||

\twolineshloka
{यष्टीश्च तोमरान्प्रासंश्चक्राणि मुसलानि च}
{उद्यम्योद्यम्य रक्षोभिर्वानरेषु निपातिताः} %||6-22-21||

\twolineshloka
{अथ सुप्तस्य रामस्य प्रहस्तेन प्रमाथिना}
{असक्तं कृतहस्तेन शिरश्छिन्नं महासिना} %||6-22-22||

\twolineshloka
{विभीषणः समुत्पत्य निगृहीतो यदृच्छया}
{दिशः प्रव्राजितः सर्वैर्लक्ष्मणः प्लवगैः सह} %||6-22-23||

\twolineshloka
{सुग्रीवो ग्रीवया शेते भग्नया प्लवगाधिपः}
{निरस्तहनुकः शेते हनूमान्राक्षसैर्हतः} %||6-22-24||

\twolineshloka
{जाम्बवानथ जानुभ्यामुत्पतन्निहतो युधि}
{पट्टसैर्बहुभिश्छिन्नो निकृत्तः पादपो यथा} %||6-22-25||

\twolineshloka
{मैन्दश्च द्विविदश्चोभौ निहतौ वानरर्षभौ}
{निःश्वसन्तौ रुदन्तौ च रुधिरेण समुक्षितौ} %||6-22-26||

\twolineshloka
{असिनाभ्याहतश्छिन्नो मध्ये रिपुनिषूदनः}
{अभिष्टनति मेदिन्यां पनसः पनसो यथा} %||6-22-27||

\twolineshloka
{नाराचैर्बहुभिश्छिन्नः शेते दर्यां दरीमुखः}
{कुमुदस्तु महातेजा निष्कूजन्सायकैर्हतः} %||6-22-28||

\twolineshloka
{अङ्गदो बहुभिश्छिन्नः शरैरासाद्य राक्षसैः}
{पातितो रुधिरोद्गारी क्षितौ निपतितोऽङ्गदः} %||6-22-29||

\twolineshloka
{हरयो मथिता नागै रथजालैस्तथापरे}
{शायिता मृदितास्तत्र वायुवेगैरिवाम्बुदाः} %||6-22-30||

\twolineshloka
{प्रद्रुताश्च परे त्रस्ता हन्यमाना जघन्यतः}
{अभिद्रुतास्तु रक्षोभिः सिंहैरिव महाद्विपाः} %||6-22-31||

\twolineshloka
{सागरे पतिताः केचित्केचिद्गगनमाश्रिताः}
{ऋक्षा वृक्षानुपारूढा वानरैस्तु विमिश्रिताः} %||6-22-32||

\twolineshloka
{सागरस्य च तीरेषु शैलेषु च वनेषु च}
{पिङ्गाक्षास्ते विरूपाक्षैर्बहुभिर्बहवो हताः} %||6-22-33||

\twolineshloka
{एवं तव हतो भर्ता ससैन्यो मम सेनया}
{क्षतजार्द्रं रजोध्वस्तमिदं चास्याहृतं शिरः} %||6-22-34||

\twolineshloka
{ततः परमदुर्धर्षो रावणो राक्षसेश्वरः}
{सीतायामुपशृण्वन्त्यां राक्षसीमिदमब्रवीत्} %||6-22-35||

\twolineshloka
{राक्षसं क्रूरकर्माणं विद्युज्जिह्वं त्वमानय}
{येन तद्राघवशिरः सङ्ग्रामात्स्वयमाहृतम्} %||6-22-36||

\twolineshloka
{विद्युज्जिह्वस्ततो गृह्य शिरस्तत्सशरासनम्}
{प्रणामं शिरसा कृत्वा रावणस्याग्रतः स्थितः} %||6-22-37||

\twolineshloka
{तमब्रवीत्ततो राजा रावणो राक्षसं स्थितम्}
{विद्युज्जिह्वं महाजिह्वं समीपपरिवर्तिनम्} %||6-22-38||

\twolineshloka
{अग्रतः कुरु सीतायाः शीघ्रं दाशरथेः शिरः}
{अवस्थां पश्चिमां भर्तुः कृपणा साधु पश्यतु} %||6-22-39||

\twolineshloka
{एवमुक्तं तु तद्रक्षः शिरस्तत्प्रियदर्शनम्}
{उपनिक्षिप्य सीतायाः क्षिप्रमन्तरधीयत} %||6-22-40||

\twolineshloka
{रावणश्चापि चिक्षेप भास्वरं कार्मुकं महत्}
{त्रिषु लोकेषु विख्यातं सीतामिदमुवाच ह} %||6-22-41||

\twolineshloka
{इदं तत्तव रामस्य कार्मुकं ज्यासमन्वितम्}
{इह प्रहस्तेनानीतं हत्वा तं निशि मानुषम्} %||6-22-42||

\fourlineindentedshloka
{स विद्युज्जिह्वेन सहैव तच्छिरो}
{धनुश्च भूमौ विनिकीर्य रावणः}
{विदेहराजस्य सुतां यशस्विनीम्}
{ततोऽब्रवीत्तां भव मे वशानुगा} %||6-23-43||


{॥इत्यार्षे श्रीमद्रामायणे वाल्मीकीये आदिकाव्ये युद्धकाण्डे द्वाविंशः सर्गः॥}

\sect{त्रयोविंशः सर्गः}

\twolineshloka
{सा सीता तच्छिरो दृष्ट्वा तच्च कार्मुकमुत्तमम्}
{सुग्रीवप्रतिसंसर्गमाख्यातं च हनूमता} %||6-23-1||

\twolineshloka
{नयने मुखवर्णं च भर्तुस्तत्सदृशं मुखम्}
{केशान्केशान्तदेशं च तं च चूडामणिं शुभम्} %||6-23-2||

\twolineshloka
{एतैः सर्वैरभिज्ञानैरभिज्ञाय सुदुःखिता}
{विजगर्हेऽथ कैकेयीं क्रोशन्ती कुररी यथा} %||6-23-3||

\twolineshloka
{सकामा भव कैकेयि हतोऽयं कुलनन्दनः}
{कुलमुत्सादितं सर्वं त्वया कलहशीलया} %||6-23-4||

\twolineshloka
{आर्येण किं नु कैकेय्याः कृतं रामेण विप्रियम्}
{यद्गृहाच्चीरवसनस्तया प्रस्थापितो वनम्} %||6-23-5||

\twolineshloka
{एवमुक्त्वा तु वैदेही वेपमाना तपस्विनी}
{जगाम जगतीं बाला छिन्ना तु कदली यथा} %||6-23-6||

\twolineshloka
{सा मुहूर्तात्समाश्वस्य प्रतिलभ्य च चेतनाम्}
{तच्छिरः समुपाघ्राय विललापायतेक्षणा} %||6-23-7||

\twolineshloka
{हा हतास्मि महाबाहो वीरव्रतमनुव्रता}
{इमां ते पश्चिमावस्थां गतास्मि विधवा कृता} %||6-23-8||

\twolineshloka
{प्रथमं मरणं नार्या भर्तुर्वैगुण्यमुच्यते}
{सुवृत्तः साधुवृत्तायाः संवृत्तस्त्वं ममाग्रतः} %||6-23-9||

\twolineshloka
{दुःखाद्दुःखं प्रपन्नाया मग्नायाः शोकसागरे}
{यो हि मामुद्यतस्त्रातुं सोऽपि त्वं विनिपातितः} %||6-23-10||

\twolineshloka
{सा श्वश्रूर्मम कौसल्या त्वया पुत्रेण राघव}
{वत्सेनेव यथा धेनुर्विवत्सा वत्सला कृता} %||6-23-11||

\twolineshloka
{आदिष्टं दीर्घमायुस्ते यैरचिन्त्यपराक्रम}
{अनृतं वचनं तेषामल्पायुरसि राघव} %||6-23-12||

\twolineshloka
{अथ वा नश्यति प्रज्ञा प्राज्ञस्यापि सतस्तव}
{पचत्येनं तथा कालो भूतानां प्रभवो ह्ययम्} %||6-23-13||

\twolineshloka
{अदृष्टं मृत्युमापन्नः कस्मात्त्वं नयशास्त्रवित्}
{व्यसनानामुपायज्ञः कुशलो ह्यसि वर्जने} %||6-23-14||

\twolineshloka
{तथा त्वं सम्परिष्वज्य रौद्रयातिनृशंसया}
{कालरात्र्या मयाच्छिद्य हृतः कमललोचनः} %||6-23-15||

\twolineshloka
{उपशेषे महाबाहो मां विहाय तपस्विनीम्}
{प्रियामिव शुभां नारीं पृथिवीं पुरुषर्षभ} %||6-23-16||

\twolineshloka
{अर्चितं सततं यत्नाद्गन्धमाल्यैर्मया तव}
{इदं ते मत्प्रियं वीर धनुः काञ्चनभूषितम्} %||6-23-17||

\twolineshloka
{पित्रा दशरथेन त्वं श्वशुरेण ममानघ}
{पूर्वैश्च पितृभिः सार्धं नूनं स्वर्गे समागतः} %||6-23-18||

\twolineshloka
{दिवि नक्षत्रभूतस्त्वं महत्कर्म कृतं प्रियम्}
{पुण्यं राजर्षिवंशं त्वमात्मनः समुपेक्षसे} %||6-23-19||

\twolineshloka
{किं मान्न प्रेक्षसे राजन्किं मां न प्रतिभाषसे}
{बालां बालेन सम्प्राप्तां भार्यां मां सहचारिणीम्} %||6-23-20||

\twolineshloka
{संश्रुतं गृह्णता पाणिं चरिष्यामीति यत्त्वया}
{स्मर तन्मम काकुत्स्थ नय मामपि दुःखिताम्} %||6-23-21||

\twolineshloka
{कस्मान्मामपहाय त्वं गतो गतिमतां वर}
{अस्माल्लोकादमुं लोकं त्यक्त्वा मामिह दुःखिताम्} %||6-23-22||

\twolineshloka
{कल्याणैरुचितं यत्तत्परिष्वक्तं मयैव तु}
{क्रव्यादैस्तच्छरीरं ते नूनं विपरिकृष्यते} %||6-23-23||

\twolineshloka
{अग्निष्टोमादिभिर्यज्ञैरिष्टवानाप्तदक्षिणैः}
{अग्निहोत्रेण संस्कारं केन त्वं तु न लप्स्यसे} %||6-23-24||

\twolineshloka
{प्रव्रज्यामुपपन्नानां त्रयाणामेकमागतम्}
{परिप्रक्ष्यति कौसल्या लक्ष्मणं शोकलालसा} %||6-23-25||

\twolineshloka
{स तस्याः परिपृच्छन्त्या वधं मित्रबलस्य ते}
{तव चाख्यास्यते नूनं निशायां राक्षसैर्वधम्} %||6-23-26||

\twolineshloka
{सा त्वां सुप्तं हतं श्रुत्वा मां च रक्षोगृहं गताम्}
{हृदयेन विदीर्णेन न भविष्यति राघव} %||6-23-27||

\twolineshloka
{साधु पातय मां क्षिप्रं रामस्योपरि रावणः}
{समानय पतिं पत्न्या कुरु कल्याणमुत्तमम्} %||6-23-28||

\threelineshloka
{शिरसा मे शिरश्चास्य कायं कायेन योजय}
{रावणानुगमिष्यामि गतिं भर्तुर्महात्मनः}
{मुहूर्तमपि नेच्छामि जीवितुं पापजीविना} %||6-23-29||

\twolineshloka
{श्रुतं मया वेदविदां ब्राह्मणानां पितुर्गृहे}
{यासां स्त्रीणां प्रियो भर्ता तासां लोका महोदयाः} %||6-23-30||

\twolineshloka
{क्षमा यस्मिन्दमस्त्यागः सत्यं धर्मः कृतज्ञता}
{अहिंसा चैव भूतानां तमृते का गतिर्मम} %||6-23-31||

\twolineshloka
{इति सा दुःखसन्तप्ता विललापायतेक्षणा}
{भर्तुः शिरो धनुस्तत्र समीक्ष्य जनकात्मजा} %||6-23-32||

\twolineshloka
{एवं लालप्यमानायां सीतायां तत्र राक्षसः}
{अभिचक्राम भर्तारमनीकस्थः कृताञ्जलिः} %||6-23-33||

\twolineshloka
{विजयस्वार्यपुत्रेति सोऽभिवाद्य प्रसाद्य च}
{न्यवेदयदनुप्राप्तं प्रहस्तं वाहिनीपतिम्} %||6-23-34||

\twolineshloka
{अमात्यैः सहितः सर्वैः प्रहस्तः समुपस्थितः}
{किञ्चिदात्ययिकं कार्यं तेषां त्वं दर्शनं कुरु} %||6-23-35||

\twolineshloka
{एतच्छ्रुत्वा दशग्रीवो राक्षसप्रतिवेदितम्}
{अशोकवनिकां त्यक्त्वा मन्त्रिणां दर्शनं ययौ} %||6-23-36||

\twolineshloka
{स तु सर्वं समर्थ्यैव मन्त्रिभिः कृत्यमात्मनः}
{सभां प्रविश्य विदधे विदित्वा रामविक्रमम्} %||6-23-37||

\twolineshloka
{अन्तर्धानं तु तच्छीर्षं तच्च कार्मुकमुत्तमम्}
{जगाम रावणस्यैव निर्याणसमनन्तरम्} %||6-23-38||

\twolineshloka
{राक्षसेन्द्रस्तु तैः सार्धं मन्त्रिभिर्भीमविक्रमैः}
{समर्थयामास तदा रामकार्यविनिश्चयम्} %||6-23-39||

\twolineshloka
{अविदूरस्थितान्सर्वान्बलाध्यक्षान्हितैषिणः}
{अब्रवीत्कालसदृशो रावणो राक्षसाधिपः} %||6-23-40||

\twolineshloka
{शीघ्रं भेरीनिनादेन स्फुटकोणाहतेन मे}
{समानयध्वं सैन्यानि वक्तव्यं च न कारणम्} %||6-23-41||

\fourlineindentedshloka
{ततस्तथेति प्रतिगृह्य तद्वचो}
{बलाधिपास्ते महदात्मनो बलम्}
{समानयंश्चैव समागतं च ते}
{न्यवेदयन्भर्तरि युद्धकाङ्क्षिणि} %||6-24-42||


{॥इत्यार्षे श्रीमद्रामायणे वाल्मीकीये आदिकाव्ये युद्धकाण्डे त्रयोविंशः सर्गः॥}

\sect{चतुर्विंशः सर्गः}

\twolineshloka
{सीतां तु मोहितां दृष्ट्वा सरमा नाम राक्षसी}
{आससादाशु वैदेहीं प्रियां प्रणयिनी सखी} %||6-24-1||

\twolineshloka
{सा हि तत्र कृता मित्रं सीतया रक्ष्यमाणया}
{रक्षन्ती रावणादिष्टा सानुक्रोशा दृढव्रता} %||6-24-2||

\twolineshloka
{सा ददर्श सखीं सीतां सरमा नष्टचेतनाम्}
{उपावृत्योत्थितां ध्वस्तां वडवामिव पांसुषु} %||6-24-3||

\twolineshloka
{तां समाश्वासयामास सखी स्नेहेन सुव्रता}
{उक्ता यद्रावणेन त्वं प्रत्युक्तं च स्वयं त्वया} %||6-24-4||

\threelineshloka
{सखीस्नेहेन तद्भीरु मया सर्वं प्रतिश्रुतम्}
{लीनया गनहे शूह्ये भयमुत्सृज्य रावणात्}
{तव हेतोर्विशालाक्षि न हि मे जीवितं प्रियम्} %||6-24-5||

\twolineshloka
{स सम्भ्रान्तश्च निष्क्रान्तो यत्कृते राक्षसाधिपः}
{तच्च मे विदितं सर्वमभिनिष्क्रम्य मैथिलि} %||6-24-6||

\twolineshloka
{न शक्यं सौप्तिकं कर्तुं रामस्य विदितात्मनः}
{वधश्च पुरुषव्याघ्रे तस्मिन्नेवोपपद्यते} %||6-24-7||

\twolineshloka
{न चैव वानरा हन्तुं शक्याः पादपयोधिनः}
{सुरा देवर्षभेणेव रामेण हि सुरक्षिताः} %||6-24-8||

\twolineshloka
{दीर्घवृत्तभुजः श्रीमान्महोरस्कः प्रतापवान्}
{धन्वी संहननोपेतो धर्मात्मा भुवि विश्रुतः} %||6-24-9||

\twolineshloka
{विक्रान्तो रक्षिता नित्यमात्मनश्च परस्य च}
{लक्ष्मणेन सह भ्रात्रा कुशली नयशास्त्रवित्} %||6-24-10||

\twolineshloka
{हन्ता परबलौघानामचिन्त्यबलपौरुषः}
{न हतो राघवः श्रीमान्सीते शत्रुनिबर्हणः} %||6-24-11||

\twolineshloka
{अयुक्तबुद्धिकृत्येन सर्वभूतविरोधिना}
{इयं प्रयुक्ता रौद्रेण माया मायाविदा त्वयि} %||6-24-12||

\twolineshloka
{शोकस्ते विगतः सर्वः कल्याणं त्वामुपस्थितम्}
{ध्रुवं त्वां भजते लक्ष्मीः प्रियं प्रीतिकरं शृणु} %||6-24-13||

\twolineshloka
{उत्तीर्य सागरं रामः सह वानरसेनया}
{संनिविष्टः समुद्रस्य तीरमासाद्य दक्षिणम्} %||6-24-14||

\twolineshloka
{दृष्टो मे परिपूर्णार्थः काकुत्स्थः सहलक्ष्मणः}
{सहितैः सागरान्तस्थैर्बलैस्तिष्ठति रक्षितः} %||6-24-15||

\twolineshloka
{अनेन प्रेषिता ये च राक्षसा लघुविक्रमः}
{राघवस्तीर्ण इत्येवं प्रवृत्तिस्तैरिहाहृता} %||6-24-16||

\twolineshloka
{स तां श्रुत्वा विशालाक्षि प्रवृत्तिं राक्षसाधिपः}
{एष मन्त्रयते सर्वैः सचिवैः सह रावणः} %||6-24-17||

\twolineshloka
{इति ब्रुवाणा सरमा राक्षसी सीतया सह}
{सर्वोद्योगेन सैन्यानां शब्दं शुश्राव भैरवम्} %||6-24-18||

\twolineshloka
{दण्डनिर्घातवादिन्याः श्रुत्वा भेर्या महास्वनम्}
{उवाच सरमा सीतामिदं मधुरभाषिणी} %||6-24-19||

\twolineshloka
{संनाहजननी ह्येषा भैरवा भीरु भेरिका}
{भेरीनादं च गम्भीरं शृणु तोयदनिस्वनम्} %||6-24-20||

\twolineshloka
{कल्प्यन्ते मत्तमातङ्गा युज्यन्ते रथवाजिनः}
{तत्र तत्र च संनद्धाः सम्पतन्ति पदातयः} %||6-24-21||

\twolineshloka
{आपूर्यन्ते राजमार्गाः सैन्यैरद्भुतदर्शनैः}
{वेगवद्भिर्नदद्भिश्च तोयौघैरिव सागरः} %||6-24-22||

\twolineshloka
{शास्त्राणां च प्रसन्नानां चर्मणां वर्मणां तथा}
{रथवाजिगजानां च भूषितानां च रक्षसाम्} %||6-24-23||

\twolineshloka
{प्रभां विसृजतां पश्य नानावर्णां समुत्थिताम्}
{वनं निर्दहतो धर्मे यथारूपं विभावसोः} %||6-24-24||

\twolineshloka
{घण्टानां शृणु निर्घोषं रथानां शृणु निस्वनम्}
{हयानां हेषमाणानां शृणु तूर्यध्वनिं यथा} %||6-24-25||

\twolineshloka
{उद्यतायुधहस्तानां राक्षसेन्द्रानुयायिनाम्}
{सम्भ्रमो रक्षसामेष तुमुलो लोमहर्षणः} %||6-24-26||

\twolineshloka
{श्रीस्त्वां भजति शोकघ्नी रक्षसां भयमागतम्}
{रामात्कमलपत्राक्षि दैत्यानामिव वासवात्} %||6-24-27||

\twolineshloka
{अवजित्य जितक्रोधस्तमचिन्त्यपराक्रमः}
{रावणं समरे हत्वा भर्ता त्वाधिगमिष्यति} %||6-24-28||

\twolineshloka
{विक्रमिष्यति रक्षःसु भर्ता ते सहलक्ष्मणः}
{यथा शत्रुषु शत्रुघ्नो विष्णुना सह वासवः} %||6-24-29||

\twolineshloka
{आगतस्य हि रामस्य क्षिप्रमङ्कगतां सतीम्}
{अहं द्रक्ष्यामि सिद्धार्थां त्वां शत्रौ विनिपातिते} %||6-24-30||

\twolineshloka
{अश्रूण्यानन्दजानि त्वं वर्तयिष्यसि शोभने}
{समागम्य परिष्वक्ता तस्योरसि महोरसः} %||6-24-31||

\twolineshloka
{अचिरान्मोक्ष्यते सीते देवि ते जघनं गताम्}
{धृतामेतां बहून्मासान्वेणीं रामो महाबलः} %||6-24-32||

\twolineshloka
{तस्य दृष्ट्वा मुखं देवि पूर्णचन्द्रमिवोदितम्}
{मोक्ष्यसे शोकजं वारि निर्मोकमिव पन्नगी} %||6-24-33||

\twolineshloka
{रावणं समरे हत्वा नचिरादेव मैथिलि}
{त्वया समग्रं प्रियया सुखार्हो लप्स्यते सुखम्} %||6-24-34||

\twolineshloka
{समागता त्वं रामेण मोदिष्यसि महात्मना}
{सुवर्षेण समायुक्ता यथा सस्येन मेदिनी} %||6-24-35||

\fourlineindentedshloka
{गिरिवरमभितोऽनुवर्तमानो}
{हय इव मण्डलमाशु यः करोति}
{तमिह शरणमभ्युपेहि देवि}
{दिवसकरं प्रभवो ह्ययं प्रजानाम्} %||6-25-36||


{॥इत्यार्षे श्रीमद्रामायणे वाल्मीकीये आदिकाव्ये युद्धकाण्डे चतुर्विंशः सर्गः॥}

\sect{पञ्चविंशः सर्गः}

\twolineshloka
{अथ तां जातसन्तापां तेन वाक्येन मोहिताम्}
{सरमा ह्लादयामास पृतिवीं द्यौरिवाम्भसा} %||6-25-1||

\twolineshloka
{ततस्तस्या हितं सख्याश्चिकीर्षन्ती सखी वचः}
{उवाच काले कालज्ञा स्मितपूर्वाभिभाषिणी} %||6-25-2||

\twolineshloka
{उत्सहेयमहं गत्वा त्वद्वाक्यमसितेक्षणे}
{निवेद्य कुशलं रामे प्रतिच्छन्ना निवर्तितुम्} %||6-25-3||

\twolineshloka
{न हि मे क्रममाणाया निरालम्बे विहायसि}
{समर्थो गतिमन्वेतुं पवनो गरुडोऽपि वा} %||6-25-4||

\twolineshloka
{एवं ब्रुवाणां तां सीता सरमां पुनरब्रवीत्}
{मधुरं श्लक्ष्णया वाचा पूर्वशोकाभिपन्नया} %||6-25-5||

\twolineshloka
{समर्था गगनं गन्तुमपि वा त्वं रसातलम्}
{अवगच्छाम्यकर्तव्यं कर्तव्यं ते मदन्तरे} %||6-25-6||

\twolineshloka
{मत्प्रियं यदि कर्तव्यं यदि बुद्धिः स्थिरा तव}
{ज्ञातुमिच्छामि तं गत्वा किं करोतीति रावणः} %||6-25-7||

\twolineshloka
{स हि मायाबलः क्रूरो रावणः शत्रुरावणः}
{मां मोहयति दुष्टात्मा पीतमात्रेव वारुणी} %||6-25-8||

\twolineshloka
{तर्जापयति मां नित्यं भर्त्सापयति चासकृत्}
{राक्षसीभिः सुघोराभिर्या मां रक्षन्ति नित्यशः} %||6-25-9||

\twolineshloka
{उद्विग्ना शङ्किता चास्मि न च स्वस्थं मनो मम}
{तद्भयाच्चाहमुद्विग्ना अशोकवनिकां गताः} %||6-25-10||

\twolineshloka
{यदि नाम कथा तस्य निश्चितं वापि यद्भवेत्}
{निवेदयेथाः सर्वं तत्परो मे स्यादनुग्रहः} %||6-25-11||

\twolineshloka
{सा त्वेवं ब्रुवतीं सीतां सरमा वल्गुभाषिणी}
{उवाच वचनं तस्याः स्पृशन्ती बाष्पविक्लवम्} %||6-25-12||

\twolineshloka
{एष ते यद्यभिप्रायस्तस्माद्गच्छामि जानकि}
{गृह्य शत्रोरभिप्रायमुपावृत्तां च पश्य माम्} %||6-25-13||

\twolineshloka
{एवमुक्त्वा ततो गत्वा समीपं तस्य रक्षसः}
{शुश्राव कथितं तस्य रावणस्य समन्त्रिणः} %||6-25-14||

\twolineshloka
{सा श्रुत्वा निश्चयं तस्य निश्चयज्ञा दुरात्मनः}
{पुनरेवागमत्क्षिप्रमशोकवनिकां तदा} %||6-25-15||

\twolineshloka
{सा प्रविष्टा पुनस्तत्र ददर्श जनकात्मजाम्}
{प्रतीक्षमाणां स्वामेव भ्रष्टपद्मामिव श्रियम्} %||6-25-16||

\twolineshloka
{तां तु सीता पुनः प्राप्तां सरमां वल्गुभाषिणीम्}
{परिष्वज्य च सुस्निग्धं ददौ च स्वयमासनम्} %||6-25-17||

\twolineshloka
{इहासीना सुखं सर्वमाख्याहि मम तत्त्वतः}
{क्रूरस्य निश्चयं तस्य रावणस्य दुरात्मनः} %||6-25-18||

\twolineshloka
{एवमुक्ता तु सरमा सीतया वेपमानया}
{कथितं सर्वमाचष्ट रावणस्य समन्त्रिणः} %||6-25-19||

\twolineshloka
{जनन्या राक्षसेन्द्रो वै त्वन्मोक्षार्थं बृहद्वचः}
{अविद्धेन च वैदेहि मन्त्रिवृद्धेन बोधितः} %||6-25-20||

\twolineshloka
{दीयतामभिसत्कृत्य मनुजेन्द्राय मैथिली}
{निदर्शनं ते पर्याप्तं जनस्थाने यदद्भुतम्} %||6-25-21||

\twolineshloka
{लङ्घनं च समुद्रस्य दर्शनं च हनूमतः}
{वधं च रक्षसां युद्धे कः कुर्यान्मानुषो भुवि} %||6-25-22||

\twolineshloka
{एवं स मन्त्रिवृद्धैश्च मात्रा च बहु भाषितः}
{न त्वामुत्सहते मोक्तुमर्तह्मर्थपरो यथा} %||6-25-23||

\twolineshloka
{नोत्सहत्यमृतो मोक्तुं युद्धे त्वामिति मैथिलि}
{सामात्यस्य नृशंसस्य निश्चयो ह्येष वर्तते} %||6-25-24||

\threelineshloka
{तदेषा सुस्थिरा बुद्धिर्मृत्युलोभादुपस्थिता}
{भयान्न शक्तस्त्वां मोक्तुमनिरस्तस्तु संयुगे}
{राक्षसानां च सर्वेषामात्मनश्च वधेन हि} %||6-25-25||

\twolineshloka
{निहत्य रावणं सङ्ख्ये सर्वथा निशितैः शरैः}
{प्रतिनेष्यति रामस्त्वामयोध्यामसितेक्षणे} %||6-25-26||

\twolineshloka
{एतस्मिन्नन्तरे शब्दो भेरीशङ्खसमाकुलः}
{श्रुतो वै सर्वसैन्यानां कम्पयन्धरणीतलम्} %||6-25-27||

\fourlineindentedshloka
{श्रुत्वा तु तं वानरसैन्यशब्दम्}
{लङ्कागता राक्षसराजभृत्याः}
{नष्टौजसो दैन्यपरीतचेष्टाः}
{श्रेयो न पश्यन्ति नृपस्य दोषैः} %||6-26-28||


{॥इत्यार्षे श्रीमद्रामायणे वाल्मीकीये आदिकाव्ये युद्धकाण्डे पञ्चविंशः सर्गः॥}

\sect{षड्विंशः सर्गः}

\twolineshloka
{तेन शङ्खविमिश्रेण भेरीशब्देन राघवः}
{उपयतो महाबाहू रामः परपुरंजयः} %||6-26-1||

\twolineshloka
{तं निनादं निशम्याथ रावणो राक्षसेश्वरः}
{मुहूर्तं ध्यानमास्थाय सचिवानभ्युदैक्षत} %||6-26-2||

\twolineshloka
{अथ तान्सचिवांस्तत्र सर्वानाभाष्य रावणः}
{सभां संनादयन्सर्वामित्युवाच महाबलः} %||6-26-3||

\threelineshloka
{तरणं सागरस्यापि विक्रमं बलसञ्चयम्}
{यदुक्तवन्तो रामस्य भवन्तस्तन्मया श्रुतम्}
{भवतश्चाप्यहं वेद्मि युद्धे सत्यपराक्रमान्} %||6-26-4||

\twolineshloka
{ततस्तु सुमहाप्राज्ञो माल्यवान्नाम राक्षसः}
{रावणस्य वचः श्रुत्वा मातुः पैतामहोऽब्रवीत्} %||6-26-5||

\twolineshloka
{विद्यास्वभिविनीतो यो राजा राजन्नयानुगः}
{स शास्ति चिरमैश्वर्यमरींश्च कुरुते वशे} %||6-26-6||

\twolineshloka
{सन्दधानो हि कालेन विगृह्णंश्चारिभिः सह}
{स्वपक्षवर्धनं कुर्वन्महदैश्वर्यमश्नुते} %||6-26-7||

\twolineshloka
{हीयमानेन कर्तव्यो राज्ञा सन्धिः समेन च}
{न शत्रुमवमन्येत ज्यायान्कुर्वीत विग्रहम्} %||6-26-8||

\twolineshloka
{तन्मह्यं रोचते सन्धिः सह रामेण रावण}
{यदर्थमभियुक्ताः स्म सीता तस्मै प्रदीयताम्} %||6-26-9||

\twolineshloka
{तस्य देवर्षयः सर्वे गन्धर्वाश्च जयैषिणः}
{विरोधं मा गमस्तेन सन्धिस्ते तेन रोचताम्} %||6-26-10||

\twolineshloka
{असृजद्भगवान्पक्षौ द्वावेव हि पितामहः}
{सुराणामसुराणां च धर्माधर्मौ तदाश्रयौ} %||6-26-11||

\twolineshloka
{धर्मो हि श्रूयते पक्षः सुराणां च महात्मनाम्}
{अधर्मो रक्षसं पक्षो ह्यसुराणां च रावण} %||6-26-12||

\twolineshloka
{धर्मो वै ग्रसतेऽधर्मं ततः कृतमभूद्युगम्}
{अधर्मो ग्रसते धर्मं ततस्तिष्यः प्रवर्तते} %||6-26-13||

\twolineshloka
{तत्त्वया चरता लोकान्धर्मो विनिहतो महान्}
{अधर्मः प्रगृहीतश्च तेनास्मद्बलिनः परे} %||6-26-14||

\twolineshloka
{स प्रमादाद्विवृद्धस्तेऽधर्मोऽहिर्ग्रसते हि नः}
{विवर्धयति पक्षं च सुराणां सुरभावनः} %||6-26-15||

\threelineshloka
{विषयेषु प्रसक्तेन यत्किञ्चित्कारिणा त्वया}
{ऋषीणामग्निकल्पानामुद्वेगो जनितो महान्}
{तेषां प्रभावो दुर्धर्षः प्रदीप्त इव पावकः} %||6-26-16||

\twolineshloka
{तपसा भावितात्मानो धर्मस्यानुग्रहे रताः}
{मुख्यैर्यज्ञैर्यजन्त्येते नित्यं तैस्तैर्द्विजातयः} %||6-26-17||

\threelineshloka
{जुह्वत्यग्नींश्च विधिवद्वेदांश्चोच्चैरधीयते}
{अभिभूय च रक्षांसि ब्रह्मघोषानुदैरयन्}
{दिशो विप्रद्रुताः सर्वे स्तनयित्नुरिवोष्णगे} %||6-26-18||

\twolineshloka
{ऋषीणामग्निकल्पानामग्निहोत्रसमुत्थितः}
{आदत्ते रक्षसां तेजो धूमो व्याप्य दिशो दश} %||6-26-19||

\twolineshloka
{तेषु तेषु च देशेषु पुण्येषु च दृढव्रतैः}
{चर्यमाणं तपस्तीव्रं सन्तापयति राक्षसान्} %||6-26-20||

\twolineshloka
{उत्पातान्विविधान्दृष्ट्वा घोरान्बहुविधांस्तथा}
{विनाशमनुपश्यामि सर्वेषां रक्षसामहम्} %||6-26-21||

\twolineshloka
{खराभिस्तनिता घोरा मेघाः प्रतिभयङ्करः}
{शोणितेनाभिवर्षन्ति लङ्कामुष्णेन सर्वतः} %||6-26-22||

\twolineshloka
{रुदतां वाहनानां च प्रपतन्त्यस्रबिन्दवः}
{ध्वजा ध्वस्ता विवर्णाश्च न प्रभान्ति यथापुरम्} %||6-26-23||

\twolineshloka
{व्याला गोमायवो गृध्रा वाशन्ति च सुभैरवम्}
{प्रविश्य लङ्कामनिशं समवायांश्च कुर्वते} %||6-26-24||

\twolineshloka
{कालिकाः पाण्डुरैर्दन्तैः प्रहसन्त्यग्रतः स्थिताः}
{स्त्रियः स्वप्नेषु मुष्णन्त्यो गृहाणि प्रतिभाष्य च} %||6-26-25||

\twolineshloka
{गृहाणां बलिकर्माणि श्वानः पर्युपभुञ्जते}
{खरा गोषु प्रजायन्ते मूषिका नकुलैः सह} %||6-26-26||

\twolineshloka
{मार्जारा द्वीपिभिः सार्धं सूकराः शुनकैः सह}
{किंनरा राक्षसैश्चापि समेयुर्मानुषैः सह} %||6-26-27||

\twolineshloka
{पाण्डुरा रक्तपादाश्च विहगाः कालचोदिताः}
{राक्षसानां विनाशाय कपोता विचरन्ति च} %||6-26-28||

\twolineshloka
{चीकी कूचीति वाशन्त्यः शारिका वेश्मसु स्थिताः}
{पतन्ति ग्रथिताश्चापि निर्जिताः कलहैषिणः} %||6-26-29||

\threelineshloka
{करालो विकटो मुण्डः पुरुषः कृष्णपिङ्गलः}
{कालो गृहाणि सर्वेषां काले कालेऽन्ववेक्षते}
{एतान्यन्यानि दुष्टानि निमित्तान्युत्पतन्ति च} %||6-26-30||

\twolineshloka
{विष्णुं मन्यामहे रामं मानुषं देहमास्थितम्}
{न हि मानुषमात्रोऽसौ राघवो दृढविक्रमः} %||6-26-31||

\twolineshloka
{येन बद्धः समुद्रस्य स सेतुः परमाद्भुतः}
{कुरुष्व नरराजेन सन्धिं रामेण रावण} %||6-26-32||

\fourlineindentedshloka
{इदं वचस्तत्र निगद्य माल्यवन्}
{परीक्ष्य रक्षोऽधिपतेर्मनः पुनः}
{अनुत्तमेषूत्तमपौरुषो बली}
{बभूव तूष्णीं समवेक्ष्य रावणम्} %||6-27-33||


{॥इत्यार्षे श्रीमद्रामायणे वाल्मीकीये आदिकाव्ये युद्धकाण्डे षड्विंशः सर्गः॥}

\sect{सप्तविंशः सर्गः}

\twolineshloka
{तत्तु माल्यवतो वाक्यं हितमुक्तं दशाननः}
{न मर्षयति दुष्टात्मा कालस्य वशमागतः} %||6-27-1||

\twolineshloka
{स बद्ध्वा भ्रुकुटिं वक्त्रे क्रोधस्य वशमागतः}
{अमर्षात्परिवृत्ताक्षो माल्यवन्तमथाब्रवीत्} %||6-27-2||

\twolineshloka
{हितबुद्ध्या यदहितं वचः परुषमुच्यते}
{परपक्षं प्रविश्यैव नैतच्छ्रोत्रगतं मम} %||6-27-3||

\twolineshloka
{मानुषं कृपणं राममेकं शाखामृगाश्रयम्}
{समर्थं मन्यसे केन त्यक्तं पित्रा वनालयम्} %||6-27-4||

\twolineshloka
{रक्षसामीश्वरं मां च देवतानां भयङ्करम्}
{हीनं मां मन्यसे केन अहीनं सर्वविक्रमैः} %||6-27-5||

\twolineshloka
{वीरद्वेषेण वा शङ्के पक्षपातेन वा रिपोः}
{त्वयाहं परुषाण्युक्तः परप्रोत्साहनेन वा} %||6-27-6||

\twolineshloka
{प्रभवन्तं पदस्थं हि परुषं कोऽह्बिधास्यति}
{पण्डितः शास्त्रतत्त्वज्ञो विना प्रोत्साहनाद्रिपोः} %||6-27-7||

\twolineshloka
{आनीय च वनात्सीतां पद्महीनामिव श्रियम्}
{किमर्थं प्रतिदास्यामि राघवस्य भयादहम्} %||6-27-8||

\twolineshloka
{वृतं वानरकोटीभिः ससुग्रीवं सलक्ष्मणम्}
{पश्य कैश्चिदहोभिस्त्वं राघवं निहतं मया} %||6-27-9||

\twolineshloka
{द्वन्द्वे यस्य न तिष्ठन्ति दैवतान्यपि संयुगे}
{स कस्माद्रावणो युद्धे भयमाहारयिष्यति} %||6-27-10||

\twolineshloka
{द्विधा भज्येयमप्येवं न नमेयं तु कस्यचित्}
{एष मे सहजो दोषः स्वभावो दुरतिक्रमः} %||6-27-11||

\twolineshloka
{यदि तावत्समुद्रे तु सेतुर्बद्धो यदृच्छया}
{रामेण विस्मयः कोऽत्र येन ते भयमागतम्} %||6-27-12||

\twolineshloka
{स तु तीर्त्वार्णवं रामः सह वानरसेनया}
{प्रतिजानामि ते सत्यं न जीवन्प्रतियास्यति} %||6-27-13||

\twolineshloka
{एवं ब्रुवाणं संरब्धं रुष्टं विज्ञाय रावणम्}
{व्रीडितो माल्यवान्वाक्यं नोत्तरं प्रत्यपद्यत} %||6-27-14||

\twolineshloka
{जयाशिषा च राजानं वर्धयित्वा यथोचितम्}
{माल्यवानभ्यनुज्ञातो जगाम स्वं निवेशनम्} %||6-27-15||

\twolineshloka
{रावणस्तु सहामात्यो मन्त्रयित्वा विमृश्य च}
{लङ्कायामतुलां गुप्तिं कारयामास राक्षसः} %||6-27-16||

\twolineshloka
{व्यादिदेश च पूर्वस्यां प्रहस्तं द्वारि राक्षसम्}
{दक्षिणस्यां महावीर्यौ महापार्श्व महोदरौ} %||6-27-17||

\twolineshloka
{पश्चिमायामथो द्वारि पुत्रमिन्द्रजितं तथा}
{व्यादिदेश महामायं राक्षसैर्बहुभिर्वृतम्} %||6-27-18||

\twolineshloka
{उत्तरस्यां पुरद्वारि व्यादिश्य शुकसारणौ}
{स्वयं चात्र भविष्यामि मन्त्रिणस्तानुवाच ह} %||6-27-19||

\twolineshloka
{राक्षसं तु विरूपाक्षं महावीर्यपराक्रमम्}
{मध्यमेऽस्थापयद्गुल्मे बहुभिः सह राक्षसैः} %||6-27-20||

\twolineshloka
{एवंविधानं लङ्कायां कृत्वा राक्षसपुङ्गवः}
{मेने कृतार्थमात्मानं कृतान्तवशमागतः} %||6-27-21||

\fourlineindentedshloka
{विसर्जयामास ततः स मन्त्रिणो}
{विधानमाज्ञाप्य पुरस्य पुष्कलम्}
{जयाशिषा मन्त्रगणेन पूजितो}
{विवेश सोऽन्तःपुरमृद्धिमन्महत्} %||6-28-22||


{॥इत्यार्षे श्रीमद्रामायणे वाल्मीकीये आदिकाव्ये युद्धकाण्डे सप्तविंशः सर्गः॥}

\sect{अष्टाविंशः सर्गः}

\twolineshloka
{नरवानरराजौ तौ स च वायुसुतः कपिः}
{जाम्बवानृक्षराजश्च राक्षसश्च विभीषणः} %||6-28-1||

\twolineshloka
{अङ्गदो वालिपुत्रश्च सौमित्रिः शरभः कपिः}
{सुषेणः सहदायादो मैन्दो द्विविद एव च} %||6-28-2||

\twolineshloka
{गजो गवाक्षो कुमुदो नलोऽथ पनसस्तथा}
{अमित्रविषयं प्राप्ताः समवेताः समर्थयन्} %||6-28-3||

\twolineshloka
{इयं सा लक्ष्यते लङ्का पुरी रावणपालिता}
{सासुरोरगगन्धर्वैरमरैरपि दुर्जया} %||6-28-4||

\twolineshloka
{कार्यसिद्धिं पुरस्कृत्य मन्त्रयध्वं विनिर्णये}
{नित्यं संनिहितो ह्यत्र रावणो राक्षसाधिपः} %||6-28-5||

\twolineshloka
{तथा तेषु ब्रुवाणेषु रावणावरजोऽब्रवीत्}
{वाक्यमग्राम्यपदवत्पुष्कलार्थं विभीषणः} %||6-28-6||

\twolineshloka
{अनलः शरभश्चैव सम्पातिः प्रघसस्तथा}
{गत्वा लङ्कां ममामात्याः पुरीं पुनरिहागताः} %||6-28-7||

\twolineshloka
{भूत्वा शकुनयः सर्वे प्रविष्टाश्च रिपोर्बलम्}
{विधानं विहितं यच्च तद्दृष्ट्वा समुपस्थिताः} %||6-28-8||

\twolineshloka
{संविधानं यथाहुस्ते रावणस्य दुरात्मनः}
{राम तद्ब्रुवतः सर्वं यथातथ्येन मे शृणु} %||6-28-9||

\twolineshloka
{पूर्वं प्रहस्तः सबलो द्वारमासाद्य तिष्ठति}
{दक्षिणं च महावीर्यौ महापार्श्वमहोदरौ} %||6-28-10||

\twolineshloka
{इन्द्रजित्पश्चिमद्वारं राक्षसैर्बहुभिर्वृतः}
{पट्टसासिधनुष्मद्भिः शूलमुद्गरपाणिभिः} %||6-28-11||

\twolineshloka
{नानाप्रहरणैः शूरैरावृतो रावणात्मजः}
{राक्षसानां सहस्रैस्तु बहुभिः शस्त्रपाणिभिः} %||6-28-12||

\twolineshloka
{युक्तः परमसंविग्नो राक्षसैर्बहुभिर्वृतः}
{उत्तरं नगरद्वारं रावणः स्वयमास्थितः} %||6-28-13||

\twolineshloka
{विरूपाक्षस्तु महता शूलखड्गधनुष्मता}
{बलेन राक्षसैः सार्धं मध्यमं गुल्ममास्थितः} %||6-28-14||

\twolineshloka
{एतानेवंविधान्गुल्माँल्लङ्कायां समुदीक्ष्य ते}
{मामकाः सचिवाः सर्वे शीघ्रं पुनरिहागताः} %||6-28-15||

\twolineshloka
{गजानां च सहस्रं च रथानामयुतं पुरे}
{हयानामयुते द्वे च साग्रकोटी च रक्षसाम्} %||6-28-16||

\twolineshloka
{विक्रान्ता बलवन्तश्च संयुगेष्वाततायिनः}
{इष्टा राक्षसराजस्य नित्यमेते निशाचराः} %||6-28-17||

\twolineshloka
{एकैकस्यात्र युद्धार्थे राक्षसस्य विशां पते}
{परिवारः सहस्राणां सहस्रमुपतिष्ठते} %||6-28-18||

\twolineshloka
{एतां प्रवृत्तिं लङ्कायां मन्त्रिप्रोक्तं विभीषणः}
{रामं कमलपत्राक्षमिदमुत्तरमब्रवीत्} %||6-28-19||

\twolineshloka
{कुबेरं तु यदा राम रावणः प्रत्ययुध्यत}
{षष्टिः शतसहस्राणि तदा निर्यान्ति राक्षसाः} %||6-28-20||

\twolineshloka
{पराक्रमेण वीर्येण तेजसा सत्त्वगौरवात्}
{सदृशा योऽत्र दर्पेण रावणस्य दुरात्मनः} %||6-28-21||

\twolineshloka
{अत्र मन्युर्न कर्तव्यो रोषये त्वां न भीषये}
{समर्थो ह्यसि वीर्येण सुराणामपि निग्रहे} %||6-28-22||

\twolineshloka
{तद्भवांश्चतुरङ्गेण बलेन महता वृतः}
{व्यूह्येदं वानरानीकं निर्मथिष्यसि रावणम्} %||6-28-23||

\twolineshloka
{रावणावरजे वाक्यमेवं ब्रुवति राघवः}
{शत्रूणां प्रतिघातार्थमिदं वचनमब्रवीत्} %||6-28-24||

\twolineshloka
{पूर्वद्वारे तु लङ्काया नीलो वानरपुङ्गवः}
{प्रहस्तं प्रतियोद्धा स्याद्वानरैर्बहुभिर्वृतः} %||6-28-25||

\twolineshloka
{अङ्गदो वालिपुत्रस्तु बलेन महता वृतः}
{दक्षिणे बाधतां द्वारे महापार्श्वमहोदरौ} %||6-28-26||

\twolineshloka
{हनूमान्पश्चिमद्वारं निपीड्य पवनात्मजः}
{प्रविशत्वप्रमेयात्मा बहुभिः कपिभिर्वृतः} %||6-28-27||

\twolineshloka
{दैत्यदानवसङ्घानामृषीणां च महात्मनाम्}
{विप्रकारप्रियः क्षुद्रो वरदानबलान्वितः} %||6-28-28||

\twolineshloka
{परिक्रामति यः सर्वाँल्लोकान्सन्तापयन्प्रजाः}
{तस्याहं राक्षसेन्द्रस्य स्वयमेव वधे धृतः} %||6-28-29||

\twolineshloka
{उत्तरं नगरद्वारमहं सौमित्रिणा सह}
{निपीड्याभिप्रवेक्ष्यामि सबलो यत्र रावणः} %||6-28-30||

\twolineshloka
{वानरेन्द्रश्च बलवानृक्षराजश्च जाम्बवान्}
{राक्षसेन्द्रानुजश्चैव गुल्मे भवतु मध्यमे} %||6-28-31||

\twolineshloka
{न चैव मानुषं रूपं कार्यं हरिभिराहवे}
{एषा भवतु नः संज्ञा युद्धेऽस्मिन्वानरे बले} %||6-28-32||

\twolineshloka
{वानरा एव निश्चिह्नं स्वजनेऽस्मिन्भविष्यति}
{वयं तु मानुषेणैव सप्त योत्स्यामहे परान्} %||6-28-33||

\twolineshloka
{अहमेव सह भ्रात्रा लक्ष्मणेन महौजसा}
{आत्मना पञ्चमश्चायं सखा मम विभीषणः} %||6-28-34||

\twolineshloka
{स रामः कार्यसिद्ध्यर्थमेवमुक्त्वा विभीषणम्}
{सुवेलारोहणे बुद्धिं चकार मतिमान्मतिम्} %||6-28-35||

\fourlineindentedshloka
{ततस्तु रामो महता बलेन}
{प्रच्छाद्य सर्वां पृथिवीं महात्मा}
{प्रहृष्टरूपोऽभिजगाम लङ्काम्}
{कृत्वा मतिं सोऽरिवधे महात्मा} %||6-29-36||


{॥इत्यार्षे श्रीमद्रामायणे वाल्मीकीये आदिकाव्ये युद्धकाण्डे अष्टाविंशः सर्गः॥}

\sect{एकोनत्रिंशः सर्गः}

\twolineshloka
{स तु कृत्वा सुवेलस्य मतिमारोहणं प्रति}
{लक्ष्मणानुगतो रामः सुग्रीवमिदमब्रवीत्} %||6-29-1||

\twolineshloka
{विभीषणं च धर्मज्ञमनुरक्तं निशाचरम्}
{मन्त्रज्ञं च विधिज्ञं च श्लक्ष्णया परया गिरा} %||6-29-2||

\twolineshloka
{सुवेलं साधु शैलेन्द्रमिमं धातुशतैश्चितम्}
{अध्यारोहामहे सर्वे वत्स्यामोऽत्र निशामिमाम्} %||6-29-3||

\twolineshloka
{लङ्कां चालोकयिष्यामो निलयं तस्य रक्षसः}
{येन मे मरणान्ताय हृता भार्या दुरात्मना} %||6-29-4||

\twolineshloka
{येन धर्मो न विज्ञातो न वृत्तं न कुलं तथा}
{राक्षस्या नीचया बुद्ध्या येन तद्गर्हितं कृतम्} %||6-29-5||

\twolineshloka
{यस्मिन्मे वर्धते रोषः कीर्तिते राक्षसाधमे}
{यस्यापराधान्नीचस्य वधं द्रक्ष्यामि रक्षसाम्} %||6-29-6||

\twolineshloka
{एको हि कुरुते पापं कालपाशवशं गतः}
{नीचेनात्मापचारेण कुलं तेन विनश्यति} %||6-29-7||

\twolineshloka
{एवं सम्मन्त्रयन्नेव सक्रोधो रावणं प्रति}
{रामः सुवेलं वासाय चित्रसानुमुपारुहत्} %||6-29-8||

\twolineshloka
{पृष्ठतो लक्ष्मण चैनमन्वगच्छत्समाहितः}
{सशरं चापमुद्यम्य सुमहद्विक्रमे रतः} %||6-29-9||

\twolineshloka
{तमन्वरोहत्सुग्रीवः सामात्यः सविभीषणः}
{हनूमानङ्गदो नीलो मैन्दो द्विविद एव च} %||6-29-10||

\twolineshloka
{गजो गवाक्षो गवयः शरभो गन्धमादनः}
{पनसः कुमुदश्चैव हरो रम्भश्च यूथपः} %||6-29-11||

\threelineshloka
{एते चान्ये च बहवो वानराः शीघ्रगामिनः}
{ते वायुवेगप्रवणास्तं गिरिं गिरिचारिणः}
{अध्यारोहन्त शतशः सुवेलं यत्र राघवः} %||6-29-12||

\twolineshloka
{ते त्वदीर्घेण कालेन गिरिमारुह्य सर्वतः}
{ददृशुः शिखरे तस्य विषक्तामिव खे पुरीम्} %||6-29-13||

\twolineshloka
{तां शुभां प्रवरद्वारां प्राकारवरशोभिताम्}
{लङ्कां राक्षससम्पूर्णां ददृशुर्हरियूथपाः} %||6-29-14||

\twolineshloka
{प्राकारचयसंस्थैश्च तथा नीलैर्निशाचरैः}
{ददृशुस्ते हरिश्रेष्ठाः प्राकारमपरं कृतम्} %||6-29-15||

\twolineshloka
{ते दृष्ट्वा वानराः सर्वे राक्षसान्युद्धकाङ्क्षिणः}
{मुमुचुर्विपुलान्नादांस्तत्र रामस्य पश्यतः} %||6-29-16||

\twolineshloka
{ततोऽस्तमगमत्सूर्यः सन्ध्यया प्रतिरञ्जितः}
{पूर्णचन्द्रप्रदीपा च क्षपा समभिवर्तते} %||6-29-17||

\fourlineindentedshloka
{ततः स रामो हरिवाहिनीपतिर्-}
{विभीषणेन प्रतिनन्द्य सत्कृतः}
{सलक्ष्मणो यूथपयूथसंवृतः}
{सुवेल पृष्ठे न्यवसद्यथासुखम्} %||6-30-18||


{॥इत्यार्षे श्रीमद्रामायणे वाल्मीकीये आदिकाव्ये युद्धकाण्डे एकोनत्रिंशः सर्गः॥}

\sect{त्रिंशः सर्गः}

\twolineshloka
{तां रात्रिमुषितास्तत्र सुवेले हरिपुङ्गवाः}
{लङ्कायां ददृशुर्वीरा वनान्युपवनानि च} %||6-30-1||

\twolineshloka
{समसौम्यानि रम्याणि विशालान्यायतानि च}
{दृष्टिरम्याणि ते दृष्ट्वा बभूवुर्जातविस्मयाः} %||6-30-2||

\twolineshloka
{चम्पकाशोकपुंनागसालतालसमाकुला}
{तमालवनसञ्छन्ना नागमालासमावृता} %||6-30-3||

\twolineshloka
{हिन्तालैरर्जुनैर्नीपैः सप्तपर्णैश्च पुष्पितैः}
{तिलकैः कर्णिकारैश्च पटालैश्च समन्ततः} %||6-30-4||

\twolineshloka
{शुशुभे पुष्पिताग्रैश्च लतापरिगतैर्द्रुमैः}
{लङ्का बहुविधैर्दिव्यैर्यथेन्द्रस्यामरावती} %||6-30-5||

\twolineshloka
{विचित्रकुसुमोपेतै रक्तकोमलपल्लवैः}
{शाद्वलैश्च तथा नीलैश्चित्राभिर्वनराजिभिः} %||6-30-6||

\twolineshloka
{गन्धाढ्यान्यभिरम्याणि पुष्पाणि च फलानि च}
{धारयन्त्यगमास्तत्र भूषणानीव मानवाः} %||6-30-7||

\twolineshloka
{तच्चैत्ररथसङ्काशं मनोज्ञं नन्दनोपमम्}
{वनं सर्वर्तुकं रम्यं शुशुभे षट्पदायुतम्} %||6-30-8||

\twolineshloka
{नत्यूहकोयष्टिभकैर्नृत्यमानैश्च बर्हिभिः}
{रुतं परभृतानां च शुश्रुवे वननिर्झरे} %||6-30-9||

\twolineshloka
{नित्यमत्तविहङ्गानि भ्रमराचरितानि च}
{कोकिलाकुलषण्डानि विहगाभिरुतानि च} %||6-30-10||

\twolineshloka
{भृङ्गराजाभिगीतानि भ्रमरैः सेवितानि च}
{कोणालकविघुष्टानि सारसाभिरुतानि च} %||6-30-11||

\twolineshloka
{विविशुस्ते ततस्तानि वनान्युपवनानि च}
{हृष्टाः प्रमुदिता वीरा हरयः कामरूपिणः} %||6-30-12||

\twolineshloka
{तेषां प्रविशतां तत्र वानराणां महौजसाम्}
{पुष्पसंसर्गसुरभिर्ववौ घ्राणसुखोऽनिलः} %||6-30-13||

\twolineshloka
{अन्ये तु हरिवीराणां यूथान्निष्क्रम्य यूथपाः}
{सुग्रीवेणाभ्यनुज्ञाता लङ्कां जग्मुः पताकिनीम्} %||6-30-14||

\twolineshloka
{वित्रासयन्तो विहगांस्त्रासयन्तो मृगद्विपान्}
{कम्पयन्तश्च तां लङ्कां नादैः स्वैर्नदतां वराः} %||6-30-15||

\twolineshloka
{कुर्वन्तस्ते महावेगा महीं चारणपीडिताम्}
{रजश्च सहसैवोर्ध्वं जगाम चरणोद्धतम्} %||6-30-16||

\twolineshloka
{ऋक्षाः सिंहा वराहाश्च महिषा वारणा मृगाः}
{तेन शब्देन वित्रस्ता जग्मुर्भीता दिशो दश} %||6-30-17||

\twolineshloka
{शिखरं तु त्रिकूटस्य प्रांशु चैकं दिविस्पृशम्}
{समन्तात्पुष्पसञ्छन्नं महारजतसंनिभम्} %||6-30-18||

\twolineshloka
{शतयोजनविस्तीर्णं विमलं चारुदर्शनम्}
{श्लक्ष्णं श्रीमन्महच्चैव दुष्प्रापं शकुनैरपि} %||6-30-19||

\twolineshloka
{मनसापि दुरारोहं किं पुनः कर्मणा जनैः}
{निविष्टा तत्र शिखरे लङ्का रावणपालिता} %||6-30-20||

\twolineshloka
{सा पुरी गोपुरैरुच्चैः पाण्डुराम्बुदसंनिभैः}
{काञ्चनेन च सालेन राजतेन च शोभिता} %||6-30-21||

\twolineshloka
{प्रासादैश्च विमानैश्च लङ्का परमभूषिता}
{घनैरिवातपापाये मध्यमं वैष्णवं पदम्} %||6-30-22||

\twolineshloka
{यस्यां स्तम्भसहस्रेण प्रासादः समलङ्कृतः}
{कैलासशिखराकारो दृश्यते खमिवोल्लिखन्} %||6-30-23||

\twolineshloka
{चैत्यः स राक्षसेन्द्रस्य बभूव पुरभूषणम्}
{शतेन रक्षसां नित्यं यः समग्रेण रक्ष्यते} %||6-30-24||

\twolineshloka
{तां समृद्धां समृद्धार्थो लक्ष्मीवाँल्लक्ष्मणाग्रजः}
{रावणस्य पुरीं रामो ददर्श सह वानरैः} %||6-30-25||

\fourlineindentedshloka
{तां रत्नपूर्णां बहुसंविधानाम्}
{प्रासादमालाभिरलङ्कृतां च}
{पुरीं महायन्त्रकवाटमुख्याम्}
{ददर्श रामो महता बलेन} %||6-31-26||


{॥इत्यार्षे श्रीमद्रामायणे वाल्मीकीये आदिकाव्ये युद्धकाण्डे त्रिंशः सर्गः॥}

\sect{एकत्रिंशः सर्गः}

\twolineshloka
{अथ तस्मिन्निमित्तानि दृष्ट्वा लक्ष्मणपूर्वजः}
{लक्ष्मणं लक्ष्मिसम्पन्नमिदं वचनमब्रवीत्} %||6-31-1||

\twolineshloka
{परिगृह्योदकं शीतं वनानि फलवन्ति च}
{बलौघं संविभज्येमं व्यूह्य तिष्ठेम लक्ष्मण} %||6-31-2||

\twolineshloka
{लोकक्षयकरं भीमं भयं पश्याम्युपस्थितम्}
{निबर्हणं प्रवीराणामृक्षवानररक्षसाम्} %||6-31-3||

\twolineshloka
{वाताश्च परुषं वान्ति कम्पते च वसुन्धरा}
{पर्वताग्राणि वेपन्ते पतन्ति धरणीधराः} %||6-31-4||

\twolineshloka
{मेघाः क्रव्यादसङ्काशाः परुषाः परुषस्वनाः}
{क्रूराः क्रूरं प्रवर्षन्ति मिश्रं शोणितबिन्दुभिः} %||6-31-5||

\twolineshloka
{रक्तचन्दनसङ्काशा सन्ध्यापरमदारुणा}
{ज्वलच्च निपतत्येतदादित्यादग्निमण्डलम्} %||6-31-6||

\twolineshloka
{आदित्यमभिवाश्यन्ते जनयन्तो महद्भयम्}
{दीना दीनस्वरा घोरा अप्रशस्ता मृगद्विजाः} %||6-31-7||

\twolineshloka
{रजन्यामप्रकाशश्च सन्तापयति चन्द्रमाः}
{कृष्णरक्तांशुपर्यन्तो यथा लोकस्य सङ्क्षये} %||6-31-8||

\twolineshloka
{ह्रस्वो रूक्षोऽप्रशस्तश्च परिवेषः सुलोहितः}
{आदित्यमण्डले नीलं लक्ष्म लक्ष्मण दृश्यते} %||6-31-9||

\twolineshloka
{दृश्यन्ते न यथावच्च नक्षत्राण्यभिवर्तते}
{युगान्तमिव लोकस्य पश्य लक्ष्मण शंसति} %||6-31-10||

\twolineshloka
{काकाः श्येनास्तथा गृध्रा नीचैः परिपतन्ति च}
{शिवाश्चाप्यशिवा वाचः प्रवदन्ति महास्वनाः} %||6-31-11||

\twolineshloka
{क्षिप्रमद्य दुराधर्षां पुरीं रावणपालिताम्}
{अभियाम जवेनैव सर्वतो हरिभिर्वृताः} %||6-31-12||

\twolineshloka
{इत्येवं तु वदन्वीरो लक्ष्मणं लक्ष्मणाग्रजः}
{तस्मादवातरच्छीघ्रं पर्वताग्रान्महाबलः} %||6-31-13||

\twolineshloka
{अवतीर्य तु धर्मात्मा तस्माच्छैलात्स राघवः}
{परैः परमदुर्धर्षं ददर्श बलमात्मनः} %||6-31-14||

\twolineshloka
{संनह्य तु ससुग्रीवः कपिराजबलं महत्}
{कालज्ञो राघवः काले संयुगायाभ्यचोदयत्} %||6-31-15||

\twolineshloka
{ततः काले महाबाहुर्बलेन महता वृतः}
{प्रस्थितः पुरतो धन्वी लङ्कामभिमुखः पुरीम्} %||6-31-16||

\twolineshloka
{तं विभीषण सुग्रीवौ हनूमाञ्जाम्बवान्नलः}
{ऋक्षराजस्तथा नीलो लक्ष्मणश्चान्ययुस्तदा} %||6-31-17||

\twolineshloka
{ततः पश्चात्सुमहती पृतनर्क्षवनौकसाम्}
{प्रच्छाद्य महतीं भूमिमनुयाति स्म राघवम्} %||6-31-18||

\twolineshloka
{शैलशृङ्गाणि शतशः प्रवृद्धांश्च महीरुहाम्}
{जगृहुः कुञ्जरप्रख्या वानराः परवारणाः} %||6-31-19||

\twolineshloka
{तौ त्वदीर्घेण कालेन भ्रातरौ रामलक्ष्मणौ}
{रावणस्य पुरीं लङ्कामासेदतुररिन्दमौ} %||6-31-20||

\twolineshloka
{पताकामालिनीं रम्यामुद्यानवनशोभिताम्}
{चित्रवप्रां सुदुष्प्रापामुच्चप्राकारतोरणाम्} %||6-31-21||

\twolineshloka
{तां सुरैरपि दुर्धर्षां रामवाक्यप्रचोदिताः}
{यथानिदेशं सम्पीड्य न्यविशन्त वनौकसः} %||6-31-22||

\twolineshloka
{लङ्कायास्तूत्तरद्वारं शैलशृङ्गमिवोन्नतम्}
{रामः सहानुजो धन्वी जुगोप च रुरोध च} %||6-31-23||

\twolineshloka
{लङ्कामुपनिविष्टश्च रामो दशरथात्मजः}
{लक्ष्मणानुचरो वीरः पुरीं रावणपालिताम्} %||6-31-24||

\twolineshloka
{उत्तरद्वारमासाद्य यत्र तिष्ठति रावणः}
{नान्यो रामाद्धि तद्द्वारं समर्थः परिरक्षितुम्} %||6-31-25||

\threelineshloka
{रावणाधिष्ठितं भीमं वरुणेनेव सागरम्}
{सायुधौ राक्षसैर्भीमैरभिगुप्तं समन्ततः}
{लघूनां त्रासजननं पातालमिव दानवैः} %||6-31-26||

\twolineshloka
{विन्यस्तानि च योधानां बहूनि विविधानि च}
{ददर्शायुधजालानि तथैव कवचानि च} %||6-31-27||

\twolineshloka
{पूर्वं तु द्वारमासाद्य नीलो हरिचमूपतिः}
{अतिष्ठत्सह मैन्देन द्विविदेन च वीर्यवान्} %||6-31-28||

\twolineshloka
{अङ्गदो दक्षिणद्वारं जग्राह सुमहाबलः}
{ऋषभेण गवाक्षेण गजेन गवयेन च} %||6-31-29||

\twolineshloka
{हनूमान्पश्चिमद्वारं ररक्ष बलवान्कपिः}
{प्रमाथि प्रघसाभ्यां च वीरैरन्यैश्च सङ्गतः} %||6-31-30||

\twolineshloka
{मध्यमे च स्वयं गुल्मे सुग्रीवः समतिष्ठत}
{सह सर्वैर्हरिश्रेष्ठैः सुपर्णश्वसनोपमैः} %||6-31-31||

\twolineshloka
{वानराणां तु षट्त्रिंशत्कोट्यः प्रख्यातयूथपाः}
{निपीड्योपनिविष्टाश्च सुग्रीवो यत्र वानरः} %||6-31-32||

\twolineshloka
{शासनेन तु रामस्य लक्ष्मणः सविभीषणः}
{द्वारे द्वारे हरीणां तु कोटिं कोटिं न्यवेशयत्} %||6-31-33||

\twolineshloka
{पश्चिमेन तु रामस्य सुग्रीवः सह जाम्बवान्}
{अदूरान्मध्यमे गुल्मे तस्थौ बहुबलानुगः} %||6-31-34||

\twolineshloka
{ते तु वानरशार्दूलाः शार्दूला इव दंष्ट्रिणः}
{गृहीत्वा द्रुमशैलाग्रान्हृष्टा युद्धाय तस्थिरे} %||6-31-35||

\twolineshloka
{सर्वे विकृतलाङ्गूलाः सर्वे दंष्ट्रानखायुधाः}
{सर्वे विकृतचित्राङ्गाः सर्वे च विकृताननाः} %||6-31-36||

\twolineshloka
{दशनागबलाः केचित्केचिद्दशगुणोत्तराः}
{केचिन्नागसहस्रस्य बभूवुस्तुल्यविक्रमाः} %||6-31-37||

\twolineshloka
{सन्ति चौघा बलाः केचित्केचिच्छतगुणोत्तराः}
{अप्रमेयबलाश्चान्ये तत्रासन्हरियूथपाः} %||6-31-38||

\twolineshloka
{अद्भुतश्च विचित्रश्च तेषामासीत्समागमः}
{तत्र वानरसैन्यानां शलभानामिवोद्गमः} %||6-31-39||

\twolineshloka
{परिपूर्णमिवाकाशं सञ्छन्नेव च मेदिनी}
{लङ्कामुपनिविष्टैश्च सम्पतद्भिश्च वानरैः} %||6-31-40||

\twolineshloka
{शतं शतसहस्राणां पृथगृक्षवनौकसाम्}
{लङ्का द्वाराण्युपाजग्मुरन्ये योद्धुं समन्ततः} %||6-31-41||

\twolineshloka
{आवृतः स गिरिः सर्वैस्तैः समन्तात्प्लवङ्गमैः}
{अयुतानां सहस्रं च पुरीं तामभ्यवर्तत} %||6-31-42||

\twolineshloka
{वानरैर्बलवद्भिश्च बभूव द्रुमपाणिभिः}
{सर्वतः संवृता लङ्का दुष्प्रवेशापि वायुना} %||6-31-43||

\twolineshloka
{राक्षसा विस्मयं जग्मुः सहसाभिनिपीडिताः}
{वानरैर्मेघसङ्काशैः शक्रतुल्यपराक्रमैः} %||6-31-44||

\twolineshloka
{महाञ्शब्दोऽभवत्तत्र बलौघस्याभिवर्ततः}
{सागरस्येव भिन्नस्य यथा स्यात्सलिलस्वनः} %||6-31-45||

\twolineshloka
{तेन शब्देन महता सप्राकारा सतोरणा}
{लङ्का प्रचलिता सर्वा सशैलवनकानना} %||6-31-46||

\twolineshloka
{रामलक्ष्मणगुप्ता सा सुग्रीवेण च वाहिनी}
{बभूव दुर्धर्षतरा सर्वैरपि सुरासुरैः} %||6-31-47||

\twolineshloka
{राघवः संनिवेश्यैवं सैन्यं स्वं रक्षसां वधे}
{सम्मन्त्र्य मन्त्रिभिः सार्धं निश्चित्य च पुनः पुनः} %||6-31-48||

\threelineshloka
{आनन्तर्यमभिप्रेप्सुः क्रमयोगार्थतत्त्ववित्}
{विभीषणस्यानुमते राजधर्ममनुस्मरन्}
{अङ्गदं वालितनयं समाहूयेदमब्रवीत्} %||6-31-49||

\twolineshloka
{गत्वा सौम्य दशग्रीवं ब्रूहि मद्वचनात्कपे}
{लङ्घयित्वा पुरीं लङ्कां भयं त्यक्त्वा गतव्यथः} %||6-31-50||

\twolineshloka
{भ्रष्टश्रीकगतैश्वर्यमुमूर्षो नष्टचेतनः}
{ऋषीणां देवतानां च गन्धर्वाप्सरसां तथा} %||6-31-51||

\twolineshloka
{नागानामथ यक्षाणां राज्ञां च रजनीचर}
{यच्च पापं कृतं मोहादवलिप्तेन राक्षस} %||6-31-52||

\threelineshloka
{नूनमद्य गतो दर्पः स्वयम्भू वरदानजः}
{यस्य दण्डधरस्तेऽहं दाराहरणकर्शितः}
{दण्डं धारयमाणस्तु लङ्काद्वरे व्यवस्थितः} %||6-31-53||

\twolineshloka
{पदवीं देवतानां च महर्षीणां च राक्षस}
{राजर्षीणां च सर्वेणां गमिष्यसि मया हतः} %||6-31-54||

\twolineshloka
{बलेन येन वै सीतां मायया राक्षसाधम}
{मामतिक्रामयित्वा त्वं हृतवांस्तद्विदर्शय} %||6-31-55||

\twolineshloka
{अराक्षसमिमं लोकं कर्तास्मि निशितैः शरैः}
{न चेच्छरणमभ्येषि मामुपादाय मैथिलीम्} %||6-31-56||

\twolineshloka
{धर्मात्मा रक्षसां श्रेष्ठः सम्प्राप्तोऽयं विभीषणः}
{लङ्कैश्वर्यं ध्रुवं श्रीमानयं प्राप्नोत्यकण्टकम्} %||6-31-57||

\twolineshloka
{न हि राज्यमधर्मेण भोक्तुं क्षणमपि त्वया}
{शक्यं मूर्खसहायेन पापेनाविजितात्मना} %||6-31-58||

\twolineshloka
{युध्यस्व वा धृतिं कृत्वा शौर्यमालम्ब्य राक्षस}
{मच्छरैस्त्वं रणे शान्तस्ततः पूतो भविष्यसि} %||6-31-59||

\twolineshloka
{यद्याविशसि लोकांस्त्रीन्पक्षिभूतो मनोजवः}
{मम चक्षुष्पथं प्राप्य न जीवन्प्रतियास्यसि} %||6-31-60||

\twolineshloka
{ब्रवीमि त्वां हितं वाक्यं क्रियतामौर्ध्वदेकिकम्}
{सुदृष्टा क्रियतां लङ्का जीवितं ते मयि स्थितम्} %||6-31-61||

\twolineshloka
{इत्युक्तः स तु तारेयो रामेणाक्लिष्टकर्मणा}
{जगामाकाशमाविश्य मूर्तिमानिव हव्यवाट्} %||6-31-62||

\twolineshloka
{सोऽतिपत्य मुहूर्तेन श्रीमान्रावणमन्दिरम्}
{ददर्शासीनमव्यग्रं रावणं सचिवैः सह} %||6-31-63||

\twolineshloka
{ततस्तस्याविदूरेण निपत्य हरिपुङ्गवः}
{दीप्ताग्निसदृशस्तस्थावङ्गदः कनकाङ्गदः} %||6-31-64||

\twolineshloka
{तद्रामवचनं सर्वमन्यूनाधिकमुत्तमम्}
{सामात्यं श्रावयामास निवेद्यात्मानमात्मना} %||6-31-65||

\twolineshloka
{दूतोऽहं कोसलेन्द्रस्य रामस्याक्लिष्टकर्मणः}
{वालिपुत्रोऽङ्गदो नाम यदि ते श्रोत्रमागतः} %||6-31-66||

\twolineshloka
{आह त्वां राघवो रामः कौसल्यानन्दवर्धनः}
{निष्पत्य प्रतियुध्यस्व नृशंसं पुरुषाधम} %||6-31-67||

\twolineshloka
{हन्तास्मि त्वां सहामात्यं सपुत्रज्ञातिबान्धवम्}
{निरुद्विग्नास्त्रयो लोका भविष्यन्ति हते त्वयि} %||6-31-68||

\twolineshloka
{देवदानवयक्षाणां गन्धर्वोरगरक्षसाम्}
{शत्रुमद्योद्धरिष्यामि त्वामृषीणां च कण्टकम्} %||6-31-69||

\twolineshloka
{विभीषणस्य चैश्वर्यं भविष्यति हते त्वयि}
{न चेत्सत्कृत्य वैदेहीं प्रणिपत्य प्रदास्यसि} %||6-31-70||

\twolineshloka
{इत्येवं परुषं वाक्यं ब्रुवाणे हरिपुङ्गवे}
{अमर्षवशमापन्नो निशाचरगणेश्वरः} %||6-31-71||

\twolineshloka
{ततः स रोषताम्राक्षः शशास सचिवांस्तदा}
{गृह्यतामेष दुर्मेधा वध्यतामिति चासकृत्} %||6-31-72||

\twolineshloka
{रावणस्य वचः श्रुत्वा दीप्ताग्निसमतेजसः}
{जगृहुस्तं ततो घोराश्चत्वारो रजनीचराः} %||6-31-73||

\twolineshloka
{ग्राहयामास तारेयः स्वयमात्मानमात्मना}
{बलं दर्शयितुं वीरो यातुधानगणे तदा} %||6-31-74||

\twolineshloka
{स तान्बाहुद्वये सक्तानादाय पतगानिव}
{प्रासादं शैलसङ्काशमुत्पापाताङ्गदस्तदा} %||6-31-75||

\twolineshloka
{तेऽन्तरिक्षाद्विनिर्धूतास्तस्य वेगेन राक्षसाः}
{भुमौ निपतिताः सर्वे राक्षसेन्द्रस्य पश्यतः} %||6-31-76||

\twolineshloka
{ततः प्रासादशिखरं शैलशृङ्गमिवोन्नतम्}
{तत्पफाल तदाक्रान्तं दशग्रीवस्य पश्यतः} %||6-31-77||

\twolineshloka
{भङ्क्त्वा प्रासादशिखरं नाम विश्राव्य चात्मनः}
{विनद्य सुमहानादमुत्पपात विहायसा} %||6-31-78||

\twolineshloka
{रावणस्तु परं चक्रे क्रोधं प्रासादधर्षणात्}
{विनाशं चात्मनः पश्यन्निःश्वासपरमोऽभवत्} %||6-31-79||

\twolineshloka
{रामस्तु बहुभिर्हृष्टैर्निनदद्भिः प्लवङ्गमैः}
{वृतो रिपुवधाकाङ्क्षी युद्धायैवाभ्यवर्तत} %||6-31-80||

\twolineshloka
{सुषेणस्तु महावीर्यो गिरिकूटोपमो हरिः}
{बहुभिः संवृतस्तत्र वानरैः कामरूपिभिः} %||6-31-81||

\twolineshloka
{चतुर्द्वाराणि सर्वाणि सुग्रीववचनात्कपिः}
{पर्याक्रमत दुर्धर्षो नक्षत्राणीव चन्द्रमाः} %||6-31-82||

\twolineshloka
{तेषामक्षौहिणिशतं समवेक्ष्य वनौकसाम्}
{लङ्कामुपनिविष्टानां सागरं चातिवर्तताम्} %||6-31-83||

\twolineshloka
{राक्षसा विस्मयं जग्मुस्त्रासं जग्मुस्तथापरे}
{अपरे समरोद्धर्षाद्धर्षमेवोपपेदिरे} %||6-31-84||

\twolineshloka
{कृत्स्नं हि कपिभिर्व्याप्तं प्राकारपरिखान्तरम्}
{ददृशू राक्षसा दीनाः प्राकारं वानरीकृतम्} %||6-31-85||

\fourlineindentedshloka
{तस्मिन्महाभीषणके प्रवृत्ते}
{कोलाहले राक्षसराजधान्याम्}
{प्रगृह्य रक्षांसि महायुधानि}
{युगान्तवाता इव संविचेरुः} %||6-32-86||


{॥इत्यार्षे श्रीमद्रामायणे वाल्मीकीये आदिकाव्ये युद्धकाण्डे एकत्रिंशः सर्गः॥}

\sect{द्वात्रिंशः सर्गः}

\twolineshloka
{ततस्ते राक्षसास्तत्र गत्वा रावणमन्दिरम्}
{न्यवेदयन्पुरीं रुद्धां रामेण सह वानरैः} %||6-32-1||

\twolineshloka
{रुद्धां तु नगरीं श्रुत्वा जातक्रोधो निशाचरः}
{विधानं द्विगुणं श्रुत्वा प्रासादं सोऽध्यरोहत} %||6-32-2||

\twolineshloka
{स ददर्शावृतां लङ्कां सशैलवनकाननाम्}
{असङ्ख्येयैर्हरिगणैः सर्वतो युद्धकाङ्क्षिभिः} %||6-32-3||

\twolineshloka
{स दृष्ट्वा वानरैः सर्वां वसुधां कवलीकृताम्}
{कथं क्षपयितव्याः स्युरिति चिन्तापरोऽभवत्} %||6-32-4||

\twolineshloka
{स चिन्तयित्वा सुचिरं धैर्यमालम्ब्य रावणः}
{राघवं हरियूथांश्च ददर्शायतलोचनः} %||6-32-5||

\twolineshloka
{प्रेक्षतो राक्षसेन्द्रस्य तान्यनीकानि भागशः}
{राघवप्रियकामार्थं लङ्कामारुरुहुस्तदा} %||6-32-6||

\twolineshloka
{ते ताम्रवक्त्रा हेमाभा रामार्थे त्यक्तजीविताः}
{लङ्कामेवाह्यवर्तन्त सालतालशिलायुधाः} %||6-32-7||

\twolineshloka
{ते द्रुमैः पर्वताग्रैश्च मुष्टिभिश्च प्लवङ्गमाः}
{प्रासादाग्राणि चोच्चानि ममन्तुस्तोरणानि च} %||6-32-8||

\twolineshloka
{पारिखाः पूरयन्ति स्म प्रसन्नसलिलायुताः}
{पांसुभिः पर्वताग्रैश्च तृणैः काष्ठैश्च वानराः} %||6-32-9||

\twolineshloka
{ततः सहस्रयूथाश्च कोटियूथाश्च यूथपाः}
{कोटीशतयुताश्चान्ये लङ्कामारुरुहुस्तदा} %||6-32-10||

\twolineshloka
{काञ्चनानि प्रमृद्नन्तस्तोरणानि प्लवङ्गमाः}
{कैलासशिखराभानि गोपुराणि प्रमथ्य च} %||6-32-11||

\twolineshloka
{आप्लवन्तः प्लवन्तश्च गर्जन्तश्च प्लवङ्गमाः}
{लङ्कां तामभ्यवर्तन्त महावारणसंनिभाः} %||6-32-12||

\twolineshloka
{जयत्यतिबलो रामो लक्ष्मणश्च महाबलः}
{राजा जयति सुग्रीवो राघवेणाभिपालितः} %||6-32-13||

\twolineshloka
{इत्येवं घोषयन्तश्च गर्जन्तश्च प्लवङ्गमाः}
{अभ्यधावन्त लङ्कायाः प्राकारं कामरूपिणः} %||6-32-14||

\twolineshloka
{वीरबाहुः सुबाहुश्च नलश्च वनगोचरः}
{निपीड्योपनिविष्टास्ते प्राकारं हरियूथपाः} %||6-32-15||

{एतस्मिन्नन्तरे चक्रुः स्कन्धावारनिवेशनम्॥\devanumber{16}॥} %||6-32-16|| (Check)

\addtocounter{shlokacount}{1}

\twolineshloka
{पूर्वद्वारं तु कुमुदः कोटिभिर्दशभिर्वृतः}
{आवृत्य बलवांस्तस्थौ हरिभिर्जितकाशिभिः} %||6-32-17||

\twolineshloka
{दक्षिणद्वारमागम्य वीरः शतबलिः कपिः}
{आवृत्य बलवांस्तस्थौ विंशत्या कोटिभिर्वृतः} %||6-32-18||

\twolineshloka
{सुषेणः पश्चिमद्वारं गतस्तारा पिता हरिः}
{आवृत्य बलवांस्तस्थौ षष्टि कोटिभिरावृतः} %||6-32-19||

\twolineshloka
{उत्तरद्वारमासाद्य रामः सौमित्रिणा सह}
{आवृत्य बलवांस्तस्थौ सुग्रीवश्च हरीश्वरः} %||6-32-20||

\twolineshloka
{गोलाङ्गूलो महाकायो गवाक्षो भीमदर्शनः}
{वृतः कोट्या महावीर्यस्तस्थौ रामस्य पार्वतः} %||6-32-21||

\twolineshloka
{ऋष्काणां भीमवेगानां धूम्रः शत्रुनिबर्हणः}
{वृतः कोट्या महावीर्यस्तस्थौ रामस्य पार्श्वतः} %||6-32-22||

\twolineshloka
{संनद्धस्तु महावीर्यो गदापाणिर्विभीषणः}
{वृतो यस्तैस्तु सचिवैस्तस्थौ तत्र महाबलः} %||6-32-23||

\twolineshloka
{गजो गवाक्षो गवयः शरभो गन्धमादनः}
{समन्तात्परिघावन्तो ररक्षुर्हरिवाहिनीम्} %||6-32-24||

\twolineshloka
{ततः कोपपरीतात्मा रावणो राक्षसेश्वरः}
{निर्याणं सर्वसैन्यानां द्रुतमाज्ञापयत्तदा} %||6-32-25||

\twolineshloka
{निष्पतन्ति ततः सैन्या हृष्टा रावणचोदिताः}
{समये पूर्यमाणस्य वेगा इव महोदधेः} %||6-32-26||

\twolineshloka
{एतस्मिन्नन्तरे घोरः सङ्ग्रामः समपद्यत}
{रक्षसां वानराणां च यथा देवासुरे पुरा} %||6-32-27||

\twolineshloka
{ते गदाभिः प्रदीप्ताभिः शक्तिशूलपरश्वधैः}
{निजघ्नुर्वानरान्घोराः कथयन्तः स्वविक्रमान्} %||6-32-28||

\twolineshloka
{तथा वृक्षैर्महाकायाः पर्वताग्रैश्च वानराः}
{राक्षसास्तानि रक्षांसि नखैर्दन्तैश्च वेगिताः} %||6-32-29||

\twolineshloka
{राक्षसास्त्वपरे भीमाः प्राकारस्था महीगतान्}
{भिण्डिपालैश्च खड्गैश्च शूलैश्चैव व्यदारयन्} %||6-32-30||

\twolineshloka
{वानराश्चापि सङ्क्रुद्धाः प्राकारस्थान्महीगताः}
{राक्षसान्पातयामासुः समाप्लुत्य प्लवङ्गमाः} %||6-32-31||

\twolineshloka
{स सम्प्रहारस्तुमुलो मांसशोणितकर्दमः}
{रक्षसां वानराणां च सम्बभूवाद्भुतोपमाः} %||6-33-32||


{॥इत्यार्षे श्रीमद्रामायणे वाल्मीकीये आदिकाव्ये युद्धकाण्डे द्वात्रिंशः सर्गः॥}

\sect{त्रयस्त्रिंशः सर्गः}

\twolineshloka
{युध्यतां तु ततस्तेषां वानराणां महात्मनाम्}
{रक्षसां सम्बभूवाथ बलकोपः सुदारुणः} %||6-33-1||

\twolineshloka
{ते हयैः काञ्चनापीडैर्ध्वजैश्चाग्निशिखोपमैः}
{रथैश्चादित्यसङ्काशैः कवचैश्च मनोरमैः} %||6-33-2||

\twolineshloka
{निर्ययू राक्षसव्याघ्रा नादयन्तो दिशो दश}
{राक्षसा भीमकर्माणो रावणस्य जयैषिणः} %||6-33-3||

\twolineshloka
{वानराणामपि चमूर्महती जयमिच्चताम्}
{अभ्यधावत तां सेनां रक्षसां कामरूपिणाम्} %||6-33-4||

\twolineshloka
{एतस्मिन्नन्तरे तेषामन्योन्यमभिधावताम्}
{रक्षसां वानराणां च द्वन्द्वयुद्धमवर्तत} %||6-33-5||

\twolineshloka
{अङ्गदेनेन्द्रजित्सार्धं वालिपुत्रेण राक्षसः}
{अयुध्यत महातेजास्त्र्यम्बकेण यथान्धकः} %||6-33-6||

\twolineshloka
{प्रजङ्घेन च सम्पातिर्नित्यं दुर्मर्षणो रणे}
{जम्बूमालिनमारब्धो हनूमानपि वानरः} %||6-33-7||

\twolineshloka
{सङ्गतः सुमहाक्रोधो राक्षसो रावणानुजः}
{समरे तीक्ष्णवेगेन मित्रघ्नेन विभीषणः} %||6-33-8||

\twolineshloka
{तपनेन गजः सार्धं राक्षसेन महाबलः}
{निकुम्भेन महातेजा नीलोऽपि समयुध्यत} %||6-33-9||

\twolineshloka
{वानरेन्द्रस्तु सुग्रीवः प्रघसेन समागतः}
{सङ्गतः समरे श्रीमान्विरूपाक्षेण लक्ष्मणः} %||6-33-10||

\twolineshloka
{अग्निकेतुश्च दुर्धर्षो रश्मिकेतुश्च राक्षसः}
{सुप्तघ्नो यज्ञकोपश्च रामेण सह सङ्गताः} %||6-33-11||

\twolineshloka
{वज्रमुष्टिस्तु मैन्देन द्विविदेनाशनिप्रभः}
{राक्षसाभ्यां सुघोराभ्यां कपिमुख्यौ समागतौ} %||6-33-12||

\twolineshloka
{वीरः प्रतपनो घोरो राक्षसो रणदुर्धरः}
{समरे तीक्ष्णवेगेन नलेन समयुध्यत} %||6-33-13||

\twolineshloka
{धर्मस्य पुत्रो बलवान्सुषेण इति विश्रुतः}
{स विद्युन्मालिना सार्धमयुध्यत महाकपिः} %||6-33-14||

\twolineshloka
{वानराश्चापरे भीमा राक्षसैरपरैः सह}
{द्वन्द्वं समीयुर्बहुधा युद्धाय बहुभिः सह} %||6-33-15||

\twolineshloka
{तत्रासीत्सुमहद्युद्धं तुमुलं लोमहर्षणम्}
{रक्षसां वानराणां च वीराणां जयमिच्छताम्} %||6-33-16||

\twolineshloka
{हरिराक्षसदेहेभ्यः प्रसृताः केशशाड्वलाः}
{शरीरसङ्घाटवहाः प्रसुस्रुः शोणितापगाः} %||6-33-17||

\twolineshloka
{आजघानेन्द्रजित्क्रुद्धो वज्रेणेव शतक्रतुः}
{अङ्गदं गदया वीरं शत्रुसैन्यविदारणम्} %||6-33-18||

\twolineshloka
{तस्य काञ्चनचित्राङ्गं रथं साश्वं ससारथिम्}
{जघान समरे श्रीमानङ्गदो वेगवान्कपिः} %||6-33-19||

\twolineshloka
{सम्पातिस्तु त्रिभिर्बाणैः प्रजङ्घेन समाहतः}
{निजघानाश्वकर्णेन प्रजङ्घं रणमूर्धनि} %||6-33-20||

\twolineshloka
{जम्बूमाली रथस्थस्तु रथशक्त्या महाबलः}
{बिभेद समरे क्रुद्धो हनूमन्तं स्तनान्तरे} %||6-33-21||

\twolineshloka
{तस्य तं रथमास्थाय हनूमान्मारुतात्मजः}
{प्रममाथ तलेनाशु सह तेनैव रक्षसा} %||6-33-22||

\twolineshloka
{भिन्नगात्रः शरैस्तीक्ष्णैः क्षिप्रहस्तेन रक्षसा}
{प्रजघानाद्रिशृङ्गेण तपनं मुष्टिना गजः} %||6-33-23||

\twolineshloka
{ग्रसन्तमिव सैन्यानि प्रघसं वानराधिपः}
{सुग्रीवः सप्तपर्णेन निर्बिभेद जघान च} %||6-33-24||

\twolineshloka
{प्रपीड्य शरवर्षेण राक्षसं भीमदर्शनम्}
{निजघान विरूपाक्षं शरेणैकेन लक्ष्मणः} %||6-33-25||

\twolineshloka
{अग्निकेतुश्च दुर्धर्षो रश्मिकेतुश्च राक्षसः}
{सुप्तिघ्नो यज्ञकोपश्च रामं निर्बिभिदुः शरैः} %||6-33-26||

\twolineshloka
{तेषां चतुर्णां रामस्तु शिरांसि समरे शरैः}
{क्रुद्धश्चतुर्भिश्चिच्छेद घोरैरग्निशिखोपमैः} %||6-33-27||

\twolineshloka
{वज्रमुष्टिस्तु मैन्देन मुष्टिना निहतो रणे}
{पपात सरथः साश्वः पुराट्ट इव भूतले} %||6-33-28||

\twolineshloka
{वज्राशनिसमस्पर्शो द्विविदोऽप्यशनिप्रभम्}
{जघान गिरिशृङ्गेण मिषतां सर्वरक्षसाम्} %||6-33-29||

\twolineshloka
{द्विविदं वानरेन्द्रं तु द्रुमयोधिनमाहवे}
{शरैरशनिसङ्काशैः स विव्याधाशनिप्रभः} %||6-33-30||

\twolineshloka
{स शरैरतिविद्धाङ्गो द्विविदः क्रोधमूर्छितः}
{सालेन सरथं साश्वं निजघानाशनिप्रभम्} %||6-33-31||

\twolineshloka
{निकुम्भस्तु रणे नीलं नीलाञ्जनचयप्रभम्}
{निर्बिभेद शरैस्तीक्ष्णैः करैर्मेघमिवांशुमान्} %||6-33-32||

\twolineshloka
{पुनः शरशतेनाथ क्षिप्रहस्तो निशाचरः}
{बिभेद समरे नीलं निकुम्भः प्रजहास च} %||6-33-33||

\twolineshloka
{तस्यैव रथचक्रेण नीलो विष्णुरिवाहवे}
{शिरश्चिच्छेद समरे निकुम्भस्य च सारथेः} %||6-33-34||

\twolineshloka
{विद्युन्माली रथस्थस्तु शरैः काञ्चनभूषणैः}
{सुषेणं ताडयामास ननाद च मुहुर्मुहुः} %||6-33-35||

\twolineshloka
{तं रथस्थमथो दृष्ट्वा सुषेणो वानरोत्तमः}
{गिरिशृङ्गेण महता रथमाशु न्यपातयत्} %||6-33-36||

\twolineshloka
{लाघवेन तु संयुक्तो विद्युन्माली निशाचरः}
{अपक्रम्य रथात्तूर्णं गदापाणिः क्षितौ स्थितः} %||6-33-37||

\twolineshloka
{ततः क्रोधसमाविष्टः सुषेणो हरिपुङ्गवः}
{शिलां सुमहतीं गृह्य निशाचरमभिद्रवत्} %||6-33-38||

\twolineshloka
{तमापतन्तं गदया विद्युन्माली निशाचरः}
{वक्षस्यभिजग्नानाशु सुषेणं हरिसत्तमम्} %||6-33-39||

\twolineshloka
{गदाप्रहारं तं घोरमचिन्त्यप्लवगोत्तमः}
{तां शिलां पातयामास तस्योरसि महामृधे} %||6-33-40||

\twolineshloka
{शिलाप्रहाराभिहतो विद्युन्माली निशाचरः}
{निष्पिष्टहृदयो भूमौ गतासुर्निपपात ह} %||6-33-41||

\twolineshloka
{एवं तैर्वानरैः शूरैः शूरास्ते रजनीचराः}
{द्वन्द्वे विमृदितास्तत्र दैत्या इव दिवौकसैः} %||6-33-42||

\twolineshloka
{भल्लैः खड्गैर्गदाभिश्च शक्तितोमर पट्टसैः}
{अपविद्धश्च भिन्नश्च रथैः साङ्ग्रामिकैर्हयैः} %||6-33-43||

\threelineshloka
{निहतैः कुञ्जरैर्मत्तैस्तथा वानरराक्षसैः}
{चक्राक्षयुगदण्डैश्च भग्नैर्धरणिसंश्रितैः}
{बभूवायोधनं घोरं गोमायुगणसेवितम्} %||6-33-44||

\twolineshloka
{कबन्धानि समुत्पेतुर्दिक्षु वानररक्षसाम्}
{विमर्दे तुमुले तस्मिन्देवासुररणोपमे} %||6-33-45||

\fourlineindentedshloka
{विदार्यमाणा हरिपुङ्गवैस्तदा}
{निशाचराः शोणितदिग्धगात्राः}
{पुनः सुयुद्धं तरसा समाश्रिता}
{दिवाकरस्यास्तमयाभिकाङ्क्षिणः} %||6-34-46||


{॥इत्यार्षे श्रीमद्रामायणे वाल्मीकीये आदिकाव्ये युद्धकाण्डे त्रयस्त्रिंशः सर्गः॥}

\sect{चतुस्त्रिंशः सर्गः}

\twolineshloka
{युध्यतामेव तेषां तु तदा वानररक्षसाम्}
{रविरस्तं गतो रात्रिः प्रवृत्ता प्राणहारिणी} %||6-34-1||

\twolineshloka
{अन्योन्यं बद्धवैराणां घोराणां जयमिच्छताम्}
{सम्प्रवृत्तं निशायुद्धं तदा वारणरक्षसाम्} %||6-34-2||

\twolineshloka
{राक्षसोऽसीति हरयो हरिश्चासीति राक्षसाः}
{अन्योन्यं समरे जघ्नुस्तस्मिंस्तमसि दारुणे} %||6-34-3||

\twolineshloka
{जहि दारय चैतीति कथं विद्रवसीति च}
{एवं सुतुमुलः शब्दस्तस्मिंस्तमसि शुश्रुवे} %||6-34-4||

\twolineshloka
{कालाः काञ्चनसंनाहास्तस्मिंस्तमसि राक्षसाः}
{सम्प्रादृश्यन्त शैलेन्द्रा दीप्तौषधिवना इव} %||6-34-5||

\twolineshloka
{तस्मिंस्तमसि दुष्पारे राक्षसाः क्रोधमूर्छिताः}
{परिपेतुर्महावेगा भक्षयन्तः प्लवङ्गमान्} %||6-34-6||

\twolineshloka
{ते हयान्काञ्चनापीडन्ध्वजांश्चाग्निशिखोपमान्}
{आप्लुत्य दशनैस्तीक्ष्णैर्भीमकोपा व्यदारयन्} %||6-34-7||

\twolineshloka
{कुञ्जरान्कुञ्जरारोहान्पताकाध्वजिनो रथान्}
{चकर्षुश्च ददंशुश्च दशनैः क्रोधमूर्छिताः} %||6-34-8||

\twolineshloka
{लक्ष्मणश्चापि रामश्च शरैराशीविषोमपैः}
{दृश्यादृश्यानि रक्षांसि प्रवराणि निजघ्नतुः} %||6-34-9||

\twolineshloka
{तुरङ्गखुरविध्वस्तं रथनेमिसमुद्धतम्}
{रुरोध कर्णनेत्राणिण्युध्यतां धरणीरजः} %||6-34-10||

\twolineshloka
{वर्तमाने तथा घोरे सङ्ग्रामे लोमहर्षणे}
{रुधिरोदा महावेगा नद्यस्तत्र प्रसुस्रुवुः} %||6-34-11||

\twolineshloka
{ततो भेरीमृदङ्गानां पणवानां च निस्वनः}
{शङ्खवेणुस्वनोन्मिश्रः सम्बभूवाद्भुतोपमः} %||6-34-12||

\twolineshloka
{हतानां स्तनमानानां राक्षसानां च निस्वनः}
{शस्त्राणां वानराणां च सम्बभूवातिदारुणः} %||6-34-13||

\twolineshloka
{शस्त्रपुष्पोपहारा च तत्रासीद्युद्धमेदिनी}
{दुर्ज्ञेया दुर्निवेशा च शोणितास्रवकर्दमा} %||6-34-14||

\twolineshloka
{सा बभूव निशा घोरा हरिराक्षसहारिणी}
{कालरात्रीव भूतानां सर्वेषां दुरतिक्रमा} %||6-34-15||

\twolineshloka
{ततस्ते राक्षसास्तत्र तस्मिंस्तमसि दारुणे}
{राममेवाभ्यधावन्त संहृष्टा शरवृष्टिभिः} %||6-34-16||

\twolineshloka
{तेषामापततां शब्दः क्रुद्धानामभिगर्जताम्}
{उद्वर्त इव सप्तानां समुद्राणामभूत्स्वनः} %||6-34-17||

\twolineshloka
{तेषां रामः शरैः षड्भिः षड्जघान निशाचरान्}
{निमेषान्तरमात्रेण शितैरग्निशिखोपमैः} %||6-34-18||

\twolineshloka
{यज्ञशत्रुश्च दुर्धर्षो महापार्श्वमहोदरौ}
{वज्रदंष्ट्रो महाकायस्तौ चोभौ शुकसारणौ} %||6-34-19||

\twolineshloka
{ते तु रामेण बाणौघः सर्वमर्मसु ताडिताः}
{युद्धादपसृतास्तत्र सावशेषायुषोऽभवन्} %||6-34-20||

\twolineshloka
{ततः काञ्चनचित्राङ्गैः शरैरग्निशिखोपमैः}
{दिशश्चकार विमलाः प्रदिशश्च महाबलः} %||6-34-21||

\twolineshloka
{ये त्वन्ये राक्षसा वीरा रामस्याभिमुखे स्थिताः}
{तेऽपि नष्टाः समासाद्य पतङ्गा इव पावकम्} %||6-34-22||

\twolineshloka
{सुवर्णपुङ्खैर्विशिखैः सम्पतद्भिः सहस्रशः}
{बभूव रजनी चित्रा खद्योतैरिव शारदी} %||6-34-23||

\twolineshloka
{राक्षसानां च निनदैर्हरीणां चापि गर्जितैः}
{सा बभूव निशा घोरा भूयो घोरतरा तदा} %||6-34-24||

\twolineshloka
{तेन शब्देन महता प्रवृद्धेन समन्ततः}
{त्रिकूटः कन्दराकीर्णः प्रव्याहरदिवाचलः} %||6-34-25||

\twolineshloka
{गोलाङ्गूला महाकायास्तमसा तुल्यवर्चसः}
{सम्परिष्वज्य बाहुभ्यां भक्षयन्रजनीचरान्} %||6-34-26||

\twolineshloka
{अङ्गदस्तु रणे शत्रुं निहन्तुं समुपस्थितः}
{रावणेर्निजघानाशु सारथिं च हयानपि} %||6-34-27||

\twolineshloka
{इन्द्रजित्तु रथं त्यक्त्वा हताश्वो हतसारथिः}
{अङ्गदेन महामायस्तत्रैवान्तरधीयत} %||6-34-28||

\threelineshloka
{सोऽन्तर्धान गतः पापो रावणी रणकर्कशः}
{ब्रह्मदत्तवरो वीरो रावणिः क्रोधमूर्छितः}
{अदृश्यो निशितान्बाणान्मुमोचाशनिवर्चसः} %||6-34-29||

\twolineshloka
{स रामं लक्ष्मणं चैव घोरैर्नागमयैः शरैः}
{बिभेद समरे क्रुद्धः सर्वगात्रेषु राक्षसः} %||6-35-30||


{॥इत्यार्षे श्रीमद्रामायणे वाल्मीकीये आदिकाव्ये युद्धकाण्डे चतुस्त्रिंशः सर्गः॥}

\sect{पञ्चत्रिंशः सर्गः}

\twolineshloka
{स तस्य गतिमन्विच्छन्राजपुत्रः प्रतापवान्}
{दिदेशातिबलो रामो दशवानरयूथपान्} %||6-35-1||

\twolineshloka
{द्वौ सुषेणस्य दायादौ नीलं च प्लवगर्षभम्}
{अङ्गदं वालिपुत्रं च शरभं च तरस्विनम्} %||6-35-2||

\twolineshloka
{विनतं जाम्बवन्तं च सानुप्रस्थं महाबलम्}
{ऋषभं चर्षभस्कन्धमादिदेश परन्तपः} %||6-35-3||

\twolineshloka
{ते सम्प्रहृष्टा हरयो भीमानुद्यम्य पादपान्}
{आकाशं विविशुः सर्वे मार्गामाणा दिशो दश} %||6-35-4||

\twolineshloka
{तेषां वेगवतां वेगमिषुभिर्वेगवत्तरैः}
{अस्त्रवित्परमास्त्रेण वारयामास रावणिः} %||6-35-5||

\twolineshloka
{तं भीमवेगा हरयो नाराचैः क्षतविक्षताः}
{अन्धकारे न ददृशुर्मेघैः सूर्यमिवावृतम्} %||6-35-6||

\twolineshloka
{रामलक्ष्मणयोरेव सर्वमर्मभिदः शरान्}
{भृशमावेशयामास रावणिः समितिंजयः} %||6-35-7||

\twolineshloka
{निरन्तरशरीरौ तु भ्रातरौ रामलक्ष्मणौ}
{क्रुद्धेनेन्द्रजिता वीरौ पन्नगैः शरतां गतैः} %||6-35-8||

\twolineshloka
{तयोः क्षतजमार्गेण सुस्राव रुधिरं बहु}
{तावुभौ च प्रकाशेते पुष्पिताविव किंशुकौ} %||6-35-9||

\twolineshloka
{ततः पर्यन्तरक्ताक्षो भिन्नाञ्जनचयोपमः}
{रावणिर्भ्रातरौ वाक्यमन्तर्धानगतोऽब्रवीत्} %||6-35-10||

\twolineshloka
{युध्यमानमनालक्ष्यं शक्रोऽपि त्रिदशेश्वरः}
{द्रष्टुमासादितुं वापि न शक्तः किं पुनर्युवाम्} %||6-35-11||

\twolineshloka
{प्रावृताविषुजालेन राघवौ कङ्कपत्रिणा}
{एष रोषपरीतात्मा नयामि यमसादनम्} %||6-35-12||

\twolineshloka
{एवमुक्त्वा तु धर्मज्ञौ भ्रातरौ रामलक्ष्मणौ}
{निर्बिभेद शितैर्बाणैः प्रजहर्ष ननाद च} %||6-35-13||

\twolineshloka
{भिन्नाञ्जनचयश्यामो विस्फार्य विपुलं धनुः}
{भूयो भूयः शरान्घोरान्विससर्ज महामृधे} %||6-35-14||

\twolineshloka
{ततो मर्मसु मर्मज्ञो मज्जयन्निशिताञ्शरान्}
{रामलक्ष्मणयोर्वीरो ननाद च मुहुर्मुहुः} %||6-35-15||

\twolineshloka
{बद्धौ तु शरबन्धेन तावुभौ रणमूर्धनि}
{निमेषान्तरमात्रेण न शेकतुरुदीक्षितुम्} %||6-35-16||

\twolineshloka
{ततो विभिन्नसर्वाङ्गौ शरशल्याचितावुभौ}
{ध्वजाविव महेन्द्रस्य रज्जुमुक्तौ प्रकम्पितौ} %||6-35-17||

\twolineshloka
{तौ सम्प्रचलितौ वीरौ मर्मभेदेन कर्शितौ}
{निपेततुर्महेष्वासौ जगत्यां जगतीपती} %||6-35-18||

\twolineshloka
{तौ वीरशयने वीरौ शयानौ रुधिरोक्षितौ}
{शरवेष्टितसर्वाङ्गावार्तौ परमपीडितौ} %||6-35-19||

\twolineshloka
{न ह्यविद्धं तयोर्गात्रं बभूवाङ्गुलमन्तरम्}
{नानिर्भिन्नं न चास्तब्धमा कराग्रादजिह्मगैः} %||6-35-20||

\twolineshloka
{तौ तु क्रूरेण निहतौ रक्षसा कामरूपिणा}
{असृक्सुस्रुवतुस्तीव्रं जलं प्रस्रवणाविव} %||6-35-21||

\twolineshloka
{पपात प्रथमं रामो विद्धो मर्मसु मार्गणैः}
{क्रोधादिन्द्रजिता येन पुरा शक्रो विनिर्जितः} %||6-35-22||

\twolineshloka
{नारचैरर्धनाराचैर्भल्लैरञ्जलिकैरपि}
{विव्याध वत्सदन्तैश्च सिंहदंष्ट्रैः क्षुरैस्तथा} %||6-35-23||

\twolineshloka
{स वीरशयने शिश्ये विज्यमादाय कार्मुकम्}
{भिन्नमुष्टिपरीणाहं त्रिणतं रुक्मभूषितम्} %||6-35-24||

\twolineshloka
{बाणपातान्तरे रामं पतितं पुरुषर्षभम्}
{स तत्र लक्ष्मणो दृष्ट्वा निराशो जीवितेऽभवत्} %||6-35-25||

\fourlineindentedshloka
{बद्धौ तु वीरौ पतितौ शयानौ}
{तौ वानराः सम्परिवार्य तस्थुः}
{समागता वायुसुतप्रमुख्या}
{विषदमार्ताः परमं च जग्मुः} %||6-36-26||


{॥इत्यार्षे श्रीमद्रामायणे वाल्मीकीये आदिकाव्ये युद्धकाण्डे पञ्चत्रिंशः सर्गः॥}

\sect{षट्त्रिंशः सर्गः}

\twolineshloka
{ततो द्यां पृथिवीं चैव वीक्षमाणा वनौकसः}
{ददृशुः सन्ततौ बाणैर्भ्रातरौ रामलक्ष्मणौ} %||6-36-1||

\twolineshloka
{वृष्ट्वेवोपरते देवे कृतकर्मणि राक्षसे}
{आजगामाथ तं देशं ससुग्रीवो विभीषणः} %||6-36-2||

\twolineshloka
{नीलद्विविदमैन्दाश्च सुषेणसुमुखाङ्गदाः}
{तूर्णं हनुमता सार्धमन्वशोचन्त राघवौ} %||6-36-3||

\twolineshloka
{निश्चेष्टौ मन्दनिःश्वासौ शोणितौघपरिप्लुतौ}
{शरजालाचितौ स्तब्धौ शयानौ शरतल्पयोः} %||6-36-4||

\twolineshloka
{निःश्वसन्तौ यथा सर्पौ निश्चेष्टौ मन्दविक्रमौ}
{रुधिरस्रावदिग्धाङ्गौ तापनीयाविव ध्वजौ} %||6-36-5||

\twolineshloka
{तौ वीरशयने वीरौ शयानौ मन्दचेष्टितौ}
{यूथपैस्तैः परिवृतौ बाष्पव्याकुललोचनैः} %||6-36-6||

\twolineshloka
{राघवौ पतितौ दृष्ट्वा शरजालसमावृतौ}
{बभूवुर्व्यथिताः सर्वे वानराः सविभीषणाः} %||6-36-7||

\twolineshloka
{अन्तरिक्षं निरीक्षन्तो दिशः सर्वाश्च वानराः}
{न चैनं मायया छन्नं ददृशू रावणिं रणे} %||6-36-8||

\twolineshloka
{तं तु मायाप्रतिच्छिन्नं माययैव विभीषणः}
{वीक्षमाणो ददर्शाथ भ्रातुः पुत्रमवस्थितम्} %||6-36-9||

\twolineshloka
{तमप्रतिम कर्माणमप्रतिद्वन्द्वमाहवे}
{ददर्शान्तर्हितं वीरं वरदानाद्विभीषणः} %||6-36-10||

\twolineshloka
{इन्द्रजित्त्वात्मनः कर्म तौ शयानौ समीक्ष्य च}
{उवाच परमप्रीतो हर्षयन्सर्वनैरृतान्} %||6-36-11||

\twolineshloka
{दूषणस्य च हन्तारौ खरस्य च महाबलौ}
{सादितौ मामकैर्बाणैर्भ्रातरौ रामलक्ष्मणौ} %||6-36-12||

\twolineshloka
{नेमौ मोक्षयितुं शक्यावेतस्मादिषुबन्धनात्}
{सर्वैरपि समागम्य सर्षिसङ्घैः सुरासुरैः} %||6-36-13||

\twolineshloka
{यत्कृते चिन्तयानस्य शोकार्तस्य पितुर्मम}
{अस्पृष्ट्वा शयनं गात्रैस्त्रियामा याति शर्वती} %||6-36-14||

\twolineshloka
{कृत्स्नेयं यत्कृते लङ्का नदी वर्षास्विवाकुला}
{सोऽयं मूलहरोऽनर्थः सर्वेषां निहतो मया} %||6-36-15||

\twolineshloka
{रामस्य लक्ष्मणस्यैव सर्वेषां च वनौकसाम्}
{विक्रमा निष्फलाः सर्वे यथा शरदि तोयदाः} %||6-36-16||

\twolineshloka
{एवमुक्त्वा तु तान्सर्वान्राक्षसान्परिपार्श्वगान्}
{यूथपानपि तान्सर्वांस्ताडयामास रावणिः} %||6-36-17||

\twolineshloka
{तानर्दयित्वा बाणौघैस्त्रासयित्वा च वानरान्}
{प्रजहास महाबाहुर्वचनं चेदमब्रवीत्} %||6-36-18||

\twolineshloka
{शरबन्धेन घोरेण मया बद्धौ चमूमुखे}
{सहितौ भ्रातरावेतौ निशामयत राक्षसाः} %||6-36-19||

\twolineshloka
{एवमुक्तास्तु ते सर्वे राक्षसाः कूटयोधिनः}
{परं विस्मयमाजग्मुः कर्मणा तेन तोषिताः} %||6-36-20||

\twolineshloka
{विनेदुश्च महानादान्सर्वे ते जलदोपमाः}
{हतो राम इति ज्ञात्वा रावणिं समपूजयन्} %||6-36-21||

\twolineshloka
{निष्पन्दौ तु तदा दृष्ट्वा तावुभौ रामलक्ष्मणौ}
{वसुधायां निरुच्छ्वासौ हतावित्यन्वमन्यत} %||6-36-22||

\twolineshloka
{हर्षेण तु समाविष्ट इन्द्रजित्समितिंजयः}
{प्रविवेश पुरीं लङ्कां हर्षयन्सर्वनैरृतान्} %||6-36-23||

\twolineshloka
{रामलक्ष्मणयोर्दृष्ट्वा शरीरे सायकैश्चिते}
{सर्वाणि चाङ्गोपाङ्गानि सुग्रीवं भयमाविशत्} %||6-36-24||

\twolineshloka
{तमुवाच परित्रस्तं वानरेन्द्रं विभीषणः}
{सबाष्पवदनं दीनं शोकव्याकुललोचनम्} %||6-36-25||

\twolineshloka
{अलं त्रासेन सुग्रीव बाष्पवेगो निगृह्यताम्}
{एवं प्रायाणि युद्धानि विजयो नास्ति नैष्ठिकः} %||6-36-26||

\twolineshloka
{सशेषभाग्यतास्माकं यदि वीर भविष्यति}
{मोहमेतौ प्रहास्येते भ्रातरौ रामलक्ष्मणौ} %||6-36-27||

\twolineshloka
{पर्यवस्थापयात्मानमनाथं मां च वानर}
{सत्यधर्मानुरक्तानां नास्ति मृत्युकृतं भयम्} %||6-36-28||

\twolineshloka
{एवमुक्त्वा ततस्तस्य जलक्लिन्नेन पाणिना}
{सुग्रीवस्य शुभे नेत्रे प्रममार्ज विभीषणः} %||6-36-29||

\twolineshloka
{प्रमृज्य वदनं तस्य कपिराजस्य धीमतः}
{अब्रवीत्कालसम्प्रातमसम्भ्रान्तमिदं वचः} %||6-36-30||

\twolineshloka
{न कालः कपिराजेन्द्र वैक्लव्यमनुवर्तितुम्}
{अतिस्नेहोऽप्यकालेऽस्मिन्मरणायोपपद्यते} %||6-36-31||

\twolineshloka
{तस्मादुत्सृज्य वैक्लव्यं सर्वकार्यविनाशनम्}
{हितं रामपुरोगाणां सैन्यानामनुचिन्त्यताम्} %||6-36-32||

\twolineshloka
{अथ वा रक्ष्यतां रामो यावत्संज्ञा विपर्ययः}
{लब्धसंज्ञौ तु काकुत्स्थौ भयं नो व्यपनेष्यतः} %||6-36-33||

\twolineshloka
{नैतत्किञ्चन रामस्य न च रामो मुमूर्षति}
{न ह्येनं हास्यते लक्ष्मीर्दुर्लभा या गतायुषाम्} %||6-36-34||

\twolineshloka
{तस्मादाश्वासयात्मानं बलं चाश्वासय स्वकम्}
{यावत्सर्वाणि सैन्यानि पुनः संस्थापयाम्यहम्} %||6-36-35||

\twolineshloka
{एते ह्युत्फुल्लनयनास्त्रासादागतसाध्वसाः}
{कर्णे कर्णे प्रकथिता हरयो हरिपुङ्गव} %||6-36-36||

\twolineshloka
{मां तु दृष्ट्वा प्रधावन्तमनीकं सम्प्रहर्षितुम्}
{त्यजन्तु हरयस्त्रासं भुक्तपूर्वामिव स्रजम्} %||6-36-37||

\twolineshloka
{समाश्वास्य तु सुग्रीवं राक्षसेन्द्रो विभीषणः}
{विद्रुतं वानरानीकं तत्समाश्वासयत्पुनः} %||6-36-38||

\twolineshloka
{इन्द्रजित्तु महामायः सर्वसैन्यसमावृतः}
{विवेश नगरीं लङ्कां पितरं चाभ्युपागमत्} %||6-36-39||

\twolineshloka
{तत्र रावणमासीनमभिवाद्य कृताञ्जलिः}
{आचचक्षे प्रियं पित्रे निहतौ रामलक्ष्मणौ} %||6-36-40||

\twolineshloka
{उत्पपात ततो हृष्टः पुत्रं च परिषस्वजे}
{रावणो रक्षसां मध्ये श्रुत्वा शत्रू निपातितौ} %||6-36-41||

\twolineshloka
{उपाघ्राय स मूर्ध्न्येनं पप्रच्छ प्रीतमानसः}
{पृच्छते च यथावृत्तं पित्रे सर्वं न्यवेदयत्} %||6-36-42||

\fourlineindentedshloka
{स हर्षवेगानुगतान्तरात्मा}
{श्रुत्वा वचस्तस्य महारथस्य}
{जहौ ज्वरं दाशरथेः समुत्थितम्}
{प्रहृष्य वाचाभिननन्द पुत्रम्} %||6-37-43||


{॥इत्यार्षे श्रीमद्रामायणे वाल्मीकीये आदिकाव्ये युद्धकाण्डे षट्त्रिंशः सर्गः॥}

\sect{सप्तत्रिंशः सर्गः}

\twolineshloka
{प्रतिप्रविष्टे लङ्कां तु कृतार्थे रावणात्मजे}
{राघवं परिवार्यार्ता ररक्षुर्वानरर्षभाः} %||6-37-1||

\twolineshloka
{हनूमानङ्गदो नीलः सुषेणः कुमुदो नलः}
{गजो गवाक्षो गवयः शरभो गन्धमादनः} %||6-37-2||

\twolineshloka
{जाम्बवानृषभः सुन्दो रम्भः शतबलिः पृथुः}
{व्यूढानीकाश्च यत्ताश्च द्रुमानादाय सर्वतः} %||6-37-3||

\twolineshloka
{वीक्षमाणा दिशः सर्वास्तिर्यगूर्ध्वं च वानराः}
{तृणेष्वपि च चेष्टत्सु राक्षसा इति मेनिरे} %||6-37-4||

\twolineshloka
{रावणश्चापि संहृष्टो विसृज्येन्द्रजितं सुतम्}
{आजुहाव ततः सीता रक्षणी राक्षसीस्तदा} %||6-37-5||

\twolineshloka
{राक्षस्यस्त्रिजटा चापि शासनात्तमुपस्थिताः}
{ता उवाच ततो हृष्टो राक्षसी राक्षसेश्वरः} %||6-37-6||

\twolineshloka
{हताविन्द्रजिताख्यात वैदेह्या रामलक्ष्मणौ}
{पुष्पकं च समारोप्य दर्शयध्वं हतौ रणे} %||6-37-7||

\twolineshloka
{यदाश्रयादवष्टब्धो नेयं मामुपतिष्ठति}
{सोऽस्या भर्ता सह भ्रात्रा निरस्तो रणमूर्धनि} %||6-37-8||

\twolineshloka
{निर्विशङ्का निरुद्विग्ना निरपेक्षा च मैथिली}
{मामुपस्थास्यते सीता सर्वाभरणभूषिता} %||6-37-9||

\twolineshloka
{अद्य कालवशं प्राप्तं रणे रामं सलक्ष्मणम्}
{अवेक्ष्य विनिवृत्ताशा नान्यां गतिमपश्यती} %||6-37-10||

\twolineshloka
{तस्य तद्वचनं श्रुत्वा रावणस्य दुरात्मनः}
{राक्षस्यस्तास्तथेत्युक्त्वा प्रजग्मुर्यत्र पुष्पकम्} %||6-37-11||

\twolineshloka
{ततः पुष्पकमादय राक्षस्यो रावणाज्ञया}
{अशोकवनिकास्थां तां मैथिलीं समुपानयन्} %||6-37-12||

\twolineshloka
{तामादाय तु राक्षस्यो भर्तृशोकपरायणाम्}
{सीतामारोपयामासुर्विमानं पुष्पकं तदा} %||6-37-13||

\twolineshloka
{ततः पुष्पकमारोप्य सीतां त्रिजटया सह}
{रावणोऽकारयल्लङ्कां पताकाध्वजमालिनीम्} %||6-37-14||

\twolineshloka
{प्राघोषयत हृष्टश्च लङ्कायां राक्षसेश्वरः}
{राघवो लक्ष्मणश्चैव हताविन्द्रजिता रणे} %||6-37-15||

\twolineshloka
{विमानेनापि सीता तु गत्वा त्रिजटया सह}
{ददर्श वानराणां तु सर्वं सिन्यं निपातितम्} %||6-37-16||

\twolineshloka
{प्रहृष्टमनसश्चापि ददर्श पिशिताशनान्}
{वानरांश्चापि दुःखार्तान्रामलक्ष्मणपार्श्वतः} %||6-37-17||

\twolineshloka
{ततः सीता ददर्शोभौ शयानौ शततल्पयोः}
{लक्ष्मणं चैव रामं च विसंज्ञौ शरपीडितौ} %||6-37-18||

\twolineshloka
{विध्वस्तकवचौ वीरौ विप्रविद्धशरासनौ}
{सायकैश्छिन्नसर्वाङ्गौ शरस्तम्भमयौ क्षितौ} %||6-37-19||

\twolineshloka
{तौ दृष्ट्वा भ्रातरौ तत्र वीरौ सा पुरुषर्षभौ}
{दुःखार्ता सुभृशं सीता करुणं विललाप ह} %||6-37-20||

\fourlineindentedshloka
{सा बाष्पशोकाभिहता समीक्ष्य}
{तौ भ्रातरौ देवसमप्रभावौ}
{वितर्कयन्ती निधनं तयोः सा}
{दुःखान्विता वाक्यमिदं जगाद} %||6-38-21||


{॥इत्यार्षे श्रीमद्रामायणे वाल्मीकीये आदिकाव्ये युद्धकाण्डे सप्तत्रिंशः सर्गः॥}

\sect{अष्टात्रिंशः सर्गः}

\twolineshloka
{भर्तारं निहतं दृष्ट्वा लक्ष्मणं च महाबलम्}
{विललाप भृशं सीता करुणं शोककर्शिता} %||6-38-1||

\twolineshloka
{ऊचुर्लक्षणिका ये मां पुत्रिण्यविधवेति च}
{तेऽस्य सर्वे हते रामेऽज्ञानिनोऽनृतवादिनः} %||6-38-2||

\twolineshloka
{यज्वनो महिषीं ये मामूचुः पत्नीं च सत्रिणः}
{तेऽद्य सर्वे हते रामेऽज्ञानिनोऽनृतवादिनः} %||6-38-3||

\twolineshloka
{वीरपार्थिवपत्नी त्वं ये धन्येति च मां विदुः}
{तेऽद्य सर्वे हते रामेऽज्ञानिनोऽनृतवादिनः} %||6-38-4||

\twolineshloka
{ऊचुः संश्रवणे ये मां द्विजाः कार्तान्तिकाः शुभाम्}
{तेऽद्य सर्वे हते रामेऽज्ञानिनोऽनृतवादिनः} %||6-38-5||

\twolineshloka
{इमानि खलु पद्मानि पादयोर्यैः किल स्त्रियः}
{अधिराज्येऽभिषिच्यन्ते नरेन्द्रैः पतिभिः सह} %||6-38-6||

\twolineshloka
{वैधव्यं यान्ति यैर्नार्योऽलक्षणैर्भाग्यदुर्लभाः}
{नात्मनस्तानि पश्यामि पश्यन्ती हतलक्षणा} %||6-38-7||

\twolineshloka
{सत्यानीमानि पद्मानि स्त्रीणामुक्त्वानि लक्षणे}
{तान्यद्य निहते रामे वितथानि भवन्ति मे} %||6-38-8||

\twolineshloka
{केशाः सूक्ष्माः समा नीला भ्रुवौ चासङ्गते मम}
{वृत्ते चालोमशे जङ्घे दन्ताश्चाविरला मम} %||6-38-9||

\twolineshloka
{शङ्खे नेत्रे करौ पादौ गुल्फावूरू च मे चितौ}
{अनुवृत्ता नखाः स्निग्धाः समाश्चाङ्गुलयो मम} %||6-38-10||

\twolineshloka
{स्तनौ चाविरलौ पीनौ ममेमौ मग्नचूचुकौ}
{मग्ना चोत्सङ्गिनी नाभिः पार्श्वोरस्कं च मे चितम्} %||6-38-11||

\twolineshloka
{मम वर्णो मणिनिभो मृदून्यङ्गरुहाणि च}
{प्रतिष्ठितां द्वदशभिर्मामूचुः शुभलक्षणाम्} %||6-38-12||

\twolineshloka
{समग्रयवमच्छिद्रं पाणिपादं च वर्णवत्}
{मन्दस्मितेत्येव च मां कन्यालक्षणिका विदुः} %||6-38-13||

\twolineshloka
{अधिराज्येऽभिषेको मे ब्राह्मणैः पतिना सह}
{कृतान्तकुशलैरुक्तं तत्सर्वं वितथीकृतम्} %||6-38-14||

\twolineshloka
{शोधयित्वा जनस्थानं प्रवृत्तिमुपलभ्य च}
{तीर्त्वा सागरमक्षोभ्यं भ्रातरौ गोष्पदे हतौ} %||6-38-15||

\twolineshloka
{ननु वारुणमाग्नेयमैन्द्रं वायव्यमेव च}
{अस्त्रं ब्रह्मशिरश्चैव राघवौ प्रत्यपद्यताम्} %||6-38-16||

\twolineshloka
{अदृश्यमानेन रणे मायया वासवोपमौ}
{मम नाथावनाथाया निहतौ रामलक्ष्मणौ} %||6-38-17||

\twolineshloka
{न हि दृष्टिपथं प्राप्य राघवस्य रणे रिपुः}
{जीवन्प्रतिनिवर्तेत यद्यपि स्यान्मनोजवः} %||6-38-18||

\twolineshloka
{न कालस्यातिभारोऽस्ति कृतान्तश्च सुदुर्जयः}
{यत्र रामः सह भ्रात्रा शेते युधि निपाथितः} %||6-38-19||

\twolineshloka
{नाहं शोचामि भर्तारं निहतं न च लक्ष्मणम्}
{नात्मानं जननी चापि यथा श्वश्रूं तपस्विनीम्} %||6-38-20||

\twolineshloka
{सा हि चिन्तयते नित्यं समाप्तव्रतमागतम्}
{कदा द्रक्ष्यामि सीतां च रामं च सहलक्ष्मणम्} %||6-38-21||

\twolineshloka
{परिदेवयमानां तां राक्षसी त्रिजटाब्रवीत्}
{मा विषादं कृथा देवि भर्तायं तव जीवति} %||6-38-22||

\twolineshloka
{कारणानि च वक्ष्यामि महान्ति सदृशानि च}
{यथेमौ जीवतो देवि भ्रातरौ रामलक्ष्मणौ} %||6-38-23||

\twolineshloka
{न हि कोपपरीतानि हर्षपर्युत्सुकानि च}
{भवन्ति युधि योधानां मुखानि निहते पतौ} %||6-38-24||

\twolineshloka
{इदं विमानं वैदेहि पुष्पकं नाम नामतः}
{दिव्यं त्वां धारयेन्नेदं यद्येतौ गजजीवितौ} %||6-38-25||

\twolineshloka
{हतवीरप्रधाना हि हतोत्साहा निरुद्यमा}
{सेना भ्रमति सङ्ख्येषु हतकर्णेव नौर्जले} %||6-38-26||

\twolineshloka
{इयं पुनरसम्भ्रान्ता निरुद्विग्ना तरस्विनी}
{सेना रक्षति काकुत्स्थौ मायया निर्जितौ रणे} %||6-38-27||

\twolineshloka
{सा त्वं भव सुविस्रब्धा अनुमानैः सुखोदयैः}
{अहतौ पश्य काकुत्स्थौ स्नेहादेतद्ब्रवीमि ते} %||6-38-28||

\twolineshloka
{अनृतं नोक्तपूर्वं मे न च वक्ष्ये कदाचन}
{चारित्रसुखशीलत्वात्प्रविष्टासि मनो मम} %||6-38-29||

\twolineshloka
{नेमौ शक्यौ रणे जेतुं सेन्द्रैरपि सुरासुरैः}
{एतयोराननं दृष्ट्वा मया चावेदितं तव} %||6-38-30||

\twolineshloka
{इदं च सुमहच्चिह्नं शनैः पश्यस्व मैथिलि}
{निःसंज्ञावप्युभावेतौ नैव लक्ष्मीर्वियुज्यते} %||6-38-31||

\twolineshloka
{प्रायेण गतसत्त्वानां पुरुषाणां गतायुषाम्}
{दृश्यमानेषु वक्त्रेषु परं भवति वैकृतम्} %||6-38-32||

\twolineshloka
{त्यज शोकं च दुःखं च मोहं च जनकात्मजे}
{रामलक्ष्मणयोरर्थे नाद्य शक्यमजीवितुम्} %||6-38-33||

\twolineshloka
{श्रुत्वा तु वचनं तस्याः सीता सुरसुतोपमा}
{कृताञ्जलिरुवाचेदमेवमस्त्विति मैथिली} %||6-38-34||

\twolineshloka
{विमानं पुष्पकं तत्तु समिवर्त्य मनोजवम्}
{दीना त्रिजटया सीता लङ्कामेव प्रवेशिता} %||6-38-35||

\twolineshloka
{ततस्त्रिजटया सार्धं पुष्पकादवरुह्य सा}
{अशोकवनिकामेव रक्षसीभिः प्रवेशिता} %||6-38-36||

\fourlineindentedshloka
{प्रविश्य सीता बहुवृक्षषण्डाम्}
{तां राक्षसेन्द्रस्य विहारभूमिम्}
{सम्प्रेक्ष्य सञ्चिन्त्य च राजपुत्रौ}
{परं विषादं समुपाजगाम} %||6-39-37||


{॥इत्यार्षे श्रीमद्रामायणे वाल्मीकीये आदिकाव्ये युद्धकाण्डे अष्टात्रिंशः सर्गः॥}

\sect{एकोनचत्वारिंशः सर्गः}

\twolineshloka
{घोरेण शरबन्धेन बद्धौ दशरथात्मजौ}
{निश्वसन्तौ यथा नागौ शयानौ रुधिरोक्षितौ} %||6-39-1||

\twolineshloka
{सर्वे ते वानरश्रेष्ठाः ससुग्रीवा महाबलाः}
{परिवार्य महात्मानौ तस्थुः शोकपरिप्लुताः} %||6-39-2||

\twolineshloka
{एतस्मिन्नन्तेरे रामः प्रत्यबुध्यत वीर्यवान्}
{स्थिरत्वात्सत्त्वयोगाच्च शरैः सन्दानितोऽपि सन्} %||6-39-3||

\twolineshloka
{ततो दृष्ट्वा सरुधिरं विषण्णं गाढमर्पितम्}
{भ्रातरं दीनवदनं पर्यदेवयदातुरः} %||6-39-4||

\twolineshloka
{किं नु मे सीतया कार्यं किं कार्यं जीवितेन वा}
{शयानं योऽद्य पश्यामि भ्रातरं युधि निर्जितम्} %||6-39-5||

\twolineshloka
{शक्या सीता समा नारी प्राप्तुं लोके विचिन्वता}
{न लक्ष्मणसमो भ्राता सचिवः साम्परायिकः} %||6-39-6||

\twolineshloka
{परित्यक्ष्याम्यहं प्राणान्वानराणां तु पश्यताम्}
{यदि पञ्चत्वमापन्नः सुमित्रानन्दवर्धनः} %||6-39-7||

\twolineshloka
{किं नु वक्ष्यामि कौसल्यां मातरं किं नु कैकयीम्}
{कथमम्बां सुमित्राञ्च पुत्रदर्शनलालसाम्} %||6-39-8||

\twolineshloka
{विवत्सां वेपमानां च क्रोशन्तीं कुररीमिव}
{कथमाश्वासयिष्यामि यदि यास्यामि तं विना} %||6-39-9||

\twolineshloka
{कथं वक्ष्यामि शत्रुघ्नं भरतं च यशस्विनम्}
{मया सह वनं यातो विना तेनागतः पुनः} %||6-39-10||

\twolineshloka
{उपालम्भं न शक्ष्यामि सोढुं बत सुमित्रया}
{इहैव देहं त्यक्ष्यामि न हि जीवितुमुत्सहे} %||6-39-11||

\twolineshloka
{धिङ्मां दुष्कृतकर्माणमनार्यं यत्कृते ह्यसौ}
{लक्ष्मणः पतितः शेते शरतल्पे गतासुवत्} %||6-39-12||

\twolineshloka
{त्वं नित्यं सुविषण्णं मामाश्वासयसि लक्ष्मण}
{गतासुर्नाद्य शक्नोषि मामार्तमभिभाषितुम्} %||6-39-13||

\twolineshloka
{येनाद्य बहवो युद्धे राक्षसा निहताः क्षितौ}
{तस्यामेव क्षितौ वीरः स शेते निहतः परैः} %||6-39-14||

\twolineshloka
{शयानः शरतल्पेऽस्मिन्स्वशोणितपरिप्लुतः}
{शरजालैश्चितो भाति भास्करोऽस्तमिव व्रजन्} %||6-39-15||

\twolineshloka
{बाणाभिहतमर्मत्वान्न शक्नोत्यभिवीक्षितुम्}
{रुजा चाब्रुवतो ह्यस्य दृष्टिरागेण सूच्यते} %||6-39-16||

\twolineshloka
{यथैव मां वनं यान्तमनुयातो महाद्युतिः}
{अहमप्यनुयास्यामि तथैवैनं यमक्षयम्} %||6-39-17||

\twolineshloka
{इष्टबन्धुजनो नित्यं मां च नित्यमनुव्रतः}
{इमामद्य गतोऽवस्थां ममानार्यस्य दुर्नयैः} %||6-39-18||

\twolineshloka
{सुरुष्टेनापि वीरेण लक्ष्मणेना न संस्मरे}
{परुषं विप्रियं वापि श्रावितं न कदाचन} %||6-39-19||

\twolineshloka
{विससर्जैकवेगेन पञ्चबाणशतानि यः}
{इष्वस्त्रेष्वधिकस्तस्मात्कार्तवीर्याच्च लक्ष्मणः} %||6-39-20||

\twolineshloka
{अस्त्रैरस्त्राणि यो हन्याच्छक्रस्यापि महात्मनः}
{सोऽयमुर्व्यांहतः शेते महार्हशयनोचितः} %||6-39-21||

\twolineshloka
{तच्च मिथ्या प्रलप्तं मां प्रधक्ष्यति न संशयः}
{यन्मया न कृतो राजा राक्षसानां विभीषणः} %||6-39-22||

\twolineshloka
{अस्मिन्मुहूर्ते सुग्रीव प्रतियातुमितोऽर्हसि}
{मत्वा हीनं मया राजन्रावणोऽभिद्रवेद्बली} %||6-39-23||

\twolineshloka
{अङ्गदं तु पुरस्कृत्य ससैन्यः ससुहृज्जनः}
{सागरं तर सुग्रीव पुनस्तेनैव सेतुना} %||6-39-24||

\twolineshloka
{कृतं हनुमता कार्यं यदन्यैर्दुष्करं रणे}
{ऋक्षराजेन तुष्यामि गोलाङ्गूलाधिपेन च} %||6-39-25||

\twolineshloka
{अङ्गदेन कृतं कर्म मैन्देन द्विविदेन च}
{युद्धं केसरिणा सङ्ख्ये घोरं सम्पातिना कृतम्} %||6-39-26||

\twolineshloka
{गवयेन गवाक्षेण शरभेण गजेन च}
{अन्यैश्च हरिभिर्युद्धं मदार्थे त्यक्तजीवितैः} %||6-39-27||

\threelineshloka
{न चातिक्रमितुं शक्यं दैवं सुग्रीव मानुषैः}
{यत्तु शक्यं वयस्येन सुहृदा वा परन्तप}
{कृतं सुग्रीव तत्सर्वं भवताधर्मभीरुणा} %||6-39-28||

\twolineshloka
{मित्रकार्यं कृतमिदं भवद्भिर्वानरर्षभाः}
{अनुज्ञाता मया सर्वे यथेष्टं गन्तुमर्हथ} %||6-39-29||

\twolineshloka
{शुश्रुवुस्तस्य ते सर्वे वानराः परिदेवितम्}
{वर्तयां चक्रुरश्रूणि नेत्रैः कृष्णेतरेक्षणाः} %||6-39-30||

\twolineshloka
{ततः सर्वाण्यनीकानि स्थापयित्वा विभीषणः}
{आजगाम गदापाणिस्त्वरितो यत्र राघवः} %||6-39-31||

\twolineshloka
{तं दृष्ट्वा त्वरितं यान्तं नीलाञ्जनचयोपमम्}
{वानरा दुद्रुवुः सर्वे मन्यमानास्तु रावणिम्} %||6-40-32||


{॥इत्यार्षे श्रीमद्रामायणे वाल्मीकीये आदिकाव्ये युद्धकाण्डे एकोनचत्वारिंशः सर्गः॥}

\sect{चत्वारिंशः सर्गः}

\twolineshloka
{अथोवाच महातेजा हरिराजो महाबलः}
{किमियं व्यथिता सेना मूढवातेव नौर्जले} %||6-40-1||

\twolineshloka
{सुग्रीवस्य वचः श्रुत्वा वालिपुत्रोऽङ्गदोऽब्रवीत्}
{न त्वं पश्यसि रामं च लक्ष्मणं च महाबलम्} %||6-40-2||

\twolineshloka
{शरजालाचितौ वीरावुभौ दशरथात्मजौ}
{शरतल्पे महात्मानौ शयानौ रुधिरोक्षितौ} %||6-40-3||

\twolineshloka
{अथाब्रवीद्वानरेन्द्रः सुग्रीवः पुत्रमङ्गदम्}
{नानिमित्तमिदं मन्ये भवितव्यं भयेन तु} %||6-40-4||

\twolineshloka
{विषण्णवदना ह्येते त्यक्तप्रहरणा दिशः}
{प्रपलायन्ति हरयस्त्रासादुत्फुल्ललोचनाः} %||6-40-5||

\twolineshloka
{अन्योन्यस्य न लज्जन्ते न निरीक्षन्ति पृष्ठतः}
{विप्रकर्षन्ति चान्योन्यं पतितं लङ्घयन्ति च} %||6-40-6||

\twolineshloka
{एतस्मिन्नन्तरे वीरो गदापाणिर्विभीषणः}
{सुग्रीवं वर्धयामास राघवं च निरैक्षत} %||6-40-7||

\twolineshloka
{विभीषणं तं सुग्रीवो दृष्ट्वा वानरभीषणम्}
{ऋक्षराजं समीपस्थं जाम्बवन्तमुवाच ह} %||6-40-8||

\twolineshloka
{विभीषणोऽयं सम्प्राप्तो यं दृष्ट्वा वानरर्षभाः}
{विद्रवन्ति परित्रस्ता रावणात्मजशङ्कया} %||6-40-9||

\twolineshloka
{शीघ्रमेतान्सुवित्रस्तान्बहुधा विप्रधावितान्}
{पर्यवस्थापयाख्याहि विभीषणमुपस्थितम्} %||6-40-10||

\twolineshloka
{सुग्रीवेणैवमुक्तस्तु जाम्बवानृक्षपार्थिवः}
{वानरान्सान्त्वयामास संनिवर्त्य प्रहावतः} %||6-40-11||

\twolineshloka
{ते निवृत्ताः पुनः सर्वे वानरास्त्यक्तसम्भ्रमाः}
{ऋक्षराजवचः श्रुत्वा तं च दृष्ट्वा विभीषणम्} %||6-40-12||

\twolineshloka
{विभीषणस्तु रामस्य दृष्ट्वा गात्रं शरैश्चितम्}
{लक्ष्मणस्य च धर्मात्मा बभूव व्यथितेन्द्रियः} %||6-40-13||

\twolineshloka
{जलक्लिन्नेन हस्तेन तयोर्नेत्रे प्रमृज्य च}
{शोकसम्पीडितमना रुरोद विललाप च} %||6-40-14||

\twolineshloka
{इमौ तौ सत्त्वसम्पन्नौ विक्रान्तौ प्रियसंयुगौ}
{इमामवस्थां गमितौ राकसैः कूटयोधिभिः} %||6-40-15||

\twolineshloka
{भ्रातुः पुत्रेण मे तेन दुष्पुत्रेण दुरात्मना}
{राक्षस्या जिह्मया बुद्ध्या छलितावृजुविक्रमौ} %||6-40-16||

\twolineshloka
{शरैरिमावलं विद्धौ रुधिरेण समुक्षितौ}
{वसुधायामिम सुप्तौ दृश्येते शल्यकाविव} %||6-40-17||

\twolineshloka
{ययोर्वीर्यमुपाश्रित्य प्रतिष्ठा काङ्क्षिता मया}
{तावुभौ देहनाशाय प्रसुप्तौ पुरुषर्षभौ} %||6-40-18||

\twolineshloka
{जीवन्नद्य विपन्नोऽस्मि नष्टराज्यमनोरथः}
{प्राप्तप्रतिज्ञश्च रिपुः सकामो रावणः कृतः} %||6-40-19||

\twolineshloka
{एवं विलपमानं तं परिष्वज्य विभीषणम्}
{सुग्रीवः सत्त्वसम्पन्नो हरिराजोऽब्रवीदिदम्} %||6-40-20||

\twolineshloka
{राज्यं प्राप्स्यसि धर्मज्ञ लङ्कायां नात्र संशयः}
{रावणः सह पुत्रेण स राज्यं नेह लप्स्यते} %||6-40-21||

\twolineshloka
{शरसम्पीडितावेतावुभौ राघवलक्ष्मणौ}
{त्यक्त्वा मोहं वधिष्येते सगणं रावणं रणे} %||6-40-22||

\twolineshloka
{तमेवं सान्त्वयित्वा तु समाश्वास्य च राक्षसम्}
{सुषेणं श्वशुरं पार्श्वे सुग्रीवस्तमुवाच ह} %||6-40-23||

\twolineshloka
{सह शूरैर्हरिगणैर्लब्धसंज्ञावरिन्दमौ}
{गच्छ त्वं भ्रातरौ गृह्य किष्किन्धां रामलक्ष्मणौ} %||6-40-24||

\twolineshloka
{अहं तु रावणं हत्वा सपुत्रं सहबान्धवम्}
{मैथिलीमानयिष्यामि शक्रो नष्टामिव श्रियम्} %||6-40-25||

\twolineshloka
{श्रुत्वैतद्वानरेन्द्रस्य सुषेणो वाक्यमब्रवीत्}
{देवासुरं महायुद्धमनुभूतं सुदारुणम्} %||6-40-26||

\twolineshloka
{तदा स्म दानवा देवाञ्शरसंस्पर्शकोविदाः}
{निजघ्नुः शस्त्रविदुषश्छादयन्तो मुहुर्मुहुः} %||6-40-27||

\twolineshloka
{तानार्तान्नष्टसंज्ञांश्च परासूंश्च बृहस्पतिः}
{विध्याभिर्मन्त्रयुक्ताभिरोषधीभिश्चिकित्सति} %||6-40-28||

\twolineshloka
{तान्यौषधान्यानयितुं क्षीरोदं यान्तु सागरम्}
{जवेन वानराः शीघ्रं सम्पाति पनसादयः} %||6-40-29||

\twolineshloka
{हरयस्तु विजानन्ति पार्वती ते महौषधी}
{संजीवकरणीं दिव्यां विशल्यां देवनिर्मिताम्} %||6-40-30||

\twolineshloka
{चन्द्रश्च नाम द्रोणश्च पर्वतौ सागरोत्तमे}
{अमृतं यत्र मथितं तत्र ते परमौषधी} %||6-40-31||

\twolineshloka
{ते तत्र निहिते देवैः पर्वते परमौषधी}
{अयं वायुसुतो राजन्हनूमांस्तत्र गच्छतु} %||6-40-32||

\twolineshloka
{एतस्मिन्नन्तरे वायुर्मेघांश्चापि सविद्युतः}
{पर्यस्यन्सागरे तोयं कम्पयन्निव पर्वतान्} %||6-40-33||

\twolineshloka
{महता पक्षवातेन सर्वे द्वीपमहाद्रुमाः}
{निपेतुर्भग्नविटपाः समूला लवणाम्भसि} %||6-40-34||

\twolineshloka
{अभवन्पन्नगास्त्रस्ता भोगिनस्तत्रवासिनः}
{शीघ्रं सर्वाणि यादांसि जग्मुश्च लवणार्णवम्} %||6-40-35||

\twolineshloka
{ततो मुहूर्तद्गरुडं वैनतेयं महाबलम्}
{वानरा ददृशुः सर्वे ज्वलन्तमिव पावकम्} %||6-40-36||

\twolineshloka
{तमागतमभिप्रेक्ष्य नागास्ते विप्रदुद्रुवुः}
{यैस्तौ सत्पुरुषौ बद्धौ शरभूतैर्महाबलौ} %||6-40-37||

\twolineshloka
{ततः सुपर्णः काकुत्स्थौ दृष्ट्वा प्रत्यभिनन्द्य च}
{विममर्श च पाणिभ्यां मुखे चन्द्रसमप्रभे} %||6-40-38||

\twolineshloka
{वैनतेयेन संस्पृष्टास्तयोः संरुरुहुर्व्रणाः}
{सुवर्णे च तनू स्निग्धे तयोराशु बभूवतुः} %||6-40-39||

\twolineshloka
{तेजो वीर्यं बलं चौज उत्साहश्च महागुणाः}
{प्रदर्शनं च बुद्धिश्च स्मृतिश्च द्विगुणं तयोः} %||6-40-40||

\twolineshloka
{तावुत्थाप्य महावीर्यौ गरुडो वासवोपमौ}
{उभौ तौ सस्वजे हृष्टौ रामश्चैनमुवाच ह} %||6-40-41||

\twolineshloka
{भवत्प्रसादाद्व्यसनं रावणिप्रभवं महत्}
{आवामिह व्यतिक्रान्तौ शीघ्रं च बलिनौ कृतौ} %||6-40-42||

\twolineshloka
{यथा तातं दशरथं यथाजं च पितामहम्}
{तथा भवन्तमासाद्य हृषयं मे प्रसीदति} %||6-40-43||

\twolineshloka
{को भवान्रूपसम्पन्नो दिव्यस्रगनुलेपनः}
{वसानो विरजे वस्त्रे दिव्याभरणभूषितः} %||6-40-44||

\twolineshloka
{तमुवाच महातेजा वैनतेयो महाबलः}
{पतत्रिराजः प्रीतात्मा हर्षपर्याकुलेक्षणः} %||6-40-45||

\twolineshloka
{अहं सखा ते काकुत्स्थ प्रियः प्राणो बहिश्चरः}
{गरुत्मानिह सम्प्राप्तो युवयोः साह्यकारणात्} %||6-40-46||

\twolineshloka
{असुरा वा महावीर्या दानवा वा महाबलाः}
{सुराश्चापि सगन्धर्वाः पुरस्कृत्य शतक्रतुम्} %||6-40-47||

\twolineshloka
{नेमं मोक्षयितुं शक्ताः शरबन्धं सुदारुणम्}
{माया बलादिन्द्रजिता निर्मितं क्रूरकर्मणा} %||6-40-48||

\twolineshloka
{एते नागाः काद्रवेयास्तीक्ष्णदंष्ट्राविषोल्बणाः}
{रक्षोमाया प्रभावेन शरा भूत्वा त्वदाश्रिताः} %||6-40-49||

\twolineshloka
{सभाग्यश्चासि धर्मज्ञ राम सत्यपराक्रम}
{लक्ष्मणेन सह भ्रात्रा समरे रिपुघातिना} %||6-40-50||

\twolineshloka
{इमं श्रुत्वा तु वृत्तान्तं त्वरमाणोऽहमागतः}
{सहसा युवयोः स्नेहात्सखित्वमनुपालयन्} %||6-40-51||

\twolineshloka
{मोक्षितौ च महाघोरादस्मात्सायकबन्धनात्}
{अप्रमादश्च कर्तव्यो युवाभ्यां नित्यमेव हि} %||6-40-52||

\twolineshloka
{प्रकृत्या राक्षसाः सर्वे सङ्ग्रामे कूटयोधिनः}
{शूराणां शुद्धभावानां भवतामार्जवं बलम्} %||6-40-53||

\twolineshloka
{तन्न विश्वसितव्यं वो राक्षसानां रणाजिरे}
{एतेनैवोपमानेन नित्यजिह्मा हि राक्षसाः} %||6-40-54||

\twolineshloka
{एवमुक्त्वा ततो रामं सुपर्णः सुमहाबलः}
{परिष्वज्य सुहृत्स्निग्धमाप्रष्टुमुपचक्रमे} %||6-40-55||

\twolineshloka
{सखे राघव धर्मज्ञ रिपूणामपि वत्सल}
{अभ्यनुज्ञातुमिच्छामि गमिष्यामि यथागतम्} %||6-40-56||

\twolineshloka
{बालवृद्धावशेषां तु लङ्कां कृत्वा शरोर्मिभिः}
{रावणं च रिपुं हत्वा सीतां समुपलप्स्यसे} %||6-40-57||

\twolineshloka
{इत्येवमुक्त्वा वचनं सुपर्णः शीघ्रविक्रमः}
{रामं च विरुजं कृत्वा मध्ये तेषां वनौकसाम्} %||6-40-58||

\twolineshloka
{प्रदक्षिणं ततः कृत्वा परिष्वज्य च वीर्यवान्}
{जगामाकाशमाविश्य सुपर्णः पवनो यथा} %||6-40-59||

\twolineshloka
{विरुजौ राघवौ दृष्ट्वा ततो वानरयूथपाः}
{सिंहनादांस्तदा नेदुर्लाङ्गूलं दुधुवुश्च ते} %||6-40-60||

\twolineshloka
{ततो भेरीः समाजघ्नुर्मृदङ्गांश्च व्यनादयन्}
{दध्मुः शङ्खान्सम्प्रहृष्टाः क्ष्वेलन्त्यपि यथापुरम्} %||6-40-61||

\twolineshloka
{आस्फोट्यास्फोट्य विक्रान्ता वानरा नगयोधिनः}
{द्रुमानुत्पाट्य विविधांस्तस्थुः शतसहस्रशः} %||6-40-62||

\twolineshloka
{विसृजन्तो महानादांस्त्रासयन्तो निशाचरान्}
{लङ्काद्वाराण्युपाजग्मुर्योद्धुकामाः प्लवङ्गमाः} %||6-40-63||

\fourlineindentedshloka
{ततस्तु भीमस्तुमुलो निनादो}
{बभूव शाखामृगयूथपानाम्}
{क्षये निदाघस्य यथा घनानाम्}
{नादः सुभीमो नदतां निशीथे} %||6-41-64||


{॥इत्यार्षे श्रीमद्रामायणे वाल्मीकीये आदिकाव्ये युद्धकाण्डे चत्वारिंशः सर्गः॥}

\sect{एकचत्वारिंशः सर्गः}

\twolineshloka
{तेषां सुतुमुलं शब्दं वानराणां तरस्विनाम्}
{नर्दतां राक्षसैः सार्धं तदा शुश्राव रावणः} %||6-41-1||

\twolineshloka
{स्निग्धगम्भीरनिर्घोषं श्रुत्वा स निनदं भृशम्}
{सचिवानां ततस्तेषां मध्ये वचनमब्रवीत्} %||6-41-2||

\twolineshloka
{यथासौ सम्प्रहृष्टानां वानराणां समुत्थितः}
{बहूनां सुमहान्नादो मेघानामिव गर्जताम्} %||6-41-3||

\twolineshloka
{व्यक्तं सुमहती प्रीतिरेतेषां नात्र संशयः}
{तथा हि विपुलैर्नादैश्चुक्षुभे वरुणालयः} %||6-41-4||

\twolineshloka
{तौ तु बद्धौ शरैस्तीष्क्णैर्भ्रातरौ रामलक्ष्मणौ}
{अयं च सुमहान्नादः शङ्कां जनयतीव मे} %||6-41-5||

\twolineshloka
{एतत्तु वचनं चोक्त्वा मन्त्रिणो राक्षसेश्वरः}
{उवाच नैरृतांस्तत्र समीपपरिवर्तिनः} %||6-41-6||

\twolineshloka
{ज्ञायतां तूर्णमेतषां सर्वेषां वनचारिणाम्}
{शोककाले समुत्पन्ने हर्षकारणमुत्थितम्} %||6-41-7||

\twolineshloka
{तथोक्तास्तेन सम्भ्रान्ताः प्राकारमधिरुह्य ते}
{ददृशुः पालितां सेनां सुग्रीवेण महात्मना} %||6-41-8||

\twolineshloka
{तौ च मुक्तौ सुघोरेण शरबन्धेन राघवौ}
{समुत्थितौ महाभागौ विषेदुः प्रेक्ष्य राक्षसाः} %||6-41-9||

\twolineshloka
{सन्त्रस्तहृदया सर्वे प्राकारादवरुह्य ते}
{विषण्णवदनाः सर्वे राक्षसेन्द्रमुपस्थिताः} %||6-41-10||

\twolineshloka
{तदप्रियं दीनमुखा रावणस्य निशाचराः}
{कृत्स्नं निवेदयामासुर्यथावद्वाक्यकोविदाः} %||6-41-11||

\twolineshloka
{यौ ताविन्द्रजिता युद्धे भ्रातरौ रामलक्ष्मणौ}
{निबद्धौ शरबन्धेन निष्प्रकम्पभुजौ कृतौ} %||6-41-12||

\twolineshloka
{विमुक्तौ शरबन्धेन तौ दृश्येते रणाजिरे}
{पाशानिव गजौ छित्त्वा गजेन्द्रसमविक्रमौ} %||6-41-13||

\twolineshloka
{तच्छ्रुत्वा वचनं तेषां राक्षसेन्द्रो महाबलः}
{चिन्ताशोकसमाक्रान्तो विषण्णवदनोऽब्रवीत्} %||6-41-14||

\twolineshloka
{घोरैर्दत्तवरैर्बद्धौ शरैराशीविषोमपैः}
{अमोघैः सूर्यसङ्काशैः प्रमथ्येन्द्रजिता युधि} %||6-41-15||

\twolineshloka
{तमस्त्रबन्धमासाद्य यदि मुक्तौ रिपू मम}
{संशयस्थमिदं सर्वमनुपश्याम्यहं बलम्} %||6-41-16||

\twolineshloka
{निष्फलाः खलु संवृत्ताः शरा वासुकितेजसः}
{आदत्तं यैस्तु सङ्ग्रामे रिपूणां मम जीवितम्} %||6-41-17||

\twolineshloka
{एवमुक्त्वा तु सङ्क्रुद्धो निश्वसन्नुरगो यथा}
{अब्रवीद्रक्षसां मध्ये धूम्राक्षं नाम राकसम्} %||6-41-18||

\twolineshloka
{बलेन महता युक्तो रक्षसां भीमकर्मणाम्}
{त्वं वधायाभिनिर्याहि रामस्य सह वानरैः} %||6-41-19||

\twolineshloka
{एवमुक्तस्तु धूम्राक्षो राक्षसेन्द्रेण धीमता}
{कृत्वा प्रणामं संहृष्टो निर्जगाम नृपालयात्} %||6-41-20||

\twolineshloka
{अभिनिष्क्रम्य तद्द्वारं बलाध्यक्षमुवाच ह}
{त्वरयस्व बलं तूर्णं किं चिरेण युयुत्सतः} %||6-41-21||

\twolineshloka
{धूम्राक्षस्य वचः श्रुत्वा बलाध्यक्षो बलानुगः}
{बलमुद्योजयामास रावणस्याज्ञया द्रुतम्} %||6-41-22||

\twolineshloka
{ते बद्धघण्टा बलिनो घोररूपा निशाचराः}
{विनर्दमानाः संहृष्टा धूम्राक्षं पर्यवारयन्} %||6-41-23||

\twolineshloka
{विविधायुधहस्ताश्च शूलमुद्गरपाणयः}
{गदाभिः पट्टसैर्दण्डैरायसैर्मुसलैर्भृशम्} %||6-41-24||

\twolineshloka
{परिघैर्भिण्डिपालैश्च भल्लैः प्रासैः परश्वधैः}
{निर्ययू राक्षसा घोरा नर्दन्तो जलदा यथा} %||6-41-25||

\twolineshloka
{रथैः कवचिनस्त्वन्ये ध्वजैश्च समलङ्कृतैः}
{सुवर्णजालविहितैः खरैश्च विविधाननैः} %||6-41-26||

\twolineshloka
{हयैः परमशीघ्रैश्च गजेन्द्रैश्च मदोत्कटैः}
{निर्ययू राक्षसव्याघ्रा व्याघ्रा इव दुरासदाः} %||6-41-27||

\twolineshloka
{वृकसिंहमुखैर्युक्तं खरैः कनकभूषणैः}
{आरुरोह रथं दिव्यं धूम्राक्षः खरनिस्वनः} %||6-41-28||

\twolineshloka
{स निर्यातो महावीर्यो धूम्राक्षो राक्षसैर्वृतः}
{प्रहसन्पश्चिमद्वारं हनूमान्यत्र यूथपः} %||6-41-29||

\twolineshloka
{प्रयान्तं तु महाघोरं राक्षसं भीमदर्शनम्}
{अन्तरिक्षगताः क्रूराः शकुनाः प्रत्यवारयन्} %||6-41-30||

\twolineshloka
{रथशीर्षे महाभीमो गृध्रश्च निपपात ह}
{ध्वजाग्रे ग्रथिताश्चैव निपेतुः कुणपाशनाः} %||6-41-31||

\twolineshloka
{रुधिरार्द्रो महाञ्श्वेतः कबन्धः पतितो भुवि}
{विस्वरं चोत्सृजन्नादं धूम्राक्षस्य समीपतः} %||6-41-32||

\threelineshloka
{ववर्ष रुधिरं देवः सञ्चचाल च मेदिनी}
{प्रतिलोमं ववौ वायुर्निर्घातसमनिस्वनः}
{तिमिरौघावृतास्तत्र दिशश्च न चकाशिरे} %||6-41-33||

\twolineshloka
{स तूत्पातांस्ततो दृष्ट्वा राक्षसानां भयावहान्}
{प्रादुर्भूतान्सुघोरांश्च धूम्राक्षो व्यथितोऽभवत्} %||6-41-34||

\fourlineindentedshloka
{ततः सुभीमो बहुभिर्निशाचरैर्-}
{वृतोऽभिनिष्क्रम्य रणोत्सुको बली}
{ददर्श तां राघवबाहुपालिताम्}
{समुद्रकल्पां बहुवानरीं चमूम्} %||6-42-35||


{॥इत्यार्षे श्रीमद्रामायणे वाल्मीकीये आदिकाव्ये युद्धकाण्डे एकचत्वारिंशः सर्गः॥}

\sect{द्विचत्वारिंशः सर्गः}

\twolineshloka
{धूम्राक्षं प्रेक्ष्य निर्यान्तं राक्षसं भीमनिस्वनम्}
{विनेदुर्वानराः सर्वे प्रहृष्टा युद्धकाङ्क्षिणः} %||6-42-1||

\twolineshloka
{तेषां तु तुमुलं युद्धं संजज्ञे हरिरक्षसाम्}
{अन्योन्यं पादपैर्घोरैर्निघ्नतं शूलमुद्गरैः} %||6-42-2||

\twolineshloka
{राक्षसैर्वानरा घोरा विनिकृत्ताः समन्ततः}
{वानरै राक्षसाश्चापि द्रुमैर्भूमौ समीकृताः} %||6-42-3||

\twolineshloka
{राक्षसाश्चापि सङ्क्रुद्धा वानरान्निशितैः शरैः}
{विव्यधुर्घोरसङ्काशैः कङ्कपत्रैरजिह्मगैः} %||6-42-4||

\twolineshloka
{ते गदाभिश्च भीमाभिः पट्टसैः कूटमुद्गरैः}
{घोरैश्च परिघैश्चित्रैस्त्रिशूलैश्चापि संशितैः} %||6-42-5||

\twolineshloka
{विदार्यमाणा रक्षोभिर्वानरास्ते महाबलाः}
{अमर्षाज्जनितोद्धर्षाश्चक्रुः कर्माण्यभीतवत्} %||6-42-6||

\twolineshloka
{शरनिर्भिन्नगात्रास्ते शूलनिर्भिन्नदेहिनः}
{जगृहुस्ते द्रुमांस्तत्र शिलाश्च हरियूथपाः} %||6-42-7||

\twolineshloka
{ते भीमवेगा हरयो नर्दमानास्ततस्ततः}
{ममन्थू राक्षसान्भीमान्नामानि च बभाषिरे} %||6-42-8||

\twolineshloka
{तद्बभूवाद्भुतं घोरं युद्धं वानररक्षसाम्}
{शिलाभिर्विविधाभिश्च बहुशाखैश्च पादपैः} %||6-42-9||

\twolineshloka
{राक्षसा मथिताः केचिद्वानरैर्जितकाशिभिः}
{ववर्षू रुधिरं केचिन्मुखै रुधिरभोजनाः} %||6-42-10||

\twolineshloka
{पार्श्वेषु दारिताः केचित्केचिद्राशीकृता द्रुमैः}
{शिलाभिश्चूर्णिताः केचित्केचिद्दन्तैर्विदारिताः} %||6-42-11||

\twolineshloka
{ध्वजैर्विमथितैर्भग्नैः खरैश्च विनिपातितैः}
{रथैर्विध्वंसितैश्चापि पतितै रजनीचरैः} %||6-42-12||

\twolineshloka
{वानरैर्भीमविक्रान्तैराप्लुत्याप्लुत्य वेगितैः}
{राक्षसाः करजैस्तीक्ष्णैर्मुखेषु विनिकर्तिताः} %||6-42-13||

\twolineshloka
{विवर्णवदना भूयो विप्रकीर्णशिरोरुहाः}
{मूढाः शोणितगन्धेन निपेतुर्धरणीतले} %||6-42-14||

\twolineshloka
{नये तु परमक्रुद्धा राक्षसा भीमविक्रमाः}
{तलैरेवाभिधावन्ति वज्रस्पर्शसमैर्हरीन्} %||6-42-15||

\twolineshloka
{वनरैरापतन्तस्ते वेगिता वेगवत्तरैः}
{मुष्टिभिश्चरणैर्दन्तैः पादपैश्चापपोथिताः} %||6-42-16||

\twolineshloka
{सन्यं तु विद्रुतं दृष्ट्वा धूम्राक्षो राक्षसर्षभः}
{क्रोधेन कदनं चक्रे वानराणां युयुत्सताम्} %||6-42-17||

\twolineshloka
{प्रासैः प्रमथिताः केचिद्वानराः शोणितस्रवाः}
{मुद्गरैराहताः केचित्पतिता धरणीतले} %||6-42-18||

\twolineshloka
{परिघैर्मथितः केचिद्भिण्डिपालैर्विदारिताः}
{पट्टसैराहताः केचिद्विह्वलन्तो गतासवः} %||6-42-19||

\twolineshloka
{केचिद्विनिहता भूमौ रुधिरार्द्रा वनौकसः}
{केचिद्विद्राविता नष्टाः सङ्क्रुद्धै राक्षसैर्युधि} %||6-42-20||

\twolineshloka
{विभिन्नहृदयाः केचिदेकपार्श्वेन शायिताः}
{विदारितास्त्रशूलै च केचिदान्त्रैर्विनिस्रुताः} %||6-42-21||

\twolineshloka
{तत्सुभीमं महद्युद्धं हरिराकस सङ्कुलम्}
{प्रबभौ शस्त्रबहुलं शिलापादपसङ्कुलम्} %||6-42-22||

\twolineshloka
{धनुर्ज्यातन्त्रिमधुरं हिक्कातालसमन्वितम्}
{मन्द्रस्तनितसङ्गीतं युद्धगान्धर्वमाबभौ} %||6-42-23||

\twolineshloka
{धूम्राक्षस्तु धनुष्पाणिर्वानरान्रणमूर्धनि}
{हसन्विद्रावयामास दिशस्ताञ्शरवृष्टिभिः} %||6-42-24||

\twolineshloka
{धूम्राक्षेणार्दितं सैन्यं व्यथितं दृश्य मारुतिः}
{अभ्यवर्तत सङ्क्रुद्धः प्रगृह्य विपुलां शिलाम्} %||6-42-25||

\twolineshloka
{क्रोधाद्द्विगुणताम्राक्षः पितृतुल्यपराक्रमः}
{शिलां तां पातयामास धूम्राक्षस्य रथं प्रति} %||6-42-26||

\twolineshloka
{आपतन्तीं शिलां दृष्ट्वा गदामुद्यम्य सम्भ्रमात्}
{रथादाप्लुत्य वेगेन वसुधायां व्यतिष्ठत} %||6-42-27||

\twolineshloka
{सा प्रमथ्य रथं तस्य निपपात शिलाभुवि}
{सचक्रकूबरं साश्वं सध्वजं सशरासनम्} %||6-42-28||

\twolineshloka
{स भङ्क्त्वा तु रथं तस्य हनूमान्मारुतात्मजः}
{रक्षसां कदनं चक्रे सस्कन्धविटपैर्द्रुमैः} %||6-42-29||

\twolineshloka
{विभिन्नशिरसो भूत्वा राक्षसाः शोणितोक्षिताः}
{द्रुमैः प्रमथिताश्चान्ये निपेतुर्धरणीतले} %||6-42-30||

\twolineshloka
{विद्राव्य राक्षसं सैन्यं हनूमान्मारुतात्मजः}
{गिरेः शिखरमादाय धूम्राक्षमभिदुद्रुवे} %||6-42-31||

\twolineshloka
{तमापतन्तं धूम्राक्षो गदामुद्यम्य वीर्यवान्}
{विनर्दमानः सहसा हनूमन्तमभिद्रवत्} %||6-42-32||

\twolineshloka
{ततः क्रुद्धस्तु वेगेन गदां तां बहुकण्टकाम्}
{पातयामास धूम्राक्षो मस्तके तु हनूमतः} %||6-42-33||

\threelineshloka
{ताडितः स तया तत्र गदया भीमरूपया}
{स कपिर्मारुतबलस्तं प्रहारमचिन्तयन्}
{धूम्राक्षस्य शिरो मध्ये गिरिशृङ्गमपातयत्} %||6-42-34||

\twolineshloka
{स विह्वलितसर्वाङ्गो गिरिशृङ्गेण ताडितः}
{पपात सहसा भूमौ विकीर्ण इव पर्वतः} %||6-42-35||

\twolineshloka
{धूम्राक्षं निहतं दृष्ट्वा हतशेषा निशाचराः}
{त्रस्ताः प्रविविशुर्लङ्कां वध्यमानाः प्लवङ्गमैः} %||6-42-36||

\fourlineindentedshloka
{स तु पवनसुतो निहत्य शत्रुम्}
{क्षतजवहाः सरितश्च संविकीर्य}
{रिपुवधजनितश्रमो महात्मा}
{मुदमगमत्कपिभिश्च पूज्यमानः} %||6-43-37||


{॥इत्यार्षे श्रीमद्रामायणे वाल्मीकीये आदिकाव्ये युद्धकाण्डे द्विचत्वारिंशः सर्गः॥}

\sect{त्रिचत्वारिंशः सर्गः}

\twolineshloka
{धूम्राक्षं निहतं श्रुत्वा रावणो राक्षसेश्वरः}
{बलाध्यक्षमुवाचेदं कृताञ्जलिमुपस्थितम्} %||6-43-1||

\twolineshloka
{शीघ्रं निर्यान्तु दुर्धर्षा राक्षसा भीमविक्रमाः}
{अकम्पनं पुरस्कृत्य सर्वशस्त्रप्रकोविदम्} %||6-43-2||

\twolineshloka
{ततो नानाप्रहरणा भीमाक्षा भीमदर्शनाः}
{निष्पेतू राक्षसा मुख्या बलाध्यक्षप्रचोदिताः} %||6-43-3||

\twolineshloka
{रथमास्थाय विपुलं तप्तकाञ्चनकुण्डलः}
{राकसैः संवृतो घोरैस्तदा निर्यात्यकम्पनः} %||6-43-4||

\twolineshloka
{न हि कम्पयितुं शक्यः सुरैरपि महामृधे}
{अकम्पनस्ततस्तेषामादित्य इव तेजसा} %||6-43-5||

\twolineshloka
{तस्य निधावमानस्य संरब्धस्य युयुत्सया}
{अकस्माद्दैन्यमागच्छद्धयानां रथवाहिनाम्} %||6-43-6||

\twolineshloka
{व्यस्फुरन्नयनं चास्य सव्यं युद्धाभिनन्दिनः}
{विवर्णो मुखवर्णश्च गद्गदश्चाभवत्स्वरः} %||6-43-7||

\twolineshloka
{अभवत्सुदिने चापि दुर्दिने रूक्षमारुतम्}
{ऊचुः खगा मृगाः सर्वे वाचः क्रूरा भयावहाः} %||6-43-8||

\twolineshloka
{स सिंहोपचितस्कन्धः शार्दूलसमविक्रमः}
{तानुत्पातानचिन्त्यैव निर्जगाम रणाजिरम्} %||6-43-9||

\twolineshloka
{तदा निर्गच्छतस्तस्य रक्षसः सह राक्षसैः}
{बभूव सुमहान्नादः क्षोभयन्निव सागरम्} %||6-43-10||

\twolineshloka
{तेन शब्देन वित्रस्ता वानराणां महाचमूः}
{द्रुमशैलप्रहरणा योद्धुं समवतिष्ठत} %||6-43-11||

\twolineshloka
{तेषां युद्धं महारौद्रं संजज्ञे कपिरक्षसाम्}
{रामरावणयोरर्थे समभित्यक्तजीविनाम्} %||6-43-12||

\twolineshloka
{सर्वे ह्यतिबलाः शूराः सर्वे पर्वतसंनिभाः}
{हरयो राक्षसाश्चैव परस्परजिघंसवः} %||6-43-13||

\twolineshloka
{तेषां विनर्दातां शब्दः संयुगेऽतितरस्विनाम्}
{शुश्रुवे सुमहान्क्रोधादन्योन्यमभिगर्जताम्} %||6-43-14||

\twolineshloka
{रजश्चारुणवर्णाभं सुभीममभवद्भृशम्}
{उद्धूतं हरिरक्षोभिः संरुरोध दिशो दश} %||6-43-15||

\twolineshloka
{अन्योन्यं रजसा तेन कौशेयोद्धूतपाण्डुना}
{संवृतानि च भूतानि ददृशुर्न रणाजिरे} %||6-43-16||

\twolineshloka
{न ध्वजो न पताकावा वर्म वा तुरगोऽपि वा}
{आयुधं स्यन्दनं वापि ददृशे तेन रेणुना} %||6-43-17||

\twolineshloka
{शब्दश्च सुमहांस्तेषां नर्दतामभिधावताम्}
{श्रूयते तुमुले युद्धे न रूपाणि चकाशिरे} %||6-43-18||

\twolineshloka
{हरीनेव सुसङ्क्रुद्धा हरयो जघ्नुराहवे}
{राक्षसाश्चापि रक्षांसि निजघ्नुस्तिमिरे तदा} %||6-43-19||

\twolineshloka
{परांश्चैव विनिघ्नन्तः स्वांश्च वानरराक्षसाः}
{रुधिरार्द्रं तदा चक्रुर्महीं पङ्कानुलेपनाम्} %||6-43-20||

\twolineshloka
{ततस्तु रुधिरौघेण सिक्तं व्यपगतं रजः}
{शरीरशवसङ्कीर्णा बभूव च वसुन्धरा} %||6-43-21||

\twolineshloka
{द्रुमशक्तिशिलाप्रासैर्गदापरिघतोमरैः}
{हरयो राक्षसास्तूर्णं जघ्नुरन्योन्यमोजसा} %||6-43-22||

\twolineshloka
{बाहुभिः परिघाकारैर्युध्यन्तः पर्वतोपमाः}
{हरयो भीमकर्माणो राक्षसाञ्जघ्नुराहवे} %||6-43-23||

\twolineshloka
{राक्षसाश्चापि सङ्क्रुद्धाः प्रासतोमरपाणयः}
{कपीन्निजघ्निरे तत्र शस्त्रैः परमदारुणैः} %||6-43-24||

\twolineshloka
{हरयस्त्वपि रक्षांसि महाद्रुममहाश्मभिः}
{विदारयन्त्यभिक्रम्य शस्त्राण्याच्छिद्य वीर्यतः} %||6-43-25||

\twolineshloka
{एतस्मिन्नन्तरे वीरा हरयः कुमुदो नलः}
{मैन्दश्च परमक्रुद्धश्चक्रुर्वेगमनुत्तमम्} %||6-43-26||

\twolineshloka
{ते तु वृक्षैर्महावेगा राक्षसानां चमूमुखे}
{कदनं सुमह चक्रुर्लीलया हरियूथपाः} %||6-44-27||


{॥इत्यार्षे श्रीमद्रामायणे वाल्मीकीये आदिकाव्ये युद्धकाण्डे त्रिचत्वारिंशः सर्गः॥}

\sect{चतुश्चत्वारिंशः सर्गः}

\twolineshloka
{तद्दृष्ट्वा सुमहत्कर्म कृतं वानरसत्तमैः}
{क्रोधमाहारयामास युधि तीव्रमकम्पनः} %||6-44-1||

\twolineshloka
{क्रोधमूर्छितरूपस्तु ध्नुवन्परमकार्मुकम्}
{दृष्ट्वा तु कर्म शत्रूणां सारथिं वाक्यमब्रवीत्} %||6-44-2||

\twolineshloka
{तत्रैव तावत्त्वरितं रथं प्रापय सारथे}
{एतेऽत्र बहवो घ्नन्ति सुबहून्राक्षसान्रणे} %||6-44-3||

\twolineshloka
{एतेऽत्र बलवन्तो हि भीमकायाश्च वानराः}
{द्रुमशैलप्रहरणास्तिष्ठन्ति प्रमुखे मम} %||6-44-4||

\twolineshloka
{एतान्निहन्तुमिच्छामि समरश्लाघिनो ह्यहम्}
{एतैः प्रमथितं सर्वं दृश्यते राक्षसं बलम्} %||6-44-5||

\twolineshloka
{ततः प्रजविताश्वेन रथेन रथिनां वरः}
{हरीनभ्यहनत्क्रोधाच्छरजालैरकम्पनः} %||6-44-6||

\twolineshloka
{न स्थातुं वानराः शेकुः किं पुनर्योद्धुमाहवे}
{अकम्पनशरैर्भग्नाः सर्व एव प्रदुद्रुवुः} %||6-44-7||

\twolineshloka
{तान्मृत्युवशमापन्नानकम्पनवशं गतान्}
{समीक्ष्य हनुमाञ्ज्ञातीनुपतस्थे महाबलः} %||6-44-8||

\twolineshloka
{तं महाप्लवगं दृष्ट्वा सर्वे प्लवगयूथपाः}
{समेत्य समरे वीराः सहिताः पर्यवारयन्} %||6-44-9||

\twolineshloka
{व्यवस्थितं हनूमन्तं ते दृष्ट्वा हरियूथपाः}
{बभूवुर्बलवन्तो हि बलवन्तमुपाश्रिताः} %||6-44-10||

\twolineshloka
{अकम्पनस्तु शैलाभं हनूमन्तमवस्थितम्}
{महेन्द्र इव धाराभिः शरैरभिववर्ष ह} %||6-44-11||

\twolineshloka
{अचिन्तयित्वा बाणौघाञ्शरीरे पतिताञ्शितान्}
{अकम्पनवधार्थाय मनो दध्रे महाबलः} %||6-44-12||

\twolineshloka
{स प्रहस्य महातेजा हनूमान्मारुतात्मजः}
{अभिदुद्राव तद्रक्षः कम्पयन्निव मेदिनीम्} %||6-44-13||

\twolineshloka
{तस्याभिनर्दमानस्य दीप्यमानस्य तेजसा}
{बभूव रूपं दुर्धर्षं दीप्तस्येव विभावसोः} %||6-44-14||

\twolineshloka
{आत्मानं त्वप्रहरणं ज्ञात्वा क्रोधसमन्वितः}
{शैलमुत्पाटयामास वेगेन हरिपुङ्गवः} %||6-44-15||

\twolineshloka
{तं गृहीत्वा महाशैलं पाणिनैकेन मारुतिः}
{विनद्य सुमहानादं भ्रामयामास वीर्यवान्} %||6-44-16||

\twolineshloka
{ततस्तमभिदुद्राव राक्षसेन्द्रमकम्पनम्}
{यथा हि नमुचिं सङ्ख्ये वज्रेणेव पुरन्दरः} %||6-44-17||

\twolineshloka
{अकम्पनस्तु तद्दृष्ट्वा गिरिशृङ्गं समुद्यतम्}
{दूरादेव महाबाणैरर्धचन्द्रैर्व्यदारयत्} %||6-44-18||

\twolineshloka
{तत्पर्वताग्रमाकाशे रक्षोबाणविदारितम्}
{विकीर्णं पतितं दृष्ट्वा हनूमान्क्रोधमूर्छितः} %||6-44-19||

\twolineshloka
{सोऽश्वकर्णं समासाद्य रोषदर्पान्वितो हरिः}
{तूर्णमुत्पाटयामास महागिरिमिवोच्छ्रितम्} %||6-44-20||

\twolineshloka
{तं गृहीत्वा महास्कन्धं सोऽश्वकर्णं महाद्युतिः}
{प्रहस्य परया प्रीत्या भ्रामयामास संयुगे} %||6-44-21||

\twolineshloka
{प्रधावन्नुरुवेगेन प्रभञ्जंस्तरसा द्रुमान्}
{हनूमान्परमक्रुद्धश्चरणैर्दारयत्क्षितिम्} %||6-44-22||

\twolineshloka
{गजांश्च सगजारोहान्सरथान्रथिनस्तथा}
{जघान हनुमान्धीमान्राक्षसांश्च पदातिकान्} %||6-44-23||

\twolineshloka
{तमन्तकमिव क्रुद्धं समरे प्राणहारिणम्}
{हनूमन्तमभिप्रेक्ष्य राक्षसा विप्रदुद्रुवुः} %||6-44-24||

\twolineshloka
{तमापतन्तं सङ्क्रुद्धं राक्षसानां भयावहम्}
{ददर्शाकम्पनो वीरश्चुक्रोध च ननाद च} %||6-44-25||

\twolineshloka
{स चतुर्दशभिर्बाणैः शितैर्देहविदारणैः}
{निर्बिभेद हनूमन्तं महावीर्यमकम्पनः} %||6-44-26||

\twolineshloka
{स तथा प्रतिविद्धस्तु बह्वीभिः शरवृष्टिभिः}
{हनूमान्ददृशे वीरः प्ररूढ इव सानुमान्} %||6-44-27||

\twolineshloka
{ततोऽन्यं वृक्षमुत्पाट्य कृत्वा वेगमनुत्तमम्}
{शिरस्यभिजघानाशु राक्षसेन्द्रमकम्पनम्} %||6-44-28||

\twolineshloka
{स वृक्षेण हतस्तेन सक्रोधेन महात्मना}
{राक्षसो वानरेन्द्रेण पपात स ममार च} %||6-44-29||

\twolineshloka
{तं दृष्ट्वा निहतं भूमौ राक्षसेन्द्रमकम्पनम्}
{व्यथिता राक्षसाः सर्वे क्षितिकम्प इव द्रुमाः} %||6-44-30||

\twolineshloka
{त्यक्तप्रहरणाः सर्वे राक्षसास्ते पराजिताः}
{लङ्कामभिययुस्त्रस्ता वानरैस्तैरभिद्रुताः} %||6-44-31||

\twolineshloka
{ते मुक्तकेशाः सम्भ्रान्ता भग्नमानाः पराजिताः}
{स्रवच्छ्रमजलैरङ्गैः श्वसन्तो विप्रदुद्रुवुः} %||6-44-32||

\twolineshloka
{अन्योन्यं प्रममन्तुस्ते विविशुर्नगरं भयात्}
{पृष्ठतस्ते सुसम्मूढाः प्रेक्षमाणा मुहुर्मुहुः} %||6-44-33||

\twolineshloka
{तेषु लङ्कां प्रविष्टेषु राक्षसेषु महाबलाः}
{समेत्य हरयः सर्वे हनूमन्तमपूजयन्} %||6-44-34||

\twolineshloka
{सोऽपि प्रहृष्टस्तान्सर्वान्हरीन्सम्प्रत्यपूजयत्}
{हनूमान्सत्त्वसम्पन्नो यथार्हमनुकूलतः} %||6-44-35||

\twolineshloka
{विनेदुश्च यथा प्राणं हरयो जितकाशिनः}
{चकर्षुश्च पुनस्तत्र सप्राणानेव राक्षसान्} %||6-44-36||

\fourlineindentedshloka
{स वीरशोभामभजन्महाकपिः}
{समेत्य रक्षांसि निहत्य मारुतिः}
{महासुरं भीमममित्रनाशनम्}
{यथैव विष्णुर्बलिनं चमूमुखे} %||6-44-37||

\fourlineindentedshloka
{अपूजयन्देवगणास्तदा कपिम्}
{स्वयं च रामोऽतिबलश्च लक्ष्मणः}
{तथैव सुग्रीवमुखाः प्लवङ्गमा}
{विभीषणश्चैव महाबलस्तदा} %||6-45-38||


{॥इत्यार्षे श्रीमद्रामायणे वाल्मीकीये आदिकाव्ये युद्धकाण्डे चतुश्चत्वारिंशः सर्गः॥}

\sect{पञ्चचत्वारिंशः सर्गः}

\twolineshloka
{अकम्पनवधं श्रुत्वा क्रुद्धो वै राक्षसेश्वरः}
{किञ्चिद्दीनमुखश्चापि सचिवांस्तानुदैक्षत} %||6-45-1||

\twolineshloka
{स तु ध्यात्वा मुहूर्तं तु मन्त्रिभिः संविचार्य च}
{पुरीं परिययौ लङ्कां सर्वान्गुल्मानवेक्षितुम्} %||6-45-2||

\twolineshloka
{तां राक्षसगणैर्गुप्तां गुल्मैर्बहुभिरावृताम्}
{ददर्श नगरीं लङ्कां पताकाध्वजमालिनीम्} %||6-45-3||

\twolineshloka
{रुद्धां तु नगरीं दृष्ट्वा रावणो राक्षसेश्वरः}
{उवाचामर्षितः काले प्रहस्तं युद्धकोविदम्} %||6-45-4||

\twolineshloka
{पुरस्योपनिविष्टस्य सहसा पीडितस्य च}
{नान्यं युद्धात्प्रपश्यामि मोक्षं युद्धविशारद} %||6-45-5||

\twolineshloka
{अहं वा कुम्भकर्णो वा त्वं वा सेनापतिर्मम}
{इन्द्रजिद्वा निकुम्भो वा वहेयुर्भारमीदृशम्} %||6-45-6||

\twolineshloka
{स त्वं बलमितः शीघ्रमादाय परिगृह्य च}
{विजयायाभिनिर्याहि यत्र सर्वे वनौकसः} %||6-45-7||

\twolineshloka
{निर्याणादेव ते नूनं चपला हरिवाहिनी}
{नर्दतां राक्षसेन्द्राणां श्रुत्वा नादं द्रविष्यति} %||6-45-8||

\twolineshloka
{चपला ह्यविनीताश्च चलचित्ताश्च वानराः}
{न सहिष्यन्ति ते नादं सिंहनादमिव द्विपाः} %||6-45-9||

\twolineshloka
{विद्रुते च बले तस्मिन्रामः सौमित्रिणा सह}
{अवशस्ते निरालम्बः प्रहस्तवशमेष्यति} %||6-45-10||

\twolineshloka
{आपत्संशयिता श्रेयो नात्र निःसंशयीकृता}
{प्रतिलोमानुलोमं वा यद्वा नो मन्यसे हितम्} %||6-45-11||

\twolineshloka
{रावणेनैवमुक्तस्तु प्रहस्तो वाहिनीपतिः}
{राक्षसेन्द्रमुवाचेदमसुरेन्द्रमिवोशना} %||6-45-12||

\twolineshloka
{राजन्मन्त्रितपूर्वं नः कुशलैः सह मन्त्रिभिः}
{विवादश्चापि नो वृत्तः समवेक्ष्य परस्परम्} %||6-45-13||

\twolineshloka
{प्रदानेन तु सीतायाः श्रेयो व्यवसितं मया}
{अप्रदाने पुनर्युद्धं दृष्टमेतत्तथैव नः} %||6-45-14||

\twolineshloka
{सोऽहं दानैश्च मानैश्च सततं पूजितस्त्वया}
{सान्त्वैश्च विविधैः काले किं न कुर्यां प्रियं तव} %||6-45-15||

\twolineshloka
{न हि मे जीवितं रक्ष्यं पुत्रदारधनानि वा}
{त्वं पश्य मां जुहूषन्तं त्वदर्थे जीवितं युधि} %||6-45-16||

\twolineshloka
{एवमुक्त्वा तु भर्तारं रावणं वाहिनीपतिः}
{समानयत मे शीघ्रं राक्षसानां महद्बलम्} %||6-45-17||

\twolineshloka
{मद्बाणाशनिवेगेन हतानां तु रणाजिरे}
{अद्य तृप्यन्तु मांसेन पक्षिणः काननौकसाम्} %||6-45-18||

\twolineshloka
{इत्युक्तास्ते प्रहस्तेन बलाध्यक्षाः कृतत्वराः}
{बलमुद्योजयामासुस्तस्मिन्राक्षसमन्दिरे} %||6-45-19||

\twolineshloka
{सा बभूव मुहूर्तेन तिग्मनानाविधायुधैः}
{लङ्का राक्षसवीरैस्तैर्गजैरिव समाकुला} %||6-45-20||

\twolineshloka
{हुताशनं तर्पयतां ब्राह्मणांश्च नमस्यताम्}
{आज्यगन्धप्रतिवहः सुरभिर्मारुतो ववौ} %||6-45-21||

\twolineshloka
{स्रजश्च विविधाकारा जगृहुस्त्वभिमन्त्रिताः}
{सङ्ग्रामसज्जाः संहृष्टा धारयन्राक्षसास्तदा} %||6-45-22||

\twolineshloka
{सधनुष्काः कवचिनो वेगादाप्लुत्य राक्षसाः}
{रावणं प्रेक्ष्य राजानं प्रहस्तं पर्यवारयन्} %||6-45-23||

\twolineshloka
{अथामन्त्र्य च राजानं भेरीमाहत्य भैरवाम्}
{आरुरोह रथं दिव्यं प्रहस्तः सज्जकल्पितम्} %||6-45-24||

\twolineshloka
{हयैर्महाजवैर्युक्तं सम्यक्सूतसुसंयुतम्}
{महाजलदनिर्घोषं साक्षाच्चन्द्रार्कभास्वरम्} %||6-45-25||

\twolineshloka
{उरगध्वजदुर्धर्षं सुवरूथं स्वपस्करम्}
{सुवर्णजालसंयुक्तं प्रहसन्तमिव श्रिया} %||6-45-26||

\twolineshloka
{ततस्तं रथमास्थाय रावणार्पितशासनः}
{लङ्काया निर्ययौ तूर्णं बलेन महता वृतः} %||6-45-27||

\twolineshloka
{ततो दुन्दुभिनिर्घोषः पर्जन्यनिनदोपमः}
{शुश्रुवे शङ्खशब्दश्च प्रयाते वाहिनीपतौ} %||6-45-28||

\twolineshloka
{निनदन्तः स्वरान्घोरान्राक्षसा जग्मुरग्रतः}
{भीमरूपा महाकायाः प्रहस्तस्य पुरःसराः} %||6-45-29||

\twolineshloka
{व्यूढेनैव सुघोरेण पूर्वद्वारात्स निर्ययौ}
{गजयूथनिकाशेन बलेन महता वृतः} %||6-45-30||

\twolineshloka
{सागरप्रतिमौघेन वृतस्तेन बलेन सः}
{प्रहस्तो निर्ययौ तूर्णं क्रुद्धः कालान्तकोपमः} %||6-45-31||

\twolineshloka
{तस्य निर्याण घोषेण राक्षसानां च नर्दताम्}
{लङ्कायां सर्वभूतानि विनेदुर्विकृतैः स्वरैः} %||6-45-32||

\twolineshloka
{व्यभ्रमाकाशमाविश्य मांसशोणितभोजनाः}
{मण्डलान्यपसव्यानि खगाश्चक्रू रथं प्रति} %||6-45-33||

{वमन्त्यः पावकज्वालाः शिवा घोरा ववाशिरे॥\devanumber{34}॥} %||6-45-34|| (Check)

\addtocounter{shlokacount}{1}

\twolineshloka
{अन्तरिक्षात्पपातोल्का वायुश्च परुषो ववौ}
{अन्योन्यमभिसंरब्धा ग्रहाश्च न चकाशिरे} %||6-45-35||

\twolineshloka
{ववर्षू रुधिरं चास्य सिषिचुश्च पुरःसरान्}
{केतुमूर्धनि गृध्रोऽस्य विलीनो दक्षिणामुखः} %||6-45-36||

\twolineshloka
{सारथेर्बहुशश्चास्य सङ्ग्राममवगाहतः}
{प्रतोदो न्यपतद्धस्तात्सूतस्य हयसादिनः} %||6-45-37||

\twolineshloka
{निर्याण श्रीश्च यास्यासीद्भास्वरा च सुदुर्लभा}
{सा ननाश मुहूर्तेन समे च स्खलिता हयाः} %||6-45-38||

\twolineshloka
{प्रहस्तं त्वभिनिर्यान्तं प्रख्यात बलपौरुषम्}
{युधि नानाप्रहरणा कपिसेनाभ्यवर्तत} %||6-45-39||

\twolineshloka
{अथ घोषः सुतुमुलो हरीणां समजायत}
{वृक्षानारुजतां चैव गुर्वीश्चागृह्णतां शिलाः} %||6-45-40||

\threelineshloka
{उभे प्रमुदिते सैन्ये रक्षोगणवनौकसाम्}
{वेगितानां समर्थानामन्योन्यवधकाङ्क्षिणाम्}
{परस्परं चाह्वयतां निनादः श्रूयते महान्} %||6-45-41||

\fourlineindentedshloka
{ततः प्रहस्तः कपिराजवाहिनीम्}
{अभिप्रतस्थे विजयाय दुर्मतिः}
{विवृद्धवेगां च विवेश तां चमूम्}
{यथा मुमूर्षुः शलभो विभावसुम्} %||6-46-42||


{॥इत्यार्षे श्रीमद्रामायणे वाल्मीकीये आदिकाव्ये युद्धकाण्डे पञ्चचत्वारिंशः सर्गः॥}

\sect{षट्चत्वारिंशः सर्गः}

\twolineshloka
{ततः प्रहस्तं निर्यान्तं भीमं भीमपराक्रमम्}
{गर्जन्तं सुमहाकायं राक्षसैरभिसंवृतम्} %||6-46-1||

\twolineshloka
{ददर्श महती सेना वानराणां बलीयसाम्}
{अतिसंजातरोषाणां प्रहस्तमभिगर्जताम्} %||6-46-2||

\twolineshloka
{खड्गशक्त्यष्टिबाणाश्च शूलानि मुसलानि च}
{गदाश्च परिघाः प्रासा विविधाश्च परश्वधाः} %||6-46-3||

\twolineshloka
{धनूंषि च विचित्राणि राक्षसानां जयैषिणाम्}
{प्रगृहीतान्यशोभन्त वानरानभिधावताम्} %||6-46-4||

\twolineshloka
{जगृहुः पादपांश्चापि पुष्पितान्वानरर्षभाः}
{शिलाश्च विपुला दीर्घा योद्धुकामाः प्लवङ्गमाः} %||6-46-5||

\twolineshloka
{तेषामन्योन्यमासाद्य सङ्ग्रामः सुमहानभूत्}
{बहूनामश्मवृष्टिं च शरवृष्टिं च वर्षताम्} %||6-46-6||

\twolineshloka
{बहवो राक्षसा युद्धे बहून्वानरयूथपान्}
{वानरा राक्षसांश्चापि निजघ्नुर्बहवो बहून्} %||6-46-7||

\twolineshloka
{शूलैः प्रमथिताः केचित्केचित्तु परमायुधैः}
{परिघैराहताः केचित्केचिच्छिन्नाः परश्वधैः} %||6-46-8||

\twolineshloka
{निरुच्छ्वासाः पुनः केचित्पतिता धरणीतले}
{विभिन्नहृदयाः केचिदिषुसन्तानसन्दिताः} %||6-46-9||

\twolineshloka
{केचिद्द्विधाकृताः खड्गैः स्फुरन्तः पतिता भुवि}
{वानरा राक्षसैः शूलैः पार्श्वतश्च विदारिताः} %||6-46-10||

\twolineshloka
{वानरैश्चापि सङ्क्रुद्धै राक्षसौघाः समन्ततः}
{पादपैर्गिरिशृङ्गैश्च सम्पिष्टा वसुधातले} %||6-46-11||

\twolineshloka
{वज्रस्पर्शतलैर्हस्तैर्मुष्टिभिश्च हता भृशम्}
{वेमुः शोणितमास्येभ्यो विशीर्णदशनेक्षणः} %||6-46-12||

\twolineshloka
{आर्तस्वरं च स्वनतां सिंहनादं च नर्दताम्}
{बभूव तुमुलः शब्दो हरीणां रक्षसां युधि} %||6-46-13||

\twolineshloka
{वानरा राक्षसाः क्रुद्धा वीरमार्गमनुव्रताः}
{विवृत्तनयनाः क्रूराश्चक्रुः कर्माण्यभीतवत्} %||6-46-14||

\twolineshloka
{नरान्तकः कुम्भहनुर्महानादः समुन्नतः}
{एते प्रहस्तसचिवाः सर्वे जघ्नुर्वनौकसः} %||6-46-15||

\twolineshloka
{तेषामापततां शीघ्रं निघ्नतां चापि वानरान्}
{द्विविदो गिरिशृङ्गेण जघानैकं नरान्तकम्} %||6-46-16||

\twolineshloka
{दुर्मुखः पुनरुत्पाट्य कपिः स विपुलद्रुमम्}
{राक्षसं क्षिप्रहस्तस्तु समुन्नतमपोथयत्} %||6-46-17||

\twolineshloka
{जाम्बवांस्तु सुसङ्क्रुद्धः प्रगृह्य महतीं शिलाम्}
{पातयामास तेजस्वी महानादस्य वक्षसि} %||6-46-18||

\twolineshloka
{अथ कुम्भहनुस्तत्र तारेणासाद्य वीर्यवान्}
{वृक्षेणाभिहतो मूर्ध्नि प्राणांस्तत्याज राक्षसः} %||6-46-19||

\twolineshloka
{अमृष्यमाणस्तत्कर्म प्रहस्तो रथमास्थितः}
{चकार कदनं घोरं धनुष्पाणिर्वनौकसाम्} %||6-46-20||

\twolineshloka
{आवर्त इव संजज्ञे उभयोः सेनयोस्तदा}
{क्षुभितस्याप्रमेयस्य सागरस्येव निस्वनः} %||6-46-21||

\twolineshloka
{महता हि शरौघेण प्रहस्तो युद्धकोविदः}
{अर्दयामास सङ्क्रुद्धो वानरान्परमाहवे} %||6-46-22||

\twolineshloka
{वानराणां शरीरैस्तु राक्षसानां च मेदिनी}
{बभूव निचिता घोरा पतितैरिव पर्वतैः} %||6-46-23||

\twolineshloka
{सा महीरुधिरौघेण प्रच्छन्ना सम्प्रकाशते}
{सञ्छन्ना माधवे मासि पलाशैरिव पुष्पितैः} %||6-46-24||

\twolineshloka
{हतवीरौघवप्रां तु भग्नायुधमहाद्रुमाम्}
{शोणितौघमहातोयां यमसागरगामिनीम्} %||6-46-25||

\twolineshloka
{यकृत्प्लीहमहापङ्कां विनिकीर्णान्त्रशैवलाम्}
{भिन्नकायशिरोमीनामङ्गावयवशाड्वलाम्} %||6-46-26||

\twolineshloka
{गृध्रहंसगणाकीर्णां कङ्कसारससेविताम्}
{मेधःफेनसमाकीर्णामार्तस्तनितनिस्वनाम्} %||6-46-27||

\twolineshloka
{तां कापुरुषदुस्तारां युद्धभूमिमयीं नदीम्}
{नदीमिव घनापाये हंससारससेविताम्} %||6-46-28||

\twolineshloka
{राक्षसाः कपिमुख्याश्च तेरुस्तां दुस्तरां नदीम्}
{यथा पद्मरजोध्वस्तां नलिनीं गजयूथपाः} %||6-46-29||

\twolineshloka
{ततः सृजन्तं बाणौघान्प्रहस्तं स्यन्दने स्थितम्}
{ददर्श तरसा नीलो विनिघ्नन्तं प्लवङ्गमान्} %||6-46-30||

\twolineshloka
{स तं परमदुर्धर्षमापतन्तं महाकपिः}
{प्रहस्तं ताडयामास वृक्षमुत्पाट्य वीर्यवान्} %||6-46-31||

\twolineshloka
{स तेनाभिहतः क्रुद्धो नदन्राक्षसपुङ्गवः}
{ववर्ष शरवर्षाणि प्लवगानां चमूपतौ} %||6-46-32||

\twolineshloka
{अपारयन्वारयितुं प्रत्यगृह्णान्निमीलितः}
{यथैव गोवृषो वर्षं शारदं शीघ्रमागतम्} %||6-46-33||

\twolineshloka
{एवमेव प्रहस्तस्य शरवर्षं दुरासदम्}
{निमीलिताक्षः सहसा नीलः सेहे सुदारुणम्} %||6-46-34||

\twolineshloka
{रोषितः शरवर्षेण सालेन महता महान्}
{प्रजघान हयान्नीलः प्रहस्तस्य मनोजवान्} %||6-46-35||

\twolineshloka
{विधनुस्तु कृतस्तेन प्रहस्तो वाहिनीपतिः}
{प्रगृह्य मुसलं घोरं स्यन्दनादवपुप्लुवे} %||6-46-36||

\twolineshloka
{तावुभौ वाहिनीमुख्यौ जातरोषौ तरस्विनौ}
{स्थितौ क्षतजदिग्धाङ्गौ प्रभिन्नाविव कुञ्जरौ} %||6-46-37||

\twolineshloka
{उल्लिखन्तौ सुतीक्ष्णाभिर्दंष्ट्राभिरितरेतरम्}
{सिंहशार्दूलसदृशौ सिंहशार्दूलचेष्टितौ} %||6-46-38||

\twolineshloka
{विक्रान्तविजयौ वीरौ समरेष्वनिवर्तिनौ}
{काङ्क्षमाणौ यशः प्राप्तुं वृत्रवासवयोः समौ} %||6-46-39||

\twolineshloka
{आजघान तदा नीलं ललाटे मुसलेन सः}
{प्रहस्तः परमायस्तस्तस्य सुस्राव शोणितम्} %||6-46-40||

\twolineshloka
{ततः शोणितदिग्धाङ्गः प्रगृह्य सुमहातरुम्}
{प्रहस्तस्योरसि क्रुद्धो विससर्ज महाकपिः} %||6-46-41||

\twolineshloka
{तमचिन्त्यप्रहारं स प्रगृह्य मुसलं महत्}
{अभिदुद्राव बलिनं बली नीलं प्लवङ्गमम्} %||6-46-42||

\twolineshloka
{तमुग्रवेगं संरब्धमापतन्तं महाकपिः}
{ततः सम्प्रेक्ष्य जग्राह महावेगो महाशिलाम्} %||6-46-43||

\twolineshloka
{तस्य युद्धाभिकामस्य मृधे मुसलयोधिनः}
{प्रहस्तस्य शिलां नीलो मूर्ध्नि तूर्णमपातयत्} %||6-46-44||

\twolineshloka
{सा तेन कपिमुख्येन विमुक्ता महती शिला}
{बिभेद बहुधा घोरा प्रहस्तस्य शिरस्तदा} %||6-46-45||

\twolineshloka
{स गतासुर्गतश्रीको गतसत्त्वो गतेन्द्रियः}
{पपात सहसा भूमौ छिन्नमूल इव द्रुमः} %||6-46-46||

\twolineshloka
{विभिन्नशिरसस्तस्य बहु सुस्रावशोणितम्}
{शरीरादपि सुस्राव गिरेः प्रस्रवणं यथा} %||6-46-47||

\twolineshloka
{हते प्रहस्ते नीलेन तदकम्प्यं महद्बलम्}
{रक्षसामप्रहृष्टानां लङ्कामभिजगाम ह} %||6-46-48||

\twolineshloka
{न शेकुः समवस्थातुं निहते वाहिनीपतौ}
{सेतुबन्धं समासाद्य विशीर्णं सलिलं यथा} %||6-46-49||

\twolineshloka
{हते तस्मिंश्चमूमुख्ये राक्षसस्ते निरुद्यमाः}
{रक्षःपतिगृहं गत्वा ध्यानमूकत्वमागताः} %||6-46-50||

\fourlineindentedshloka
{ततस्तु नीलो विजयी महाबलः}
{प्रशस्यमानः स्वकृतेन कर्मणा}
{समेत्य रामेण सलक्ष्मणेन}
{प्रहृष्टरूपस्तु बभूव यूथपः} %||6-47-51||


{॥इत्यार्षे श्रीमद्रामायणे वाल्मीकीये आदिकाव्ये युद्धकाण्डे षट्चत्वारिंशः सर्गः॥}

\sect{सप्तचत्वारिंशः सर्गः}

\fourlineindentedshloka
{तस्मिन्हते राक्षससैन्यपाले}
{प्लवङ्गमानामृषभेण युद्धे}
{भीमायुधं सागरतुल्यवेगम्}
{प्रदुद्रुवे राक्षसराजसैन्यम्} %||6-47-1||

\fourlineindentedshloka
{गत्वा तु रक्षोऽधिपतेः शशंसुः}
{सेनापतिं पावकसूनुशस्तम्}
{तच्चापि तेषां वचनं निशम्य}
{रक्षोऽधिपः क्रोधवशं जगाम} %||6-47-2||

\fourlineindentedshloka
{सङ्ख्ये प्रहस्तं निहतं निशम्य}
{शोकार्दितः क्रोधपरीतचेताः}
{उवाच तान्नैरृतयोधमुख्यान्}
{इन्द्रो यथा चामरयोधमुख्यान्} %||6-47-3||

\twolineshloka
{नावज्ञा रिपवे कार्या यैरिन्द्रबलसूदनः}
{सूदितः सैन्यपालो मे सानुयात्रः सकुञ्जरः} %||6-47-4||

\twolineshloka
{सोऽहं रिपुविनाशाय विजयायाविचारयन्}
{स्वयमेव गमिष्यामि रणशीर्षं तदद्भुतम्} %||6-47-5||

\twolineshloka
{अद्य तद्वानरानीकं रामं च सहलक्ष्मणम्}
{निर्दहिष्यामि बाणौघैर्वनं दीप्तैरिवाग्निभिः} %||6-47-6||

\fourlineindentedshloka
{स एवमुक्त्वा ज्वलनप्रकाशम्}
{रथं तुरङ्गोत्तमराजियुक्तम्}
{प्रकाशमानं वपुषा ज्वलन्तम्}
{समारुरोहामरराजशत्रुः} %||6-47-7||

\fourlineindentedshloka
{स शङ्खभेरीपटह प्रणादैः}
{आस्फोटितक्ष्वेडितसिंहनादैः}
{पुण्यैः स्तवैश्चाप्यभिपूज्यमानः}
{तदा ययौ राक्षसराजमुख्यः} %||6-47-8||

\fourlineindentedshloka
{स शैलजीमूतनिकाश रूपैर्-}
{मांसाशनैः पावकदीप्तनेत्रैः}
{बभौ वृतो राक्षसराजमुख्यैर्-}
{भूतैर्वृतो रुद्र इवामरेशः} %||6-47-9||

\fourlineindentedshloka
{ततो नगर्याः सहसा महौजा}
{निष्क्रम्य तद्वानरसैन्यमुग्रम्}
{महार्णवाभ्रस्तनितं ददर्श}
{समुद्यतं पादपशैलहस्तम्} %||6-47-10||

\fourlineindentedshloka
{तद्राक्षसानीकमतिप्रचण्डम्}
{आलोक्य रामो भुजगेन्द्रबाहुः}
{विभीषणं शस्त्रभृतां वरिष्ठम्}
{उवाच सेनानुगतः पृथुश्रीः} %||6-47-11||

\fourlineindentedshloka
{नानापताकाध्वजशस्त्रजुष्टम्}
{प्रासासिशूलायुधचक्रजुष्टम्}
{सैन्यं नगेन्द्रोपमनागजुष्टम्}
{कस्येदमक्षोभ्यमभीरुजुष्टम्} %||6-47-12||

\fourlineindentedshloka
{ततस्तु रामस्य निशम्य वाक्यम्}
{विभीषणः शक्रसमानवीर्यः}
{शशंस रामस्य बलप्रवेकम्}
{महात्मनां राक्षसपुङ्गवानाम्} %||6-47-13||

\fourlineindentedshloka
{योऽसौ गजस्कन्धगतो महात्मा}
{नवोदितार्कोपमताम्रवक्त्रः}
{प्रकम्पयन्नागशिरोऽभ्युपैति}
{ह्यकम्पनं त्वेनमवेहि राजन्} %||6-47-14||

\fourlineindentedshloka
{योऽसौ रथस्थो मृगराजकेतुर्-}
{धून्वन्धनुः शक्रधनुःप्रकाशम्}
{करीव भात्युग्रविवृत्तदंष्ट्रः}
{स इन्द्रजिन्नाम वरप्रधानः} %||6-47-15||

\fourlineindentedshloka
{यश्चैष विन्ध्यास्तमहेन्द्रकल्पो}
{धन्वी रथस्थोऽतिरथोऽतिवीर्यः}
{विस्फारयंश्चापमतुल्यमानम्}
{नाम्नातिकायोऽतिविवृद्धकायः} %||6-47-16||

\fourlineindentedshloka
{योऽसौ नवार्कोदितताम्रचक्षुः}
{आरुह्य घण्टानिनदप्रणादम्}
{गजं खरं गर्जति वै महात्मा}
{महोदरो नाम स एष वीरः} %||6-47-17||

\fourlineindentedshloka
{योऽसौ हयं काञ्चनचित्रभाण्डम्}
{आरुह्य सन्ध्याभ्रगिरिप्रकाशम्}
{प्रासं समुद्यम्य मरीचिनद्धम्}
{पिशाच एषाशनितुल्यवेगः} %||6-47-18||

\fourlineindentedshloka
{यश्चैष शूलं निशितं प्रगृह्य}
{विद्युत्प्रभं किङ्करवज्रवेगम्}
{वृषेन्द्रमास्थाय गिरिप्रकाशम्}
{आयाति सोऽसौ त्रिशिरा यशस्वी} %||6-47-19||

\fourlineindentedshloka
{असौ च जीमूतनिकाश रूपः}
{कुम्भः पृथुव्यूढसुजातवक्षाः}
{समाहितः पन्नगराजकेतुर्-}
{विस्फारयन्भाति धनुर्विधून्वन्} %||6-47-20||

\fourlineindentedshloka
{यश्चैष जाम्बूनदवज्रजुष्टम्}
{दीप्तं सधूमं परिघं प्रगृह्य}
{आयाति रक्षोबलकेतुभूतः}
{सोऽसौ निकुम्भोऽद्भुतघोरकर्मा} %||6-47-21||

\fourlineindentedshloka
{यश्चैष चापासिशरौघजुष्टम्}
{पताकिनं पावकदीप्तरूपम्}
{रथं समास्थाय विभात्युदग्रो}
{नरान्तकोऽसौ नगशृङ्गयोधी} %||6-47-22||

\fourlineindentedshloka
{यश्चैष नानाविधघोररूपैर्-}
{व्याघ्रोष्ट्रनागेन्द्रमृगेन्द्रवक्त्रैः}
{भूतैर्वृतो भाति विवृत्तनेत्रैः}
{सोऽसौ सुराणामपि दर्पहन्ता} %||6-47-23||

\fourlineindentedshloka
{यत्रैतदिन्दुप्रतिमं विभातिं}
{छत्त्रं सितं सूक्ष्मशलाकमग्र्यम्}
{अत्रैष रक्षोऽधिपतिर्महात्मा}
{भूतैर्वृतो रुद्र इवावभाति} %||6-47-24||

\fourlineindentedshloka
{असौ किरीटी चलकुण्डलास्यो}
{नागेन्द्रविन्ध्योपमभीमकायः}
{महेन्द्रवैवस्वतदर्पहन्ता}
{रक्षोऽधिपः सूर्य इवावभाति} %||6-47-25||

\twolineshloka
{प्रत्युवाच ततो रामो विभीषणमरिन्दमम्}
{अहो दीप्तो महातेजा रावणो राक्षसेश्वरः} %||6-47-26||

\twolineshloka
{आदित्य इव दुष्प्रेक्ष्यो रश्मिभिर्भाति रावणः}
{सुव्यक्तं लक्षये ह्यस्य रूपं तेजःसमावृतम्} %||6-47-27||

\twolineshloka
{देवदानववीराणां वपुर्नैवंविधं भवेत्}
{यादृशं राक्षसेन्द्रस्य वपुरेतत्प्रकाशते} %||6-47-28||

\twolineshloka
{सर्वे पर्वतसङ्काशाः सर्वे पर्वतयोधिनः}
{सर्वे दीप्तायुधधरा योधश्चास्य महौजसः} %||6-47-29||

\twolineshloka
{भाति राक्षसराजोऽसौ प्रदीप्तैर्भीमविक्रमैः}
{भूतैः परिवृतस्तीक्ष्णैर्देहवद्भिरिवान्तकः} %||6-47-30||

\twolineshloka
{एवमुक्त्वा ततो रामो धनुरादाय वीर्यवान्}
{लक्ष्मणानुचरस्तस्थौ समुद्धृत्य शरोत्तमम्} %||6-47-31||

\fourlineindentedshloka
{ततः स रक्षोऽधिपतिर्महात्मा}
{रक्षांसि तान्याह महाबलानि}
{द्वारेषु चर्यागृहगोपुरेषु}
{सुनिर्वृतास्तिष्ठत निर्विशङ्काः} %||6-47-32||

\fourlineindentedshloka
{विसर्जयित्वा सहसा ततस्तान्}
{गतेषु रक्षःसु यथानियोगम्}
{व्यदारयद्वानरसागरौघम्}
{महाझषः पूर्ममिवार्णवौघम्} %||6-47-33||

\fourlineindentedshloka
{तमापतन्तं सहसा समीक्ष्य}
{दीप्तेषुचापं युधि राक्षसेन्द्रम्}
{महत्समुत्पाट्य महीधराग्रम्}
{दुद्राव रक्षोऽधिपतिं हरीशः} %||6-47-34||

\fourlineindentedshloka
{तच्छैलशृङ्गं बहुवृक्षसानुम्}
{प्रगृह्य चिक्षेप निशाचराय}
{तमापतन्तं सहसा समीक्ष्य}
{बिभेद बाणैस्तपनीयपुङ्खैः} %||6-47-35||

\fourlineindentedshloka
{तस्मिन्प्रवृद्धोत्तमसानुवृक्षे}
{शृङ्गे विकीर्णे पतिते पृथिव्याम्}
{महाहिकल्पं शरमन्तकाभम्}
{समाददे राक्षसलोकनाथः} %||6-47-36||

\fourlineindentedshloka
{स तं गृहीत्वानिलतुल्यवेगम्}
{सविस्फुलिङ्गज्वलनप्रकाशम्}
{बाणं महेन्द्राशनितुल्यवेगम्}
{चिक्षेप सुग्रीववधाय रुष्टः} %||6-47-37||

\fourlineindentedshloka
{स सायको रावणबाहुमुक्तः}
{शक्राशनिप्रख्यवपुः शिताग्रः}
{सुग्रीवमासाद्य बिभेद वेगाद्-}
{गुहेरिता क्रौचमिवोग्रशक्तिः} %||6-47-38||

\fourlineindentedshloka
{स सायकार्तो विपरीतचेताः}
{कूजन्पृथिव्यां निपपात वीरः}
{तं प्रेक्ष्य भूमौ पतितं विसञ्ज्मम्}
{नेदुः प्रहृष्टा युधि यातुधानाः} %||6-47-39||

\fourlineindentedshloka
{ततो गवाक्षो गवयः सुदंष्ट्रः}
{तथर्षभो ज्योतिमुखो नलश्च}
{शैलान्समुद्यम्य विवृद्धकायाः}
{प्रदुद्रुवुस्तं प्रति राक्षसेन्द्रम्} %||6-47-40||

\fourlineindentedshloka
{तेषां प्रहारान्स चकार मेघान्}
{रक्षोऽधिपो बाणगणैः शिताग्रैः}
{तान्वानरेन्द्रानपि बाणजालैर्-}
{बिभेद जाम्बूनदचित्रपुङ्खैः} %||6-47-41||

\fourlineindentedshloka
{ते वानरेन्द्रास्त्रिदशारिबाणैर्-}
{भिन्ना निपेतुर्भुवि भीमरूपाः}
{ततस्तु तद्वानरसैन्यमुग्रम्}
{प्रच्छादयामास स बाणजालैः} %||6-47-42||

\fourlineindentedshloka
{ते वध्यमानाः पतिताग्र्यवीरा}
{नानद्यमाना भयशल्यविद्धाः}
{शाखामृगा रावणसायकार्ता}
{जग्मुः शरण्यं शरणं स्म रामम्} %||6-47-43||

\fourlineindentedshloka
{ततो महात्मा स धनुर्धनुष्मा}
{नादाय रामः सहरा जगाम}
{तं लक्ष्मणः प्राञ्जलिरभ्युपेत्य}
{उवाच वाक्यं परमार्थयुक्तम्} %||6-47-44||

\twolineshloka
{काममार्यः सुपर्याप्तो वधायास्य दुरात्मनः}
{विधमिष्याम्यहं नीचमनुजानीहि मां विभो} %||6-47-45||

\twolineshloka
{तमब्रवीन्महातेजा रामः सत्यपराक्रमः}
{गच्छ यत्नपरश्चापि भव लक्ष्मण संयुगे} %||6-47-46||

\twolineshloka
{रावणो हि महावीर्यो रणेऽद्भुतपराक्रमः}
{त्रैलोक्येनापि सङ्क्रुद्धो दुष्प्रसह्यो न संशयः} %||6-47-47||

\twolineshloka
{तस्य च्छिद्राणि मार्गस्व स्वच्छिद्राणि च गोपय}
{चक्षुषा धनुषा यत्नाद्रक्षात्मानं समाहितः} %||6-47-48||

\twolineshloka
{राघवस्य वचः श्रुत्वा सम्परिष्वज्य पूज्य च}
{अभिवाद्य ततो रामं ययौ सौमित्रिराहवम्} %||6-47-49||

\fourlineindentedshloka
{स रावणं वारणहस्तबाहुर्-}
{ददर्श दीप्तोद्यतभीमचापम्}
{प्रच्छादयन्तं शरवृष्टिजालैः}
{तान्वानरान्भिन्नविकीर्णदेहान्} %||6-47-50||

\twolineshloka
{तमालोक्य महातेजा हनूमान्मारुतात्मजा}
{निवार्य शरजालानि प्रदुद्राव स रावणम्} %||6-47-51||

\twolineshloka
{रथं तस्य समासाद्य भुजमुद्यम्य दक्षिणम्}
{त्रासयन्रावणं धीमान्हनूमान्वाक्यमब्रवीत्} %||6-47-52||

\twolineshloka
{देवदानवगन्धर्वा यक्षाश्च सह राक्षसैः}
{अवध्यत्वात्त्वया भग्ना वानरेभ्यस्तु ते भयम्} %||6-47-53||

\twolineshloka
{एष मे दक्षिणो बाहुः पञ्चशाखः समुद्यतः}
{विधमिष्यति ते देहाद्भूतात्मानं चिरोषितम्} %||6-47-54||

\twolineshloka
{श्रुत्वा हनूमतो वाक्यं रावणो भीमविक्रमः}
{संरक्तनयनः क्रोधादिदं वचनमब्रवीत्} %||6-47-55||

\twolineshloka
{क्षिप्रं प्रहर निःशङ्कं स्थिरां कीर्तिमवाप्नुहि}
{ततस्त्वां ज्ञातिविक्रान्तं नाशयिष्यामि वानर} %||6-47-56||

\twolineshloka
{रावणस्य वचः श्रुत्वा वायुसूनुर्वचोऽब्रवीत्}
{प्रहृतं हि मया पूर्वमक्षं स्मर सुतं तव} %||6-47-57||

\twolineshloka
{एवमुक्तो महातेजा रावणो राक्षसेश्वरः}
{आजघानानिलसुतं तलेनोरसि वीर्यवान्} %||6-47-58||

\twolineshloka
{स तलाभिहतस्तेन चचाल च मुहुर्मुहुः}
{आजघानाभिसङ्क्रुद्धस्तलेनैवामरद्विषम्} %||6-47-59||

\twolineshloka
{ततस्तलेनाभिहतो वानरेण महात्मना}
{दशग्रीवः समाधूतो यथा भूमिचलेऽचलः} %||6-47-60||

\twolineshloka
{सङ्ग्रामे तं तथा दृष्ट्व रावणं तलताडितम्}
{ऋषयो वानराः सिद्धा नेदुर्देवाः सहासुराः} %||6-47-61||

\twolineshloka
{अथाश्वस्य महातेजा रावणो वाक्यमब्रवीत्}
{साधु वानरवीर्येण श्लाघनीयोऽसि मे रिपुः} %||6-47-62||

\twolineshloka
{रावणेनैवमुक्तस्तु मारुतिर्वाक्यमब्रवीत्}
{धिगस्तु मम वीर्यं तु यत्त्वं जीवसि रावण} %||6-47-63||

\threelineshloka
{सकृत्तु प्रहरेदानीं दुर्बुद्धे किं विकत्थसे}
{ततस्त्वां मामको मुष्टिर्नयिष्यामि यथाक्षयम्}
{ततो मारुतिवाक्येन क्रोधस्तस्य तदाज्वलत्} %||6-47-64||

\threelineshloka
{संरक्तनयनो यत्नान्मुष्टिमुद्यम्य दक्षिणम्}
{पातयामास वेगेन वानरोरसि वीर्यवान्}
{हनूमान्वक्षसि व्यूधे सञ्चचाल हतः पुनः} %||6-47-65||

\twolineshloka
{विह्वलं तं तदा दृष्ट्वा हनूमन्तं महाबलम्}
{रथेनातिरथः शीघ्रं नीलं प्रति समभ्यगात्} %||6-47-66||

\twolineshloka
{पन्नगप्रतिमैर्भीमैः परमर्मातिभेदिभिः}
{शरैरादीपयामास नीलं हरिचमूपतिम्} %||6-47-67||

\twolineshloka
{स शरौघसमायस्तो नीलः कपिचमूपतिः}
{करेणैकेन शैलाग्रं रक्षोऽधिपतयेऽसृजत्} %||6-47-68||

\twolineshloka
{हनूमानपि तेजस्वी समाश्वस्तो महामनाः}
{विप्रेक्षमाणो युद्धेप्सुः सरोषमिदमब्रवीत्} %||6-47-69||

\twolineshloka
{नीलेन सह संयुक्तं रावणं राक्षसेश्वरम्}
{अन्येन युध्यमानस्य न युक्तमभिधावनम्} %||6-47-70||

\twolineshloka
{रावणोऽपि महातेजास्तच्छृङ्गं सप्तभिः शरैः}
{आजघान सुतीक्ष्णाग्रैस्तद्विकीर्णं पपात ह} %||6-47-71||

\twolineshloka
{तद्विकीर्णं गिरेः शृङ्गं दृष्ट्वा हरिचमूपतिः}
{कालाग्निरिव जज्वाल क्रोधेन परवीरहा} %||6-47-72||

\twolineshloka
{सोऽश्वकर्णान्धवान्सालांश्चूतांश्चापि सुपुष्पितान्}
{अन्यांश्च विविधान्वृक्षान्नीलश्चिक्षेप संयुगे} %||6-47-73||

\twolineshloka
{स तान्वृक्षान्समासाद्य प्रतिचिच्छेद रावणः}
{अभ्यवर्षत्सुघोरेण शरवर्षेण पावकिम्} %||6-47-74||

\twolineshloka
{अभिवृष्टः शरौघेण मेघेनेव महाचलः}
{ह्रस्वं कृत्वा तदा रूपं ध्वजाग्रे निपपात ह} %||6-47-75||

\twolineshloka
{पावकात्मजमालोक्य ध्वजाग्रे समवस्थितम्}
{जज्वाल रावणः क्रोधात्ततो नीलो ननाद ह} %||6-47-76||

\twolineshloka
{ध्वजाग्रे धनुषश्चाग्रे किरीटाग्रे च तं हरिम्}
{लक्ष्मणोऽथ हनूमांश्च दृष्ट्वा रामश्च विस्मिताः} %||6-47-77||

\twolineshloka
{रावणोऽपि महातेजाः कपिलाघवविस्मितः}
{अस्त्रमाहारयामास दीप्तमाग्नेयमद्भुतम्} %||6-47-78||

\twolineshloka
{ततस्ते चुक्रुशुर्हृष्टा लब्धलक्ष्याः प्लवङ्गमाः}
{नीललाघवसम्भ्रान्तं दृष्ट्वा रावणमाहवे} %||6-47-79||

\twolineshloka
{वानराणां च नादेन संरब्धो रावणस्तदा}
{सम्भ्रमाविष्टहृदयो न किञ्चित्प्रत्यपद्यत} %||6-47-80||

\twolineshloka
{आग्नेयेनाथ संयुक्तं गृहीत्वा रावणः शरम्}
{ध्वजशीर्षस्थितं नीलमुदैक्षत निशाचरः} %||6-47-81||

\twolineshloka
{ततोऽब्रवीन्महातेजा रावणो राक्षसेश्वरः}
{कपे लाघवयुक्तोऽसि मायया परयानया} %||6-47-82||

\twolineshloka
{जीवितं खलु रक्षस्व यदि शक्नोषि वानर}
{तानि तान्यात्मरूपाणि सृजसे त्वमनेकशः} %||6-47-83||

\twolineshloka
{तथापि त्वां मया मुक्तः सायकोऽस्त्रप्रयोजितः}
{जीवितं परिरक्षन्तं जीविताद्भ्रंशयिष्यति} %||6-47-84||

\twolineshloka
{एवमुक्त्वा महाबाहू रावणो राक्षसेश्वरः}
{सन्धाय बाणमस्त्रेण चमूपतिमताडयत्} %||6-47-85||

\twolineshloka
{सोऽस्त्रयुक्तेन बाणेन नीलो वक्षसि ताडितः}
{निर्दह्यमानः सहसा निपपात महीतले} %||6-47-86||

\twolineshloka
{पितृमाहात्म्य संयोगादात्मनश्चापि तेजसा}
{जानुभ्यामपतद्भूमौ न च प्राणैर्व्ययुज्यत} %||6-47-87||

\twolineshloka
{विसंज्ञं वानरं दृष्ट्वा दशग्रीवो रणोत्सुकः}
{रथेनाम्बुदनादेन सौमित्रिमभिदुद्रुवे} %||6-47-88||

\fourlineindentedshloka
{तमाह सौमित्रिरदीनसत्त्वो}
{विस्फारयन्तं धनुरप्रमेयम्}
{अन्वेहि मामेव निशाचरेन्द्र}
{न वानरांस्त्वं प्रति योद्धुमर्हसि} %||6-47-89||

\fourlineindentedshloka
{स तस्य वाक्यं परिपूर्णघोषम्}
{ज्याशब्दमुग्रं च निशम्य राजा}
{आसाद्य सौमित्रिमवस्थितं तम्}
{कोपान्वितं वाक्यमुवाच रक्षः} %||6-47-90||

\fourlineindentedshloka
{दिष्ट्यासि मे राघव दृष्टिमार्गम्}
{प्राप्तोऽन्तगामी विपरीतबुद्धिः}
{अस्मिन्क्षणे यास्यसि मृत्युदेशम्}
{संसाद्यमानो मम बाणजालैः} %||6-47-91||

\fourlineindentedshloka
{तमाह सौमित्रिरविस्मयानो}
{गर्जन्तमुद्वृत्तसिताग्रदंष्ट्रम्}
{राजन्न गर्जन्ति महाप्रभावा}
{विकत्थसे पापकृतां वरिष्ठ} %||6-47-92||

\fourlineindentedshloka
{जानामि वीर्यं तव राक्षसेन्द्र}
{बलं प्रतापं च पराक्रमं च}
{अवस्थितोऽहं शरचापपाणिः}
{आगच्छ किं मोघविकत्थनेन} %||6-47-93||

\fourlineindentedshloka
{स एवमुक्तः कुपितः ससर्ज}
{रक्षोऽधिपः सप्तशरान्सुपुङ्खान्}
{ताँल्लक्ष्मणः काञ्चनचित्रपुङ्खैः}
{चिच्छेद बाणैर्निशिताग्रधारैः} %||6-47-94||

\fourlineindentedshloka
{तान्प्रेक्षमाणः सहसा निकृत्तान्}
{निकृत्तभोगानिव पन्नगेन्द्रान्}
{लङ्केश्वरः क्रोधवशं जगाम}
{ससर्ज चान्यान्निशितान्पृषत्कान्} %||6-47-95||

\fourlineindentedshloka
{स बाणवर्षं तु ववर्ष तीव्रम्}
{रामानुजः कार्मुकसम्प्रयुक्तम्}
{क्षुरार्धचन्द्रोत्तमकर्णिभल्लैः}
{शरांश्च चिच्छेद न चुक्षुभे च} %||6-47-96||

\fourlineindentedshloka
{स लक्ष्मणश्चाशु शराञ्शिताग्रान्}
{महेन्द्रवज्राशनितुल्यवेगान्}
{सन्धाय चापे ज्वलनप्रकाशान्}
{ससर्ज रक्षोऽधिपतेर्वधाय} %||6-47-97||

\fourlineindentedshloka
{स तान्प्रचिच्छेद हि राक्षसेन्द्रः}
{छित्त्वा च ताँल्लक्ष्मणमाजघान}
{शरेण कालाग्निसमप्रभेण}
{स्वयम्भुदत्तेन ललाटदेशे} %||6-47-98||

\fourlineindentedshloka
{स लक्ष्मणो रावणसायकार्तः}
{चचाल चापं शिथिलं प्रगृह्य}
{पुनश्च संज्ञां प्रतिलभ्य कृच्छ्रात्}
{चिच्छेद चापं त्रिदशेन्द्रशत्रोः} %||6-47-99||

\fourlineindentedshloka
{निकृत्तचापं त्रिभिराजघान}
{बाणैस्तदा दाशरथिः शिताग्रैः}
{स सायकार्तो विचचाल राजा}
{कृच्छ्राच्च संज्ञां पुनराससाद} %||6-47-100||

\fourlineindentedshloka
{स कृत्तचापः शरताडितश्च}
{स्वेदार्द्रगात्रो रुधिरावसिक्तः}
{जग्राह शक्तिं समुदग्रशक्तिः}
{स्वयम्भुदत्तां युधि देवशत्रुः} %||6-47-101||

\fourlineindentedshloka
{स तां विधूमानलसंनिकाशाम्}
{वित्रासनीं वानरवाहिनीनाम्}
{चिक्षेप शक्तिं तरसा ज्वलन्तीम्}
{सौमित्रये राक्षसराष्ट्रनाथः} %||6-47-102||

\fourlineindentedshloka
{तामापतन्तीं भरतानुजोऽस्त्रैर्-}
{जघान बाणैश्च हुताग्निकल्पैः}
{तथापि सा तस्य विवेश शक्तिर्-}
{भुजान्तरं दाशरथेर्विशालम्} %||6-47-103||

\twolineshloka
{शक्त्या ब्राम्या तु सौमित्रिस्ताडितस्तु स्तनान्तरे}
{विष्णोरचिन्त्यं स्वं भागमात्मानं प्रत्यनुस्मरत्} %||6-47-104||

\twolineshloka
{ततो दानवदर्पघ्नं सौमित्रिं देवकण्टकः}
{तं पीडयित्वा बाहुभ्यामप्रभुर्लङ्घनेऽभवत्} %||6-47-105||

\twolineshloka
{हिमवान्मन्दरो मेरुस्त्रैलोक्यं वा सहामरैः}
{शक्यं भुजाभ्यामुद्धर्तुं न सङ्ख्ये भरतानुजः} %||6-47-106||

\twolineshloka
{अथैनं वैष्णवं भागं मानुषं देहमास्थितम्}
{विसंज्ञं लक्ष्मणं दृष्ट्वा रावणो विस्मितोऽभवत्} %||6-47-107||

\twolineshloka
{अथ वायुसुतः क्रुद्धो रावणं समभिद्रवत्}
{आजघानोरसि क्रुद्धो वज्रकल्पेन मुष्टिना} %||6-47-108||

\twolineshloka
{तेन मुष्टिप्रहारेण रावणो राक्षसेश्वरः}
{जानुभ्यामपतद्भूमौ चचाल च पपात च} %||6-47-109||

\twolineshloka
{विसंज्ञं रावणं दृष्ट्वा समरे भीमविक्रमम्}
{ऋषयो वानराश्चैव नेदुर्देवाः सवासवाः} %||6-47-110||

\twolineshloka
{हनूमानपि तेजस्वी लक्ष्मणं रावणार्दितम्}
{अनयद्राघवाभ्याशं बाहुभ्यां परिगृह्य तम्} %||6-47-111||

\twolineshloka
{वायुसूनोः सुहृत्त्वेन भक्त्या परमया च सः}
{शत्रूणामप्रकम्प्योऽपि लघुत्वमगमत्कपेः} %||6-47-112||

\twolineshloka
{तं समुत्सृज्य सा शक्तिः सौमित्रिं युधि दुर्जयम्}
{रावणस्य रथे तस्मिन्स्थानं पुनरुपागमत्} %||6-47-113||

\twolineshloka
{रावणोऽपि महातेजाः प्राप्य संज्ञां महाहवे}
{आददे निशितान्बाणाञ्जग्राह च महद्धनुः} %||6-47-114||

\twolineshloka
{आश्वस्तश्च विशल्यश्च लक्ष्मणः शत्रुसूदनः}
{विष्णोर्भागममीमांस्यमात्मानं प्रत्यनुस्मरन्} %||6-47-115||

\twolineshloka
{निपातितमहावीरां वानराणां महाचमूम्}
{राघवस्तु रणे दृष्ट्वा रावणं समभिद्रवत्} %||6-47-116||

\twolineshloka
{अथैनमुपसङ्गम्य हनूमान्वाक्यमब्रवीत्}
{मम पृष्ठं समारुह्य रक्षसं शास्तुमर्हसि} %||6-47-117||

\threelineshloka
{तच्छ्रुत्वा राघवो वाक्यं वायुपुत्रेण भाषितम्}
{आरोहत्सहसा शूरो हनूमन्तं महाकपिम्}
{रथस्थं रावणं सङ्ख्ये ददर्श मनुजाधिपः} %||6-47-118||

\twolineshloka
{तमालोक्य महातेजाः प्रदुद्राव स राघवः}
{वैरोचनमिव क्रुद्धो विष्णुरभ्युद्यतायुधः} %||6-47-119||

\twolineshloka
{ज्याशब्दमकरोत्तीव्रं वज्रनिष्पेषनिस्वनम्}
{गिरा गम्भीरया रामो राक्षसेन्द्रमुवाच ह} %||6-47-120||

\twolineshloka
{तिष्ठ तिष्ठ मम त्वं हि कृत्वा विप्रियमीदृशम्}
{क्व नु राक्षसशार्दूल गतो मोक्षमवाप्स्यसि} %||6-47-121||

\fourlineindentedshloka
{यदीन्द्रवैवस्वत भास्करान्वा}
{स्वयम्भुवैश्वानरशङ्करान्वा}
{गमिष्यसि त्वं दश वा दिशो वा}
{तथापि मे नाद्य गतो विमोक्ष्यसे} %||6-47-122||

\fourlineindentedshloka
{यश्चैष शक्त्याभिहतस्त्वयाद्य}
{इच्छन्विषादं सहसाभ्युपेतः}
{स एष रक्षोगणराज मृत्युः}
{सपुत्रदारस्य तवाद्य युद्धे} %||6-47-123||

\twolineshloka
{राघवस्य वचः श्रुत्वा राक्षसेन्द्रो महाकपिम्}
{आजघान शरैस्तीक्ष्णैः कालानलशिखोपमैः} %||6-47-124||

\twolineshloka
{राक्षसेनाहवे तस्य ताडितस्यापि सायकैः}
{स्वभावतेजोयुक्तस्य भूयस्तेजो व्यवर्धत} %||6-47-125||

\twolineshloka
{ततो रामो महातेजा रावणेन कृतव्रणम्}
{दृष्ट्वा प्लवगशार्दूलं क्रोधस्य वशमेयिवान्} %||6-47-126||

\fourlineindentedshloka
{तस्याभिसङ्क्रम्य रथं सचक्रम्}
{साश्वध्वजच्छत्रमहापताकम्}
{ससारथिं साशनिशूलखड्गम्}
{रामः प्रचिच्छेद शरैः सुपुङ्खैः} %||6-47-127||

\fourlineindentedshloka
{अथेन्द्रशत्रुं तरसा जघान}
{बाणेन वज्राशनिसंनिभेन}
{भुजान्तरे व्यूढसुजातरूपे}
{वज्रेण मेरुं भगवानिवेन्द्रः} %||6-47-128||

\fourlineindentedshloka
{यो वज्रपाताशनिसंनिपातान्}
{न चुक्षुभे नापि चचाल राजा}
{स रामबाणाभिहतो भृशार्तः}
{चचाल चापं च मुमोच वीरः} %||6-47-129||

\fourlineindentedshloka
{तं विह्वलन्तं प्रसमीक्ष्य रामः}
{समाददे दीप्तमथार्धचन्द्रम्}
{तेनार्कवर्णं सहसा किरीटम्}
{चिच्छेद रक्षोऽधिपतेर्महात्माः} %||6-47-130||

\fourlineindentedshloka
{तं निर्विषाशीविषसंनिकाशम्}
{शान्तार्चिषं सूर्यमिवाप्रकाशम्}
{गतश्रियं कृत्तकिरीटकूटम्}
{उवाच रामो युधि राक्षसेन्द्रम्} %||6-47-131||

\fourlineindentedshloka
{कृतं त्वया कर्म महत्सुभीमम्}
{हतप्रवीरश्च कृतस्त्वयाहम्}
{तस्मात्परिश्रान्त इति व्यवस्य}
{न त्वं शरैर्मृत्युवशं नयामि} %||6-47-132||

\fourlineindentedshloka
{स एवमुक्तो हतदर्पहर्षो}
{निकृत्तचापः स हताश्वसूतः}
{शरार्दितः कृत्तमहाकिरीटो}
{विवेश लङ्कां सहसा स्म राजा} %||6-47-133||

\fourlineindentedshloka
{तस्मिन्प्रविष्टे रजनीचरेन्द्रे}
{महाबले दानवदेवशत्रौ}
{हरीन्विशल्यान्सहलक्ष्मणेन}
{चकार रामः परमाहवाग्रे} %||6-47-134||

\fourlineindentedshloka
{तस्मिन्प्रभग्ने त्रिदशेन्द्रशत्रौ}
{सुरासुरा भूतगणा दिशश्च}
{ससागराः सर्षिमहोरगाश्च}
{तथैव भूम्यम्बुचराश्च हृष्टाः} %||6-48-135||


{॥इत्यार्षे श्रीमद्रामायणे वाल्मीकीये आदिकाव्ये युद्धकाण्डे सप्तचत्वारिंशः सर्गः॥}

\sect{अष्टचत्वारिंशः सर्गः}

\twolineshloka
{स प्रविश्य पुरीं लङ्कां रामबाणभयार्दितः}
{भग्नदर्पस्तदा राजा बभूव व्यथितेन्द्रियः} %||6-48-1||

\twolineshloka
{मातङ्ग इव सिंहेन गरुडेनेव पन्नगः}
{अभिभूतोऽभवद्राजा राघवेण महात्मना} %||6-48-2||

\twolineshloka
{ब्रह्मदण्डप्रकाशानां विद्युत्सदृशवर्चसाम्}
{स्मरन्राघवबाणानां विव्यथे राक्षसेश्वरः} %||6-48-3||

\twolineshloka
{स काञ्चनमयं दिव्यमाश्रित्य परमासनम्}
{विक्प्रेक्षमाणो रक्षांसि रावणो वाक्यमब्रवीत्} %||6-48-4||

\twolineshloka
{सर्वं तत्खलु मे मोघं यत्तप्तं परमं तपः}
{यत्समानो महेन्द्रेण मानुषेणास्मि निर्जितः} %||6-48-5||

\twolineshloka
{इदं तद्ब्रह्मणो घोरं वाक्यं मामभ्युपस्थितम्}
{मानुषेभ्यो विजानीहि भयं त्वमिति तत्तथा} %||6-48-6||

\twolineshloka
{देवदानवगन्धर्वैर्यक्षराक्षसपन्नगैः}
{अवध्यत्वं मया प्राप्तं मानुषेभ्यो न याचितम्} %||6-48-7||

\twolineshloka
{एतदेवाभ्युपागम्य यत्नं कर्तुमिहार्हथ}
{राक्षसाश्चापि तिष्ठन्तु चर्यागोपुरमूर्धसु} %||6-48-8||

\twolineshloka
{स चाप्रतिमगम्भीरो देवदानवदर्पहा}
{ब्रह्मशापाभिभूतस्तु कुम्भकर्णो विबोध्यताम्} %||6-48-9||

\twolineshloka
{स पराजितमात्मानं प्रहस्तं च निषूदितम्}
{ज्ञात्वा रक्षोबलं भीममादिदेश महाबलः} %||6-48-10||

\twolineshloka
{द्वारेषु यत्नः क्रियतां प्राकाराश्चाधिरुह्यताम्}
{निद्रावशसमाविष्टः कुम्भकर्णो विबोध्यताम्} %||6-48-11||

\twolineshloka
{नव षट्सप्त चाष्टौ च मासान्स्वपिति राक्षसः}
{तं तु बोधयत क्षिप्रं कुम्भकर्णं महाबलम्} %||6-48-12||

\twolineshloka
{स हि सङ्ख्ये महाबाहुः ककुदं सर्वरक्षसाम्}
{वानरान्राजपुत्रौ च क्षिप्रमेव वधिष्यति} %||6-48-13||

\twolineshloka
{कुम्भकर्णः सदा शेते मूढो ग्राम्यसुखे रतः}
{रामेणाभिनिरस्तस्य सङ्ग्रामोऽस्मिन्सुदारुणे} %||6-48-14||

\twolineshloka
{भविष्यति न मे शोकः कुम्भकर्णे विबोधिते}
{किं करिष्याम्यहं तेन शक्रतुल्यबलेन हि} %||6-48-15||

\twolineshloka
{ईदृशे व्यसने प्राप्ते यो न साह्याय कल्पते}
{ते तु तद्वचनं श्रुत्वा राक्षसेन्द्रस्य राक्षसाः} %||6-48-16||

\twolineshloka
{जग्मुः परमसम्भ्रान्ताः कुम्भकर्णनिवेशनम्}
{ते रावणसमादिष्टा मांसशोणितभोजनाः} %||6-48-17||

\twolineshloka
{गन्धमाल्यांस्तथा भक्ष्यानादाय सहसा ययुः}
{तां प्रविश्य महाद्वारां सर्वतो योजनायताम्} %||6-48-18||

\twolineshloka
{कुम्भकर्णगुहां रम्यां सर्वगन्धप्रवाहिनीम्}
{प्रतिष्ठमानाः कृच्छ्रेण यत्नात्प्रविविशुर्गुहाम्} %||6-48-19||

\twolineshloka
{तां प्रविश्य गुहां रम्यां शुभां काञ्चनकुट्टिमाम्}
{ददृशुर्नैरृतव्याघ्रं शयानं भीमदर्शनम्} %||6-48-20||

\twolineshloka
{ते तु तं विकृतं सुप्तं विकीर्णमिव पर्वतम्}
{कुम्भकर्णं महानिद्रं सहिताः प्रत्यबोधयन्} %||6-48-21||

\twolineshloka
{ऊर्ध्वरोमाञ्चिततनुं श्वसन्तमिव पन्नगम्}
{त्रासयन्तं महाश्वासैः शयानं भीमदर्शनम्} %||6-48-22||

\twolineshloka
{भीमनासापुटं तं तु पातालविपुलाननम्}
{ददृशुर्नैरृतव्याघ्रं कुम्भकर्णं महाबलम्} %||6-48-23||

\twolineshloka
{ततश्चक्रुर्महात्मानः कुम्भकर्णाग्रतस्तदा}
{मांसानां मेरुसङ्काशं राशिं परमतर्पणम्} %||6-48-24||

\twolineshloka
{मृगाणां महिषाणां च वराहाणां च सञ्चयान्}
{चक्रुर्नैरृतशार्दूला राशिमन्नस्य चाद्भुतम्} %||6-48-25||

\twolineshloka
{ततः शोणितकुम्भांश्च मद्यानि विविधानि च}
{पुरस्तात्कुम्भकर्णस्य चक्रुस्त्रिदशशत्रवः} %||6-48-26||

\twolineshloka
{लिलिपुश्च परार्ध्येन चन्दनेन परन्तपम्}
{दिव्यैराच्छादयामासुर्माल्यैर्गन्धैः सुगन्धिभिः} %||6-48-27||

\twolineshloka
{धूपं सुगन्धं ससृजुस्तुष्टुवुश्च परन्तपम्}
{जलदा इव चोनेदुर्यातुधानाः सहस्रशः} %||6-48-28||

\twolineshloka
{शङ्खानापूरयामासुः शशाङ्कसदृशप्रभान्}
{तुमुलं युगपच्चापि विनेदुश्चाप्यमर्षिताः} %||6-48-29||

\twolineshloka
{नेदुरास्फोटयामासुश्चिक्षिपुस्ते निशाचराः}
{कुम्भकर्णविबोधार्थं चक्रुस्ते विपुलं स्वनम्} %||6-48-30||

\fourlineindentedshloka
{सशङ्खभेरीपटहप्रणादम्}
{आस्फोटितक्ष्वेडितसिंहनादम्}
{दिशो द्रवन्तस्त्रिदिवं किरन्तः}
{श्रुत्वा विहङ्गाः सहसा निपेतुः} %||6-48-31||

\fourlineindentedshloka
{यदा भृशं तैर्निनदैर्महात्मा}
{न कुम्भकर्णो बुबुधे प्रसुप्तः}
{ततो मुसुण्डीमुसलानि सर्वे}
{रक्षोगणास्ते जगृहुर्गदाश्च} %||6-48-32||

\fourlineindentedshloka
{तं शैलशृङ्गैर्मुसलैर्गदाभिर्-}
{वृक्षैस्तलैर्मुद्गरमुष्टिभिश्च}
{सुखप्रसुप्तं भुवि कुम्भकर्णम्}
{रक्षांस्युदग्राणि तदा निजघ्नुः} %||6-48-33||

\twolineshloka
{तस्य निश्वासवातेन कुम्भकर्णस्य रक्षसः}
{राक्षसा बलवन्तोऽपि स्थातुं नाशक्नुवन्पुरः} %||6-48-34||

\threelineshloka
{ततोऽस्य पुरतो गाढं राक्षसा भीमविक्रमाः}
{मृदङ्गपणवान्भेरीः शङ्खकुम्भगणांस्तथा}
{दशराक्षससाहस्रं युगपत्पर्यवादयन्} %||6-48-35||

\twolineshloka
{नीलाञ्जनचयाकारं ते तु तं प्रत्यबोधयन्}
{अभिघ्नन्तो नदन्तश्च नैव संविविदे तु सः} %||6-48-36||

\twolineshloka
{यदा चैनं न शेकुस्ते प्रतिबोधयितुं तदा}
{ततो गुरुतरं यत्नं दारुणं समुपाक्रमन्} %||6-48-37||

\twolineshloka
{अश्वानुष्ट्रान्खरान्नागाञ्जघ्नुर्दण्डकशाङ्कुशैः}
{भेरीशङ्खमृदङ्गांश्च सर्वप्राणैरवादयन्} %||6-48-38||

\twolineshloka
{निजघ्नुश्चास्य गात्राणि महाकाष्ठकटं करैः}
{मुद्गरैर्मुसलैश्चैव सर्वप्राणसमुद्यतैः} %||6-48-39||

\twolineshloka
{तेन शब्देन महता लङ्का समभिपूरिता}
{सपर्वतवना सर्वा सोऽपि नैव प्रबुध्यते} %||6-48-40||

\twolineshloka
{ततः सहस्रं भेरीणां युगपत्समहन्यत}
{मृष्टकाञ्चनकोणानामसक्तानां समन्ततः} %||6-48-41||

\twolineshloka
{एवमप्यतिनिद्रस्तु यदा नैव प्रबुध्यत}
{शापस्य वशमापन्नस्ततः क्रुद्धा निशाचराः} %||6-48-42||

\twolineshloka
{महाक्रोधसमाविष्टाः सर्वे भीमपराक्रमाः}
{तद्रक्षोबोधयिष्यन्तश्चक्रुरन्ये पराक्रमम्} %||6-48-43||

\threelineshloka
{अन्ये भेरीः समाजघ्नुरन्ये चक्रुर्महास्वनम्}
{केशानन्ये प्रलुलुपुः कर्णावन्ये दशन्ति च}
{न कुम्भकर्णः पस्पन्दे महानिद्रावशं गतः} %||6-48-44||

\twolineshloka
{अन्ये च बलिनस्तस्य कूटमुद्गरपाणयः}
{मूर्ध्नि वक्षसि गात्रेषु पातयन्कूटमुद्गरान्} %||6-48-45||

\twolineshloka
{रज्जुबन्धनबद्धाभिः शतघ्नीभिश्च सर्वतः}
{वध्यमानो महाकायो न प्राबुध्यत राक्षसः} %||6-48-46||

\twolineshloka
{वारणानां सहस्रं तु शरीरेऽस्य प्रधावितम्}
{कुम्भकर्णस्ततो बुद्धः स्पर्शं परमबुध्यत} %||6-48-47||

\fourlineindentedshloka
{स पात्यमानैर्गिरिशृङ्गवृक्षैः}
{अचिन्तयंस्तान्विपुलान्प्रहारान्}
{निद्राक्षयात्क्षुद्भयपीडितश्च}
{विजृम्भमाणः सहसोत्पपात} %||6-48-48||

\fourlineindentedshloka
{स नागभोगाचलशृङ्गकल्पौ}
{विक्षिप्य बाहू गिरिशृङ्गसारौ}
{विवृत्य वक्त्रं वडवामुखाभम्}
{निशाचरोऽसौ विकृतं जजृम्भे} %||6-48-49||

\twolineshloka
{तस्य जाजृम्भमाणस्य वक्त्रं पातालसंनिभम्}
{ददृशे मेरुशृङ्गाग्रे दिवाकर इवोदितः} %||6-48-50||

\twolineshloka
{विजृम्भमाणोऽतिबलः प्रतिबुद्धो निशाचरः}
{निश्वासश्चास्य संजज्ञे पर्वतादिव मारुतः} %||6-48-51||

\twolineshloka
{रूपमुत्तिष्ठतस्तस्य कुम्भकर्णस्य तद्बभौ}
{तपान्ते सबलाकस्य मेघस्येव विवर्षतः} %||6-48-52||

\twolineshloka
{तस्य दीप्ताग्निसदृशे विद्युत्सदृशवर्चसी}
{ददृशाते महानेत्रे दीप्ताविव महाग्रहौ} %||6-48-53||

\twolineshloka
{आदद्बुभुक्षितो मांसं शोणितं तृषितोऽपिबत्}
{मेदः कुम्भं च मद्यं च पपौ शक्ररिपुस्तदा} %||6-48-54||

\twolineshloka
{ततस्तृप्त इति ज्ञात्वा समुत्पेतुर्निशाचराः}
{शिरोभिश्च प्रणम्यैनं सर्वतः पर्यवारयन्} %||6-48-55||

\twolineshloka
{स सर्वान्सान्त्वयामास नैरृतान्नैरृतर्षभः}
{बोधनाद्विस्मितश्चापि राक्षसानिदमब्रवीत्} %||6-48-56||

\twolineshloka
{किमर्थमहमाहत्य भवद्भिः प्रतिबोधितः}
{कच्चित्सुकुशलं राज्ञो भयं वा नेह किञ्चन} %||6-48-57||

\twolineshloka
{अथ वा ध्रुवमन्येभ्यो भयं परमुपस्थितम्}
{यदर्थमेव त्वरितैर्भवद्भिः प्रतिबोधितः} %||6-48-58||

\twolineshloka
{अद्य राक्षसराजस्य भयमुत्पाटयाम्यहम्}
{पातयिष्ये महेन्द्रं वा शातयिष्ये तथानलम्} %||6-48-59||

\twolineshloka
{न ह्यल्पकारणे सुप्तं बोधयिष्यति मां भृशम्}
{तदाख्यातार्थतत्त्वेन मत्प्रबोधनकारणम्} %||6-48-60||

\twolineshloka
{एवं ब्रुवाणं संरब्धं कुम्भकर्णमरिन्दमम्}
{यूपाक्षः सचिवो राज्ञः कृताञ्जलिरुवाच ह} %||6-48-61||

\threelineshloka
{न नो देवकृतं किञ्चिद्भयमस्ति कदाचन}
{न दैत्यदानवेभ्यो वा भयमस्ति हि तादृशम्}
{यादृशं मानुषं राजन्भयमस्मानुपस्थितम्} %||6-48-62||

\twolineshloka
{वानरैः पर्वताकारैर्लङ्केयं परिवारिता}
{सीताहरणसन्तप्ताद्रामान्नस्तुमुलं भयम्} %||6-48-63||

\twolineshloka
{एकेन वानरेणेयं पूर्वं दग्धा महापुरी}
{कुमारो निहतश्चाक्षः सानुयात्रः सकुञ्जरः} %||6-48-64||

\twolineshloka
{स्वयं रक्षोऽधिपश्चापि पौलस्त्यो देवकण्टकः}
{मृतेति संयुगे मुक्तारामेणादित्यतेजसा} %||6-48-65||

\twolineshloka
{यन्न देवैः कृतो राजा नापि दैत्यैर्न दानवैः}
{कृतः स इह रामेण विमुक्तः प्राणसंशयात्} %||6-48-66||

\twolineshloka
{स यूपाक्षवचः श्रुत्वा भ्रातुर्युधि पराजयम्}
{कुम्भकर्णो विवृत्ताक्षो यूपाक्षमिदमब्रवीत्} %||6-48-67||

\twolineshloka
{सर्वमद्यैव यूपाक्ष हरिसैन्यं सलक्ष्मणम्}
{राघवं च रणे हत्वा पश्चाद्द्रक्ष्यामि रावणम्} %||6-48-68||

\twolineshloka
{राक्षसांस्तर्पयिष्यामि हरीणां मांसशोणितैः}
{रामलक्ष्मणयोश्चापि स्वयं पास्यामि शोणितम्} %||6-48-69||

\fourlineindentedshloka
{तत्तस्य वाक्यं ब्रुवतो निशम्य}
{सगर्वितं रोषविवृद्धदोषम्}
{महोदरो नैरृतयोधमुख्यः}
{कृताञ्जलिर्वाक्यमिदं बभाषे} %||6-48-70||

\twolineshloka
{रावणस्य वचः श्रुत्वा गुणदोषु विमृश्य च}
{पश्चादपि महाबाहो शत्रून्युधि विजेष्यसि} %||6-48-71||

\twolineshloka
{महोदरवचः श्रुत्वा राक्षसैः परिवारितः}
{कुम्भकर्णो महातेजाः सम्प्रतस्थे महाबलः} %||6-48-72||

\twolineshloka
{तं समुत्थाप्य भीमाक्षं भीमरूपपराक्रमम्}
{राक्षसास्त्वरिता जग्मुर्दशग्रीवनिवेशनम्} %||6-48-73||

\twolineshloka
{ततो गत्वा दशग्रीवमासीनं परमासने}
{ऊचुर्बद्धाञ्जलिपुटाः सर्व एव निशाचराः} %||6-48-74||

\twolineshloka
{प्रबुद्धः कुम्भकर्णोऽसौ भ्राता ते राक्षसर्षभ}
{कथं तत्रैव निर्यातु द्रक्ष्यसे तमिहागतम्} %||6-48-75||

\twolineshloka
{रावणस्त्वब्रवीद्धृष्टो यथान्यायं च पूजितम्}
{द्रष्टुमेनमिहेच्छामि यथान्यायं च पूजितम्} %||6-48-76||

\twolineshloka
{तथेत्युक्त्वा तु ते सर्वे पुनरागम्य राक्षसाः}
{कुम्भकर्णमिदं वाक्यमूचू रावणचोदिताः} %||6-48-77||

\twolineshloka
{द्रष्टुं त्वां काङ्क्षते राजा सर्वराक्षसपुङ्गवः}
{गमने क्रियतां बुद्धिर्भ्रातरं सम्प्रहर्षय} %||6-48-78||

\twolineshloka
{कुम्भकर्णस्तु दुर्धर्षो भ्रातुराज्ञाय शासनम्}
{तथेत्युक्त्वा महावीर्यः शयनादुत्पपात ह} %||6-48-79||

\twolineshloka
{प्रक्षाल्य वदनं हृष्टः स्नातः परमभूषितः}
{पिपासुस्त्वरयामास पानं बलसमीरणम्} %||6-48-80||

\twolineshloka
{ततस्ते त्वरितास्तस्य राज्षसा रावणाज्ञया}
{मद्यं भक्ष्यांश्च विविधान्क्षिप्रमेवोपहारयन्} %||6-48-81||

{पीत्वा घटसहस्रं स गमनायोपचक्रमे॥\devanumber{82}॥} %||6-48-82|| (Check)

\addtocounter{shlokacount}{1}

\twolineshloka
{ईषत्समुत्कटो मत्तस्तेजोबलसमन्वितः}
{कुम्भकर्णो बभौ हृष्टः कालान्तकयमोपमः} %||6-48-83||

\twolineshloka
{भ्रातुः स भवनं गच्छन्रक्षोबलसमन्वितः}
{कुम्भकर्णः पदन्यासैरकम्पयत मेदिनीम्} %||6-48-84||

\fourlineindentedshloka
{स राजमार्गं वपुषा प्रकाशयन्}
{सहस्ररश्मिर्धरणीमिवांशुभिः}
{जगाम तत्राञ्जलिमालया वृतः}
{शतक्रतुर्गेहमिव स्वयम्भुवः} %||6-48-85||

\fourlineindentedshloka
{केचिच्छरण्यं शरणं स्म रामम्}
{व्रजन्ति केचिद्व्यथिताः पतन्ति}
{केचिद्दिशः स्म व्यथिताः प्रयान्ति}
{केचिद्भयार्ता भुवि शेरते स्म} %||6-48-86||

\fourlineindentedshloka
{तमद्रिशृङ्गप्रतिमं किरीटिनम्}
{स्पृशन्तमादित्यमिवात्मतेजसा}
{वनौकसः प्रेक्ष्य विवृद्धमद्भुतम्}
{भयार्दिता दुद्रुविरे ततस्ततः} %||6-49-87||


{॥इत्यार्षे श्रीमद्रामायणे वाल्मीकीये आदिकाव्ये युद्धकाण्डे अष्टचत्वारिंशः सर्गः॥}

\sect{एकोनपञ्चाशत्तमः सर्गः}

\twolineshloka
{ततो रामो महातेजा धनुरादाय वीर्यवान्}
{किरीटिनं महाकायं कुम्भकर्णं ददर्श ह} %||6-49-1||

\twolineshloka
{तं दृष्ट्वा राक्षसश्रेष्ठं पर्वताकारदर्शनम्}
{क्रममाणमिवाकाशं पुरा नारायणं प्रभुम्} %||6-49-2||

\twolineshloka
{सतोयाम्बुदसङ्काशं काञ्चनाङ्गदभूषणम्}
{दृष्ट्वा पुनः प्रदुद्राव वानराणां महाचमूः} %||6-49-3||

\twolineshloka
{विद्रुतां वाहिनीं दृष्ट्वा वर्धमानं च राक्षसम्}
{सविस्मयमिदं रामो विभीषणमुवाच ह} %||6-49-4||

\twolineshloka
{कोऽसौ पर्वतसङ्कशः किरीटी हरिलोचनः}
{लङ्कायां दृश्यते वीरः सविद्युदिव तोयदः} %||6-49-5||

\twolineshloka
{पृथिव्याः केतुभूतोऽसौ महानेकोऽत्र दृश्यते}
{यं दृष्ट्वा वानराः सर्वे विद्रवन्ति ततस्ततः} %||6-49-6||

\twolineshloka
{आचक्ष्व मे महान्कोऽसौ रक्षो वा यदि वासुरः}
{न मयैवंविधं भूतं दृष्टपूर्वं कदाचन} %||6-49-7||

\twolineshloka
{स पृष्टो राजपुत्रेण रामेणाक्लिष्टकारिणा}
{विभीषणो महाप्राज्ञः काकुत्स्थमिदमब्रवीत्} %||6-49-8||

\twolineshloka
{येन वैवस्वतो युद्धे वासवश्च पराजितः}
{सैष विश्रवसः पुत्रः कुम्भकर्णः प्रतापवान्} %||6-49-9||

\fourlineindentedshloka
{एतेन देवा युधि दानवाश्च}
{यक्षा भुजङ्गाः पिशिताशनाश्च}
{गन्धर्वविद्याधरकिंनराश्च}
{सहस्रशो राघव सम्प्रभग्नाः} %||6-49-10||

\twolineshloka
{शूलपाणिं विरूपाक्षं कुम्भकर्णं महाबलम्}
{हन्तुं न शेकुस्त्रिदशाः कालोऽयमिति मोहिताः} %||6-49-11||

\twolineshloka
{प्रकृत्या ह्येष तेजस्वी कुम्भकर्णो महाबलः}
{अन्येषां राक्षसेन्द्राणां वरदानकृतं बलम्} %||6-49-12||

\twolineshloka
{एतेन जातमात्रेण क्षुधार्तेन महात्मना}
{भक्षितानि सहस्राणि सत्त्वानां सुबहून्यपि} %||6-49-13||

\twolineshloka
{तेषु सम्भक्ष्यमाणेषु प्रजा भयनिपीडिताः}
{यान्ति स्म शरणं शक्रं तमप्यर्थं न्यवेदयन्} %||6-49-14||

\fourlineindentedshloka
{स कुम्भकर्णं कुपितो महेन्द्रो}
{जघान वज्रेण शितेन वज्री}
{स शक्रवज्राभिहतो महात्मा}
{चचाल कोपाच्च भृशं ननाद} %||6-49-15||

\twolineshloka
{तस्य नानद्यमानस्य कुम्भकर्णस्य धीमतः}
{श्रुत्वा निनादं वित्रस्ता भूयो भूमिर्वितत्रसे} %||6-49-16||

\twolineshloka
{ततः कोपान्महेन्द्रस्य कुम्भकर्णो महाबलः}
{विकृष्यैरावताद्दन्तं जघानोरसि वासवम्} %||6-49-17||

\twolineshloka
{कुम्भकर्णप्रहारार्तो विचचाल स वासवः}
{ततो विषेदुः सहसा देवब्रह्मर्षिदानवाः} %||6-49-18||

\threelineshloka
{प्रजाभिः सह शक्रश्च ययौ स्थानं स्वयम्भुवः}
{कुम्भकर्णस्य दौरात्म्यं शशंसुस्ते प्रजापतेः}
{प्रजानां भक्षणं चापि देवानां चापि धर्षणम्} %||6-49-19||

\twolineshloka
{एवं प्रजा यदि त्वेष भक्षयिष्यति नित्यशः}
{अचिरेणैव कालेन शून्यो लोको भविष्यति} %||6-49-20||

\twolineshloka
{वासवस्य वचः श्रुत्वा सर्वलोकपितामहः}
{रक्षांस्यावाहयामास कुम्भकर्णं ददर्श ह} %||6-49-21||

\twolineshloka
{कुम्भकर्णं समीक्ष्यैव वितत्रास प्रजापतिः}
{दृष्ट्वा निश्वस्य चैवेदं स्वयम्भूरिदमब्रवीत्} %||6-49-22||

\threelineshloka
{ध्रुवं लोकविनाशाय पौरस्त्येनासि निर्मितः}
{तस्मात्त्वमद्य प्रभृति मृतकल्पः शयिष्यसि}
{ब्रह्मशापाभिभूतोऽथ निपपाताग्रतः प्रभोः} %||6-49-23||

\twolineshloka
{ततः परमसम्भ्रान्तो रावणो वाक्यमब्रवीत्}
{विवृद्धः काञ्चनो वृक्षः फलकाले निकृत्यते} %||6-49-24||

\threelineshloka
{न नप्तारं स्वकं न्याय्यं शप्तुमेवं प्रजापते}
{न मिथ्यावचनश्च त्वं स्वप्स्यत्येष न संशयः}
{कालस्तु क्रियतामस्य शयने जागरे तथा} %||6-49-25||

\twolineshloka
{रावणस्य वचः श्रुत्वा स्वयम्भूरिदमब्रवीत्}
{शयिता ह्येष षण्मासानेकाहं जागरिष्यति} %||6-49-26||

\twolineshloka
{एकेनाह्ना त्वसौ वीरश्चरन्भूमिं बुभुक्षितः}
{व्यात्तास्यो भक्षयेल्लोकान्सङ्क्रुद्ध इव पावकः} %||6-49-27||

\twolineshloka
{सोऽसौ व्यसनमापन्नः कुम्भकर्णमबोधयत्}
{त्वत्पराक्रमभीतश्च राजा सम्प्रति रावणः} %||6-49-28||

\twolineshloka
{स एष निर्गतो वीरः शिबिराद्भीमविक्रमः}
{वानरान्भृशसङ्क्रुद्धो भक्षयन्परिधावति} %||6-49-29||

\twolineshloka
{कुम्भकर्णं समीक्ष्यैव हरयो विप्रदुद्रुवुः}
{कथमेनं रणे क्रुद्धं वारयिष्यन्ति वानराः} %||6-49-30||

\twolineshloka
{उच्यन्तां वानराः सर्वे यन्त्रमेतत्समुच्छ्रितम्}
{इति विज्ञाय हरयो भविष्यन्तीह निर्भयाः} %||6-49-31||

\twolineshloka
{विभीषणवचः श्रुत्वा हेतुमत्सुमुखोद्गतम्}
{उवाच राघवो वाक्यं नीलं सेनापतिं तदा} %||6-49-32||

\twolineshloka
{गच्छ सैन्यानि सर्वाणि व्यूह्य तिष्ठस्व पावके}
{द्वाराण्यादाय लङ्कायाश्चर्याश्चाप्यथ सङ्क्रमान्} %||6-49-33||

\twolineshloka
{शैलशृङ्गाणि वृक्षांश्च शिलाश्चाप्युपसंहरन्}
{तिष्ठन्तु वानराः सर्वे सायुधाः शैलपाणयः} %||6-49-34||

\twolineshloka
{राघवेण समादिष्टो नीलो हरिचमूपतिः}
{शशास वानरानीकं यथावत्कपिकुञ्जरः} %||6-49-35||

\twolineshloka
{ततो गवाक्षः शरभो हनुमानङ्गदो नलः}
{शैलशृङ्गाणि शैलाभा गृहीत्वा द्वारमभ्ययुः} %||6-49-36||

\fourlineindentedshloka
{ततो हरीणां तदनीकमुग्रम्}
{रराज शैलोद्यतवृक्षहस्तम्}
{गिरेः समीपानुगतं यथैव}
{महन्महाम्भोधरजालमुग्रम्} %||6-50-37||


{॥इत्यार्षे श्रीमद्रामायणे वाल्मीकीये आदिकाव्ये युद्धकाण्डे एकोनपञ्चाशत्तमः सर्गः॥}

\sect{पञ्चाशत्तमः सर्गः}

\twolineshloka
{स तु राक्षसशार्दूलो निद्रामदसमाकुलः}
{राजमार्गं श्रिया जुष्टं ययौ विपुलविक्रमः} %||6-50-1||

\twolineshloka
{राक्षसानां सहस्रैश्च वृतः परमदुर्जयः}
{गृहेभ्यः पुष्पवर्षेण कार्यमाणस्तदा ययौ} %||6-50-2||

\twolineshloka
{स हेमजालविततं भानुभास्वरदर्शनम्}
{ददर्श विपुलं रम्यं राक्षसेन्द्रनिवेशनम्} %||6-50-3||

\fourlineindentedshloka
{स तत्तदा सूर्य इवाभ्रजालम्}
{प्रविश्य रक्षोऽधिपतेर्निवेशनम्}
{ददर्श दूरेऽग्रजमासनस्थम्}
{स्वयम्भुवं शक्र इवासनस्थम्} %||6-50-4||

\twolineshloka
{सोऽभिगम्य गृहं भ्रातुः कक्ष्यामभिविगाह्य च}
{ददर्शोद्विग्नमासीनं विमाने पुष्पके गुरुम्} %||6-50-5||

\twolineshloka
{अथ दृष्ट्वा दशग्रीवः कुम्भकर्णमुपस्थितम्}
{तूर्णमुत्थाय संहृष्टः संनिकर्षमुपानयत्} %||6-50-6||

\threelineshloka
{अथासीनस्य पर्यङ्के कुम्भकर्णो महाबलः}
{भ्रातुर्ववन्दे चरणां किं कृत्यमिति चाब्रवीत्}
{उत्पत्य चैनं मुदितो रावणः परिषस्वजे} %||6-50-7||

\twolineshloka
{स भ्रात्रा सम्परिष्वक्तो यथावच्चाभिनन्दितः}
{कुम्भकर्णः शुभं दिव्यं प्रतिपेदे वरासनम्} %||6-50-8||

\twolineshloka
{स तदासनमाश्रित्य कुम्भकर्णो महाबलः}
{संरक्तनयनः कोपाद्रावणं वाक्यमब्रवीत्} %||6-50-9||

\twolineshloka
{किमर्थमहमादृत्य त्वया राजन्प्रबोधितः}
{शंस कस्माद्भयं तेऽस्ति कोऽद्य प्रेतो भविष्यति} %||6-50-10||

\twolineshloka
{भ्रातरं रावणः क्रुद्धं कुम्भकर्णमवस्थितम्}
{ईषत्तु परिवृत्ताभ्यां नेत्राभ्यां वाक्यमब्रवीत्} %||6-50-11||

\twolineshloka
{अद्य ते सुमहान्कालः शयानस्य महाबल}
{सुखितस्त्वं न जानीषे मम रामकृतं भयम्} %||6-50-12||

\twolineshloka
{एष दाशरथी रामः सुग्रीवसहितो बली}
{समुद्रं सबलस्तीर्त्वा मूलं नः परिकृन्तति} %||6-50-13||

\twolineshloka
{हन्त पश्यस्व लङ्काया वनान्युपवनानि च}
{सेतुना सुखमागम्य वानरैकार्णवं कृतम्} %||6-50-14||

\twolineshloka
{ये राक्षसा मुख्यतमा हतास्ते वानरैर्युधि}
{वानराणां क्षयं युद्धे न पश्यामि कदाचन} %||6-50-15||

\twolineshloka
{सर्वक्षपितकोशं च स त्वमभ्यवपद्य माम्}
{त्रायस्वेमां पुरीं लङ्कां बालवृद्धावशेषिताम्} %||6-50-16||

\threelineshloka
{भ्रातुरर्थे महाबाहो कुरु कर्म सुदुष्करम्}
{मयैवं नोक्तपूर्वो हि कश्चिद्भ्रातः परन्तप}
{त्वय्यस्ति मम च स्नेहः परा सम्भावना च मे} %||6-50-17||

\threelineshloka
{देवासुरविमर्देषु बहुशो राक्षसर्षभ}
{त्वया देवाः प्रतिव्यूह्य निर्जिताश्चासुरा युधि}
{न हि ते सर्वभूतेषु दृश्यते सदृशो बली} %||6-50-18||

\fourlineindentedshloka
{कुरुष्व मे प्रियहितमेतदुत्तमम्}
{यथाप्रियं प्रियरणबान्धवप्रिय}
{स्वतेजसा विधम सपत्नवाहिनीम्}
{शरद्घनं पवन इवोद्यतो महान्} %||6-51-19||


{॥इत्यार्षे श्रीमद्रामायणे वाल्मीकीये आदिकाव्ये युद्धकाण्डे पञ्चाशत्तमः सर्गः॥}

\sect{एकपञ्चाशत्तमः सर्गः}

\twolineshloka
{तस्य राक्षसराजस्य निशम्य परिदेवितम्}
{कुम्भकर्णो बभाषेऽथ वचनं प्रजहास च} %||6-51-1||

\twolineshloka
{दृष्टो दोषो हि योऽस्माभिः पुरा मन्त्रविनिर्णये}
{हितेष्वनभियुक्तेन सोऽयमासादितस्त्वया} %||6-51-2||

\twolineshloka
{शीघ्रं खल्वभ्युपेतं त्वां फलं पापस्य कर्मणः}
{निरयेष्वेव पतनं यथा दुष्कृतकर्मणः} %||6-51-3||

\twolineshloka
{प्रथमं वै महाराज कृत्यमेतदचिन्तितम्}
{केवलं वीर्यदर्पेण नानुबन्धो विचारितः} %||6-51-4||

\twolineshloka
{यः पश्चात्पूर्वकार्याणि कुर्यादैश्वर्यमास्थितः}
{पूर्वं चोत्तरकार्याणि न स वेद नयानयौ} %||6-51-5||

\twolineshloka
{देशकालविहीनानि कर्माणि विपरीतवत्}
{क्रियमाणानि दुष्यन्ति हवींष्यप्रयतेष्विव} %||6-51-6||

\twolineshloka
{त्रयाणां पञ्चधा योगं कर्मणां यः प्रपश्यति}
{सचिवैः समयं कृत्वा स सभ्ये वर्तते पथि} %||6-51-7||

\twolineshloka
{यथागमं च यो राजा समयं विचिकीर्षति}
{बुध्यते सचिवान्बुद्ध्या सुहृदश्चानुपश्यति} %||6-51-8||

\twolineshloka
{धर्ममर्थं च कामं च सर्वान्वा रक्षसां पते}
{भजते पुरुषः काले त्रीणि द्वन्द्वानि वा पुनः} %||6-51-9||

\twolineshloka
{त्रिषु चैतेषु यच्छ्रेष्ठं श्रुत्वा तन्नावबुध्यते}
{राजा वा राजमात्रो वा व्यर्थं तस्य बहुश्रुतम्} %||6-51-10||

\twolineshloka
{उपप्रदानं सान्त्वं वा भेदं काले च विक्रमम्}
{योगं च रक्षसां श्रेष्ठ तावुभौ च नयानयौ} %||6-51-11||

\twolineshloka
{काले धर्मार्थकामान्यः सम्मन्त्र्य सचिवैः सह}
{निषेवेतात्मवाँल्लोके न स व्यसनमाप्नुयात्} %||6-51-12||

\twolineshloka
{हितानुबन्धमालोक्य कार्याकार्यमिहात्मनः}
{राजा सहार्थतत्त्वज्ञैः सचिवैः सह जीवति} %||6-51-13||

\twolineshloka
{अनभिज्ञाय शास्त्रार्थान्पुरुषाः पशुबुद्धयः}
{प्रागल्भ्याद्वक्तुमिच्छन्ति मन्त्रेष्वभ्यन्तरीकृताः} %||6-51-14||

\twolineshloka
{अशास्त्रविदुषां तेषां न कार्यमहितं वचः}
{अर्थशास्त्रानभिज्ञानां विपुलां श्रियमिच्छताम्} %||6-51-15||

\twolineshloka
{अहितं च हिताकारं धार्ष्ट्याज्जल्पन्ति ये नराः}
{अवेक्ष्य मन्त्रबाह्यास्ते कर्तव्याः कृत्यदूषणाः} %||6-51-16||

\twolineshloka
{विनाशयन्तो भर्तारं सहिताः शत्रुभिर्बुधैः}
{विपरीतानि कृत्यानि कारयन्तीह मन्त्रिणः} %||6-51-17||

\twolineshloka
{तान्भर्ता मित्रसङ्काशानमित्रान्मन्त्रनिर्णये}
{व्यवहारेण जानीयात्सचिवानुपसंहितान्} %||6-51-18||

\twolineshloka
{चपलस्येह कृत्यानि सहसानुप्रधावतः}
{छिद्रमन्ये प्रपद्यन्ते क्रौञ्चस्य खमिव द्विजाः} %||6-51-19||

\twolineshloka
{यो हि शत्रुमवज्ञाय नात्मानमभिरक्षति}
{अवाप्नोति हि सोऽनर्थान्स्थानाच्च व्यवरोप्यते} %||6-51-20||

\twolineshloka
{तत्तु श्रुत्वा दशग्रीवः कुम्भकर्णस्य भाषितम्}
{भ्रुकुटिं चैव सञ्चक्रे क्रुद्धश्चैनमुवाच ह} %||6-51-21||

\twolineshloka
{मान्यो गुरुरिवाचार्यः किं मां त्वमनुशासति}
{किमेवं वाक्श्रमं कृत्वा काले युक्तं विधीयताम्} %||6-51-22||

\twolineshloka
{विभ्रमाच्चित्तमोहाद्वा बलवीर्याश्रयेण वा}
{नाभिपन्नमिदानीं यद्व्यर्थास्तस्य पुनः कृथाः} %||6-51-23||

\twolineshloka
{अस्मिन्काले तु यद्युक्तं तदिदानीं विधीयताम्}
{ममापनयजं दोषं विक्रमेण समीकुरु} %||6-51-24||

\twolineshloka
{यदि खल्वस्ति मे स्नेहो भ्रातृत्वं वावगच्छसि}
{यदि वा कार्यमेतत्ते हृदि कार्यतमं मतम्} %||6-51-25||

\twolineshloka
{स सुहृद्यो विपन्नार्थं दीनमभ्यवपद्यते}
{स बन्धुर्योऽपनीतेषु साहाय्यायोपकल्पते} %||6-51-26||

\twolineshloka
{तमथैवं ब्रुवाणं तु वचनं धीरदारुणम्}
{रुष्टोऽयमिति विज्ञाय शनैः श्लक्ष्णमुवाच ह} %||6-51-27||

\twolineshloka
{अतीव हि समालक्ष्य भ्रातरं क्षुभितेन्द्रियम्}
{कुम्भकर्णः शनैर्वाक्यं बभाषे परिसान्त्वयन्} %||6-51-28||

\twolineshloka
{अलं राक्षसराजेन्द्र सन्तापमुपपद्य ते}
{रोषं च सम्परित्यज्य स्वस्थो भवितुमर्हसि} %||6-51-29||

\twolineshloka
{नैतन्मनसि कर्तव्व्यं मयि जीवति पार्थिव}
{तमहं नाशयिष्यामि यत्कृते परितप्यसे} %||6-51-30||

\twolineshloka
{अवश्यं तु हितं वाच्यं सर्वावस्थं मया तव}
{बन्धुभावादभिहितं भ्रातृस्नेहाच्च पार्थिव} %||6-51-31||

\twolineshloka
{सदृशं यत्तु कालेऽस्मिन्कर्तुं स्निग्धेन बन्धुना}
{शत्रूणां कदनं पश्य क्रियमाणं मया रणे} %||6-51-32||

\twolineshloka
{अद्य पश्य महाबाहो मया समरमूर्धनि}
{हते रामे सह भ्रात्रा द्रवन्तीं हरिवाहिनीम्} %||6-51-33||

\twolineshloka
{अद्य रामस्य तद्दृष्ट्वा मयानीतं रणाच्छिरः}
{सुखीभव महाबाहो सीता भवतु दुःखिता} %||6-51-34||

\twolineshloka
{अद्य रामस्य पश्यन्तु निधनं सुमहत्प्रियम्}
{लङ्कायां राक्षसाः सर्वे ये ते निहतबान्धवाः} %||6-51-35||

\twolineshloka
{अद्य शोकपरीतानां स्वबन्धुवधकारणात्}
{शत्रोर्युधि विनाशेन करोम्यस्रप्रमार्जनम्} %||6-51-36||

\twolineshloka
{अद्य पर्वतसङ्काशं ससूर्यमिव तोयदम्}
{विकीर्णं पश्य समरे सुग्रीवं प्लवगेश्वरम्} %||6-51-37||

\twolineshloka
{न परः प्रेषणीयस्ते युद्धायातुल विक्रम}
{अहमुत्सादयिष्यामि शत्रूंस्तव महाबल} %||6-51-38||

\twolineshloka
{यदि शक्रो यदि यमो यदि पावकमारुतौ}
{तानहं योधयिष्यामि कुबेर वरुणावपि} %||6-51-39||

\twolineshloka
{गिरिमात्रशरीरस्य शितशूलधरस्य मे}
{नर्दतस्तीक्ष्णदंष्ट्रस्य बिभीयाच्च पुरन्दरः} %||6-51-40||

\twolineshloka
{अथ वा त्यक्तशस्त्रस्य मृद्गतस्तरसा रिपून्}
{न मे प्रतिमुखे कश्चिच्छक्तः स्थातुं जिजीविषुः} %||6-51-41||

\twolineshloka
{नैव शक्त्या न गदया नासिना न शितैः शरैः}
{हस्ताभ्यामेव संरब्धो हनिष्याम्यपि वज्रिणम्} %||6-51-42||

\twolineshloka
{यदि मे मुष्टिवेगं स राघवोऽद्य सहिष्यति}
{ततः पास्यन्ति बाणौघा रुधिरं राघवस्य ते} %||6-51-43||

\twolineshloka
{चिन्तया बाध्यसे राजन्किमर्थं मयि तिष्ठति}
{सोऽहं शत्रुविनाशाय तव निर्यातुमुद्यतः} %||6-51-44||

\threelineshloka
{मुञ्च रामाद्भयं राजन्हनिष्यामीह संयुगे}
{राघवं लक्ष्मणं चैव सुग्रीवं च महाबलम्}
{असाधारणमिच्छामि तव दातुं महद्यशः} %||6-51-45||

\fourlineindentedshloka
{वधेन ते दाशरथेः सुखावहम्}
{सुखं समाहर्तुमहं व्रजामि}
{निहत्य रामं सहलक्ष्मणेन}
{खादामि सर्वान्हरियूथमुख्यान्} %||6-51-46||

\fourlineindentedshloka
{रमस्व कामं पिब चाग्र्यवारुणीम्}
{कुरुष्व कृत्यानि विनीयतां ज्वरः}
{मयाद्य रामे गमिते यमक्षयम्}
{चिराय सीता वशगा भविष्यति} %||6-52-47||


{॥इत्यार्षे श्रीमद्रामायणे वाल्मीकीये आदिकाव्ये युद्धकाण्डे एकपञ्चाशत्तमः सर्गः॥}

\sect{द्विपञ्चाशत्तमः सर्गः}

\twolineshloka
{तदुक्तमतिकायस्य बलिनो बाहुशालिनः}
{कुम्भकर्णस्य वचनं श्रुत्वोवाच महोदरः} %||6-52-1||

\twolineshloka
{कुम्भकर्णकुले जातो धृष्टः प्राकृतदर्शनः}
{अवलिप्तो न शक्नोषि कृत्यं सर्वत्र वेदितुम्} %||6-52-2||

\twolineshloka
{न हि राजा न जानीते कुम्भकर्ण नयानयौ}
{त्वं तु कैशोरकाद्धृष्टः केवलं वक्तुमिच्छसि} %||6-52-3||

\twolineshloka
{स्थानं वृद्धिं च हानिं च देशकालविभागवित्}
{आत्मनश्च परेषां च बुध्यते राक्षसर्षभ} %||6-52-4||

\twolineshloka
{यत्तु शक्यं बलवता कर्तुं प्राकृतबुद्धिना}
{अनुपासितवृद्धेन कः कुर्यात्तादृशं बुधः} %||6-52-5||

\twolineshloka
{यांस्तु धर्मार्थकामांस्त्वं ब्रवीषि पृथगाश्रयान्}
{अनुबोद्धुं स्वभावेन न हि लक्षणमस्ति ते} %||6-52-6||

\twolineshloka
{कर्म चैव हि सर्वेषां कारणानां प्रयोजनम्}
{श्रेयः पापीयसां चात्र फलं भवति कर्मणाम्} %||6-52-7||

\twolineshloka
{निःश्रेयस फलावेव धर्मार्थावितरावपि}
{अधर्मानर्थयोः प्राप्तिः फलं च प्रत्यवायिकम्} %||6-52-8||

\twolineshloka
{ऐहलौकिकपारत्र्यं कर्म पुम्भिर्निषेव्यते}
{कर्माण्यपि तु कल्प्यानि लभते काममास्थितः} %||6-52-9||

\twolineshloka
{तत्र कॢप्तमिदं राज्ञा हृदि कार्यं मतं च नः}
{शत्रौ हि साहसं यत्स्यात्किमिवात्रापनीयते} %||6-52-10||

\twolineshloka
{एकस्यैवाभियाने तु हेतुर्यः प्रकृतस्त्वया}
{तत्राप्यनुपपन्नं ते वक्ष्यामि यदसाधु च} %||6-52-11||

\twolineshloka
{येन पूर्वं जनस्थाने बहवोऽतिबला हताः}
{राक्षसा राघवं तं त्वं कथमेको जयिष्यसि} %||6-52-12||

\twolineshloka
{ये पुरा निर्जितास्तेन जनस्थाने महौजसः}
{राक्षसांस्तान्पुरे सर्वान्भीतानद्यापि पश्यसि} %||6-52-13||

\twolineshloka
{तं सिंहमिव सङ्क्रुद्धं रामं दशरथात्मजम्}
{सर्पं सुप्तमिवाबुद्ध्या प्रबोधयितुमिच्छसि} %||6-52-14||

\twolineshloka
{ज्वलन्तं तेजसा नित्यं क्रोधेन च दुरासदम्}
{कस्तं मृत्युमिवासह्यमासादयितुमर्हति} %||6-52-15||

\twolineshloka
{संशयस्थमिदं सर्वं शत्रोः प्रतिसमासने}
{एकस्य गमनं तत्र न हि मे रोचते तव} %||6-52-16||

\twolineshloka
{हीनार्थस्तु समृद्धार्थं को रिपुं प्राकृतो यथा}
{निश्चितं जीवितत्यागे वशमानेतुमिच्छति} %||6-52-17||

\twolineshloka
{यस्य नास्ति मनुष्येषु सदृशो राक्षसोत्तम}
{कथमाशंससे योद्धुं तुल्येनेन्द्रविवस्वतोः} %||6-52-18||

\twolineshloka
{एवमुक्त्वा तु संरब्धं कुम्भकर्णं महोदरः}
{उवाच रक्षसां मध्ये रावणो लोकरावणम्} %||6-52-19||

\twolineshloka
{लब्ध्वा पुनस्तां वैदेहीं किमर्थं त्वं प्रजल्पसि}
{यदेच्छसि तदा सीता वशगा ते भविष्यति} %||6-52-20||

\twolineshloka
{दृष्टः कश्चिदुपायो मे सीतोपस्थानकारकः}
{रुचितश्चेत्स्वया बुद्ध्या राक्षसेश्वर तं शृणु} %||6-52-21||

\twolineshloka
{अहं द्विजिह्वः संह्रादी कुम्भकर्णो वितर्दनः}
{पञ्चरामवधायैते निर्यान्तीत्यवघोषय} %||6-52-22||

\twolineshloka
{ततो गत्वा वयं युद्धं दास्यामस्तस्य यत्नतः}
{जेष्यामो यदि ते शत्रून्नोपायैः कृत्यमस्ति नः} %||6-52-23||

\twolineshloka
{अथ जीवति नः शत्रुर्वयं च कृतसंयुगाः}
{ततः समभिपत्स्यामो मनसा यत्समीक्षितुम्} %||6-52-24||

\twolineshloka
{वयं युद्धादिहैष्यामो रुधिरेण समुक्षिताः}
{विदार्य स्वतनुं बाणै रामनामाङ्कितैः शितैः} %||6-52-25||

\twolineshloka
{भक्षितो राघवोऽस्माभिर्लक्ष्मणश्चेति वादिनः}
{तव पादौ ग्रहीष्यामस्त्वं नः काम प्रपूरय} %||6-52-26||

\twolineshloka
{ततोऽवघोषय पुरे गजस्कन्धेन पार्थिव}
{हतो रामः सह भ्रात्रा ससैन्य इति सर्वतः} %||6-52-27||

\twolineshloka
{प्रीतो नाम ततो भूत्वा भृत्यानां त्वमरिन्दम}
{भोगांश्च परिवारांश्च कामांश्च वसुदापय} %||6-52-28||

\twolineshloka
{ततो माल्यानि वासांसि वीराणामनुलेपनम्}
{पेयं च बहु योधेभ्यः स्वयं च मुदितः पिब} %||6-52-29||

\threelineshloka
{ततोऽस्मिन्बहुलीभूते कौलीने सर्वतो गते}
{प्रविश्याश्वास्य चापि त्वं सीतां रहसि सान्त्वय}
{धनधान्यैश्च कामैश्च रत्नैश्चैनां प्रलोभय} %||6-52-30||

\twolineshloka
{अनयोपधया राजन्भयशोकानुबन्धया}
{अकामा त्वद्वशं सीता नष्टनाथा गमिष्यति} %||6-52-31||

\twolineshloka
{रञ्जनीयं हि भर्तारं विनष्टमवगम्य सा}
{नैराश्यात्स्त्रीलघुत्वाच्च त्वद्वशं प्रतिपत्स्यते} %||6-52-32||

\twolineshloka
{सा पुरा सुखसंवृद्धा सुखार्हा दुःखकर्षिता}
{त्वय्यधीनः सुखं ज्ञात्वा सर्वथोपगमिष्यति} %||6-52-33||

\fourlineindentedshloka
{एतत्सुनीतं मम दर्शनेन}
{रामं हि दृष्ट्वैव भवेदनर्थः}
{इहैव ते सेत्स्यति मोत्सुको भूर्-}
{महानयुद्धेन सुखस्य लाभः} %||6-52-34||

\fourlineindentedshloka
{अनष्टसैन्यो ह्यनवाप्तसंशयो}
{रिपूनयुद्धेन जयञ्जनाधिप}
{यशश्च पुण्यं च महन्महीपते}
{श्रियं च कीर्तिं च चिरं समश्नुते} %||6-53-35||


{॥इत्यार्षे श्रीमद्रामायणे वाल्मीकीये आदिकाव्ये युद्धकाण्डे द्विपञ्चाशत्तमः सर्गः॥}

\sect{त्रिपञ्चाशत्तमः सर्गः}

\twolineshloka
{स तथोक्तस्तु निर्भर्त्स्य कुम्भकर्णो महोदरम्}
{अब्रवीद्राक्षसश्रेष्ठं भ्रातरं रावणं ततः} %||6-53-1||

\twolineshloka
{सोऽहं तव भयं घोरं वधात्तस्य दुरात्मनः}
{रामस्याद्य प्रमार्जामि निर्वैरस्त्वं सुखीभव} %||6-53-2||

\twolineshloka
{गर्जन्ति न वृथा शूर निर्जला इव तोयदाः}
{पश्य सम्पाद्यमानं तु गर्जितं युधि कर्मणा} %||6-53-3||

\twolineshloka
{न मर्षयति चात्मानं सम्भावयति नात्मना}
{अदर्शयित्वा शूरास्तु कर्म कुर्वन्ति दुष्करम्} %||6-53-4||

\twolineshloka
{विक्लवानामबुद्धीनां राज्ञां पण्डितमानिनाम्}
{शृण्वतामादित इदं त्वद्विधानां महोदर} %||6-53-5||

\twolineshloka
{युद्धे कापुरुषैर्नित्यं भवद्भिः प्रियवादिभिः}
{राजानमनुगच्छद्भिः कृत्यमेतद्विनाशितम्} %||6-53-6||

\twolineshloka
{राजशेषा कृता लङ्का क्षीणः कोशो बलं हतम्}
{राजानमिममासाद्य सुहृच्चिह्नममित्रकम्} %||6-53-7||

\twolineshloka
{एष निर्याम्यहं युद्धमुद्यतः शत्रुनिर्जये}
{दुर्नयं भवतामद्य समीकर्तुं महाहवे} %||6-53-8||

\twolineshloka
{एवमुक्तवतो वाक्यं कुम्भकर्णस्य धीमतः}
{प्रत्युवाच ततो वाक्यं प्रहसन्राक्षसाधिपः} %||6-53-9||

\twolineshloka
{महोदरोऽयं रामात्तु परित्रस्तो न संशयः}
{न हि रोचयते तात युद्धं युद्धविशारद} %||6-53-10||

\twolineshloka
{कश्चिन्मे त्वत्समो नास्ति सौहृदेन बलेन च}
{गच्छ शत्रुवधाय त्वं कुम्भकर्णजयाय च} %||6-53-11||

\twolineshloka
{आददे निशितं शूलं वेगाच्छत्रुनिबर्हणः}
{सर्वकालायसं दीप्तं तप्तकाञ्चनभूषणम्} %||6-53-12||

\twolineshloka
{इन्द्राशनिसमं भीमं वज्रप्रतिमगौरवम्}
{देवदानवगन्धर्वयक्षकिंनरसूदनम्} %||6-53-13||

\threelineshloka
{रक्तमाल्य महादाम स्वतश्चोद्गतपावकम्}
{आदाय निशितं शूलं शत्रुशोणितरञ्जितम्}
{कुम्भकर्णो महातेजा रावणं वाक्यमब्रवीत्} %||6-53-14||

\twolineshloka
{गमिष्याम्यहमेकाकी तिष्ठत्विह बलं महत्}
{अद्य तान्क्षुधितः क्रुद्धो भक्षयिष्यामि वानरान्} %||6-53-15||

\twolineshloka
{कुम्भकर्णवचः श्रुत्वा रावणो वाक्यमब्रवीत्}
{सैन्यैः परिवृतो गच्छ शूलमुद्गलपाणिभिः} %||6-53-16||

\twolineshloka
{वानरा हि महात्मानः शीघ्राश्च व्यवसायिनः}
{एकाकिनं प्रमत्तं वा नयेयुर्दशनैः क्षयम्} %||6-53-17||

\twolineshloka
{तस्मात्परमदुर्धर्षैः सैन्यैः परिवृतो व्रज}
{रक्षसामहितं सर्वं शत्रुपक्षं निसूदय} %||6-53-18||

\twolineshloka
{अथासनात्समुत्पत्य स्रजं मणिकृतान्तराम्}
{आबबन्ध महातेजाः कुम्भकर्णस्य रावणः} %||6-53-19||

\twolineshloka
{अङ्गदानङ्गुलीवेष्टान्वराण्याभरणानि च}
{हारं च शशिसङ्काशमाबबन्ध महात्मनः} %||6-53-20||

\twolineshloka
{दिव्यानि च सुगन्धीनि माल्यदामानि रावणः}
{श्रोत्रे चासज्जयामास श्रीमती चास्य कुण्डले} %||6-53-21||

\twolineshloka
{काञ्चनाङ्गदकेयूरो निष्काभरणभूषितः}
{कुम्भकर्णो बृहत्कर्णः सुहुतोऽग्निरिवाबभौ} %||6-53-22||

\twolineshloka
{श्रोणीसूत्रेण महता मेचकेन विराजितः}
{अमृतोत्पादने नद्धो भुजङ्गेनेव मन्दरः} %||6-53-23||

\fourlineindentedshloka
{स काञ्चनं भारसहं निवातम्}
{विद्युत्प्रभं दीप्तमिवात्मभासा}
{आबध्यमानः कवचं रराज}
{सन्ध्याभ्रसंवीत इवाद्रिराजः} %||6-53-24||

\twolineshloka
{सर्वाभरणनद्धाङ्गः शूलपाणिः स राक्षसः}
{त्रिविक्रमकृतोत्साहो नारायण इवाबभौ} %||6-53-25||

\threelineshloka
{भ्रातरं सम्परिष्वज्य कृत्वा चापि प्रदक्षिणम्}
{प्रणम्य शिरसा तस्मै सम्प्रतस्थे महाबलिः}
{तमाशीर्भिः प्रशस्ताभिः प्रेषयामास रावणः} %||6-53-26||

\threelineshloka
{शङ्खदुन्दुभिनिर्घोषैः सैन्यैश्चापि वरायुधैः}
{तं गजैश्च तुरङ्गैश्च स्यन्दनैश्चाम्बुदस्वनैः}
{अनुजग्मुर्महात्मानं रथिनो रथिनां वरम्} %||6-53-27||

\twolineshloka
{सर्पैरुष्ट्रैः खरैरश्वैः सिंहद्विपमृगद्विजैः}
{अनुजग्मुश्च तं घोरं कुम्भकर्णं महाबलम्} %||6-53-28||

\fourlineindentedshloka
{स पुष्पवर्णैरवकीर्यमाणो}
{धृतातपत्रः शितशूलपाणिः}
{मदोत्कटः शोणितगन्धमत्तो}
{विनिर्ययौ दानवदेवशत्रुः} %||6-53-29||

\twolineshloka
{पदातयश बहवो महानादा महाबलाः}
{अन्वयू राक्षसा भीमा भीमाक्षाः शस्त्रपाणयः} %||6-53-30||

\twolineshloka
{रक्ताक्षाः सुमहाकाया नीलाञ्जनचयोपमाः}
{शूरानुद्यम्य खड्गांश्च निशितांश्च परश्वधान्} %||6-53-31||

\twolineshloka
{बहुव्यामांश्च विपुलान्क्षेपणीयान्दुरासदान्}
{तालस्कन्धांश्च विपुलान्क्षेपणीयान्दुरासदान्} %||6-53-32||

\twolineshloka
{अथान्यद्वपुरादाय दारुणं लोमहर्षणम्}
{निष्पपात महातेजाः कुम्भकर्णो महाबलः} %||6-53-33||

\twolineshloka
{धनुःशतपरीणाहः स षट्शतसमुच्छितः}
{रौद्रः शकटचक्राक्षो महापर्वतसंनिभः} %||6-53-34||

\twolineshloka
{संनिपत्य च रक्षांसि दग्धशैलोपमो महान्}
{कुम्भकर्णो महावक्त्रः प्रहसन्निदमब्रवीत्} %||6-53-35||

\twolineshloka
{अद्य वानरमुख्यानां तानि यूथानि भागशः}
{निर्दहिष्यामि सङ्क्रुद्धः शलभानिव पावकः} %||6-53-36||

\twolineshloka
{नापराध्यन्ति मे कामं वानरा वनचारिणः}
{जातिरस्मद्विधानां सा पुरोद्यानविभूषणम्} %||6-53-37||

\twolineshloka
{पुररोधस्य मूलं तु राघवः सहलक्ष्मणः}
{हते तस्मिन्हतं सर्वं तं वधिष्यामि संयुगे} %||6-53-38||

\twolineshloka
{एवं तस्य ब्रुवाणस्य कुम्भकर्णस्य राक्षसाः}
{नादं चक्रुर्महाघोरं कम्पयन्त इवार्णवम्} %||6-53-39||

\twolineshloka
{तस्य निष्पततस्तूर्णं कुम्भकर्णस्य धीमतः}
{बभूवुर्घोररूपाणि निमित्तानि समन्ततः} %||6-53-40||

\twolineshloka
{उल्काशनियुता मेघा विनेदुश्च सुदारुणाः}
{ससागरवना चैव वसुधा समकम्पत} %||6-53-41||

\twolineshloka
{घोररूपाः शिवा नेदुः सज्वालकवलैर्मुखैः}
{मण्डलान्यपसव्यानि बबन्धुश्च विहङ्गमाः} %||6-53-42||

\twolineshloka
{निष्पपात च गृध्रेऽस्य शूले वै पथि गच्छतः}
{प्रास्फुरन्नयनं चास्य सव्यो बाहुरकम्पत} %||6-53-43||

\twolineshloka
{निष्पपात तदा चोक्ला ज्वलन्ती भीमनिस्वना}
{आदित्यो निष्प्रभश्चासीन्न प्रवाति सुखोऽनिलः} %||6-53-44||

\twolineshloka
{अचिन्तयन्महोत्पातानुत्थिताँल्लोमहर्षणान्}
{निर्ययौ कुम्भकर्णस्तु कृतान्तबलचोदितः} %||6-53-45||

\twolineshloka
{स लङ्घयित्वा प्राकारं पद्भ्यां पर्वतसंनिभः}
{ददर्शाभ्रघनप्रख्यं वानरानीकमद्भुतम्} %||6-53-46||

\twolineshloka
{ते दृष्ट्वा राक्षसश्रेष्ठं वानराः पर्वतोपमम्}
{वायुनुन्ना इव घना ययुः सर्वा दिशस्तदा} %||6-53-47||

\fourlineindentedshloka
{तद्वानरानीकमतिप्रचण्डम्}
{दिशो द्रवद्भिन्नमिवाभ्रजालम्}
{स कुम्भकर्णः समवेक्ष्य हर्षात्}
{ननाद भूयो घनवद्घनाभः} %||6-53-48||

\fourlineindentedshloka
{ते तस्य घोरं निनदं निशम्य}
{यथा निनादं दिवि वारिदस्य}
{पेतुर्धरण्यां बहवः प्लवङ्गा}
{निकृत्तमूला इव सालवृक्षाः} %||6-53-49||

\fourlineindentedshloka
{विपुलपरिघवान्स कुम्भकर्णो}
{रिपुनिधनाय विनिःसृतो महात्मा}
{कपि गणभयमाददत्सुभीमम्}
{प्रभुरिव किङ्करदण्डवान्युगान्ते} %||6-54-50||


{॥इत्यार्षे श्रीमद्रामायणे वाल्मीकीये आदिकाव्ये युद्धकाण्डे त्रिपञ्चाशत्तमः सर्गः॥}

\sect{चतुःपञ्चाशत्तमः सर्गः}

\twolineshloka
{स ननाद महानादं समुद्रमभिनादयन्}
{जनयन्निव निर्घातान्विधमन्निव पर्वतान्} %||6-54-1||

\twolineshloka
{तमवध्यं मघवता यमेन वरुणेन च}
{प्रेक्ष्य भीमाक्षमायान्तं वानरा विप्रदुद्रुवुः} %||6-54-2||

\twolineshloka
{तांस्तु विद्रवतो दृष्ट्वा वालिपुत्रोऽङ्गदोऽब्रवीत्}
{नलं नीलं गवाक्षं च कुमुदं च महाबलम्} %||6-54-3||

\twolineshloka
{आत्मानमत्र विस्मृत्य वीर्याण्यभिजनानि च}
{क्व गच्छत भयत्रस्ताः प्राकृता हरयो यथा} %||6-54-4||

\twolineshloka
{साधु सौम्या निवर्तध्वं किं प्राणान्परिरक्षथ}
{नालं युद्धाय वै रक्षो महतीयं विभीषिकाः} %||6-54-5||

\twolineshloka
{महतीमुत्थितामेनां राक्षसानां विभीषिकाम्}
{विक्रमाद्विधमिष्यामो निवर्तध्वं प्लवङ्गमाः} %||6-54-6||

\twolineshloka
{कृच्छ्रेण तु समाश्वास्य सङ्गम्य च ततस्ततः}
{वृक्षाद्रिहस्ता हरयः सम्प्रतस्थू रणाजिरम्} %||6-54-7||

\threelineshloka
{ते निवृत्य तु सङ्क्रुद्धाः कुम्भकर्णं वनौकसः}
{निजघ्नुः परमक्रुद्धाः समदा इव कुञ्जराः}
{प्रांशुभिर्गिरिशृङ्गैश्च शिलाभिश्च महाबलाः} %||6-54-8||

\threelineshloka
{पादपैः पुष्पिताग्रैश्च हन्यमानो न कम्पते}
{तस्य गात्रेषु पतिता भिद्यन्ते शतशः शिलाः}
{पादपाः पुष्पिताग्राश्च भग्नाः पेतुर्महीतले} %||6-54-9||

\twolineshloka
{सोऽपि सैन्यानि सङ्क्रुद्धो वानराणां महौजसाम्}
{ममन्थ परमायत्तो वनान्यग्निरिवोत्थितः} %||6-54-10||

\twolineshloka
{लोहितार्द्रास्तु बहवः शेरते वानरर्षभाः}
{निरस्ताः पतिता भूमौ ताम्रपुष्पा इव द्रुमाः} %||6-54-11||

\twolineshloka
{लङ्घयन्तः प्रधावन्तो वानरा नावलोकयन्}
{केचित्समुद्रे पतिताः केचिद्गगनमाश्रिताः} %||6-54-12||

\twolineshloka
{वध्यमानास्तु ते वीरा राक्षसेन बलीयसा}
{सागरं येन ते तीर्णाः पथा तेनैव दुद्रुवुः} %||6-54-13||

\twolineshloka
{ते स्थलानि तथा निम्नं विषण्णवदना भयात्}
{ऋक्षा वृक्षान्समारूढाः केचित्पर्वतमाश्रिताः} %||6-54-14||

\twolineshloka
{ममज्जुरर्णवे केचिद्गुहाः केचित्समाश्रिताः}
{निषेदुः प्लवगाः केचित्केचिन्नैवावतस्थिरे} %||6-54-15||

\twolineshloka
{तान्समीक्ष्याङ्गदो भङ्गान्वानरानिदमब्रवीत्}
{अवतिष्ठत युध्यामो निवर्तध्वं प्लवङ्गमाः} %||6-54-16||

\twolineshloka
{भग्नानां वो न पश्यामि परिगम्य महीमिमाम्}
{स्थानं सर्वे निवर्तध्वं किं प्राणान्परिरक्षथ} %||6-54-17||

\twolineshloka
{निरायुधानां द्रवतामसङ्गगतिपौरुषाः}
{दारा ह्यपहसिष्यन्ति स वै घातस्तु जीविताम्} %||6-54-18||

\twolineshloka
{कुलेषु जाताः सर्वे स्म विस्तीर्णेषु महत्सु च}
{अनार्याः खलु यद्भीतास्त्यक्त्वा वीर्यं प्रधावत} %||6-54-19||

\twolineshloka
{विकत्थनानि वो यानि यदा वै जनसंसदि}
{तानि वः क्व च यतानि सोदग्राणि महान्ति च} %||6-54-20||

\twolineshloka
{भीरुप्रवादाः श्रूयन्ते यस्तु जीवति धिक्कृतः}
{मार्गः सत्पुरुषैर्जुष्टः सेव्यतां त्यज्यतां भयम्} %||6-54-21||

\threelineshloka
{शयामहे वा निहताः पृथिव्यामल्पजीविताः}
{दुष्प्रापं ब्रह्मलोकं वा प्राप्नुमो युधि सूदिताः}
{सम्प्राप्नुयामः कीर्तिं वा निहत्य शत्रुमाहवे} %||6-54-22||

\twolineshloka
{न कुम्भकर्णः काकुत्स्थं दृष्ट्वा जीवन्गमिष्यति}
{दीप्यमानमिवासाद्य पतङ्गो ज्वलनं यथा} %||6-54-23||

\twolineshloka
{पलायनेन चोद्दिष्टाः प्राणान्रक्षामहे वयम्}
{एकेन बहवो भग्ना यशो नाशं गमिष्यति} %||6-54-24||

\twolineshloka
{एवं ब्रुवाणं तं शूरमङ्गदं कनकाङ्गदम्}
{द्रवमाणास्ततो वाक्यमूचुः शूरविगर्हितम्} %||6-54-25||

\twolineshloka
{कृतं नः कदनं घोरं कुम्भकर्णेन रक्षसा}
{न स्थानकालो गच्छामो दयितं जीवितं हि नः} %||6-54-26||

\twolineshloka
{एतावदुक्त्वा वचनं सर्वे ते भेजिरे दिशः}
{भीमं भीमाक्षमायान्तं दृष्ट्वा वानरयूथपाः} %||6-54-27||

\twolineshloka
{द्रवमाणास्तु ते वीरा अङ्गदेन वलीमुखाः}
{सान्त्वैश्च बहुमानैश्च ततः सर्वे निवर्तिताः} %||6-54-28||

\fourlineindentedshloka
{ऋषभशरभमैन्दधूम्रनीलाः}
{कुमुदसुषेणगवाक्षरम्भताराः}
{द्विविदपनसवायुपुत्रमुख्याः}
{त्वरिततराभिमुखं रणं प्रयाताः} %||6-55-29||


{॥इत्यार्षे श्रीमद्रामायणे वाल्मीकीये आदिकाव्ये युद्धकाण्डे चतुःपञ्चाशत्तमः सर्गः॥}

\sect{पञ्चपञ्चाशत्तमः सर्गः}

\twolineshloka
{ते निवृत्ता महाकायाः श्रुत्वाङ्गदवचस्तदा}
{नैष्ठिकीं बुद्धिमास्थाय सर्वे सङ्ग्रामकाङ्क्षिणः} %||6-55-1||

\twolineshloka
{समुदीरितवीर्यास्ते समारोपितविक्रमाः}
{पर्यवस्थापिता वाक्यैरङ्गदेन वलीमुखाः} %||6-55-2||

\twolineshloka
{प्रयाताश्च गता हर्षं मरणे कृतनिश्चयाः}
{चक्रुः सुतुमुलं युद्धं वानरास्त्यक्तजीविताः} %||6-55-3||

\twolineshloka
{अथ वृक्षान्महाकायाः सानूनि सुमहान्ति च}
{वानरास्तूर्णमुद्यम्य कुम्भकर्णमभिद्रवन्} %||6-55-4||

\twolineshloka
{स कुम्भकर्णः सङ्क्रुद्धो गदामुद्यम्य वीर्यवान्}
{अर्दयन्सुमहाकायः समन्ताद्व्याक्षिपद्रिपून्} %||6-55-5||

\twolineshloka
{शतानि सप्त चाष्टौ च सहस्राणि च वानराः}
{प्रकीर्णाः शेरते भूमौ कुम्भकर्णेन पोथिताः} %||6-55-6||

\threelineshloka
{षोडशाष्टौ च दश च विंशत्त्रिंशत्तथैव च}
{परिक्षिप्य च बाहुभ्यां खादन्विपरिधावति}
{भक्षयन्भृशसङ्क्रुद्धो गरुडः पन्नगानिव} %||6-55-7||

\twolineshloka
{हनूमाञ्शैलशृङ्गाणि वृक्षांश्च विविधान्बहून्}
{ववर्ष कुम्भकर्णस्य शिरस्यम्बरमास्थितः} %||6-55-8||

\twolineshloka
{तानि पर्वतशृङ्गाणि शूलेन तु बिभेद ह}
{बभञ्ज वृक्षवर्षं च कुम्भकर्णो महाबलः} %||6-55-9||

\fourlineindentedshloka
{ततो हरीणां तदनीकमुग्रम्}
{दुद्राव शूलं निशितं प्रगृह्य}
{तस्थौ ततोऽस्यापततः पुरस्तात्}
{महीधराग्रं हनुमान्प्रगृह्य} %||6-55-10||

\fourlineindentedshloka
{स कुम्भकर्णं कुपितो जघान}
{वेगेन शैलोत्तमभीमकायम्}
{स चुक्षुभे तेन तदाभिबूतो}
{मेदार्द्रगात्रो रुधिरावसिक्तः} %||6-55-11||

\fourlineindentedshloka
{स शूलमाविध्य तडित्प्रकाशम्}
{गिरिं यथा प्रज्वलिताग्रशृङ्गम्}
{बाह्वन्तरे मारुतिमाजघान}
{गुहोऽचलं क्रौञ्चमिवोग्रशक्त्या} %||6-55-12||

\fourlineindentedshloka
{स शूलनिर्भिन्न महाभुजान्तरः}
{प्रविह्वलः शोणितमुद्वमन्मुखात्}
{ननाद भीमं हनुमान्महाहवे}
{युगान्तमेघस्तनितस्वनोपमम्} %||6-55-13||

\fourlineindentedshloka
{ततो विनेदुः सहसा प्रहृष्टा}
{रक्षोगणास्तं व्यथितं समीक्ष्य}
{प्लवङ्गमास्तु व्यथिता भयार्ताः}
{प्रदुद्रुवुः संयति कुम्भकर्णात्} %||6-55-14||

\twolineshloka
{नीलश्चिक्षेप शैलाग्रं कुम्भकर्णाय धीमते}
{तमापतन्तं सम्प्रेक्ष्य मुष्टिनाभिजघान ह} %||6-55-15||

\twolineshloka
{मुष्टिप्रहाराभिहतं तच्छैलाग्रं व्यशीर्यत}
{सविस्फुलिङ्गं सज्वालं निपपात महीतले} %||6-55-16||

\twolineshloka
{ऋषभः शरभो नीलो गवाक्षो गन्धमादनः}
{पञ्चवानरशार्दूलाः कुम्भकर्णमुपाद्रवन्} %||6-55-17||

\twolineshloka
{शैलैर्वृक्षैस्तलैः पादैर्मुष्टिभिश्च महाबलाः}
{कुम्भकर्णं महाकायं सर्वतोऽभिनिजघ्निरे} %||6-55-18||

\twolineshloka
{स्पर्शानिव प्रहारांस्तान्वेदयानो न विव्यथे}
{ऋषभं तु महावेगं बाहुभ्यां परिषस्वजे} %||6-55-19||

\twolineshloka
{कुम्भकर्णभुजाभ्यां तु पीडितो वानरर्षभः}
{निपपातर्षभो भीमः प्रमुखागतशोणितः} %||6-55-20||

\twolineshloka
{मुष्टिना शरभं हत्वा जानुना नीलमाहवे}
{आजघान गवाक्षं च तलेनेन्द्ररिपुस्तदा} %||6-55-21||

\twolineshloka
{दत्तप्रहरव्यथिता मुमुहुः शोणितोक्षिताः}
{निपेतुस्ते तु मेदिन्यां निकृत्ता इव किंशुकाः} %||6-55-22||

\twolineshloka
{तेषु वानरमुख्येषु पतितेषु महात्मसु}
{वानराणां सहस्राणि कुम्भकर्णं प्रदुद्रुवुः} %||6-55-23||

\twolineshloka
{तं शैलमिव शैलाभाः सर्वे तु प्लवगर्षभाः}
{समारुह्य समुत्पत्य ददंशुश्च महाबलाः} %||6-55-24||

\twolineshloka
{तं नखैर्दशनैश्चापि मुष्टिभिर्जानुभिस्तथा}
{कुम्भकर्णं महाकायं ते जघ्नुः प्लवगर्षभाः} %||6-55-25||

\twolineshloka
{स वानरसहस्रैस्तैराचितः पर्वतोपमः}
{रराज राक्षसव्याघ्रो गिरिरात्मरुहैरिव} %||6-55-26||

\twolineshloka
{बाहुभ्यां वानरान्सर्वान्प्रगृह्य स महाबलः}
{भक्षयामास सङ्क्रुद्धो गरुडः पन्नगानिव} %||6-55-27||

\twolineshloka
{प्रक्षिप्ताः कुम्भकर्णेन वक्त्रे पातालसंनिभे}
{नासा पुटाभ्यां निर्जग्मुः कर्णाभ्यां चैव वानराः} %||6-55-28||

\twolineshloka
{भक्षयन्भृशसङ्क्रुद्धो हरीन्पर्वतसंनिभः}
{बभञ्ज वानरान्सर्वान्सङ्क्रुद्धो राक्षसोत्तमः} %||6-55-29||

\twolineshloka
{मांसशोणितसङ्क्लेदां भूमिं कुर्वन्स राक्षसः}
{चचार हरिसैन्येषु कालाग्निरिव मूर्छितः} %||6-55-30||

\twolineshloka
{वज्रहस्तो यथा शक्रः पाशहस्त इवान्तकः}
{शूलहस्तो बभौ तस्मिन्कुम्भकर्णो महाबलः} %||6-55-31||

\twolineshloka
{यथा शुष्काण्यरण्यानि ग्रीष्मे दहति पावकः}
{तथा वानरसैन्यानि कुम्भकर्णो विनिर्दहत्} %||6-55-32||

\twolineshloka
{ततस्ते वध्यमानास्तु हतयूथा विनायकाः}
{वानरा भयसंविग्ना विनेदुर्विस्वरं भृशम्} %||6-55-33||

\twolineshloka
{अनेकशो वध्यमानाः कुम्भकर्णेन वानराः}
{राघवं शरणं जग्मुर्व्यथिताः खिन्नचेतसः} %||6-55-34||

\twolineshloka
{तमापतन्तं सम्प्रेक्ष्य कुम्भकर्णं महाबलम्}
{उत्पपात तदा वीरः सुग्रीवो वानराधिपः} %||6-55-35||

\twolineshloka
{स पर्वताग्रमुत्क्षिप्य समाविध्य महाकपिः}
{अभिदुद्राव वेगेन कुम्भकर्णं महाबलम्} %||6-55-36||

\twolineshloka
{तमापतन्तं सम्प्रेक्ष्य कुम्भकर्णः प्लवङ्गमम्}
{तस्थौ विवृतसर्वाङ्गो वानरेन्द्रस्य सम्मुखः} %||6-55-37||

\twolineshloka
{कपिशोणितदिग्धाङ्गं भक्षयन्तं महाकपीन्}
{कुम्भकर्णं स्थितं दृष्ट्वा सुग्रीवो वाक्यमब्रवीत्} %||6-55-38||

\twolineshloka
{पातिताश्च त्वया वीराः कृतं कर्म सुदुष्करम्}
{भक्षितानि च सैन्यानि प्राप्तं ते परमं यशः} %||6-55-39||

\twolineshloka
{त्यज तद्वानरानीकं प्राकृतैः किं करिष्यसि}
{सहस्वैकं निपातं मे पर्वतस्यास्य राक्षस} %||6-55-40||

\twolineshloka
{तद्वाक्यं हरिराजस्य सत्त्वधैर्यसमन्वितम्}
{श्रुत्वा राक्षसशार्दूलः कुम्भकर्णोऽब्रवीद्वचः} %||6-55-41||

\twolineshloka
{प्रजापतेस्तु पौत्रस्त्वं तथैवर्क्षरजःसुतः}
{श्रुतपौरुषसम्पन्नस्तस्माद्गर्जसि वानर} %||6-55-42||

\fourlineindentedshloka
{स कुम्भकर्णस्य वचो निशम्य}
{व्याविध्य शैलं सहसा मुमोच}
{तेनाजघानोरसि कुम्भकर्णम्}
{शैलेन वज्राशनिसंनिभेन} %||6-55-43||

\fourlineindentedshloka
{तच्छैलशृङ्गं सहसा विकीर्णम्}
{भुजान्तरे तस्य तदा विशाले}
{ततो विषेदुः सहसा प्लवङ्गमा}
{रक्षोगणाश्चापि मुदा विनेदुः} %||6-55-44||

\fourlineindentedshloka
{स शैलशृङ्गाभिहतश्चुकोप}
{ननाद कोपाच्च विवृत्य वक्त्रम्}
{व्याविध्य शूलं च तडित्प्रकाशम्}
{चिक्षेप हर्यृक्षपतेर्वधाय} %||6-55-45||

\fourlineindentedshloka
{तत्कुम्भकर्णस्य भुजप्रविद्धम्}
{शूलं शितं काञ्चनदामजुष्टम्}
{क्षिप्रं समुत्पत्य निगृह्य दोर्भ्याम्}
{बभञ्ज वेगेन सुतोऽनिलस्य} %||6-55-46||

\twolineshloka
{कृतं भारसहस्रस्य शूलं कालायसं महत्}
{बभञ्ज जनौमारोप्य प्रहृष्टः प्लवगर्षभः} %||6-55-47||

\fourlineindentedshloka
{स तत्तदा भग्नमवेक्ष्य शूलम्}
{चुकोप रक्षोऽधिपतिर्महात्मा}
{उत्पाट्य लङ्कामलयात्स शृङ्गम्}
{जघान सुग्रीवमुपेत्य तेन} %||6-55-48||

\fourlineindentedshloka
{स शैलशृङ्गाभिहतो विसंज्ञः}
{पपात भूमौ युधि वानरेन्द्रः}
{तं प्रेक्ष्य भूमौ पतितं विसंज्ञम्}
{नेदुः प्रहृष्टा युधि यातुधानाः} %||6-55-49||

\fourlineindentedshloka
{तमभ्युपेत्याद्भुतघोरवीर्यम्}
{स कुम्भकर्णो युधि वानरेन्द्रम्}
{जहार सुग्रीवमभिप्रगृह्य}
{यथानिलो मेघमतिप्रचण्डः} %||6-55-50||

\fourlineindentedshloka
{स तं महामेघनिकाशरूपम्}
{उत्पाट्य गच्छन्युधि कुम्भकर्णः}
{रराज मेरुप्रतिमानरूपो}
{मेरुर्यथात्युच्छ्रितघोरशृङ्गः} %||6-55-51||

\fourlineindentedshloka
{ततः समुत्पाट्य जगाम वीरः}
{संस्तूयमानो युधि राक्षसेन्द्रैः}
{शृण्वन्निनादं त्रिदशालयानाम्}
{प्लवङ्गराजग्रहविस्मितानाम्} %||6-55-52||

\fourlineindentedshloka
{ततस्तमादाय तदा स मेने}
{हरीन्द्रमिन्द्रोपममिन्द्रवीर्यः}
{अस्मिन्हृते सर्वमिदं हृतं स्यात्}
{सराघवं सैन्यमितीन्द्रशत्रुः} %||6-55-53||

\twolineshloka
{विद्रुतां वाहिनीं दृष्ट्वा वानराणां ततस्ततः}
{कुम्भकर्णेन सुग्रीवं गृहीतं चापि वानरम्} %||6-55-54||

\twolineshloka
{हनूमांश्चिन्तयामास मतिमान्मारुतात्मजः}
{एवं गृहीते सुग्रीवे किं कर्तव्यं मया भवेत्} %||6-55-55||

\twolineshloka
{यद्वै न्याय्यं मया कर्तुं तत्करिष्यामि सर्वथा}
{भूत्वा पर्वतसङ्काशो नाशयिष्यामि राक्षसम्} %||6-55-56||

\fourlineindentedshloka
{मया हते संयति कुम्भकर्णे}
{महाबले मुष्टिविशीर्णदेहे}
{विमोचिते वानरपार्थिवे च}
{भवन्तु हृष्टाः प्रवगाः समग्राः} %||6-55-57||

\twolineshloka
{अथ वा स्वयमप्येष मोक्षं प्राप्स्यति पार्थिवः}
{गृहीतोऽयं यदि भवेत्त्रिदशैः सासुरोरगैः} %||6-55-58||

\twolineshloka
{मन्ये न तावदात्मानं बुध्यते वानराधिपः}
{शैलप्रहाराभिहतः कुम्भकर्णेन संयुगे} %||6-55-59||

\twolineshloka
{अयं मुहूर्तात्सुग्रीवो लब्धसंज्ञो महाहवे}
{आत्मनो वानराणां च यत्पथ्यं तत्करिष्यति} %||6-55-60||

\twolineshloka
{मया तु मोक्षितस्यास्य सुग्रीवस्य महात्मनः}
{अप्रीतश्च भवेत्कष्टा कीर्तिनाशश्च शाश्वतः} %||6-55-61||

\twolineshloka
{तस्मान्मुहूर्तं काङ्क्षिष्ये विक्रमं पार्थिवस्य नः}
{भिन्नं च वानरानीकं तावदाश्वासयाम्यहम्} %||6-55-62||

\twolineshloka
{इत्येवं चिन्तयित्वा तु हनूमान्मारुतात्मजः}
{भूयः संस्तम्भयामास वानराणां महाचमूम्} %||6-55-63||

\fourlineindentedshloka
{स कुम्भकर्णोऽथ विवेश लङ्काम्}
{स्फुरन्तमादाय महाहरिं तम्}
{विमानचर्यागृहगोपुरस्थैः}
{पुष्पाग्र्यवर्षैरवकीर्यमाणः} %||6-55-64||

\fourlineindentedshloka
{ततः स संज्ञामुपलभ्य कृच्छ्राद्-}
{बलीयसस्तस्य भुजान्तरस्थः}
{अवेक्षमाणः पुरराजमार्गम्}
{विचिन्तयामास मुहुर्महात्मा} %||6-55-65||

\fourlineindentedshloka
{एवं गृहीतेन कथं नु नाम}
{शक्यं मया सम्प्रति कर्तुमद्य}
{तथा करिष्यामि यथा हरीणाम्}
{भविष्यतीष्टं च हितं च कार्यम्} %||6-55-66||

\fourlineindentedshloka
{ततः कराग्रैः सहसा समेत्य}
{राजा हरीणाममरेन्द्रशत्रोः}
{नखैश्च कर्णौ दशनैश्च नासाम्}
{ददंश पार्श्वेषु च कुम्भकर्णम्} %||6-55-67||

\fourlineindentedshloka
{स कुम्भकर्णौ हृतकर्णनासो}
{विदारितस्तेन विमर्दितश्च}
{रोषाभिभूतः क्षतजार्द्रगात्रः}
{सुग्रीवमाविध्य पिपेष भूमौ} %||6-55-68||

\fourlineindentedshloka
{स भूतले भीमबलाभिपिष्टः}
{सुरारिभिस्तैरभिहन्यमानः}
{जगाम खं वेगवदभ्युपेत्य}
{पुनश्च रामेण समाजगाम} %||6-55-69||

\twolineshloka
{कर्णनासा विहीनस्य कुम्भकर्णो महाबलः}
{रराज शोणितोत्सिक्तो गिरिः प्रस्रवणैरिव} %||6-55-70||

\fourlineindentedshloka
{ततः स पुर्याः सहसा महात्मा}
{निष्क्रम्य तद्वानरसैन्यमुग्रम्}
{बभक्ष रक्षो युधि कुम्भकर्णः}
{प्रजा युगान्ताग्निरिव प्रदीप्तः} %||6-55-71||

\fourlineindentedshloka
{बुभुक्षितः शोणितमांसगृध्नुः}
{प्रविश्य तद्वानरसैन्यमुग्रम्}
{चखाद रक्षांसि हरीन्पिशाचान्}
{ऋक्षांश्च मोहाद्युधि कुम्भकर्णः} %||6-55-72||

\twolineshloka
{एकं द्वौ त्रीन्बहून्क्रुद्धो वानरान्सह राक्षसैः}
{समादायैकहस्तेन प्रचिक्षेप त्वरन्मुखे} %||6-55-73||

\threelineshloka
{सम्प्रस्रवंस्तदा मेदः शोणितं च महाबलः}
{वध्यमानो नगेन्द्राग्रैर्भक्षयामास वानरान्}
{ते भक्ष्यमाणा हरयो रामं जग्मुस्तदा गतिम्} %||6-55-74||

\twolineshloka
{तस्मिन्काले सुमित्रायाः पुत्रः परबलार्दनः}
{चकार लक्ष्मणः क्रुद्धो युद्धं परपुरंजयः} %||6-55-75||

\twolineshloka
{स कुम्भकर्णस्य शराञ्शरीरे सप्त वीर्यवान्}
{निचखानाददे चान्यान्विससर्ज च लक्ष्मणः} %||6-55-76||

\twolineshloka
{अतिक्रम्य च सौमित्रिं कुम्भकर्णो महाबलः}
{राममेवाभिदुद्राव दारयन्निव मेदिनीम्} %||6-55-77||

\twolineshloka
{अथ दाशरथी रामो रौद्रमस्त्रं प्रयोजयन्}
{कुम्भकर्णस्य हृदये ससर्ज निशिताञ्शरान्} %||6-55-78||

\twolineshloka
{तस्य रामेण विद्धस्य सहसाभिप्रधावतः}
{अङ्गारमिश्राः क्रुद्धस्य मुखान्निश्चेरुरर्चिषः} %||6-55-79||

\twolineshloka
{तस्योरसि निमग्नाश्च शरा बर्हिणवाससः}
{हस्ताच्चास्य परिभ्रष्टा पपातोर्व्यां महागदा} %||6-55-80||

\twolineshloka
{स निरायुधमात्मानं यदा मेने महाबलः}
{मुष्टिभ्यां चारणाभ्यां च चकार कदनं महत्} %||6-55-81||

\twolineshloka
{स बाणैरतिविद्धाङ्गः क्षतजेन समुक्षितः}
{रुधिरं परिसुस्राव गिरिः प्रस्रवणानिव} %||6-55-82||

\twolineshloka
{स तीव्रेण च कोपेन रुधिरेण च मूर्छितः}
{वानरान्राक्षसानृक्षान्खादन्विपरिधावति} %||6-55-83||

\twolineshloka
{तस्मिन्काले स धर्मात्मा लक्ष्मणो राममब्रवीत्}
{कुम्भकर्णवधे युक्तो योगान्परिमृशन्बहून्} %||6-55-84||

\twolineshloka
{नैवायं वानरान्राजन्न विजानाति राक्षसान्}
{मत्तः शोणितगन्धेन स्वान्परांश्चैव खादति} %||6-55-85||

\twolineshloka
{साध्वेनमधिरोहन्तु सर्वतो वानरर्षभाः}
{यूथपाश्च यथामुख्यास्तिष्ठन्त्वस्य समन्ततः} %||6-55-86||

\twolineshloka
{अप्ययं दुर्मतिः काले गुरुभारप्रपीडितः}
{प्रपतन्राक्षसो भूमौ नान्यान्हन्यात्प्लवङ्गमान्} %||6-55-87||

\twolineshloka
{तस्य तद्वचनं श्रुत्वा राजपुत्रस्य धीमतः}
{ते समारुरुहुर्हृष्टाः कुम्भकर्णं प्लवङ्गमाः} %||6-55-88||

\twolineshloka
{कुम्भकर्णस्तु सङ्क्रुद्धः समारूढः प्लवङ्गमैः}
{व्यधूनयत्तान्वेगेन दुष्टहस्तीव हस्तिपान्} %||6-55-89||

\twolineshloka
{तान्दृष्ट्वा निर्धूतान्रामो रुष्टोऽयमिति राक्षसः}
{समुत्पपात वेगेन धनुरुत्तममाददे} %||6-55-90||

\fourlineindentedshloka
{स चापमादाय भुजङ्गकल्पम्}
{दृढज्यमुग्रं तपनीयचित्रम्}
{हरीन्समाश्वास्य समुत्पपात}
{रामो निबद्धोत्तमतूणबाणः} %||6-55-91||

\twolineshloka
{स वानरगणैस्तैस्तु वृतः परमदुर्जयः}
{लक्ष्मणानुचरो रामः सम्प्रतस्थे महाबलः} %||6-55-92||

\twolineshloka
{स ददर्श महात्मानं किरीटिनमरिन्दमम्}
{शोणिताप्लुतसर्वाङ्गं कुम्भकर्णं महाबलम्} %||6-55-93||

\twolineshloka
{सर्वान्समभिधावन्तं यथारुष्टं दिशा गजम्}
{मार्गमाणं हरीन्क्रुद्धं राक्षसैः परिवारितम्} %||6-55-94||

\twolineshloka
{विन्ध्यमन्दरसङ्काशं काञ्चनाङ्गदभूषणम्}
{स्रवन्तं रुधिरं वक्त्राद्वर्षमेघमिवोत्थितम्} %||6-55-95||

\twolineshloka
{जिह्वया परिलिह्यन्तं शोणितं शोणितोक्षितम्}
{मृद्नन्तं वानरानीकं कालान्तकयमोपमम्} %||6-55-96||

\twolineshloka
{तं दृष्ट्वा राक्षसश्रेष्ठं प्रदीप्तानलवर्चसम्}
{विस्फारयामास तदा कार्मुकं पुरुषर्षभः} %||6-55-97||

\twolineshloka
{स तस्य चापनिर्घोषात्कुपितो नैरृतर्षभः}
{अमृष्यमाणस्तं घोषमभिदुद्राव राघवम्} %||6-55-98||

\fourlineindentedshloka
{ततस्तु वातोद्धतमेघकल्पम्}
{भुजङ्गराजोत्तमभोगबाहुम्}
{तमापतन्तं धरणीधराभम्}
{उवाच रामो युधि कुम्भकर्णम्} %||6-55-99||

\fourlineindentedshloka
{आगच्छ रक्षोऽधिपमा विषादम्}
{अवस्थितोऽहं प्रगृहीतचापः}
{अवेहि मां शक्रसपत्न रामम्}
{अयं मुहूर्ताद्भविता विचेताः} %||6-55-100||

\twolineshloka
{रामोऽयमिति विज्ञाय जहास विकृतस्वनम्}
{पातयन्निव सर्वेषां हृदयानि वनौकसाम्} %||6-55-101||

\twolineshloka
{प्रहस्य विकृतं भीमं स मेघस्वनितोपमम्}
{कुम्भकर्णो महातेजा राघवं वाक्यमब्रवीत्} %||6-55-102||

\twolineshloka
{नाहं विराधो विज्ञेयो न कबन्धः खरो न च}
{न वाली न च मारीचः कुम्भकर्णोऽहमागतः} %||6-55-103||

\twolineshloka
{पश्य मे मुद्गरं घोरं सर्वकालायसं महत्}
{अनेन निर्जिता देवा दानवाश्च मया पुरा} %||6-55-104||

\twolineshloka
{विकर्णनास इति मां नावज्ञातुं त्वमर्हसि}
{स्वल्पापि हि न मे पीडा कर्णनासाविनाशनात्} %||6-55-105||

\twolineshloka
{दर्शयेक्ष्वाकुशार्दूल वीर्यं गात्रेषु मे लघु}
{ततस्त्वां भक्षयिष्यामि दृष्टपौरुषविक्रमम्} %||6-55-106||

\fourlineindentedshloka
{स कुम्भकर्णस्य वचो निशम्य}
{रामः सुपुङ्खान्विससर्ज बाणान्}
{तैराहतो वज्रसमप्रवेगैर्-}
{न चुक्षुभे न व्यथते सुरारिः} %||6-55-107||

\fourlineindentedshloka
{यैः सायकैः सालवरा निकृत्ता}
{वाली हतो वानरपुङ्गवश्च}
{ते कुम्भकर्णस्य तदा शरीरम्}
{वज्रोपमा न व्यथयां प्रचक्रुः} %||6-55-108||

\fourlineindentedshloka
{स वारिधारा इव सायकांस्तान्}
{पिबञ्शरीरेण महेन्द्रशत्रुः}
{जघान रामस्य शरप्रवेगम्}
{व्याविध्य तं मुद्गरमुग्रवेगम्} %||6-55-109||

\fourlineindentedshloka
{ततस्तु रक्षः क्षतजानुलिप्तम्}
{वित्रासनं देवमहाचमूनाम्}
{व्याविध्य तं मुद्गरमुग्रवेगम्}
{विद्रावयामास चमूं हरीणाम्} %||6-55-110||

\fourlineindentedshloka
{वायव्यमादाय ततो वरास्त्रम्}
{रामः प्रचिक्षेप निशाचराय}
{समुद्गरं तेन जहार बाहुम्}
{स कृत्तबाहुस्तुमुलं ननाद} %||6-55-111||

\fourlineindentedshloka
{स तस्य बाहुर्गिरिशृङ्गकल्पः}
{समुद्गरो राघवबाणकृत्तः}
{पपात तस्मिन्हरिराजसैन्ये}
{जघान तां वानरवाहिनीं च} %||6-55-112||

\fourlineindentedshloka
{ते वानरा भग्नहतावशेषाः}
{पर्यन्तमाश्रित्य तदा विषण्णाः}
{प्रवेपिताङ्गा ददृशुः सुघोरम्}
{नरेन्द्ररक्षोऽधिपसंनिपातम्} %||6-55-113||

\fourlineindentedshloka
{स कुम्भकर्णोऽस्त्रनिकृत्तबाहुर्-}
{महान्निकृत्ताग्र इवाचलेन्द्रः}
{उत्पाटयामास करेण वृक्षम्}
{ततोऽभिदुद्राव रणे नरेन्द्रम्} %||6-55-114||

\fourlineindentedshloka
{तं तस्य बाहुं सह सालवृक्षम्}
{समुद्यतं पन्नगभोगकल्पम्}
{ऐन्द्रास्त्रयुक्तेन जहार रामो}
{बाणेन जाम्बूनदचित्रितेन} %||6-55-115||

\fourlineindentedshloka
{स कुम्भकर्णस्य भुजो निकृत्तः}
{पपात भूमौ गिरिसंनिकाशः}
{विवेष्टमानो निजघान वृक्षान्}
{शैलाञ्शिलावानरराक्षसांश्च} %||6-55-116||

\fourlineindentedshloka
{तं छिन्नबाहुं समवेक्ष्य रामः}
{समापतन्तं सहसा नदन्तम्}
{द्वावर्धचन्द्रौ निशितौ प्रगृह्य}
{चिच्छेद पादौ युधि राक्षसस्य} %||6-55-117||

\fourlineindentedshloka
{निकृत्तबाहुर्विनिकृत्तपादो}
{विदार्य वक्त्रं वडवामुखाभम्}
{दुद्राव रामं सहसाभिगर्जन्}
{राहुर्यथा चन्द्रमिवान्तरिक्षे} %||6-55-118||

\fourlineindentedshloka
{अपूरयत्तस्य मुखं शिताग्रै}
{रामः शरैर्हेमपिनद्धपुङ्खैः}
{स पूर्णवक्त्रो न शशाक वक्तुम्}
{चुकूज कृच्छ्रेण मुमोह चापि} %||6-55-119||

\fourlineindentedshloka
{अथाददे सूर्यमरीचिकल्पम्}
{स ब्रह्मदण्डान्तककालकल्पम्}
{अरिष्टमैन्द्रं निशितं सुपुङ्खम्}
{रामः शरं मारुततुल्यवेगम्} %||6-55-120||

\fourlineindentedshloka
{तं वज्रजाम्बूनदचारुपुङ्खम्}
{प्रदीप्तसूर्यज्वलनप्रकाशम्}
{महेन्द्रवज्राशनितुल्यवेगम्}
{रामः प्रचिक्षेप निशाचराय} %||6-55-121||

\fourlineindentedshloka
{स सायको राघवबाहुचोदितो}
{दिशः स्वभासा दश सम्प्रकाशयन्}
{विधूमवैश्वानरदीप्तदर्शनो}
{जगाम शक्राशनितुल्यविक्रमः} %||6-55-122||

\fourlineindentedshloka
{स तन्महापर्वतकूटसंनिभम्}
{विवृत्तदंष्ट्रं चलचारुकुण्डलम्}
{चकर्त रक्षोऽधिपतेः शिरस्तदा}
{यथैव वृत्रस्य पुरा पुरन्दरः} %||6-55-123||

\fourlineindentedshloka
{तद्रामबाणाभिहतं पपात}
{रक्षःशिरः पर्वतसंनिकाशम्}
{बभञ्ज चर्यागृहगोपुराणि}
{प्राकारमुच्चं तमपातयच्च} %||6-55-124||

\fourlineindentedshloka
{तच्चातिकायं हिमवत्प्रकाशम्}
{रक्षस्तदा तोयनिधौ पपात}
{ग्राहान्महामीनचयान्भुजङ्गमान्}
{ममर्द भूमिं च तथा विवेश} %||6-55-125||

\fourlineindentedshloka
{तस्मिर्हते ब्राह्मणदेवशत्रौ}
{महाबले संयति कुम्भकर्णे}
{चचाल भूर्भूमिधराश्च सर्वे}
{हर्षाच्च देवास्तुमुलं प्रणेदुः} %||6-55-126||

\fourlineindentedshloka
{ततस्तु देवर्षिमहर्षिपन्नगाः}
{सुराश्च भूतानि सुपर्णगुह्यकाः}
{सयक्षगन्धर्वगणा नभोगताः}
{प्रहर्षिता राम पराक्रमेण} %||6-55-127||

\fourlineindentedshloka
{प्रहर्षमीयुर्बहवस्तु वानराः}
{प्रबुद्धपद्मप्रतिमैरिवाननैः}
{अपूजयन्राघवमिष्टभागिनम्}
{हते रिपौ भीमबले दुरासदे} %||6-55-128||

\fourlineindentedshloka
{स कुम्भकर्णं सुरसैन्यमर्दनम्}
{महत्सु युद्धेष्वपराजितश्रमम्}
{ननन्द हत्वा भरताग्रजो रणे}
{महासुरं वृत्रमिवामराधिपः} %||6-56-129||


{॥इत्यार्षे श्रीमद्रामायणे वाल्मीकीये आदिकाव्ये युद्धकाण्डे पञ्चपञ्चाशत्तमः सर्गः॥}

\sect{षट्पञ्चाशत्तमः सर्गः}

\twolineshloka
{कुम्भकर्णं हतं दृष्ट्वा राघवेण महात्मना}
{राक्षसा राक्षसेन्द्राय रावणाय न्यवेदयन्} %||6-56-1||

\twolineshloka
{श्रुत्वा विनिहतं सङ्ख्ये कुम्भकर्णं महाबलम्}
{रावणः शोकसन्तप्तो मुमोह च पपात च} %||6-56-2||

\twolineshloka
{पितृव्यं निहतं श्रुत्वा देवान्तकनरान्तकौ}
{त्रिशिराश्चातिकायश्च रुरुदुः शोकपीडिताः} %||6-56-3||

\twolineshloka
{भ्रातरं निहतं श्रुत्वा रामेणाक्लिष्टकर्मणा}
{महोदरमहापार्श्वौ शोकाक्रान्तौ बभूवतुः} %||6-56-4||

\twolineshloka
{ततः कृच्छ्रात्समासाद्य संज्ञां राक्षसपुङ्गवः}
{कुम्भकर्णवधाद्दीनो विललाप स रावणः} %||6-56-5||

\twolineshloka
{हा वीर रिपुदर्पघ्न कुम्भकर्ण महाबल}
{शत्रुसैन्यं प्रताप्यैकः क्व मां सन्त्यज्य गच्छसि} %||6-56-6||

\twolineshloka
{इदानीं खल्वहं नास्मि यस्य मे पतितो भुजः}
{दक्षिणो यं समाश्रित्य न बिभेमि सुरासुरान्} %||6-56-7||

\twolineshloka
{कथमेवंविधो वीरो देवदानवदर्पहा}
{कालाग्निप्रतिमो ह्यद्य राघवेण रणे हतः} %||6-56-8||

\twolineshloka
{यस्य ते वज्रनिष्पेषो न कुर्याद्व्यसनं सदा}
{स कथं रामबाणार्तः प्रसुप्तोऽसि महीतले} %||6-56-9||

\twolineshloka
{एते देवगणाः सार्धमृषिभिर्गगने स्थिताः}
{निहतं त्वां रणे दृष्ट्वा निनदन्ति प्रहर्षिताः} %||6-56-10||

\twolineshloka
{ध्रुवमद्यैव संहृष्टा लब्धलक्ष्याः प्लवङ्गमाः}
{आरोक्ष्यन्तीह दुर्गाणि लङ्काद्वाराणि सर्वशः} %||6-56-11||

\twolineshloka
{राज्येन नास्ति मे कार्यं किं करिष्यामि सीतया}
{कुम्भकर्णविहीनस्य जीविते नास्ति मे रतिः} %||6-56-12||

\twolineshloka
{यद्यहं भ्रातृहन्तारं न हन्मि युधि राघवम्}
{ननु मे मरणं श्रेयो न चेदं व्यर्थजीवितम्} %||6-56-13||

\twolineshloka
{अद्यैव तं गमिष्यामि देशं यत्रानुजो मम}
{न हि भ्रातॄन्समुत्सृज्य क्षणं जीवितुमुत्सहे} %||6-56-14||

\twolineshloka
{देवा हि मां हसिष्यन्ति दृष्ट्वा पूर्वापकारिणम्}
{कथमिन्द्रं जयिष्यामि कुम्भकर्णहते त्वयि} %||6-56-15||

\twolineshloka
{तदिदं मामनुप्राप्तं विभीषणवचः शुभम्}
{यदज्ञानान्मया तस्य न गृहीतं महात्मनः} %||6-56-16||

\twolineshloka
{विभीषणवचो यावत्कुम्भकर्णप्रहस्तयोः}
{विनाशोऽयं समुत्पन्नो मां व्रीडयति दारुणः} %||6-56-17||

\twolineshloka
{तस्यायं कर्मणः प्रातो विपाको मम शोकदः}
{यन्मया धार्मिकः श्रीमान्स निरस्तो विभीषणः} %||6-56-18||

\fourlineindentedshloka
{इति बहुविधमाकुलान्तरात्मा}
{कृपणमतीव विलप्य कुम्भकर्णम्}
{न्यपतदथ दशाननो भृशार्तः}
{तमनुजमिन्द्ररिपुं हतं विदित्वा} %||6-57-19||


{॥इत्यार्षे श्रीमद्रामायणे वाल्मीकीये आदिकाव्ये युद्धकाण्डे षट्पञ्चाशत्तमः सर्गः॥}

\sect{सप्तपञ्चाशत्तमः सर्गः}

\twolineshloka
{एवं विलपमानस्य रावणस्य दुरात्मनः}
{श्रुत्वा शोकाभितप्तस्य त्रिशिरा वाक्यमब्रवीत्} %||6-57-1||

\twolineshloka
{एवमेव महावीर्यो हतो नस्तात मध्यमः}
{न तु सत्पुरुषा राजन्विलपन्ति यथा भवान्} %||6-57-2||

\twolineshloka
{नूनं त्रिभुवणस्यापि पर्याप्तस्त्वमसि प्रभो}
{स कस्मात्प्राकृत इव शोकस्यात्मानमीदृशम्} %||6-57-3||

\twolineshloka
{ब्रह्मदत्तास्ति ते शक्तिः कवचः सायको धनुः}
{सहस्रखरसंयुक्तो रथो मेघसमस्वनः} %||6-57-4||

\twolineshloka
{त्वयासकृद्विशस्त्रेण विशस्ता देवदानवाः}
{स सर्वायुधसम्पन्नो राघवं शास्तुमर्हसि} %||6-57-5||

\twolineshloka
{कामं तिष्ठ महाराजनिर्गमिष्याम्यहं रणम्}
{उद्धरिष्यामि ते शत्रून्गरुडः पन्नगानिह} %||6-57-6||

\twolineshloka
{शम्बरो देवराजेन नरको विष्णुना यथा}
{तथाद्य शयिता रामो मया युधि निपातितः} %||6-57-7||

\twolineshloka
{श्रुत्वा त्रिशिरसो वाक्यं रावणो राक्षसाधिपः}
{पुनर्जातमिवात्मानं मन्यते कालचोदितः} %||6-57-8||

\twolineshloka
{श्रुत्वा त्रिशिरसो वाक्यं देवान्तकनरान्तकौ}
{अतिकायश्च तेजस्वी बभूवुर्युद्धहर्षिताः} %||6-57-9||

\twolineshloka
{ततोऽहमहमित्येवं गर्जन्तो नैरृतर्षभाः}
{रावणस्य सुता वीराः शक्रतुल्यपराक्रमाः} %||6-57-10||

\twolineshloka
{अन्तरिक्षचराः सर्वे सर्वे माया विशारदाः}
{सर्वे त्रिदशदर्पघ्नाः सर्वे च रणदुर्मदाः} %||6-57-11||

\twolineshloka
{सर्वेऽस्त्रबलसम्पन्नाः सर्वे विस्तीर्ण कीर्तयः}
{सर्वे समरमासाद्य न श्रूयन्ते स्म निर्जिताः} %||6-57-12||

\twolineshloka
{सर्वेऽस्त्रविदुषो वीराः सर्वे युद्धविशारदाः}
{सर्वे प्रवरजिज्ञानाः सर्वे लब्धवरास्तथा} %||6-57-13||

\fourlineindentedshloka
{स तैस्तथा भास्करतुल्यवर्चसैः}
{सुतैर्वृतः शत्रुबलप्रमर्दनैः}
{रराज राजा मघवान्यथामरैर्-}
{वृतो महादानवदर्पनाशनैः} %||6-57-14||

\twolineshloka
{स पुत्रान्सम्परिष्वज्य भूषयित्वा च भूषणैः}
{आशीर्भिश्च प्रशस्ताभिः प्रेषयामास संयुगे} %||6-57-15||

\twolineshloka
{महोदरमहापार्श्वौ भ्रातरौ चापि रावणः}
{रक्षणार्थं कुमाराणां प्रेषयामास संयुगे} %||6-57-16||

\twolineshloka
{तेऽभिवाद्य महात्मानं रावणं रिपुरावणम्}
{कृत्वा प्रदक्षिणं चैव महाकायाः प्रतस्थिरे} %||6-57-17||

\twolineshloka
{सर्वौषधीभिर्गन्धैश्च समालभ्य महाबलाः}
{निर्जग्मुर्नैरृतश्रेष्ठाः षडेते युद्धकाङ्क्षिणः} %||6-57-18||

\twolineshloka
{ततः सुदर्शनं नाम नीलजीमूतसंनिभम्}
{ऐरावतकुले जातमारुरोह महोदरः} %||6-57-19||

\twolineshloka
{सर्वायुधसमायुक्तं तूणीभिश्च स्वलङ्कृतम्}
{रराज गजमास्थाय सवितेवास्तमूर्धनि} %||6-57-20||

\twolineshloka
{हयोत्तमसमायुक्तं सर्वायुधसमाकुलम्}
{आरुरोह रथश्रेष्ठं त्रिशिरा रावणात्मजः} %||6-57-21||

\twolineshloka
{त्रिशिरा रथमास्थाय विरराज धनुर्धरः}
{सविद्युदुल्कः सज्वालः सेन्द्रचाप इवाम्बुदः} %||6-57-22||

\twolineshloka
{त्रिभिः किरीटैस्त्रिशिराः शुशुभे स रथोत्तमे}
{हिमवानिव शैलेन्द्रस्त्रिभिः काञ्चनपर्वतैः} %||6-57-23||

\twolineshloka
{अतिकायोऽपि तेजस्वी राक्षसेन्द्रसुतस्तदा}
{आरुरोह रथश्रेष्ठं श्रेष्ठः सर्वधनुष्मताम्} %||6-57-24||

\twolineshloka
{सुचक्राक्षं सुसंयुक्तं सानुकर्षं सकूबरम्}
{तूणीबाणासनैर्दीप्तं प्रासासि परिघाकुलम्} %||6-57-25||

\twolineshloka
{स काञ्चनविचित्रेण किरीटेन विराजता}
{भूषणैश्च बभौ मेरुः प्रभाभिरिव भास्वरः} %||6-57-26||

\twolineshloka
{स रराज रथे तस्मिन्राजसूनुर्महाबलः}
{वृतो नैरृतशार्दूलैर्वज्रपाणिरिवामरैः} %||6-57-27||

\twolineshloka
{हयमुच्चैःश्रवः प्रख्यं श्वेतं कनकभूषणम्}
{मनोजवं महाकायमारुरोह नरान्तकः} %||6-57-28||

\twolineshloka
{गृहीत्वा प्रासमुक्लाभं विरराज नरान्तकः}
{शक्तिमादाय तेजस्वी गुहः शत्रुष्विवाहवे} %||6-57-29||

\twolineshloka
{देवान्तकः समादाय परिघं वज्रभूषणम्}
{परिगृह्य गिरिं दोर्भ्यां वपुर्विष्णोर्विडम्बयन्} %||6-57-30||

\twolineshloka
{महापार्श्वो महातेजा गदामादाय वीर्यवान्}
{विरराज गदापाणिः कुबेर इव संयुगे} %||6-57-31||

\twolineshloka
{ते प्रतस्थुर्महात्मानो बलैरप्रतिमैर्वृताः}
{सुरा इवामरावत्यां बलैरप्रतिमैर्वृताः} %||6-57-32||

\twolineshloka
{तान्गजैश्च तुरङ्गैश्च रथैश्चाम्बुदनिस्वनैः}
{अनुजग्मुर्महात्मानो राक्षसाः प्रवरायुधाः} %||6-57-33||

\twolineshloka
{ते विरेजुर्महात्मानो कुमाराः सूर्यवर्चसः}
{किरीटिनः श्रिया जुष्टा ग्रहा दीप्ता इवाम्बरे} %||6-57-34||

\twolineshloka
{प्रगृहीता बभौ तेषां छत्राणामावलिः सिता}
{शारदाभ्रप्रतीकाशां हंसावलिरिवाम्बरे} %||6-57-35||

\twolineshloka
{मरणं वापि निश्चित्य शत्रूणां वा पराजयम्}
{इति कृत्वा मतिं वीरा निर्जग्मुः संयुगार्थिनः} %||6-57-36||

\twolineshloka
{जगर्जुश्च प्रणेदुश्च चिक्षिपुश्चापि सायकान्}
{जहृषुश्च महात्मानो निर्यान्तो युद्धदुर्मदाः} %||6-57-37||

\twolineshloka
{क्ष्वेडितास्फोटनिनदैः सञ्चचालेव मेदिनी}
{रक्षसां सिंहनादैश्च पुस्फोटेव तदाम्बरम्} %||6-57-38||

\twolineshloka
{तेऽभिनिष्क्रम्य मुदिता राक्षसेन्द्रा महाबलाः}
{ददृशुर्वानरानीकं समुद्यतशिलानगम्} %||6-57-39||

\twolineshloka
{हरयोऽपि महात्मानो ददृशुर्नैरृतं बलम्}
{हस्त्यश्वरथसम्बाधं किङ्किणीशतनादितम्} %||6-57-40||

\twolineshloka
{नीलजीमूतसङ्काशं समुद्यतमहायुधम्}
{दीप्तानलरविप्रख्यैर्नैरृतैः सर्वतो वृतम्} %||6-57-41||

\twolineshloka
{तद्दृष्ट्वा बलमायान्तं लब्धलक्ष्याः प्लवङ्गमाः}
{समुद्यतमहाशैलाः सम्प्रणेदुर्मुहुर्मुहुः} %||6-57-42||

\fourlineindentedshloka
{ततः समुद्घुष्टरवं निशम्य}
{रक्षोगणा वानरयूथपानाम्}
{अमृष्यमाणाः परहर्षमुग्रम्}
{महाबला भीमतरं विनेदुः} %||6-57-43||

\twolineshloka
{ते राक्षसबलं घोरं प्रविश्य हरियूथपाः}
{विचेरुरुद्यतैः शैलैर्नगाः शिखरिणो यथा} %||6-57-44||

\twolineshloka
{केचिदाकाशमाविश्य केचिदुर्व्यां प्लवङ्गमाः}
{रक्षःसैन्येषु सङ्क्रुद्धाश्चेरुर्द्रुमशिलायुधाः} %||6-57-45||

\twolineshloka
{ते पादपशिलाशैलैश्चक्रुर्वृष्टिमनुत्तमाम्}
{बाणौघैर्वार्यमाणाश्च हरयो भीमविक्रमाः} %||6-57-46||

\twolineshloka
{सिंहनादान्विनेदुश्च रणे राक्षसवानराः}
{शिलाभिश्चूर्णयामासुर्यातुधानान्प्लवङ्गमाः} %||6-57-47||

\twolineshloka
{निजघ्नुः संयुगे क्रुद्धाः कवचाभरणावृतान्}
{केचिद्रथगतान्वीरान्गजवाजिगतानपि} %||6-57-48||

\threelineshloka
{निजघ्नुः सहसाप्लुत्य यातुधानान्प्लवङ्गमाः}
{शैलशृङ्गनिपातैश्च मुष्टिभिर्वान्तलोचनाः}
{चेलुः पेतुश्च नेदुश्च तत्र राक्षसपुङ्गवाः} %||6-57-49||

\twolineshloka
{ततः शैलैश्च खड्गैश्च विसृष्टैर्हरिराक्षसैः}
{मुहूर्तेनावृता भूमिरभवच्छोणिताप्लुता} %||6-57-50||

\twolineshloka
{विकीर्णपर्वताकारै रक्षोभिररिमर्दनैः}
{आक्षिप्ताः क्षिप्यमाणाश्च भग्नशूलाश्च वानरैः} %||6-57-51||

\twolineshloka
{वानरान्वानरैरेव जग्नुस्ते रजनीचराः}
{राक्षसान्राक्षसैरेव जघ्नुस्ते वानरा अपि} %||6-57-52||

\twolineshloka
{आक्षिप्य च शिलास्तेषां निजघ्नू राक्षसा हरीन्}
{तेषां चाच्छिद्य शस्त्राणि जघ्नू रक्षांसि वानराः} %||6-57-53||

\twolineshloka
{निजघ्नुः शैलशूलास्त्रैर्विभिदुश्च परस्परम्}
{सिंहनादान्विनेदुश्च रणे वानरराक्षसाः} %||6-57-54||

\twolineshloka
{छिन्नवर्मतनुत्राणा राक्षसा वानरैर्हताः}
{रुधिरं प्रस्रुतास्तत्र रससारमिव द्रुमाः} %||6-57-55||

\twolineshloka
{रथेन च रथं चापि वारणेन च वारणम्}
{हयेन च हयं केचिन्निजघ्नुर्वानरा रणे} %||6-57-56||

\twolineshloka
{क्षुरप्रैरर्धचन्द्रैश्च भल्लैश्च निशितैः शरैः}
{राक्षसा वानरेन्द्राणां चिच्छिदुः पादपाञ्शिलाः} %||6-57-57||

\twolineshloka
{विकीर्णैः पर्वताग्रैश्च द्रुमैश्छिन्नैश्च संयुगे}
{हतैश्च कपिरक्षोभिर्दुर्गमा वसुधाभवत्} %||6-57-58||

\fourlineindentedshloka
{तस्मिन्प्रवृत्ते तुमुले विमर्दे}
{प्रहृष्यमाणेषु वली मुखेषु}
{निपात्यमानेषु च राक्षसेषु}
{महर्षयो देवगणाश्च नेदुः} %||6-57-59||

\fourlineindentedshloka
{ततो हयं मारुततुल्यवेगम्}
{आरुह्य शक्तिं निशितां प्रगृह्य}
{नरान्तको वानरराजसैन्यम्}
{महार्णवं मीन इवाविवेश} %||6-57-60||

\fourlineindentedshloka
{स वानरान्सप्तशतानि वीरः}
{प्रासेन दीप्तेन विनिर्बिभेद}
{एकः क्षणेनेन्द्ररिपुर्महात्मा}
{जघान सैन्यं हरिपुङ्गवानाम्} %||6-57-61||

\twolineshloka
{ददृशुश्च महात्मानं हयपृष्ठे प्रतिष्ठितम्}
{चरन्तं हरिसैन्येषु विद्याधरमहर्षयः} %||6-57-62||

\twolineshloka
{स तस्य ददृशे मार्गो मांसशोणितकर्दमः}
{पतितैः पर्वताकारैर्वानरैरभिसंवृतः} %||6-57-63||

\twolineshloka
{यावद्विक्रमितुं बुद्धिं चक्रुः प्लवगपुङ्गवाः}
{तावदेतानतिक्रम्य निर्बिभेद नरान्तकः} %||6-57-64||

\twolineshloka
{ज्वलन्तं प्रासमुद्यम्य सङ्ग्रामान्ते नरान्तकः}
{ददाह हरिसैन्यानि वनानीव विभावसुः} %||6-57-65||

\twolineshloka
{यावदुत्पाटयामासुर्वृक्षाञ्शैलान्वनौकसः}
{तावत्प्रासहताः पेतुर्वज्रकृत्ता इवाचलाः} %||6-57-66||

\twolineshloka
{दिक्षु सर्वासु बलवान्विचचार नरान्तकः}
{प्रमृद्नन्सर्वतो युद्धे प्रावृट्काले यथानिलः} %||6-57-67||

\twolineshloka
{न शेकुर्धावितुं वीरा न स्थातुं स्पन्दितुं कुतः}
{उत्पतन्तं स्थितं यान्तं सर्वान्विव्याध वीर्यवान्} %||6-57-68||

\twolineshloka
{एकेनान्तककल्पेन प्रासेनादित्यतेजसा}
{भिन्नानि हरिसैन्यानि निपेतुर्धरणीतले} %||6-57-69||

\twolineshloka
{वज्रनिष्पेषसदृशं प्रासस्याभिनिपातनम्}
{न शेकुर्वानराः सोढुं ते विनेदुर्महास्वनम्} %||6-57-70||

\twolineshloka
{पततां हरिवीराणां रूपाणि प्रचकाशिरे}
{वज्रभिन्नाग्रकूटानां शैलानां पततामिव} %||6-57-71||

\twolineshloka
{ये तु पूर्वं महात्मानः कुम्भकर्णेन पातिताः}
{तेऽस्वस्था वानरश्रेष्ठाः सुग्रीवमुपतस्थिरे} %||6-57-72||

\twolineshloka
{विप्रेक्षमाणः सुग्रीवो ददर्श हरिवाहिनीम्}
{नरान्तकभयत्रस्तां विद्रवन्तीमितस्ततः} %||6-57-73||

\twolineshloka
{विद्रुतां वाहिनीं दृष्ट्वा स ददर्श नरान्तकम्}
{गृहीतप्रासमायान्तं हयपृष्ठे प्रतिष्ठितम्} %||6-57-74||

\twolineshloka
{अथोवाच महातेजाः सुग्रीवो वानराधिपः}
{कुमारमङ्गदं वीरं शक्रतुल्यपराक्रमम्} %||6-57-75||

\twolineshloka
{गच्छैनं राक्षसं वीर योऽसौ तुरगमास्थितः}
{क्षोभयन्तं हरिबलं क्षिप्रं प्राणैर्वियोजय} %||6-57-76||

\twolineshloka
{स भर्तुर्वचनं श्रुत्वा निष्पपाताङ्गदस्तदा}
{अनीकान्मेघसङ्काशान्मेघानीकादिवांशुमान्} %||6-57-77||

\twolineshloka
{शैलसङ्घातसङ्काशो हरीणामुत्तमोऽङ्गदः}
{रराजाङ्गदसंनद्धः सधातुरिव पर्वतः} %||6-57-78||

\twolineshloka
{निरायुधो महातेजाः केवलं नखदंष्ट्रवान्}
{नरान्तकमभिक्रम्य वालिपुत्रोऽब्रवीद्वचः} %||6-57-79||

\twolineshloka
{तिष्ठ किं प्राकृतैरेभिर्हरिभिस्त्वं करिष्यसि}
{अस्मिन्वज्रसमस्पर्शे प्रासं क्षिप ममोरसि} %||6-57-80||

\twolineshloka
{अङ्गदस्य वचः श्रुत्वा प्रचुक्रोध नरान्तकः}
{सन्दश्य दशनैरोष्ठं निश्वस्य च भुजङ्गवत्} %||6-57-81||

\fourlineindentedshloka
{स प्रासमाविध्य तदाङ्गदाय}
{समुज्ज्वलन्तं सहसोत्ससर्ज}
{स वालिपुत्रोरसि वज्रकल्पे}
{बभूव भग्नो न्यपतच्च भूमौ} %||6-57-82||

\fourlineindentedshloka
{तं प्रासमालोक्य तदा विभग्नम्}
{सुपर्णकृत्तोरगभोगकल्पम्}
{तलं समुद्यम्य स वालिपुत्रः}
{तुरङ्गमस्याभिजघान मूर्ध्नि} %||6-57-83||

\fourlineindentedshloka
{निमग्नपादः स्फुटिताक्षि तारो}
{निष्क्रान्तजिह्वोऽचलसंनिकाशः}
{स तस्य वाजी निपपात भूमौ}
{तलप्रहारेण विकीर्णमूर्धा} %||6-57-84||

\fourlineindentedshloka
{नरान्तकः क्रोधवशं जगाम}
{हतं तुरगं पतितं निरीक्ष्य}
{स मुष्टिमुद्यम्य महाप्रभावो}
{जघान शीर्षे युधि वालिपुत्रम्} %||6-57-85||

\fourlineindentedshloka
{अथाङ्गदो मुष्टिविभिन्नमूर्धा}
{सुस्राव तीव्रं रुधिरं भृशोष्णम्}
{मुहुर्विजज्वाल मुमोह चापि}
{संज्ञां समासाद्य विसिष्मिये च} %||6-57-86||

\fourlineindentedshloka
{अथाङ्गदो वज्रसमानवेगम्}
{संवर्त्य मुष्टिं गिरिशृङ्गकल्पम्}
{निपातयामास तदा महात्मा}
{नरान्तकस्योरसि वालिपुत्रः} %||6-57-87||

\fourlineindentedshloka
{स मुष्टिनिष्पिष्टविभिन्नवक्षा}
{ज्वालां वमञ्शोणितदिग्धगात्रः}
{नरान्तको भूमितले पपात}
{यथाचलो वज्रनिपातभग्नः} %||6-57-88||

\fourlineindentedshloka
{अथान्तरिक्षे त्रिदशोत्तमानाम्}
{वनौकसां चैव महाप्रणादः}
{बभूव तस्मिन्निहतेऽग्र्यवीरे}
{नरान्तके वालिसुतेन सङ्ख्ये} %||6-57-89||

\fourlineindentedshloka
{अथाङ्गदो राममनः प्रहर्षणम्}
{सुदुष्करं तं कृतवान्हि विक्रमम्}
{विसिष्मिये सोऽप्यतिवीर्य विक्रमः}
{पुनश्च युद्धे स बभूव हर्षितः} %||6-58-90||


{॥इत्यार्षे श्रीमद्रामायणे वाल्मीकीये आदिकाव्ये युद्धकाण्डे सप्तपञ्चाशत्तमः सर्गः॥}

\sect{अष्टपञ्चाशत्तमः सर्गः}

\twolineshloka
{नरान्तकं हतं दृष्ट्वा चुक्रुशुर्नैरृतर्षभाः}
{देवान्तकस्त्रिमूर्धा च पौलस्त्यश्च महोदरः} %||6-58-1||

\twolineshloka
{आरूढो मेघसङ्काशं वारणेन्द्रं महोदरः}
{वालिपुत्रं महावीर्यमभिदुद्राव वीर्यवान्} %||6-58-2||

\twolineshloka
{भ्रातृव्यसनसन्तप्तस्तदा देवान्तको बली}
{आदाय परिघं दीप्तमङ्गदं समभिद्रवत्} %||6-58-3||

\twolineshloka
{रथमादित्यसङ्काशं युक्तं परमवाजिभिः}
{आस्थाय त्रिशिरा वीरो वालिपुत्रमथाभ्ययात्} %||6-58-4||

\twolineshloka
{स त्रिभिर्देवदर्पघ्नैर्नैरृतेन्द्रैरभिद्रुतः}
{वृक्षमुत्पाटयामास महाविटपमङ्गदः} %||6-58-5||

\twolineshloka
{देवान्तकाय तं वीरश्चिक्षेप सहसाङ्गदः}
{महावृक्षं महाशाखं शक्रो दीप्तमिवाशनिम्} %||6-58-6||

\twolineshloka
{त्रिशिरास्तं प्रचिच्छेद शरैराशीविषोपमैः}
{स वृक्षं कृत्तमालोक्य उत्पपात ततोऽङ्गदः} %||6-58-7||

\twolineshloka
{स ववर्ष ततो वृक्षाञ्शिलाश्च कपिकुञ्जरः}
{तान्प्रचिच्छेद सङ्क्रुद्धस्त्रिशिरा निशितैः शरैः} %||6-58-8||

\twolineshloka
{परिघाग्रेण तान्वृक्षान्बभञ्ज च सुरान्तकः}
{त्रिशिराश्चाङ्गदं वीरमभिदुद्राव सायकैः} %||6-58-9||

\twolineshloka
{गजेन समभिद्रुत्य वालिपुत्रं महोदरः}
{जघानोरसि सङ्क्रुद्धस्तोमरैर्वज्रसंनिभैः} %||6-58-10||

\twolineshloka
{देवान्तकश्च सङ्क्रुद्धः परिघेण तदाङ्गदम्}
{उपगम्याभिहत्याशु व्यपचक्राम वेगवान्} %||6-58-11||

\twolineshloka
{स त्रिभिर्नैरृतश्रेष्ठैर्युगपत्समभिद्रुतः}
{न विव्यथे महातेजा वालिपुत्रः प्रतापवान्} %||6-58-12||

\twolineshloka
{तलेन भृशमुत्पत्य जघानास्य महागजम्}
{पेततुर्लोचने तस्य विननाद स वारणः} %||6-58-13||

\twolineshloka
{विषाणं चास्य निष्कृष्य वालिपुत्रो महाबलः}
{देवान्तकमभिद्रुत्य ताडयामास संयुगे} %||6-58-14||

\twolineshloka
{स विह्वलितसर्वाङ्गो वातोद्धत इव द्रुमः}
{लाक्षारससवर्णं च सुस्राव रुधिरं मुखात्} %||6-58-15||

\twolineshloka
{अथाश्वास्य महातेजाः कृच्छ्राद्देवान्तको बली}
{आविध्य परिघं घोरमाजघान तदाङ्गदम्} %||6-58-16||

\twolineshloka
{परिघाभिहतश्चापि वानरेन्द्रात्मजस्तदा}
{जानुभ्यां पतितो भूमौ पुनरेवोत्पपात ह} %||6-58-17||

\twolineshloka
{समुत्पतन्तं त्रिशिरास्त्रिभिराशीविषोपमैः}
{घोरैर्हरिपतेः पुत्रं ललाटेऽभिजघान ह} %||6-58-18||

\twolineshloka
{ततोऽङ्गदं परिक्षिप्तं त्रिभिर्नैरृतपुङ्गवैः}
{हनूमानपि विज्ञाय नीलश्चापि प्रतस्थतुः} %||6-58-19||

\twolineshloka
{ततश्चिक्षेप शैलाग्रं नीलस्त्रिशिरसे तदा}
{तद्रावणसुतो धीमान्बिभेद निशितैः शरैः} %||6-58-20||

\twolineshloka
{तद्बाणशतनिर्भिन्नं विदारितशिलातलम्}
{सविस्फुलिङ्गं सज्वालं निपपात गिरेः शिरः} %||6-58-21||

\twolineshloka
{ततो जृम्भितमालोक्य हर्षाद्देवान्तकस्तदा}
{परिघेणाभिदुद्राव मारुतात्मजमाहवे} %||6-58-22||

\twolineshloka
{तमापतन्तमुत्पत्य हनूमान्मारुतात्मजः}
{आजघान तदा मूर्ध्नि वज्रवेगेन मुष्टिना} %||6-58-23||

\fourlineindentedshloka
{स मुष्टिनिष्पिष्टविकीर्णमूर्धा}
{निर्वान्तदन्ताक्षिविलम्बिजिह्वः}
{देवान्तको राक्षसराजसूनुर्-}
{गतासुरुर्व्यां सहसा पपात} %||6-58-24||

\fourlineindentedshloka
{तस्मिन्हते राक्षसयोधमुख्ये}
{महाबले संयति देवशत्रौ}
{क्रुद्धस्त्रिमूर्धा निशिताग्रमुग्रम्}
{ववर्ष नीलोरसि बाणवर्षम्} %||6-58-25||

\fourlineindentedshloka
{स तैः शरौघैरभिवर्ष्यमाणो}
{विभिन्नगात्रः कपिसैन्यपालः}
{नीलो बभूवाथ विसृष्टगात्रो}
{विष्टम्भितस्तेन महाबलेन} %||6-58-26||

\fourlineindentedshloka
{ततस्तु नीलः प्रतिलभ्य संज्ञाम्}
{शैलं समुत्पाट्य सवृक्षषण्डम्}
{ततः समुत्पत्य भृशोग्रवेगो}
{महोदरं तेन जघान मूर्ध्नि} %||6-58-27||

\fourlineindentedshloka
{ततः स शैलाभिनिपातभग्नो}
{महोदरस्तेन सह द्विपेन}
{विपोथितो भूमितले गतासुः}
{पपात वर्जाभिहतो यथाद्रिः} %||6-58-28||

\twolineshloka
{पितृव्यं निहतं दृष्ट्वा त्रिशिराश्चापमाददे}
{हनूमन्तं च सङ्क्रुद्धो विव्याध निशितैः शरैः} %||6-58-29||

\twolineshloka
{हनूमांस्तु समुत्पत्य हयांस्त्रिशिरसस्तदा}
{विददार नखैः क्रुद्धो गजेन्द्रं मृगराडिव} %||6-58-30||

\twolineshloka
{अथ शक्तिं समादाय कालरात्रिमिवान्तकः}
{चिक्षेपानिलपुत्राय त्रिशिरा रावणात्मजः} %||6-58-31||

\twolineshloka
{दिवि क्षिप्तामिवोल्कां तां शक्तिं क्षिप्तामसङ्गताम्}
{गृहीत्वा हरिशार्दूलो बभञ्ज च ननाद च} %||6-58-32||

\twolineshloka
{तां दृष्ट्वा घोरसङ्काशां शक्तिं भग्नां हनूमता}
{प्रहृष्टा वानरगणा विनेदुर्जलदा इव} %||6-58-33||

\twolineshloka
{ततः खड्गं समुद्यम्य त्रिशिरा राक्षसोत्तमः}
{निचखान तदा रोषाद्वानरेन्द्रस्य वक्षसि} %||6-58-34||

\twolineshloka
{खड्गप्रहाराभिहतो हनूमान्मारुतात्मजः}
{आजघान त्रिमूर्धानं तलेनोरसि वीर्यवान्} %||6-58-35||

\twolineshloka
{स तलभिहतस्तेन स्रस्तहस्ताम्बरो भुवि}
{निपपात महातेजास्त्रिशिरास्त्यक्तचेतनः} %||6-58-36||

\twolineshloka
{स तस्य पततः खड्गं समाच्छिद्य महाकपिः}
{ननाद गिरिसङ्काशस्त्रासयन्सर्वनैरृतान्} %||6-58-37||

\twolineshloka
{अमृष्यमाणस्तं घोषमुत्पपात निशाचरः}
{उत्पत्य च हनूमन्तं ताडयामास मुष्टिना} %||6-58-38||

\twolineshloka
{तेन मुष्टिप्रहारेण सञ्चुकोप महाकपिः}
{कुपितश्च निजग्राह किरीटे राक्षसर्षभम्} %||6-58-39||

\fourlineindentedshloka
{स तस्य शीर्षाण्यसिना शितेन}
{किरीटजुष्टानि सकुण्डलानि}
{क्रुद्धः प्रचिच्छेद सुतोऽनिलस्य}
{त्वष्टुः सुतस्येव शिरांसि शक्रः} %||6-58-40||

\fourlineindentedshloka
{तान्यायताक्षाण्यगसंनिभानि}
{प्रदीप्तवैश्वानरलोचनानि}
{पेतुः शिरांसीन्द्ररिपोर्धरण्याम्}
{ज्योतींषि मुक्तानि यथार्कमार्गात्} %||6-58-41||

\fourlineindentedshloka
{तस्मिन्हते देवरिपौ त्रिशीर्षे}
{हनूमत शक्रपराक्रमेण}
{नेदुः प्लवङ्गाः प्रचचाल भूमी}
{रक्षांस्यथो दुद्रुविरे समन्तात्} %||6-58-42||

\twolineshloka
{हतं त्रिशिरसं दृष्ट्वा तथैव च महोदरम्}
{हतौ प्रेक्ष्य दुराधर्षौ देवान्तकनरान्तकौ} %||6-58-43||

\twolineshloka
{चुकोप परमामर्षी महापार्श्वो महाबलः}
{जग्राहार्चिष्मतीं चापि गदां सर्वायसीं शुभाम्} %||6-58-44||

\twolineshloka
{हेमपट्टपरिक्षिप्तां मांसशोणितलेपनाम्}
{विराजमानां वपुषा शत्रुशोणितरञ्जिताम्} %||6-58-45||

\twolineshloka
{तेजसा सम्प्रदीप्ताग्रां रक्तमाल्यविभूषिताम्}
{ऐरावतमहापद्मसार्वभौम भयावहाम्} %||6-58-46||

\twolineshloka
{गदामादाय सङ्क्रुद्धो महापार्श्वो महाबलः}
{हरीन्समभिदुद्राव युगान्ताग्निरिव ज्वलन्} %||6-58-47||

\twolineshloka
{अथर्षयः समुत्पत्य वानरो रवणानुजम्}
{महापार्श्वमुपागम्य तस्थौ तस्याग्रतो बली} %||6-58-48||

\twolineshloka
{तं पुरस्तात्स्थितं दृष्ट्वा वानरं पर्वतोपमम्}
{आजघानोरसि क्रुद्धो गदया वज्रकल्पया} %||6-58-49||

\twolineshloka
{स तयाभिहतस्तेन गदया वानरर्षभः}
{भिन्नवक्षाः समाधूतः सुस्राव रुधिरं बहु} %||6-58-50||

\twolineshloka
{स सम्प्राप्य चिरात्संज्ञामृषभो वानरर्षभः}
{क्रुद्धो विस्फुरमाणौष्ठो महापार्श्वमुदैक्षत} %||6-58-51||

\twolineshloka
{तां गृहीत्वा गदां भीमामाविध्य च पुनः पुनः}
{मत्तानीकं महापार्श्वं जघान रणमूर्धनि} %||6-58-52||

\twolineshloka
{स स्वया गदया भिन्नो विकीर्णदशनेक्षणः}
{निपपात महापार्श्वो वज्राहत इवाचलः} %||6-58-53||

\fourlineindentedshloka
{तस्मिन्हते भ्रातरि रावणस्य}
{तन्नैरृतानां बलमर्णवाभम्}
{त्यक्तायुधं केवलजीवितार्थम्}
{दुद्राव भिन्नार्णवसंनिकाशम्} %||6-59-54||


{॥इत्यार्षे श्रीमद्रामायणे वाल्मीकीये आदिकाव्ये युद्धकाण्डे अष्टपञ्चाशत्तमः सर्गः॥}

\sect{एकोनषष्टितमः सर्गः}

\twolineshloka
{स्वबलं व्यथितं दृष्ट्वा तुमुलं लोमहर्षणम्}
{भ्रातॄंश्च निहतान्दृष्ट्वा शक्रतुल्यपराक्रमान्} %||6-59-1||

\twolineshloka
{पितृव्यौ चापि सन्दृश्य समरे संनिषूदितौ}
{महोदरमहापार्श्वौ भ्रातरौ राक्षसर्षभौ} %||6-59-2||

\twolineshloka
{चुकोप च महातेजा ब्रह्मदत्तवरो युधि}
{अतिकायोऽद्रिसङ्काशो देवदानवदर्पहा} %||6-59-3||

\twolineshloka
{स भास्करसहस्रस्य सङ्घातमिव भास्वरम्}
{रथमास्थाय शक्रारिरभिदुद्राव वानरान्} %||6-59-4||

\twolineshloka
{स विस्फार्य महच्चापं किरीटी मृष्टकुण्डलः}
{नाम विश्रावयामास ननाद च महास्वनम्} %||6-59-5||

\twolineshloka
{तेन सिंहप्रणादेन नामविश्रावणेन च}
{ज्याशब्देन च भीमेन त्रासयामास वानरान्} %||6-59-6||

\twolineshloka
{ते तस्य रूपमालोक्य यथा विष्णोस्त्रिविक्रमे}
{भयार्ता वानराः सर्वे विद्रवन्ति दिशो दश} %||6-59-7||

\twolineshloka
{तेऽतिकायं समासाद्य वानरा मूढचेतसः}
{शरण्यं शरणं जग्मुर्लक्ष्मणाग्रजमाहवे} %||6-59-8||

\twolineshloka
{ततोऽतिकायं काकुत्स्थो रथस्थं पर्वतोपमम्}
{ददर्श धन्विनं दूराद्गर्जन्तं कालमेघवत्} %||6-59-9||

\twolineshloka
{स तं दृष्ट्वा महात्मानं राघवस्तु सुविस्मितः}
{वानरान्सान्त्वयित्वा तु विभीषणमुवाच ह} %||6-59-10||

\twolineshloka
{कोऽसौ पर्वतसङ्काशो धनुष्मान्हरिलोचनः}
{युक्ते हयसहस्रेण विशाले स्यन्दने स्थितः} %||6-59-11||

\twolineshloka
{य एष निशितैः शूलैः सुतीक्ष्णैः प्रासतोमरैः}
{अर्चिष्मद्भिर्वृतो भाति भूतैरिव महेश्वरः} %||6-59-12||

\twolineshloka
{कालजिह्वाप्रकाशाभिर्य एषोऽभिविराजते}
{आवृतो रथशक्तीभिर्विद्युद्भिरिव तोयदः} %||6-59-13||

\twolineshloka
{धनूंसि चास्य सज्यानि हेमपृष्ठानि सर्वशः}
{शोभयन्ति रथश्रेष्ठं शक्रपातमिवाम्बरम्} %||6-59-14||

\twolineshloka
{क एष रक्षः शार्दूलो रणभूमिं विराजयन्}
{अभ्येति रथिनां श्रेष्ठो रथेनादित्यतेजसा} %||6-59-15||

\twolineshloka
{ध्वजशृङ्गप्रतिष्ठेन राहुणाभिविराजते}
{सूर्यरश्मिप्रभैर्बाणैर्दिशो दश विराजयन्} %||6-59-16||

\twolineshloka
{त्रिणतं मेघनिर्ह्रादं हेमपृष्ठमलङ्कृतम्}
{शतक्रतुधनुःप्रख्यं धनुश्चास्य विराजते} %||6-59-17||

\twolineshloka
{सध्वजः सपताकश्च सानुकर्षो महारथः}
{चतुःसादिसमायुक्तो मेघस्तनितनिस्वनः} %||6-59-18||

\twolineshloka
{विंशतिर्दश चाष्टौ च तूणीररथमास्थिताः}
{कार्मुकाणि च भीमानि ज्याश्च काञ्चनपिङ्गलाः} %||6-59-19||

\twolineshloka
{द्वौ च खड्गौ रथगतौ पार्श्वस्थौ पार्श्वशोभिनौ}
{चतुर्हस्तत्सरुचितौ व्यक्तहस्तदशायतौ} %||6-59-20||

\twolineshloka
{रक्तकण्ठगुणो धीरो महापर्वतसंनिभः}
{कालः कालमहावक्त्रो मेघस्थ इव भास्करः} %||6-59-21||

\twolineshloka
{काञ्चनाङ्गदनद्धाभ्यां भुजाभ्यामेष शोभते}
{शृङ्गाभ्यामिव तुङ्गाभ्यां हिमवान्पर्वतोत्तमः} %||6-59-22||

\twolineshloka
{कुण्डलाभ्यां तु यस्यैतद्भाति वक्त्रं शुभेक्षणम्}
{पुनर्वस्वन्तरगतं पूर्णबिम्बमिवैन्दवम्} %||6-59-23||

\twolineshloka
{आचक्ष्व मे महाबाहो त्वमेनं राक्षसोत्तमम्}
{यं दृष्ट्वा वानराः सर्वे भयार्ता विद्रुता दिशः} %||6-59-24||

\twolineshloka
{स पृष्ठो राजपुत्रेण रामेणामिततेजसा}
{आचचक्षे महातेजा राघवाय विभीषणः} %||6-59-25||

\twolineshloka
{दशग्रीवो महातेजा राजा वैश्रवणानुजः}
{भीमकर्मा महोत्साहो रावणो राक्षसाधिपः} %||6-59-26||

\twolineshloka
{तस्यासीद्वीर्यवान्पुत्रो रावणप्रतिमो रणे}
{वृद्धसेवी श्रुतधरः सर्वास्त्रविदुषां वरः} %||6-59-27||

\twolineshloka
{अश्वपृष्ठे रथे नागे खड्गे धनुषि कर्षणे}
{भेदे सान्त्वे च दाने च नये मन्त्रे च सम्मतः} %||6-59-28||

\twolineshloka
{यस्य बाहुं समाश्रित्य लङ्का भवति निर्भया}
{तनयं धान्यमालिन्या अतिकायमिमं विदुः} %||6-59-29||

\twolineshloka
{एतेनाराधितो ब्रह्मा तपसा भावितात्मना}
{अस्त्राणि चाप्यवाप्तानि रिपवश्च पराजिताः} %||6-59-30||

\twolineshloka
{सुरासुरैरवध्यत्वं दत्तमस्मै स्वयम्भुवा}
{एतच्च कवचं दिव्यं रथश्चैषोऽर्कभास्करः} %||6-59-31||

\twolineshloka
{एतेन शतशो देवा दानवाश्च पराजिताः}
{रक्षितानि च रक्षामि यक्षाश्चापि निषूदिताः} %||6-59-32||

\twolineshloka
{वज्रं विष्टम्भितं येन बाणैरिन्द्रस्य धीमतः}
{पाशः सलिलराजस्य युद्धे प्रतिहतस्तथा} %||6-59-33||

\twolineshloka
{एषोऽतिकायो बलवान्राक्षसानामथर्षभः}
{रावणस्य सुतो धीमान्देवदनव दर्पहा} %||6-59-34||

\twolineshloka
{तदस्मिन्क्रियतां यत्नः क्षिप्रं पुरुषपुङ्गव}
{पुरा वानरसैन्यानि क्षयं नयति सायकैः} %||6-59-35||

\twolineshloka
{ततोऽतिकायो बलवान्प्रविश्य हरिवाहिनीम्}
{विस्फारयामास धनुर्ननाद च पुनः पुनः} %||6-59-36||

\twolineshloka
{तं भीमवपुषं दृष्ट्वा रथस्थं रथिनां वरम्}
{अभिपेतुर्महात्मानो ये प्रधानाः प्लवङ्गमाः} %||6-59-37||

\twolineshloka
{कुमुदो द्विविदो मैन्दो नीलः शरभ एव च}
{पादपैर्गिरिशृङ्गैश्च युगपत्समभिद्रवन्} %||6-59-38||

\twolineshloka
{तेषां वृक्षांश्च शैलांश्च शरैः काञ्चनभूषणैः}
{अतिकायो महातेजाश्चिच्छेदास्त्रविदां वरः} %||6-59-39||

\twolineshloka
{तांश्चैव सरान्स हरीञ्शरैः सर्वायसैर्बली}
{विव्याधाभिमुखः सङ्ख्ये भीमकायो निशाचरः} %||6-59-40||

\twolineshloka
{तेऽर्दिता बाणबर्षेण भिन्नगात्राः प्लवङ्गमाः}
{न शेकुरतिकायस्य प्रतिकर्तुं महारणे} %||6-59-41||

\twolineshloka
{तत्सैन्यं हरिवीराणां त्रासयामास राक्षसः}
{मृगयूथमिव क्रुद्धो हरिर्यौवनमास्थितः} %||6-59-42||

\fourlineindentedshloka
{स राषसेन्द्रो हरिसैन्यमध्ये}
{नायुध्यमानं निजघान कञ्चित्}
{उपेत्य रामं सधनुः कलापी}
{सगर्वितं वाक्यमिदं बभाषे} %||6-59-43||

\fourlineindentedshloka
{रथे स्थितोऽहं शरचापपाणिर्-}
{न प्राकृतं कञ्चन योधयामि}
{यस्यास्ति शक्तिर्व्यवसाय युक्ता}
{ददातुं मे क्षिप्रमिहाद्य युद्धम्} %||6-59-44||

\fourlineindentedshloka
{तत्तस्य वाक्यं ब्रुवतो निशम्य}
{चुकोप सौमित्रिरमित्रहन्ता}
{अमृष्यमाणश्च समुत्पपात}
{जग्राह चापं च ततः स्मयित्वा} %||6-59-45||

\twolineshloka
{क्रुद्धः सौमित्रिरुत्पत्य तूणादाक्षिप्य सायकम्}
{पुरस्तादतिकायस्य विचकर्ष महद्धनुः} %||6-59-46||

\twolineshloka
{पूरयन्स महीं शैलानाकाशं सागरं दिशः}
{ज्याशब्दो लक्ष्मणस्योग्रस्त्रासयन्रजनीचरान्} %||6-59-47||

\twolineshloka
{सौमित्रेश्चापनिर्घोषं श्रुत्वा प्रतिभयं तदा}
{विसिष्मिये महातेजा राक्षसेन्द्रात्मजो बली} %||6-59-48||

\twolineshloka
{अथातिकायः कुपितो दृष्ट्वा लक्ष्मणमुत्थितम्}
{आदाय निशितं बाणमिदं वचनमब्रवीत्} %||6-59-49||

\twolineshloka
{बालस्त्वमसि सौमित्रे विक्रमेष्वविचक्षणः}
{गच्छ किं कालसदृशं मां योधयितुमिच्छसि} %||6-59-50||

\twolineshloka
{न हि मद्बाहुसृष्टानामस्त्राणां हिमवानपि}
{सोढुमुत्सहते वेगमन्तरिक्षमथो मही} %||6-59-51||

\twolineshloka
{सुखप्रसुप्तं कालाग्निं प्रबोधयितुमिच्छसि}
{न्यस्य चापं निवर्तस्व मा प्राणाञ्जहि मद्गतः} %||6-59-52||

\twolineshloka
{अथ वा त्वं प्रतिष्टब्धो न निवर्तितुमिच्छसि}
{तिष्ठ प्राणान्परित्यज्य गमिष्यसि यमक्षयम्} %||6-59-53||

\twolineshloka
{पश्य मे निशितान्बाणानरिदर्पनिषूदनान्}
{ईश्वरायुधसङ्काशांस्तप्तकाञ्चनभूषणान्} %||6-59-54||

\twolineshloka
{एष ते सर्पसङ्काशो बाणः पास्यति शोणितम्}
{मृगराज इव क्रुद्धो नागराजस्य शोणितम्} %||6-59-55||

\fourlineindentedshloka
{श्रुत्वातिकायस्य वचः सरोषम्}
{सगर्वितं संयति राजपुत्रः}
{स सञ्चुकोपातिबलो बृहच्छ्रीः}
{उवाच वाक्यं च ततो महार्थम्} %||6-59-56||

\fourlineindentedshloka
{न वाक्यमात्रेण भवान्प्रधानो}
{न कत्थनात्सत्पुरुषा भवन्ति}
{मयि स्थिते धन्विनि बाणपाणौ}
{विदर्शयस्वात्मबलं दुरात्मन्} %||6-59-57||

\twolineshloka
{कर्मणा सूचयात्मानं न विकत्थितुमर्हसि}
{पौरुषेण तु यो युक्तः स तु शूर इति स्मृतः} %||6-59-58||

\twolineshloka
{सर्वायुधसमायुक्तो धन्वी त्वं रथमास्थितः}
{शरैर्वा यदि वाप्यस्त्रैर्दर्शयस्व पराक्रमम्} %||6-59-59||

\twolineshloka
{ततः शिरस्ते निशितैः पातयिष्याम्यहं शरैः}
{मारुतः कालसम्पक्वं वृन्तात्तालफलं यथा} %||6-59-60||

\twolineshloka
{अद्य ते मामका बाणास्तप्तकाञ्चनभूषणाः}
{पास्यन्ति रुधिरं गात्राद्बाणशल्यान्तरोत्थितम्} %||6-59-61||

\twolineshloka
{बालोऽयमिति विज्ञाय न मावज्ञातुमर्हसि}
{बालो वा यदि वा वृद्धो मृत्युं जानीहि संयुगे} %||6-59-62||

\twolineshloka
{लक्ष्मणस्य वचः श्रुत्वा हेतुमत्परमार्थवत्}
{अतिकायः प्रचुक्रोध बाणं चोत्तममाददे} %||6-59-63||

\twolineshloka
{ततो विद्याधरा भूता देवा दैत्या महर्षयः}
{गुह्यकाश्च महात्मानस्तद्युद्धं ददृशुस्तदा} %||6-59-64||

\twolineshloka
{ततोऽतिकायः कुपितश्चापमारोप्य सायकम्}
{लक्ष्मणस्य प्रचिक्षेप सङ्क्षिपन्निव चाम्बरम्} %||6-59-65||

\twolineshloka
{तमापतन्तं निशितं शरमाशीविषोपमम्}
{अर्धचन्द्रेण चिच्छेद लक्ष्मणः परवीरहा} %||6-59-66||

\twolineshloka
{तं निकृत्तं शरं दृष्ट्वा कृत्तभोगमिवोरगम्}
{अतिकायो भृशं क्रुद्धः पञ्चबाणान्समाददे} %||6-59-67||

\twolineshloka
{ताञ्शरान्सम्प्रचिक्षेप लक्ष्मणाय निशाचरः}
{तानप्राप्ताञ्शरैस्तीक्ष्णैश्चिच्छेद भरतानुजः} %||6-59-68||

\twolineshloka
{स तांश्छित्त्वा शरैस्तीक्ष्णैर्लक्ष्मणः परवीरहा}
{आददे निशितं बाणं ज्वलन्तमिव तेजसा} %||6-59-69||

\twolineshloka
{तमादाय धनुः श्रेष्ठे योजयामास लक्ष्मणः}
{विचकर्ष च वेगेन विससर्ज च सायकम्} %||6-59-70||

\twolineshloka
{पूर्णायतविसृष्टेन शरेणानत पर्वणा}
{ललाटे राक्षसश्रेष्ठमाजघान स वीर्यवान्} %||6-59-71||

\twolineshloka
{स ललाटे शरो मग्नस्तस्य भीमस्य रक्षसः}
{ददृशे शोणितेनाक्तः पन्नगेन्द्र इवाहवे} %||6-59-72||

\twolineshloka
{राक्षसः प्रचकम्पे च लक्ष्मणेषु प्रकम्पितः}
{रुद्रबाणहतं भीमं यथा त्रिपुरगोपुरम्} %||6-59-73||

\twolineshloka
{चिन्तयामास चाश्वस्य विमृश्य च महाबलः}
{साधु बाणनिपातेन श्वाघनीयोऽसि मे रिपुः} %||6-59-74||

\twolineshloka
{विचार्यैवं विनम्यास्यं विनम्य च भुजावुभौ}
{स रथोपस्थमास्थाय रथेन प्रचचार ह} %||6-59-75||

\twolineshloka
{एकं त्रीन्पञ्च सप्तेति सायकान्राक्षसर्षभः}
{आददे सन्दधे चापि विचकर्षोत्ससर्ज च} %||6-59-76||

\twolineshloka
{ते बाणाः कालसङ्काशा राक्षसेन्द्रधनुश्च्युताः}
{हेमपुङ्खा रविप्रख्याश्चक्रुर्दीप्तमिवाम्बरम्} %||6-59-77||

\twolineshloka
{ततस्तान्राक्षसोत्सृष्टाञ्शरौघान्रावणानुजः}
{असम्भ्रान्तः प्रचिच्छेद निशितैर्बहुभिः शरैः} %||6-59-78||

\twolineshloka
{ताञ्शरान्युधि सम्प्रेक्ष्य निकृत्तान्रावणात्मजः}
{चुकोप त्रिदशेन्द्रारिर्जग्राह निशितं शरम्} %||6-59-79||

\twolineshloka
{स सन्धाय महातेजास्तं बाणं सहसोत्सृजत्}
{ततः सौमित्रिमायान्तमाजघान स्तनान्तरे} %||6-59-80||

\twolineshloka
{अतिकायेन सौमित्रिस्ताडितो युधि वक्षसि}
{सुस्राव रुधिरं तीव्रं मदं मत्त इव द्विपः} %||6-59-81||

\twolineshloka
{स चकार तदात्मानं विशल्यं सहसा विभुः}
{जग्राह च शरं तीष्णमस्त्रेणापि समादधे} %||6-59-82||

\twolineshloka
{आग्नेयेन तदास्त्रेण योजयामास सायकम्}
{स जज्वाल तदा बाणो धनुश्चास्य महात्मनः} %||6-59-83||

\twolineshloka
{अतिकायोऽतितेजस्वी सौरमस्त्रं समाददे}
{तेन बाणं भुजङ्गाभं हेमपुङ्खमयोजयत्} %||6-59-84||

\twolineshloka
{ततस्तं ज्वलितं घोरं लक्ष्मणः शरमाहितम्}
{अतिकायाय चिक्षेप कालदण्डमिवान्तकः} %||6-59-85||

\twolineshloka
{आग्नेयेनाभिसंयुक्तं दृष्ट्वा बाणं निशाचरः}
{उत्ससर्ज तदा बाणं दीप्तं सूर्यास्त्रयोजितम्} %||6-59-86||

\twolineshloka
{तावुभावम्बरे बाणावन्योन्यमभिजघ्नतुः}
{तेजसा सम्प्रदीप्ताग्रौ क्रुद्धाविव भुजं गमौ} %||6-59-87||

\twolineshloka
{तावन्योन्यं विनिर्दह्य पेततुर्धरणीतले}
{निरर्चिषौ भस्मकृतौ न भ्राजेते शरोत्तमौ} %||6-59-88||

\twolineshloka
{ततोऽतिकायः सङ्क्रुद्धस्त्वस्त्रमैषीकमुत्सृजत्}
{तत्प्रचिच्छेद सौमित्रिरस्त्रमैन्द्रेण वीर्यवान्} %||6-59-89||

\twolineshloka
{ऐषीकं निहतं दृष्ट्वा कुमारो रावणात्मजः}
{याम्येनास्त्रेण सङ्क्रुद्धो योजयामास सायकम्} %||6-59-90||

\twolineshloka
{ततस्तदस्त्रं चिक्षेप लक्ष्मणाय निशाचरः}
{वायव्येन तदस्त्रं तु निजघान स लक्ष्मणः} %||6-59-91||

\twolineshloka
{अथैनं शरधाराभिर्धाराभिरिव तोयदः}
{अभ्यवर्षत सङ्क्रुद्धो लक्ष्मणो रावणात्मजम्} %||6-59-92||

\twolineshloka
{तेऽतिकायं समासाद्य कवचे वज्रभूषिते}
{भग्नाग्रशल्याः सहसा पेतुर्बाणा महीतले} %||6-59-93||

\twolineshloka
{तान्मोघानभिसम्प्रेक्ष्य लक्ष्मणः परवीरहा}
{अभ्यवर्षत बाणानां सहस्रेण महायशाः} %||6-59-94||

\twolineshloka
{स वर्ष्यमाणो बाणौघैरतिकायो महाबलः}
{अवध्यकवचः सङ्ख्ये राक्षसो नैव विव्यथे} %||6-59-95||

\twolineshloka
{न शशाक रुजं कर्तुं युधि तस्य नरोत्तमः}
{अथैनमभ्युपागम्य वायुर्वाक्यमुवाच ह} %||6-59-96||

\twolineshloka
{ब्रह्मदत्तवरो ह्येष अवध्य कवचावृतः}
{ब्राह्मेणास्त्रेण भिन्ध्येनमेष वध्यो हि नान्यथा} %||6-59-97||

\fourlineindentedshloka
{ततः स वायोर्वचनं निशम्य}
{सौमित्रिरिन्द्रप्रतिमानवीर्यः}
{समाददे बाणममोघवेगम्}
{तद्ब्राह्ममस्त्रं सहसा नियोज्य} %||6-59-98||

\fourlineindentedshloka
{तस्मिन्वरास्त्रे तु नियुज्यमाने}
{सौमित्रिणा बाणवरे शिताग्रे}
{दिशः सचन्द्रार्कमहाग्रहाश्च}
{नभश्च तत्रास ररास चोर्वी} %||6-59-99||

\fourlineindentedshloka
{तं ब्रह्मणोऽस्त्रेण नियुज्य चापे}
{शरं सुपुङ्खं यमदूतकल्पम्}
{सौमित्रिरिन्द्रारिसुतस्य तस्य}
{ससर्ज बाणं युधि वज्रकल्पम्} %||6-59-100||

\fourlineindentedshloka
{तं लक्ष्मणोत्सृष्टममोघवेगम्}
{समापतन्तं ज्वलनप्रकाशम्}
{सुवर्णवज्रोत्तमचित्रपुङ्खम्}
{तदातिकायः समरे ददर्श} %||6-59-101||

\fourlineindentedshloka
{तं प्रेक्षमाणः सहसातिकायो}
{जघान बाणैर्निशितैरनेकैः}
{स सायकस्तस्य सुपर्णवेगः}
{तदातिवेगेन जगाम पार्श्वम्} %||6-59-102||

\fourlineindentedshloka
{तमागतं प्रेक्ष्य तदातिकायो}
{बाणं प्रदीप्तान्तककालकल्पम्}
{जघान शक्त्यृष्टिगदाकुठारैः}
{शूलैर्हलैश्चाप्यविपन्नचेष्टः} %||6-59-103||

\fourlineindentedshloka
{तान्यायुधान्यद्भुतविग्रहाणि}
{मोघानि कृत्वा स शरोऽग्निदीप्तः}
{प्रसह्य तस्यैव किरीटजुष्टम्}
{तदातिकायस्य शिरो जहार} %||6-59-104||

\twolineshloka
{तच्छिरः सशिरस्त्राणं लक्ष्मणेषुप्रपीडितम्}
{पपात सहसा भूमौ शृङ्गं हिमवतो यथा} %||6-59-105||

\fourlineindentedshloka
{प्रहर्षयुक्ता बहवस्तु वानरा}
{प्रबुद्धपद्मप्रतिमाननास्तदा}
{अपूजयँल्लक्ष्मणमिष्टभागिनम्}
{हते रिपौ भीमबले दुरासदे} %||6-60-106||


{॥इत्यार्षे श्रीमद्रामायणे वाल्मीकीये आदिकाव्ये युद्धकाण्डे एकोनषष्टितमः सर्गः॥}

\sect{षष्टितमः सर्गः}

\fourlineindentedshloka
{ततो हतान्राक्षसपुङ्गवांस्तान्}
{देवान्तकादित्रिशिरोऽतिकायान्}
{रक्षोगणास्तत्र हतावशिष्टाः}
{ते रावणाय त्वरितं शशंसुः} %||6-60-1||

\fourlineindentedshloka
{ततो हतांस्तान्सहसा निशम्य}
{राजा मुमोहाश्रुपरिप्लुताक्षः}
{पुत्रक्षयं भ्रातृवधं च घोरम्}
{विचिन्त्य राजा विपुलं प्रदध्यौ} %||6-60-2||

\fourlineindentedshloka
{ततस्तु राजानमुदीक्ष्य दीनम्}
{शोकार्णवे सम्परिपुप्लुवानम्}
{अथर्षभो राक्षसराजसूनुः}
{अथेन्द्रजिद्वाक्यमिदं बभाषे} %||6-60-3||

\fourlineindentedshloka
{न तात मोहं प्रतिगन्तुमर्हसि}
{यत्रेन्द्रजिज्जीवति राक्षसेन्द्र}
{नेन्द्रारिबाणाभिहतो हि कश्चित्}
{प्राणान्समर्थः समरेऽभिधर्तुम्} %||6-60-4||

\fourlineindentedshloka
{पश्याद्य रामं सहलक्ष्मणेन}
{मद्बाणनिर्भिन्नविकीर्णदेहम्}
{गतायुषं भूमितले शयानम्}
{शरैः शितैराचितसर्वगात्रम्} %||6-60-5||

\fourlineindentedshloka
{इमां प्रतिज्ञां शृणु शक्रशत्रोः}
{सुनिश्चितां पौरुषदैवयुक्ताम्}
{अद्यैव रामं सहलक्ष्मणेन}
{सन्तापयिष्यामि शरैरमोघैः} %||6-60-6||

\fourlineindentedshloka
{अद्येन्द्रवैवस्वतविष्णुमित्र}
{साध्याश्विवैश्वानरचन्द्रसूर्याः}
{द्रक्ष्यन्ति मे विक्रममप्रमेयम्}
{विष्णोरिवोग्रं बलियज्ञवाटे} %||6-60-7||

\fourlineindentedshloka
{स एवमुक्त्वा त्रिदशेन्द्रशत्रुः}
{आपृच्छ्य राजानमदीनसत्त्वः}
{समारुरोहानिलतुल्यवेगम्}
{रथं खरश्रेष्ठसमाधियुक्तम्} %||6-60-8||

\twolineshloka
{समास्थाय महातेजा रथं हरिरथोपमम्}
{जगाम सहसा तत्र यत्र युद्धमरिन्दम} %||6-60-9||

\twolineshloka
{तं प्रस्थितं महात्मानमनुजग्मुर्महाबलाः}
{संहर्षमाणा बहवो धनुःप्रवरपाणयः} %||6-60-10||

\twolineshloka
{गजस्कन्धगताः केचित्केचित्परमवाजिभिः}
{प्रासमुद्गरनिस्त्रिंश परश्वधगदाधराः} %||6-60-11||

\twolineshloka
{स शङ्खनिनदैर्भीमैर्भेरीणां च महास्वनैः}
{जगाम त्रिदशेन्द्रारिः स्तूयमानो निशाचरैः} %||6-60-12||

\twolineshloka
{स शङ्खशशिवर्णेन छत्रेण रिपुसादनः}
{रराज परिपूर्णेन नभश्चन्द्रमसा यथा} %||6-60-13||

\twolineshloka
{अवीज्यत ततो वीरो हैमैर्हेमविभूषितैः}
{चारुचामरमुख्यैश्च मुख्यः सर्वधनुष्मताम्} %||6-60-14||

\twolineshloka
{ततस्त्विन्द्रजिता लङ्का सूर्यप्रतिमतेजसा}
{रराजाप्रतिवीर्येण द्यौरिवार्केण भास्वता} %||6-60-15||

\twolineshloka
{स तु दृष्ट्वा विनिर्यान्तं बलेन महता वृतम्}
{राक्षसाधिपतिः श्रीमान्रावणः पुत्रमब्रवीत्} %||6-60-16||

\twolineshloka
{त्वमप्रतिरथः पुत्र जितस्ते युधि वासवः}
{किं पुनर्मानुषं धृष्यं न वधिष्यसि राघवम्} %||6-60-17||

\twolineshloka
{तथोक्तो राक्षसेन्द्रेण प्रतिगृह्य महाशिषः}
{रथेनाश्वयुजा वीरः शीघ्रं गत्वा निकुम्भिलाम्} %||6-60-18||

\twolineshloka
{स सम्प्राप्य महातेजा युद्धभूमिमरिन्दमः}
{स्थापयामास रक्षांसि रथं प्रति समन्ततः} %||6-60-19||

\twolineshloka
{ततस्तु हुतभोक्तारं हुतभुक्सदृशप्रभः}
{जुहुवे राक्षसश्रेष्ठो मन्त्रवद्विधिवत्तदा} %||6-60-20||

\twolineshloka
{स हविर्जालसंस्कारैर्माल्यगन्धपुरस्कृतैः}
{जुहुवे पावकं तत्र राक्षसेन्द्रः प्रतापवान्} %||6-60-21||

\twolineshloka
{शस्त्राणि शरपत्राणि समिधोऽथ विभीतकः}
{लोहितानि च वासांसि स्रुवं कार्ष्णायसं तथा} %||6-60-22||

\twolineshloka
{स तत्राग्निं समास्तीर्य शरपत्रैः सतोमरैः}
{छागस्य सर्वकृष्णस्य गलं जग्राह जीवतः} %||6-60-23||

\twolineshloka
{सकृदेव समिद्धस्य विधूमस्य महार्चिषः}
{बभूवुस्तानि लिङ्गानि विजयं यान्यदर्शयन्} %||6-60-24||

\twolineshloka
{प्रदक्षिणावर्तशिखस्तप्तकाञ्चनसंनिभः}
{हविस्तत्प्रतिजग्राह पावकः स्वयमुत्थितः} %||6-60-25||

\twolineshloka
{सोऽस्त्रमाहारयामास ब्राह्ममस्त्रविदां वरः}
{धनुश्चात्मरथं चैव सर्वं तत्राभ्यमन्त्रयत्} %||6-60-26||

\twolineshloka
{तस्मिन्नाहूयमानेऽस्त्रे हूयमाने च पावके}
{सार्कग्रहेन्दु नक्षत्रं वितत्रास नभस्तलम्} %||6-60-27||

\fourlineindentedshloka
{स पावकं पावकदीप्ततेजा}
{हुत्वा महेन्द्रप्रतिमप्रभावः}
{सचापबाणासिरथाश्वसूतः}
{खेऽन्तर्दध आत्मानमचिन्त्यरूपः} %||6-60-28||

\fourlineindentedshloka
{स सैन्यमुत्सृज्य समेत्य तूर्णम्}
{महारणे वानरवाहिनीषु}
{अदृश्यमानः शरजालमुग्रम्}
{ववर्ष नीलाम्बुधरो यथाम्बु} %||6-60-29||

\fourlineindentedshloka
{ते शक्रजिद्बाणविशीर्णदेहा}
{मायाहता विस्वरमुन्नदन्तः}
{रणे निपेतुर्हरयोऽद्रिकल्पा}
{यथेन्द्रवज्राभिहता नगेन्द्राः} %||6-60-30||

\fourlineindentedshloka
{ते केवलं सन्ददृशुः शिताग्रान्}
{बाणान्रणे वानरवाहिनीषु}
{माया निगूढं च सुरेन्द्रशत्रुम्}
{न चात्र तं राक्षसमभ्यपश्यन्} %||6-60-31||

\fourlineindentedshloka
{ततः स रक्षोऽधिपतिर्महात्मा}
{सर्वा दिशो बाणगणैः शिताग्रैः}
{प्रच्छादयामास रविप्रकाशैर्-}
{विषादयामास च वानरेन्द्रान्} %||6-60-32||

\fourlineindentedshloka
{स शूलनिस्त्रिंश परश्वधानि}
{व्याविध्य दीप्तानलसंनिभानि}
{सविस्फुलिङ्गोज्ज्वलपावकानि}
{ववर्ष तीव्रं प्लवगेन्द्रसैन्ये} %||6-60-33||

\twolineshloka
{ततो ज्वलनसङ्काशैः शितैर्वानरयूथपाः}
{ताडिताः शक्रजिद्बाणैः प्रफुल्ला इव किंशुकाः} %||6-60-34||

\twolineshloka
{अन्योन्यमभिसर्पन्तो निनदन्तश्च विस्वरम्}
{राक्षसेन्द्रास्त्रनिर्भिन्ना निपेतुर्वानरर्षभाः} %||6-60-35||

\twolineshloka
{उदीक्षमाणा गगनं केचिन्नेत्रेषु ताडिताः}
{शरैर्विविशुरन्योन्यं पेतुश्च जगतीतले} %||6-60-36||

\twolineshloka
{हनूमन्तं च सुग्रीवमङ्गदं गन्धमादनम्}
{जाम्बवन्तं सुषेणं च वेगदर्शिनमेव च} %||6-60-37||

\twolineshloka
{मैन्दं च द्विविदं नीलं गवाक्षं गजगोमुखौ}
{केसरिं हरिलोमानं विद्युद्दंष्ट्रं च वानरम्} %||6-60-38||

\twolineshloka
{सूर्याननं ज्योतिमुखं तथा दधिमुखं हरिम्}
{पावकाक्षं नलं चैव कुमुदं चैव वानरम्} %||6-60-39||

\twolineshloka
{प्रासैः शूलैः शितैर्बाणैरिन्द्रजिन्मन्त्रसंहितैः}
{विव्याध हरिशार्दूलान्सर्वांस्तान्राक्षसोत्तमः} %||6-60-40||

\fourlineindentedshloka
{स वै गदाभिर्हरियूथमुख्यान्}
{निर्भिद्य बाणैस्तपनीयपुङ्खैः}
{ववर्ष रामं शरवृष्टिजालैः}
{सलक्ष्मणं भास्कररश्मिकल्पैः} %||6-60-41||

\fourlineindentedshloka
{स बाणवर्षैरभिवर्ष्यमाणो}
{धारानिपातानिव तान्विचिन्त्य}
{समीक्षमाणः परमाद्भुतश्री}
{रामस्तदा लक्ष्मणमित्युवाच} %||6-60-42||

\fourlineindentedshloka
{असौ पुनर्लक्ष्मण राक्षसेन्द्रो}
{ब्रह्मास्त्रमाश्रित्य सुरेन्द्रशत्रुः}
{निपातयित्वा हरिसैन्यमुग्रम्}
{अस्माञ्शरैरर्दयति प्रसक्तम्} %||6-60-43||

\fourlineindentedshloka
{स्वयम्भुवा दत्तवरो महात्मा}
{खमास्थितोऽन्तर्हितभीमकायः}
{कथं नु शक्यो युधि नष्टदेहो}
{निहन्तुमद्येन्द्रजिदुद्यतास्त्रः} %||6-60-44||

\fourlineindentedshloka
{मन्ये स्वयम्भूर्भगवानचिन्त्यो}
{यस्यैतदस्त्रं प्रभवश्च योऽस्य}
{बाणावपातांस्त्वमिहाद्य धीमन्}
{मया सहाव्यग्रमनाः सहस्व} %||6-60-45||

\fourlineindentedshloka
{प्रच्छादयत्येष हि राक्षसेन्द्रः}
{सर्वा दिशः सायकवृष्टिजालैः}
{एतच्च सर्वं पतिताग्र्यवीरम्}
{न भ्राजते वानरराजसैन्यम्} %||6-60-46||

\fourlineindentedshloka
{आवां तु दृष्ट्वा पतितौ विसंज्ञौ}
{निवृत्तयुद्धौ हतरोषहर्षौ}
{ध्रुवं प्रवेक्ष्यत्यमरारिवासम्}
{असौ समादाय रणाग्रलक्ष्मीम्} %||6-60-47||

\fourlineindentedshloka
{ततस्तु ताविन्द्रजिदस्त्रजालैर्-}
{बभूवतुस्तत्र तदा विशस्तौ}
{स चापि तौ तत्र विषादयित्वा}
{ननाद हर्षाद्युधि राक्षसेन्द्रः} %||6-60-48||

\fourlineindentedshloka
{स तत्तदा वानरराजसैन्यम्}
{रामं च सङ्ख्ये सहलक्ष्मणेन}
{विषादयित्वा सहसा विवेश}
{पुरीं दशग्रीवभुजाभिगुप्ताम्} %||6-61-49||


{॥इत्यार्षे श्रीमद्रामायणे वाल्मीकीये आदिकाव्ये युद्धकाण्डे षष्टितमः सर्गः॥}

\sect{एकषष्टितमः सर्गः}

\fourlineindentedshloka
{तयोस्तदा सादितयो रणाग्रे}
{मुमोह सैन्यं हरियूथपानाम्}
{सुग्रीवनीलाङ्गदजाम्बवन्तो}
{न चापि किञ्चित्प्रतिपेदिरे ते} %||6-61-1||

\fourlineindentedshloka
{ततो विषण्णं समवेक्ष्य सैन्यम्}
{विभीषणो बुद्धिमतां वरिष्ठः}
{उवाच शाखामृगराजवीरान्}
{आश्वासयन्नप्रतिमैर्वचोभिः} %||6-61-2||

\fourlineindentedshloka
{मा भैष्ट नास्त्यत्र विषादकालो}
{यदार्यपुत्राववशौ विषण्णौ}
{स्वयम्भुवो वाक्यमथोद्वहन्तौ}
{यत्सादिताविन्द्रजिदस्त्रजालैः} %||6-61-3||

\fourlineindentedshloka
{तस्मै तु दत्तं परमास्त्रमेतत्}
{स्वयम्भुवा ब्राह्मममोघवेगम्}
{तन्मानयन्तौ यदि राजपुत्रौ}
{निपातितौ कोऽत्र विषादकालः} %||6-61-4||

\twolineshloka
{ब्राह्ममस्त्रं तदा धीमान्मानयित्वा तु मारुतिः}
{विभीषणवचः श्रुत्वा हनूमांस्तमथाब्रवीत्} %||6-61-5||

\twolineshloka
{एतस्मिन्निहते सैन्ये वानराणां तरस्विनाम्}
{यो यो धारयते प्राणांस्तं तमाश्वासयावहे} %||6-61-6||

\twolineshloka
{तावुभौ युगपद्वीरौ हनूमद्राक्षसोत्तमौ}
{उल्काहस्तौ तदा रात्रौ रणशीर्षे विचेरतुः} %||6-61-7||

\twolineshloka
{छिन्नलाङ्गूलहस्तोरुपादाङ्गुलि शिरो धरैः}
{स्रवद्भिः क्षतजं गात्रैः प्रस्रवद्भिः समन्ततः} %||6-61-8||

\twolineshloka
{पतितैः पर्वताकारैर्वानरैरभिसङ्कुलाम्}
{शस्त्रैश्च पतितैर्दीप्तैर्ददृशाते वसुन्धराम्} %||6-61-9||

\twolineshloka
{सुग्रीवमङ्गदं नीलं शरभं गन्धमादनम्}
{जाम्बवन्तं सुषेणं च वेगदर्शनमाहुकम्} %||6-61-10||

\twolineshloka
{मैन्दं नलं ज्योतिमुखं द्विविदं पनसं तथा}
{विभीषणो हनूमांश्च ददृशाते हतान्रणे} %||6-61-11||

\twolineshloka
{सप्तषष्टिर्हताः कोट्यो वानराणां तरस्विनाम्}
{अह्नः पञ्चमशेषेण वल्लभेन स्वयम्भुवः} %||6-61-12||

\twolineshloka
{सागरौघनिभं भीमं दृष्ट्वा बाणार्दितं बलम्}
{मार्गते जाम्बवन्तं स्म हनूमान्सविभीषणः} %||6-61-13||

\twolineshloka
{स्वभावजरया युक्तं वृद्धं शरशतैश्चितम्}
{प्रजापतिसुतं वीरं शाम्यन्तमिव पावकम्} %||6-61-14||

\twolineshloka
{दृष्ट्वा तमुपसङ्गम्य पौलस्त्यो वाक्यमब्रवीत्}
{कच्चिदार्यशरैस्तीर्ष्णैर्न प्राणा ध्वंसितास्तव} %||6-61-15||

\twolineshloka
{विभीषणवचः श्रुत्वा जाम्बवानृक्षपुङ्गवः}
{कृच्छ्रादभ्युद्गिरन्वाक्यमिदं वचनमब्रवीत्} %||6-61-16||

\twolineshloka
{नैरृतेन्द्रमहावीर्यस्वरेण त्वाभिलक्षये}
{पीड्यमानः शितैर्बाणैर्न त्वां पश्यामि चक्षुषा} %||6-61-17||

\twolineshloka
{अञ्जना सुप्रजा येन मातरिश्वा च नैरृत}
{हनूमान्वानरश्रेष्ठः प्राणान्धारयते क्वचित्} %||6-61-18||

\twolineshloka
{श्रुत्वा जाम्बवतो वाक्यमुवाचेदं विभीषणः}
{आर्यपुत्रावतिक्रम्य कस्मात्पृच्छसि मारुतिम्} %||6-61-19||

\twolineshloka
{नैव राजनि सुग्रीवे नाङ्गदे नापि राघवे}
{आर्य सन्दर्शितः स्नेहो यथा वायुसुते परः} %||6-61-20||

\twolineshloka
{विभीषणवचः श्रुत्वा जाम्बवान्वाक्यमब्रवीत्}
{शृणु नैरृतशार्दूल यस्मात्पृच्छामि मारुतिम्} %||6-61-21||

\twolineshloka
{तस्मिञ्जीवति वीरे तु हतमप्यहतं बलम्}
{हनूमत्युज्झितप्राणे जीवन्तोऽपि वयं हताः} %||6-61-22||

\twolineshloka
{ध्रियते मारुतिस्तात मारुतप्रतिमो यदि}
{वैश्वानरसमो वीर्ये जीविताशा ततो भवेत्} %||6-61-23||

\twolineshloka
{ततो वृद्धमुपागम्य नियमेनाभ्यवादयत्}
{गृह्य जाम्बवतः पादौ हनूमान्मारुतात्मजः} %||6-61-24||

\twolineshloka
{श्रुत्वा हनुमतो वाक्यं तथापि व्यथितेन्द्रियः}
{पुनर्जातमिवात्मानं स मेने ऋक्षपुङ्गवः} %||6-61-25||

\twolineshloka
{ततोऽब्रवीन्महातेजा हनूमन्तं स जाम्बवान्}
{आगच्छ हरिशार्दूलवानरांस्त्रातुमर्हसि} %||6-61-26||

\twolineshloka
{नान्यो विक्रमपर्याप्तस्त्वमेषां परमः सखा}
{त्वत्पराक्रमकालोऽयं नान्यं पश्यामि कञ्चन} %||6-61-27||

\twolineshloka
{ऋक्षवानरवीराणामनीकानि प्रहर्षय}
{विशल्यौ कुरु चाप्येतौ सादितौ रामलक्ष्मणौ} %||6-61-28||

\twolineshloka
{गत्वा परममध्वानमुपर्युपरि सागरम्}
{हिमवन्तं नगश्रेष्ठं हनूमन्गन्तुमर्हसि} %||6-61-29||

\twolineshloka
{ततः काञ्चनमत्युग्रमृषभं पर्वतोत्तमम्}
{कैलासशिखरं चापि द्रक्ष्यस्यरिनिषूदन} %||6-61-30||

\twolineshloka
{तयोः शिखरयोर्मध्ये प्रदीप्तमतुलप्रभम्}
{सर्वौषधियुतं वीर द्रक्ष्यस्यौषधिपर्वतम्} %||6-61-31||

\twolineshloka
{तस्य वानरशार्दूलचतस्रो मूर्ध्नि सम्भवाः}
{द्रक्ष्यस्योषधयो दीप्ता दीपयन्त्यो दिशो दश} %||6-61-32||

\twolineshloka
{मृतसंजीवनीं चैव विशल्यकरणीमपि}
{सौवर्णकरणीं चैव सन्धानीं च महौषधीम्} %||6-61-33||

\twolineshloka
{ताः सर्वा हनुमन्गृह्य क्षिप्रमागन्तुमर्हसि}
{आश्वासय हरीन्प्राणैर्योज्य गन्धवहात्मजः} %||6-61-34||

\twolineshloka
{श्रुत्वा जाम्बवतो वाक्यं हनूमान्हरिपुङ्गवः}
{आपूर्यत बलोद्धर्षैस्तोयवेगैरिवार्णवः} %||6-61-35||

\twolineshloka
{स पर्वततटाग्रस्थः पीडयन्पर्वतोत्तरम्}
{हनूमान्दृश्यते वीरो द्वितीय इव पर्वतः} %||6-61-36||

\twolineshloka
{हरिपादविनिर्भिन्नो निषसाद स पर्वतः}
{न शशाक तदात्मानं सोढुं भृशनिपीडितः} %||6-61-37||

\twolineshloka
{तस्य पेतुर्नगा भूमौ हरिवेगाच्च जज्वलुः}
{शृङ्गाणि च व्यकीर्यन्त पीडितस्य हनूमता} %||6-61-38||

\twolineshloka
{तस्मिन्सम्पीड्यमाने तु भग्नद्रुमशिलातले}
{न शेकुर्वानराः स्थातुं घूर्णमाने नगोत्तमे} %||6-61-39||

\twolineshloka
{स घूर्णितमहाद्वारा प्रभग्नगृहगोपुरा}
{लङ्का त्रासाकुला रात्रौ प्रनृत्तेवाभवत्तदा} %||6-61-40||

\twolineshloka
{पृथिवीधरसङ्काशो निपीड्य धरणीधरम्}
{पृथिवीं क्षोभयामास सार्णवां मारुतात्मजः} %||6-61-41||

\twolineshloka
{पद्भ्यां तु शैलमापीड्य वडवामुखवन्मुखम्}
{विवृत्योग्रं ननादोच्चैस्त्रासयन्निव राक्षसान्} %||6-61-42||

\twolineshloka
{तस्य नानद्यमानस्य श्रुत्वा निनदमद्भुतम्}
{लङ्कास्था राक्षसाः सर्वे न शेकुः स्पन्दितुं भयात्} %||6-61-43||

\twolineshloka
{नमस्कृत्वाथ रामाय मारुतिर्भीमविक्रमः}
{राघवार्थे परं कर्म समैहत परन्तपः} %||6-61-44||

\fourlineindentedshloka
{स पुच्छमुद्यम्य भुजङ्गकल्पम्}
{विनम्य पृष्ठं श्रवणे निकुञ्च्य}
{विवृत्य वक्त्रं वडवामुखाभम्}
{आपुप्लुवे व्योम्नि स चण्डवेगः} %||6-61-45||

\fourlineindentedshloka
{स वृक्षषण्डांस्तरसा जहार}
{शैलाञ्शिलाः प्राकृतवानरांश्च}
{बाहूरुवेगोद्धतसम्प्रणुन्नाः}
{ते क्षीणवेगाः सलिले निपेतुः} %||6-61-46||

\fourlineindentedshloka
{स तौ प्रसार्योरगभोगकल्पौ}
{भुजौ भुजङ्गारिनिकाशवीर्यः}
{जगाम मेरुं नगराजमग्र्यम्}
{दिशः प्रकर्षन्निव वायुसूनुः} %||6-61-47||

\fourlineindentedshloka
{स सागरं घूर्णितवीचिमालम्}
{तदा भृशं भ्रामितसर्वसत्त्वम्}
{समीक्षमाणः सहसा जगाम}
{चक्रं यथा विष्णुकराग्रमुक्तम्} %||6-61-48||

\fourlineindentedshloka
{स पर्वतान्वृक्षगणान्सरांसि}
{नदीस्तटाकानि पुरोत्तमानि}
{स्फीताञ्जनांस्तानपि सम्प्रपश्यन्}
{जगाम वेगात्पितृतुल्यवेगः} %||6-61-49||

\twolineshloka
{आदित्यपथमाश्रित्य जगाम स गतश्रमः}
{स ददर्श हरिश्रेष्ठो हिमवन्तं नगोत्तमम्} %||6-61-50||

\twolineshloka
{नानाप्रस्रवणोपेतं बहुकन्दरनिर्झरम्}
{श्वेताभ्रचयसङ्काशैः शिखरैश्चारुदर्शनैः} %||6-61-51||

\fourlineindentedshloka
{स तं समासाद्य महानगेन्द्रम्}
{अतिप्रवृद्धोत्तमघोरशृङ्गम्}
{ददर्श पुण्यानि महाश्रमाणि}
{सुरर्षिसङ्घोत्तमसेवितानि} %||6-61-52||

\fourlineindentedshloka
{स ब्रह्मकोशं रजतालयं च}
{शक्रालयं रुद्रशरप्रमोक्षम्}
{हयाननं ब्रह्मशिरश्च दीप्तम्}
{ददर्श वैवस्वत किङ्करांश्च} %||6-61-53||

\fourlineindentedshloka
{वज्रालयं वैश्वरणालयं च}
{सूर्यप्रभं सूर्यनिबन्धनं च}
{ब्रह्मासनं शङ्करकार्मुकं च}
{ददर्श नाभिं च वसुन्धरायाः} %||6-61-54||

\fourlineindentedshloka
{कैलासमग्र्यं हिमवच्छिलां च}
{तथर्षभं काञ्चनशैलमग्र्यम्}
{स दीप्तसर्वौषधिसम्प्रदीप्तम्}
{ददर्श सर्वौषधिपर्वतेन्द्रम्} %||6-61-55||

\fourlineindentedshloka
{स तं समीक्ष्यानलरश्मिदीप्तम्}
{विसिष्मिये वासवदूतसूनुः}
{आप्लुत्य तं चौषधिपर्वतेन्द्रम्}
{तत्रौषधीनां विचयं चकार} %||6-61-56||

\twolineshloka
{स योजनसहस्राणि समतीत्य महाकपिः}
{दिव्यौषधिधरं शैलं व्यचरन्मारुतात्मजः} %||6-61-57||

\twolineshloka
{महौषध्यस्तु ताः सर्वास्तस्मिन्पर्वतसत्तमे}
{विज्ञायार्थिनमायान्तं ततो जग्मुरदर्शनम्} %||6-61-58||

\fourlineindentedshloka
{स ता महात्मा हनुमानपश्यम्ः}
{चुकोप कोपाच्च भृशं ननाद}
{अमृष्यमाणोऽग्निनिकाशचक्षुर्-}
{महीधरेन्द्रं तमुवाच वाक्यम्} %||6-61-59||

\fourlineindentedshloka
{किमेतदेवं सुविनिश्चितं ते}
{यद्राघवे नासि कृतानुकम्पः}
{पश्याद्य मद्बाहुबलाभिभूतो}
{विकीर्णमात्मानमथो नगेन्द्र} %||6-61-60||

\fourlineindentedshloka
{स तस्य शृङ्गं सनगं सनागम्}
{सकाञ्चनं धातुसहस्रजुष्टम्}
{विकीर्णकूटं चलिताग्रसानुम्}
{प्रगृह्य वेगात्सहसोन्ममाथ} %||6-61-61||

\fourlineindentedshloka
{स तं समुत्पाट्य खमुत्पपात}
{वित्रास्य लोकान्ससुरान्सुरेन्द्रान्}
{संस्तूयमानः खचरैरनेकैर्-}
{जगाम वेगाद्गरुडोग्रवीर्यः} %||6-61-62||

\fourlineindentedshloka
{स भास्कराध्वानमनुप्रपन्नः}
{तद्भास्कराभं शिखरं प्रगृह्य}
{बभौ तदा भास्करसंनिकाशो}
{रवेः समीपे प्रतिभास्कराभः} %||6-61-63||

\fourlineindentedshloka
{स तेन शैलेन भृशं रराज}
{शैलोपमो गन्धवहात्मजस्तु}
{सहस्रधारेण सपावकेन}
{चक्रेण खे विष्णुरिवोद्धृतेन} %||6-61-64||

\fourlineindentedshloka
{तं वानराः प्रेक्ष्य तदा विनेदुः}
{स तानपि प्रेक्ष्य मुदा ननाद}
{तेषां समुद्घुष्टरवं निशम्य}
{लङ्कालया भीमतरं विनेदुः} %||6-61-65||

\fourlineindentedshloka
{ततो महात्मा निपपात तस्मिन्}
{शैलोत्तमे वानरसैन्यमध्ये}
{हर्युत्तमेभ्यः शिरसाभिवाद्य}
{विभीषणं तत्र च सस्वजे सः} %||6-61-66||

\fourlineindentedshloka
{तावप्युभौ मानुषराजपुत्रौ}
{तं गन्धमाघ्राय महौषधीनाम्}
{बभूवतुस्तत्र तदा विशल्यौ}
{उत्तस्थुरन्ये च हरिप्रवीराः} %||6-61-67||

\fourlineindentedshloka
{ततो हरिर्गन्धवहात्मजस्तु}
{तमोषधीशैलमुदग्रवीर्यः}
{निनाय वेगाद्धिमवन्तमेव}
{पुनश्च रामेण समाजगाम} %||6-62-68||


{॥इत्यार्षे श्रीमद्रामायणे वाल्मीकीये आदिकाव्ये युद्धकाण्डे एकषष्टितमः सर्गः॥}

\sect{द्विषष्टितमः सर्गः}

\twolineshloka
{ततोऽब्रवीन्महातेजाः सुग्रीवो वानराधिपः}
{अर्थ्यं विजापयंश्चापि हनूमन्तं महाबलम्} %||6-62-1||

\twolineshloka
{यतो हतः कुम्भकर्णः कुमाराश्च निषूदिताः}
{नेदानीमुपनिर्हारं रावणो दातुमर्हति} %||6-62-2||

\twolineshloka
{ये ये महाबलाः सन्ति लघवश्च प्लवङ्गमाः}
{लङ्कामभ्युत्पतन्त्वाशु गृह्योल्काः प्लवगर्षभाः} %||6-62-3||

\twolineshloka
{ततोऽस्तं गत आदित्ये रौद्रे तस्मिन्निशामुखे}
{लङ्कामभिमुखाः सोल्का जग्मुस्ते प्लवगर्षभाः} %||6-62-4||

\twolineshloka
{उल्काहस्तैर्हरिगणैः सर्वतः समभिद्रुताः}
{आरक्षस्था विरूपाक्षाः सहसा विप्रदुद्रुवुः} %||6-62-5||

\twolineshloka
{गोपुराट्ट प्रतोलीषु चर्यासु विविधासु च}
{प्रासादेषु च संहृष्टाः ससृजुस्ते हुताशनम्} %||6-62-6||

\twolineshloka
{तेषां गृहसहस्राणि ददाह हुतभुक्तदा}
{आवासान्राक्षसानां च सर्वेषां गृहमेधिनाम्} %||6-62-7||

\twolineshloka
{हेमचित्रतनुत्राणां स्रग्दामाम्बरधारिणाम्}
{सीधुपानचलाक्षाणां मदविह्वलगामिनाम्} %||6-62-8||

\twolineshloka
{कान्तालम्बितवस्त्राणां शत्रुसंजातमन्युनाम्}
{गदाशूलासि हस्तानां खादतां पिबतामपि} %||6-62-9||

\twolineshloka
{शयनेषु महार्हेषु प्रसुप्तानां प्रियैः सह}
{त्रस्तानां गच्छतां तूर्णं पुत्रानादाय सर्वतः} %||6-62-10||

\twolineshloka
{तेषां गृहसहस्राणि तदा लङ्कानिवासिनाम्}
{अदहत्पावकस्तत्र जज्वाल च पुनः पुनः} %||6-62-11||

\twolineshloka
{सारवन्ति महार्हाणि गम्भीरगुणवन्ति च}
{हेमचन्द्रार्धचन्द्राणि चन्द्रशालोन्नतानि च} %||6-62-12||

\twolineshloka
{रत्नचित्रगवाक्षाणि साधिष्ठानानि सर्वशः}
{मणिविद्रुमचित्राणि स्पृशन्तीव च भास्करम्} %||6-62-13||

\twolineshloka
{क्रौञ्चबर्हिणवीणानां भूषणानां च निस्वनैः}
{नादितान्यचलाभानि वेश्मान्यग्निर्ददाह सः} %||6-62-14||

\twolineshloka
{ज्वलनेन परीतानि तोरणानि चकाशिरे}
{विद्युद्भिरिव नद्धानि मेघजालानि घर्मगे} %||6-62-15||

\twolineshloka
{विमानेषु प्रसुप्ताश्च दह्यमाना वराङ्गनाः}
{त्यक्ताभरणसंयोगा हाहेत्युच्चैर्विचुक्रुशः} %||6-62-16||

\twolineshloka
{तत्र चाग्निपरीतानि निपेतुर्भवनान्यपि}
{वज्रिवज्रहतानीव शिखराणि महागिरेः} %||6-62-17||

\twolineshloka
{तानि निर्दह्यमानानि दूरतः प्रचकाशिरे}
{हिमवच्छिखराणीव दीप्तौषधिवनानि च} %||6-62-18||

\twolineshloka
{हर्म्याग्रैर्दह्यमानैश्च ज्वालाप्रज्वलितैरपि}
{रात्रौ सा दृश्यते लङ्का पुष्पितैरिव किंशुकैः} %||6-62-19||

\twolineshloka
{हस्त्यध्यक्षैर्गजैर्मुक्तैर्मुक्तैश्च तुरगैरपि}
{बभूव लङ्का लोकान्ते भ्रान्तग्राह इवार्णवः} %||6-62-20||

\twolineshloka
{अश्वं मुक्तं गजो दृष्ट्वा कच्चिद्भीतोऽपसर्पति}
{भीतो भीतं गजं दृष्ट्वा क्वचिदश्वो निवर्तते} %||6-62-21||

\twolineshloka
{सा बभूव मुहूर्तेन हरिभिर्दीपिता पुरी}
{लोकस्यास्य क्षये घोरे प्रदीप्तेव वसुन्धरा} %||6-62-22||

\twolineshloka
{नारी जनस्य धूमेन व्याप्तस्योच्चैर्विनेदुषः}
{स्वनो ज्वलनतप्तस्य शुश्रुवे दशयोजनम्} %||6-62-23||

\twolineshloka
{प्रदग्धकायानपरान्राक्षसान्निर्गतान्बहिः}
{सहसाभ्युत्पतन्ति स्म हरयोऽथ युयुत्सवः} %||6-62-24||

\twolineshloka
{उद्घुष्टं वानराणां च राक्षसानां च निस्वनः}
{दिशो दश समुद्रं च पृथिवीं चान्वनादयत्} %||6-62-25||

\twolineshloka
{विशल्यौ तु महात्मानौ तावुभौ रामलक्ष्मणौ}
{असम्भ्रान्तौ जगृहतुस्तावुभौ धनुषी वरे} %||6-62-26||

\twolineshloka
{ततो विस्फारयाणस्य रामस्य धनुरुत्तमम्}
{बभूव तुमुलः शब्दो राक्षसानां भयावहः} %||6-62-27||

\twolineshloka
{अशोभत तदा रामो धनुर्विस्फारयन्महत्}
{भगवानिव सङ्क्रुद्धो भवो वेदमयं धनुः} %||6-62-28||

\twolineshloka
{वानरोद्घुष्टघोषश्च राक्षसानां च निस्वनः}
{ज्याशब्दश्चापि रामस्य त्रयं व्याप दिशो दश} %||6-62-29||

\twolineshloka
{तस्य कार्मुकमुक्तैश्च शरैस्तत्पुरगोपुरम्}
{कैलासशृङ्गप्रतिमं विकीर्णमपतद्भुवि} %||6-62-30||

\twolineshloka
{ततो रामशरान्दृष्ट्वा विमानेषु गृहेषु च}
{संनाहो राक्षसेन्द्राणां तुमुलः समपद्यत} %||6-62-31||

\twolineshloka
{तेषां संनह्यमानानां सिंहनादं च कुर्वताम्}
{शर्वरी राक्षसेन्द्राणां रौद्रीव समपद्यत} %||6-62-32||

\twolineshloka
{आदिष्टा वानरेन्द्रास्ते सुग्रीवेण महात्मना}
{आसन्ना द्वारमासाद्य युध्यध्वं प्लवगर्षभाः} %||6-62-33||

\twolineshloka
{यश्च वो वितथं कुर्यात्तत्र तत्र व्यवस्थितः}
{स हन्तव्योऽभिसम्प्लुत्य राजशासनदूषकः} %||6-62-34||

\twolineshloka
{तेषु वानरमुख्येषु दीप्तोल्कोज्ज्वलपाणिषु}
{स्थितेषु द्वारमासाद्य रावणं मन्युराविशत्} %||6-62-35||

\twolineshloka
{तस्य जृम्भितविक्षेपाद्व्यामिश्रा वै दिशो दश}
{रूपवानिव रुद्रस्य मन्युर्गात्रेष्वदृश्यत} %||6-62-36||

\twolineshloka
{स निकुम्भं च कुम्भं च कुम्भकर्णात्मजावुभौ}
{प्रेषयामास सङ्क्रुद्धो राक्षसैर्बहुभिः सह} %||6-62-37||

\twolineshloka
{शशास चैव तान्सर्वान्राक्षसान्राक्षसेश्वरः}
{राक्षसा गच्छतात्रैव सिंहनादं च नादयन्} %||6-62-38||

\twolineshloka
{ततस्तु चोदितास्तेन राक्षसा ज्वलितायुधाः}
{लङ्काया निर्ययुर्वीराः प्रणदन्तः पुनः पुनः} %||6-62-39||

\twolineshloka
{भीमाश्वरथमातङ्गं नानापत्ति समाकुलम्}
{दीप्तशूलगदाखड्गप्रासतोमरकार्मुकम्} %||6-62-40||

\twolineshloka
{तद्राक्षसबलं घोरं भीमविक्रमपौरुषम्}
{ददृशे ज्वलितप्रासं किङ्किणीशतनादितम्} %||6-62-41||

\twolineshloka
{हेमजालाचितभुजं व्यावेष्टितपरश्वधम्}
{व्याघूर्णितमहाशस्त्रं बाणसंसक्तकार्मुकम्} %||6-62-42||

\twolineshloka
{गन्धमाल्यमधूत्सेकसम्मोदित महानिलम्}
{घोरं शूरजनाकीर्णं महाम्बुधरनिस्वनम्} %||6-62-43||

\twolineshloka
{तं दृष्ट्वा बलमायान्तं राक्षसानां सुदारुणम्}
{सञ्चचाल प्लवङ्गानां बलमुच्चैर्ननाद च} %||6-62-44||

\twolineshloka
{जवेनाप्लुत्य च पुनस्तद्राक्षसबलं महत्}
{अभ्ययात्प्रत्यरिबलं पतङ्ग इव पावकम्} %||6-62-45||

\twolineshloka
{तेषां भुजपरामर्शव्यामृष्टपरिघाशनि}
{राक्षसानां बलं श्रेष्ठं भूयस्तरमशोभत} %||6-62-46||

\twolineshloka
{तथैवाप्यपरे तेषां कपीनामसिभिः शितैः}
{प्रवीरानभितो जघ्नुर्घोररूपा निशाचराः} %||6-62-47||

\twolineshloka
{घ्नन्तमन्यं जघानान्यः पातयन्तमपातयत्}
{गर्हमाणं जगर्हान्ये दशन्तमपरेऽदशत्} %||6-62-48||

\twolineshloka
{देहीत्यन्ये ददात्यन्यो ददामीत्यपरः पुनः}
{किं क्लेशयसि तिष्ठेति तत्रान्योन्यं बभाषिरे} %||6-62-49||

\twolineshloka
{समुद्यतमहाप्रासं मुष्टिशूलासिसङ्कुलम्}
{प्रावर्तत महारौद्रं युद्धं वानररक्षसाम्} %||6-62-50||

\twolineshloka
{वानरान्दश सप्तेति राक्षसा अभ्यपातयन्}
{राक्षसान्दशसप्तेति वानरा जघ्नुराहवे} %||6-62-51||

\twolineshloka
{विस्रस्तकेशरसनं विमुक्तकवचध्वजम्}
{बलं राक्षसमालम्ब्य वानराः पर्यवारयन्} %||6-63-52||


{॥इत्यार्षे श्रीमद्रामायणे वाल्मीकीये आदिकाव्ये युद्धकाण्डे द्विषष्टितमः सर्गः॥}

\sect{त्रिषष्टितमः सर्गः}

\twolineshloka
{प्रवृत्ते सङ्कुले तस्मिन्घोरे वीरजनक्षये}
{अङ्गदः कम्पनं वीरमाससाद रणोत्सुकः} %||6-63-1||

\twolineshloka
{आहूय सोऽङ्गदं कोपात्ताडयामास वेगितः}
{गदया कम्पनः पूर्वं स चचाल भृशाहतः} %||6-63-2||

\twolineshloka
{स संज्ञां प्राप्य तेजस्वी चिक्षेप शिखरं गिरेः}
{अर्दितश्च प्रहारेण कम्पनः पतितो भुवि} %||6-63-3||

\threelineshloka
{हतप्रवीरा व्यथिता राक्षसेन्द्रचमूस्तदा}
{जगामाभिमुखी सा तु कुम्भकर्णसुतो यतः}
{आपतन्तीं च वेगेन कुम्भस्तां सान्त्वयच्चमूम्} %||6-63-4||

\twolineshloka
{स धनुर्धन्विनां श्रेष्ठः प्रगृह्य सुसमाहितः}
{मुमोचाशीविषप्रख्याञ्शरान्देहविदारणान्} %||6-63-5||

\twolineshloka
{तस्य तच्छुशुभे भूयः सशरं धनुरुत्तमम्}
{विद्युदैरावतार्चिष्मद्द्वितीयेन्द्रधनुर्यथा} %||6-63-6||

\twolineshloka
{आकर्णकृष्टमुक्तेन जघान द्विविदं तदा}
{तेन हाटकपुङ्खेन पत्रिणा पत्रवाससा} %||6-63-7||

\twolineshloka
{सहसाभिहतस्तेन विप्रमुक्तपदः स्फुरन्}
{निपपाताद्रिकूटाभो विह्वलः प्लवगोत्तमः} %||6-63-8||

\twolineshloka
{मैन्दस्तु भ्रातरं दृष्ट्वा भग्नं तत्र महाहवे}
{अभिदुद्राव वेगेन प्रगृह्य महतीं शिलाम्} %||6-63-9||

\twolineshloka
{तां शिलां तु प्रचिक्षेप राक्षसाय महाबलः}
{बिभेद तां शिलां कुम्भः प्रसन्नैः पञ्चभिः शरैः} %||6-63-10||

\twolineshloka
{सन्धाय चान्यं सुमुखं शरमाशीविषोपमम्}
{आजघान महातेजा वक्षसि द्विविदाग्रजम्} %||6-63-11||

\twolineshloka
{स तु तेन प्रहारेण मैन्दो वानरयूथपः}
{मर्मण्यभिहतस्तेन पपात भुवि मूर्छितः} %||6-63-12||

\twolineshloka
{अङ्गदो मातुलौ दृष्ट्वा पतितौ तौ महाबलौ}
{अभिदुद्राव वेगेन कुम्भमुद्यतकार्मुकम्} %||6-63-13||

\twolineshloka
{तमापतन्तं विव्याध कुम्भः पञ्चभिरायसैः}
{त्रिभिश्चान्यैः शितैर्बाणैर्मातङ्गमिव तोमरैः} %||6-63-14||

\twolineshloka
{सोऽङ्गदं विविधैर्बाणैः कुम्भो विव्याध वीर्यवान्}
{अकुण्ठधारैर्निशितैस्तीक्ष्णैः कनकभूषणैः} %||6-63-15||

\twolineshloka
{अङ्गदः प्रतिविद्धाङ्गो वालिपुत्रो न कम्पते}
{शिलापादपवर्षाणि तस्य मूर्ध्नि ववर्ष ह} %||6-63-16||

\twolineshloka
{स प्रचिच्छेद तान्सर्वान्बिभेद च पुनः शिलाः}
{कुम्भकर्णात्मजः श्रीमान्वालिपुत्रसमीरितान्} %||6-63-17||

\twolineshloka
{आपतन्तं च सम्प्रेक्ष्य कुम्भो वानरयूथपम्}
{भ्रुवोर्विव्याध बाणाभ्यामुल्काभ्यामिव कुञ्जरम्} %||6-63-18||

\twolineshloka
{अङ्गदः पाणिना नेत्रे पिधाय रुधिरोक्षिते}
{सालमासन्नमेकेन परिजग्राह पाणिना} %||6-63-19||

\twolineshloka
{तमिन्द्रकेतुप्रतिमं वृक्षं मन्दरसंनिभम्}
{समुत्सृजन्तं वेगेन पश्यतां सर्वरक्षसाम्} %||6-63-20||

\twolineshloka
{स चिच्छेद शितैर्बाणैः सप्तभिः कायभेदनैः}
{अङ्गदो विव्यथेऽभीक्ष्णं ससाद च मुमोह च} %||6-63-21||

\twolineshloka
{अङ्गदं व्यथितं दृष्ट्वा सीदन्तमिव सागरे}
{दुरासदं हरिश्रेष्ठा राघवाय न्यवेदयन्} %||6-63-22||

\twolineshloka
{रामस्तु व्यथितं श्रुत्वा वालिपुत्रं महाहवे}
{व्यादिदेश हरिश्रेष्ठाञ्जाम्बवत्प्रमुखांस्ततः} %||6-63-23||

\twolineshloka
{ते तु वानरशार्दूलाः श्रुत्वा रामस्य शासनम्}
{अभिपेतुः सुसङ्क्रुद्धाः कुम्भमुद्यतकार्मुकम्} %||6-63-24||

\twolineshloka
{ततो द्रुमशिलाहस्ताः कोपसंरक्तलोचनाः}
{रिरक्षिषन्तोऽभ्यपतन्नङ्गदं वानरर्षभाः} %||6-63-25||

\twolineshloka
{जाम्बवांश्च सुषेणश्च वेगदर्शी च वानरः}
{कुम्भकर्णात्मजं वीरं क्रुद्धाः समभिदुद्रुवुः} %||6-63-26||

\twolineshloka
{समीक्ष्यातततस्तांस्तु वानरेन्द्रान्महाबलान्}
{आववार शरौघेण नगेनेव जलाशयम्} %||6-63-27||

\twolineshloka
{तस्य बाणचयं प्राप्य न शोकेरतिवर्तितुम्}
{वानरेन्द्रा महात्मानो वेलामिव महोदधिः} %||6-63-28||

\twolineshloka
{तांस्तु दृष्ट्वा हरिगणाञ्शरवृष्टिभिरर्दितान्}
{अङ्गदं पृष्ठतः कृत्वा भ्रातृजं प्लवगेश्वरः} %||6-63-29||

\twolineshloka
{अभिदुद्राव वेगेन सुग्रीवः कुम्भमाहवे}
{शैलसानु चरं नागं वेगवानिव केसरी} %||6-63-30||

\twolineshloka
{उत्पाट्य च महाशैलानश्वकर्णान्धवान्बहून्}
{अन्यांश्च विविधान्वृक्षांश्चिक्षेप च महाबलः} %||6-63-31||

\twolineshloka
{तां छादयन्तीमाकाशं वृक्षवृष्टिं दुरासदाम्}
{कुम्भकर्णात्मजः श्रीमांश्चिच्छेद निशितैः शरैः} %||6-63-32||

\twolineshloka
{अभिलक्ष्येण तीव्रेण कुम्भेन निशितैः शरैः}
{आचितास्ते द्रुमा रेजुर्यथा घोराः शतघ्नयः} %||6-63-33||

\twolineshloka
{द्रुमवर्षं तु तच्छिन्नं दृष्ट्वा कुम्भेन वीर्यवान्}
{वानराधिपतिः श्रीमान्महासत्त्वो न विव्यथे} %||6-63-34||

\twolineshloka
{निर्भिद्यमानः सहसा सहमानश्च ताञ्शरान्}
{कुम्भस्य धनुराक्षिप्य बभञ्जेन्द्रधनुःप्रभम्} %||6-63-35||

\twolineshloka
{अवप्लुत्य ततः शीघ्रं कृत्वा कर्म सुदुष्करम्}
{अब्रवीत्कुपितः कुम्भं भग्नशृङ्गमिव द्विपम्} %||6-63-36||

\twolineshloka
{निकुम्भाग्रज वीर्यं ते बाणवेगं तदद्भुतम्}
{संनतिश्च प्रभावश्च तव वा रावणस्य वा} %||6-63-37||

\twolineshloka
{प्रह्रादबलिवृत्रघ्नकुबेरवरुणोपम}
{एकस्त्वमनुजातोऽसि पितरं बलवत्तरः} %||6-63-38||

\twolineshloka
{त्वामेवैकं महाबाहुं शूलहस्तमरिन्दमम्}
{त्रिदशा नातिवर्तन्ते जितेन्द्रियमिवाधयः} %||6-63-39||

\twolineshloka
{वरदानात्पितृव्यस्ते सहते देवदानवान्}
{कुम्भकर्णस्तु वीर्येण सहते च सुरासुरान्} %||6-63-40||

\twolineshloka
{धनुषीन्द्रजितस्तुल्यः प्रतापे रावणस्य च}
{त्वमद्य रक्षसां लोके श्रेष्ठोऽसि बलवीर्यतः} %||6-63-41||

\twolineshloka
{महाविमर्दं समरे मया सह तवाद्भुतम्}
{अद्य भूतानि पश्यन्तु शक्रशम्बरयोरिव} %||6-63-42||

\twolineshloka
{कृतमप्रतिमं कर्म दर्शितं चास्त्रकौशलम्}
{पातिता हरिवीराश्च त्वयैते भीमविक्रमाः} %||6-63-43||

\twolineshloka
{उपालम्भभयाच्चापि नासि वीर मया हतः}
{कृतकर्मा परिश्रान्तो विश्रान्तः पश्य मे बलम्} %||6-63-44||

\twolineshloka
{तेन सुग्रीववाक्येन सावमानेन मानितः}
{अग्नेराज्यहुतस्येव तेजस्तस्याभ्यवर्धत} %||6-63-45||

\twolineshloka
{ततः कुम्भः समुत्पत्य सुग्रीवमभिपद्य च}
{आजघानोरसि क्रुद्धो वज्रवेगेन मुष्टिना} %||6-63-46||

\twolineshloka
{तस्य चर्म च पुस्फोट संजज्ञे चास्य शोणितम्}
{स च मुष्टिर्महावेगः प्रतिजघ्नेऽस्थिमण्डले} %||6-63-47||

\twolineshloka
{तदा वेगेन तत्रासीत्तेजः प्रज्वालितं मुहुः}
{वज्रनिष्पेषसंजातज्वाला मेरौ यथा गिरौ} %||6-63-48||

\twolineshloka
{स तत्राभिहतस्तेन सुग्रीवो वानरर्षभः}
{मुष्टिं संवर्तयामास वज्रकल्पं महाबलः} %||6-63-49||

\twolineshloka
{अर्चिःसहस्रविकचं रविमण्डलसप्रभम्}
{स मुष्टिं पातयामास कुम्भस्योरसि वीर्यवान्} %||6-63-50||

\twolineshloka
{मुष्टिनाभिहतस्तेन निपपाताशु राक्षसः}
{लोहिताङ्ग इवाकाशाद्दीप्तरश्मिर्यदृच्छया} %||6-63-51||

\twolineshloka
{कुम्भस्य पततो रूपं भग्नस्योरसि मुष्टिना}
{बभौ रुद्राभिपन्नस्य यथारूपं गवां पतेः} %||6-63-52||

\fourlineindentedshloka
{तस्मिन्हते भीमपराक्रमेण}
{प्लवङ्गमानामृषभेण युद्धे}
{मही सशैला सवना चचाल}
{भयं च रक्षांस्यधिकं विवेश} %||6-64-53||


{॥इत्यार्षे श्रीमद्रामायणे वाल्मीकीये आदिकाव्ये युद्धकाण्डे त्रिषष्टितमः सर्गः॥}

\sect{चतुःषष्टितमः सर्गः}

\twolineshloka
{निकुम्भो भ्रातरं दृष्ट्वा सुग्रीवेण निपातितम्}
{प्रदहन्निव कोपेन वानरेन्द्रमवैक्षत} %||6-64-1||

\twolineshloka
{ततः स्रग्दामसंनद्धं दत्तपञ्चाङ्गुलं शुभम्}
{आददे परिघं वीरो नगेन्द्रशिखरोपमम्} %||6-64-2||

\twolineshloka
{हेमपट्टपरिक्षिप्तं वज्रविद्रुमभूषितम्}
{यमदण्डोपमं भीमं रक्षसां भयनाशनम्} %||6-64-3||

\twolineshloka
{तमाविध्य महातेजाः शक्रध्वजसमं रणे}
{विननाद विवृत्तास्यो निकुम्भो भीमविक्रमः} %||6-64-4||

\twolineshloka
{उरोगतेन निष्केण भुजस्थैरङ्गदैरपि}
{कुण्डलाभ्यां च मृष्टाभ्यां मालया च विचित्रया} %||6-64-5||

\twolineshloka
{निकुम्भो भूषणैर्भाति तेन स्म परिघेण च}
{यथेन्द्रधनुषा मेघः सविद्युत्स्तनयित्नुमान्} %||6-64-6||

\twolineshloka
{परिघाग्रेण पुस्फोट वातग्रन्थिर्महात्मनः}
{प्रजज्वाल सघोषश्च विधूम इव पावकः} %||6-64-7||

\twolineshloka
{नगर्या विटपावत्या गन्धर्वभवनोत्तमैः}
{सह चैवामरावत्या सर्वैश्च भवनैः सह} %||6-64-8||

\twolineshloka
{सतारागणनक्षत्रं सचन्द्रं समहाग्रहम्}
{निकुम्भपरिघाघूर्णं भ्रमतीव नभस्तलम्} %||6-64-9||

\twolineshloka
{दुरासदश्च संजज्ञे परिघाभरणप्रभः}
{क्रोधेन्धनो निकुम्भाग्निर्युगान्ताग्निरिवोत्थितः} %||6-64-10||

\twolineshloka
{राक्षसा वानराश्चापि न शेकुः स्पन्दितुं भयात्}
{हनूमंस्तु विवृत्योरस्तस्थौ प्रमुखतो बली} %||6-64-11||

\twolineshloka
{परिघोपमबाहुस्तु परिघं भास्करप्रभम्}
{बली बलवतस्तस्य पातयामास वक्षसि} %||6-64-12||

\twolineshloka
{स्थिरे तस्योरसि व्यूढे परिघः शतधा कृतः}
{विशीर्यमाणः सहसा उल्का शतमिवाम्बरे} %||6-64-13||

\twolineshloka
{स तु तेन प्रहारेण चचाल च महाकपिः}
{परिघेण समाधूतो यथा भूमिचलेऽचलः} %||6-64-14||

\twolineshloka
{स तथाभिहतस्तेन हनूमान्प्लवगोत्तमः}
{मुष्टिं संवर्तयामास बलेनातिमहाबलः} %||6-64-15||

\twolineshloka
{तमुद्यम्य महातेजा निकुम्भोरसि वीर्यवान्}
{अभिचिक्षेप वेगेन वेगवान्वायुविक्रमः} %||6-64-16||

\twolineshloka
{ततः पुस्फोट चर्मास्य प्रसुस्राव च शोणितम्}
{मुष्टिना तेन संजज्ञे ज्वाला विद्युदिवोत्थिता} %||6-64-17||

\twolineshloka
{स तु तेन प्रहारेण निकुम्भो विचचाल ह}
{स्वस्थश्चापि निजग्राह हनूमन्तं महाबलम्} %||6-64-18||

\twolineshloka
{विचुक्रुशुस्तदा सङ्ख्ये भीमं लङ्कानिवासिनः}
{निकुम्भेनोद्धृतं दृष्ट्वा हनूमन्तं महाबलम्} %||6-64-19||

\twolineshloka
{स तथा ह्रियमाणोऽपि कुम्भकर्णात्मजेन हि}
{आजघानानिलसुतो वज्रवेगेन मुष्टिना} %||6-64-20||

\twolineshloka
{आत्मानं मोचयित्वाथ क्षितावभ्यवपद्यत}
{हनूमानुन्ममथाशु निकुम्भं मारुतात्मजः} %||6-64-21||

\twolineshloka
{निक्षिप्य परमायत्तो निकुम्भं निष्पिपेष च}
{उत्पत्य चास्य वेगेन पपातोरसि वीर्यवान्} %||6-64-22||

\twolineshloka
{परिगृह्य च बाहुभ्यां परिवृत्य शिरोधराम्}
{उत्पाटयामास शिरो भैरवं नदतो महत्} %||6-64-23||

\fourlineindentedshloka
{अथ विनदति सादिते निकुम्भे}
{पवनसुतेन रणे बभूव युद्धम्}
{दशरथसुतराक्षसेन्द्रचम्वोर्-}
{भृशतरमागतरोषयोः सुभीमम्} %||6-65-24||


{॥इत्यार्षे श्रीमद्रामायणे वाल्मीकीये आदिकाव्ये युद्धकाण्डे चतुःषष्टितमः सर्गः॥}

\sect{पञ्चषष्टितमः सर्गः}

\twolineshloka
{निकुम्भं च हतं श्रुत्वा कुम्भं च विनिपातितम्}
{रावणः परमामर्षी प्रजज्वालानलो यथा} %||6-65-1||

\twolineshloka
{नैरृतः क्रोधशोकाभ्यां द्वाभ्यां तु परिमूर्छितः}
{खरपुत्रं विशालाक्षं मकराक्षमचोदयत्} %||6-65-2||

\twolineshloka
{गच्छ पुत्र मयाज्ञप्तो बलेनाभिसमन्वितः}
{राघवं लक्ष्मणं चैव जहि तौ सवनौकसौ} %||6-65-3||

\twolineshloka
{रावणस्य वचः श्रुत्वा शूरो मानी खरात्मजः}
{बाढमित्यब्रवीद्धृष्टो मकराक्षो निशाचरः} %||6-65-4||

\twolineshloka
{सोऽभिवाद्य दशग्रीवं कृत्वा चापि प्रदक्षिणम्}
{निर्जगाम गृहाच्छुभ्राद्रावणस्याज्ञया बली} %||6-65-5||

\twolineshloka
{समीपस्थं बलाध्यक्षं खरपुत्रोऽब्रवीदिदम्}
{रथमानीयतां शीघ्रं सैन्यं चानीयतां त्वरात्} %||6-65-6||

\twolineshloka
{तस्य तद्वचनं श्रुत्वा बलाध्यक्षो निशाचरः}
{स्यन्दनं च बलं चैव समीपं प्रत्यपादयत्} %||6-65-7||

\twolineshloka
{प्रदक्षिणं रथं कृत्वा आरुरोह निशाचरः}
{सूतं सञ्चोदयामास शीघ्रं मे रथमावह} %||6-65-8||

\twolineshloka
{अथ तान्राक्षसान्सर्वान्मकराक्षोऽब्रवीदिदम्}
{यूयं सर्वे प्रयुध्यध्वं पुरस्तान्मम राक्षसाः} %||6-65-9||

\twolineshloka
{अहं राक्षसराजेन रावणेन महात्मना}
{आज्ञप्तः समरे हन्तुं तावुभौ रामलक्ष्मणौ} %||6-65-10||

\twolineshloka
{अद्य रामं वधिष्यामि लक्ष्मणं च निशाचराः}
{शाखामृगं च सुग्रीवं वानरांश्च शरोत्तमैः} %||6-65-11||

\twolineshloka
{अद्य शूलनिपातैश्च वानराणां महाचमूम्}
{प्रदहिष्यामि सम्प्राप्तां शुष्केन्धनमिवानलः} %||6-65-12||

\twolineshloka
{मकराक्षस्य तच्छ्रुत्वा वचनं ते निशाचराः}
{सर्वे नानायुधोपेता बलवन्तः समाहिताः} %||6-65-13||

\twolineshloka
{ते कामरूपिणः शूरा दंष्ट्रिणः पिङ्गलेक्षणाः}
{मातङ्गा इव नर्दन्तो ध्वस्तकेशा भयानकाः} %||6-65-14||

\twolineshloka
{परिवार्य महाकाया महाकायं खरात्मजम्}
{अभिजग्मुस्तदा हृष्टाश्चालयन्तो वसुन्धराम्} %||6-65-15||

\twolineshloka
{शङ्खभेरीसहस्राणामाहतानां समन्ततः}
{क्ष्वेडितास्फोटितानां च ततः शब्दो महानभूत्} %||6-65-16||

\twolineshloka
{प्रभ्रष्टोऽथ करात्तस्य प्रतोदः सारथेस्तदा}
{पपात सहसा चैव ध्वजस्तस्य च रक्षसः} %||6-65-17||

\twolineshloka
{तस्य ते रथसंयुक्ता हया विक्रमवर्जिताः}
{चरणैराकुलैर्गत्वा दीनाः सास्रमुखा ययुः} %||6-65-18||

\twolineshloka
{प्रवाति पवनस्तस्य सपांसुः खरदारुणः}
{निर्याणे तस्य रौद्रस्य मकराक्षस्य दुर्मतेः} %||6-65-19||

\twolineshloka
{तानि दृष्ट्वा निमित्तानि राक्षसा वीर्यवत्तमाः}
{अचिन्त्यनिर्गताः सर्वे यत्र तौ रामलक्ष्मणौ} %||6-65-20||

\fourlineindentedshloka
{घनगजमहिषाङ्गतुल्यवर्णाः}
{समरमुखेष्वसकृद्गदासिभिन्नाः}
{अहमहमिति युद्धकौशलास्ते}
{रजनिचराः परिबभ्रमुर्नदन्तः} %||6-66-21||


{॥इत्यार्षे श्रीमद्रामायणे वाल्मीकीये आदिकाव्ये युद्धकाण्डे पञ्चषष्टितमः सर्गः॥}

\sect{षट्षष्टितमः सर्गः}

\twolineshloka
{निर्गतं मकराक्षं ते दृष्ट्वा वानरपुङ्गवाः}
{आप्लुत्य सहसा सर्वे योद्धुकामा व्यवस्थिताः} %||6-66-1||

\twolineshloka
{ततः प्रवृत्तं सुमहत्तद्युद्धं लोमहर्षणम्}
{निशाचरैः प्लवङ्गानां देवानां दानवैरिव} %||6-66-2||

\twolineshloka
{वृक्षशूलनिपातैश्च शिलापरिघपातनैः}
{अन्योन्यं मर्दयन्ति स्म तदा कपिनिशाचराः} %||6-66-3||

\twolineshloka
{शक्तिशूलगदाखड्गैस्तोमरैश्च निशाचराः}
{पट्टसैर्भिन्दिपालैश्च बाणपातैः समन्ततः} %||6-66-4||

\twolineshloka
{पाशमुद्गरदण्डैश्च निर्घातैश्चापरैस्तथा}
{कदनं कपिसिंहानां चक्रुस्ते रजनीचराः} %||6-66-5||

\twolineshloka
{बाणौघैरर्दिताश्चापि खरपुत्रेण वानराः}
{सम्भ्रान्तमनसः सर्वे दुद्रुवुर्भयपीडिताः} %||6-66-6||

\twolineshloka
{तान्दृष्ट्वा राक्षसाः सर्वे द्रवमाणान्वनौकसः}
{नेदुस्ते सिंहवद्धृष्टा राक्षसा जितकाशिनः} %||6-66-7||

\twolineshloka
{विद्रवत्सु तदा तेषु वानरेषु समन्ततः}
{रामस्तान्वारयामास शरवर्षेण राक्षसान्} %||6-66-8||

\twolineshloka
{वारितान्राक्षसान्दृष्ट्वा मकराक्षो निशाचरः}
{क्रोधानलसमाविष्टो वचनं चेदमब्रवीत्} %||6-66-9||

\twolineshloka
{तिष्ठ राम मया सार्धं द्वन्द्वयुद्धं ददामि ते}
{त्याजयिष्यामि ते प्राणान्धनुर्मुक्तैः शितैः शरैः} %||6-66-10||

\twolineshloka
{यत्तदा दण्डकारण्ये पितरं हतवान्मम}
{मदग्रतः स्वकर्मस्थं स्मृत्वा रोषोऽभिवर्धते} %||6-66-11||

\twolineshloka
{दह्यन्ते भृशमङ्गानि दुरात्मन्मम राघव}
{यन्मयासि न दृष्टस्त्वं तस्मिन्काले महावने} %||6-66-12||

\twolineshloka
{दिष्ट्यासि दर्शनं राम मम त्वं प्राप्तवानिह}
{काङ्क्षितोऽसि क्षुधार्तस्य सिंहस्येवेतरो मृगः} %||6-66-13||

\twolineshloka
{अद्य मद्बाणवेगेन प्रेतराड्विषयं गतः}
{ये त्वया निहताः शूराः सह तैस्त्वं समेष्यसि} %||6-66-14||

\twolineshloka
{बहुनात्र किमुक्तेन शृणु राम वचो मम}
{पश्यन्तु सकला लोकास्त्वां मां चैव रणाजिरे} %||6-66-15||

\twolineshloka
{अस्त्रैर्वा गदया वापि बाहुभ्यां वा महाहवे}
{अभ्यस्तं येन वा राम तेन वा वर्ततां युधि} %||6-66-16||

\twolineshloka
{मकराक्षवचः श्रुत्वा रामो दशरथात्मजः}
{अब्रवीत्प्रहसन्वाक्यमुत्तरोत्तरवादिनम्} %||6-66-17||

\twolineshloka
{चतुर्दशसहस्राणि रक्षसां त्वत्पिता च यः}
{त्रिशिरा दूषणश्चापि दण्डके निहता मया} %||6-66-18||

\twolineshloka
{स्वाशितास्तव मांसेन गृध्रगोमायुवायसाः}
{भविष्यन्त्यद्य वै पाप तीक्ष्णतुण्डनखाङ्कुशाः} %||6-66-19||

\twolineshloka
{एवमुक्तस्तु रामेण खरपुत्रो निशाचरः}
{बाणौघानसृजत्तस्मै राघवाय रणाजिरे} %||6-66-20||

\twolineshloka
{ताञ्शराञ्शरवर्षेण रामश्चिच्छेद नैकधा}
{निपेतुर्भुवि ते छिन्ना रुक्मपुङ्खाः सहस्रशः} %||6-66-21||

\twolineshloka
{तद्युद्धमभवत्तत्र समेत्यान्योन्यमोजसा}
{खर राक्षसपुत्रस्य सूनोर्दशरथस्य च} %||6-66-22||

\twolineshloka
{जीमूतयोरिवाकाशे शब्दो ज्यातलयोस्तदा}
{धनुर्मुक्तः स्वनोत्कृष्टः श्रूयते च रणाजिरे} %||6-66-23||

\twolineshloka
{देवदानवगन्धर्वाः किंनराश्च महोरगाः}
{अन्तरिक्षगताः सर्वे द्रष्टुकामास्तदद्भुतम्} %||6-66-24||

\twolineshloka
{विद्धमन्योन्यगात्रेषु द्विगुणं वर्धते बलम्}
{कृतप्रतिकृतान्योन्यं कुर्वाते तौ रणाजिरे} %||6-66-25||

\twolineshloka
{राममुक्तास्तु बाणौघान्राक्षसस्त्वच्छिनद्रणे}
{रक्षोमुक्तांस्तु रामो वै नैकधा प्राच्छिनच्छरैः} %||6-66-26||

\twolineshloka
{बाणौघवितताः सर्वा दिशश्च विदिशस्तथा}
{सञ्छन्ना वसुधा चैव समन्तान्न प्रकाशते} %||6-66-27||

\threelineshloka
{ततः क्रुद्धो महाबाहुर्धनुश्चिच्छेद रक्षसः}
{अष्टाभिरथ नाराचैः सूतं विव्याध राघवः}
{भित्त्वा शरै रथं रामो रथाश्वान्समपातयत्} %||6-66-28||

\threelineshloka
{विरथो वसुधां तिष्ठन्मकराक्षो निशाचरः}
{अतिष्ठद्वसुधां रक्षः शूलं जग्राह पाणिना}
{त्रासनं सर्वभूतानां युगान्ताग्निसमप्रभम्} %||6-66-29||

\twolineshloka
{विभ्राम्य च महच्छूलं प्रज्वलन्तं निशाचरः}
{स क्रोधात्प्राहिणोत्तस्मै राघवाय महाहवे} %||6-66-30||

\twolineshloka
{तमापतन्तं ज्वलितं खरपुत्रकराच्च्युतम्}
{बाणैस्तु त्रिभिराकाशे शूलं चिच्छेद राघवः} %||6-66-31||

\twolineshloka
{सच्छिन्नो नैकधा शूलो दिव्यहाटकमण्डितः}
{व्यशीर्यत महोक्लेव रामबाणार्दितो भुवि} %||6-66-32||

\twolineshloka
{तच्छूलं निहतं दृष्ट्वा रामेणाद्भुतकर्मणा}
{साधु साध्विति भूतानि व्याहरन्ति नभोगताः} %||6-66-33||

\twolineshloka
{तद्दृष्ट्वा निहतं शूलं मकराक्षो निशाचरः}
{मुष्टिमुद्यम्य काकुत्स्थं तिष्ठ तिष्ठेति चाब्रवीत्} %||6-66-34||

\twolineshloka
{स तं दृष्ट्वा पतन्तं वै प्रहस्य रघुनन्दनः}
{पावकास्त्रं ततो रामः सन्दधे स्वशरासने} %||6-66-35||

\twolineshloka
{तेनास्त्रेण हतं रक्षः काकुत्स्थेन तदा रणे}
{सञ्छिन्नहृदयं तत्र पपात च ममार च} %||6-66-36||

\twolineshloka
{दृष्ट्वा ते राक्षसाः सर्वे मकराक्षस्य पातनम्}
{लङ्कामेव प्रधावन्त रामबालार्दितास्तदा} %||6-66-37||

\fourlineindentedshloka
{दशरथनृपपुत्रबाणवेगै}
{रजनिचरं निहतं खरात्मजं तम्}
{ददृशुरथ च देवताः प्रहृष्टा}
{गिरिमिव वज्रहतं यथा विशीर्णम्} %||6-67-38||


{॥इत्यार्षे श्रीमद्रामायणे वाल्मीकीये आदिकाव्ये युद्धकाण्डे षट्षष्टितमः सर्गः॥}

\sect{सप्तषष्टितमः सर्गः}

\twolineshloka
{मकराक्षं हतं श्रुत्वा रावणः समितिंजयः}
{आदिदेशाथ सङ्क्रुद्धो रणायेन्द्रजितं सुतम्} %||6-67-1||

\twolineshloka
{जहि वीर महावीर्यौ भ्रातरौ रामलक्ष्मणौ}
{अदृश्यो दृश्यमानो वा सर्वथा त्वं बलाधिकः} %||6-67-2||

\twolineshloka
{त्वमप्रतिमकर्माणमिन्द्रं जयसि संयुगे}
{किं पुनर्मानुषौ दृष्ट्वा न वधिष्यसि संयुगे} %||6-67-3||

\twolineshloka
{तथोक्तो राक्षसेन्द्रेण प्रतिगृह्य पितुर्वचः}
{यज्ञभूमौ स विधिवत्पावकं जुहुवे न्द्रजित्} %||6-67-4||

\twolineshloka
{जुह्वतश्चापि तत्राग्निं रक्तोष्णीषधराः स्त्रियः}
{आजग्मुस्तत्र सम्भ्रान्ता राक्षस्यो यत्र रावणिः} %||6-67-5||

\twolineshloka
{शस्त्राणि शरपत्राणि समिधोऽथ विभीतकाः}
{लोहितानि च वासांसि स्रुवं कार्ष्णायसं तथा} %||6-67-6||

\twolineshloka
{सर्वतोऽग्निं समास्तीर्य शरपत्रैः समन्ततः}
{छागस्य सर्वकृष्णस्य गलं जग्राह जीवतः} %||6-67-7||

\twolineshloka
{चरुहोमसमिद्धस्य विधूमस्य महार्चिषः}
{बभूवुस्तानि लिङ्गानि विजयं दर्शयन्ति च} %||6-67-8||

\twolineshloka
{प्रदक्षिणावर्तशिखस्तप्तहाटकसंनिभः}
{हविस्तत्प्रतिजग्राह पावकः स्वयमुत्थितः} %||6-67-9||

\twolineshloka
{हुत्वाग्निं तर्पयित्वाथ देवदानवराक्षसान्}
{आरुरोह रथश्रेष्ठमन्तर्धानगतं शुभम्} %||6-67-10||

\twolineshloka
{स वाजिभिश्चतुर्भिस्तु बाणैश्च निशितैर्युतः}
{आरोपितमहाचापः शुशुभे स्यन्दनोत्तमे} %||6-67-11||

\twolineshloka
{जाज्वल्यमानो वपुषा तपनीयपरिच्छदः}
{शरैश्चन्द्रार्धचन्द्रैश्च स रथः समलङ्कृतः} %||6-67-12||

\twolineshloka
{जाम्बूनदमहाकम्बुर्दीप्तपावकसंनिभः}
{बभूवेन्द्रजितः केतुर्वैदूर्यसमलङ्कृतः} %||6-67-13||

\twolineshloka
{तेन चादित्यकल्पेन ब्रह्मास्त्रेण च पालितः}
{स बभूव दुराधर्षो रावणिः सुमहाबलः} %||6-67-14||

\twolineshloka
{सोऽभिनिर्याय नगरादिन्द्रजित्समितिंजयः}
{हुत्वाग्निं राक्षसैर्मन्त्रैरन्तर्धानगतोऽब्रवीत्} %||6-67-15||

\twolineshloka
{अद्य हत्वाहवे यौ तौ मिथ्या प्रव्रजितौ वने}
{जयं पित्रे प्रदास्यामि रावणाय रणाधिकम्} %||6-67-16||

\twolineshloka
{कृत्वा निर्वानरामुर्वीं हत्वा रामं सलक्ष्मणम्}
{करिष्ये परमां प्रीतिमित्युक्त्वान्तरधीयत} %||6-67-17||

\twolineshloka
{आपपाताथ सङ्क्रुद्धो दशग्रीवेण चोदितः}
{तीक्ष्णकार्मुकनाराचैस्तीक्ष्णस्त्विन्द्ररिपू रणे} %||6-67-18||

\twolineshloka
{स ददर्श महावीर्यौ नागौ त्रिशिरसाविव}
{सृजन्ताविषुजालानि वीरौ वानरमध्यगौ} %||6-67-19||

\twolineshloka
{इमौ ताविति सञ्चिन्त्य सज्यं कृत्वा च कार्मुकम्}
{सन्ततानेषुधाराभिः पर्जन्य इव वृष्टिमान्} %||6-67-20||

\twolineshloka
{स तु वैहायसं प्राप्य सरथो रामलक्ष्मणौ}
{अचक्षुर्विषये तिष्ठन्विव्याध निशितैः शरैः} %||6-67-21||

\twolineshloka
{तौ तस्य शरवेगेन परीतौ रामलक्ष्मणौ}
{धनुषी सशरे कृत्वा दिव्यमस्त्रं प्रचक्रतुः} %||6-67-22||

\twolineshloka
{प्रच्छादयन्तौ गगनं शरजालैर्महाबलौ}
{तमस्त्रैः सुरसङ्काशौ नैव पस्पर्शतुः शरैः} %||6-67-23||

\twolineshloka
{स हि धूमान्धकारं च चक्रे प्रच्छादयन्नभः}
{दिशश्चान्तर्दधे श्रीमान्नीहारतमसावृतः} %||6-67-24||

\twolineshloka
{नैव ज्यातलनिर्घोषो न च नेमिखुरस्वनः}
{शुश्रुवे चरतस्तस्य न च रूपं प्रकाशते} %||6-67-25||

\twolineshloka
{घनान्धकारे तिमिरे शरवर्षमिवाद्भुतम्}
{स ववर्ष महाबाहुर्नाराचशरवृष्टिभिः} %||6-67-26||

\twolineshloka
{स रामं सूर्यसङ्काशैः शरैर्दत्तवरो भृशम्}
{विव्याध समरे क्रुद्धः सर्वगात्रेषु रावणिः} %||6-67-27||

\twolineshloka
{तौ हन्यमानौ नाराचैर्धाराभिरिव पर्वतौ}
{हेमपुङ्खान्नरव्याघ्रौ तिग्मान्मुमुचतुः शरान्} %||6-67-28||

\twolineshloka
{अन्तरिक्षं समासाद्य रावणिं कङ्कपत्रिणः}
{निकृत्य पतगा भूमौ पेतुस्ते शोणितोक्षिताः} %||6-67-29||

\twolineshloka
{अतिमात्रं शरौघेण पीड्यमानौ नरोत्तमौ}
{तानिषून्पततो भल्लैरनेकैर्निचकर्ततुः} %||6-67-30||

\twolineshloka
{यतो हि ददृशाते तौ शरान्निपतिताञ्शितान्}
{ततस्ततो दाशरथी ससृजातेऽस्त्रमुत्तमम्} %||6-67-31||

\twolineshloka
{रावणिस्तु दिशः सर्वा रथेनातिरथः पतन्}
{विव्याध तौ दाशरथी लघ्वस्त्रो निशितैः शरैः} %||6-67-32||

\twolineshloka
{तेनातिविद्धौ तौ वीरौ रुक्मपुङ्खैः सुसंहतैः}
{बभूवतुर्दाशरथी पुष्पिताविव किंशुकौ} %||6-67-33||

\twolineshloka
{नास्य वेद गतिं कश्चिन्न च रूपं धनुः शरान्}
{न चान्यद्विदितं किञ्चित्सूर्यस्येवाभ्रसम्प्लवे} %||6-67-34||

\twolineshloka
{तेन विद्धाश्च हरयो निहताश्च गतासवः}
{बभूवुः शतशस्तत्र पतिता धरणीतले} %||6-67-35||

\twolineshloka
{लक्ष्मणस्तु सुसङ्क्रुद्धो भ्रातरं वाक्यमब्रवीत्}
{ब्राह्ममस्त्रं प्रयोक्ष्यामि वधार्थं सर्वरक्षसाम्} %||6-67-36||

\twolineshloka
{तमुवाच ततो रामो लक्ष्मणं शुभलक्षणम्}
{नैकस्य हेतो रक्षांसि पृथिव्यां हन्तुमर्हसि} %||6-67-37||

\twolineshloka
{अयुध्यमानं प्रच्छन्नं प्राञ्जलिं शरणागतम्}
{पलायन्तं प्रमत्तं वा न त्वं हन्तुमिहार्हसि} %||6-67-38||

\twolineshloka
{अस्यैव तु वधे यत्नं करिष्यावो महाबल}
{आदेक्ष्यावो महावेगानस्त्रानाशीविषोपमान्} %||6-67-39||

\twolineshloka
{तमेनं मायिनं क्षुद्रमन्तर्हितरथं बलात्}
{राक्षसं निहनिष्यन्ति दृष्ट्वा वानरयूथपाः} %||6-67-40||

\fourlineindentedshloka
{यद्येष भूमिं विशते दिवं वा}
{रसातलं वापि नभस्तलं वा}
{एवं निगूढोऽपि ममास्त्रदग्धः}
{पतिष्यते भूमितले गतासुः} %||6-67-41||

\fourlineindentedshloka
{इत्येवमुक्त्वा वचनं महात्मा}
{रघुप्रवीरः प्लवगर्षभैर्वृतः}
{वधाय रौद्रस्य नृशंसकर्मणः}
{तदा महात्मा त्वरितं निरीक्षते} %||6-68-42||


{॥इत्यार्षे श्रीमद्रामायणे वाल्मीकीये आदिकाव्ये युद्धकाण्डे सप्तषष्टितमः सर्गः॥}

\sect{अष्टषष्टितमः सर्गः}

\twolineshloka
{विज्ञाय तु मनस्तस्य राघवस्य महात्मनः}
{संनिवृत्याहवात्तस्मात्प्रविवेश पुरं ततः} %||6-68-1||

\twolineshloka
{सोऽनुस्मृत्य वधं तेषां राक्षसानां तरस्विनाम्}
{क्रोधताम्रेक्षणः शूरो निर्जगाम महाद्युतिः} %||6-68-2||

\twolineshloka
{स पश्चिमेन द्वारेण निर्ययौ राक्षसैर्वृतः}
{इन्द्रजित्तु महावीर्यः पौलस्त्यो देवकण्टकः} %||6-68-3||

\twolineshloka
{इन्द्रजित्तु ततो दृष्ट्वा भ्रातरौ रामलक्ष्मणौ}
{रणायाभ्युद्यतौ वीरौ मायां प्रादुष्करोत्तदा} %||6-68-4||

\twolineshloka
{इन्द्रजित्तु रथे स्थाप्य सीतां मायामयीं तदा}
{बलेन महतावृत्य तस्या वधमरोचयत्} %||6-68-5||

\twolineshloka
{मोहनार्थं तु सर्वेषां बुद्धिं कृत्वा सुदुर्मतिः}
{हन्तुं सीतां व्यवसितो वानराभिमुखो ययौ} %||6-68-6||

\twolineshloka
{तं दृष्ट्वा त्वभिनिर्यान्तं नगर्याः काननौकसः}
{उत्पेतुरभिसङ्क्रुद्धाः शिलाहस्ता युयुत्सवः} %||6-68-7||

\twolineshloka
{हनूमान्पुरतस्तेषां जगाम कपिकुञ्जरः}
{प्रगृह्य सुमहच्छृङ्गं पर्वतस्य दुरासदम्} %||6-68-8||

\twolineshloka
{स ददर्श हतानन्दां सीतामिन्द्रजितो रथे}
{एकवेणीधरां दीनामुपवासकृशाननाम्} %||6-68-9||

\twolineshloka
{परिक्लिष्टैकवसनाममृजां राघवप्रियाम्}
{रजोमलाभ्यामालिप्तैः सर्वगात्रैर्वरस्त्रियम्} %||6-68-10||

\twolineshloka
{तां निरीक्ष्य मुहूर्तं तु मैथिलीमध्यवस्य च}
{बाष्पपर्याकुलमुखो हनूमान्व्यथितोऽभवत्} %||6-68-11||

\twolineshloka
{अब्रवीत्तां तु शोकार्तां निरानन्दां तपस्विनाम्}
{दृष्ट्वा रथे स्तितां सीतां राक्षसेन्द्रसुताश्रिताम्} %||6-68-12||

\twolineshloka
{किं समर्थितमस्येति चिन्तयन्स महाकपिः}
{सह तैर्वानरश्रेष्ठैरभ्यधावत रावणिम्} %||6-68-13||

\twolineshloka
{तद्वानरबलं दृष्ट्वा रावणिः क्रोधमूर्छितः}
{कृत्वा विशोकं निस्त्रिंशं मूर्ध्नि सीतां परामृशत्} %||6-68-14||

\twolineshloka
{तं स्त्रियं पश्यतां तेषां ताडयामास रावणिः}
{क्रोशन्तीं राम रामेति मायया योजितां रथे} %||6-68-15||

\threelineshloka
{गृहीतमूर्धजां दृष्ट्वा हनूमान्दैन्यमागतः}
{दुःखजं वारिनेत्राभ्यामुत्सृजन्मारुतात्मजः}
{अब्रवीत्परुषं वाक्यं क्रोधाद्रक्षोऽधिपात्मजम्} %||6-68-16||

\threelineshloka
{दुरात्मन्नात्मनाशाय केशपक्षे परामृशः}
{ब्रह्मर्षीणां कुले जातो राक्षसीं योनिमाश्रितः}
{धिक्त्वां पापसमाचारं यस्य ते मतिरीदृशी} %||6-68-17||

\twolineshloka
{नृशंसानार्य दुर्वृत्त क्षुद्र पापपराक्रम}
{अनार्यस्येदृशं कर्म घृणा ते नास्ति निर्घृण} %||6-68-18||

\twolineshloka
{च्युता गृहाच्च राज्याच्च रामहस्ताच्च मैथिली}
{किं तवैषापराद्धा हि यदेनां हन्तुमिच्छसि} %||6-68-19||

\twolineshloka
{सीतां च हत्वा न चिरं जीविष्यसि कथञ्चन}
{वधार्हकर्मणानेन मम हस्तगतो ह्यसि} %||6-68-20||

\twolineshloka
{ये च स्त्रीघातिनां लोका लोकवध्यैश्च कुत्सिताः}
{इह जीवितमुत्सृज्य प्रेत्य तान्प्रतिलप्स्यसे} %||6-68-21||

\twolineshloka
{इति ब्रुवाणो हनुमान्सायुधैर्हरिभिर्वृतः}
{अभ्यधावत सङ्क्रुद्धो राक्षसेन्द्रसुतं प्रति} %||6-68-22||

\twolineshloka
{आपतन्तं महावीर्यं तदनीकं वनौकसाम्}
{रक्षसां भीमवेगानामनीकेन न्यवारयत्} %||6-68-23||

\twolineshloka
{स तां बाणसहस्रेण विक्षोभ्य हरिवाहिनीम्}
{हरिश्रेष्ठं हनूमन्तमिन्द्रजित्प्रत्युवाच ह} %||6-68-24||

\twolineshloka
{सुग्रीवस्त्वं च रामश्च यन्निमित्तमिहागताः}
{तां हनिष्यामि वैदेहीमद्यैव तव पश्यतः} %||6-68-25||

\twolineshloka
{इमां हत्वा ततो रामं लक्ष्मणं त्वां च वानर}
{सुग्रीवं च वधिष्यामि तं चानार्यं विभीषणम्} %||6-68-26||

\twolineshloka
{न हन्तव्याः स्त्रियश्चेति यद्ब्रवीषि प्लवङ्गम}
{पीडा करममित्राणां यत्स्यात्कर्तव्यमेत तत्} %||6-68-27||

\twolineshloka
{तमेवमुक्त्वा रुदतीं सीतां मायामयीं ततः}
{शितधारेण खड्गेन निजघानेन्द्रजित्स्वयम्} %||6-68-28||

\twolineshloka
{यज्ञोपवीतमार्गेण छिन्ना तेन तपस्विनी}
{सा पृथिव्यां पृथुश्रोणी पपात प्रियदर्शना} %||6-68-29||

\twolineshloka
{तामिन्द्रजित्स्त्रियं हत्वा हनूमन्तमुवाच ह}
{मया रामस्य पश्येमां कोपेन च निषूदिताम्} %||6-68-30||

\twolineshloka
{ततः खड्गेन महता हत्वा तामिन्द्रजित्स्वयम्}
{हृष्टः स रथमास्थाय विननाद महास्वनम्} %||6-68-31||

\twolineshloka
{वानराः शुश्रुवुः शब्दमदूरे प्रत्यवस्थिताः}
{व्यादितास्यस्य नदतस्तद्दुर्गं संश्रितस्य तु} %||6-68-32||

\fourlineindentedshloka
{तथा तु सीतां विनिहत्य दुर्मतिः}
{प्रहृष्टचेताः स बभूव रावणिः}
{तं हृष्टरूपं समुदीक्ष्य वानरा}
{विषण्णरूपाः समभिप्रदुद्रुवुः} %||6-69-33||


{॥इत्यार्षे श्रीमद्रामायणे वाल्मीकीये आदिकाव्ये युद्धकाण्डे अष्टषष्टितमः सर्गः॥}

\sect{एकोनसप्ततितमः सर्गः}

\twolineshloka
{श्रुत्वा तं भीमनिर्ह्रादं शक्राशनिसमस्वनम्}
{वीक्षमाणा दिशः सर्वा दुद्रुवुर्वानरर्षभाः} %||6-69-1||

\twolineshloka
{तानुवाच ततः सर्वान्हनूमान्मारुतात्मजः}
{विषण्णवदनान्दीनांस्त्रस्तान्विद्रवतः पृथक्} %||6-69-2||

\twolineshloka
{कस्माद्विषण्णवदना विद्रवध्वं प्लवङ्गमाः}
{त्यक्तयुद्धसमुत्साहाः शूरत्वं क्व नु वो गतम्} %||6-69-3||

\twolineshloka
{पृष्ठतोऽनुव्रजध्वं मामग्रतो यान्तमाहवे}
{शूरैरभिजनोपेतैरयुक्तं हि निवर्तितुम्} %||6-69-4||

\twolineshloka
{एवमुक्ताः सुसङ्क्रुद्धा वायुपुत्रेण धीमता}
{शैलशृङ्गान्द्रुमांश्चैव जगृहुर्हृष्टमानसाः} %||6-69-5||

\twolineshloka
{अभिपेतुश्च गर्जन्तो राक्षसान्वानरर्षभाः}
{परिवार्य हनूमन्तमन्वयुश्च महाहवे} %||6-69-6||

\twolineshloka
{स तैर्वानरमुख्यैस्तु हनूमान्सर्वतो वृतः}
{हुताशन इवार्चिष्मानदहच्छत्रुवाहिनीम्} %||6-69-7||

\twolineshloka
{स राक्षसानां कदनं चकार सुमहाकपिः}
{वृतो वानरसैन्येन कालान्तकयमोपमः} %||6-69-8||

\twolineshloka
{स तु शोकेन चाविष्टः क्रोधेन च महाकपिः}
{हनूमान्रावणि रथे महतीं पातयच्छिलाम्} %||6-69-9||

\twolineshloka
{तामापतन्तीं दृष्ट्वैव रथः सारथिना तदा}
{विधेयाश्व समायुक्तः सुदूरमपवाहितः} %||6-69-10||

\twolineshloka
{तमिन्द्रजितमप्राप्य रथथं सहसारथिम्}
{विवेश धरणीं भित्त्वा सा शिलाव्यर्थमुद्यता} %||6-69-11||

\twolineshloka
{पतितायां शिलायां तु रक्षसां व्यथिता चमूः}
{तमभ्यधावञ्शतशो नदन्तः काननौकसः} %||6-69-12||

\twolineshloka
{ते द्रुमांश्च महाकाया गिरिशृङ्गाणि चोद्यताः}
{चिक्षिपुर्द्विषतां मध्ये वानरा भीमविक्रमाः} %||6-69-13||

\twolineshloka
{वानरैर्तैर्महावीर्यैर्घोररूपा निशाचराः}
{वीर्यादभिहता वृक्षैर्व्यवेष्टन्त रणक्षितौ} %||6-69-14||

\twolineshloka
{स्वसैन्यमभिवीक्ष्याथ वानरार्दितमिन्द्रजित्}
{प्रगृहीतायुधः क्रुद्धः परानभिमुखो ययौ} %||6-69-15||

\twolineshloka
{स शरौघानवसृजन्स्वसैन्येनाभिसंवृतः}
{जघान कपिशार्दूलान्सुबहून्दृष्टविक्रमः} %||6-69-16||

\twolineshloka
{शूलैरशनिभिः खड्गैः पट्टसैः कूटमुद्गरैः}
{ते चाप्यनुचरांस्तस्य वानरा जघ्नुराहवे} %||6-69-17||

\twolineshloka
{सस्कन्धविटपैः सालैः शिलाभिश्च महाबलैः}
{हनूमान्कदनं चक्रे रक्षसां भीमकर्मणाम्} %||6-69-18||

\twolineshloka
{स निवार्य परानीकमब्रवीत्तान्वनौकसः}
{हनूमान्संनिवर्तध्वं न नः साध्यमिदं बलम्} %||6-69-19||

\twolineshloka
{त्यक्त्वा प्राणान्विचेष्टन्तो राम प्रियचिकीर्षवः}
{यन्निमित्तं हि युध्यामो हता सा जनकात्मजा} %||6-69-20||

\twolineshloka
{इममर्थं हि विज्ञाप्य रामं सुग्रीवमेव च}
{तौ यत्प्रतिविधास्येते तत्करिष्यामहे वयम्} %||6-69-21||

\twolineshloka
{इत्युक्त्वा वानरश्रेष्ठो वारयन्सर्ववानरान्}
{शनैः शनैरसन्त्रस्तः सबलः स न्यवर्तत} %||6-69-22||

\twolineshloka
{स तु प्रेक्ष्य हनूमन्तं व्रजन्तं यत्र राघवः}
{निकुम्भिलामधिष्ठाय पावकं जुहुवे न्द्रजित्} %||6-69-23||

\twolineshloka
{यज्ञभूम्यां तु विधिवत्पावकस्तेन रक्षसा}
{हूयमानः प्रजज्वाल होमशोणितभुक्तदा} %||6-69-24||

\twolineshloka
{सोऽर्चिः पिनद्धो ददृशे होमशोणिततर्पितः}
{सन्ध्यागत इवादित्यः स तीव्राग्निः समुत्थितः} %||6-69-25||

\fourlineindentedshloka
{अथेन्द्रजिद्राक्षसभूतये तु}
{जुहाव हव्यं विधिना विधानवत्}
{दृष्ट्वा व्यतिष्ठन्त च राक्षसास्ते}
{महासमूहेषु नयानयज्ञाः} %||6-70-26||


{॥इत्यार्षे श्रीमद्रामायणे वाल्मीकीये आदिकाव्ये युद्धकाण्डे एकोनसप्ततितमः सर्गः॥}

\sect{सप्ततितमः सर्गः}

\twolineshloka
{राघवश्चापि विपुलं तं राक्षसवनौकसाम्}
{श्रुत्वा सङ्ग्रामनिर्घोषं जाम्बवन्तमुवाच ह} %||6-70-1||

\twolineshloka
{सौम्य नूनं हनुमता कृतं कर्म सुदुष्करम्}
{श्रूयते हि यथा भीमः सुमहानायुधस्वनः} %||6-70-2||

\twolineshloka
{तद्गच्छ कुरु साहाय्यं स्वबलेनाभिसंवृतः}
{क्षिप्रमृष्कपते तस्य कपिश्रेष्ठस्य युध्यतः} %||6-70-3||

\twolineshloka
{ऋक्षराजस्तथेत्युक्त्वा स्वेनानीकेन संवृतः}
{आगच्छत्पश्चिमद्वारं हनूमान्यत्र वानरः} %||6-70-4||

\twolineshloka
{अथायान्तं हनूमन्तं ददर्शर्क्षपतिः पथि}
{वानरैः कृतसङ्ग्रामैः श्वसद्भिरभिसंवृतम्} %||6-70-5||

\twolineshloka
{दृष्ट्वा पथि हनूमांश्च तदृष्कबलमुद्यतम्}
{नीलमेघनिभं भीमं संनिवार्य न्यवर्तत} %||6-70-6||

\twolineshloka
{स तेन हरिसैन्येन संनिकर्षं महायशाः}
{शीघ्रमागम्य रामाय दुःखितो वाक्यमब्रवीत्} %||6-70-7||

\twolineshloka
{समरे युध्यमानानामस्माकं प्रेक्षतां च सः}
{जघान रुदतीं सीतामिन्द्रजिद्रावणात्मजः} %||6-70-8||

\twolineshloka
{उद्भ्रान्तचित्तस्तां दृष्ट्वा विषण्णोऽहमरिन्दम}
{तदहं भवतो वृत्तं विज्ञापयितुमागतः} %||6-70-9||

\twolineshloka
{तस्य तद्वचनं श्रुत्वा राघवः शोकमूर्छितः}
{निपपात तदा भूमौ छिन्नमूल इव द्रुमः} %||6-70-10||

\twolineshloka
{तं भूमौ देवसङ्काशं पतितं दृश्य राघवम्}
{अभिपेतुः समुत्पत्य सर्वतः कपिसत्तमाः} %||6-70-11||

\twolineshloka
{असिञ्चन्सलिलैश्चैनं पद्मोत्पलसुगन्धिभिः}
{प्रदहन्तमसह्यं च सहसाग्निमिवोत्थितम्} %||6-70-12||

\twolineshloka
{तं लक्ष्मणोऽथ बाहुभ्यां परिष्वज्य सुदुःखितः}
{उवाच राममस्वस्थं वाक्यं हेत्वर्थसंहितम्} %||6-70-13||

\twolineshloka
{शुभे वर्त्मनि तिष्ठन्तं त्वामार्यविजितेन्द्रियम्}
{अनर्थेभ्यो न शक्नोति त्रातुं धर्मो निरर्थकः} %||6-70-14||

\twolineshloka
{भूतानां स्थावराणां च जङ्गमानां च दर्शनम्}
{यथास्ति न तथा धर्मस्तेन नास्तीति मे मतिः} %||6-70-15||

\twolineshloka
{यथैव स्थावरं व्यक्तं जङ्गमं च तथाविधम्}
{नायमर्थस्तथा युक्तस्त्वद्विधो न विपद्यते} %||6-70-16||

\twolineshloka
{यद्यधर्मो भवेद्भूतो रावणो नरकं व्रजेत्}
{भवांश्च धर्मसंयुक्तो नैवं व्यसनमाप्नुयात्} %||6-70-17||

\twolineshloka
{तस्य च व्यसनाभावाद्व्यसनं च गते त्वयि}
{धर्मेणोपलभेद्धर्ममधर्मं चाप्यधर्मतः} %||6-70-18||

\twolineshloka
{यदि धर्मेण युज्येरन्नाधर्मरुचयो जनाः}
{धर्मेण चरतां धर्मस्तथा चैषां फलं भवेत्} %||6-70-19||

\twolineshloka
{यस्मादर्था विवर्धन्ते येष्वधर्मः प्रतिष्ठितः}
{क्लिश्यन्ते धर्मशीलाश्च तस्मादेतौ निरर्थकौ} %||6-70-20||

\twolineshloka
{वध्यन्ते पापकर्माणो यद्यधर्मेण राघव}
{वधकर्महतो धर्मः स हतः कं वधिष्यति} %||6-70-21||

\twolineshloka
{अथ वा विहितेनायं हन्यते हन्ति वा परम्}
{विधिरालिप्यते तेन न स पापेन कर्मणा} %||6-70-22||

\twolineshloka
{अदृष्टप्रतिकारेण अव्यक्तेनासता सता}
{कथं शक्यं परं प्राप्तुं धर्मेणारिविकर्शन} %||6-70-23||

\twolineshloka
{यदि सत्स्यात्सतां मुख्य नासत्स्यात्तव किञ्चन}
{त्वया यदीदृशं प्राप्तं तस्मात्सन्नोपपद्यते} %||6-70-24||

\twolineshloka
{अथ वा दुर्बलः क्लीबो बलं धर्मोऽनुवर्तते}
{दुर्बलो हृतमर्यादो न सेव्य इति मे मतिः} %||6-70-25||

\twolineshloka
{बलस्य यदि चेद्धर्मो गुणभूतः पराक्रमे}
{धर्ममुत्सृज्य वर्तस्व यथा धर्मे तथा बले} %||6-70-26||

\twolineshloka
{अथ चेत्सत्यवचनं धर्मः किल परन्तप}
{अनृतस्त्वय्यकरुणः किं न बद्धस्त्वया पिता} %||6-70-27||

\twolineshloka
{यदि धर्मो भवेद्भूत अधर्मो वा परन्तप}
{न स्म हत्वा मुनिं वज्री कुर्यादिज्यां शतक्रतुः} %||6-70-28||

\twolineshloka
{अधर्मसंश्रितो धर्मो विनाशयति राघव}
{सर्वमेतद्यथाकामं काकुत्स्थ कुरुते नरः} %||6-70-29||

\twolineshloka
{मम चेदं मतं तात धर्मोऽयमिति राघव}
{धर्ममूलं त्वया छिन्नं राज्यमुत्सृजता तदा} %||6-70-30||

\twolineshloka
{अर्थेभ्यो हि विवृद्धेभ्यः संवृद्धेभ्यस्ततस्ततः}
{क्रियाः सर्वाः प्रवर्तन्ते पर्वतेभ्य इवापगाः} %||6-70-31||

\twolineshloka
{अर्थेन हि वियुक्तस्य पुरुषस्याल्पतेजसः}
{व्युच्छिद्यन्ते क्रियाः सर्वा ग्रीष्मे कुसरितो यथा} %||6-70-32||

\twolineshloka
{सोऽयमर्थं परित्यज्य सुखकामः सुखैधितः}
{पापमारभते कर्तुं तथा दोषः प्रवर्तते} %||6-70-33||

\twolineshloka
{यस्यार्थास्तस्य मित्राणि यस्यार्थास्तस्य बान्धवः}
{यस्यार्थाः स पुमाँल्लोके यस्यार्थाः स च पण्डितः} %||6-70-34||

\twolineshloka
{यस्यार्थाः स च विक्रान्तो यस्यार्थाः स च बुद्धिमान्}
{यस्यार्थाः स महाभागो यस्यार्थाः स महागुणः} %||6-70-35||

\twolineshloka
{अर्थस्यैते परित्यागे दोषाः प्रव्याहृता मया}
{राज्यमुत्सृजता वीर येन बुद्धिस्त्वया कृता} %||6-70-36||

\twolineshloka
{यस्यार्था धर्मकामार्थास्तस्य सर्वं प्रदक्षिणम्}
{अधनेनार्थकामेन नार्थः शक्यो विचिन्वता} %||6-70-37||

\twolineshloka
{हर्षः कामश्च दर्पश्च धर्मः क्रोधः शमो दमः}
{अर्थादेतानि सर्वाणि प्रवर्तन्ते नराधिप} %||6-70-38||

\twolineshloka
{येषां नश्यत्ययं लोकश्चरतां धर्मचारिणाम्}
{तेऽर्थास्त्वयि न दृश्यन्ते दुर्दिनेषु यथा ग्रहाः} %||6-70-39||

\twolineshloka
{त्वयि प्रव्रजिते वीर गुरोश्च वचने स्थिते}
{रक्षसापहृता भार्या प्राणैः प्रियतरा तव} %||6-70-40||

\twolineshloka
{तदद्य विपुलं वीर दुःखमिन्द्रजिता कृतम्}
{कर्मणा व्यपनेष्यामि तस्मादुत्तिष्ठ राघव} %||6-70-41||

\fourlineindentedshloka
{अयमनघ तवोदितः प्रियार्थम्}
{जनकसुता निधनं निरीक्ष्य रुष्टः}
{सहयगजरथां सराक्षसेन्द्राम्}
{भृशमिषुभिर्विनिपातयामि लङ्काम्} %||6-71-42||


{॥इत्यार्षे श्रीमद्रामायणे वाल्मीकीये आदिकाव्ये युद्धकाण्डे सप्ततितमः सर्गः॥}

\sect{एकसप्ततितमः सर्गः}

\twolineshloka
{राममाश्वासयाने तु लक्ष्मणे भ्रातृवत्सले}
{निक्षिप्य गुल्मान्स्वस्थाने तत्रागच्छद्विभीषणः} %||6-71-1||

\twolineshloka
{नानाप्रहरणैर्वीरैश्चतुर्भिः सचिवैर्वृतः}
{नीलाञ्जनचयाकारैर्मातङ्गैरिव यूथपः} %||6-71-2||

\twolineshloka
{सोऽभिगम्य महात्मानं राघवं शोकलालसम्}
{वानरांश्चैव ददृशे बाष्पपर्याकुलेक्षणान्} %||6-71-3||

\twolineshloka
{राघवं च महात्मानमिक्ष्वाकुकुलनन्दनम्}
{ददर्श मोहमापन्नं लक्ष्मणस्याङ्कमाश्रितम्} %||6-71-4||

\twolineshloka
{व्रीडितं शोकसन्तप्तं दृष्ट्वा रामं विभीषणः}
{अन्तर्दुःखेन दीनात्मा किमेतदिति सोऽब्रवीत्} %||6-71-5||

\twolineshloka
{विभीषण मुखं दृष्ट्वा सुग्रीवं तांश्च वानरान्}
{उवाच लक्ष्मणो वाक्यमिदं बाष्पपरिप्लुतः} %||6-71-6||

\twolineshloka
{हतामिन्द्रजिता सीतामिह श्रुत्वैव राघवः}
{हनूमद्वचनात्सौम्य ततो मोहमुपागतः} %||6-71-7||

\twolineshloka
{कथयन्तं तु सौमित्रिं संनिवार्य विभीषणः}
{पुष्कलार्थमिदं वाक्यं विसंज्ञं राममब्रवीत्} %||6-71-8||

\twolineshloka
{मनुजेन्द्रार्तरूपेण यदुक्तस्त्वं हनूमता}
{तदयुक्तमहं मन्ये सागरस्येव शोषणम्} %||6-71-9||

\twolineshloka
{अभिप्रायं तु जानामि रावणस्य दुरात्मनः}
{सीतां प्रति महाबाहो न च घातं करिष्यति} %||6-71-10||

\twolineshloka
{याच्यमानः सुबहुशो मया हितचिकीर्षुणा}
{वैदेहीमुत्सृजस्वेति न च तत्कृतवान्वचः} %||6-71-11||

\twolineshloka
{नैव साम्ना न भेदेन न दानेन कुतो युधा}
{सा द्रष्टुमपि शक्येत नैव चान्येन केनचित्} %||6-71-12||

\twolineshloka
{वानरान्मोहयित्वा तु प्रतियातः स राक्षसः}
{चैत्यं निकुम्भिलां नाम यत्र होमं करिष्यति} %||6-71-13||

\twolineshloka
{हुतवानुपयातो हि देवैरपि सवासवैः}
{दुराधर्षो भवत्येष सङ्ग्रामे रावणात्मजः} %||6-71-14||

\threelineshloka
{तेन मोहयता नूनमेषा माया प्रयोजिता}
{विघ्नमन्विच्छता तात वानराणां पराक्रमे}
{ससैन्यास्तत्र गच्छामो यावत्तन्न समाप्यते} %||6-71-15||

\twolineshloka
{त्यजेमं नरशार्दूलमिथ्या सन्तापमागतम्}
{सीदते हि बलं सर्वं दृष्ट्वा त्वां शोककर्शितम्} %||6-71-16||

\twolineshloka
{इह त्वं स्वस्थ हृदयस्तिष्ठ सत्त्वसमुच्छ्रितः}
{लक्ष्मणं प्रेषयास्माभिः सह सैन्यानुकर्षिभिः} %||6-71-17||

\twolineshloka
{एष तं नरशार्दूलो रावणिं निशितैः शरैः}
{त्याजयिष्यति तत्कर्म ततो वध्यो भविष्यति} %||6-71-18||

\twolineshloka
{तस्यैते निशितास्तीक्ष्णाः पत्रिपत्राङ्गवाजिनः}
{पतत्रिण इवासौम्याः शराः पास्यन्ति शोणितम्} %||6-71-19||

\twolineshloka
{तत्सन्दिश महाबाहो लक्ष्मणं शुभलक्षणम्}
{राक्षसस्य विनाशाय वज्रं वज्रधरो यथा} %||6-71-20||

\fourlineindentedshloka
{मनुजवर न कालविप्रकर्षो}
{रिपुनिधनं प्रति यत्क्षमोऽद्य कर्तुम्}
{त्वमतिसृज रिपोर्वधाय बाणीम्}
{असुरपुरोन्मथने यथा महेन्द्रः} %||6-71-21||

\fourlineindentedshloka
{समाप्तकर्मा हि स राक्षसेन्द्रो}
{भवत्यदृश्यः समरे सुरासुरैः}
{युयुत्सता तेन समाप्तकर्मणा}
{भवेत्सुराणामपि संशयो महान्} %||6-72-22||


{॥इत्यार्षे श्रीमद्रामायणे वाल्मीकीये आदिकाव्ये युद्धकाण्डे एकसप्ततितमः सर्गः॥}

\sect{द्विसप्ततितमः सर्गः}

\twolineshloka
{तस्य तद्वचनं श्रुत्वा राघवः शोककर्शितः}
{नोपधारयते व्यक्तं यदुक्तं तेन रक्षसा} %||6-72-1||

\twolineshloka
{ततो धैर्यमवष्टभ्य रामः परपुरंजयः}
{विभीषणमुपासीनमुवाच कपिसंनिधौ} %||6-72-2||

\twolineshloka
{नैरृताधिपते वाक्यं यदुक्तं ते विभीषण}
{भूयस्तच्छ्रोतुमिच्छामि ब्रूहि यत्ते विवक्षितम्} %||6-72-3||

\twolineshloka
{राघवस्य वचः श्रुत्वा वाक्यं वाक्यविशारदः}
{यत्तत्पुनरिदं वाक्यं बभाषे स विभीषणः} %||6-72-4||

\twolineshloka
{यथाज्ञप्तं महाबाहो त्वया गुल्मनिवेशनम्}
{तत्तथानुष्ठितं वीर त्वद्वाक्यसमनन्तरम्} %||6-72-5||

\twolineshloka
{तान्यनीकानि सर्वाणि विभक्तानि समन्ततः}
{विन्यस्ता यूथपाश्चैव यथान्यायं विभागशः} %||6-72-6||

\twolineshloka
{भूयस्तु मम विजाप्यं तच्छृणुष्व महायशः}
{त्वय्यकारणसन्तप्ते सन्तप्तहृदया वयम्} %||6-72-7||

\twolineshloka
{त्यज राजन्निमं शोकं मिथ्या सन्तापमागतम्}
{तदियं त्यज्यतां चिन्ता शत्रुहर्षविवर्धनी} %||6-72-8||

\twolineshloka
{उद्यमः क्रियतां वीर हर्षः समुपसेव्यताम्}
{प्राप्तव्या यदि ते सीता हन्तव्यश्व्च निशाचराः} %||6-72-9||

\threelineshloka
{रघुनन्दन वक्ष्यामि श्रूयतां मे हितं वचः}
{साध्वयं यातु सौमित्रिर्बलेन महता वृतः}
{निकुम्भिलायां सम्प्राप्य हन्तुं रावणिमाहवे} %||6-72-10||

\twolineshloka
{धनुर्मण्डलनिर्मुक्तैराशीविषविषोपमैः}
{शरैर्हन्तुं महेष्वासो रावणिं समितिंजयः} %||6-72-11||

\twolineshloka
{तेन वीरेण तपसा वरदानात्स्वयम्भुतः}
{अस्त्रं ब्रह्मशिरः प्राप्तं कामगाश्च तुरङ्गमाः} %||6-72-12||

\threelineshloka
{निकुम्भिलामसम्प्राप्तमहुताग्निं च यो रिपुः}
{त्वामाततायिनं हन्यादिन्द्रशत्रो स ते वधः}
{इत्येवं विहितो राजन्वधस्तस्यैव धीमतः} %||6-72-13||

\twolineshloka
{वधायेन्द्रजितो राम तं दिशस्व महाबलम्}
{हते तस्मिन्हतं विद्धि रावणं ससुहृज्जनम्} %||6-72-14||

\twolineshloka
{विभीषणवचः श्रुत्व रामो वाक्यमथाब्रवीत्}
{जानामि तस्य रौद्रस्य मायां सत्यपराक्रम} %||6-72-15||

\twolineshloka
{स हि ब्रह्मास्त्रवित्प्राज्ञो महामायो महाबलः}
{करोत्यसंज्ञान्सङ्ग्रामे देवान्सवरुणानपि} %||6-72-16||

\twolineshloka
{तस्यान्तरिक्षे चरतो रथस्थस्य महायशः}
{न गतिर्ज्ञायते वीरसूर्यस्येवाभ्रसम्प्लवे} %||6-72-17||

\twolineshloka
{राघवस्तु रिपोर्ज्ञात्वा मायावीर्यं दुरात्मनः}
{लक्ष्मणं कीर्तिसम्पन्नमिदं वचनमब्रवीत्} %||6-72-18||

\twolineshloka
{यद्वानरेन्द्रस्य बलं तेन सर्वेण संवृतः}
{हनूमत्प्रमुखैश्चैव यूथपैः सहलक्ष्मण} %||6-72-19||

\twolineshloka
{जाम्बवेनर्क्षपतिना सह सैन्येन संवृतः}
{जहि तं राक्षससुतं मायाबलविशारदम्} %||6-72-20||

\twolineshloka
{अयं त्वां सचिवैः सार्धं महात्मा रजनीचरः}
{अभिज्ञस्तस्य देशस्य पृष्ठतोऽनुगमिष्यति} %||6-72-21||

\twolineshloka
{राघवस्य वचः श्रुत्वा लक्ष्मणः सविभीषणः}
{जग्राह कार्मुकं श्रेष्ठमन्यद्भीमपराक्रमः} %||6-72-22||

\twolineshloka
{संनद्धः कवची खड्गी स शरी हेमचापधृक्}
{रामपादावुपस्पृश्य हृष्टः सौमित्रिरब्रवीत्} %||6-72-23||

\twolineshloka
{अद्य मत्कार्मुकोन्मुखाः शरा निर्भिद्य रावणिम्}
{लङ्कामभिपतिष्यन्ति हंसाः पुष्करिणीमिव} %||6-72-24||

\twolineshloka
{अद्यैव तस्य रौद्रस्य शरीरं मामकाः शराः}
{विधमिष्यन्ति हत्वा तं महाचापगुणच्युताः} %||6-72-25||

\twolineshloka
{स एवमुक्त्वा द्युतिमान्वचनं भ्रातुरग्रतः}
{स रावणिवधाकाङ्क्षी लक्ष्मणस्त्वरितो ययौ} %||6-72-26||

\twolineshloka
{सोऽभिवाद्य गुरोः पादौ कृत्वा चापि प्रदक्षिणम्}
{निकुम्भिलामभिययौ चैत्यं रावणिपालितम्} %||6-72-27||

\twolineshloka
{विभीषणेन सहितो राजपुत्रः प्रतापवान्}
{कृतस्वस्त्ययनो भ्रात्रा लक्ष्मणस्त्वरितो ययौ} %||6-72-28||

\twolineshloka
{वानराणां सहस्रैस्तु हनूमान्बहुभिर्वृतः}
{विभीषणः सहामात्यस्तदा लक्ष्मणमन्वगात्} %||6-72-29||

\twolineshloka
{महता हरिसैन्येन सवेगमभिसंवृतः}
{ऋक्षराजबलं चैव ददर्श पथि विष्ठितम्} %||6-72-30||

\twolineshloka
{स गत्वा दूरमध्वानं सौमित्रिर्मित्रनन्दनः}
{राक्षसेन्द्रबलं दूरादपश्यद्व्यूहमास्थितम्} %||6-72-31||

\twolineshloka
{स सम्प्राप्य धनुष्पाणिर्मायायोगमरिन्दम}
{तस्थौ ब्रह्मविधानेन विजेतुं रघुनन्दनः} %||6-72-32||

\fourlineindentedshloka
{विविधममलशस्त्रभास्वरं तद्-}
{ध्वजगहनं विपुलं महारथैश्च}
{प्रतिभयतममप्रमेयवेगम्}
{तिमिरमिव द्विषतां बलं विवेश} %||6-73-33||


{॥इत्यार्षे श्रीमद्रामायणे वाल्मीकीये आदिकाव्ये युद्धकाण्डे द्विसप्ततितमः सर्गः॥}

\sect{त्रिसप्ततितमः सर्गः}

\twolineshloka
{अथ तस्यामवस्थायां लक्ष्मणं रावणानुजः}
{परेषामहितं वाक्यमर्थसाधकमब्रवीत्} %||6-73-1||

\twolineshloka
{अस्यानीकस्य महतो भेदने यतलक्ष्मण}
{राक्षसेन्द्रसुतोऽप्यत्र भिन्ने दृश्यो भविष्यति} %||6-73-2||

\twolineshloka
{स त्वमिन्द्राशनिप्रख्यैः शरैरवकिरन्परान्}
{अभिद्रवाशु यावद्वै नैतत्कर्म समाप्यते} %||6-73-3||

\twolineshloka
{जहि वीरदुरात्मानं मायापरमधार्मिकम्}
{रावणिं क्रूरकर्माणं सर्वलोकभयावहम्} %||6-73-4||

\twolineshloka
{विभीषणवचः श्रुत्वा लक्ष्मणः शुभलक्षणः}
{ववर्ष शरवर्षाणि राक्षसेन्द्रसुतं प्रति} %||6-73-5||

\twolineshloka
{ऋक्षाः शाखामृगाश्चैव द्रुमाद्रिवरयोधिनः}
{अभ्यधावन्त सहितास्तदनीकमवस्थितम्} %||6-73-6||

\twolineshloka
{राक्षसाश्च शितैर्बाणैरसिभिः शक्तितोमरैः}
{उद्यतैः समवर्तन्त कपिसैन्यजिघांसवः} %||6-73-7||

\twolineshloka
{स सम्प्रहारस्तुमुलः संजज्ञे कपिरक्षसाम्}
{शब्देन महता लङ्कां नादयन्वै समन्ततः} %||6-73-8||

\twolineshloka
{शस्त्रैर्बहुविधाकारैः शितैर्बाणैश्च पादपैः}
{उद्यतैर्गिरिशृङ्गैश्च घोरैराकाशमावृतम्} %||6-73-9||

\twolineshloka
{ते राक्षसा वानरेषु विकृताननबाहवः}
{निवेशयन्तः शस्त्राणि चक्रुस्ते सुमहद्भयम्} %||6-73-10||

\twolineshloka
{तथैव सकलैर्वृक्षैर्गिरिशृङ्गैश्च वानराः}
{अभिजघ्नुर्निजघ्नुश्च समरे राक्षसर्षभान्} %||6-73-11||

\twolineshloka
{ऋक्षवानरमुख्यैश्च महाकायैर्महाबलैः}
{रक्षसां वध्यमानानां महद्भयमजायत} %||6-73-12||

\twolineshloka
{स्वमनीकं विषण्णं तु श्रुत्वा शत्रुभिरर्दितम्}
{उदतिष्ठत दुर्धर्षस्तत्कर्मण्यननुष्ठिते} %||6-73-13||

\twolineshloka
{वृक्षान्धकारान्निष्क्रम्य जातक्रोधः स रावणिः}
{आरुरोह रथं सज्जं पूर्वयुक्तं स राक्षसः} %||6-73-14||

\twolineshloka
{स भीमकार्मुकशरः कृष्णाञ्जनचयोपमः}
{रक्तास्यनयनः क्रूरो बभौ मृत्युरिवान्तकः} %||6-73-15||

\twolineshloka
{दृष्ट्वैव तु रथस्थं तं पर्यवर्तत तद्बलम्}
{रक्षसां भीमवेगानां लक्ष्मणेन युयुत्सताम्} %||6-73-16||

\twolineshloka
{तस्मिन्काले तु हनुमानुद्यम्य सुदुरासदम्}
{धरणीधरसङ्काशी महावृक्षमरिन्दमः} %||6-73-17||

\twolineshloka
{स राक्षसानां तत्सैन्यं कालाग्निरिव निर्दहन्}
{चकार बहुभिर्वृक्षैर्निःसंज्ञं युधि वानरः} %||6-73-18||

\twolineshloka
{विध्वंसयन्तं तरसा दृष्ट्वैव पवनात्मजम्}
{राक्षसानां सहस्राणि हनूमन्तमवाकिरन्} %||6-73-19||

\twolineshloka
{शितशूलधराः शूलैरसिभिश्चासिपाणयः}
{शक्तिभिः शक्तिहस्ताश्च पट्टसैः पट्टसायुधाः} %||6-73-20||

\twolineshloka
{परिघैश्च गदाभिश्च कुन्तैश्च शुभदर्शनैः}
{शतशश्च शतघ्नीभिरायसैरपि मुद्गरैः} %||6-73-21||

\twolineshloka
{घोरैः परशुभिश्चैव भिण्डिपालैश्च राक्षसाः}
{मुष्टिभिर्वज्रवेगैश्च तलैरशनिसंनिभैः} %||6-73-22||

\twolineshloka
{अभिजघ्नुः समासाद्य समन्तात्पर्वतोपमम्}
{तेषामपि च सङ्क्रुद्धश्चकार कदनं महत्} %||6-73-23||

\twolineshloka
{स ददर्श कपिश्रेष्ठमचलोपममिन्द्रजित्}
{सूदयानममित्रघ्नममित्रान्पवनात्मजम्} %||6-73-24||

\twolineshloka
{स सारथिमुवाचेदं याहि यत्रैष वानरः}
{क्षयमेव हि नः कुर्याद्राक्षसानामुपेक्षितः} %||6-73-25||

\twolineshloka
{इत्युक्तः सारथिस्तेन ययौ यत्र स मारुतिः}
{वहन्परमदुर्धर्षं स्थितमिन्द्रजितं रथे} %||6-73-26||

\twolineshloka
{सोऽभ्युपेत्य शरान्खड्गान्पट्टसासिपरश्वधान्}
{अभ्यवर्षत दुर्धर्षः कपिमूर्ध्नि स राक्षसः} %||6-73-27||

\twolineshloka
{तानि शस्त्राणि घोराणि प्रतिगृह्य स मारुतिः}
{रोषेण महताविषो वाक्यं चेदमुवाच ह} %||6-73-28||

\twolineshloka
{युध्यस्व यदि शूरोऽसि रावणात्मज दुर्मते}
{वायुपुत्रं समासाद्य न जीवन्प्रतियास्यसि} %||6-73-29||

\twolineshloka
{बाहुभ्यां सम्प्रयुध्यस्व यदि मे द्वन्द्वमाहवे}
{वेगं सहस्व दुर्बुद्धे ततस्त्वं रक्षसां वरः} %||6-73-30||

\twolineshloka
{हनूमन्तं जिघांसन्तं समुद्यतशरासनम्}
{रावणात्मजमाचष्टे लक्ष्मणाय विभीषणः} %||6-73-31||

\twolineshloka
{यस्तु वासवनिर्जेता रावणस्यात्मसम्भवः}
{स एष रथमास्थाय हनूमन्तं जिघांसति} %||6-73-32||

\twolineshloka
{तमप्रतिमसंस्थानैः शरैः शत्रुविदारणैः}
{जीवितान्तकरैर्घोरैः सौमित्रे रावणिं जहि} %||6-73-33||

\fourlineindentedshloka
{इत्येवमुक्तस्तु तदा महात्मा}
{विभीषणेनारिविभीषणेन}
{ददर्श तं पर्वतसंनिकाशम्}
{रथस्थितं भीमबलं दुरासदम्} %||6-74-34||


{॥इत्यार्षे श्रीमद्रामायणे वाल्मीकीये आदिकाव्ये युद्धकाण्डे त्रिसप्ततितमः सर्गः॥}

\sect{चतुःसप्ततितमः सर्गः}

\twolineshloka
{एवमुक्त्वा तु सौमित्रिं जातहर्षो विभीषणः}
{धनुष्पाणिनमादाय त्वरमाणो जगाम सः} %||6-74-1||

\twolineshloka
{अविदूरं ततो गत्वा प्रविश्य च महद्वनम्}
{दर्शयामास तत्कर्म लक्ष्मणाय विभीषणः} %||6-74-2||

\twolineshloka
{नीलजीमूतसङ्काशं न्यग्रोधं भीमदर्शनम्}
{तेजस्वी रावणभ्राता लक्ष्मणाय न्यवेदयत्} %||6-74-3||

\twolineshloka
{इहोपहारं भूतानां बलवान्रावणातजः}
{उपहृत्य ततः पश्चात्सङ्ग्राममभिवर्तते} %||6-74-4||

\twolineshloka
{अदृश्यः सर्वभूतानां ततो भवति राक्षसः}
{निहन्ति समरे शत्रून्बध्नाति च शरोत्तमैः} %||6-74-5||

\twolineshloka
{तमप्रविष्टं न्यग्रोधं बलिनं रावणात्मजम्}
{विध्वंसय शरैस्तीक्ष्णैः सरथं साश्वसारथिम्} %||6-74-6||

\twolineshloka
{तथेत्युक्त्वा महातेजाः सौमित्रिर्मित्रनन्दनः}
{बभूवावस्थितस्तत्र चित्रं विस्फारयन्धनुः} %||6-74-7||

\twolineshloka
{स रथेनाग्निवर्णेन बलवान्रावणात्मजः}
{इन्द्रजित्कवची खड्गी सध्वजः प्रत्यदृश्यत} %||6-74-8||

\twolineshloka
{तमुवाच महातेजाः पौलस्त्यमपराजितम्}
{समाह्वये त्वां समरे सम्यग्युद्धं प्रयच्छ मे} %||6-74-9||

\twolineshloka
{एवमुक्तो महातेजा मनस्वी रावणात्मजः}
{अब्रवीत्परुषं वाक्यं तत्र दृष्ट्वा विभीषणम्} %||6-74-10||

\twolineshloka
{इह त्वं जातसंवृद्धः साक्षाद्भ्राता पितुर्मम}
{कथं द्रुह्यसि पुत्रस्य पितृव्यो मम राक्षस} %||6-74-11||

\twolineshloka
{न ज्ञातित्वं न सौहार्दं न जातिस्तव दुर्मते}
{प्रमाणं न च सोदर्यं न धर्मो धर्मदूषण} %||6-74-12||

\twolineshloka
{शोच्यस्त्वमसि दुर्बुद्धे निन्दनीयश्च साधुभिः}
{यस्त्वं स्वजनमुत्सृज्य परभृत्यत्वमागतः} %||6-74-13||

\twolineshloka
{नैतच्छिथिलया बुद्ध्या त्वं वेत्सि महदन्तरम्}
{क्व च स्वजनसंवासः क्व च नीचपराश्रयः} %||6-74-14||

\twolineshloka
{गुणवान्वा परजनः स्वजनो निर्गुणोऽपि वा}
{निर्गुणः स्वजनः श्रेयान्यः परः पर एव सः} %||6-74-15||

\twolineshloka
{निरनुक्रोशता चेयं यादृशी ते निशाचर}
{स्वजनेन त्वया शक्यं परुषं रावणानुज} %||6-74-16||

\twolineshloka
{इत्युक्तो भ्रातृपुत्रेण प्रत्युवाच विभीषणः}
{अजानन्निव मच्छीलं किं राक्षस विकत्थसे} %||6-74-17||

\threelineshloka
{राक्षसेन्द्रसुतासाधो पारुष्यं त्यज गौरवात्}
{कुले यद्यप्यहं जातो रक्षसां क्रूरकर्मणाम्}
{गुणोऽयं प्रथमो नॄणां तन्मे शीलमराक्षसम्} %||6-74-18||

\twolineshloka
{न रमे दारुणेनाहं न चाधर्मेण वै रमे}
{भ्रात्रा विषमशीलेन कथं भ्राता निरस्यते} %||6-74-19||

\twolineshloka
{परस्वानां च हरणं परदाराभिमर्शनम्}
{सुहृदामतिशङ्कां च त्रयो दोषाः क्षयावहाः} %||6-74-20||

\twolineshloka
{महर्षीणां वधो घोरः सर्वदेवैश्च विग्रहः}
{अभिमानश्च कोपश्च वैरित्वं प्रतिकूलता} %||6-74-21||

\twolineshloka
{एते दोषा मम भ्रातुर्जीवितैश्वर्यनाशनाः}
{गुणान्प्रच्छादयामासुः पर्वतानिव तोयदाः} %||6-74-22||

\twolineshloka
{दोषैरेतैः परित्यक्तो मया भ्राता पिता तव}
{नेयमस्ति पुरी लङ्का न च त्वं न च ते पिता} %||6-74-23||

\twolineshloka
{अतिमानी च बालश्च दुर्विनीतश्च राक्षस}
{बद्धस्त्वं कालपाशेन ब्रूहि मां यद्यदिच्छसि} %||6-74-24||

\twolineshloka
{अद्य ते व्यसनं प्राप्तं किमिह त्वं तु वक्ष्यसि}
{प्रवेष्टुं न त्वया शक्यो न्यग्रोधो राक्षसाधम} %||6-74-25||

\threelineshloka
{धर्षयित्वा तु काकुत्स्थौ न शक्यं जीवितुं त्वया}
{युध्यस्व नरदेवेन लक्ष्मणेन रणे सह}
{हतस्त्वं देवता कार्यं करिष्यसि यमक्षये} %||6-74-26||

\fourlineindentedshloka
{निदर्शयस्वात्मबलं समुद्यतम्}
{कुरुष्व सर्वायुधसायकव्ययम्}
{न लक्ष्मणस्यैत्य हि बाणगोचरम्}
{त्वमद्य जीवन्सबलो गमिष्यसि} %||6-75-27||


{॥इत्यार्षे श्रीमद्रामायणे वाल्मीकीये आदिकाव्ये युद्धकाण्डे चतुःसप्ततितमः सर्गः॥}

\sect{पञ्चसप्ततितमः सर्गः}

\twolineshloka
{विभीषण वचः श्रुत्वा रावणिः क्रोधमूर्छितः}
{अब्रवीत्परुषं वाक्यं वेगेनाभ्युत्पपात ह} %||6-75-1||

\twolineshloka
{उद्यतायुधनिस्त्रिंशो रथे तु समलङ्कृते}
{कालाश्वयुक्ते महति स्थितः कालान्तकोपमः} %||6-75-2||

\twolineshloka
{महाप्रमाणमुद्यम्य विपुलं वेगवद्दृढम्}
{धनुर्भीमं परामृश्य शरांश्चामित्रनाशनान्} %||6-75-3||

\twolineshloka
{उवाचैनं समारब्धः सौमित्रिं सविभीषणम्}
{तांश्च वानरशार्दूलान्पश्यध्वं मे पराक्रमम्} %||6-75-4||

\twolineshloka
{अद्य मत्कार्मुकोत्सृष्टं शरवर्षं दुरासदम्}
{मुक्तं वर्षमिवाकाशे वारयिष्यथ संयुगे} %||6-75-5||

\twolineshloka
{अद्य वो मामका बाणा महाकार्मुकनिःसृताः}
{विधमिष्यन्ति गात्राणि तूलराशिमिवानलः} %||6-75-6||

\twolineshloka
{तीक्ष्णसायकनिर्भिन्नाञ्शूलशक्त्यृष्टितोमरैः}
{अद्य वो गमयिष्यामि सर्वानेव यमक्षयम्} %||6-75-7||

\twolineshloka
{क्षिपतः शरवर्षाणि क्षिप्रहस्तस्य मे युधि}
{जीमूतस्येव नदतः कः स्थास्यति ममाग्रतः} %||6-75-8||

\twolineshloka
{तच्छ्रुत्वा राक्षसेन्द्रस्य गर्जितं लक्ष्मणस्तदा}
{अभीतवदनः क्रुद्धो रावणिं वाक्यमब्रवीत्} %||6-75-9||

\twolineshloka
{उक्तश्च दुर्गमः पारः कार्याणां राक्षस त्वया}
{कार्याणां कर्मणा पारं यो गच्छति स बुद्धिमान्} %||6-75-10||

\twolineshloka
{स त्वमर्थस्य हीनार्थो दुरवापस्य केनचित्}
{वचो व्याहृत्य जानीषे कृतार्थोऽस्मीति दुर्मते} %||6-75-11||

\twolineshloka
{अन्तर्धानगतेनाजौ यस्त्वयाचरितस्तदा}
{तस्कराचरितो मार्गो नैष वीरनिषेवितः} %||6-75-12||

\twolineshloka
{यथा बाणपथं प्राप्य स्थितोऽहं तव राक्षस}
{दर्शयस्वाद्य तत्तेजो वाचा त्वं किं विकत्थसे} %||6-75-13||

\twolineshloka
{एवमुक्तो धनुर्भीमं परामृश्य महाबलः}
{ससर्जे निशितान्बाणानिन्द्रजित्समिजिंजय} %||6-75-14||

\twolineshloka
{ते निसृष्टा महावेगाः शराः सर्पविषोपमाः}
{सम्प्राप्य लक्ष्मणं पेतुः श्वसन्त इव पन्नगाः} %||6-75-15||

\twolineshloka
{शरैरतिमहावेगैर्वेगवान्रावणात्मजः}
{सौमित्रिमिन्द्रजिद्युद्धे विव्याध शुभलक्षणम्} %||6-75-16||

\twolineshloka
{स शरैरतिविद्धाङ्गो रुधिरेण समुक्षितः}
{शुशुभे लक्ष्मणः श्रीमान्विधूम इव पावकः} %||6-75-17||

\twolineshloka
{इन्द्रजित्त्वात्मनः कर्म प्रसमीक्ष्याधिगम्य च}
{विनद्य सुमहानादमिदं वचनमब्रवीत्} %||6-75-18||

\twolineshloka
{पत्रिणः शितधारास्ते शरा मत्कार्मुकच्युताः}
{आदास्यन्तेऽद्य सौमित्रे जीवितं जीवितान्तगाः} %||6-75-19||

\twolineshloka
{अद्य गोमायुसङ्घाश्च श्येनसङ्घाश्च लक्ष्मण}
{गृध्राश्च निपतन्तु त्वां गतासुं निहतं मया} %||6-75-20||

\twolineshloka
{क्षत्रबन्धुः सदानार्यो रामः परमदुर्मतिः}
{भक्तं भ्रातरमद्यैव त्वां द्रक्ष्यति मया हतम्} %||6-75-21||

\twolineshloka
{विशस्तकवचं भूमौ व्यपविद्धशरासनम्}
{हृतोत्तमाङ्गं सौमित्रे त्वामद्य निहतं मया} %||6-75-22||

\twolineshloka
{इति ब्रुवाणं संरब्धं परुषं रावणात्मजम्}
{हेतुमद्वाक्यमत्यर्थं लक्ष्मणः प्रत्युवाच ह} %||6-75-23||

\twolineshloka
{अकृत्वा कत्थसे कर्म किमर्थमिह राक्षस}
{कुरु तत्कर्म येनाहं श्रद्दध्यां तव कत्थनम्} %||6-75-24||

\twolineshloka
{अनुक्त्वा परुषं वाक्यं किञ्चिदप्यनवक्षिपन्}
{अविकत्थन्वधिष्यामि त्वां पश्य पुरुषादन} %||6-75-25||

\twolineshloka
{इत्युक्त्वा पञ्चनाराचानाकर्णापूरिताञ्शरान्}
{निचखान महावेगाँल्लक्ष्मणो राक्षसोरसि} %||6-75-26||

\twolineshloka
{स शरैराहतस्तेन सरोषो रावणात्मजः}
{सुप्रयुक्तैस्त्रिभिर्बाणैः प्रतिविव्याध लक्ष्मणम्} %||6-75-27||

\twolineshloka
{स बभूव महाभीमो नरराक्षससिंहयोः}
{विमर्दस्तुमुलो युद्धे परस्परवधैषिणोः} %||6-75-28||

\twolineshloka
{उभौ हि बलसम्पन्नावुभौ विक्रमशालिनौ}
{उभावपि सुविक्रान्तौ सर्वशस्त्रास्त्रकोविदौ} %||6-75-29||

\twolineshloka
{उभौ परमदुर्जेयावतुल्यबलतेजसौ}
{युयुधाते महावीरौ ग्रहाविव नभो गतौ} %||6-75-30||

\twolineshloka
{बलवृत्राविव हि तौ युधि वै दुष्प्रधर्षणौ}
{युयुधाते महात्मानौ तदा केसरिणाविव} %||6-75-31||

\twolineshloka
{बहूनवसृजन्तौ हि मार्गणौघानवस्थितौ}
{नरराक्षससिंहौ तौ प्रहृष्टावभ्ययुध्यताम्} %||6-75-32||

\fourlineindentedshloka
{सुसम्प्रहृष्टौ नरराक्षसोत्तमौ}
{जयैषिणौ मार्गणचापधारिणौ}
{परस्परं तौ प्रववर्षतुर्भृशम्}
{शरौघवर्षेण बलाहकाविव} %||6-76-33||


{॥इत्यार्षे श्रीमद्रामायणे वाल्मीकीये आदिकाव्ये युद्धकाण्डे पञ्चसप्ततितमः सर्गः॥}

\sect{षट्सप्ततितमः सर्गः}

\twolineshloka
{ततः शरं दाशरथिः सन्धायामित्रकर्शनः}
{ससर्ज राक्षसेन्द्राय क्रुद्धः सर्प इव श्वसन्} %||6-76-1||

\twolineshloka
{तस्य ज्यातलनिर्घोषं स श्रुत्वा रावणात्मजः}
{विवर्णवदनो भूत्वा लक्ष्मणं समुदैक्षत} %||6-76-2||

\twolineshloka
{तं विषण्णमुखं दृष्ट्वा राक्षसं रावणात्मजम्}
{सौमित्रिं युद्धसंसक्तं प्रत्युवाच विभीषणः} %||6-76-3||

\twolineshloka
{निमित्तान्यनुपश्यामि यान्यस्मिन्रावणात्मजे}
{त्वर तेन महाबाहो भग्न एष न संशयः} %||6-76-4||

\twolineshloka
{ततः सन्धाय सौमित्रिः शरानग्निशिखोपमान्}
{मुमोच निशितांस्तस्मै सर्वानिव विषोल्बणान्} %||6-76-5||

\twolineshloka
{शक्राशनिसमस्पर्शैर्लक्ष्मणेनाहतः शरैः}
{मुहूर्तमभवन्मूढः सर्वसङ्क्षुभितेन्द्रियः} %||6-76-6||

\twolineshloka
{उपलभ्य मुहूर्तेन संज्ञां प्रत्यागतेन्द्रियः}
{ददर्शावस्थितं वीरं वीरो दशरथात्मजम्} %||6-76-7||

\twolineshloka
{सोऽभिचक्राम सौमित्रिं रोषात्संरक्तलोचनः}
{अब्रवीच्चैनमासाद्य पुनः स परुषं वचः} %||6-76-8||

\twolineshloka
{किं न स्मरसि तद्युद्धे प्रथमे मत्पराक्रमम्}
{निबद्धस्त्वं सह भ्रात्रा यदा युधि विचेष्टसे} %||6-76-9||

\twolineshloka
{युवा खलु महायुद्धे शक्राशनिसमैः शरैः}
{शायिनौ प्रथमं भूमौ विसंज्ञौ सपुरःसरौ} %||6-76-10||

\twolineshloka
{स्मृतिर्वा नास्ति ते मन्ये व्यक्तं वा यमसादनम्}
{गन्तुमिच्छसि यस्मात्त्वं मां धर्षयितुमिच्छसि} %||6-76-11||

\twolineshloka
{यदि ते प्रथमे युद्धे न दृष्टो मत्पराक्रमः}
{अद्य त्वां दर्शयिष्यामि तिष्ठेदानीं व्यवस्थितः} %||6-76-12||

\twolineshloka
{इत्युक्त्वा सप्तभिर्बाणैरभिविव्याध लक्ष्मणम्}
{दशभिश्च हनूमन्तं तीक्ष्णधारैः शरोत्तमैः} %||6-76-13||

\twolineshloka
{ततः शरशतेनैव सुप्रयुक्तेन वीर्यवान्}
{क्रोधाद्द्विगुणसंरब्धो निर्बिभेद विभीषणम्} %||6-76-14||

\twolineshloka
{तद्दृष्ट्वेन्द्रजितः कर्म कृतं रामानुजस्तदा}
{अचिन्तयित्वा प्रहसन्नैतत्किञ्चिदिति ब्रुवन्} %||6-76-15||

\twolineshloka
{मुमोच स शरान्घोरान्सङ्गृह्य नरपुङ्गवः}
{अभीतवदनः क्रुद्धो रावणिं लक्ष्मणो युधि} %||6-76-16||

\twolineshloka
{नैवं रणगतः शूराः प्रहरन्ति निशाचर}
{लघवश्चाल्पवीर्याश्च सुखा हीमे शरास्तव} %||6-76-17||

\twolineshloka
{नैवं शूरास्तु युध्यन्ते समरे जयकाङ्क्षिणः}
{इत्येवं तं ब्रुवाणस्तु शरवर्षैरवाकिरत्} %||6-76-18||

\twolineshloka
{तस्य बाणैस्तु विध्वस्तं कवचं हेमभूषितम्}
{व्यशीर्यत रथोपस्थे ताराजालमिवाम्बरात्} %||6-76-19||

\twolineshloka
{विधूतवर्मा नाराचैर्बभूव स कृतव्रणः}
{इन्द्रजित्समरे शूरः प्ररूढ इव सानुमान्} %||6-76-20||

\twolineshloka
{अभीक्ष्णं निश्वसन्तौ हि युध्येतां तुमुलं युधि}
{शरसङ्कृत्तसर्वाङ्गो सर्वतो रुधिरोक्षितौ} %||6-76-21||

\twolineshloka
{अस्त्राण्यस्त्रविदां श्रेष्ठौ दर्शयन्तौ पुनः पुनः}
{शरानुच्चावचाकारानन्तरिक्षे बबन्धतुः} %||6-76-22||

\twolineshloka
{व्यपेतदोषमस्यन्तौ लघुचित्रं च सुष्ठु च}
{उभौ तु तुमुलं घोरं चक्रतुर्नरराक्षसौ} %||6-76-23||

\twolineshloka
{तयोः पृथक्पृथग्भीमः शुश्रुवे तलनिस्वनः}
{सुघोरयोर्निष्टनतोर्गगने मेघयोरिव} %||6-76-24||

\twolineshloka
{ते गात्रयोर्निपतिता रुक्मपुङ्खाः शरा युधि}
{असृग्दिग्धा विनिष्पेतुर्विविशुर्धरणीतलम्} %||6-76-25||

\twolineshloka
{अन्यैः सुनिशितैः शस्त्रैराकाशे संजघट्टिरे}
{बभञ्जुश्चिच्छिदुश्चापि तयोर्बाणाः सहस्रशः} %||6-76-26||

\twolineshloka
{स बभूव रणे घोरस्तयोर्बाणमयश्चयः}
{अग्निभ्यामिव दीप्ताभ्यां सत्रे कुशमयश्चयः} %||6-76-27||

\twolineshloka
{तयोः कृतव्रणौ देहौ शुशुभाते महात्मनोः}
{सपुष्पाविव निष्पत्रौ वने शाल्मलिकुंशुकौ} %||6-76-28||

\twolineshloka
{चक्रतुस्तुमुलं घोरं संनिपातं मुहुर्मुहुः}
{इन्द्रजिल्लक्ष्मणश्चैव परस्परजयैषिणौ} %||6-76-29||

\twolineshloka
{लक्ष्मणो रावणिं युद्धे रावणिश्चापि लक्ष्मणम्}
{अन्योन्यं तावभिघ्नन्तौ न श्रमं प्रत्यपद्यताम्} %||6-76-30||

\twolineshloka
{बाणजालैः शरीरस्थैरवगाढैस्तरस्विनौ}
{शुशुभाते महावीरौ विरूढाविव पर्वतौ} %||6-76-31||

\twolineshloka
{तयो रुधिरसिक्तानि संवृतानि शरैर्भृशम्}
{बभ्राजुः सर्वगात्राणि ज्वलन्त इव पावकाः} %||6-76-32||

\twolineshloka
{तयोरथ महान्कालो व्यतीयाद्युध्यमानयोः}
{न च तौ युद्धवैमुख्यं श्रमं वाप्युपजग्मतुः} %||6-76-33||

\fourlineindentedshloka
{अथ समरपरिश्रमं निहन्तुम्}
{समरमुखेष्वजितस्य लक्ष्मणस्य}
{प्रियहितमुपपादयन्महौजाः}
{समरमुपेत्य विभीषणोऽवतस्थे} %||6-77-34||


{॥इत्यार्षे श्रीमद्रामायणे वाल्मीकीये आदिकाव्ये युद्धकाण्डे षट्सप्ततितमः सर्गः॥}

\sect{सप्तसप्ततितमः सर्गः}

\twolineshloka
{युध्यमानौ तु तौ दृष्ट्वा प्रसक्तौ नरराक्षसौ}
{शूरः स रावणभ्राता तस्थौ सङ्ग्राममूर्धनि} %||6-77-1||

\twolineshloka
{ततो विस्फारयामास महद्धनुरवस्थितः}
{उत्ससर्ज च तीक्ष्णाग्रान्राक्षसेषु महाशरान्} %||6-77-2||

\twolineshloka
{ते शराः शिखिसङ्काशा निपतन्तः समाहिताः}
{राक्षसान्दारयामासुर्वज्रा इव महागिरीन्} %||6-77-3||

\twolineshloka
{विभीषणस्यानुचरास्तेऽपि शूलासिपट्टसैः}
{चिच्छेदुः समरे वीरान्राक्षसान्राक्षसोत्तमाः} %||6-77-4||

\twolineshloka
{राक्षसैस्तैः परिवृतः स तदा तु विभीषणः}
{बभौ मध्ये प्रहृष्टानां कलभानामिव द्विपः} %||6-77-5||

\twolineshloka
{ततः सञ्चोदयानो वै हरीन्रक्षोरणप्रियान्}
{उवाच वचनं काले कालज्ञो रक्षसां वरः} %||6-77-6||

\twolineshloka
{एकोऽयं राक्षसेन्द्रस्य परायणमिव स्थितः}
{एतच्छेषं बलं तस्य किं तिष्ठत हरीश्वराः} %||6-77-7||

\twolineshloka
{अस्मिन्विनिहते पापे राक्षसे रणमूर्धनि}
{रावणं वर्जयित्वा तु शेषमस्य बलं हतम्} %||6-77-8||

\twolineshloka
{प्रहस्तो निहतो वीरो निकुम्भश्च महाबलः}
{कुम्भकर्णश्च कुम्भश्च धूम्राक्षश्च निशाचरः} %||6-77-9||

\twolineshloka
{अकम्पनः सुपार्श्वश्च चक्रमाली च राक्षसः}
{कम्पनः सत्त्ववन्तश्च देवान्तकनरान्तकौ} %||6-77-10||

\twolineshloka
{एतान्निहत्यातिबलान्बहून्राक्षससत्तमान्}
{बाहुभ्यां सागरं तीर्त्वा लङ्घ्यतां गोष्पदं लघु} %||6-77-11||

\twolineshloka
{एतावदिह शेषं वो जेतव्यमिह वानराः}
{हताः सर्वे समागम्य राक्षसा बलदर्पिताः} %||6-77-12||

\twolineshloka
{अयुक्तं निधनं कर्तुं पुत्रस्य जनितुर्मम}
{घृणामपास्य रामार्थे निहन्यां भ्रातुरात्मजम्} %||6-77-13||

\threelineshloka
{हन्तुकामस्य मे बाष्पं चक्शुश्चैव निरुध्यते}
{तदेवैष महाबाहुर्लक्ष्मणः शमयिष्यति}
{वानरा घ्नन्तुं सम्भूय भृत्यानस्य समीपगान्} %||6-77-14||

\twolineshloka
{इति तेनातियशसा राक्षसेनाभिचोदिताः}
{वानरेन्द्रा जहृषिरे लाङ्गलानि च विव्यधुः} %||6-77-15||

\twolineshloka
{ततस्ते कपिशार्दूलाः क्ष्वेडन्तश्च मुहुर्मुहुः}
{मुमुचुर्विविधान्नादान्मेघान्दृष्ट्वेव बर्हिणः} %||6-77-16||

\twolineshloka
{जाम्बवानपि तैः सर्वैः स्वयूथैरभिसंवृतः}
{अश्मभिस्ताडयामास नखैर्दन्तैश्च राक्षसान्} %||6-77-17||

\twolineshloka
{निघ्नन्तमृक्षाधिपतिं राक्षसास्ते महाबलाः}
{परिवव्रुर्भयं त्यक्त्वा तमनेकविधायुधाः} %||6-77-18||

\twolineshloka
{शरैः परशुभिस्तीक्ष्णैः पट्टसैर्यष्टितोमरैः}
{जाम्बवन्तं मृधे जघ्नुर्निघ्नन्तं राक्षसीं चमूम्} %||6-77-19||

\twolineshloka
{स सम्प्रहारस्तुमुलः संजज्ञे कपिराक्षसाम्}
{देवासुराणां क्रुद्धानां यथा भीमो महास्वनः} %||6-77-20||

\twolineshloka
{हनूमानपि सङ्क्रुद्धः सालमुत्पाट्य पर्वतात्}
{रक्षसां कदनं चक्रे समासाद्य सहस्रशः} %||6-77-21||

\twolineshloka
{स दत्त्वा तुमुलं युद्धं पितृव्यस्येन्द्रजिद्युधि}
{लक्ष्मणं परवीरघ्नं पुनरेवाभ्यधावत} %||6-77-22||

\twolineshloka
{तौ प्रयुद्धौ तदा वीरौ मृधे लक्ष्मणराक्षसौ}
{शरौघानभिवर्षन्तौ जघ्नतुस्तौ परस्परम्} %||6-77-23||

\twolineshloka
{अभीक्ष्णमन्तर्दधतुः शरजालैर्महाबलौ}
{चन्द्रादित्याविवोष्णान्ते यथा मेघैस्तरस्विनौ} %||6-77-24||

\twolineshloka
{न ह्यादानं न सन्धानं धनुषो वा परिग्रहः}
{न विप्रमोक्षो बाणानां न विकर्षो न विग्रहः} %||6-77-25||

\twolineshloka
{न मुष्टिप्रतिसन्धानं न लक्ष्यप्रतिपादनम्}
{अदृश्यत तयोस्तत्र युध्यतोः पाणिलाघवात्} %||6-77-26||

\threelineshloka
{चापवेगप्रमुक्तैश्च बाणजालैः समन्ततः}
{अन्तरिक्षेऽभिसञ्छन्ने न रूपाणि चकाशिरे}
{तमसा पिहितं सर्वमासीद्भीमतरं महत्} %||6-77-27||

\threelineshloka
{न तदानीं ववौ वायुर्न जज्वाल च पावकः}
{स्वस्त्यस्तु लोकेभ्य इति जजल्पश्च महर्षयः}
{सम्पेतुश्चात्र सम्प्राप्ता गन्धर्वाः सह चारणैः} %||6-77-28||

\twolineshloka
{अथ राक्षससिंहस्य कृष्णान्कनकभूषणान्}
{शरैश्चतुर्भिः सौमित्रिर्विव्याध चतुरो हयान्} %||6-77-29||

\twolineshloka
{ततोऽपरेण भल्लेन सूतस्य विचरिष्यतः}
{लाघवाद्राघवः श्रीमाञ्शिरः कायादपाहरत्} %||6-77-30||

\twolineshloka
{निहतं सारथिं दृष्ट्वा समरे रावणात्मजः}
{प्रजहौ समरोद्धर्षं विषण्णः स बभूव ह} %||6-77-31||

\twolineshloka
{विषण्णवदनं दृष्ट्वा राक्षसं हरियूथपाः}
{ततः परमसंहृष्टो लक्ष्मणं चाभ्यपूजयन्} %||6-77-32||

\twolineshloka
{ततः प्रमाथी शरभो रभसो गन्धमादनः}
{अमृष्यमाणाश्चत्वारश्चक्रुर्वेगं हरीश्वराः} %||6-77-33||

\twolineshloka
{ते चास्य हयमुख्येषु तूर्णमुत्पत्य वानराः}
{चतुर्षु सुमहावीर्या निपेतुर्भीमविक्रमाः} %||6-77-34||

\twolineshloka
{तेषामधिष्ठितानां तैर्वानरैः पर्वतोपमैः}
{मुखेभ्यो रुधिरं व्यक्तं हयानां समवर्तत} %||6-77-35||

\twolineshloka
{ते निहत्य हयांस्तस्य प्रमथ्य च महारथम्}
{पुनरुत्पत्य वेगेन तस्थुर्लक्ष्मणपार्श्वतः} %||6-77-36||

\twolineshloka
{स हताश्वादवप्लुत्य रथान्मथितसारथेः}
{शरवर्षेण सौमित्रिमभ्यधावत रावणिः} %||6-77-37||

\fourlineindentedshloka
{ततो महेन्द्रप्रतिमंह्स लक्ष्मणः}
{पदातिनं तं निशितैः शरोत्तमैः}
{सृजन्तमादौ निशिताञ्शरोत्तमान्}
{भृशं तदा बाणगणैर्न्यवारयत्} %||6-78-38||


{॥इत्यार्षे श्रीमद्रामायणे वाल्मीकीये आदिकाव्ये युद्धकाण्डे सप्तसप्ततितमः सर्गः॥}

\sect{अष्टसप्ततितमः सर्गः}

\twolineshloka
{स हताश्वो महातेजा भूमौ तिष्ठन्निशाचरः}
{इन्द्रजित्परमक्रुद्धः सम्प्रजज्वाल तेजसा} %||6-78-1||

\twolineshloka
{तौ धन्विनौ जिघांसन्तावन्योन्यमिषुभिर्भृशम्}
{विजयेनाभिनिष्क्रान्तौ वने गजवृषाविव} %||6-78-2||

\twolineshloka
{निबर्हयन्तश्चान्योन्यं ते राक्षसवनौकसः}
{भर्तारं न जहुर्युद्धे सम्पतन्तस्ततस्ततः} %||6-78-3||

\twolineshloka
{स लक्ष्मणं समुद्दिश्य परं लाघवमास्थितः}
{ववर्ष शरवर्षाणि वर्षाणीव पुरन्दरः} %||6-78-4||

\twolineshloka
{मुक्तमिन्द्रजिता तत्तु शरवर्षमरिन्दमः}
{अवारयदसम्भ्रान्तो लक्ष्मणः सुदुरासदम्} %||6-78-5||

\threelineshloka
{अभेद्यकचनं मत्वा लक्ष्मणं रावणात्मजः}
{ललाटे लक्ष्मणं बाणैः सुपुङ्खैस्त्रिभिरिन्द्रजित्}
{अविध्यत्परमक्रुद्धः शीघ्रमस्त्रं प्रदर्शयन्} %||6-78-6||

\twolineshloka
{तैः पृषत्कैर्ललाटस्थैः शुशुभे रघुनन्दनः}
{रणाग्रे समरश्लाघी त्रिशृङ्ग इव पर्वतः} %||6-78-7||

\twolineshloka
{स तथाप्यर्दितो बाणै राक्षसेन महामृधे}
{तमाशु प्रतिविव्याध लक्ष्मणः पनभिः शरैः} %||6-78-8||

\twolineshloka
{लक्ष्मणेन्द्रजितौ वीरौ महाबलशरासनौ}
{अन्योन्यं जघ्नतुर्बाणैर्विशिखैर्भीमविक्रमौ} %||6-78-9||

\twolineshloka
{तौ परस्परमभ्येत्य सर्वगात्रेषु धन्विनौ}
{घोरैर्विव्यधतुर्बाणैः कृतभावावुभौ जये} %||6-78-10||

\twolineshloka
{तस्मै दृढतरं क्रुद्धो हताश्वाय विभीषणः}
{वज्रस्पर्शसमान्पञ्च ससर्जोरसि मार्गणान्} %||6-78-11||

\twolineshloka
{ते तस्य कायं निर्भिद्य रुक्मपुङ्खा निमित्तगाः}
{बभूवुर्लोहितादिग्धा रक्ता इव महोरगाः} %||6-78-12||

\twolineshloka
{स पितृव्यस्य सङ्क्रुद्ध इन्द्रजिच्छरमाददे}
{उत्तमं रक्षसां मध्ये यमदत्तं महाबलः} %||6-78-13||

\twolineshloka
{तं समीक्ष्य महातेजा महेषुं तेन संहितम्}
{लक्ष्मणोऽप्याददे बाणमन्यं भीमपराक्रमः} %||6-78-14||

\twolineshloka
{कुबेरेण स्वयं स्वप्ने यद्दत्तममितात्मना}
{दुर्जयं दुर्विषह्यं च सेन्द्रैरपि सुरासुरैः} %||6-78-15||

\twolineshloka
{ताभ्यां तौ धनुषि श्रेष्ठे संहितौ सायकोत्तमौ}
{विकृष्यमाणौ वीराभ्यां भृशं जज्वलतुः श्रिया} %||6-78-16||

\twolineshloka
{तौ भासयन्तावाकाशं धनुर्भ्यां विशिखौ च्युतौ}
{मुखेन मुखमाहत्य संनिपेततुरोजसा} %||6-78-17||

\twolineshloka
{तौ महाग्रहसङ्काशावन्योन्यं संनिपत्य च}
{सङ्ग्रामे शतधा यातौ मेदिन्यां विनिपेततुः} %||6-78-18||

\twolineshloka
{शरौ प्रतिहतौ दृष्ट्वा तावुभौ रणमूर्धनि}
{व्रीडितो जातरोषौ च लक्ष्मणेन्द्रजितावुभौ} %||6-78-19||

\twolineshloka
{सुसंरब्धस्तु सौमित्रिरस्त्रं वारुणमाददे}
{रौद्रं महेद्रजिद्युद्धे व्यसृजद्युधि विष्ठितः} %||6-78-20||

\twolineshloka
{तयोः सुतुमुलं युद्धं सम्बभूवाद्भुतोपमम्}
{गगनस्थानि भूतानि लक्ष्मणं पर्यवारयन्} %||6-78-21||

\twolineshloka
{भैरवाभिरुते भीमे युद्धे वानरराक्षसाम्}
{भूतैर्बहुभिराकाशं विस्मितैरावृतं बभौ} %||6-78-22||

\twolineshloka
{ऋषयः पितरो देवा गन्धर्वा गरुणोरगाः}
{शतक्रतुं पुरस्कृत्य ररक्षुर्लक्ष्मणं रणे} %||6-78-23||

\twolineshloka
{अथान्यं मार्गणश्रेष्ठं सन्दधे रावणानुजः}
{हुताशनसमस्पर्शं रावणात्मजदारुणम्} %||6-78-24||

\twolineshloka
{सुपत्रमनुवृत्ताङ्गं सुपर्वाणं सुसंस्थितम्}
{सुवर्णविकृतं वीरः शरीरान्तकरं शरम्} %||6-78-25||

\twolineshloka
{दुरावारं दुर्विषहं राक्षसानां भयावहम्}
{आशीविषविषप्रख्यं देवसङ्घैः समर्चितम्} %||6-78-26||

\twolineshloka
{येन शक्रो महातेजा दानवानजयत्प्रभुः}
{पुरा देवासुरे युद्धे वीर्यवान्हरिवाहनः} %||6-78-27||

\twolineshloka
{तदैन्द्रमस्त्रं सौमित्रिः संयुगेष्वपराजितम्}
{शरश्रेष्ठं धनुः श्रेष्ठे नरश्रेष्ठोऽभिसन्दधे} %||6-78-28||

\twolineshloka
{सन्धायामित्रदलनं विचकर्ष शरासनम्}
{सज्यमायम्य दुर्धर्शः कालो लोकक्षये यथा} %||6-78-29||

\twolineshloka
{सन्धाय धनुषि श्रेष्ठे विकर्षन्निदमब्रवीत्}
{लक्ष्मीवाँल्लक्ष्मणो वाक्यमर्थसाधकमात्मनः} %||6-78-30||

\twolineshloka
{धर्मात्मा सत्यसन्धश्च रामो दाशरथिर्यदि}
{पौरुषे चाप्रतिद्वन्द्वस्तदेनं जहि रावणिम्} %||6-78-31||

\threelineshloka
{इत्युक्त्वा बाणमाकर्णं विकृष्य तमजिह्मगम्}
{लक्ष्मणः समरे वीरः ससर्जेन्द्रजितं प्रति}
{ऐन्द्रास्त्रेण समायुज्य लक्ष्मणः परवीरहा} %||6-78-32||

\twolineshloka
{तच्छिरः सशिरस्त्राणं श्रीमज्ज्वलितकुण्डलम्}
{प्रमथ्येन्द्रजितः कायात्पपात धरणीतले} %||6-78-33||

\twolineshloka
{तद्राक्षसतनूजस्य छिन्नस्कन्धं शिरो महत्}
{तपनीयनिभं भूमौ ददृशे रुधिरोक्षितम्} %||6-78-34||

\twolineshloka
{हतस्तु निपपाताशु धरण्यां रावणात्मजः}
{कवची सशिरस्त्राणो विध्वस्तः सशरासनः} %||6-78-35||

\twolineshloka
{चुक्रुशुस्ते ततः सर्वे वानराः सविभीषणाः}
{हृष्यन्तो निहते तस्मिन्देवा वृत्रवधे यथा} %||6-78-36||

\twolineshloka
{अथान्तरिक्षे भूतानामृषीणां च महात्मनाम्}
{अभिजज्ञे च संनादो गन्धर्वाप्सरसामपि} %||6-78-37||

\twolineshloka
{पतितं समभिज्ञाय राक्षसी सा महाचमूः}
{वध्यमाना दिशो भेजे हरिभिर्जितकाशिभिः} %||6-78-38||

\twolineshloka
{वनरैर्वध्यमानास्ते शस्त्राण्युत्सृज्य राक्षसाः}
{लङ्कामभिमुखाः सर्वे नष्टसंज्ञाः प्रधाविताः} %||6-78-39||

\twolineshloka
{दुद्रुवुर्बहुधा भीता राक्षसाः शतशो दिशः}
{त्यक्त्वा प्रहरणान्सर्वे पट्टसासिपरश्वधान्} %||6-78-40||

\twolineshloka
{केचिल्लङ्कां परित्रस्ताः प्रविष्टा वानरार्दिताः}
{समुद्रे पतिताः केचित्केचित्पर्वतमाश्रिताः} %||6-78-41||

\twolineshloka
{हतमिन्द्रजितं दृष्ट्वा शयानं समरक्षितौ}
{राक्षसानां सहस्रेषु न कश्चित्प्रत्यदृश्यत} %||6-78-42||

\twolineshloka
{यथास्तं गत आदित्ये नावतिष्ठन्ति रश्मयः}
{तथा तस्मिन्निपतिते राक्षसास्ते गता दिशः} %||6-78-43||

\twolineshloka
{शान्तरक्श्मिरिवादित्यो निर्वाण इव पावकः}
{स बभूव महातेजा व्यपास्त गतजीवितः} %||6-78-44||

\twolineshloka
{प्रशान्तपीडा बहुलो विनष्टारिः प्रहर्षवान्}
{बभूव लोकः पतिते राक्षसेन्द्रसुते तदा} %||6-78-45||

\twolineshloka
{हर्षं च शक्रो भगवान्सह सर्वैः सुरर्षभैः}
{जगाम निहते तस्मिन्राक्षसे पापकर्मणि} %||6-78-46||

\twolineshloka
{शुद्धा आपो नभश्चैव जहृषुर्दैत्यदानवाः}
{आजग्मुः पतिते तस्मिन्सर्वलोकभयावहे} %||6-78-47||

\twolineshloka
{ऊचुश्च सहिताः सर्वे देवगन्धर्वदानवाः}
{विज्वराः शान्तकलुषा ब्राह्मणा विचरन्त्विति} %||6-78-48||

\twolineshloka
{ततोऽभ्यनन्दन्संहृष्टाः समरे हरियूथपाः}
{तमप्रतिबलं दृष्ट्वा हतं नैरृतपुङ्गवम्} %||6-78-49||

\twolineshloka
{विभीषणो हनूमांश्च जाम्बवांश्चर्क्षयूथपः}
{विजयेनाभिनन्दन्तस्तुष्टुवुश्चापि लक्ष्मणम्} %||6-78-50||

\twolineshloka
{क्ष्वेडन्तश्च नदन्तश्च गर्जन्तश्च प्लवङ्गमाः}
{लब्धलक्षा रघुसुतं परिवार्योपतस्थिरे} %||6-78-51||

\twolineshloka
{लाङ्गूलानि प्रविध्यन्तः स्फोटयन्तश्च वानराः}
{लक्ष्मणो जयतीत्येवं वाक्यं व्यश्रावयंस्तदा} %||6-78-52||

\twolineshloka
{अन्योन्यं च समाश्लिष्य कपयो हृष्टमानसाः}
{चक्रुरुच्चावचगुणा राघवाश्रयजाः कथाः} %||6-78-53||

\fourlineindentedshloka
{तदसुकरमथाभिवीक्ष्य हृष्टाः}
{प्रियसुहृदो युधि लक्ष्मणस्य कर्म}
{परममुपलभन्मनःप्रहर्षम्}
{विनिहतमिन्द्ररिपुं निशम्य देवाः} %||6-79-54||


{॥इत्यार्षे श्रीमद्रामायणे वाल्मीकीये आदिकाव्ये युद्धकाण्डे अष्टसप्ततितमः सर्गः॥}

\sect{एकोनाशीतितमः सर्गः}

\twolineshloka
{रुधिरक्लिन्नगात्रस्तु लक्ष्मणः शुभलक्षणः}
{बभूव हृष्टस्तं हत्वा शक्रजेतारमाहवे} %||6-79-1||

\twolineshloka
{ततः स जाम्बवन्तं च हनूमन्तं च वीर्यवान्}
{संनिवर्त्य महातेजास्तांश्च सर्वान्वनौकसः} %||6-79-2||

\twolineshloka
{आजगाम ततः शीघ्रं यत्र सुग्रीवराघवौ}
{विभीषणमवष्टभ्य हनूमन्तं च लक्ष्मणः} %||6-79-3||

\threelineshloka
{ततो राममभिक्रम्य सौमित्रिरभिवाद्य च}
{तस्थौ भ्रातृसमीपस्थः शक्रस्येन्द्रानुजो यथा}
{आचचक्षे तदा वीरो घोरमिन्द्रजितो वधम्} %||6-79-4||

\twolineshloka
{रावणस्तु शिरश्छिन्नं लक्ष्मणेन महात्मना}
{न्यवेदयत रामाय तदा हृष्टो विभीषणः} %||6-79-5||

\threelineshloka
{उपवेश्य तमुत्सङ्गे परिष्वज्यावपीडितम्}
{मूर्ध्नि चैनमुपाघ्राय भूयः संस्पृश्य च त्वरन्}
{उवाच लक्ष्मणं वाक्यमाश्वास्य पुरुषर्षभः} %||6-79-6||

\threelineshloka
{कृतं परमकल्याणं कर्म दुष्करकारिणा}
{निरमित्रः कृतोऽस्म्यद्य निर्यास्यति हि रावणः}
{बलव्यूहेन महता श्रुत्वा पुत्रं निपातितम्} %||6-79-7||

\twolineshloka
{तं पुत्रवधसन्तप्तं निर्यान्तं राक्षसाधिपम्}
{बलेनावृत्य महता निहनिष्यामि दुर्जयम्} %||6-79-8||

\twolineshloka
{त्वया लक्ष्मण नाथेन सीता च पृथिवी च मे}
{न दुष्प्रापा हते त्वद्य शक्रजेतरि चाहवे} %||6-79-9||

\twolineshloka
{स तं भ्रातरमाश्वास्य पारिष्वज्य च राघवः}
{रामः सुषेणं मुदितः समाभाष्येदमब्रवीत्} %||6-79-10||

\threelineshloka
{सशल्योऽयं महाप्राज्ञः सौमित्रिर्मित्रवत्सलः}
{यथा भवति सुस्वस्थस्तथा त्वं समुपाचर}
{विशल्यः क्रियतां क्षिप्रं सौमित्रिः सविभीषणः} %||6-79-11||

\threelineshloka
{कृष वानरसैन्यानां शूराणां द्रुमयोधिनाम्}
{ये चान्येऽत्र च युध्यन्तः सशल्या व्रणिनस्तथा}
{तेऽपि सर्वे प्रयत्नेन क्रियन्तां सुखिनस्त्वया} %||6-79-12||

\twolineshloka
{एवमुक्तः स रामेण महात्मा हरियूथपः}
{लक्ष्मणाय ददौ नस्तः सुषेणः परमौषधम्} %||6-79-13||

\twolineshloka
{स तस्य गन्धमाघ्राय विशल्यः समपद्यत}
{तदा निर्वेदनश्चैव संरूढव्रण एव च} %||6-79-14||

\twolineshloka
{विभीषण मुखानां च सुहृदां राघवाज्ञया}
{सर्ववानरमुख्यानां चिकित्सां स तदाकरोत्} %||6-79-15||

\twolineshloka
{ततः प्रकृतिमापन्नो हृतशल्यो गतव्यथः}
{सौमित्रिर्मुदितस्तत्र क्षणेन विगतज्वरः} %||6-79-16||

\fourlineindentedshloka
{तथैव रामः प्लवगाधिपस्तदा}
{विभीषणश्चर्क्षपतिश्च जाम्बवान्}
{अवेक्ष्य सौमित्रिमरोगमुत्थितम्}
{मुदा ससैन्यः सुचिरं जहर्षिरे} %||6-79-17||

\fourlineindentedshloka
{अपूजयत्कर्म स लक्ष्मणस्य}
{सुदुष्करं दाशरथिर्महात्मा}
{हृष्टा बभूवुर्युधि यूथपेन्द्रा}
{निशम्य तं शक्रजितं निपातितम्} %||6-80-18||


{॥इत्यार्षे श्रीमद्रामायणे वाल्मीकीये आदिकाव्ये युद्धकाण्डे एकोनाशीतितमः सर्गः॥}

\sect{अशीतितमः सर्गः}

\twolineshloka
{ततः पौलस्त्य सचिवाः श्रुत्वा चेन्द्रजितं हतम्}
{आचचक्षुरभिज्ञाय दशग्रीवाय सव्यथाः} %||6-80-1||

\twolineshloka
{युद्धे हतो महाराज लक्ष्मणेन तवात्मजः}
{विभीषणसहायेन मिषतां नो महाद्युते} %||6-80-2||

\twolineshloka
{शूरः शूरेण सङ्गम्य संयुगेष्वपराजितः}
{लक्ष्णनेन हतः शूरः पुत्रस्ते विबुधेन्द्रजित्} %||6-80-3||

\twolineshloka
{स तं प्रतिभयं श्रुत्वा वधं पुत्रस्य दारुणम्}
{घोरमिन्द्रजितः सङ्ख्ये कश्मलं प्राविशन्महत्} %||6-80-4||

\twolineshloka
{उपलभ्य चिरात्संज्ञां राजा राक्षसपुङ्गवः}
{पुत्रशोकार्दितो दीनो विललापाकुलेन्द्रियः} %||6-80-5||

\twolineshloka
{हा राक्षसचमूमुख्य मम वत्स महारथ}
{जित्वेन्द्रं कथमद्य त्वं लक्ष्मणस्य वशं गतः} %||6-80-6||

\twolineshloka
{ननु त्वमिषुभिः क्रुद्धो भिन्द्याः कालान्तकावपि}
{मन्दरस्यापि शृङ्गाणि किं पुनर्लक्ष्मणं रणे} %||6-80-7||

\twolineshloka
{अद्य वैवस्वतो राजा भूयो बहुमतो मम}
{येनाद्य त्वं महाबाहो संयुक्तः कालधर्मणा} %||6-80-8||

\twolineshloka
{एष पन्थाः सुयोधानां सर्वामरगणेष्वपि}
{यः कृते हन्यते भर्तुः स पुमान्स्वर्गमृच्छति} %||6-80-9||

\twolineshloka
{अद्य देवगणाः सर्वे लोकपालास्तथर्षयः}
{हतमिन्द्रजितं दृष्ट्वा सुखं स्वप्स्यन्ति निर्भयाः} %||6-80-10||

\twolineshloka
{अद्य लोकास्त्रयः कृत्स्नाः पृथिवी च सकानना}
{एकेनेन्द्रजिता हीना शूण्येव प्रतिभाति मे} %||6-80-11||

\twolineshloka
{अद्य नैरृतकन्यायां श्रोष्याम्यन्तःपुरे रवम्}
{करेणुसङ्घस्य यथा निनादं गिरिगह्वरे} %||6-80-12||

\twolineshloka
{यौवराज्यं च लङ्कां च रक्षांसि च परन्तप}
{मातरं मां च भार्यां च क्व गतोऽसि विहाय नः} %||6-80-13||

\twolineshloka
{मम नाम त्वया वीर गतस्य यमसादनम्}
{प्रेतकार्याणि कार्याणि विपरीते हि वर्तसे} %||6-80-14||

\twolineshloka
{स त्वं जीवति सुग्रीवे राघवे च सलक्ष्मणे}
{मम शल्यमनुद्धृत्य क्व गतोऽसि विहाय नः} %||6-80-15||

\twolineshloka
{एवमादिविलापार्तं रावणं राक्षसाधिपम्}
{आविवेश महान्कोपः पुत्रव्यसनसम्भवः} %||6-80-16||

\twolineshloka
{घोरं प्रकृत्या रूपं तत्तस्य क्रोधाग्निमूर्छितम्}
{बभूव रूपं रुद्रस्य क्रुद्धस्येव दुरासदम्} %||6-80-17||

\twolineshloka
{तस्य क्रुद्धस्य नेत्राभ्यां प्रापतन्नस्रबिन्दवः}
{दीप्ताभ्यामिव दीपाभ्यां सार्चिषः स्नेहबिन्दवः} %||6-80-18||

\twolineshloka
{दन्तान्विदशतस्तस्य श्रूयते दशनस्वनः}
{यन्त्रस्यावेष्ट्यमानस्य महतो दानवैरिव} %||6-80-19||

\twolineshloka
{कालाग्निरिव सङ्क्रुद्धो यां यां दिशमवैक्षत}
{तस्यां तस्यां भयत्रस्ता राक्षसाः संनिलिल्यिरे} %||6-80-20||

\twolineshloka
{तमन्तकमिव क्रुद्धं चराचरचिखादिषुम्}
{वीक्षमाणं दिशः सर्वा राक्षसा नोपचक्रमुः} %||6-80-21||

\twolineshloka
{ततः परमसङ्क्रुद्धो रावणो राक्षसाधिपः}
{अब्रवीद्रक्षसां मध्ये संस्तम्भयिषुराहवे} %||6-80-22||

\twolineshloka
{मया वर्षसहस्राणि चरित्वा दुश्चरं तपः}
{तेषु तेष्ववकाशेषु स्वयम्भूः परितोषितः} %||6-80-23||

\twolineshloka
{तस्यैव तपसो व्युष्ट्या प्रसादाच्च स्वयम्भुवः}
{नासुरेभ्यो न देवेभ्यो भयं मम कदाचन} %||6-80-24||

\twolineshloka
{कवचं ब्रह्मदत्तं मे यदादित्यसमप्रभम्}
{देवासुरविमर्देषु न भिन्नं वज्रशक्तिभिः} %||6-80-25||

\twolineshloka
{तेन मामद्य संयुक्तं रथस्थमिह संयुगे}
{प्रतीयात्कोऽद्य मामाजौ साक्षादपि पुरन्दरः} %||6-80-26||

\twolineshloka
{यत्तदाभिप्रसन्नेन सशरं कार्मुकं महत्}
{देवासुरविमर्देषु मम दत्तं स्वयम्भुवा} %||6-80-27||

\twolineshloka
{अद्य तूर्यशतैर्भीमं धनुरुत्थाप्यतां महत्}
{रामलक्ष्मणयोरेव वधाय परमाहवे} %||6-80-28||

\twolineshloka
{स पुत्रवधसन्तप्तः शूरः क्रोधवशं गतः}
{समीक्ष्य रावणो बुद्ध्या सीतां हन्तुं व्यवस्यत} %||6-80-29||

\twolineshloka
{प्रत्यवेक्ष्य तु ताम्राक्षः सुघोरो घोरदर्शनान्}
{दीनो दीनस्वरान्सर्वांस्तानुवाच निशाचरान्} %||6-80-30||

\twolineshloka
{मायया मम वत्सेन वञ्चनार्थं वनौकसाम्}
{किञ्चिदेव हतं तत्र सीतेयमिति दर्शितम्} %||6-80-31||

\threelineshloka
{तदिदं सत्यमेवाहं करिष्ये प्रियमात्मनः}
{वैदेहीं नाशयिष्यामि क्षत्रबन्धुमनुव्रताम्}
{इत्येवमुक्त्वा सचिवान्खड्गमाशु परामृशत्} %||6-80-32||

\twolineshloka
{उद्धृत्य गुणसम्पन्नं विमलाम्बरवर्चसम्}
{निष्पपात स वेगेन सभायाः सचिवैर्वृतः} %||6-80-33||

\twolineshloka
{रावणः पुत्रशोकेन भृशमाकुलचेतनः}
{सङ्क्रुद्धः खड्गमादाय सहसा यत्र मैथिली} %||6-80-34||

\twolineshloka
{व्रजन्तं राक्षसं प्रेक्ष्य सिंहनादं प्रचुक्रुशुः}
{ऊचुश्चान्योन्यमाश्लिष्य सङ्क्रुद्धं प्रेक्ष्य राक्षसाः} %||6-80-35||

\threelineshloka
{अद्यैनं तावुभौ दृष्ट्वा भ्रातरौ प्रव्यथिष्यतः}
{लोकपाला हि चत्वारः क्रुद्धेनानेन निर्जिताः}
{बहवः शत्रवश्चान्ये संयुगेष्वभिपातिताः} %||6-80-36||

\twolineshloka
{तेषां संजल्पमानानामशोकवनिकां गताम्}
{अभिदुद्राव वैदेहीं रावणः क्रोधमूर्छितः} %||6-80-37||

\twolineshloka
{वार्यमाणः सुसङ्क्रुद्धः सुहृद्भिर्हितबुद्धिभिः}
{अभ्यधावत सङ्क्रुद्धः खे ग्रहो रोहिणीमिव} %||6-80-38||

\twolineshloka
{मैथिली रक्ष्यमाणा तु राक्षसीभिरनिन्दिता}
{ददर्श राक्षसं क्रुद्धं निस्त्रिंशवरधारिणम्} %||6-80-39||

\twolineshloka
{तं निशाम्य सनिस्त्रिंशं व्यथिता जनकात्मजा}
{निवार्यमाणं बहुशः सुहृद्भिरनिवर्तिनम्} %||6-80-40||

\twolineshloka
{यथायं मामभिक्रुद्धः समभिद्रवति स्वयम्}
{वधिष्यति सनाथां मामनाथामिव दुर्मतिः} %||6-80-41||

\twolineshloka
{बहुशश्चोदयामास भर्तारं मामनुव्रताम्}
{भार्या भव रमस्येति प्रत्याख्यातोऽभवन्मया} %||6-80-42||

\twolineshloka
{सोऽयं मामनुपस्थानाद्व्यक्तं नैराश्यमागतः}
{क्रोधमोहसमाविष्टो निहन्तुं मां समुद्यतः} %||6-80-43||

\threelineshloka
{अथ वा तौ नरव्याघ्रौ भ्रातरौ रामलक्ष्मणौ}
{मन्निमित्तमनार्येण समरेऽद्य निपातितौ}
{अहो धिन्मन्निमित्तोऽयं विनाशो राजपुत्रयोः} %||6-80-44||

\threelineshloka
{हनूमतो हि तद्वाक्यं न कृतं क्षुद्रया मया}
{यद्यहं तस्य पृष्ठेन तदायासमनिन्दिता}
{नाद्यैवमनुशोचेयं भर्तुरङ्कगता सती} %||6-80-45||

\twolineshloka
{मन्ये तु हृदयं तस्याः कौसल्यायाः फलिष्यति}
{एकपुत्रा यदा पुत्रं विनष्टं श्रोष्यते युधि} %||6-80-46||

\twolineshloka
{सा हि जन्म च बाल्यं च यौवनं च महात्मनः}
{धर्मकार्याणि रूपं च रुदती संस्रमिष्यति} %||6-80-47||

\twolineshloka
{निराशा निहते पुत्रे दत्त्वा श्राद्धमचेतना}
{अग्निमारोक्ष्यते नूनमपो वापि प्रवेक्ष्यति} %||6-80-48||

\twolineshloka
{धिगस्तु कुब्जामसतीं मन्थरां पापनिश्चयाम्}
{यन्निमित्तमिदं दुःखं कौसल्या प्रतिपत्स्यते} %||6-80-49||

\twolineshloka
{इत्येवं मैथिलीं दृष्ट्वा विलपन्तीं तपस्विनीम्}
{रोहिणीमिव चन्द्रेण विना ग्रहवशं गताम्} %||6-80-50||

\twolineshloka
{सुपार्श्वो नाम मेधावी रावणं राक्षसेश्वरम्}
{निवार्यमाणं सचिवैरिदं वचनमब्रवीत्} %||6-80-51||

\twolineshloka
{कथं नाम दशग्रीव साक्षाद्वैश्रवणानुज}
{हन्तुमिच्छसि वैदेहीं क्रोधाद्धर्ममपास्य हि} %||6-80-52||

\twolineshloka
{वेद विद्याव्रत स्नातः स्वधर्मनिरतः सदा}
{स्त्रियाः कस्माद्वधं वीर मन्यसे राक्षसेश्वर} %||6-80-53||

\twolineshloka
{मैथिलीं रूपसम्पन्नां प्रत्यवेक्षस्व पार्थिव}
{त्वमेव तु सहास्माभी राघवे क्रोधमुत्सृज} %||6-80-54||

\twolineshloka
{अभ्युत्थानं त्वमद्यैव कृष्णपक्षचतुर्दशीम्}
{कृत्वा निर्याह्यमावास्यां विजयाय बलैर्वृतः} %||6-80-55||

\twolineshloka
{शूरो धीमान्रथी खड्गी रथप्रवरमास्थितः}
{हत्वा दाशरथिं रामं भवान्प्राप्स्यति मैथिलीम्} %||6-80-56||

\fourlineindentedshloka
{स तद्दुरात्मा सुहृदा निवेदितम्}
{वचः सुधर्म्यं प्रतिगृह्य रावणः}
{गृहं जगामाथ ततश्च वीर्यवान्}
{पुनः सभां च प्रययौ सुहृद्वृतः} %||6-81-57||


{॥इत्यार्षे श्रीमद्रामायणे वाल्मीकीये आदिकाव्ये युद्धकाण्डे अशीतितमः सर्गः॥}

\sect{एकाशीतितमः सर्गः}

\twolineshloka
{स प्रविश्य सभां राजा दीनः परमदुःखितः}
{निषसादासने मुख्ये सिंहः क्रुद्ध इव श्वसन्} %||6-81-1||

\twolineshloka
{अब्रवीच्च तदा सर्वान्बलमुख्यान्महाबलः}
{रावणः प्राञ्जलीन्वाक्यं पुत्रव्यसनकर्शितः} %||6-81-2||

\twolineshloka
{सर्वे भवन्तः सर्वेण हस्त्यश्वेन समावृताः}
{निर्यान्तु रथसङ्घैश्च पादातैश्चोपशोभिताः} %||6-81-3||

\twolineshloka
{एकं रामं परिक्षिप्य समरे हन्तुमर्हथ}
{प्रहृष्टा शरवर्षेण प्रावृट्काल इवाम्बुदाः} %||6-81-4||

\twolineshloka
{अथ वाहं शरैर्तीक्ष्णैर्भिन्नगात्रं महारणे}
{भवद्भिः श्वो निहन्तास्मि रामं लोकस्य पश्यतः} %||6-81-5||

\twolineshloka
{इत्येवं राक्षसेन्द्रस्य वाक्यमादाय राक्षसाः}
{निर्ययुस्ते रथैः शीघ्रं नागानीकैश्च संवृताः} %||6-81-6||

\twolineshloka
{स सङ्ग्रामो महाभीमः सूर्यस्योदयनं प्रति}
{रक्षसां वानराणां च तुमुलः समपद्यत} %||6-81-7||

\twolineshloka
{ते गदाभिर्विचित्राभिः प्रासैः खड्गैः परश्वधैः}
{अन्योन्यं समरे जघ्नुस्तदा वानरराक्षसाः} %||6-81-8||

\twolineshloka
{मातङ्गरथकूलस्य वाजिमत्स्या ध्वजद्रुमाः}
{शरीरसङ्घाटवहाः प्रसस्रुः शोणितापगाः} %||6-81-9||

\twolineshloka
{ध्वजवर्मरथानश्वान्नानाप्रहरणानि च}
{आप्लुत्याप्लुत्य समरे वानरेन्द्रा बभञ्जिरे} %||6-81-10||

\twolineshloka
{केशान्कर्णललाटांश्च नासिकाश्च प्लवङ्गमाः}
{रक्षसां दशनैस्तीक्ष्णैर्नखैश्चापि व्यकर्तयन्} %||6-81-11||

\twolineshloka
{एकैकं राक्षसं सङ्ख्ये शतं वानरपुङ्गवाः}
{अभ्यधावन्त फलिनं वृक्षं शकुनयो यथा} %||6-81-12||

\twolineshloka
{तथा गदाभिर्गुर्वीभिः प्रासैः खड्गैः परश्वधैः}
{निर्जघ्नुर्वानरान्घोरान्राक्षसाः पर्वतोपमाः} %||6-81-13||

\twolineshloka
{राक्षसैर्वध्यमानानां वानराणां महाचमूः}
{शरण्यं शरणं याता रामं दशरथात्मजम्} %||6-81-14||

\twolineshloka
{ततो रामो महातेजा धनुरादाय वीर्यवान्}
{प्रविश्य राक्षसं सैन्यं शरवर्षं ववर्ष ह} %||6-81-15||

\twolineshloka
{प्रविष्टं तु तदा रामं मेघाः सूर्यमिवाम्बरे}
{नाभिजग्मुर्महाघोरं निर्दहन्तं शराग्निना} %||6-81-16||

\twolineshloka
{कृतान्येव सुघोराणि रामेण रजनीचराः}
{रणे रामस्य ददृशुः कर्माण्यसुकराणि च} %||6-81-17||

\twolineshloka
{चालयन्तं महानीकं विधमन्तं महारथान्}
{ददृशुस्ते न वै रामं वातं वनगतं यथा} %||6-81-18||

\twolineshloka
{छिन्नं भिन्नं शरैर्दग्धं प्रभग्नं शस्त्रपीडितम्}
{बलं रामेण ददृशुर्न रमं शीघ्रकारिणम्} %||6-81-19||

\twolineshloka
{प्रहरन्तं शरीरेषु न ते पश्यन्ति राघवम्}
{इन्द्रियार्थेषु तिष्ठन्तं भूतात्मानमिव प्रजाः} %||6-81-20||

\twolineshloka
{एष हन्ति गजानीकमेष हन्ति महारथान्}
{एष हन्ति शरैस्तीक्ष्णैः पदातीन्वाजिभिः सह} %||6-81-21||

\twolineshloka
{इति ते राक्षसाः सर्वे रामस्य सदृशान्रणे}
{अन्योन्यकुपिता जघ्नुः सादृश्याद्राघवस्य ते} %||6-81-22||

\twolineshloka
{न ते ददृशिरे रामं दहन्तमरिवाहिनीम्}
{मोहिताः परमास्त्रेण गान्धर्वेण महात्मना} %||6-81-23||

\twolineshloka
{ते तु रामसहस्राणि रणे पश्यन्ति राक्षसाः}
{पुनः पश्यन्ति काकुत्स्थमेकमेव महाहवे} %||6-81-24||

\twolineshloka
{भ्रमन्तीं काञ्चनीं कोटिं कार्मुकस्य महात्मनः}
{अलातचक्रप्रतिमां ददृशुस्ते न राघवम्} %||6-81-25||

\twolineshloka
{शरीरनाभि सत्त्वार्चिः शरीरं नेमिकार्मुकम्}
{ज्याघोषतलनिर्घोषं तेजोबुद्धिगुणप्रभम्} %||6-81-26||

\twolineshloka
{दिव्यास्त्रगुणपर्यन्तं निघ्नन्तं युधि राक्षसान्}
{ददृशू रामचक्रं तत्कालचक्रमिव प्रजाः} %||6-81-27||

\twolineshloka
{अनीकं दशसाहस्रं रथानां वातरंहसाम्}
{अष्टादशसहस्राणि कुञ्जराणां तरस्विनाम्} %||6-81-28||

\twolineshloka
{चतुर्दशसहस्राणि सारोहाणां च वाजिनाम्}
{पूर्णे शतसहस्रे द्वे राक्षसानां पदातिनाम्} %||6-81-29||

\twolineshloka
{दिवसस्याष्टमे भागे शरैरग्निशिखोपमैः}
{हतान्येकेन रामेण रक्षसां कामरूपिणाम्} %||6-81-30||

\twolineshloka
{ते हताश्वा हतरथाः श्रान्ता विमथितध्वजाः}
{अभिपेतुः पुरीं लङ्कां हतशेषा निशाचराः} %||6-81-31||

\twolineshloka
{हतैर्गजपदात्यश्वैस्तद्बभूव रणाजिरम्}
{आक्रीडभूमी रुद्रस्य क्रुद्धस्येव पिनाकिनः} %||6-81-32||

\twolineshloka
{ततो देवाः सगन्धर्वाः सिद्धाश्च परमर्षयः}
{साधु साध्विति रामस्य तत्कर्म समपूजयन्} %||6-81-33||

\twolineshloka
{अब्रवीच्च तदा रामः सुग्रीवं प्रत्यनन्तरम्}
{एतदस्त्रबलं दिव्यं मम वा त्र्यम्बकस्य वा} %||6-81-34||

\fourlineindentedshloka
{निहत्य तां राक्षसवाहिनीं तु}
{रामस्तदा शक्रसमो महात्मा}
{अस्त्रेषु शस्त्रेषु जितक्लमश्च}
{संस्तूयते देवगणैः प्रहृष्टैः} %||6-82-35||


{॥इत्यार्षे श्रीमद्रामायणे वाल्मीकीये आदिकाव्ये युद्धकाण्डे एकाशीतितमः सर्गः॥}

\sect{द्व्यशीतितमः सर्गः}

\twolineshloka
{तानि नागसहस्राणि सारोहाणां च वाजिनाम्}
{रथानां चाग्निवर्णानां सध्वजानां सहस्रशः} %||6-82-1||

\twolineshloka
{राक्षसानां सहस्राणि गदापरिघयोधिनाम्}
{काञ्चनध्वजचित्राणां शूराणां कामरूपिणाम्} %||6-82-2||

\twolineshloka
{निहतानि शरैस्तीक्ष्णैस्तप्तकाञ्चनभूषणैः}
{रावणेन प्रयुक्तानि रामेणाक्लिष्टकर्मणा} %||6-82-3||

\twolineshloka
{दृष्ट्वा श्रुत्वा च सम्भ्रान्ता हतशेषा निशाचराः}
{राक्षस्यश्च समागम्य दीनाश्चिन्तापरिप्लुताः} %||6-82-4||

\twolineshloka
{विधवा हतपुत्राश्च क्रोशन्त्यो हतबान्धवाः}
{राक्षस्यः सह सङ्गम्य दुःखार्ताः पर्यदेवयन्} %||6-82-5||

\twolineshloka
{कथं शूर्पणखा वृद्धा कराला निर्णतोदरी}
{आससाद वने रामं कन्दर्पमिव रूपिणम्} %||6-82-6||

\twolineshloka
{सुकुमारं महासत्त्वं सर्वभूतहिते रतम्}
{तं दृष्ट्वा लोकवध्या सा हीनरूपा प्रकामिता} %||6-82-7||

\twolineshloka
{कथं सर्वगुणैर्हीना गुणवन्तं महौजसम्}
{सुमुखं दुर्मुखी रामं कामयामास राक्षसी} %||6-82-8||

\twolineshloka
{जनस्यास्याल्पभाग्यत्वात्पलिनी श्वेतमूर्धजा}
{अकार्यमपहास्यं च सर्वलोकविगर्हितम्} %||6-82-9||

\twolineshloka
{राक्षसानां विनाशाय दूषणस्य खरस्य च}
{चकाराप्रतिरूपा सा राघवस्य प्रधर्षणम्} %||6-82-10||

\twolineshloka
{तन्निमित्तमिदं वैरं रावणेन कृतं महत्}
{वधाय नीता सा सीता दशग्रीवेण रक्षसा} %||6-82-11||

\twolineshloka
{न च सीतां दशग्रीवः प्राप्नोति जनकात्मजाम्}
{बद्धं बलवता वैरमक्षयं राघवेण ह} %||6-82-12||

\twolineshloka
{वैदेहीं प्रार्थयानं तं विराधं प्रेक्ष्य राक्षसम्}
{हतमेकेन रामेण पर्याप्तं तन्निदर्शनम्} %||6-82-13||

\twolineshloka
{चतुर्दशसहस्राणि रक्षसां भीमकर्मणाम्}
{निहतानि जनस्थाने शरैरग्निशिखोपमैः} %||6-82-14||

\twolineshloka
{खरश्च निहतः सङ्ख्ये दूषणस्त्रिशिरास्तथा}
{शरैरादित्यसङ्काशैः पर्याप्तं तन्निदर्शनम्} %||6-82-15||

\twolineshloka
{हतो योजनबाहुश्च कबन्धो रुधिराशनः}
{क्रोधार्तो विनदन्सोऽथ पर्याप्तं तन्निदर्शनम्} %||6-82-16||

\twolineshloka
{जघान बलिनं रामः सहस्रनयनात्मजम्}
{बालिनं मेघसङ्काशं पर्याप्तं तन्निदर्शनम्} %||6-82-17||

\twolineshloka
{ऋश्यमूके वसञ्शैले दीनो भग्नमनोरथः}
{सुग्रीवः स्थापितो राज्ये पर्याप्तं तन्निदर्शनम्} %||6-82-18||

\twolineshloka
{धर्मार्थसहितं वाक्यं सर्वेषां रक्षसां हितम्}
{युक्तं विभीषणेनोक्तं मोहात्तस्य न रोचते} %||6-82-19||

\twolineshloka
{विभीषणवचः कुर्याद्यदि स्म धनदानुजः}
{श्मशानभूता दुःखार्ता नेयं लङ्का पुरी भवेत्} %||6-82-20||

\twolineshloka
{कुम्भकर्णं हतं श्रुत्वा राघवेण महाबलम्}
{प्रियं चेन्द्रजितं पुत्रं रावणो नावबुध्यते} %||6-82-21||

\twolineshloka
{मम पुत्रो मम भ्राता मम भर्ता रणे हतः}
{इत्येवं श्रूयते शब्दो राक्षसानां कुले कुले} %||6-82-22||

\twolineshloka
{रथाश्चाश्वाश्च नागाश्च हताः शतसहस्रशः}
{रणे रामेण शूरेण राक्षसाश्च पदातयः} %||6-82-23||

\twolineshloka
{रुद्रो वा यदि वा विष्णुर्महेन्द्रो वा शतक्रतुः}
{हन्ति नो रामरूपेण यदि वा स्वयमन्तकः} %||6-82-24||

\twolineshloka
{हतप्रवीरा रामेण निराशा जीविते वयम्}
{अपश्यन्त्यो भयस्यान्तमनाथा विलपामहे} %||6-82-25||

\twolineshloka
{रामहस्ताद्दशग्रीवः शूरो दत्तवरो युधि}
{इदं भयं महाघोरमुत्पन्नं नावबुध्यते} %||6-82-26||

\twolineshloka
{न देवा न च गन्धर्वा न पिशाचा न राक्षसाः}
{उपसृष्टं परित्रातुं शक्ता रामेण संयुगे} %||6-82-27||

\twolineshloka
{उत्पाताश्चापि दृश्यन्ते रावणस्य रणे रणे}
{कथयिष्यन्ति रामेण रावणस्य निबर्हणम्} %||6-82-28||

\twolineshloka
{पितामहेन प्रीतेन देवदानवराक्षसैः}
{रावणस्याभयं दत्तं मानुषेभ्यो न याचितम्} %||6-82-29||

\twolineshloka
{तदिदं मानुषान्मन्ये प्राप्तं निःसंशयं भयम्}
{जीवितान्तकरं घोरं रक्षसां रावणस्य च} %||6-82-30||

\twolineshloka
{पीड्यमानास्तु बलिना वरदानेन रक्षसा}
{दीप्तैस्तपोभिर्विबुधाः पितामहमपूजयन्} %||6-82-31||

\twolineshloka
{देवतानां हितार्थाय महात्मा वै पितामहः}
{उवाच देवताः सर्वा इदं तुष्टो महद्वचः} %||6-82-32||

\twolineshloka
{अद्य प्रभृति लोकांस्त्रीन्सर्वे दानवराक्षसाः}
{भयेन प्रावृता नित्यं विचरिष्यन्ति शाश्वतम्} %||6-82-33||

\twolineshloka
{दैवतैस्तु समागम्य सर्वैश्चेन्द्रपुरोगमैः}
{वृषध्वजस्त्रिपुरहा महादेवः प्रसादितः} %||6-82-34||

\twolineshloka
{प्रसन्नस्तु महादेवो देवानेतद्वचोऽब्रवीत्}
{उत्पत्स्यति हितार्थं वो नारी रक्षःक्षयावहा} %||6-82-35||

\twolineshloka
{एषा देवैः प्रयुक्ता तु क्षुद्यथा दानवान्पुरा}
{भक्षयिष्यति नः सीता राक्षसघ्नी सरावणान्} %||6-82-36||

\twolineshloka
{रावणस्यापनीतेन दुर्विनीतस्य दुर्मतेः}
{अयं निष्टानको घोरः शोकेन समभिप्लुतः} %||6-82-37||

\twolineshloka
{तं न पश्यामहे लोके यो नः शरणदो भवेत्}
{राघवेणोपसृष्टानां कालेनेव युगक्षये} %||6-82-38||

\fourlineindentedshloka
{इतीव सर्वा रजनीचरस्त्रियः}
{परस्परं सम्परिरभ्य बाहुभिः}
{विषेदुरार्तातिभयाभिपीडिता}
{विनेदुरुच्चैश्च तदा सुदारुणम्} %||6-83-39||


{॥इत्यार्षे श्रीमद्रामायणे वाल्मीकीये आदिकाव्ये युद्धकाण्डे द्व्यशीतितमः सर्गः॥}

\sect{त्र्यशीतितमः सर्गः}

\twolineshloka
{आर्तानां राक्षसीनां तु लङ्कायां वै कुले कुले}
{रावणः करुणं शब्दं शुश्राव परिवेदितम्} %||6-83-1||

\twolineshloka
{स तु दीर्घं विनिश्वस्य मुहूर्तं ध्यानमास्थितः}
{बभूव परमक्रुद्धो रावणो भीमदर्शनः} %||6-83-2||

\twolineshloka
{सन्दश्य दशनैरोष्ठं क्रोधसंरक्तलोचनः}
{राक्षसैरपि दुर्दर्शः कालाग्निरिव मूर्छितः} %||6-83-3||

\twolineshloka
{उवाच च समीपस्थान्राक्षसान्राक्षसेश्वरः}
{भयाव्यक्तकथांस्तत्र निर्दहन्निव चक्षुषा} %||6-83-4||

\twolineshloka
{महोदरं महापार्श्वं विरूपाक्षं च राक्षसम्}
{शीघ्रं वदत सैन्यानि निर्यातेति ममाज्ञया} %||6-83-5||

\twolineshloka
{तस्य तद्वचनं श्रुत्वा राक्षसास्ते भयार्दिताः}
{चोदयामासुरव्यग्रान्राक्षसांस्तान्नृपाज्ञया} %||6-83-6||

\twolineshloka
{ते तु सर्वे तथेत्युक्त्वा राक्षसा घोरदर्शनाः}
{कृतस्वस्त्ययनाः सर्वे रावणाभिमुखा ययुः} %||6-83-7||

\twolineshloka
{प्रतिपूज्य यथान्यायं रावणं ते महारथाः}
{तस्थुः प्राञ्जलयः सर्वे भर्तुर्विजयकाङ्क्षिणः} %||6-83-8||

\twolineshloka
{अथोवाच प्रहस्यैतान्रावणः क्रोधमूर्छितः}
{महोदरमहापार्श्वौ विरूपाक्षं च राक्षसम्} %||6-83-9||

\twolineshloka
{अद्य बाणैर्धनुर्मुक्तैर्युगान्तादित्यसंनिभैः}
{राघवं लक्ष्मणं चैव नेष्यामि यमसाधनम्} %||6-83-10||

\twolineshloka
{खरस्य कुम्भकर्णस्य प्रहस्तेन्द्रजितोस्तथा}
{करिष्यामि प्रतीकारमद्य शत्रुवधादहम्} %||6-83-11||

\twolineshloka
{नैवान्तरिक्षं न दिशो न नद्यो नापि सागरः}
{प्रकाशत्वं गमिष्यन्ति मद्बाणजलदावृताः} %||6-83-12||

\twolineshloka
{अद्य वानरयूथानां तानि यूथानि भागशः}
{धनुःसमुद्रादुद्भूतैर्मथिष्यामि शरोर्मिभिः} %||6-83-13||

\twolineshloka
{व्याकोशपद्मचक्राणि पद्मकेसरवर्चसाम्}
{अद्य यूथतटाकानि गजवत्प्रमथाम्यहम्} %||6-83-14||

\twolineshloka
{सशरैरद्य वदनैः सङ्ख्ये वानरयूथपाः}
{मण्डयिष्यन्ति वसुधां सनालैरिव पङ्कलैः} %||6-83-15||

\twolineshloka
{अद्य युद्धप्रचण्डानां हरीणां द्रुमयोधिनाम्}
{मुक्तेनैकेषुणा युद्धे भेत्स्यामि च शतंशतम्} %||6-83-16||

\twolineshloka
{हतो भर्ता हतो भ्राता यासां च तनया हताः}
{वधेनाद्य रिपोस्तासां कर्मोम्यस्रप्रमार्जनम्} %||6-83-17||

\twolineshloka
{अद्य मद्बाणनिर्भिन्नैः प्रकीर्णैर्गतचेतनैः}
{करोमि वानरैर्युद्धे यत्नावेक्ष्यतलां महीम्} %||6-83-18||

\twolineshloka
{अद्य गोमायवो गृध्रा ये च मांसाशिनोऽपरे}
{सर्वांस्तांस्तर्पयिष्यामि शत्रुमांसैः शरार्दितैः} %||6-83-19||

\twolineshloka
{कल्प्यतां मे रथशीघ्रं क्षिप्रमानीयतां धनुः}
{अनुप्रयान्तु मां युद्धे येऽवशिष्टा निशाचराः} %||6-83-20||

\twolineshloka
{तस्य तद्वचनं श्रुत्वा महापार्श्वोऽब्रवीद्वचः}
{बलाध्यक्षान्स्थितांस्तत्र बलं सन्त्वर्यतामिति} %||6-83-21||

\twolineshloka
{बलाध्यक्षास्तु संरब्धा राक्षसांस्तान्गृहाद्गृहात्}
{चोदयन्तः परिययुर्लङ्कां लघुपराक्रमाः} %||6-83-22||

\twolineshloka
{ततो मुहूर्तान्निष्पेतू राक्षसा भीमविक्रमाः}
{नर्दन्तो भीमवदना नानाप्रहरणैर्भुजैः} %||6-83-23||

\twolineshloka
{असिभिः पट्टसैः शूलैर्गलाभिर्मुसलैर्हलैः}
{शक्तिभिस्तीक्ष्णधाराभिर्महद्भिः कूटमुद्गरैः} %||6-83-24||

\twolineshloka
{यष्टिभिर्विमलैश्चक्रैर्निशितैश्च परश्वधैः}
{भिण्डिपालैः शतघ्नीभिरन्यैश्चापि वरायुधैः} %||6-83-25||

\twolineshloka
{अथानयन्बलाध्यक्षाश्चत्वारो रावणाज्ञया}
{द्रुतं सूतसमायुक्तं युक्ताष्टतुरगं रथम्} %||6-83-26||

\twolineshloka
{आरुरोह रथं दिव्यं दीप्यमानं स्वतेजसा}
{रावणः सत्त्वगाम्भीर्याद्दारयन्निव मेदिनीम्} %||6-83-27||

\twolineshloka
{रावणेनाभ्यनुज्ञातौ महापार्श्वमहोदरौ}
{विरूपाक्षश्च दुर्धर्षो रथानारुरुहुस्तदा} %||6-83-28||

\twolineshloka
{ते तु हृष्टा विनर्दन्तो भिन्दन्त इव मेदिनीम्}
{नादं घोरं विमुञ्चन्तो निर्ययुर्जयकाङ्क्षिणः} %||6-83-29||

\twolineshloka
{ततो युद्धाय तेजस्वी रक्षोगणबलैर्वृतः}
{निर्ययावुद्यतधनुः कालान्तकयमोमपः} %||6-83-30||

\twolineshloka
{ततः प्रजवनाश्वेन रथेन स महारथः}
{द्वारेण निर्ययौ तेन यत्र तौ रामलक्ष्मणौ} %||6-83-31||

\twolineshloka
{ततो नष्टप्रभः सूर्यो दिशश्च तिमिरावृताः}
{द्विजाश्च नेदुर्घोराश्च सञ्चचाल च मेदिनी} %||6-83-32||

\twolineshloka
{ववर्ष रुधिरं देवश्चस्खलुश्च तुरङ्गमाः}
{ध्वजाग्रे न्यपतद्गृध्रो विनेदुश्चाशिवं शिवाः} %||6-83-33||

\twolineshloka
{नयनं चास्फुरद्वामं सव्यो बाहुरकम्पत}
{विवर्णवदनश्चासीत्किञ्चिदभ्रश्यत स्वरः} %||6-83-34||

\twolineshloka
{ततो निष्पततो युद्धे दशग्रीवस्य रक्षसः}
{रणे निधनशंसीनि रूपाण्येतानि जज्ञिरे} %||6-83-35||

\twolineshloka
{अन्तरिक्षात्पपातोल्का निर्घातसमनिस्वना}
{विनेदुरशिवं गृध्रा वायसैरनुनादिताः} %||6-83-36||

\twolineshloka
{एतानचिन्तयन्घोरानुत्पातान्समुपस्थितान्}
{निर्ययौ रावणो मोहाद्वधार्थी कालचोदितः} %||6-83-37||

\twolineshloka
{तेषां तु रथघोषेण राक्षसानां महात्मनाम्}
{वानराणामपि चमूर्युद्धायैवाभ्यवर्तत} %||6-83-38||

\twolineshloka
{तेषां सुतुमुलं युद्धं बभूव कपिरक्षसाम्}
{अन्योन्यमाह्वयानानां क्रुद्धानां जयमिच्छताम्} %||6-83-39||

\twolineshloka
{ततः क्रुद्धो दशग्रीवः शरैः काञ्चनभूषणैः}
{वानराणामनीकेषु चकार कदनं महत्} %||6-83-40||

\threelineshloka
{निकृत्तशिरसः केचिद्रावणेन वलीमुखाः}
{निरुच्छ्वासा हताः केचित्केचित्पार्श्वेषु दारिताः}
{केचिद्विभिन्नशिरसः केचिच्चक्षुर्विवर्जिताः} %||6-83-41||

\fourlineindentedshloka
{दशाननः क्रोधविवृत्तनेत्रो}
{यतो यतोऽभ्येति रथेन सङ्ख्ये}
{ततस्ततस्तस्य शरप्रवेगम्}
{सोढुं न शेकुर्हरियूथपास्ते} %||6-84-42||


{॥इत्यार्षे श्रीमद्रामायणे वाल्मीकीये आदिकाव्ये युद्धकाण्डे त्र्यशीतितमः सर्गः॥}

\sect{चतुरशीतितमः सर्गः}

\twolineshloka
{तथा तैः कृत्तगात्रैस्तु दशग्रीवेण मार्गणैः}
{बभूव वसुधा तत्र प्रकीर्णा हरिभिर्वृता} %||6-84-1||

\twolineshloka
{रावणस्याप्रसह्यं तं शरसम्पातमेकतः}
{न शेकुः सहितुं दीप्तं पतङ्गा इव पावकम्} %||6-84-2||

\twolineshloka
{तेऽर्दिता निशितैर्बाणैः क्रोशन्तो विप्रदुद्रुवुः}
{पावकार्चिःसमाविष्टा दह्यमाना यथा गजाः} %||6-84-3||

\twolineshloka
{प्लवङ्गानामनीकानि महाभ्राणीव मारुतः}
{स ययौ समरे तस्मिन्विधमन्रावणः शरैः} %||6-84-4||

\twolineshloka
{कदनं तरसा कृत्वा राक्षसेन्द्रो वनौकसाम्}
{आससाद ततो युद्धे राघवं त्वरितस्तदा} %||6-84-5||

\twolineshloka
{सुग्रीवस्तान्कपीन्दृष्ट्वा भग्नान्विद्रवतो रणे}
{गुल्मे सुषेणं निक्षिप्य चक्रे युद्धे द्रुतं मनः} %||6-84-6||

\twolineshloka
{आत्मनः सदृशं वीरं स तं निक्षिप्य वानरम्}
{सुग्रीवोऽभिमुखः शत्रुं प्रतस्थे पादपायुधः} %||6-84-7||

\twolineshloka
{पार्श्वतः पृष्ठतश्चास्य सर्वे यूथाधिपाः स्वयम्}
{अनुजह्रुर्महाशैलान्विविधांश्च महाद्रुमान्} %||6-84-8||

\twolineshloka
{स नदन्युधि सुग्रीवः स्वरेण महता महान्}
{पातयन्विविधांश्चान्याञ्जघानोत्तमराक्षसान्} %||6-84-9||

\twolineshloka
{ममर्द च महाकायो राक्षसान्वानरेश्वरः}
{युगान्तसमये वायुः प्रवृद्धानगमानिव} %||6-84-10||

\twolineshloka
{राक्षसानामनीकेषु शैलवर्षं ववर्ष ह}
{अश्मवर्षं यथा मेघः पक्षिसङ्घेषु कानने} %||6-84-11||

\twolineshloka
{कपिराजविमुक्तैस्तैः शैलवर्षैस्तु राक्षसाः}
{विकीर्णशिरसः पेतुर्निकृत्ता इव पर्वताः} %||6-84-12||

\twolineshloka
{अथ सङ्क्षीयमाणेषु राक्षसेषु समन्ततः}
{सुग्रीवेण प्रभग्नेषु पतत्सु विनदत्सु च} %||6-84-13||

\twolineshloka
{विरूपाक्षः स्वकं नाम धन्वी विश्राव्य राक्षसः}
{रथादाप्लुत्य दुर्धर्षो गजस्कन्धमुपारुहत्} %||6-84-14||

\twolineshloka
{स तं द्विरदमारुह्य विरूपाक्षो महारथः}
{विनदन्भीमनिर्ह्रालं वानरानभ्यधावत} %||6-84-15||

\twolineshloka
{सुग्रीवे स शरान्घोरान्विससर्ज चमूमुखे}
{स्थापयामासा चोद्विग्नान्राक्षसान्सम्प्रहर्षयन्} %||6-84-16||

\twolineshloka
{सोऽतिविद्धः शितैर्बाणैः कपीन्द्रस्तेन रक्षसा}
{चुक्रोध च महाक्रोधो वधे चास्य मनो दधे} %||6-84-17||

\twolineshloka
{ततः पादपमुद्धृत्य शूरः सम्प्रधने हरिः}
{अभिपत्य जघानास्य प्रमुखे तं महागजम्} %||6-84-18||

\twolineshloka
{स तु प्रहाराभिहतः सुग्रीवेण महागजः}
{अपासर्पद्धनुर्मात्रं निषसाद ननाद च} %||6-84-19||

\twolineshloka
{गजात्तु मथितात्तूर्णमपक्रम्य स वीर्यवान्}
{राक्षसोऽभिमुखः शत्रुं प्रत्युद्गम्य ततः कपिम्} %||6-84-20||

\twolineshloka
{आर्षभं चर्मखड्गं च प्रगृह्य लघुविक्रमः}
{भर्त्सयन्निव सुग्रीवमाससाद व्यवस्थितम्} %||6-84-21||

\twolineshloka
{स हि तस्याभिसङ्क्रुद्धः प्रगृह्य महतीं शिलाम्}
{विरूपाक्षाय चिक्षेप सुग्रीवो जलदोपमाम्} %||6-84-22||

\twolineshloka
{स तां शिलामापतन्तीं दृष्ट्वा राक्षसपुङ्गवः}
{अपक्रम्य सुविक्रान्तः खड्गेन प्राहरत्तदा} %||6-84-23||

\twolineshloka
{तेन खड्गेन सङ्क्रुद्धः सुग्रीवस्य चमूमुखे}
{कवचं पातयामास स खड्गाभिहतोऽपतत्} %||6-84-24||

\twolineshloka
{स समुत्थाय पतितः कपिस्तस्य व्यसर्जयत्}
{तलप्रहारमशनेः समानं भीमनिस्वनम्} %||6-84-25||

\twolineshloka
{तलप्रहारं तद्रक्षः सुग्रीवेण समुद्यतम्}
{नैपुण्यान्मोचयित्वैनं मुष्टिनोरस्यताडयत्} %||6-84-26||

\twolineshloka
{ततस्तु सङ्क्रुद्धतरः सुग्रीवो वानरेश्वरः}
{मोक्षितं चात्मनो दृष्ट्वा प्रहारं तेन रक्षसा} %||6-84-27||

\twolineshloka
{स ददर्शान्तरं तस्य विरूपाक्षस्य वानरः}
{ततो न्यपातयत्क्रोधाच्छङ्खदेशे महातलम्} %||6-84-28||

\twolineshloka
{महेन्द्राशनिकल्पेन तलेनाभिहतः क्षितौ}
{पपात रुधिरक्लिन्नः शोणितं स समुद्वमन्} %||6-84-29||

\twolineshloka
{विवृत्तनयनं क्रोधात्सफेनरुधिराप्लुतम्}
{ददृशुस्ते विरूपाक्षं विरूपाक्षतरं कृतम्} %||6-84-30||

\twolineshloka
{स्फुरन्तं परिवर्जन्तं पार्श्वेन रुधिरोक्षितम्}
{करुणं च विनर्दान्तं ददृशुः कपयो रिपुम्} %||6-84-31||

\fourlineindentedshloka
{तथा तु तौ संयति सम्प्रयुक्तौ}
{तरस्विनौ वानरराक्षसानाम्}
{बलार्णवौ सस्वनतुः सभीमम्}
{महार्णवौ द्वाविव भिन्नवेलौ} %||6-84-32||

\fourlineindentedshloka
{विनाशितं प्रेक्ष्य विरूपनेत्रम्}
{महाबलं तं हरिपार्थिवेन}
{बलं समस्तं कपिराक्षसानाम्}
{उन्मत्तगङ्गाप्रतिमं बभूव} %||6-85-33||


{॥इत्यार्षे श्रीमद्रामायणे वाल्मीकीये आदिकाव्ये युद्धकाण्डे चतुरशीतितमः सर्गः॥}

\sect{पञ्चाशीतितमः सर्गः}

\twolineshloka
{हन्यमाने बले तूर्णमन्योन्यं ते महामृधे}
{सरसीव महाघर्मे सूपक्षीणे बभूवतुः} %||6-85-1||

\twolineshloka
{स्वबलस्य विघातेन विरूपाक्षवधेन च}
{बभूव द्विगुणं क्रुद्धो रावणो राक्षसाधिपः} %||6-85-2||

\twolineshloka
{प्रक्षीणं तु बलं दृष्ट्वा वध्यमानं वलीमुखैः}
{बभूवास्य व्यथा युद्धे प्रेक्ष्य दैवविपर्ययम्} %||6-85-3||

\twolineshloka
{उवाच च समीपस्थं महोदरमरिन्दमम्}
{अस्मिन्काले महाबाहो जयाशा त्वयि मे स्थिता} %||6-85-4||

\twolineshloka
{जहि शत्रुचमूं वीर दर्शयाद्य पराक्रमम्}
{भर्तृपिण्डस्य कालोऽयं निर्वेष्टुं साधु युध्यताम्} %||6-85-5||

\twolineshloka
{एवमुक्तस्तथेत्युक्त्वा राक्षसेन्द्रं महोदरः}
{प्रविवेशारिसेनां स पतङ्ग इव पावकम्} %||6-85-6||

\twolineshloka
{ततः स कदनं चक्रे वानराणां महाबलः}
{भर्तृवाक्येन तेजस्वी स्वेन वीर्येण चोदितः} %||6-85-7||

\twolineshloka
{प्रभग्नां समरे दृष्ट्वा वानराणां महाचमूम्}
{अभिदुद्राव सुग्रीवो महोदरमनन्तरम्} %||6-85-8||

\twolineshloka
{प्रगृह्य विपुलां घोरां महीधरसमां शिलाम्}
{चिक्षेप च महातेजास्तद्वधाय हरीश्वरः} %||6-85-9||

\twolineshloka
{तामापतन्तीं सहसा शिलां दृष्ट्वा महोदरः}
{असम्भ्रान्तस्ततो बाणैर्निर्बिभेद दुरासदाम्} %||6-85-10||

\twolineshloka
{रक्षसा तेन बाणौघैर्निकृत्ता सा सहस्रधा}
{निपपात शिला भूमौ गृध्रचक्रमिवाकुलम्} %||6-85-11||

\threelineshloka
{तां तु भिन्नां शिलां दृष्ट्वा सुग्रीवः क्रोधमूर्छितः}
{सालमुत्पाट्य चिक्षेप रक्षसे रणमूर्धनि}
{शरैश्च विददारैनं शूरः परपुरंजयः} %||6-85-12||

\threelineshloka
{स ददर्श ततः क्रुद्धः परिघं पतितं भुवि}
{आविध्य तु स तं दीप्तं परिघं तस्य दर्शयन्}
{परिघाग्रेण वेगेन जघानास्य हयोत्तमान्} %||6-85-13||

\twolineshloka
{तस्माद्धतहयाद्वीरः सोऽवप्लुत्य महारथात्}
{गदां जग्राह सङ्क्रुद्धो राक्षसोऽथ महोदरः} %||6-85-14||

\twolineshloka
{गदापरिघहस्तौ तौ युधि वीरौ समीयतुः}
{नर्दन्तौ गोवृषप्रख्यौ घनाविव सविद्युतौ} %||6-85-15||

\twolineshloka
{आजघान गदां तस्य परिघेण हरीश्वरः}
{पपात स गदोद्भिन्नः परिघस्तस्य भूतले} %||6-85-16||

\twolineshloka
{ततो जग्राह तेजस्वी सुग्रीवो वसुधातलात्}
{आयसं मुसलं घोरं सर्वतो हेमभूषितम्} %||6-85-17||

\twolineshloka
{तं समुद्यम्य चिक्षेप सोऽप्यन्यां व्याक्षिपद्गदाम्}
{भिन्नावन्योन्यमासाद्य पेततुर्धरणीतले} %||6-85-18||

\twolineshloka
{ततो भग्नप्रहरणौ मुष्टिभ्यां तौ समीयतुः}
{तेजोबलसमाविष्टौ दीप्ताविव हुताशनौ} %||6-85-19||

\twolineshloka
{जघ्नतुस्तौ तदान्योन्यं नेदतुश्च पुनः पुनः}
{तलैश्चान्योन्यमाहत्य पेततुर्धरणीतले} %||6-85-20||

\twolineshloka
{उत्पेततुस्ततस्तूर्णं जघ्नतुश्च परस्परम्}
{भुजैश्चिक्षेपतुर्वीरावन्योन्यमपराजितौ} %||6-85-21||

\twolineshloka
{आजहार तदा खड्गमदूरपरिवर्तिनम्}
{राक्षसश्चर्मणा सार्धं महावेगो महोदरः} %||6-85-22||

\twolineshloka
{तथैव च महाखड्गं चर्मणा पतितं सह}
{जग्राह वानरश्रेष्ठः सुग्रीवो वेगवत्तरः} %||6-85-23||

\twolineshloka
{तौ तु रोषपरीताङ्गौ नर्दन्तावभ्यधावताम्}
{उद्यतासी रणे हृष्टौ युधि शस्त्रविशारदौ} %||6-85-24||

\twolineshloka
{दक्षिणं मण्डलं चोभौ तौ तूर्णं सम्परीयतुः}
{अन्योन्यमभिसङ्क्रुद्धौ जये प्रणिहितावुभौ} %||6-85-25||

\twolineshloka
{स तु शूरो महावेगो वीर्यश्लाघी महोदरः}
{महाचर्मणि तं खड्गं पातयामास दुर्मतिः} %||6-85-26||

\twolineshloka
{लग्नमुत्कर्षतः खड्गं खड्गेन कपिकुञ्जरः}
{जहार सशिरस्त्राणं कुण्डलोपहितं शिरः} %||6-85-27||

\twolineshloka
{निकृत्तशिरसस्तस्य पतितस्य महीतले}
{तद्बलं राक्षसेन्द्रस्य दृष्ट्वा तत्र न तिष्ठति} %||6-85-28||

\twolineshloka
{हत्वा तं वानरैः सार्धं ननाद मुदितो हरिः}
{चुक्रोध च दशग्रीवो बभौ हृष्टश्च राघवः} %||6-86-29||


{॥इत्यार्षे श्रीमद्रामायणे वाल्मीकीये आदिकाव्ये युद्धकाण्डे पञ्चाशीतितमः सर्गः॥}

\sect{षडशीतितमः सर्गः}

\twolineshloka
{महोदरे तु निहते महापार्श्वो महाबलः}
{अङ्गदस्य चमूं भीमां क्षोभयामास सायकैः} %||6-86-1||

\twolineshloka
{स वानराणां मुख्यानामुत्तमाङ्गानि सर्वशः}
{पातयामास कायेभ्यः फलं वृन्तादिवानिलः} %||6-86-2||

\twolineshloka
{केषाञ्चिदिषुभिर्बाहून्स्कन्धांश्चिछेद राक्षसः}
{वानराणां सुसङ्क्रुद्धः पार्श्वं केषां व्यदारयत्} %||6-86-3||

\twolineshloka
{तेऽर्दिता बाणवर्षेण महापार्श्वेन वानराः}
{विषादविमुखाः सर्वे बभूवुर्गतचेतसः} %||6-86-4||

\twolineshloka
{निरीक्ष्य बलमुद्विग्नमङ्गदो राक्षसार्दितम्}
{वेगं चक्रे महाबाहुः समुद्र इव पर्वणि} %||6-86-5||

\twolineshloka
{आयसं परिघं गृह्य सूर्यरश्मिसमप्रभम्}
{समरे वानरश्रेष्ठो महापार्श्वे न्यपातयत्} %||6-86-6||

\twolineshloka
{स तु तेन प्रहारेण महापार्श्वो विचेतनः}
{ससूतः स्यन्दनात्तस्माद्विसंज्ञः प्रापतद्भुवि} %||6-86-7||

\twolineshloka
{सर्क्षराजस्तु तेजस्वी नीलाञ्जनचयोपमः}
{निष्पत्य सुमहावीर्यः स्वाद्यूथान्मेघसंनिभात्} %||6-86-8||

\twolineshloka
{प्रगृह्य गिरिशृङ्गाभां क्रुद्धः स विपुलां शिलाम्}
{अश्वाञ्जघान तरसा स्यन्दनं च बभञ्ज तम्} %||6-86-9||

\twolineshloka
{मुहूर्ताल्लब्धसंज्ञस्तु महापार्श्वो महाबलः}
{अङ्गदं बहुभिर्बाणैर्भूयस्तं प्रत्यविध्यत} %||6-86-10||

\twolineshloka
{जाम्बवन्तं त्रिभिर्बाणैराजघान स्तनान्तरे}
{ऋक्षराजं गवाक्षं च जघान बहुभिः शरैः} %||6-86-11||

\twolineshloka
{गवाक्षं जाम्बवन्तं च स दृष्ट्वा शरपीडितौ}
{जग्राह परिघं घोरमङ्गदः क्रोधमूर्छितः} %||6-86-12||

\twolineshloka
{तस्याङ्गदः प्रकुपितो राक्षसस्य तमायसम्}
{दूरस्थितस्य परिघं रविरश्मिसमप्रभम्} %||6-86-13||

\twolineshloka
{द्वाभ्यां भुजाभ्यां सङ्गृह्य भ्रामयित्वा च वेगवान्}
{महापार्श्वाय चिक्षेप वधार्थं वालिनः सुतः} %||6-86-14||

\twolineshloka
{स तु क्षिप्तो बलवता परिघस्तस्य रक्षसः}
{धनुश्च सशरं हस्ताच्छिरस्त्रं चाप्यपातयत्} %||6-86-15||

\twolineshloka
{तं समासाद्य वेगेन वालिपुत्रः प्रतापवान्}
{तलेनाभ्यहनत्क्रुद्धः कर्णमूले सकुण्डले} %||6-86-16||

\twolineshloka
{स तु क्रुद्धो महावेगो महापार्श्वो महाद्युतिः}
{करेणैकेन जग्राह सुमहान्तं परश्वधम्} %||6-86-17||

\twolineshloka
{तं तैलधौतं विमलं शैलसारमयं दृढम्}
{राक्षसः परमक्रुद्धो वालिपुत्रे न्यपातयत्} %||6-86-18||

\twolineshloka
{तेन वामांसफलके भृशं प्रत्यवपातितम्}
{अङ्गदो मोक्षयामास सरोषः स परश्वधम्} %||6-86-19||

\twolineshloka
{स वीरो वज्रसङ्काशमङ्गदो मुष्टिमात्मनः}
{संवर्तयन्सुसङ्क्रुद्धः पितुस्तुल्यपराक्रमः} %||6-86-20||

\twolineshloka
{राक्षसस्य स्तनाभ्याशे मर्मज्ञो हृदयं प्रति}
{इन्द्राशनिसमस्पर्शं स मुष्टिं विन्यपातयत्} %||6-86-21||

\twolineshloka
{तेन तस्य निपातेन राक्षसस्य महामृधे}
{पफाल हृदयं चाशु स पपात हतो भुवि} %||6-86-22||

\twolineshloka
{तस्मिन्निपतिते भूमौ तत्सैन्यं सम्प्रचुक्षुभे}
{अभवच्च महान्क्रोधः समरे रावणस्य तु} %||6-87-23||


{॥इत्यार्षे श्रीमद्रामायणे वाल्मीकीये आदिकाव्ये युद्धकाण्डे षडशीतितमः सर्गः॥}

\sect{सप्ताशीतितमः सर्गः}

\twolineshloka
{महोदरमहापार्श्वौ हतौ दृष्ट्वा तु राक्षसौ}
{तस्मिंश्च निहते वीरे विरूपाक्षे महाबले} %||6-87-1||

\twolineshloka
{आविवेश महान्क्रोधो रावणं तु महामृधे}
{सूतं सञ्चोदयामास वाक्यं चेदमुवाच ह} %||6-87-2||

\twolineshloka
{निहतानाममात्यानां रुद्धस्य नगरस्य च}
{दुःखमेषोऽपनेष्यामि हत्वा तौ रामलक्ष्मणौ} %||6-87-3||

\twolineshloka
{रामवृक्षं रणे हन्मि सीतापुष्पफलप्रदम्}
{प्रशाखा यस्य सुग्रीवो जाम्बवान्कुमुदो नलः} %||6-87-4||

\twolineshloka
{स दिशो दश घोषेण रथस्यातिरथो महान्}
{नादयन्प्रययौ तूर्णं राघवं चाभ्यवर्तत} %||6-87-5||

\twolineshloka
{पूरिता तेन शब्देन सनदीगिरिकानना}
{सञ्चचाल मही सर्वा सवराहमृगद्विपा} %||6-87-6||

\twolineshloka
{तामसं सुमहाघोरं चकारास्त्रं सुदारुणम्}
{निर्ददाह कपीन्सर्वांस्ते प्रपेतुः समन्ततः} %||6-87-7||

\twolineshloka
{तान्यनीकान्यनेकानि रावणस्य शरोत्तमैः}
{दृष्ट्वा भग्नानि शतशो राघवः पर्यवस्थितः} %||6-87-8||

\twolineshloka
{स ददर्श ततो रामं तिष्ठन्तमपराजितम्}
{लक्ष्मणेन सह भ्रात्रा विष्णुना वासवं यथा} %||6-87-9||

\twolineshloka
{आलिखन्तमिवाकाशमवष्टभ्य महद्धनुः}
{पद्मपत्रविशालाक्षं दीर्घबाहुमरिन्दमम्} %||6-87-10||

\twolineshloka
{वानरांश्च रणे भग्नानापतन्तं च रावणम्}
{समीक्ष्य राघवो हृष्टो मध्ये जग्राह कार्मुकम्} %||6-87-11||

\twolineshloka
{विस्फारयितुमारेभे ततः स धनुरुत्तमम्}
{महावेगं महानादं निर्भिन्दन्निव मेदिनीम्} %||6-87-12||

\twolineshloka
{तयोः शरपथं प्राप्य रावणो राजपुत्रयोः}
{स बभूव यथा राहुः समीपे शशिसूर्ययोः} %||6-87-13||

\twolineshloka
{रावणस्य च बाणौघै रामविस्फरितेन च}
{शब्देन राक्षसास्तेन पेतुश्च शतशस्तदा} %||6-87-14||

\twolineshloka
{तमिच्छन्प्रथमं योद्धुं लक्ष्मणो निशितैः शरैः}
{मुमोच धनुरायम्य शरानग्निशिखोपमान्} %||6-87-15||

\twolineshloka
{तान्मुक्तमात्रानाकाशे लक्ष्मणेन धनुष्मता}
{बाणान्बाणैर्महातेजा रावणः प्रत्यवारयत्} %||6-87-16||

\twolineshloka
{एकमेकेन बाणेन त्रिभिस्त्रीन्दशभिर्दश}
{लक्ष्मणस्य प्रचिच्छेद दर्शयन्पाणिलाघवम्} %||6-87-17||

\twolineshloka
{अभ्यतिक्रम्य सौमित्रिं रावणः समितिंजयः}
{आससाद ततो रामं स्थितं शैलमिवाचलम्} %||6-87-18||

\twolineshloka
{स सङ्ख्ये राममासाद्य क्रोधसंरक्तलोचनः}
{व्यसृजच्छरवर्षानि रावणो राघवोपरि} %||6-87-19||

\twolineshloka
{शरधारास्ततो रामो रावणस्य धनुश्च्युताः}
{दृष्ट्वैवापतिताः शीघ्रं भल्लाञ्जग्राह सत्वरम्} %||6-87-20||

\twolineshloka
{ताञ्शरौघांस्ततो भल्लैस्तीक्ष्णैश्चिच्छेद राघवः}
{दीप्यमानान्महावेगान्क्रुद्धानाशीविषानिव} %||6-87-21||

\twolineshloka
{राघवो रावणं तूर्णं रावणो राघवं तथा}
{अन्योन्यं विविधैस्तीक्ष्णैः शरैरभिववर्षतुः} %||6-87-22||

\twolineshloka
{चेरतुश्च चिरं चित्रं मण्डलं सव्यदक्षिणम्}
{बाणवेगान्समुदीक्ष्य समरेष्वपराजितौ} %||6-87-23||

\twolineshloka
{तयोर्भूतानि वित्रेषुर्युगपत्सम्प्रयुध्यतोः}
{रौद्रयोः सायकमुचोर्यमान्तकनिकाशयोः} %||6-87-24||

\twolineshloka
{सन्ततं विविधैर्बाणैर्बभूव गगनं तदा}
{घनैरिवातपापाये विद्युन्मालासमाकुलैः} %||6-87-25||

\twolineshloka
{गवाक्षितमिवाकाशं बभूव शूरवृष्टिभिः}
{महावेगैः सुतीक्ष्णाग्रैर्गृध्रपत्रैः सुवाजितैः} %||6-87-26||

\twolineshloka
{शरान्धकारं तौ भीमं चक्रतुः परमं तदा}
{गतेऽस्तं तपने चापि महामेघाविवोत्थितौ} %||6-87-27||

\twolineshloka
{बभूव तुमुलं युद्धमन्योन्यवधकाङ्क्षिणोः}
{अनासाद्यमचिन्त्यं च वृत्रवासवयोरिव} %||6-87-28||

\twolineshloka
{उभौ हि परमेष्वासावुभौ शस्त्रविशारदौ}
{उभौ चास्त्रविदां मुख्यावुभौ युद्धे विचेरतुः} %||6-87-29||

\twolineshloka
{उभौ हि येन व्रजतस्तेन तेन शरोर्मयः}
{ऊर्मयो वायुना विद्धा जग्मुः सागरयोरिव} %||6-87-30||

\twolineshloka
{ततः संसक्तहस्तस्तु रावणो लोकरावणः}
{नाराचमालां रामस्य ललाटे प्रत्यमुञ्चत} %||6-87-31||

\twolineshloka
{रौद्रचापप्रयुक्तां तां नीलोत्पलदलप्रभाम्}
{शिरसा धारयन्रामो न व्यथां प्रत्यपद्यत} %||6-87-32||

\twolineshloka
{अथ मन्त्रानपि जपन्रौद्रमस्त्रमुदीरयन्}
{शरान्भूयः समादाय रामः क्रोधसमन्वितः} %||6-87-33||

\twolineshloka
{मुमोच च महातेजाश्चापमायम्य वीर्यवान्}
{ताञ्शरान्राक्षसेन्द्राय चिक्षेपाच्छिन्नसायकः} %||6-87-34||

\twolineshloka
{ते महामेघसङ्काशे कवचे पतिताः शराः}
{अवध्ये राक्षसेन्द्रस्य न व्यथां जनयंस्तदा} %||6-87-35||

\twolineshloka
{पुनरेवाथ तं रामो रथस्थं राक्षसाधिपम्}
{ललाटे परमास्त्रेण सर्वास्त्रकुशलोऽभिनत्} %||6-87-36||

\twolineshloka
{ते भित्त्वा बाणरूपाणि पञ्चशीर्षा इवोरगाः}
{श्वसन्तो विविशुर्भूमिं रावणप्रतिकूलताः} %||6-87-37||

\twolineshloka
{निहत्य राघवस्यास्त्रं रावणः क्रोधमूर्छितः}
{आसुरं सुमहाघोरमन्यदस्त्रं समाददे} %||6-87-38||

\twolineshloka
{सिंहव्याघ्रमुखांश्चान्यान्कङ्ककाकमुखानपि}
{गृध्रश्येनमुखांश्चापि सृगालवदनांस्तथा} %||6-87-39||

\twolineshloka
{ईहामृगमुखांश्चान्यान्व्यादितास्यान्भयावहान्}
{पञ्चास्याँल्लेलिहानांश्च ससर्ज निशिताञ्शरान्} %||6-87-40||

\twolineshloka
{शरान्खरमुखांश्चान्यान्वराहमुखसंस्थितान्}
{श्वानकुक्कुटवक्त्रांश्च मकराशीविषाननान्} %||6-87-41||

\twolineshloka
{एतांश्चान्यांश्च मायाभिः ससर्ज निशिताञ्शरान्}
{रामं प्रति महातेजाः क्रुद्धः सर्प इव श्वसन्} %||6-87-42||

\twolineshloka
{आसुरेण समाविष्टः सोऽस्त्रेण रघुनन्दनः}
{ससर्जास्त्रं महोत्साहः पावकं पावकोपमः} %||6-87-43||

\twolineshloka
{अग्निदीप्तमुखान्बाणांस्तथा सूर्यमुखानपि}
{चन्द्रार्धचन्द्रवक्त्रांश्च धूमकेतुमुखानपि} %||6-87-44||

\twolineshloka
{ग्रहनक्षत्रवर्णांश्च महोल्कामुखसंस्थितान्}
{विद्युज्जिह्वोपमांश्चान्यान्ससर्ज निशिताञ्शरान्} %||6-87-45||

\twolineshloka
{ते रावणशरा घोरा राघवास्त्रसमाहताः}
{विलयं जग्मुराकाशे जग्मुश्चैव सहस्रशः} %||6-87-46||

\twolineshloka
{तदस्त्रं निहतं दृष्ट्वा रामेणाक्लिष्टकर्मणा}
{हृष्टा नेदुस्ततः सर्वे कपयः कामरूपिणः} %||6-88-47||


{॥इत्यार्षे श्रीमद्रामायणे वाल्मीकीये आदिकाव्ये युद्धकाण्डे सप्ताशीतितमः सर्गः॥}

\sect{अष्टाशीतितमः सर्गः}

\twolineshloka
{तस्मिन्प्रतिहतेऽस्त्रे तु रावणो राक्षसाधिपः}
{क्रोधं च द्विगुणं चक्रे क्रोधाच्चास्त्रमनन्तरम्} %||6-88-1||

\twolineshloka
{मयेन विहितं रौद्रमन्यदस्त्रं महाद्युतिः}
{उत्स्रष्टुं रावणो घोरं राघवाय प्रचक्रमे} %||6-88-2||

\twolineshloka
{ततः शूलानि निश्चेरुर्गदाश्च मुसलानि च}
{कार्मुकाद्दीप्यमानानि वज्रसाराणि सर्वशः} %||6-88-3||

\twolineshloka
{कूटमुद्गरपाशाश्च दीप्ताश्चाशनयस्तथा}
{निष्पेतुर्विविधास्तीक्ष्णा वाता इव युगक्षये} %||6-88-4||

\twolineshloka
{तदस्त्रं राघवः श्रीमानुत्तमास्त्रविदां वरः}
{जघान परमास्त्रेण गन्धर्वेण महाद्युतिः} %||6-88-5||

\twolineshloka
{तस्मिन्प्रतिहतेऽस्त्रे तु राघवेण महात्मना}
{रावणः क्रोधताम्राक्षः सौरमस्त्रमुदीरयत्} %||6-88-6||

\twolineshloka
{ततश्चक्राणि निष्पेतुर्भास्वराणि महान्ति च}
{कार्मुकाद्भीमवेगस्य दशग्रीवस्य धीमतः} %||6-88-7||

\twolineshloka
{तैरासीद्गगनं दीप्तं सम्पतद्भिरितस्ततः}
{पतद्भिश्च दिशो दीप्तैश्चन्द्रसूर्यग्रहैरिव} %||6-88-8||

\twolineshloka
{तानि चिच्छेद बाणौघैश्चक्राणि तु स राघवः}
{आयुधानि विचित्राणि रावणस्य चमूमुखे} %||6-88-9||

\twolineshloka
{तदस्त्रं तु हतं दृष्ट्वा रावणो राक्षसाधिपः}
{विव्याध दशभिर्बाणै रामं सर्वेषु मर्मसु} %||6-88-10||

\twolineshloka
{स विद्धो दशभिर्बाणैर्महाकार्मुकनिःसृतैः}
{रावणेन महातेजा न प्राकम्पत राघवः} %||6-88-11||

\twolineshloka
{ततो विव्याध गात्रेषु सर्वेषु समितिंजयः}
{राघवस्तु सुसङ्क्रुद्धो रावणं बहुभिः शरैः} %||6-88-12||

\twolineshloka
{एतस्मिन्नन्तरे क्रुद्धो राघवस्यानुजो बली}
{लक्ष्मणः सायकान्सप्त जग्राह परवीरहा} %||6-88-13||

\twolineshloka
{तैः सायकैर्महावेगै रावणस्य महाद्युतिः}
{ध्वजं मनुष्यशीर्षं तु तस्य चिच्छेद नैकधा} %||6-88-14||

\twolineshloka
{सारथेश्चापि बाणेन शिरो ज्वलितकुण्डलम्}
{जहार लक्ष्मणः श्रीमान्नैरृतस्य महाबलः} %||6-88-15||

\twolineshloka
{तस्य बाणैश्च चिच्छेद धनुर्गजकरोपमम्}
{लक्ष्मणो राक्षसेन्द्रस्य पञ्चभिर्निशितैः शरैः} %||6-88-16||

\twolineshloka
{नीलमेघनिभांश्चास्य सदश्वान्पर्वतोपमान्}
{जघानाप्लुत्य गदया रावणस्य विभीषणः} %||6-88-17||

\twolineshloka
{हताश्वाद्वेगवान्वेगादवप्लुत्य महारथात्}
{क्रोधमाहारयत्तीव्रं भ्रातरं प्रति रावणः} %||6-88-18||

\twolineshloka
{ततः शक्तिं महाशक्तिर्दीप्तां दीप्ताशनीमिव}
{विभीषणाय चिक्षेप राक्षसेन्द्रः प्रतापवान्} %||6-88-19||

\twolineshloka
{अप्राप्तामेव तां बाणैस्त्रिभिश्चिच्छेद लक्ष्मणः}
{अथोदतिष्ठत्संनादो वानराणां तदा रणे} %||6-88-20||

\twolineshloka
{सा पपात त्रिधा छिन्ना शक्तिः काञ्चनमालिनी}
{सविस्फुलिङ्गा ज्वलिता महोल्केव दिवश्च्युता} %||6-88-21||

\twolineshloka
{ततः सम्भाविततरां कालेनापि दुरासदाम्}
{जग्राह विपुलां शक्तिं दीप्यमानां स्वतेजसा} %||6-88-22||

\twolineshloka
{सा वेगिना बलवता रावणेन दुरात्मना}
{जज्वाल सुमहाघोरा शक्राशनिसमप्रभा} %||6-88-23||

\twolineshloka
{एतस्मिन्नन्तरे वीरो लक्ष्मणस्तं विभीषणम्}
{प्राणसंशयमापन्नं तूर्णमेवाभ्यपद्यत} %||6-88-24||

\twolineshloka
{तं विमोक्षयितुं वीरश्चापमायम्य लक्ष्मणः}
{रावणं शक्तिहस्तं तं शरवर्षैरवाकिरत्} %||6-88-25||

\twolineshloka
{कीर्यमाणः शरौघेण विसृष्टेन महात्मना}
{न प्रहर्तुं मनश्चक्रे विमुखीकृतविक्रमः} %||6-88-26||

\twolineshloka
{मोक्षितं भ्रातरं दृष्ट्वा लक्ष्मणेन स रावणः}
{लक्ष्मणाभिमुखस्तिष्ठन्निदं वचनमब्रवीत्} %||6-88-27||

\twolineshloka
{मोक्षितस्ते बलश्लाघिन्यस्मादेवं विभीषणः}
{विमुच्य राक्षसं शक्तिस्त्वयीयं विनिपात्यते} %||6-88-28||

\twolineshloka
{एषा ते हृदयं भित्त्वा शक्तिर्लोहितलक्षणा}
{मद्बाहुपरिघोत्सृष्टा प्राणानादाय यास्यति} %||6-88-29||

\twolineshloka
{इत्येवमुक्त्वा तां शक्तिमष्टघण्टां महास्वनाम्}
{मयेन मायाविहिताममोघां शत्रुघातिनीम्} %||6-88-30||

\twolineshloka
{लक्ष्मणाय समुद्दिश्य ज्वलन्तीमिव तेजसा}
{रावणः परमक्रुद्धश्चिक्षेप च ननाद च} %||6-88-31||

\twolineshloka
{सा क्षिप्ता भीमवेगेन शक्राशनिसमस्वना}
{शक्तिरभ्यपतद्वेगाल्लक्ष्मणं रणमूर्धनि} %||6-88-32||

\twolineshloka
{तामनुव्याहरच्छक्तिमापतन्तीं स राघवः}
{स्वस्त्यस्तु लक्ष्मणायेति मोघा भव हतोद्यमा} %||6-88-33||

\twolineshloka
{न्यपतत्सा महावेगा लक्ष्मणस्य महोरसि}
{जिह्वेवोरगराजस्य दीप्यमाना महाद्युतिः} %||6-88-34||

\twolineshloka
{ततो रावणवेगेन सुदूरमवगाढया}
{शक्त्या निर्भिन्नहृदयः पपात भुवि लक्ष्मणः} %||6-88-35||

\twolineshloka
{तदवस्थं समीपस्थो लक्ष्मणं प्रेक्ष्य राघवः}
{भ्रातृस्नेहान्महातेजा विषण्णहृदयोऽभवत्} %||6-88-36||

\twolineshloka
{स मुहूर्तमनुध्याय बाष्पव्याकुललोचनः}
{बभूव संरब्धतरो युगान्त इव पावकः} %||6-88-37||

\twolineshloka
{न विषादस्य कालोऽयमिति सञ्चिन्त्य राघवः}
{चक्रे सुतुमुलं युद्धं रावणस्य वधे धृतः} %||6-88-38||

\twolineshloka
{स ददर्श ततो रामः शक्त्या भिन्नं महाहवे}
{लक्ष्मणं रुधिरादिग्धं सपन्नगमिवाचलम्} %||6-88-39||

\threelineshloka
{तामपि प्रहितां शक्तिं रावणेन बलीयसा}
{यत्नतस्ते हरिश्रेष्ठा न शेकुरवमर्दितुम्}
{अर्दिताश्चैव बाणौघैः क्षिप्रहस्तेन रक्षसा} %||6-88-40||

\threelineshloka
{सौमित्रिं सा विनिर्भिद्य प्रविष्टा धरणीतलम्}
{तां कराभ्यां परामृश्य रामः शक्तिं भयावहाम्}
{बभञ्ज समरे क्रुद्धो बलवद्विचकर्ष च} %||6-88-41||

\twolineshloka
{तस्य निष्कर्षतः शक्तिं रावणेन बलीयसा}
{शराः सर्वेषु गात्रेषु पातिता मर्मभेदिनः} %||6-88-42||

\threelineshloka
{अचिन्तयित्वा तान्बाणान्समाश्लिष्य च लक्ष्मणम्}
{अब्रवीच्च हनूमन्तं सुग्रीवं चैव राघवः}
{लक्ष्मणं परिवार्येह तिष्ठध्वं वानरोत्तमाः} %||6-88-43||

\threelineshloka
{पराक्रमस्य कालोऽयं सम्प्राप्तो मे चिरेप्सितः}
{पापात्मायं दशग्रीवो वध्यतां पापनिश्चयः}
{काङ्क्षितः स्तोककस्येव घर्मान्ते मेघदर्शनम्} %||6-88-44||

\twolineshloka
{अस्मिन्मुहूर्ते नचिरात्सत्यं प्रतिशृणोमि वः}
{अरावणमरामं वा जगद्द्रक्ष्यथ वानराः} %||6-88-45||

\twolineshloka
{राज्यनाशं वने वासं दण्डके परिधावनम्}
{वैदेह्याश्च परामर्शं रक्षोभिश्च समागमम्} %||6-88-46||

\twolineshloka
{प्राप्तं दुःखं महद्घोरं क्लेशं च निरयोपमम्}
{अद्य सर्वमहं त्यक्ष्ये हत्वा तं रावणं रणे} %||6-88-47||

\twolineshloka
{यदर्थं वानरं सैन्यं समानीतमिदं मया}
{सुग्रीवश्च कृतो राज्ये निहत्वा वालिनं रणे} %||6-88-48||

\twolineshloka
{यदर्थं सागरः क्रान्तः सेतुर्बद्धश्च सागरे}
{सोऽयमद्य रणे पापश्चक्षुर्विषयमागतः} %||6-88-49||

\twolineshloka
{चक्षुर्विषयमागम्य नायं जीवितुमर्हति}
{दृष्टिं दृष्टिविषस्येव सर्पस्य मम रावणः} %||6-88-50||

\twolineshloka
{स्वस्थाः पश्यत दुर्धर्षा युद्धं वानरपुङ्गवाः}
{आसीनाः पर्वताग्रेषु ममेदं रावणस्य च} %||6-88-51||

\twolineshloka
{अद्य रामस्य रामत्वं पश्यन्तु मम संयुगे}
{त्रयो लोकाः सगन्धर्वाः सदेवाः सर्षिचारणाः} %||6-88-52||

\twolineshloka
{अद्य कर्म करिष्यामि यल्लोकाः सचराचराः}
{सदेवाः कथयिष्यन्ति यावद्भूमिर्धरिष्यति} %||6-88-53||

\twolineshloka
{एवमुक्त्वा शितैर्बाणैस्तप्तकाञ्चनभूषणैः}
{आजघान दशग्रीवं रणे रामः समाहितः} %||6-88-54||

\twolineshloka
{अथ प्रदीप्तैर्नाराचैर्मुसलैश्चापि रावणः}
{अभ्यवर्षत्तदा रामं धाराभिरिव तोयदः} %||6-88-55||

\twolineshloka
{रामरावणमुक्तानामन्योन्यमभिनिघ्नताम्}
{शराणां च शराणां च बभूव तुमुलः स्वनः} %||6-88-56||

\twolineshloka
{ते भिन्नाश्च विकीर्णाश्च रामरावणयोः शराः}
{अन्तरिक्षात्प्रदीप्ताग्रा निपेतुर्धरणीतले} %||6-88-57||

\twolineshloka
{तयोर्ज्यातलनिर्घोषो रामरावणयोर्महान्}
{त्रासनः सर्वबूतानां स बभूवाद्भुतोपमः} %||6-88-58||

\fourlineindentedshloka
{स कीर्यमाणः शरजालवृष्टिभिर्-}
{महात्मना दीप्तधनुष्मतार्दितः}
{भयात्प्रदुद्राव समेत्य रावणो}
{यथानिलेनाभिहतो बलाहकः} %||6-89-59||


{॥इत्यार्षे श्रीमद्रामायणे वाल्मीकीये आदिकाव्ये युद्धकाण्डे अष्टाशीतितमः सर्गः॥}

\sect{एकोननवतितमः सर्गः}

\twolineshloka
{स दत्त्वा तुमुलं युद्धं रावणस्य दुरात्मनः}
{विसृजन्नेव बाणौघान्सुषेणं वाक्यमब्रवीत्} %||6-89-1||

\twolineshloka
{एष रावणवेगेन लक्ष्मणः पतितः क्षितौ}
{सर्पवद्वेष्टते वीरो मम शोकमुदीरयन्} %||6-89-2||

\twolineshloka
{शोणितार्द्रमिमं वीरं प्राणैरिष्टतरं मम}
{पश्यतो मम का शक्तिर्योद्धुं पर्याकुलात्मनः} %||6-89-3||

\twolineshloka
{अयं स समरश्लाघी भ्राता मे शुभलक्षणः}
{यदि पञ्चत्वमापन्नः प्राणैर्मे किं सुखेन वा} %||6-89-4||

\threelineshloka
{लज्जतीव हि मे वीर्यं भ्रश्यतीव कराद्धनुः}
{सायका व्यवसीदन्ति दृष्टिर्बाष्पवशं गता}
{चिन्ता मे वर्धते तीव्रा मुमूर्षा चोपजायते} %||6-89-5||

\twolineshloka
{भ्रातरं निहतं दृष्ट्वा रावणेन दुरात्मना}
{परं विषादमापन्नो विललापाकुलेन्द्रियः} %||6-89-6||

\twolineshloka
{न हि युद्धेन मे कार्यं नैव प्राणैर्न सीतया}
{भ्रातरं निहतं दृष्ट्वा लक्ष्मणं रणपांसुषु} %||6-89-7||

\twolineshloka
{किं मे राज्येन किं प्राणैर्युद्धे कार्यं न विद्यते}
{यत्रायं निहतः शेते रणमूर्धनि लक्ष्मणः} %||6-89-8||

\twolineshloka
{राममाश्वासयन्वीरः सुषेणो वाक्यमब्रवीत्}
{न मृतोऽयं महाबाहुर्लक्ष्मणो लक्ष्मिवर्धनः} %||6-89-9||

\twolineshloka
{न चास्य विकृतं वक्त्रं नापि श्यामं न निष्प्रभम्}
{सुप्रभं च प्रसन्नं च मुखमस्याभिलक्ष्यते} %||6-89-10||

\threelineshloka
{पद्मरक्ततलौ हस्तौ सुप्रसन्ने च लोचने}
{एवं न विद्यते रूपं गतासूनां विशां पते}
{मां विषादं कृथा वीर सप्राणोऽयमरिन्दम} %||6-89-11||

\twolineshloka
{आख्यास्यते प्रसुप्तस्य स्रस्तगात्रस्य भूतले}
{सोच्छ्वासं हृदयं वीर कम्पमानं मुहुर्मुहुः} %||6-89-12||

\twolineshloka
{एवमुक्त्वा तु वाक्यज्ञः सुषेणो राघवं वचः}
{समीपस्थमुवाचेदं हनूमन्तमभित्वरन्} %||6-89-13||

\twolineshloka
{सौम्य शीघ्रमितो गत्वा शैलमोषधिपर्वतम्}
{पूर्वं हि कथितो योऽसौ वीर जाम्बवता शुभः} %||6-89-14||

\twolineshloka
{दक्षिणे शिखरे तस्य जातामोषधिमानय}
{विशल्यकरणी नाम विशल्यकरणीं शुभाम्} %||6-89-15||

\threelineshloka
{सौवर्णकरणीं चापि तथा संजीवनीमपि}
{सन्धानकरणीं चापि गत्वा शीघ्रमिहानय}
{संजीवनार्थं वीरस्य लक्ष्मणस्य महात्मनः} %||6-89-16||

\twolineshloka
{इत्येवमुक्तो हनुमान्गत्वा चौषधिपर्वतम्}
{चिन्तामभ्यगमच्छ्रीमानजानंस्ता महौषधीः} %||6-89-17||

\twolineshloka
{तस्य बुद्धिः समुत्पन्ना मारुतेरमितौजसः}
{इदमेव गमिष्यामि गृहीत्वा शिखरं गिरेः} %||6-89-18||

\twolineshloka
{अगृह्य यदि गच्छामि विशल्यकरणीमहम्}
{कालात्ययेन दोषः स्याद्वैक्लव्यं च महद्भवेत्} %||6-89-19||

\twolineshloka
{इति सञ्चिन्त्य हनुमान्गत्वा क्षिप्रं महाबलः}
{उत्पपात गृहीत्वा तु हनूमाञ्शिखरं गिरेः} %||6-89-20||

\twolineshloka
{ओषधीर्नावगछामि ता अहं हरिपुङ्गव}
{तदिदं शिखरं कृत्स्नं गिरेस्तस्याहृतं मया} %||6-89-21||

\twolineshloka
{एवं कथयमानं तं प्रशस्य पवनात्मजम्}
{सुषेणो वानरश्रेष्ठो जग्राहोत्पाट्य चौषधीः} %||6-89-22||

\twolineshloka
{ततः सङ्क्षोदयित्वा तामोषधिं वानरोत्तमः}
{लक्ष्मणस्य ददौ नस्तः सुषेणः सुमहाद्युतिः} %||6-89-23||

\twolineshloka
{सशल्यः स समाघ्राय लक्ष्मणः परवीरहा}
{विशल्यो विरुजः शीघ्रमुदतिष्ठन्महीतलात्} %||6-89-24||

\twolineshloka
{समुत्थितं ते हरयो भूतलात्प्रेक्ष्य लक्ष्मणम्}
{साधु साध्विति सुप्रीताः सुषेणं प्रत्यपूजयन्} %||6-89-25||

\twolineshloka
{एह्येहीत्यब्रवीद्रामो लक्ष्मणं परवीरहा}
{सस्वजे स्नेहगाढं च बाष्पपर्याकुलेक्षणः} %||6-89-26||

\twolineshloka
{अब्रवीच्च परिष्वज्य सौमित्रिं राघवस्तदा}
{दिष्ट्या त्वां वीर पश्यामि मरणात्पुनरागतम्} %||6-89-27||

\twolineshloka
{न हि मे जीवितेनार्थः सीतया च जयेन वा}
{को हि मे जीवितेनार्थस्त्वयि पञ्चत्वमागते} %||6-89-28||

\twolineshloka
{इत्येवं वदतस्तस्य राघवस्य महात्मनः}
{खिन्नः शिथिलया वाचा लक्ष्मणो वाक्यमब्रवीत्} %||6-89-29||

\twolineshloka
{तां प्रतिज्ञां प्रतिज्ञाय पुरा सत्यपराक्रम}
{लघुः कश्चिदिवासत्त्वो नैवं वक्तुमिहार्हसि} %||6-89-30||

\twolineshloka
{न प्रतिज्ञां हि कुर्वन्ति वितथां साधवोऽनघ}
{लक्षणं हि महत्त्वस्य प्रतिज्ञापरिपालनम्} %||6-89-31||

\twolineshloka
{नैराश्यमुपगन्तुं ते तदलं मत्कृतेऽनघ}
{वधेन रावणस्याद्य प्रतिज्ञामनुपालय} %||6-89-32||

\twolineshloka
{न जीवन्यास्यते शत्रुस्तव बाणपथं गतः}
{नर्दतस्तीक्ष्णदंष्ट्रस्य सिंहस्येव महागजः} %||6-89-33||

\twolineshloka
{अहं तु वधमिच्छामि शीघ्रमस्य दुरात्मनः}
{यावदस्तं न यात्येष कृतकर्मा दिवाकरः} %||6-90-34||


{॥इत्यार्षे श्रीमद्रामायणे वाल्मीकीये आदिकाव्ये युद्धकाण्डे एकोननवतितमः सर्गः॥}

\sect{नवतितमः सर्गः}

\twolineshloka
{लक्ष्मणेन तु तद्वाक्यमुक्तं श्रुत्वा स राघवः}
{रावणाय शरान्घोरान्विससर्ज चमूमुखे} %||6-90-1||

\twolineshloka
{दशग्रीवो रथस्थस्तु रामं वज्रोपमैः शरैः}
{आजघान महाघोरैर्धाराभिरिव तोयदः} %||6-90-2||

\twolineshloka
{दीप्तपावकसङ्काशैः शरैः काञ्चनभूषणैः}
{निर्बिभेद रणे रामो दशग्रीवं समाहितः} %||6-90-3||

\twolineshloka
{भूमिस्थितस्य रामस्य रथस्थस्य च रक्षसः}
{न समं युद्धमित्याहुर्देवगन्धर्वदानवाः} %||6-90-4||

\twolineshloka
{ततः काञ्चनचित्राङ्गः किङ्किणीशतभूषितः}
{तरुणादित्यसङ्काशो वैदूर्यमयकूबरः} %||6-90-5||

\twolineshloka
{सदश्वैः काञ्चनापीडैर्युक्तः श्वेतप्रकीर्णकैः}
{हरिभिः सूर्यसङ्काशैर्हेमजालविभूषितैः} %||6-90-6||

\twolineshloka
{रुक्मवेणुध्वजः श्रीमान्देवराजरथो वरः}
{अभ्यवर्तत काकुत्स्थमवतीर्य त्रिविष्टपात्} %||6-90-7||

\twolineshloka
{अब्रवीच्च तदा रामं सप्रतोदो रथे स्थितः}
{प्राञ्जलिर्मातलिर्वाक्यं सहस्राक्षस्य सारथिः} %||6-90-8||

\twolineshloka
{सहस्राक्षेण काकुत्स्थ रथोऽयं विजयाय ते}
{दत्तस्तव महासत्त्व श्रीमाञ्शत्रुनिबर्हणः} %||6-90-9||

\twolineshloka
{इदमैन्द्रं महच्चापं कवचं चाग्निसंनिभम्}
{शराश्चादित्यसङ्काशाः शक्तिश्च विमला शिताः} %||6-90-10||

\twolineshloka
{आरुह्येमं रथं वीर राक्षसं जहि रावणम्}
{मया सारथिना राम महेन्द्र इव दानवान्} %||6-90-11||

\twolineshloka
{इत्युक्तः स परिक्रम्य रथं तमभिवाद्य च}
{आरुरोह तदा रामो लोकाँल्लक्ष्म्या विराजयन्} %||6-90-12||

\twolineshloka
{तद्बभूवाद्भुतं युद्धं द्वैरथं लोमहर्षणम्}
{रामस्य च महाबाहो रावणस्य च रक्षसः} %||6-90-13||

\twolineshloka
{स गान्धर्वेण गान्धर्वं दैवं दैवेन राघवः}
{अस्त्रं राक्षसराजस्य जघान परमास्त्रवित्} %||6-90-14||

\twolineshloka
{अस्त्रं तु परमं घोरं राक्षसं राकसाधिपः}
{ससर्ज परमक्रुद्धः पुनरेव निशाचरः} %||6-90-15||

\twolineshloka
{ते रावणधनुर्मुक्ताः शराः काञ्चनभूषणाः}
{अभ्यवर्तन्त काकुत्स्थं सर्पा भूत्वा महाविषाः} %||6-90-16||

\twolineshloka
{ते दीप्तवदना दीप्तं वमन्तो ज्वलनं मुखैः}
{राममेवाभ्यवर्तन्त व्यादितास्या भयानकाः} %||6-90-17||

\twolineshloka
{तैर्वासुकिसमस्पर्शैर्दीप्तभोगैर्महाविषैः}
{दिशश्च सन्तताः सर्वाः प्रदिशश्च समावृताः} %||6-90-18||

\twolineshloka
{तान्दृष्ट्वा पन्नगान्रामः समापतत आहवे}
{अस्त्रं गारुत्मतं घोरं प्रादुश्चक्रे भयावहम्} %||6-90-19||

\twolineshloka
{ते राघवधनुर्मुक्ता रुक्मपुङ्खाः शिखिप्रभाः}
{सुपर्णाः काञ्चना भूत्वा विचेरुः सर्पशत्रवः} %||6-90-20||

\twolineshloka
{ते तान्सर्वाञ्शराञ्जघ्नुः सर्परूपान्महाजवान्}
{सुपर्णरूपा रामस्य विशिखाः कामरूपिणः} %||6-90-21||

\twolineshloka
{अस्त्रे प्रतिहते क्रुद्धो रावणो राक्षसाधिपः}
{अभ्यवर्षत्तदा रामं घोराभिः शरवृष्टिभिः} %||6-90-22||

\twolineshloka
{ततः शरसहस्रेण राममक्लिष्टकारिणम्}
{अर्दयित्वा शरौघेण मातलिं प्रत्यविध्यत} %||6-90-23||

\twolineshloka
{पातयित्वा रथोपस्थे रथात्केतुं च काञ्चनम्}
{ऐन्द्रानभिजघानाश्वाञ्शरजालेन रावणः} %||6-90-24||

\twolineshloka
{विषेदुर्देवगन्धर्वा दानवाश्चारणैः सह}
{राममार्तं तदा दृष्ट्वा सिद्धाश्च परमर्षयः} %||6-90-25||

\twolineshloka
{व्यथिता वानरेन्द्राश्च बभूवुः सविभीषणाः}
{रामचन्द्रमसं दृष्ट्वा ग्रस्तं रावणराहुणा} %||6-90-26||

\twolineshloka
{प्राजापत्यं च नक्षत्रं रोहिणीं शशिनः प्रियाम्}
{समाक्रम्य बुधस्तस्थौ प्रजानामशुभावहः} %||6-90-27||

\twolineshloka
{सधूमपरिवृत्तोर्मिः प्रज्वलन्निव सागरः}
{उत्पपात तदा क्रुद्धः स्पृशन्निव दिवाकरम्} %||6-90-28||

\twolineshloka
{शस्त्रवर्णः सुपरुषो मन्दरश्मिर्दिवाकरः}
{अदृश्यत कबन्धाङ्गः संसक्तो धूमकेतुना} %||6-90-29||

\twolineshloka
{कोसलानां च नक्षत्रं व्यक्तमिन्द्राग्निदैवतम्}
{आक्रम्याङ्गारकस्तस्थौ विशाखामपि चाम्बरे} %||6-90-30||

\twolineshloka
{दशास्यो विंशतिभुजः प्रगृहीतशरासनः}
{अदृश्यत दशग्रीवो मैनाक इव पर्वतः} %||6-90-31||

\twolineshloka
{निरस्यमानो रामस्तु दशग्रीवेण रक्षसा}
{नाशकदभिसन्धातुं सायकान्रणमूर्धनि} %||6-90-32||

\twolineshloka
{स कृत्वा भ्रुकुटीं क्रुद्धः किञ्चित्संरक्तलोचनः}
{जगाम सुमहाक्रोधं निर्दहन्निव चक्षुषा} %||6-91-33||


{॥इत्यार्षे श्रीमद्रामायणे वाल्मीकीये आदिकाव्ये युद्धकाण्डे नवतितमः सर्गः॥}

\sect{एकनवतितमः सर्गः}

\twolineshloka
{तस्य क्रुद्धस्य वदनं दृष्ट्वा रामस्य धीमतः}
{सर्वभूतानि वित्रेषुः प्राकम्पत च मेदिनी} %||6-91-1||

\twolineshloka
{सिंहशार्दूलवाञ्शैलः सञ्चचालाचलद्रुमः}
{बभूव चापि क्षुभितः समुद्रः सरितां पतिः} %||6-91-2||

\twolineshloka
{खगाश्च खरनिर्घोषा गगने परुषस्वनाः}
{औत्पातिका विनर्दन्तः समन्तात्परिचक्रमुः} %||6-91-3||

\twolineshloka
{रामं दृष्ट्वा सुसङ्क्रुद्धमुत्पातांश्च सुदारुणान्}
{वित्रेषुः सर्वभूतानि रावणस्याविशद्भयम्} %||6-91-4||

\twolineshloka
{विमानस्थास्तदा देवा गन्धर्वाश्च महोरगाः}
{ऋषिदानवदैत्याश्च गरुत्मन्तश्च खेचराः} %||6-91-5||

\twolineshloka
{ददृशुस्ते तदा युद्धं लोकसंवर्तसंस्थितम्}
{नानाप्रहरणैर्भीमैः शूरयोः सम्प्रयुध्यतोः} %||6-91-6||

\twolineshloka
{ऊचुः सुरासुराः सर्वे तदा विग्रहमागताः}
{प्रेक्षमाणा महायुद्धं वाक्यं भक्त्या प्रहृष्टवत्} %||6-91-7||

\twolineshloka
{दशग्रीवं जयेत्याहुरसुराः समवस्थिताः}
{देवा राममथोचुस्ते त्वं जयेति पुनः पुनः} %||6-91-8||

\twolineshloka
{एतस्मिन्नन्तरे क्रोधाद्राघवस्य स रावणः}
{प्रहर्तुकामो दुष्टात्मा स्पृशन्प्रहरणं महत्} %||6-91-9||

\twolineshloka
{वज्रसारं महानादं सर्वशत्रुनिबर्हणम्}
{शैलशृङ्गनिभैः कूटैश्चितं दृष्टिभयावहम्} %||6-91-10||

\twolineshloka
{सधूममिव तीक्ष्णाग्रं युगान्ताग्निचयोपमम्}
{अतिरौद्रमनासाद्यं कालेनापि दुरासदम्} %||6-91-11||

\twolineshloka
{त्रासनं सर्वभूतानां दारणं भेदनं तथा}
{प्रदीप्त इव रोषेण शूलं जग्राह रावणः} %||6-91-12||

\twolineshloka
{तच्छूलं परमक्रुद्धो मध्ये जग्राह वीर्यवान्}
{अनेकैः समरे शूरै राक्षसैः परिवारितः} %||6-91-13||

\twolineshloka
{समुद्यम्य महाकायो ननाद युधि भैरवम्}
{संरक्तनयनो रोषात्स्वसैन्यमभिहर्षयन्} %||6-91-14||

\twolineshloka
{पृथिवीं चान्तरिक्षं च दिशश्च प्रदिशस्तथा}
{प्राकम्पयत्तदा शब्दो राक्षसेन्द्रस्य दारुणः} %||6-91-15||

\twolineshloka
{अतिनादस्य नादेन तेन तस्य दुरात्मनः}
{सर्वभूतानि वित्रेषुः सागरश्च प्रचुक्षुभे} %||6-91-16||

\twolineshloka
{स गृहीत्वा महावीर्यः शूलं तद्रावणो महत्}
{विनद्य सुमहानादं रामं परुषमब्रवीत्} %||6-91-17||

\twolineshloka
{शूलोऽयं वज्रसारस्ते राम रोषान्मयोद्यतः}
{तव भ्रातृसहायस्य सद्यः प्राणान्हरिष्यति} %||6-91-18||

\twolineshloka
{रक्षसामद्य शूराणां निहतानां चमूमुखे}
{त्वां निहत्य रणश्लाघिन्करोमि तरसा समम्} %||6-91-19||

\twolineshloka
{तिष्ठेदानीं निहन्मि त्वामेष शूलेन राघव}
{एवमुक्त्वा स चिक्षेप तच्छूलं राक्षसाधिपः} %||6-91-20||

\twolineshloka
{आपतन्तं शरौघेण वारयामास राघवः}
{उत्पतन्तं युगान्ताग्निं जलौघैरिव वासवः} %||6-91-21||

\twolineshloka
{निर्ददाह स तान्बाणान्रामकार्मुकनिःसृतान्}
{रावणस्य महाशूलः पतङ्गानिव पावकः} %||6-91-22||

\twolineshloka
{तान्दृष्ट्वा भस्मसाद्भूताञ्शूलसंस्पर्शचूर्णितान्}
{सायकानन्तरिक्षस्थान्राघवः क्रोधमाहरत्} %||6-91-23||

\twolineshloka
{स तां मातलिनानीतां शक्तिं वासवनिर्मिताम्}
{जग्राह परमक्रुद्धो राघवो रघुनन्दनः} %||6-91-24||

\twolineshloka
{सा तोलिता बलवता शक्तिर्घण्टाकृतस्वना}
{नभः प्रज्वालयामास युगान्तोल्केव सप्रभा} %||6-91-25||

\twolineshloka
{सा क्षिप्ता राक्षसेन्द्रस्य तस्मिञ्शूले पपात ह}
{भिन्नः शक्त्या महाञ्शूलो निपपात गतद्युतिः} %||6-91-26||

\twolineshloka
{निर्बिभेद ततो बाणैर्हयानस्य महाजवान्}
{रामस्तीक्ष्णैर्महावेगैर्वज्रकल्पैः शितैः शरैः} %||6-91-27||

\twolineshloka
{निर्बिभेदोरसि तदा रावणं निशितैः शरैः}
{राघवः परमायत्तो ललाटे पत्रिभिस्त्रिभिः} %||6-91-28||

\twolineshloka
{स शरैर्भिन्नसर्वाङ्गो गात्रप्रस्रुतशोणितः}
{राक्षसेन्द्रः समूहस्थः फुल्लाशोक इवाबभौ} %||6-91-29||

\fourlineindentedshloka
{स रामबाणैरतिविद्धगात्रो}
{निशाचरेन्द्रः क्षतजार्द्रगात्रः}
{जगाम खेदं च समाजमध्ये}
{क्रोधं च चक्रे सुभृशं तदानीम्} %||6-92-30||


{॥इत्यार्षे श्रीमद्रामायणे वाल्मीकीये आदिकाव्ये युद्धकाण्डे एकनवतितमः सर्गः॥}

\sect{द्विनवतितमः सर्गः}

\twolineshloka
{स तु तेन तदा क्रोधात्काकुत्स्थेनार्दितो रणे}
{रावणः समरश्लाघी महाक्रोधमुपागमत्} %||6-92-1||

\twolineshloka
{स दीप्तनयनो रोषाच्चापमायम्य वीर्यवान्}
{अभ्यर्दयत्सुसङ्क्रुद्धो राघवं परमाहवे} %||6-92-2||

\twolineshloka
{बाणधारासहस्रैस्तु स तोयद इवाम्बरात्}
{राघवं रावणो बाणैस्तटाकमिव पूरयत्} %||6-92-3||

\twolineshloka
{पूरितः शरजालेन धनुर्मुक्तेन संयुगे}
{महागिरिरिवाकम्प्यः काकुस्थो न प्रकम्पते} %||6-92-4||

\twolineshloka
{स शरैः शरजालानि वारयन्समरे स्थितः}
{गभस्तीनिव सूर्यस्य प्रतिजग्राह वीर्यवान्} %||6-92-5||

\twolineshloka
{ततः शरसहस्राणि क्षिप्रहस्तो निशाचरः}
{निजघानोरसि क्रुद्धो राघवस्य महात्मनः} %||6-92-6||

\twolineshloka
{स शोणितसमादिग्धः समरे लक्ष्मणाग्रजः}
{दृष्टः फुल्ल इवारण्ये सुमहान्किंशुकद्रुमः} %||6-92-7||

\twolineshloka
{शराभिघातसंरब्धः सोऽपि जग्राह सायकान्}
{काकुत्स्थः सुमहातेजा युगान्तादित्यवर्चसः} %||6-92-8||

\twolineshloka
{ततोऽन्योन्यं सुसंरब्धावुभौ तौ रामरावणौ}
{शरान्धकारे समरे नोपालक्षयतां तदा} %||6-92-9||

\twolineshloka
{ततः क्रोधसमाविष्टो रामो दशरथात्मजः}
{उवाच रावणं वीरः प्रहस्य परुषं वचः} %||6-92-10||

\twolineshloka
{मम भार्या जनस्थानादज्ञानाद्राक्षसाधम}
{हृता ते विवशा यस्मात्तस्मात्त्वं नासि वीर्यवान्} %||6-92-11||

\twolineshloka
{मया विरहितां दीनां वर्तमानां महावने}
{वैदेहीं प्रसभं हृत्वा शूरोऽहमिति मन्यसे} %||6-92-12||

\twolineshloka
{स्त्रीषु शूर विनाथासु परदाराभिमर्शके}
{कृत्वा कापुरुषं कर्म शूरोऽहमिति मन्यसे} %||6-92-13||

\twolineshloka
{भिन्नमर्याद निर्लज्ज चारित्रेष्वनवस्थित}
{दर्पान्मृत्युमुपादाय शूरोऽहमिति मन्यसे} %||6-92-14||

\twolineshloka
{शूरेण धनदभ्रात्रा बलैः समुदितेन च}
{श्लाघनीयं यशस्यं च कृतं कर्म महत्त्वया} %||6-92-15||

\twolineshloka
{उत्सेकेनाभिपन्नस्य गर्हितस्याहितस्य च}
{कर्मणः प्राप्नुहीदानीं तस्याद्य सुमहत्फलम्} %||6-92-16||

\twolineshloka
{शूरोऽहमिति चात्मानमवगच्छसि दुर्मते}
{नैव लज्जास्ति ते सीतां चोरवद्व्यपकर्षतः} %||6-92-17||

\twolineshloka
{यदि मत्संनिधौ सीता धर्षिता स्यात्त्वया बलात्}
{भ्रातरं तु खरं पश्येस्तदा मत्सायकैर्हतः} %||6-92-18||

\twolineshloka
{दिष्ट्यासि मम दुष्टात्मंश्चक्षुर्विषयमागतः}
{अद्य त्वां सायकैस्तीक्ष्णैर्नयामि यमसादनम्} %||6-92-19||

\twolineshloka
{अद्य ते मच्छरैश्छिन्नं शिरो ज्वलितकुण्डलम्}
{क्रव्यादा व्यपकर्षन्तु विकीर्णं रणपांसुषु} %||6-92-20||

\twolineshloka
{निपत्योरसि गृध्रास्ते क्षितौ क्षिप्तस्य रावण}
{पिबन्तु रुधिरं तर्षाद्बाणशल्यान्तरोथितम्} %||6-92-21||

\twolineshloka
{अद्य मद्बाणाभिन्नस्य गतासोः पतितस्य ते}
{कर्षन्त्वन्त्राणि पतगा गरुत्मन्त इवोरगान्} %||6-92-22||

\twolineshloka
{इत्येवं स वदन्वीरो रामः शत्रुनिबर्हणः}
{राक्षसेन्द्रं समीपस्थं शरवर्षैरवाकिरत्} %||6-92-23||

\twolineshloka
{बभूव द्विगुणं वीर्यं बलं हर्षश्च संयुगे}
{रामस्यास्त्रबलं चैव शत्रोर्निधनकाङ्क्षिणः} %||6-92-24||

\twolineshloka
{प्रादुर्बभूवुरस्त्राणि सर्वाणि विदितात्मनः}
{प्रहर्षाच्च महातेजाः शीघ्रहस्ततरोऽभवत्} %||6-92-25||

\twolineshloka
{शुभान्येतानि चिह्नानि विज्ञायात्मगतानि सः}
{भूय एवार्दयद्रामो रावणं राक्षसान्तकृत्} %||6-92-26||

\twolineshloka
{हरीणां चाश्मनिकरैः शरवर्षैश्च राघवात्}
{हन्यमानो दशग्रीवो विघूर्णहृदयोऽभवत्} %||6-92-27||

\twolineshloka
{यदा च शस्त्रं नारेभे न व्यकर्षच्छरासनम्}
{नास्य प्रत्यकरोद्वीर्यं विक्लवेनान्तरात्मना} %||6-92-28||

\twolineshloka
{क्षिप्ताश्चापि शरास्तेन शस्त्राणि विविधानि च}
{न रणार्थाय वर्तन्ते मृत्युकालेऽभिवर्ततः} %||6-92-29||

\twolineshloka
{सूतस्तु रथनेतास्य तदवस्थं निरीक्ष्य तम्}
{शनैर्युद्धादसम्भान्तो रथं तस्यापवाहयत्} %||6-93-30||


{॥इत्यार्षे श्रीमद्रामायणे वाल्मीकीये आदिकाव्ये युद्धकाण्डे द्विनवतितमः सर्गः॥}

\sect{त्रिनवतितमः सर्गः}

\twolineshloka
{स तु मोहात्सुसङ्क्रुद्धः कृतान्तबलचोदितः}
{क्रोधसंरक्तनयनो रावणो सूतमब्रवीत्} %||6-93-1||

\twolineshloka
{हीनवीर्यमिवाशक्तं पौरुषेण विवर्जितम्}
{भीरुं लघुमिवासत्त्वं विहीनमिव तेजसा} %||6-93-2||

\twolineshloka
{विमुक्तमिव मायाभिरस्त्रैरिव बहिष्कृतम्}
{मामवज्ञाय दुर्बुद्धे स्वया बुद्ध्या विचेष्टसे} %||6-93-3||

\twolineshloka
{किमर्थं मामवज्ञाय मच्छन्दमनवेक्ष्य च}
{त्वया शत्रुसमक्षं मे रथोऽयमपवाहितः} %||6-93-4||

\twolineshloka
{त्वयाद्य हि ममानार्य चिरकालसमार्जितम्}
{यशो वीर्यं च तेजश्च प्रत्ययश्च विनाशितः} %||6-93-5||

\twolineshloka
{शत्रोः प्रख्यातवीर्यस्य रञ्जनीयस्य विक्रमैः}
{पश्यतो युद्धलुब्धोऽहं कृतः कापुरुषस्त्वया} %||6-93-6||

\twolineshloka
{यस्त्वं रथमिमं मोहान्न चोद्वहसि दुर्मते}
{सत्योऽयं प्रतितर्को मे परेण त्वमुपस्कृतः} %||6-93-7||

\twolineshloka
{न हीदं विद्यते कर्म सुहृदो हितकाङ्क्षिणः}
{रिपूणां सदृशं चैतन्न त्वयैतत्स्वनुष्ठितम्} %||6-93-8||

\twolineshloka
{निवर्तय रथं शीघ्रं यावन्नापैति मे रिपुः}
{यदि वाप्युषितोऽसि त्वं स्मर्यन्ते यदि वा गुणाः} %||6-93-9||

\twolineshloka
{एवं परुषमुक्तस्तु हितबुद्धिरबुद्धिना}
{अब्रवीद्रावणं सूतो हितं सानुनयं वचः} %||6-93-10||

\twolineshloka
{न भीतोऽस्मि न मूढोऽस्मि नोपजप्तोऽस्मि शत्रुभिः}
{न प्रमत्तो न निःस्नेहो विस्मृता न च सत्क्रिया} %||6-93-11||

\twolineshloka
{मया तु हितकामेन यशश्च परिरक्षता}
{स्नेहप्रस्कन्नमनसा प्रियमित्यप्रियं कृतम्} %||6-93-12||

\twolineshloka
{नास्मिन्नर्थे महाराज त्वं मां प्रियहिते रतम्}
{कश्चिल्लघुरिवानार्यो दोषतो गन्तुमर्हसि} %||6-93-13||

\twolineshloka
{श्रूयतामभिधास्यामि यन्निमित्तं मया रथः}
{नदीवेग इवाम्भोभिः संयुगे विनिवर्तितः} %||6-93-14||

\twolineshloka
{श्रमं तवावगच्छामि महता रणकर्मणा}
{न हि ते वीर सौमुख्यं प्रहर्षं वोपधारये} %||6-93-15||

\twolineshloka
{रथोद्वहनखिन्नाश्च त इमे रथवाजिनः}
{दीना घर्मपरिश्रान्ता गावो वर्षहता इव} %||6-93-16||

\twolineshloka
{निमित्तानि च भूयिष्ठं यानि प्रादुर्भवन्ति नः}
{तेषु तेष्वभिपन्नेषु लक्षयाम्यप्रदक्षिणम्} %||6-93-17||

\twolineshloka
{देशकालौ च विज्ञेयौ लक्षणानीङ्गितानि च}
{दैन्यं हर्षश्च खेदश्च रथिनश्च बलाबलम्} %||6-93-18||

\twolineshloka
{स्थलनिम्नानि भूमेश्च समानि विषमाणि च}
{युद्धकालश्च विज्ञेयः परस्यान्तरदर्शनम्} %||6-93-19||

\twolineshloka
{उपयानापयाने च स्थानं प्रत्यपसर्पणम्}
{सर्वमेतद्रथस्थेन ज्ञेयं रथकुटुम्बिना} %||6-93-20||

\twolineshloka
{तव विश्रामहेतोस्तु तथैषां रथवाजिनाम्}
{रौद्रं वर्जयता खेदं क्षमं कृतमिदं मया} %||6-93-21||

\twolineshloka
{न मया स्वेच्छया वीर रथोऽयमपवाहितः}
{भर्तृस्नेहपरीतेन मयेदं यत्कृतं विभो} %||6-93-22||

\twolineshloka
{आज्ञापय यथातत्त्वं वक्ष्यस्यरिनिषूदन}
{तत्करिष्याम्यहं वीरं गतानृण्येन चेतसा} %||6-93-23||

\twolineshloka
{सन्तुष्टस्तेन वाक्येन रावणस्तस्य सारथेः}
{प्रशस्यैनं बहुविधं युद्धलुब्धोऽब्रवीदिदम्} %||6-93-24||

\twolineshloka
{रथं शीघ्रमिमं सूत राघवाभिमुखं कुरु}
{नाहत्वा समरे शत्रून्निवर्तिष्यति रावणः} %||6-93-25||

\twolineshloka
{एवमुक्त्वा ततस्तुष्टो रावणो राक्षसेश्वरः}
{ददौ तस्य शुभं ह्येकं हस्ताभरणमुत्तमम्} %||6-93-26||

\fourlineindentedshloka
{ततो द्रुतं रावणवाक्यचोदितः}
{प्रचोदयामास हयान्स सारथिः}
{स राक्षसेन्द्रस्य ततो महारथः}
{क्षणेन रामस्य रणाग्रतोऽभवत्} %||6-94-27||


{॥इत्यार्षे श्रीमद्रामायणे वाल्मीकीये आदिकाव्ये युद्धकाण्डे त्रिनवतितमः सर्गः॥}

\sect{चतुर्नवतितमः सर्गः}

\twolineshloka
{तमापतन्तं सहसा स्वनवन्तं महाध्वजम्}
{रथं राक्षसराजस्य नरराजो ददर्श ह} %||6-94-1||

\threelineshloka
{कृष्णवाजिसमायुक्तं युक्तं रौद्रेण वर्चसा}
{तडित्पताकागहनं दर्शितेन्द्रायुधायुधम्}
{शरधारा विमुञ्चन्तं धारासारमिवान्बुदम्} %||6-94-2||

\threelineshloka
{तं दृष्ट्वा मेघसङ्काशमापतन्तं रथं रिपोः}
{गिरेर्वज्राभिमृष्टस्य दीर्यतः सदृशस्वनम्}
{उवाच मातलिं रामः सहस्राक्षस्य सारथिम्} %||6-94-3||

\threelineshloka
{मातले पश्य संरब्धमापतन्तं रथं रिपोः}
{यथापसव्यं पतता वेगेन महता पुनः}
{समरे हन्तुमात्मानं तथानेन कृता मतिः} %||6-94-4||

\twolineshloka
{तदप्रमादमातिष्ठ प्रत्युद्गच्छ रथं रिपोः}
{विध्वंसयितुमिच्छामि वायुर्मेघमिवोत्थितम्} %||6-94-5||

\twolineshloka
{अविक्लवमसम्भ्रान्तमव्यग्रहृदयेक्षणम्}
{रश्मिसञ्चारनियतं प्रचोदय रथं द्रुतम्} %||6-94-6||

\twolineshloka
{कामं न त्वं समाधेयः पुरन्दररथोचितः}
{युयुत्सुरहमेकाग्रः स्मारये त्वां न शिक्षये} %||6-94-7||

\twolineshloka
{परितुष्टः स रामस्य तेन वाक्येन मातलिः}
{प्रचोदयामास रथं सुरसारथिसत्तमः} %||6-94-8||

\twolineshloka
{अपसव्यं ततः कुर्वन्रावणस्य महारथम्}
{चक्रोत्क्षिप्तेन रजसा रावणं व्यवधूनयत्} %||6-94-9||

\twolineshloka
{ततः क्रुद्धो दशग्रीवस्ताम्रविस्फारितेक्षणः}
{रथप्रतिमुखं रामं सायकैरवधूनयत्} %||6-94-10||

\threelineshloka
{धर्षणामर्षितो रामो धैर्यं रोषेण लङ्घयन्}
{जग्राह सुमहावेगमैन्द्रं युधि शरासनम्}
{शरांश्च सुमहातेजाः सूर्यरश्मिसमप्रभान्} %||6-94-11||

\twolineshloka
{तदुपोढं महद्युद्धमन्योन्यवधकाङ्क्षिणोः}
{परस्पराभिमुखयोर्दृप्तयोरिव सिंहयोः} %||6-94-12||

\twolineshloka
{ततो देवाः सगन्धर्वाः सिद्धाश्च परमर्षयः}
{समीयुर्द्वैरथं द्रष्टुं रावणक्षयकाङ्क्षिणः} %||6-94-13||

\twolineshloka
{समुत्पेतुरथोत्पाता दारुणा लोमहर्षणाः}
{रावणस्य विनाशाय राघवस्य जयाय च} %||6-94-14||

\twolineshloka
{ववर्ष रुधिरं देवो रावणस्य रथोपरि}
{वाता मण्डलिनस्तीव्रा अपसव्यं प्रचक्रमुः} %||6-94-15||

\twolineshloka
{महद्गृध्रकुलं चास्य भ्रममाणं नभस्तले}
{येन येन रथो याति तेन तेन प्रधावति} %||6-94-16||

\twolineshloka
{सन्ध्यया चावृता लङ्का जपापुष्पनिकाशया}
{दृश्यते सम्प्रदीतेव दिवसेऽपि वसुन्धरा} %||6-94-17||

\twolineshloka
{सनिर्घाता महोल्काश्च सम्प्रचेतुर्महास्वनाः}
{विषादयन्त्यो रक्षांसि रावणस्य तदाहिताः} %||6-94-18||

\twolineshloka
{रावणश्च यतस्तत्र प्रचचाल वसुन्धरा}
{रक्षसां च प्रहरतां गृहीता इव बाहवः} %||6-94-19||

\twolineshloka
{ताम्राः पीताः सिताः श्वेताः पतिताः सूर्यरश्मयः}
{दृश्यन्ते रावणस्याङ्गे पर्वतस्येव धातवः} %||6-94-20||

\twolineshloka
{गृध्रैरनुगताश्चास्य वमन्त्यो ज्वलनं मुखैः}
{प्रणेदुर्मुखमीक्षन्त्यः संरब्धमशिवं शिवाः} %||6-94-21||

\twolineshloka
{प्रतिकूलं ववौ वायू रणे पांसून्समुत्किरन्}
{तस्य राक्षसराजस्य कुर्वन्दृष्टिविलोपनम्} %||6-94-22||

\twolineshloka
{निपेतुरिन्द्राशनयः सैन्ये चास्य समन्ततः}
{दुर्विषह्य स्वना घोरा विना जलधरस्वनम्} %||6-94-23||

\twolineshloka
{दिशश्च प्रदिशः सर्वा बभूवुस्तिमिरावृताः}
{पांसुवर्षेण महता दुर्दर्शं च नभोऽभवत्} %||6-94-24||

\twolineshloka
{कुर्वन्त्यः कलहं घोरं सारिकास्तद्रथं प्रति}
{निपेतुः शतशस्तत्र दारुणा दारुणस्वनाः} %||6-94-25||

\twolineshloka
{जघनेभ्यः स्फुलिङ्गांश्च नेत्रेभ्योऽश्रूणि सन्ततम्}
{मुमुचुस्तस्य तुरगास्तुल्यमग्निं च वारि च} %||6-94-26||

\twolineshloka
{एवं प्रकारा बहवः समुत्पाता भयावहाः}
{रावणस्य विनाशाय दारुणाः सम्प्रजज्ञिरे} %||6-94-27||

\twolineshloka
{रामस्यापि निमित्तानि सौम्यानि च शिवानि च}
{बभूवुर्जयशंसीनि प्रादुर्भूतानि सर्वशः} %||6-94-28||

\fourlineindentedshloka
{ततो निरीक्ष्यात्मगतानि राघवो}
{रणे निमित्तानि निमित्तकोविदः}
{जगाम हर्षं च परां च निर्वृतिम्}
{चकार युद्धेऽभ्यधिकं च विक्रमम्} %||6-95-29||


{॥इत्यार्षे श्रीमद्रामायणे वाल्मीकीये आदिकाव्ये युद्धकाण्डे चतुर्नवतितमः सर्गः॥}

\sect{पञ्चनवतितमः सर्गः}

\twolineshloka
{ततः प्रवृत्तं सुक्रूरं रामरावणयोस्तदा}
{सुमहद्द्वैरथं युद्धं सर्वलोकभयावहम्} %||6-95-1||

\twolineshloka
{ततो राक्षससैन्यं च हरीणां च महद्बलम्}
{प्रगृहीतप्रहरणं निश्चेष्टं समतिष्ठत} %||6-95-2||

\twolineshloka
{सम्प्रयुद्धौ ततो दृष्ट्वा बलवन्नरराक्षसौ}
{व्याक्षिप्तहृदयाः सर्वे परं विस्मयमागताः} %||6-95-3||

\twolineshloka
{नानाप्रहरणैर्व्यग्रैर्भुजैर्विस्मितबुद्धयः}
{तस्थुः प्रेक्ष्य च सङ्ग्रामं नाभिजघ्नुः परस्परम्} %||6-95-4||

\twolineshloka
{रक्षसां रावणं चापि वानराणां च राघवम्}
{पश्यतां विस्मिताक्षाणां सैन्यं चित्रमिवाबभौ} %||6-95-5||

\twolineshloka
{तौ तु तत्र निमित्तानि दृष्ट्वा राघवरावणौ}
{कृतबुद्धी स्थिरामर्षौ युयुधाते अभीतवत्} %||6-95-6||

\twolineshloka
{जेतव्यमिति काकुत्स्थो मर्तव्यमिति रावणः}
{धृतौ स्ववीर्यसर्वस्वं युद्धेऽदर्शयतां तदा} %||6-95-7||

\twolineshloka
{ततः क्रोधाद्दशग्रीवः शरान्सन्धाय वीर्यवान्}
{मुमोच ध्वजमुद्दिश्य राघवस्य रथे स्थितम्} %||6-95-8||

\twolineshloka
{ते शरास्तमनासाद्य पुरन्दररथध्वजम्}
{रक्तशक्तिं परामृश्य निपेतुर्धरणीतले} %||6-95-9||

\twolineshloka
{ततो रामोऽभिसङ्क्रुद्धश्चापमायम्य वीर्यवान्}
{कृतप्रतिकृतं कर्तुं मनसा सम्प्रचक्रमे} %||6-95-10||

\twolineshloka
{रावणध्वजमुद्दिश्य मुमोच निशितं शरम्}
{महासर्पमिवासह्यं ज्वलन्तं स्वेन तेजसा} %||6-95-11||

\twolineshloka
{जगाम स महीं भित्त्वा दशग्रीवध्वजं शरः}
{स निकृत्तोऽपतद्भूमौ रावणस्य रथध्वजः} %||6-95-12||

\twolineshloka
{ध्वजस्योन्मथनं दृष्ट्वा रावणः सुमहाबलः}
{क्रोधजेनाग्निना सङ्ख्ये प्रदीप्त इव चाभवत्} %||6-95-13||

\twolineshloka
{स रोषवशमापन्नः शरवर्षं महद्वमन्}
{रामस्य तुरगान्दिव्याञ्शरैर्विव्याध रावणः} %||6-95-14||

\twolineshloka
{ते विद्धा हरयस्तस्य नास्खलन्नापि बभ्रमुः}
{बभूवुः स्वस्थहृदयाः पद्मनालैरिवाहताः} %||6-95-15||

\twolineshloka
{तेषामसम्भ्रमं दृष्ट्वा वाजिनां रावणस्तदा}
{भूय एव सुसङ्क्रुद्धः शरवर्षं मुमोच ह} %||6-95-16||

\twolineshloka
{गदाश्च परिघांश्चैव चक्राणि मुसलानि च}
{गिरिशृङ्गाणि वृक्षांश्च तथा शूलपरश्वधान्} %||6-95-17||

\twolineshloka
{मायाविहितमेतत्तु शस्त्रवर्षमपातयत्}
{सहस्रशस्ततो बाणानश्रान्तहृदयोद्यमः} %||6-95-18||

\twolineshloka
{तुमुलं त्रासजननं भीमं भीमप्रतिस्वनम्}
{दुर्धर्षमभवद्युद्धे नैकशस्त्रमयं महत्} %||6-95-19||

\threelineshloka
{विमुच्य राघवरथं समन्ताद्वानरे बले}
{सायकैरन्तरिक्षं च चकाराशु निरन्तरम्}
{मुमोच च दशग्रीवो निःसङ्गेनान्तरात्मना} %||6-95-20||

\twolineshloka
{व्यायच्छमानं तं दृष्ट्वा तत्परं रावणं रणे}
{प्रहसन्निव काकुत्स्थः सन्दधे सायकाञ्शितान्} %||6-95-21||

\twolineshloka
{स मुमोच ततो बाणान्रणे शतसहस्रशः}
{तान्दृष्ट्वा रावणश्चक्रे स्वशरैः खं निरन्तरम्} %||6-95-22||

\twolineshloka
{ततस्ताभ्यां प्रयुक्तेन शरवर्षेण भास्वता}
{शरबद्धमिवाभाति द्वितीयं भास्वदम्बरम्} %||6-95-23||

\twolineshloka
{नानिमित्तोऽभवद्बाणो नातिभेत्ता न निष्फलः}
{तथा विसृजतोर्बाणान्रामरावणयोर्मृधे} %||6-95-24||

\twolineshloka
{प्रायुध्येतामविच्छिन्नमस्यन्तौ सव्यदक्षिणम्}
{चक्रतुस्तौ शरौघैस्तु निरुच्छ्वासमिवाम्बरम्} %||6-95-25||

\twolineshloka
{रावणस्य हयान्रामो हयान्रामस्य रावणः}
{जघ्नतुस्तौ तदान्योन्यं कृतानुकृतकारिणौ} %||6-96-26||


{॥इत्यार्षे श्रीमद्रामायणे वाल्मीकीये आदिकाव्ये युद्धकाण्डे पञ्चनवतितमः सर्गः॥}

\sect{षण्णवतितमः सर्गः}

\twolineshloka
{तौ तथा युध्यमानौ तु समरे रामरावणौ}
{ददृशुः सर्वभूतानि विस्मितेनान्तरात्मना} %||6-96-1||

\twolineshloka
{अर्दयन्तौ तु समरे तयोस्तौ स्यन्दनोत्तमौ}
{परस्परवधे युक्तौ घोररूपौ बभूवतुः} %||6-96-2||

\twolineshloka
{मण्डलानि च वीथीश्च गतप्रत्यागतानि च}
{दर्शयन्तौ बहुविधां सूतौ सारथ्यजां गतिम्} %||6-96-3||

\twolineshloka
{अर्दयन्रावणं रामो राघवं चापि रावणः}
{गतिवेगं समापन्नौ प्रवर्तन निवर्तने} %||6-96-4||

\twolineshloka
{क्षिपतोः शरजालानि तयोस्तौ स्यन्दनोत्तमौ}
{चेरतुः संयुगमहीं सासारौ जलदाविव} %||6-96-5||

\twolineshloka
{दर्शयित्वा तदा तौ तु गतिं बहुविधां रणे}
{परस्परस्याभिमुखौ पुनरेव च तस्थतुः} %||6-96-6||

\twolineshloka
{धुरं धुरेण रथयोर्वक्त्रं वक्त्रेण वाजिनाम्}
{पताकाश्च पताकाभिः समेयुः स्थितयोस्तदा} %||6-96-7||

\twolineshloka
{रावणस्य ततो रामो धनुर्मुक्तैः शितैः शरैः}
{चतुर्भिश्चतुरो दीप्तान्हयान्प्रत्यपसर्पयत्} %||6-96-8||

\twolineshloka
{स क्रोधवशमापन्नो हयानामपसर्पणे}
{मुमोच निशितान्बाणान्राघवाय निशाचरः} %||6-96-9||

\twolineshloka
{सोऽतिविद्धो बलवता दशग्रीवेण राघवः}
{जगाम न विकारं च न चापि व्यथितोऽभवत्} %||6-96-10||

\twolineshloka
{चिक्षेप च पुनर्बाणान्वज्रपातसमस्वनान्}
{सारथिं वज्रहस्तस्य समुद्दिश्य निशाचरः} %||6-96-11||

\twolineshloka
{मातलेस्तु महावेगाः शरीरे पतिताः शराः}
{न सूक्ष्ममपि सम्मोहं व्यथां वा प्रददुर्युधि} %||6-96-12||

\twolineshloka
{तया धर्षणया क्रोद्धो मातलेर्न तथात्मनः}
{चकार शरजालेन राघवो विमुखं रिपुम्} %||6-96-13||

\twolineshloka
{विंशतिं त्रिंशतं षष्टिं शतशोऽथ सहस्रशः}
{मुमोच राघवो वीरः सायकान्स्यन्दने रिपोः} %||6-96-14||

\twolineshloka
{गदानां मुसलानां च परिघाणां च निस्वनैः}
{शराणां पुङ्खवातैश्च क्षुभिताः सप्तसागराः} %||6-96-15||

\twolineshloka
{क्षुब्धानां सागराणां च पातालतलवासिनः}
{व्यथिताः पन्नगाः सर्वे दानवाश्च सहस्रशः} %||6-96-16||

\twolineshloka
{चकम्पे मेदिनी कृत्स्ना सशैलवनकानना}
{भास्करो निष्प्रभश्चाभून्न ववौ चापि मारुतः} %||6-96-17||

\twolineshloka
{ततो देवाः सगन्धर्वाः सिद्धाश्च परमर्षयः}
{चिन्तामापेदिरे सर्वे सकिंनरमहोरगाः} %||6-96-18||

\twolineshloka
{स्वस्ति गोब्राह्मणेभ्योऽस्तु लोकास्तिष्ठन्तु शाश्वताः}
{जयतां राघवः सङ्ख्ये रावणं राक्षसेश्वरम्} %||6-96-19||

\threelineshloka
{ततः क्रुद्धो महाबाहू रघूणां कीर्तिवर्धनः}
{सन्धाय धनुषा रामः क्षुरमाशीविषोपमम्}
{रावणस्य शिरोऽच्छिन्दच्छ्रीमज्ज्वलितकुण्डलम्} %||6-96-20||

\twolineshloka
{तच्छिरः पतितं भूमौ दृष्टं लोकैस्त्रिभिस्तदा}
{तस्यैव सदृशं चान्यद्रावणस्योत्थितं शिरः} %||6-96-21||

\twolineshloka
{तत्क्षिप्रं क्षिप्रहस्तेन रामेण क्षिप्रकारिणा}
{द्वितीयं रावणशिरश्छिन्नं संयति सायकैः} %||6-96-22||

\twolineshloka
{छिन्नमात्रं च तच्छीर्षं पुनरन्यत्स्म दृश्यते}
{तदप्यशनिसङ्काशैश्छिन्नं रामेण सायकैः} %||6-96-23||

\twolineshloka
{एवमेव शतं छिन्नं शिरसां तुल्यवर्चसाम्}
{न चैव रावणस्यान्तो दृश्यते जीवितक्षये} %||6-96-24||

\twolineshloka
{ततः सर्वास्त्रविद्वीरः कौसल्यानन्दिवर्धनः}
{मार्गणैर्बहुभिर्युक्तश्चिन्तयामास राघवः} %||6-96-25||

\twolineshloka
{मारीचो निहतो यैस्तु खरो यैस्तु सुदूषणः}
{क्रञ्चारण्ये विराधस्तु कबन्धो दण्डका वने} %||6-96-26||

\twolineshloka
{त इमे सायकाः सर्वे युद्धे प्रत्ययिका मम}
{किं नु तत्कारणं येन रावणे मन्दतेजसः} %||6-96-27||

\twolineshloka
{इति चिन्तापरश्चासीदप्रमत्तश्च संयुगे}
{ववर्ष शरवर्षाणि राघवो रावणोरसि} %||6-96-28||

\twolineshloka
{रावणोऽपि ततः क्रुद्धो रथस्थो राक्षसेश्वरः}
{गदामुसलवर्षेण रामं प्रत्यर्दयद्रणे} %||6-96-29||

\twolineshloka
{देवदानवयक्षाणां पिशाचोरगरक्षसाम्}
{पश्यतां तन्महद्युद्धं सर्वरात्रमवर्तत} %||6-96-30||

\twolineshloka
{नैव रत्रिं न दिवसं न मुहूर्तं न चक्षणम्}
{रामरावणयोर्युद्धं विराममुपगच्छति} %||6-97-31||


{॥इत्यार्षे श्रीमद्रामायणे वाल्मीकीये आदिकाव्ये युद्धकाण्डे षण्णवतितमः सर्गः॥}

\sect{सप्तनवतितमः सर्गः}

\twolineshloka
{अथ संस्मारयामास राघवं मातलिस्तदा}
{अजानन्निव किं वीर त्वमेनमनुवर्तसे} %||6-97-1||

\twolineshloka
{विसृजास्मै वधाय त्वमस्त्रं पैतामहं प्रभो}
{विनाशकालः कथितो यः सुरैः सोऽद्य वर्तते} %||6-97-2||

\twolineshloka
{ततः संस्मारितो रामस्तेन वाक्येन मातलेः}
{जग्राह स शरं दीप्तं निश्वसन्तमिवोरगम्} %||6-97-3||

\twolineshloka
{यमस्मै प्रथमं प्रादादगस्त्यो भगवानृषिः}
{ब्रह्मदत्तं महद्बाणममोघं युधि वीर्यवान्} %||6-97-4||

\twolineshloka
{ब्रह्मणा निर्मितं पूर्वमिन्द्रार्थममितौजसा}
{दत्तं सुरपतेः पूर्वं त्रिलोकजयकाङ्क्षिणः} %||6-97-5||

\twolineshloka
{यस्य वाजेषु पवनः फले पावकभास्करौ}
{शरीरमाकाशमयं गौरवे मेरुमन्दरौ} %||6-97-6||

\twolineshloka
{जाज्वल्यमानं वपुषा सुपुङ्खं हेमभूषितम्}
{तेजसा सर्वभूतानां कृतं भास्करवर्चसम्} %||6-97-7||

\twolineshloka
{सधूममिव कालाग्निं दीप्तमाशीविषं यथा}
{रथनागाश्ववृन्दानां भेदनं क्षिप्रकारिणम्} %||6-97-8||

\twolineshloka
{द्वाराणां परिघाणां च गिरीणामपि भेदनम्}
{नानारुधिरसिक्ताङ्गं मेदोदिग्धं सुदारुणम्} %||6-97-9||

\twolineshloka
{वज्रसारं महानादं नानासमितिदारुणम्}
{सर्ववित्रासनं भीमं श्वसन्तमिव पन्नगम्} %||6-97-10||

\twolineshloka
{कङ्कगृध्रबलानां च गोमायुगणरक्षसाम्}
{नित्यं भक्षप्रदं युद्धे यमरूपं भयावहम्} %||6-97-11||

\twolineshloka
{नन्दनं वानरेन्द्राणां रक्षसामवसादनम्}
{वाजितं विविधैर्वाजैश्चारुचित्रैर्गरुत्मतः} %||6-97-12||

\twolineshloka
{तमुत्तमेषुं लोकानामिक्ष्वाकुभयनाशनम्}
{द्विषतां कीर्तिहरणं प्रहर्षकरमात्मनः} %||6-97-13||

\twolineshloka
{अभिमन्त्र्य ततो रामस्तं महेषुं महाबलः}
{वेदप्रोक्तेन विधिना सन्दधे कार्मुके बली} %||6-97-14||

\twolineshloka
{स रावणाय सङ्क्रुद्धो भृशमायम्य कार्मुकम्}
{चिक्षेप परमायत्तस्तं शरं मर्मघातिनम्} %||6-97-15||

\twolineshloka
{स वज्र इव दुर्धर्षो वज्रबाहुविसर्जितः}
{कृतान्त इव चावार्यो न्यपतद्रावणोरसि} %||6-97-16||

\twolineshloka
{स विसृष्टो महावेगः शरीरान्तकरः शरः}
{बिभेद हृदयं तस्य रावणस्य दुरात्मनः} %||6-97-17||

\twolineshloka
{रुधिराक्तः स वेगेन जीवितान्तकरः शरः}
{रावणस्य हरन्प्राणान्विवेश धरणीतलम्} %||6-97-18||

\twolineshloka
{स शरो रावणं हत्वा रुधिरार्द्रकृतच्छविः}
{कृतकर्मा निभृतवत्स्वतूणीं पुनराविशत्} %||6-97-19||

\twolineshloka
{तस्य हस्ताद्धतस्याशु कार्मुकं तत्ससायकम्}
{निपपात सह प्राणैर्भ्रश्यमानस्य जीवितात्} %||6-97-20||

\twolineshloka
{गतासुर्भीमवेगस्तु नैरृतेन्द्रो महाद्युतिः}
{पपात स्यन्दनाद्भूमौ वृत्रो वज्रहतो यथा} %||6-97-21||

\twolineshloka
{तं दृष्ट्वा पतितं भूमौ हतशेषा निशाचराः}
{हतनाथा भयत्रस्ताः सर्वतः सम्प्रदुद्रुवुः} %||6-97-22||

\twolineshloka
{नर्दन्तश्चाभिपेतुस्तान्वानरा द्रुमयोधिनः}
{दशग्रीववधं दृष्ट्वा विजयं राघवस्य च} %||6-97-23||

\twolineshloka
{अर्दिता वानरैर्हृष्टैर्लङ्कामभ्यपतन्भयात्}
{हताश्रयत्वात्करुणैर्बाष्पप्रस्रवणैर्मुखैः} %||6-97-24||

\twolineshloka
{ततो विनेदुः संहृष्टा वानरा जितकाशिनः}
{वदन्तो राघवजयं रावणस्य च तं वधम्} %||6-97-25||

\twolineshloka
{अथान्तरिक्षे व्यनदत्सौम्यस्त्रिदशदुन्दुभिः}
{दिव्यगन्धवहस्तत्र मारुतः सुसुखो ववौ} %||6-97-26||

\twolineshloka
{निपपातान्तरिक्षाच्च पुष्पवृष्टिस्तदा भुवि}
{किरन्ती राघवरथं दुरवापा मनोहराः} %||6-97-27||

\twolineshloka
{राघवस्तवसंयुक्ता गगने च विशुश्रुवे}
{साधु साध्विति वागग्र्या देवतानां महात्मनाम्} %||6-97-28||

\twolineshloka
{आविवेश महान्हर्षो देवानां चारणैः सह}
{रावणे निहते रौद्रे सर्वलोकभयङ्करे} %||6-97-29||

\twolineshloka
{ततः सकामं सुग्रीवमङ्गदं च महाबलम्}
{चकार राघवः प्रीतो हत्वा राक्षसपुङ्गवम्} %||6-97-30||

\fourlineindentedshloka
{ततः प्रजग्मुः प्रशमं मरुद्गणा}
{दिशः प्रसेदुर्विमलं नभोऽभवत्}
{मही चकम्पे न च मारुता ववुः}
{स्थिरप्रभश्चाप्यभवद्दिवाकरः} %||6-97-31||

\fourlineindentedshloka
{ततस्तु सुग्रीवविभीषणादयः}
{सुहृद्विशेषाः सहलक्ष्मणास्तदा}
{समेत्य हृष्टा विजयेन राघवम्}
{रणेऽभिरामं विधिनाभ्यपूजयन्} %||6-97-32||

\fourlineindentedshloka
{स तु निहतरिपुः स्थिरप्रतिज्ञः}
{स्वजनबलाभिवृतो रणे रराज}
{रघुकुलनृपनन्दनो महौजाः}
{त्रिदशगणैरभिसंवृतो यथेन्द्रः} %||6-98-33||


{॥इत्यार्षे श्रीमद्रामायणे वाल्मीकीये आदिकाव्ये युद्धकाण्डे सप्तनवतितमः सर्गः॥}

\sect{अष्टनवतितमः सर्गः}

\twolineshloka
{रावणं निहतं श्रुत्वा राघवेण महात्मना}
{अन्तःपुराद्विनिष्पेतू राक्षस्यः शोककर्शिताः} %||6-98-1||

\twolineshloka
{वार्यमाणाः सुबहुशो वृष्टन्त्यः क्षितिपांसुषु}
{विमुक्तकेश्यो दुःखार्ता गावो वत्सहता यथा} %||6-98-2||

\twolineshloka
{उत्तरेण विनिष्क्रम्य द्वारेण सह राक्षसैः}
{प्रविश्यायोधनं घोरं विचिन्वन्त्यो हतं पतिम्} %||6-98-3||

\twolineshloka
{आर्यपुत्रेति वादिन्यो हा नाथेति च सर्वशः}
{परिपेतुः कबन्धाङ्कां महीं शोणितकर्दमाम्} %||6-98-4||

\twolineshloka
{ता बाष्पपरिपूर्णाक्ष्यो भर्तृशोकपराजिताः}
{करेण्व इव नर्दन्त्यो विनेदुर्हतयूथपाः} %||6-98-5||

\twolineshloka
{ददृशुस्ता महाकायं महावीर्यं महाद्युतिम्}
{रावणं निहतं भूमौ नीलाञ्जनचयोपमम्} %||6-98-6||

\twolineshloka
{ताः पतिं सहसा दृष्ट्वा शयानं रणपांसुषु}
{निपेतुस्तस्य गात्रेषु छिन्ना वनलता इव} %||6-98-7||

\twolineshloka
{बहुमानात्परिष्वज्य काचिदेनं रुरोद ह}
{चरणौ काचिदालिङ्ग्य काचित्कण्ठेऽवलम्ब्य च} %||6-98-8||

\twolineshloka
{उद्धृत्य च भुजौ काचिद्भूमौ स्म परिवर्तते}
{हतस्य वदनं दृष्ट्वा काचिन्मोहमुपागमत्} %||6-98-9||

\twolineshloka
{काचिदङ्के शिरः कृत्वा रुरोद मुखमीक्षती}
{स्नापयन्ती मुखं बाष्पैस्तुषारैरिव पङ्कजम्} %||6-98-10||

\twolineshloka
{एवमार्ताः पतिं दृष्ट्वा रावणं निहतं भुवि}
{चुक्रुशुर्बहुधा शोकाद्भूयस्ताः पर्यदेवयन्} %||6-98-11||

\twolineshloka
{येन वित्रासितः शक्रो येन वित्रासितो यमः}
{येन वैश्रवणो राजा पुष्पकेण वियोजितः} %||6-98-12||

\twolineshloka
{गन्धर्वाणामृषीणां च सुराणां च महात्मनाम्}
{भयं येन महद्दत्तं सोऽयं शेते रणे हतः} %||6-98-13||

\twolineshloka
{असुरेभ्यः सुरेभ्यो वा पन्नगेभ्योऽपि वा तथा}
{न भयं यो विजानाति तस्येदं मानुषाद्भयम्} %||6-98-14||

\twolineshloka
{अवध्यो देवतानां यस्तथा दानवरक्षसाम्}
{हतः सोऽयं रणे शेते मानुषेण पदातिना} %||6-98-15||

\twolineshloka
{यो न शक्यः सुरैर्हन्तुं न यक्षैर्नासुरैस्तथा}
{सोऽयं कश्चिदिवासत्त्वो मृत्युं मर्त्येन लम्भितः} %||6-98-16||

\twolineshloka
{एवं वदन्त्यो बहुधा रुरुदुस्तस्य ताः स्त्रियः}
{भूय एव च दुःखार्ता विलेपुश्च पुनः पुनः} %||6-98-17||

\twolineshloka
{अशृण्वता तु सुहृदां सततं हितवादिनाम्}
{एताः सममिदानीं ते वयमात्मा च पातिताः} %||6-98-18||

\twolineshloka
{ब्रुवाणोऽपि हितं वाक्यमिष्टो भ्राता विभीषणः}
{धृष्टं परुषितो मोहात्त्वयात्मवधकाङ्क्षिणा} %||6-98-19||

\twolineshloka
{यदि निर्यातिता ते स्यात्सीता रामाय मैथिली}
{न नः स्याद्व्यसनं घोरमिदं मूलहरं महत्} %||6-98-20||

\twolineshloka
{वृत्तकामो भवेद्भ्राता रामो मित्रकुलं भवेत्}
{वयं चाविधवाः सर्वाः सकामा न च शत्रवः} %||6-98-21||

\twolineshloka
{त्वया पुनर्नृशंसेन सीतां संरुन्धता बलात्}
{राक्षसा वयमात्मा च त्रयं तुलं निपातितम्} %||6-98-22||

\twolineshloka
{न कामकारः कामं वा तव राक्षसपुङ्गव}
{दैवं चेष्टयते सर्वं हतं दैवेन हन्यते} %||6-98-23||

\twolineshloka
{वानराणां विनाशोऽयं राक्षसानां च ते रणे}
{तव चैव महाबाहो दैवयोगादुपागतः} %||6-98-24||

\twolineshloka
{नैवार्थेन न कामेन विक्रमेण न चाज्ञया}
{शक्या दैवगतिर्लोके निवर्तयितुमुद्यता} %||6-98-25||

\twolineshloka
{विलेपुरेवं दीनास्ता राक्षसाधिपयोषितः}
{कुरर्य इव दुःखार्ता बाष्पपर्याकुलेक्षणाः} %||6-99-26||


{॥इत्यार्षे श्रीमद्रामायणे वाल्मीकीये आदिकाव्ये युद्धकाण्डे अष्टनवतितमः सर्गः॥}

\sect{एकोनशततमः सर्गः}

\twolineshloka
{तासां विलपमानानां तथा राक्षसयोषिताम्}
{ज्येष्ठा पत्नी प्रिया दीना भर्तारं समुदैक्षत} %||6-99-1||

\twolineshloka
{दशग्रीवं हतं दृष्ट्वा रामेणाचिन्त्यकर्मणा}
{पतिं मन्दोदरी तत्र कृपणा पर्यदेवयत्} %||6-99-2||

\twolineshloka
{ननु नाम महाबाहो तव वैश्रवणानुज}
{क्रुद्धस्य प्रमुखे स्थातुं त्रस्यत्यपि पुरन्दरः} %||6-99-3||

\twolineshloka
{ऋषयश्च महीदेवा गन्धर्वाश्च यशस्विनः}
{ननु नाम तवोद्वेगाच्चारणाश्च दिशो गताः} %||6-99-4||

\twolineshloka
{स त्वं मानुषमात्रेण रामेण युधि निर्जितः}
{न व्यपत्रपसे राजन्किमिदं राक्षसर्षभ} %||6-99-5||

\twolineshloka
{कथं त्रैलोक्यमाक्रम्य श्रिया वीर्येण चान्वितम्}
{अविषह्यं जघान त्वां मानुषो वनगोचरः} %||6-99-6||

\twolineshloka
{मानुषाणामविषये चरतः कामरूपिणः}
{विनाशस्तव रामेण संयुगे नोपपद्यते} %||6-99-7||

\twolineshloka
{न चैतत्कर्म रामस्य श्रद्दधामि चमूमुखे}
{सर्वतः समुपेतस्य तव तेनाभिमर्शनम्} %||6-99-8||

\twolineshloka
{इन्द्रियाणि पुरा जित्वा जितं त्रिभुवणं त्वया}
{स्मरद्भिरिव तद्वैरमिन्द्रियैरेव निर्जितः} %||6-99-9||

\twolineshloka
{अथ वा रामरूपेण वासवः स्वयमागतः}
{मायां तव विनाशाय विधायाप्रतितर्किताम्} %||6-99-10||

\twolineshloka
{यदैव हि जनस्थाने राक्षसैर्बहुभिर्वृतः}
{खरस्तव हतो भ्राता तदैवासौ न मानुषः} %||6-99-11||

\twolineshloka
{यदैव नगरीं लङ्कां दुष्प्रवेषां सुरैरपि}
{प्रविष्टो हनुमान्वीर्यात्तदैव व्यथिता वयम्} %||6-99-12||

\twolineshloka
{क्रियतामविरोधश्च राघवेणेति यन्मया}
{उच्यमानो न गृह्णासि तस्येयं व्युष्टिरागता} %||6-99-13||

\twolineshloka
{अकस्माच्चाभिकामोऽसि सीतां राक्षसपुङ्गव}
{ऐश्वर्यस्य विनाशाय देहस्य स्वजनस्य च} %||6-99-14||

\twolineshloka
{अरुन्धत्या विशिष्टां तां रोहिण्याश्चापि दुर्मते}
{सीतां धर्षयता मान्यां त्वया ह्यसदृशं कृतम्} %||6-99-15||

\twolineshloka
{न कुलेन न रूपेण न दाक्षिण्येन मैथिली}
{मयाधिका वा तुल्या वा त्वं तु मोहान्न बुध्यसे} %||6-99-16||

\twolineshloka
{सर्वथा सर्वभूतानां नास्ति मृत्युरलक्षणः}
{तव तावदयं मृत्युर्मैथिलीकृतलक्षणः} %||6-99-17||

\twolineshloka
{मैथिली सह रामेण विशोका विहरिष्यति}
{अल्पपुण्या त्वहं घोरे पतिता शोकसागरे} %||6-99-18||

\twolineshloka
{कैलासे मन्दरे मेरौ तथा चैत्ररथे वने}
{देवोद्यानेषु सर्वेषु विहृत्य सहिता त्वया} %||6-99-19||

\threelineshloka
{विमानेनानुरूपेण या याम्यतुलया श्रिया}
{पश्यन्ती विविधान्देशांस्तांस्तांश्चित्रस्रगम्बरा}
{भ्रंशिता कामभोगेभ्यः सास्मि वीरवधात्तव} %||6-99-20||

\twolineshloka
{सत्यवाक्स महाभागो देवरो मे यदब्रवीत्}
{अयं राक्षसमुख्यानां विनाशः पर्युपस्थितः} %||6-99-21||

\twolineshloka
{कामक्रोधसमुत्थेन व्यसनेन प्रसङ्गिना}
{त्वया कृतमिदं सर्वमनाथं रक्षसां कुलम्} %||6-99-22||

\twolineshloka
{न हि त्वं शोचितव्यो मे प्रख्यातबलपौरुषः}
{स्त्रीस्वभावात्तु मे बुद्धिः कारुण्ये परिवर्तते} %||6-99-23||

\twolineshloka
{सुकृतं दुष्कृतं च त्वं गृहीत्वा स्वां गतिं गतः}
{आत्मानमनुशोचामि त्वद्वियोगेन दुःखिताम्} %||6-99-24||

\threelineshloka
{नीलजीमूतसङ्काशः पीताम्बरशुभाङ्गदः}
{सर्वगात्राणि विक्षिप्य किं शेषे रुधिराप्लुतः}
{प्रसुप्त इव शोकार्तां किं मां न प्रतिभाषसे} %||6-99-25||

\twolineshloka
{महावीर्यस्य दक्षस्य संयुगेष्वपलायिनः}
{यातुधानस्य दौहित्रीं किं त्वं मां नाभ्युदीक्षसे} %||6-99-26||

\twolineshloka
{येन सूदयसे शत्रून्समरे सूर्यवर्चसा}
{वज्रो वज्रधरस्येव सोऽयं ते सततार्चितः} %||6-99-27||

\twolineshloka
{रणे शत्रुप्रहरणो हेमजालपरिष्कृतः}
{परिघो व्यवकीर्णस्ते बाणैश्छिन्नः सहस्रधा} %||6-99-28||

\twolineshloka
{धिगस्तु हृदयं यस्या ममेदं न सहस्रधा}
{त्वयि पञ्चत्वमापन्ने फलते शोकपीडितम्} %||6-99-29||

\twolineshloka
{एतस्मिन्नन्तरे रामो विभीषणमुवाच ह}
{संस्कारः क्रियतां भ्रातुः स्त्रियश्चैता निवर्तय} %||6-99-30||

\threelineshloka
{तं प्रश्रितस्ततो रामं श्रुतवाक्यो विभीषणः}
{विमृश्य बुद्ध्या धर्मज्ञो धर्मार्थसहितं वचः}
{रामस्यैवानुवृत्त्यर्थमुत्तरं प्रत्यभाषत} %||6-99-31||

\twolineshloka
{त्यक्तधर्मव्रतं क्रूरं नृशंसमनृतं तथा}
{नाहमर्होऽस्मि संस्कर्तुं परदाराभिमर्शकम्} %||6-99-32||

\twolineshloka
{भ्रातृरूपो हि मे शत्रुरेष सर्वाहिते रतः}
{रावणो नार्हते पूजां पूज्योऽपि गुरुगौरवात्} %||6-99-33||

\twolineshloka
{नृशंस इति मां राम वक्ष्यन्ति मनुजा भुवि}
{श्रुत्वा तस्य गुणान्सर्वे वक्ष्यन्ति सुकृतं पुनः} %||6-99-34||

\twolineshloka
{तच्छ्रुत्वा परमप्रीतो रामो धर्मभृतां वरः}
{विभीषणमुवाचेदं वाक्यज्ञो वाक्यकोविदम्} %||6-99-35||

\twolineshloka
{तवापि मे प्रियं कार्यं त्वत्प्रभवाच्च मे जितम्}
{अवश्यं तु क्षमं वाच्यो मया त्वं राक्षसेश्वर} %||6-99-36||

\twolineshloka
{अधर्मानृतसंयुक्तः काममेष निशाचरः}
{तेजस्वी बलवाञ्शूरः सङ्ग्रामेषु च नित्यशः} %||6-99-37||

\twolineshloka
{शतक्रतुमुखैर्देवैः श्रूयते न पराजितः}
{महात्मा बलसम्पन्नो रावणो लोकरावणः} %||6-99-38||

\twolineshloka
{मरणान्तानि वैराणि निर्वृत्तं नः प्रयोजनम्}
{क्रियतामस्य संस्कारो ममाप्येष यथा तव} %||6-99-39||

\twolineshloka
{त्वत्सकाशान्महाबाहो संस्कारं विधिपूर्वकम्}
{क्षिप्रमर्हति धर्मज्ञ त्वं यशोभाग्भविष्यसि} %||6-99-40||

\twolineshloka
{राघवस्य वचः श्रुत्वा त्वरमाणो विभीषणः}
{संस्कारेणानुरूपेण योजयामास रावणम्} %||6-99-41||

\twolineshloka
{स ददौ पावकं तस्य विधियुक्तं विभीषणः}
{ताः स्त्रियोऽनुनयामास सान्त्वमुक्त्वा पुनः पुनः} %||6-99-42||

\twolineshloka
{प्रविष्टासु च सर्वासु राक्षसीषु विभीषणः}
{रामपार्श्वमुपागम्य तदातिष्ठद्विनीतवत्} %||6-99-43||

\twolineshloka
{रामोऽपि सह सैन्येन ससुग्रीवः सलक्ष्मणः}
{हर्षं लेभे रिपुं हत्वा यथा वृत्रं शतक्रतुः} %||6-100-44||


{॥इत्यार्षे श्रीमद्रामायणे वाल्मीकीये आदिकाव्ये युद्धकाण्डे एकोनशततमः सर्गः॥}

\sect{शततमः सर्गः}

\twolineshloka
{ते रावणवधं दृष्ट्वा देवगन्धर्वदानवाः}
{जग्मुस्तैस्तैर्विमानैः स्वैः कथयन्तः शुभाः कथाः} %||6-100-1||

\twolineshloka
{रावणस्य वधं घोरं राघवस्य पराक्रमम्}
{सुयुद्धं वानराणां च सुग्रीवस्य च मन्त्रितम्} %||6-100-2||

\twolineshloka
{अनुरागं च वीर्यं च सौमित्रेर्लक्ष्मणस्य च}
{कथयन्तो महाभागा जग्मुर्हृष्टा यथागतम्} %||6-100-3||

\twolineshloka
{राघवस्तु रथं दिव्यमिन्द्रदत्तं शिखिप्रभम्}
{अनुज्ञाय महाभागो मातलिं प्रत्यपूजयत्} %||6-100-4||

\twolineshloka
{राघवेणाभ्यनुज्ञातो मातलिः शक्रसारथिः}
{दिव्यं तं रथमास्थाय दिवमेवारुरोह सः} %||6-100-5||

\twolineshloka
{तस्मिंस्तु दिवमारूढे सुरसारथिसत्तमे}
{राघवः परमप्रीतः सुग्रीवं परिषस्वजे} %||6-100-6||

\twolineshloka
{परिष्वज्य च सुग्रीवं लक्ष्मणेनाभिवादितः}
{पूज्यमानो हरिश्रेष्ठैराजगाम बलालयम्} %||6-100-7||

\twolineshloka
{अब्रवीच्च तदा रामः समीपपरिवर्तिनम्}
{सौमित्रिं सत्त्वसम्पन्नं लक्ष्मणं दीप्ततेजसम्} %||6-100-8||

\twolineshloka
{विभीषणमिमं सौम्य लङ्कायामभिषेचय}
{अनुरक्तं च भक्तं च मम चैवोपकारिणम्} %||6-100-9||

\twolineshloka
{एष मे परमः कामो यदिमं रावणानुजम्}
{लङ्कायां सौम्य पश्येयमभिषिक्तं विभीषणम्} %||6-100-10||

\twolineshloka
{एवमुक्तस्तु सौमित्री राघवेण महात्मना}
{तथेत्युक्त्वा तु संहृष्टः सौवर्णं घटमाददे} %||6-100-11||

\twolineshloka
{घटेन तेन सौमित्रिरभ्यषिञ्चद्विभीषणम्}
{लङ्कायां रक्षसां मध्ये राजानं रामशासनात्} %||6-100-12||

\twolineshloka
{अभ्यषिञ्चत्स धर्मात्मा शुद्धात्मानं विभीषणम्}
{तस्यामात्या जहृषिरे भक्ता ये चास्य राक्षसाः} %||6-100-13||

\twolineshloka
{दृष्ट्वाभिषिक्तं लङ्कायां राक्षसेन्द्रं विभीषणम्}
{राघवः परमां प्रीतिं जगाम सहलक्ष्मणः} %||6-100-14||

\twolineshloka
{स तद्राज्यं महत्प्राप्य रामदत्तं विभीषणः}
{प्रकृतीः सान्त्वयित्वा च ततो राममुपागमत्} %||6-100-15||

\twolineshloka
{अक्षतान्मोदकाँल्लाजान्दिव्याः सुमनसस्तथा}
{आजह्रुरथ संहृष्टाः पौरास्तस्मै निशाचराः} %||6-100-16||

\twolineshloka
{स तान्गृहीत्वा दुर्धर्षो राघवाय न्यवेदयत्}
{मङ्गल्यं मङ्गलं सर्वं लक्ष्मणाय च वीर्यवान्} %||6-100-17||

\twolineshloka
{कृतकार्यं समृद्धार्थं दृष्ट्वा रामो विभीषणम्}
{प्रतिजग्राह तत्सर्वं तस्यैव प्रियकाम्यया} %||6-100-18||

\twolineshloka
{ततः शैलोपमं वीरं प्राञ्जलिं पार्श्वतः स्थितम्}
{अब्रवीद्राघवो वाक्यं हनूमन्तं प्लवङ्गमम्} %||6-100-19||

\twolineshloka
{अनुमान्य महाराजमिमं सौम्य विभीषणम्}
{प्रविश्य रावणगृहं विनयेनोपसृत्य च} %||6-100-20||

\twolineshloka
{वैदेह्या मां कुशलिनं ससुग्रीवं सलक्ष्मणम्}
{आचक्ष्व जयतां श्रेष्ठ रावणं च मया हतम्} %||6-100-21||

\twolineshloka
{प्रियमेतदुदाहृत्य मैथिल्यास्त्वं हरीश्वर}
{प्रतिगृह्य च सन्देशमुपावर्तितुमर्हसि} %||6-101-22||


{॥इत्यार्षे श्रीमद्रामायणे वाल्मीकीये आदिकाव्ये युद्धकाण्डे शततमः सर्गः॥}

\sect{एकाधिकशततमः सर्गः}

\twolineshloka
{इति प्रतिसमादिष्टो हनूमान्मारुतात्मजः}
{प्रविवेश पुरीं लङ्कां पूज्यमानो निशाचरैः} %||6-101-1||

\twolineshloka
{प्रविश्य तु महातेजा रावणस्य निवेशनम्}
{ददर्श शशिना हीनां सातङ्कामिव रोहिणीम्} %||6-101-2||

\twolineshloka
{निभृतः प्रणतः प्रह्वः सोऽभिगम्याभिवाद्य च}
{रामस्य वचनं सर्वमाख्यातुमुपचक्रमे} %||6-101-3||

\twolineshloka
{वैदेहि कुशली रामः ससुग्रीवः सलक्ष्मणः}
{कुशलं चाह सिद्धार्थो हतशत्रुररिन्दमः} %||6-101-4||

\twolineshloka
{विभीषणसहायेन रामेण हरिभिः सह}
{निहतो रावणो देवि लक्ष्मणस्य नयेन च} %||6-101-5||

\twolineshloka
{पृष्ट्वा च कुशलं रामो वीरस्त्वां रघुनन्दनः}
{अब्रवीत्परमप्रीतः कृतार्थेनान्तरात्मना} %||6-101-6||

\twolineshloka
{प्रियमाख्यामि ते देवि त्वां तु भूयः सभाजये}
{दिष्ट्या जीवसि धर्मज्ञे जयेन मम संयुगे} %||6-101-7||

\twolineshloka
{लब्धो नो विजयः सीते स्वस्था भव गतव्यथा}
{रावणः स हतः शत्रुर्लङ्का चेयं वशे स्थिता} %||6-101-8||

\twolineshloka
{मया ह्यलब्धनिद्रेण धृतेन तव निर्जये}
{प्रतिज्ञैषा विनिस्तीर्णा बद्ध्वा सेतुं महोदधौ} %||6-101-9||

\twolineshloka
{सम्भ्रमश्च न कर्तव्यो वर्तन्त्या रावणालये}
{विभीषणविधेयं हि लङ्कैश्वर्यमिदं कृतम्} %||6-101-10||

\twolineshloka
{तदाश्वसिहि विश्वस्ता स्वगृहे परिवर्तसे}
{अयं चाभ्येति संहृष्टस्त्वद्दर्शनसमुत्सुकः} %||6-101-11||

\twolineshloka
{एवमुक्ता समुत्पत्य सीता शशिनिभानना}
{प्रहर्षेणावरुद्धा सा व्याजहार न किञ्चन} %||6-101-12||

\twolineshloka
{अब्रवीच्च हरिश्रेष्ठः सीतामप्रतिजल्पतीम्}
{किं त्वं चिन्तयसे देवि किं च मां नाभिभाषसे} %||6-101-13||

\twolineshloka
{एवमुक्ता हनुमता सीता धर्मे व्यवस्थिता}
{अब्रवीत्परमप्रीता हर्षगद्गदया गिरा} %||6-101-14||

\twolineshloka
{प्रियमेतदुपश्रुत्य भर्तुर्विजयसंश्रितम्}
{प्रहर्षवशमापन्ना निर्वाक्यास्मि क्षणान्तरम्} %||6-101-15||

\twolineshloka
{न हि पश्यामि सदृशं चिन्तयन्ती प्लवङ्गम}
{मत्प्रियाख्यानकस्येह तव प्रत्यभिनन्दनम्} %||6-101-16||

\twolineshloka
{न च पश्यामि तत्सौम्य पृथिव्यामपि वानर}
{सदृशं मत्प्रियाख्याने तव दातुं भवेत्समम्} %||6-101-17||

\twolineshloka
{हिरण्यं वा सुवर्णं वा रत्नानि विविधानि च}
{राज्यं वा त्रिषु लोकेषु नैतदर्हति भाषितुम्} %||6-101-18||

\twolineshloka
{एवमुक्तस्तु वैदेह्या प्रत्युवाच प्लवङ्गमः}
{प्रगृहीताञ्जलिर्वाक्यं सीतायाः प्रमुखे स्थितः} %||6-101-19||

\twolineshloka
{भर्तुः प्रियहिते युक्ते भर्तुर्विजयकाङ्क्षिणि}
{स्निग्धमेवंविधं वाक्यं त्वमेवार्हसि भाषितुम्} %||6-101-20||

\twolineshloka
{तवैतद्वचनं सौम्ये सारवत्स्निग्धमेव च}
{रत्नौघाद्विविधाच्चापि देवराज्याद्विशिष्यते} %||6-101-21||

\twolineshloka
{अर्थतश्च मया प्राप्ता देवराज्यादयो गुणाः}
{हतशत्रुं विजयिनं रामं पश्यामि यत्स्थितम्} %||6-101-22||

\twolineshloka
{इमास्तु खलु राक्षस्यो यदि त्वमनुमन्यसे}
{हन्तुमिच्छाम्यहं सर्वा याभिस्त्वं तर्जिता पुरा} %||6-101-23||

\twolineshloka
{क्लिश्यन्तीं पतिदेवां त्वामशोकवनिकां गताम्}
{घोररूपसमाचाराः क्रूराः क्रूरतरेक्षणाः} %||6-101-24||

\twolineshloka
{राक्षस्यो दारुणकथा वरमेतं प्रयच्छ मे}
{इच्छामि विविधैर्घातैर्हन्तुमेताः सुदारुणाः} %||6-101-25||

\twolineshloka
{मुष्टिभिः पाणिभिश्चैव चरणैश्चैव शोभने}
{घोरैर्जानुप्रहारैश्च दशनानां च पातनैः} %||6-101-26||

\twolineshloka
{भक्षणैः कर्णनासानां केशानां लुञ्चनैस्तथा}
{भृशं शुष्कमुखीभिश्च दारुणैर्लङ्घनैर्हतैः} %||6-101-27||

\twolineshloka
{एवम्प्रकारैर्बहुभिर्विप्रकारैर्यशस्विनि}
{हन्तुमिच्छाम्यहं देवि तवेमाः कृतकिल्बिषाः} %||6-101-28||

\twolineshloka
{एवमुक्ता महुमता वैदेही जनकात्मजा}
{उवाच धर्मसहितं हनूमन्तं यशस्विनी} %||6-101-29||

\twolineshloka
{राजसंश्रयवश्यानां कुर्वतीनां पराज्ञया}
{विधेयानां च दासीनां कः कुप्येद्वानरोत्तम} %||6-101-30||

\twolineshloka
{भाग्यवैषम्ययोगेन पुरा दुश्चरितेन च}
{मयैतत्प्राप्यते सर्वं स्वकृतं ह्युपभुज्यते} %||6-101-31||

\twolineshloka
{प्राप्तव्यं तु दशायोगान्मयैतदिति निश्चितम्}
{दासीनां रावणस्याहं मर्षयामीह दुर्बला} %||6-101-32||

\twolineshloka
{आज्ञप्ता रावणेनैता राक्षस्यो मामतर्जयन्}
{हते तस्मिन्न कुर्युर्हि तर्जनं वानरोत्तम} %||6-101-33||

\twolineshloka
{अयं व्याघ्रसमीपे तु पुराणो धर्मसंहितः}
{ऋक्षेण गीतः श्लोको मे तं निबोध प्लवङ्गम} %||6-101-34||

\twolineshloka
{न परः पापमादत्ते परेषां पापकर्मणाम्}
{समयो रक्षितव्यस्तु सन्तश्चारित्रभूषणाः} %||6-101-35||

\twolineshloka
{पापानां वा शुभानां वा वधार्हाणां प्लवङ्गम}
{कार्यं कारुण्यमार्येण न कश्चिन्नापराध्यति} %||6-101-36||

\twolineshloka
{लोकहिंसाविहाराणां रक्षसां कामरूपिणम्}
{कुर्वतामपि पापानि नैव कार्यमशोभनम्} %||6-101-37||

\twolineshloka
{एवमुक्तस्तु हनुमान्सीतया वाक्यकोविदः}
{प्रत्युवाच ततः सीतां रामपत्नीं यशस्विनीम्} %||6-101-38||

\twolineshloka
{युक्ता रामस्य भवती धर्मपत्नी यशस्विनी}
{प्रतिसन्दिश मां देवि गमिष्ये यत्र राघवः} %||6-101-39||

\twolineshloka
{एवमुक्ता हनुमता वैदेही जनकात्मजा}
{अब्रवीद्द्रष्टुमिच्छामि भर्तारं वानरोत्तम} %||6-101-40||

\twolineshloka
{तस्यास्तद्वचनं श्रुत्वा हनुमान्पवनात्मजः}
{हर्षयन्मैथिलीं वाक्यमुवाचेदं महाद्युतिः} %||6-101-41||

\twolineshloka
{पूर्णचन्द्राननं रामं द्रक्ष्यस्यार्ये सलक्ष्मणम्}
{स्थिरमित्रं हतामित्रं शचीव त्रिदशेश्वरम्} %||6-101-42||

\twolineshloka
{तामेवमुक्त्वा राजन्तीं सीतां साक्षादिव श्रियम्}
{आजगाम महावेगो हनूमान्यत्र राघवः} %||6-102-43||


{॥इत्यार्षे श्रीमद्रामायणे वाल्मीकीये आदिकाव्ये युद्धकाण्डे एकाधिकशततमः सर्गः॥}

\sect{द्व्यधिकशततमः सर्गः}

\twolineshloka
{स उवाच महाप्रज्ञमभिगम्य प्लवङ्गमः}
{रामं वचनमर्थज्ञो वरं सर्वधनुष्मताम्} %||6-102-1||

\twolineshloka
{यन्निमित्तोऽयमारम्भः कर्मणां च फलोदयः}
{तां देवीं शोकसन्तप्तां मैथिलीं द्रष्टुमर्हसि} %||6-102-2||

\twolineshloka
{सा हि शोकसमाविष्टा बाष्पपर्याकुलेक्षणा}
{मैथिली विजयं श्रुत्वा तव हर्षमुपागमत्} %||6-102-3||

\twolineshloka
{पूर्वकात्प्रत्ययाच्चाहमुक्तो विश्वस्तया तया}
{भर्तारं द्रष्टुमिच्छामि कृतार्थं सहलक्ष्मणम्} %||6-102-4||

\twolineshloka
{एवमुक्तो हनुमता रामो धर्मभृतां वरः}
{अगच्छत्सहसा ध्यानमासीद्बाष्पपरिप्लुतः} %||6-102-5||

\twolineshloka
{दीर्घमुष्णं च निश्वस्य मेदिनीमवलोकयन्}
{उवाच मेघसङ्काशं विभीषणमुपस्थितम्} %||6-102-6||

\twolineshloka
{दिव्याङ्गरागां वैदेहीं दिव्याभरणभूषिताम्}
{इह सीतां शिरःस्नातामुपस्थापय माचिरम्} %||6-102-7||

\twolineshloka
{एवमुक्तस्तु रामेण त्वरमाणो विभीषणः}
{प्रविश्यान्तःपुरं सीतां स्त्रीभिः स्वाभिरचोदयत्} %||6-102-8||

\twolineshloka
{दिव्याङ्गरागा वैदेही दिव्याभरणभूषिता}
{यानमारोह भद्रं ते भर्ता त्वां द्रष्टुमिच्छति} %||6-102-9||

\twolineshloka
{एवमुक्ता तु वैदेही प्रत्युवाच विभीषणम्}
{अस्नाता द्रष्टुमिच्छामि भर्तारं राक्षसाधिप} %||6-102-10||

\twolineshloka
{तस्यास्तद्वचनं श्रुत्वा प्रत्युवाच विभीषणः}
{यथाह रामो भर्ता ते तत्तथा कर्तुमर्हसि} %||6-102-11||

\twolineshloka
{तस्य तद्वचनं श्रुत्वा मैथिली भर्तृदेवता}
{भर्तृभक्तिव्रता साध्वी तथेति प्रत्यभाषत} %||6-102-12||

\twolineshloka
{ततः सीतां शिरःस्नातां युवतीभिरलङ्कृताम्}
{महार्हाभरणोपेतां महार्हाम्बरधारिणीम्} %||6-102-13||

\twolineshloka
{आरोप्य शिबिकां दीप्तां परार्ध्याम्बरसंवृताम्}
{रक्षोभिर्बहुभिर्गुप्तामाजहार विभीषणः} %||6-102-14||

\twolineshloka
{सोऽभिगम्य महात्मानं ज्ञात्वाभिध्यानमास्थितम्}
{प्रणतश्च प्रहृष्टश्च प्राप्तां सीतां न्यवेदयत्} %||6-102-15||

\twolineshloka
{तामागतामुपश्रुत्य रक्षोगृहचिरोषिताम्}
{हर्षो दैन्यं च रोषश्च त्रयं राघवमाविशत्} %||6-102-16||

\twolineshloka
{ततः पार्श्वगतं दृष्ट्वा सविमर्शं विचारयन्}
{विभीषणमिदं वाक्यमहृष्टो राघवोऽब्रवीत्} %||6-102-17||

\twolineshloka
{राक्षसाधिपते सौम्य नित्यं मद्विजये रत}
{वैदेही संनिकर्षं मे शीघ्रं समुपगच्छतु} %||6-102-18||

\twolineshloka
{स तद्वचनमाज्ञाय राघवस्य विभीषणः}
{तूर्णमुत्सारणे यत्नं कारयामास सर्वतः} %||6-102-19||

\twolineshloka
{कञ्चुकोष्णीषिणस्तत्र वेत्रझर्झरपाणयः}
{उत्सारयन्तः पुरुषाः समन्तात्परिचक्रमुः} %||6-102-20||

\twolineshloka
{ऋक्षाणां वानराणां च राक्षसानां च सर्वतः}
{वृन्दान्युत्सार्यमाणानि दूरमुत्ससृजुस्ततः} %||6-102-21||

\twolineshloka
{तेषामुत्सार्यमाणानां सर्वेषां ध्वनिरुत्थितः}
{वायुनोद्वर्तमानस्य सागरस्येव निस्वनः} %||6-102-22||

\twolineshloka
{उत्सार्यमाणांस्तान्दृष्ट्वा समन्ताज्जातसम्भ्रमान्}
{दाक्षिण्यात्तदमर्षाच्च वारयामास राघवः} %||6-102-23||

\twolineshloka
{संरब्धश्चाब्रवीद्रामश्चक्षुषा प्रदहन्निव}
{विभीषणं महाप्राज्ञं सोपालम्भमिदं वचः} %||6-102-24||

\twolineshloka
{किमर्थं मामनादृत्य कृश्यतेऽयं त्वया जनः}
{निवर्तयैनमुद्योगं जनोऽयं स्वजनो मम} %||6-102-25||

\twolineshloka
{न गृहाणि न वस्त्राणि न प्राकारास्तिरस्क्रियाः}
{नेदृशा राजसत्कारा वृत्तमावरणं स्त्रियः} %||6-102-26||

\twolineshloka
{व्यसनेषु न कृच्छ्रेषु न युद्धे न स्वयंवरे}
{न क्रतौ नो विवाहे च दर्शनं दुष्यते स्त्रियः} %||6-102-27||

\twolineshloka
{सैषा युद्धगता चैव कृच्छ्रे महति च स्थिता}
{दर्शनेऽस्या न दोषः स्यान्मत्समीपे विशेषतः} %||6-102-28||

\twolineshloka
{तदानय समीपं मे शीघ्रमेनां विभीषण}
{सीता पश्यतु मामेषा सुहृद्गणवृतं स्थितम्} %||6-102-29||

\twolineshloka
{एवमुक्तस्तु रामेण सविमर्शो विभीषणः}
{रामस्योपानयत्सीतां संनिकर्षं विनीतवत्} %||6-102-30||

\twolineshloka
{ततो लक्ष्मणसुग्रीवौ हनूमांश्च प्लवङ्गमः}
{निशम्य वाक्यं रामस्य बभूवुर्व्यथिता भृशम्} %||6-102-31||

\twolineshloka
{कलत्रनिरपेक्षैश्च इङ्गितैरस्य दारुणैः}
{अप्रीतमिव सीतायां तर्कयन्ति स्म राघवम्} %||6-102-32||

\twolineshloka
{लज्जया त्ववलीयन्ती स्वेषु गात्रेषु मैथिली}
{विभीषणेनानुगता भर्तारं साभ्यवर्तत} %||6-102-33||

\twolineshloka
{सा वस्त्रसंरुद्धमुखी लज्जया जनसंसदि}
{रुरोदासाद्य भर्तारमार्यपुत्रेति भाषिणी} %||6-102-34||

\twolineshloka
{विस्मयाच्च प्रहर्षाच्च स्नेहाच्च परिदेवता}
{उदैक्षत मुखं भर्तुः सौम्यं सौम्यतरानना} %||6-102-35||

\fourlineindentedshloka
{अथ समपनुदन्मनःक्लमं सा}
{सुचिरमदृष्टमुदीक्ष्य वै प्रियस्य}
{वदनमुदितपूर्णचन्द्रकान्तम्}
{विमलशशाङ्कनिभानना तदासीत्} %||6-103-36||


{॥इत्यार्षे श्रीमद्रामायणे वाल्मीकीये आदिकाव्ये युद्धकाण्डे द्व्यधिकशततमः सर्गः॥}

\sect{त्र्यधिकशततमः सर्गः}

\twolineshloka
{तां तु पार्श्वे स्थितां प्रह्वां रामः सम्प्रेक्ष्य मैथिलीम्}
{हृदयान्तर्गतक्रोधो व्याहर्तुमुपचक्रमे} %||6-103-1||

\twolineshloka
{एषासि निर्जिता भद्रे शत्रुं जित्वा मया रणे}
{पौरुषाद्यदनुष्ठेयं तदेतदुपपादितम्} %||6-103-2||

\twolineshloka
{गतोऽस्म्यन्तममर्षस्य धर्षणा सम्प्रमार्जिता}
{अवमानश्च शत्रुश्च मया युगपदुद्धृतौ} %||6-103-3||

\twolineshloka
{अद्य मे पौरुषं दृष्टमद्य मे सफलः श्रमः}
{अद्य तीर्णप्रतिज्ञत्वात्प्रभवामीह चात्मनः} %||6-103-4||

\twolineshloka
{या त्वं विरहिता नीता चलचित्तेन रक्षसा}
{दैवसम्पादितो दोषो मानुषेण मया जितः} %||6-103-5||

\twolineshloka
{सम्प्राप्तमवमानं यस्तेजसा न प्रमार्जति}
{कस्तस्य पुरुषार्थोऽस्ति पुरुषस्याल्पतेजसः} %||6-103-6||

\twolineshloka
{लङ्घनं च समुद्रस्य लङ्कायाश्चावमर्दनम्}
{सफलं तस्य तच्छ्लाघ्यमद्य कर्म हनूमतः} %||6-103-7||

\twolineshloka
{युद्धे विक्रमतश्चैव हितं मन्त्रयतश्च मे}
{सुग्रीवस्य ससैन्यस्य सफलोऽद्य परिश्रमः} %||6-103-8||

\twolineshloka
{निर्गुणं भ्रातरं त्यक्त्वा यो मां स्वयमुपस्थितः}
{विभीषणस्य भक्तस्य सफलोऽद्य परिश्रमः} %||6-103-9||

\twolineshloka
{इत्येवं ब्रुवतस्तस्य सीता रामस्य तद्वचः}
{मृगीवोत्फुल्लनयना बभूवाश्रुपरिप्लुता} %||6-103-10||

\twolineshloka
{पश्यतस्तां तु रामस्य भूयः क्रोधोऽभ्यवर्तत}
{प्रभूताज्यावसिक्तस्य पावकस्येव दीप्यतः} %||6-103-11||

\twolineshloka
{स बद्ध्वा भ्रुकुटिं वक्त्रे तिर्यक्प्रेक्षितलोचनः}
{अब्रवीत्परुषं सीतां मध्ये वानररक्षसाम्} %||6-103-12||

\twolineshloka
{यत्कर्तव्यं मनुष्येण धर्षणां परिमार्जता}
{तत्कृतं सकलं सीते शत्रुहस्तादमर्षणात्} %||6-103-13||

\twolineshloka
{निर्जिता जीवलोकस्य तपसा भावितात्मना}
{अगस्त्येन दुराधर्षा मुनिना दक्षिणेव दिक्} %||6-103-14||

\twolineshloka
{विदितश्चास्तु भद्रं ते योऽयं रणपरिश्रमः}
{स तीर्णः सुहृदां वीर्यान्न त्वदर्थं मया कृतः} %||6-103-15||

\twolineshloka
{रक्षता तु मया वृत्तमपवादं च सर्वशः}
{प्रख्यातस्यात्मवंशस्य न्यङ्गं च परिमार्जता} %||6-103-16||

\twolineshloka
{प्राप्तचारित्रसन्देहा मम प्रतिमुखे स्थिता}
{दीपो नेत्रातुरस्येव प्रतिकूलासि मे दृढम्} %||6-103-17||

\twolineshloka
{तद्गच्छ ह्यभ्यनुज्ञाता यतेष्टं जनकात्मजे}
{एता दश दिशो भद्रे कार्यमस्ति न मे त्वया} %||6-103-18||

\twolineshloka
{कः पुमान्हि कुले जातः स्त्रियं परगृहोषिताम्}
{तेजस्वि पुनरादद्यात्सुहृल्लेखेन चेतसा} %||6-103-19||

\twolineshloka
{रावणाङ्कपरिभ्रष्टां दृष्टां दुष्टेन चक्षुषा}
{कथं त्वां पुनरादद्यां कुलं व्यपदिशन्महत्} %||6-103-20||

\twolineshloka
{तदर्थं निर्जिता मे त्वं यशः प्रत्याहृतं मया}
{नास्ति मे त्वय्यभिष्वङ्गो यथेष्टं गम्यतामितः} %||6-103-21||

\twolineshloka
{इति प्रव्याहृतं भद्रे मयैतत्कृतबुद्धिना}
{लक्ष्मणे भरते वा त्वं कुरु बुद्धिं यथासुखम्} %||6-103-22||

\twolineshloka
{सुग्रीवे वानरेन्द्रे वा राक्षसेन्द्रे विभीषणे}
{निवेशय मनः सीते यथा वा सुखमात्मनः} %||6-103-23||

\twolineshloka
{न हि त्वां रावणो दृष्ट्वा दिव्यरूपां मनोरमाम्}
{मर्षयते चिरं सीते स्वगृहे परिवर्तिनीम्} %||6-103-24||

\fourlineindentedshloka
{ततः प्रियार्हश्वरणा तदप्रियम्}
{प्रियादुपश्रुत्य चिरस्य मैथिली}
{मुमोच बाष्पं सुभृशं प्रवेपिता}
{गजेन्द्रहस्ताभिहतेव वल्लरी} %||6-104-25||


{॥इत्यार्षे श्रीमद्रामायणे वाल्मीकीये आदिकाव्ये युद्धकाण्डे त्र्यधिकशततमः सर्गः॥}

\sect{चतुरधिकशततमः सर्गः}

\twolineshloka
{एवमुक्ता तु वैदेही परुषं लोमहर्षणम्}
{राघवेण सरोषेण भृशं प्रव्यथिताभवत्} %||6-104-1||

\twolineshloka
{सा तदश्रुतपूर्वं हि जने महति मैथिली}
{श्रुत्वा भर्तृवचो रूक्षं लज्जया व्रीडिताभवत्} %||6-104-2||

\twolineshloka
{प्रविशन्तीव गात्राणि स्वान्येव जनकात्मजा}
{वाक्शल्यैस्तैः सशल्येव भृशमश्रूण्यवर्तयत्} %||6-104-3||

\twolineshloka
{ततो बाष्पपरिक्लिष्टं प्रमार्जन्ती स्वमाननम्}
{शनैर्गद्गदया वाचा भर्तारमिदमब्रवीत्} %||6-104-4||

\twolineshloka
{किं मामसदृशं वाक्यमीदृशं श्रोत्रदारुणम्}
{रूक्षं श्रावयसे वीर प्राकृतः प्राकृतामिव} %||6-104-5||

\twolineshloka
{न तथास्मि महाबाहो यथा त्वमवगच्छसि}
{प्रत्ययं गच्छ मे स्वेन चारित्रेणैव ते शपे} %||6-104-6||

\twolineshloka
{पृथक्स्त्रीणां प्रचारेण जातिं त्वं परिशङ्कसे}
{परित्यजेमां शङ्कां तु यदि तेऽहं परीक्षिता} %||6-104-7||

\twolineshloka
{यद्यहं गात्रसंस्पर्शं गतास्मि विवशा प्रभो}
{कामकारो न मे तत्र दैवं तत्रापराध्यति} %||6-104-8||

\twolineshloka
{मदधीनं तु यत्तन्मे हृदयं त्वयि वर्तते}
{पराधीनेषु गात्रेषु किं करिष्याम्यनीश्वरा} %||6-104-9||

\twolineshloka
{सहसंवृद्धभावाच्च संसर्गेण च मानद}
{यद्यहं ते न विज्ञाता हता तेनास्मि शाश्वतम्} %||6-104-10||

\twolineshloka
{प्रेषितस्ते यदा वीरो हनूमानवलोककः}
{लङ्कास्थाहं त्वया वीर किं तदा न विसर्जिता} %||6-104-11||

\twolineshloka
{प्रत्यक्षं वानरेन्द्रस्य त्वद्वाक्यसमनन्तरम्}
{त्वया सन्त्यक्तया वीर त्यक्तं स्याज्जीवितं मया} %||6-104-12||

\twolineshloka
{न वृथा ते श्रमोऽयं स्यात्संशये न्यस्य जीवितम्}
{सुहृज्जनपरिक्लेशो न चायं निष्फलस्तव} %||6-104-13||

\twolineshloka
{त्वया तु नरशार्दूल क्रोधमेवानुवर्तता}
{लघुनेव मनुष्येण स्त्रीत्वमेव पुरस्कृतम्} %||6-104-14||

\twolineshloka
{अपदेशेन जनकान्नोत्पत्तिर्वसुधातलात्}
{मम वृत्तं च वृत्तज्ञ बहु ते न पुरस्कृतम्} %||6-104-15||

\twolineshloka
{न प्रमाणीकृतः पाणिर्बाल्ये बालेन पीडितः}
{मम भक्तिश्च शीलं च सर्वं ते पृष्ठतः कृतम्} %||6-104-16||

\twolineshloka
{एवं ब्रुवाणा रुदती बाष्पगद्गदभाषिणी}
{अब्रवील्लक्ष्मणं सीता दीनं ध्यानपरं स्थितम्} %||6-104-17||

\twolineshloka
{चितां मे कुरु सौमित्रे व्यसनस्यास्य भेषजम्}
{मिथ्यापवादोपहता नाहं जीवितुमुत्सहे} %||6-104-18||

\twolineshloka
{अप्रीतस्य गुणैर्भर्तुस्त्यक्तया जनसंसदि}
{या क्षमा मे गतिर्गन्तुं प्रवेक्ष्ये हव्यवाहनम्} %||6-104-19||

\twolineshloka
{एवमुक्तस्तु वैदेह्या लक्ष्मणः परवीरहा}
{अमर्षवशमापन्नो राघवाननमैक्षत} %||6-104-20||

\twolineshloka
{स विज्ञाय मनश्छन्दं रामस्याकारसूचितम्}
{चितां चकार सौमित्रिर्मते रामस्य वीर्यवान्} %||6-104-21||

\twolineshloka
{अधोमुखं ततो रामं शनैः कृत्वा प्रदक्षिणम्}
{उपासर्पत वैदेही दीप्यमानं हुताशनम्} %||6-104-22||

\twolineshloka
{प्रणम्य देवताभ्यश्च ब्राह्मणेभ्यश्च मैथिली}
{बद्धाञ्जलिपुटा चेदमुवाचाग्निसमीपतः} %||6-104-23||

\twolineshloka
{यथा मे हृदयं नित्यं नापसर्पति राघवात्}
{तथा लोकस्य साक्षी मां सर्वतः पातु पावकः} %||6-104-24||

\twolineshloka
{एवमुक्त्वा तु वैदेही परिक्रम्य हुताशनम्}
{विवेश ज्वलनं दीप्तं निःसङ्गेनान्तरात्मना} %||6-104-25||

\twolineshloka
{जनः स सुमहांस्तत्र बालवृद्धसमाकुलः}
{ददर्श मैथिलीं तत्र प्रविशन्तीं हुताशनम्} %||6-104-26||

\twolineshloka
{तस्यामग्निं विशन्त्यां तु हाहेति विपुलः स्वनः}
{रक्षसां वानराणां च सम्बभूवाद्भुतोपमः} %||6-105-27||


{॥इत्यार्षे श्रीमद्रामायणे वाल्मीकीये आदिकाव्ये युद्धकाण्डे चतुरधिकशततमः सर्गः॥}

\sect{पञ्चाधिकशततमः सर्गः}

\twolineshloka
{ततो वैश्रवणो राजा यमश्चामित्रकर्शनः}
{सहस्राक्षो महेन्द्रश्च वरुणश्च परन्तपः} %||6-105-1||

\twolineshloka
{षडर्धनयनः श्रीमान्महादेवो वृषध्वजः}
{कर्ता सर्वस्य लोकस्य ब्रह्मा ब्रह्मविदां वरः} %||6-105-2||

\twolineshloka
{एते सर्वे समागम्य विमानैः सूर्यसंनिभैः}
{आगम्य नगरीं लङ्कामभिजग्मुश्च राघवम्} %||6-105-3||

\twolineshloka
{ततः सहस्ताभरणान्प्रगृह्य विपुलान्भुजान्}
{अब्रुवंस्त्रिदशश्रेष्ठाः प्राञ्जलिं राघवं स्थितम्} %||6-105-4||

\threelineshloka
{कर्ता सर्वस्य लोकस्य श्रेष्ठो ज्ञानवतां वरः}
{उपेक्षसे कथं सीतां पतन्तीं हव्यवाहने}
{कथं देवगणश्रेष्ठमात्मानं नावबुध्यसे} %||6-105-5||

\twolineshloka
{ऋतधामा वसुः पूर्वं वसूनां च प्रजापतिः}
{त्वं त्रयाणां हि लोकानामादिकर्ता स्वयम्प्रभुः} %||6-105-6||

\twolineshloka
{रुद्राणामष्टमो रुद्रः साध्यानामपि पञ्चमः}
{अश्विनौ चापि ते कर्णौ चन्द्रसूर्यौ च चक्षुषी} %||6-105-7||

\twolineshloka
{अन्ते चादौ च लोकानां दृश्यसे त्वं परन्तप}
{उपेक्षसे च वैदेहीं मानुषः प्राकृतो यथा} %||6-105-8||

\twolineshloka
{इत्युक्तो लोकपालैस्तैः स्वामी लोकस्य राघवः}
{अब्रवीत्त्रिदशश्रेष्ठान्रामो धर्मभृतां वरः} %||6-105-9||

\twolineshloka
{आत्मानं मानुषं मन्ये रामं दशरथात्मजम्}
{योऽहं यस्य यतश्चाहं भगवांस्तद्ब्रवीतु मे} %||6-105-10||

\twolineshloka
{इति ब्रुवाणं काकुत्स्थं ब्रह्मा ब्रह्मविदां वरः}
{अब्रवीच्छृणु मे राम सत्यं सत्यपराक्रम} %||6-105-11||

\twolineshloka
{भवान्नारायणो देवः श्रीमांश्चक्रायुधो विभुः}
{एकशृङ्गो वराहस्त्वं भूतभव्यसपत्नजित्} %||6-105-12||

\twolineshloka
{अक्षरं ब्रह्मसत्यं च मध्ये चान्ते च राघव}
{लोकानां त्वं परो धर्मो विष्वक्सेनश्चतुर्भुजः} %||6-105-13||

\twolineshloka
{शार्ङ्गधन्वा हृषीकेशः पुरुषः पुरुषोत्तमः}
{अजितः खड्गधृग्विष्णुः कृष्णश्चैव बृहद्बलः} %||6-105-14||

\twolineshloka
{सेनानीर्ग्रामणीश्च त्वं बुद्धिः सत्त्वं क्षमा दमः}
{प्रभवश्चाप्ययश्च त्वमुपेन्द्रो मधुसूदनः} %||6-105-15||

\twolineshloka
{इन्द्रकर्मा महेन्द्रस्त्वं पद्मनाभो रणान्तकृत्}
{शरण्यं शरणं च त्वामाहुर्दिव्या महर्षयः} %||6-105-16||

\twolineshloka
{सहस्रशृङ्गो वेदात्मा शतजिह्वो महर्षभः}
{त्वं यज्ञस्त्वं वषट्कारस्त्वमोङ्कारः परन्तप} %||6-105-17||

\twolineshloka
{प्रभवं निधनं वा ते न विदुः को भवानिति}
{दृश्यसे सर्वभूतेषु ब्राह्मणेषु च गोषु च} %||6-105-18||

\twolineshloka
{दिक्षु सर्वासु गगने पर्वतेषु वनेषु च}
{सहस्रचरणः श्रीमाञ्शतशीर्षः सहस्रधृक्} %||6-105-19||

\twolineshloka
{त्वं धारयसि भूतानि वसुधां च सपर्वताम्}
{अन्ते पृथिव्याः सलिले दृश्यसे त्वं महोरगः} %||6-105-20||

\twolineshloka
{त्रीँल्लोकान्धारयन्राम देवगन्धर्वदानवान्}
{अहं ते हृदयं राम जिह्वा देवी सरस्वती} %||6-105-21||

\twolineshloka
{देवा गात्रेषु लोमानि निर्मिता ब्रह्मणा प्रभो}
{निमेषस्तेऽभवद्रात्रिरुन्मेषस्तेऽभवद्दिवा} %||6-105-22||

\twolineshloka
{संस्कारास्तेऽभवन्वेदा न तदस्ति त्वया विना}
{जगत्सर्वं शरीरं ते स्थैर्यं ते वसुधातलम्} %||6-105-23||

\twolineshloka
{अग्निः कोपः प्रसादस्ते सोमः श्रीवत्सलक्षण}
{त्वया लोकास्त्रयः क्रान्ताः पुराणे विक्रमैस्त्रिभिः} %||6-105-24||

\twolineshloka
{महेन्द्रश्च कृतो राजा बलिं बद्ध्वा महासुरम्}
{सीता लक्ष्मीर्भवान्विष्णुर्देवः कृष्णः प्रजापतिः} %||6-105-25||

\twolineshloka
{वधार्थं रावणस्येह प्रविष्टो मानुषीं तनुम्}
{तदिदं नः कृतं कार्यं त्वया धर्मभृतां वर} %||6-105-26||

\twolineshloka
{निहतो रावणो राम प्रहृष्टो दिवमाक्रम}
{अमोघं बलवीर्यं ते अमोघस्ते पराक्रमः} %||6-105-27||

\twolineshloka
{अमोघास्ते भविष्यन्ति भक्तिमन्तश्च ये नराः}
{ये त्वां देवं ध्रुवं भक्ताः पुराणं पुरुषोत्तमम्} %||6-105-28||

{ये नराः कीर्तयिष्यन्ति नास्ति तेषां पराभवः॥\devanumber{29}॥} %||6-106-29|| (Check)

\addtocounter{shlokacount}{1}


{॥इत्यार्षे श्रीमद्रामायणे वाल्मीकीये आदिकाव्ये युद्धकाण्डे पञ्चाधिकशततमः सर्गः॥}

\sect{षष्ठाधिकशततमः सर्गः}

\twolineshloka
{एतच्छ्रुत्वा शुभं वाक्यं पितामहसमीरितम्}
{अङ्केनादाय वैदेहीमुत्पपात विभावसुः} %||6-106-1||

\twolineshloka
{तरुणादित्यसङ्काशां तप्तकाञ्चनभूषणाम्}
{रक्ताम्बरधरां बालां नीलकुञ्चितमूर्धजाम्} %||6-106-2||

\twolineshloka
{अक्लिष्टमाल्याभरणां तथा रूपां मनस्विनीम्}
{ददौ रामाय वैदेहीमङ्के कृत्वा विभावसुः} %||6-106-3||

\twolineshloka
{अब्रवीच्च तदा रामं साक्षी लोकस्य पावकः}
{एषा ते राम वैदेही पापमस्या न विद्यते} %||6-106-4||

\twolineshloka
{नैव वाचा न मनसा नानुध्यानान्न चक्षुषा}
{सुवृत्ता वृत्तशौण्डीरा न त्वामतिचचार ह} %||6-106-5||

\twolineshloka
{रावणेनापनीतैषा वीर्योत्सिक्तेन रक्षसा}
{त्वया विरहिता दीना विवशा निर्जनाद्वनात्} %||6-106-6||

\twolineshloka
{रुद्धा चान्तःपुरे गुप्ता त्वक्चित्ता त्वत्परायणा}
{रक्षिता राक्षसी सङ्घैर्विकृतैर्घोरदर्शनैः} %||6-106-7||

\twolineshloka
{प्रलोभ्यमाना विविधं भर्त्स्यमाना च मैथिली}
{नाचिन्तयत तद्रक्षस्त्वद्गतेनान्तरात्मना} %||6-106-8||

\twolineshloka
{विशुद्धभावां निष्पापां प्रतिगृह्णीष्व राघव}
{न किञ्चिदभिधातव्यमहमाज्ञापयामि ते} %||6-106-9||

\twolineshloka
{एवमुक्तो महातेजा धृतिमान्दृढविक्रमः}
{अब्रवीत्त्रिदशश्रेष्ठं रामो धर्मभृतां वरः} %||6-106-10||

\twolineshloka
{अवश्यं त्रिषु लोकेषु सीता पावनमर्हति}
{दीर्घकालोषिता चेयं रावणान्तःपुरे शुभा} %||6-106-11||

\twolineshloka
{बालिशः खलु कामात्मा रामो दशरथात्मजः}
{इति वक्ष्यन्ति मां सन्तो जानकीमविशोध्य हि} %||6-106-12||

\twolineshloka
{अनन्यहृदयां भक्तां मच्चित्तपरिरक्षणीम्}
{अहमप्यवगच्छामि मैथिलीं जनकात्मजाम्} %||6-106-13||

\twolineshloka
{प्रत्ययार्थं तु लोकानां त्रयाणां सत्यसंश्रयः}
{उपेक्षे चापि वैदेहीं प्रविशन्तीं हुताशनम्} %||6-106-14||

\twolineshloka
{इमामपि विशालाक्षीं रक्षितां स्वेन तेजसा}
{रावणो नातिवर्तेत वेलामिव महोदधिः} %||6-106-15||

\twolineshloka
{न हि शक्तः स दुष्टात्मा मनसापि हि मैथिलीम्}
{प्रधर्षयितुमप्राप्तां दीप्तामग्निशिखामिव} %||6-106-16||

\twolineshloka
{नेयमर्हति चैश्वर्यं रावणान्तःपुरे शुभा}
{अनन्या हि मया सीतां भास्करेण प्रभा यथा} %||6-106-17||

\twolineshloka
{विशुद्धा त्रिषु लोकेषु मैथिली जनकात्मजा}
{न हि हातुमियं शक्या कीर्तिरात्मवता यथा} %||6-106-18||

\twolineshloka
{अवश्यं च मया कार्यं सर्वेषां वो वचो हितम्}
{स्निग्धानां लोकमान्यानामेवं च ब्रुवतां हितम्} %||6-106-19||

\fourlineindentedshloka
{इतीदमुक्त्वा वचनं महाबलैः}
{प्रशस्यमानः स्वकृतेन कर्मणा}
{समेत्य रामः प्रियया महाबलः}
{सुखं सुखार्होऽनुबभूव राघवः} %||6-107-20||


{॥इत्यार्षे श्रीमद्रामायणे वाल्मीकीये आदिकाव्ये युद्धकाण्डे षष्ठाधिकशततमः सर्गः॥}

\sect{सप्तमाधिकशततमः सर्गः}

\twolineshloka
{एतच्छ्रुत्वा शुभं वाक्यं राघवेण सुभाषितम्}
{इदं शुभतरं वाक्यं व्याजहार महेश्वरः} %||6-107-1||

\twolineshloka
{पुष्कराक्ष महाबाहो महावक्षः परन्तप}
{दिष्ट्या कृतमिदं कर्म त्वया शस्त्रभृतां वर} %||6-107-2||

\twolineshloka
{दिष्ट्या सर्वस्य लोकस्य प्रवृद्धं दारुणं तमः}
{अपावृत्तं त्वया सङ्ख्ये राम रावणजं भयम्} %||6-107-3||

\twolineshloka
{आश्वास्य भरतं दीनं कौसल्यां च यशस्विनीम्}
{कैकेयीं च सुमित्रां च दृष्ट्वा लक्ष्मणमातरम्} %||6-107-4||

\twolineshloka
{प्राप्य राज्यमयोध्यायां नन्दयित्वा सुहृज्जनम्}
{इक्ष्वाकूणां कुले वंशं स्थापयित्वा महाबल} %||6-107-5||

\twolineshloka
{इष्ट्वा तुरगमेधेन प्राप्य चानुत्तमं यशः}
{ब्राह्मणेभ्यो धनं दत्त्वा त्रिदिवं गन्तुमर्हसि} %||6-107-6||

\twolineshloka
{एष राजा विमानस्थः पिता दशरथस्तव}
{काकुत्स्थ मानुषे लोके गुरुस्तव महायशाः} %||6-107-7||

\twolineshloka
{इन्द्रलोकं गतः श्रीमांस्त्वया पुत्रेण तारितः}
{लक्ष्मणेन सह भ्रात्रा त्वमेनमभिवादय} %||6-107-8||

\twolineshloka
{महादेववचः श्रुत्वा काकुत्स्थः सहलक्ष्मणः}
{विमानशिखरस्थस्य प्रणाममकरोत्पितुः} %||6-107-9||

\twolineshloka
{दीप्यमानं स्वयां लक्ष्म्या विरजोऽम्बरधारिणम्}
{लक्ष्मणेन सह भ्रात्रा ददर्श पितरं प्रभुः} %||6-107-10||

\twolineshloka
{हर्षेण महताविष्टो विमानस्थो महीपतिः}
{प्राणैः प्रियतरं दृष्ट्वा पुत्रं दशरथस्तदा} %||6-107-11||

\twolineshloka
{आरोप्याङ्कं महाबाहुर्वरासनगतः प्रभुः}
{बाहुभ्यां सम्परिष्वज्य ततो वाक्यं समाददे} %||6-107-12||

\twolineshloka
{न मे स्वर्गो बहुमतः सम्मानश्च सुरर्षिभिः}
{त्वया राम विहीनस्य सत्यं प्रतिशृणोमि ते} %||6-107-13||

\twolineshloka
{कैकेय्या यानि चोक्तानि वाक्यानि वदतां वर}
{तव प्रव्राजनार्थानि स्थितानि हृदये मम} %||6-107-14||

\twolineshloka
{त्वां तु दृष्ट्वा कुशलिनं परिष्वज्य सलक्ष्मणम्}
{अद्य दुःखाद्विमुक्तोऽस्मि नीहारादिव भास्करः} %||6-107-15||

\twolineshloka
{तारितोऽहं त्वया पुत्र सुपुत्रेण महात्मना}
{अष्टावक्रेण धर्मात्मा तारितो ब्राह्मणो यथा} %||6-107-16||

\twolineshloka
{इदानीं च विजानामि यथा सौम्य सुरेश्वरैः}
{वधार्थं रावणस्येह विहितं पुरुषोत्तमम्} %||6-107-17||

\twolineshloka
{सिद्धार्था खलु कौसल्या या त्वां राम गृहं गतम्}
{वनान्निवृत्तं संहृष्टा द्रक्ष्यते शत्रुसूदन} %||6-107-18||

\twolineshloka
{सिद्धार्थाः खलु ते राम नरा ये त्वां पुरीं गतम्}
{जलार्द्रमभिषिक्तं च द्रक्ष्यन्ति वसुधाधिपम्} %||6-107-19||

\twolineshloka
{अनुरक्तेन बलिना शुचिना धर्मचारिणा}
{इच्छेयं त्वामहं द्रष्टुं भरतेन समागतम्} %||6-107-20||

\twolineshloka
{चतुर्दशसमाः सौम्य वने निर्यापितास्त्वया}
{वसता सीतया सार्धं लक्ष्मणेन च धीमता} %||6-107-21||

\twolineshloka
{निवृत्तवनवासोऽसि प्रतिज्ञा सफला कृता}
{रावणं च रणे हत्वा देवास्ते परितोषिताः} %||6-107-22||

\twolineshloka
{कृतं कर्म यशः श्लाघ्यं प्राप्तं ते शत्रुसूदन}
{भ्रातृभिः सह राज्यस्थो दीर्घमायुरवाप्नुहि} %||6-107-23||

\twolineshloka
{इति ब्रुवाणं राजानं रामः प्राञ्जलिरब्रवीत्}
{कुरु प्रसादं धर्मज्ञ कैकेय्या भरतस्य च} %||6-107-24||

\twolineshloka
{सपुत्रां त्वां त्यजामीति यदुक्ता कैकयी त्वया}
{स शापः कैकयीं घोरः सपुत्रां न स्पृशेत्प्रभो} %||6-107-25||

\twolineshloka
{स तथेति महाराजो राममुक्त्वा कृताञ्जलिम्}
{लक्ष्मणं च परिष्वज्य पुनर्वाक्यमुवाच ह} %||6-107-26||

\twolineshloka
{रामं शुश्रूषता भक्त्या वैदेह्या सह सीतया}
{कृता मम महाप्रीतिः प्राप्तं धर्मफलं च ते} %||6-107-27||

\twolineshloka
{धर्मं प्राप्स्यसि धर्मज्ञ यशश्च विपुलं भुवि}
{रामे प्रसन्ने स्वर्गं च महिमानं तथैव च} %||6-107-28||

\twolineshloka
{रामं शुश्रूष भद्रं ते सुमित्रानन्दवर्धन}
{रामः सर्वस्य लोकस्य शुभेष्वभिरतः सदा} %||6-107-29||

\twolineshloka
{एते सेन्द्रास्त्रयो लोकाः सिद्धाश्च परमर्षयः}
{अभिगम्य महात्मानमर्चन्ति पुरुषोत्तमम्} %||6-107-30||

\twolineshloka
{एतत्तदुक्तमव्यक्तमक्षरं ब्रह्मनिर्मितम्}
{देवानां हृदयं सौम्य गुह्यं रामः परन्तपः} %||6-107-31||

\twolineshloka
{अवाप्तं धर्मचरणं यशश्च विपुलं त्वया}
{रामं शुश्रूषता भक्त्या वैदेह्या सह सीतया} %||6-107-32||

\twolineshloka
{स तथोक्त्वा महाबाहुर्लक्ष्मणं प्राञ्जलिं स्थितम्}
{उवाच राजा धर्मात्मा वैदेहीं वचनं शुभम्} %||6-107-33||

\twolineshloka
{कर्तव्यो न तु वैदेहि मन्युस्त्यागमिमं प्रति}
{रामेण त्वद्विशुद्ध्यर्थं कृतमेतद्धितैषिणा} %||6-107-34||

\twolineshloka
{न त्वं सुभ्रु समाधेया पतिशुश्रूवणं प्रति}
{अवश्यं तु मया वाच्यमेष ते दैवतं परम्} %||6-107-35||

\twolineshloka
{इति प्रतिसमादिश्य पुत्रौ सीतां तथा स्नुषाम्}
{इन्द्रलोकं विमानेन ययौ दशरथो ज्वलन्} %||6-108-36||


{॥इत्यार्षे श्रीमद्रामायणे वाल्मीकीये आदिकाव्ये युद्धकाण्डे सप्तमाधिकशततमः सर्गः॥}

\sect{अष्टमाधिकशततमः सर्गः}

\twolineshloka
{प्रतिप्रयाते काकुत्स्थे महेन्द्रः पाकशासनः}
{अब्रवीत्परमप्रीतो राघवं प्राञ्जलिं स्थितम्} %||6-108-1||

\twolineshloka
{अमोघं दर्शनं राम तवास्माकं परन्तप}
{प्रीतियुक्तोऽस्मि तेन त्वं ब्रूहि यन्मनसेच्छसि} %||6-108-2||

\twolineshloka
{एवमुक्तस्तु काकुत्स्थः प्रत्युवाच कृताञ्जलिः}
{लक्ष्मणेन सह भ्रात्रा सीतया चापि भार्यया} %||6-108-3||

\twolineshloka
{यदि प्रीतिः समुत्पन्ना मयि सर्वसुरेश्वर}
{वक्ष्यामि कुरु मे सत्यं वचनं वदतां वर} %||6-108-4||

\twolineshloka
{मम हेतोः पराक्रान्ता ये गता यमसादनम्}
{ते सर्वे जीवितं प्राप्य समुत्तिष्ठन्तु वानराः} %||6-108-5||

\twolineshloka
{मत्प्रियेष्वभिरक्ताश्च न मृत्युं गणयन्ति च}
{त्वत्प्रसादात्समेयुस्ते वरमेतदहं वृणे} %||6-108-6||

\twolineshloka
{नीरुजान्निर्व्रणांश्चैव सम्पन्नबलपौरुषान्}
{गोलाङ्गूलांस्तथैवर्क्षान्द्रष्टुमिच्छामि मानद} %||6-108-7||

\twolineshloka
{अकाले चापि मुख्यानि मूलानि च फलानि च}
{नद्यश्च विमलास्तत्र तिष्ठेयुर्यत्र वानराः} %||6-108-8||

\twolineshloka
{श्रुत्वा तु वचनं तस्य राघवस्य महात्मनः}
{महेन्द्रः प्रत्युवाचेदं वचनं प्रीतिलक्षणम्} %||6-108-9||

\twolineshloka
{महानयं वरस्तात त्वयोक्तो रघुनन्दन}
{समुत्थास्यन्ति हरयः सुप्ता निद्राक्षये यथा} %||6-108-10||

\twolineshloka
{सुहृद्भिर्बान्धवैश्चैव ज्ञातिभिः स्वजनेन च}
{सर्व एव समेष्यन्ति संयुक्ताः परया मुदा} %||6-108-11||

\twolineshloka
{अकाले पुष्पशबलाः फलवन्तश्च पादपाः}
{भविष्यन्ति महेष्वास नद्यश्च सलिलायुताः} %||6-108-12||

\twolineshloka
{सव्रणैः प्रथमं गात्रैः संवृतैर्निव्रणैः पुनः}
{बभूवुर्वानराः सर्वे किमेतदिति विस्मितः} %||6-108-13||

\twolineshloka
{काकुत्स्थं परिपूर्णार्थं दृष्ट्वा सर्वे सुरोत्तमाः}
{ऊचुस्ते प्रथमं स्तुत्वा स्तवार्हं सहलक्ष्मणम्} %||6-108-14||

\twolineshloka
{गच्छायोध्यामितो वीर विसर्जय च वानरान्}
{मैथिलीं सान्त्वयस्वैनामनुरक्तां तपस्विनीम्} %||6-108-15||

\twolineshloka
{भ्रातरं पश्य भरतं त्वच्छोकाद्व्रतचारिणम्}
{अभिषेचय चात्मानं पौरान्गत्वा प्रहर्षय} %||6-108-16||

\twolineshloka
{एवमुक्त्वा तमामन्त्र्य रामं सौमित्रिणा सह}
{विमानैः सूर्यसङ्काशैर्हृष्टा जग्मुः सुरा दिवम्} %||6-108-17||

\twolineshloka
{अभिवाद्य च काकुत्स्थः सर्वांस्तांस्त्रिदशोत्तमान्}
{लक्ष्मणेन सह भ्रात्रा वासमाज्ञापयत्तदा} %||6-108-18||

\fourlineindentedshloka
{ततस्तु सा लक्ष्मणरामपालिता}
{महाचमूर्हृष्टजना यशस्विनी}
{श्रिया ज्वलन्ती विरराज सर्वतो}
{निशाप्रणीतेव हि शीतरश्मिना} %||6-109-19||


{॥इत्यार्षे श्रीमद्रामायणे वाल्मीकीये आदिकाव्ये युद्धकाण्डे अष्टमाधिकशततमः सर्गः॥}

\sect{नवमाधिकशततमः सर्गः}

\twolineshloka
{तां रात्रिमुषितं रामं सुखोत्थितमरिन्दमम्}
{अब्रवीत्प्राञ्जलिर्वाक्यं जयं पृष्ट्वा विभीषणः} %||6-109-1||

\twolineshloka
{स्नानानि चाङ्गरागाणि वस्त्राण्याभरणानि च}
{चन्दनानि च दिव्यानि माल्यानि विविधानि च} %||6-109-2||

\twolineshloka
{अलङ्कारविदश्चेमा नार्यः पद्मनिभेक्षणाः}
{उपस्थितास्त्वां विधिवत्स्नापयिष्यन्ति राघव} %||6-109-3||

\twolineshloka
{एवमुक्तस्तु काकुत्स्थः प्रत्युवाच विभीषणम्}
{हरीन्सुग्रीवमुख्यांस्त्वं स्नानेनोपनिमन्त्रय} %||6-109-4||

\twolineshloka
{स तु ताम्यति धर्मात्मा ममहेतोः सुखोचितः}
{सुकुमारो महाबाहुः कुमारः सत्यसंश्रवः} %||6-109-5||

\twolineshloka
{तं विना कैकेयीपुत्रं भरतं धर्मचारिणम्}
{न मे स्नानं बहुमतं वस्त्राण्याभरणानि च} %||6-109-6||

\twolineshloka
{इत एव पथा क्षिप्रं प्रतिगच्छाम तां पुरीम्}
{अयोध्यामायतो ह्येष पन्थाः परमदुर्गमः} %||6-109-7||

\twolineshloka
{एवमुक्तस्तु काकुत्स्थं प्रत्युवाच विभीषणः}
{अह्ना त्वां प्रापयिष्यामि तां पुरीं पार्थिवात्मज} %||6-109-8||

\twolineshloka
{पुष्पकं नाम भद्रं ते विमानं सूर्यसंनिभम्}
{मम भ्रातुः कुबेरस्य रावणेनाहृतं बलात्} %||6-109-9||

\twolineshloka
{तदिदं मेघसङ्काशं विमानमिह तिष्ठति}
{तेन यास्यसि यानेन त्वमयोध्यां गजज्वरः} %||6-109-10||

\twolineshloka
{अहं ते यद्यनुग्राह्यो यदि स्मरसि मे गुणान्}
{वस तावदिह प्राज्ञ यद्यस्ति मयि सौहृदम्} %||6-109-11||

\twolineshloka
{लक्ष्मणेन सह भ्रात्रा वैदेह्या चापि भार्यया}
{अर्चितः सर्वकामैस्त्वं ततो राम गमिष्यसि} %||6-109-12||

\twolineshloka
{प्रीतियुक्तस्तु मे राम ससैन्यः ससुहृद्गणः}
{सत्क्रियां विहितां तावद्गृहाण त्वं मयोद्यताम्} %||6-109-13||

\twolineshloka
{प्रणयाद्बहुमानाच्च सौहृदेन च राघव}
{प्रसादयामि प्रेष्योऽहं न खल्वाज्ञापयामि ते} %||6-109-14||

\twolineshloka
{एवमुक्तस्ततो रामः प्रत्युवाच विभीषणम्}
{रक्षसां वानराणां च सर्वेषां चोपशृण्वताम्} %||6-109-15||

\twolineshloka
{पूजितोऽहं त्वया वीर साचिव्येन परन्तप}
{सर्वात्मना च चेष्टिभिः सौहृदेनोत्तमेन च} %||6-109-16||

\twolineshloka
{न खल्वेतन्न कुर्यां ते वचनं राक्षसेश्वर}
{तं तु मे भ्रातरं द्रष्टुं भरतं त्वरते मनः} %||6-109-17||

\twolineshloka
{मां निवर्तयितुं योऽसौ चित्रकूटमुपागतः}
{शिरसा याचतो यस्य वचनं न कृतं मया} %||6-109-18||

\twolineshloka
{कौसल्यां च सुमित्रां च कैकेयीं च यशस्विनीम्}
{गुरूंश्च सुहृदश्चैव पौरांश्च तनयैः सह} %||6-109-19||

\twolineshloka
{उपस्थापय मे क्षिप्रं विमानं राक्षसेश्वर}
{कृतकार्यस्य मे वासः कथञ्चिदिह सम्मतः} %||6-109-20||

\twolineshloka
{अनुजानीहि मां सौम्य पूजितोऽस्मि विभीषण}
{मन्युर्न खलु कर्तव्यस्त्वरितस्त्वानुमानये} %||6-109-21||

\twolineshloka
{ततः काञ्चनचित्राङ्गं वैदूर्यमणिवेदिकम्}
{कूटागारैः परिक्षिप्तं सर्वतो रजतप्रभम्} %||6-109-22||

\twolineshloka
{पाण्डुराभिः पताकाभिर्ध्वजैश्च समलङ्कृतम्}
{शोभितं काञ्चनैर्हर्म्यैर्हेमपद्मविभूषितम्} %||6-109-23||

\twolineshloka
{प्रकीर्णं किङ्किणीजालैर्मुक्तामणिगवाक्षितम्}
{घण्टाजालैः परिक्षिप्तं सर्वतो मधुरस्वनम्} %||6-109-24||

\twolineshloka
{तन्मेरुशिखराकारं निर्मितं विश्वकर्मणा}
{बहुभिर्भूषितं हर्म्यैर्मुक्तारजतसंनिभौ} %||6-109-25||

\twolineshloka
{तलैः स्फटिकचित्राङ्गैर्वैदूर्यैश्च वरासनैः}
{महार्हास्तरणोपेतैरुपपन्नं महाधनैः} %||6-109-26||

\twolineshloka
{उपस्थितमनाधृष्यं तद्विमानं मनोजवम्}
{निवेदयित्वा रामाय तस्थौ तत्र विभीषणः} %||6-110-27||


{॥इत्यार्षे श्रीमद्रामायणे वाल्मीकीये आदिकाव्ये युद्धकाण्डे नवमाधिकशततमः सर्गः॥}

\sect{दशमाधिकशततमः सर्गः}

\twolineshloka
{उपस्थितं तु तं दृष्ट्वा पुष्पकं पुष्पभूषितम्}
{अविदूरे स्थितं रामं प्रत्युवाच विभीषणः} %||6-110-1||

\twolineshloka
{स तु बद्धाञ्जलिः प्रह्वो विनीतो राक्षसेश्वरः}
{अब्रवीत्त्वरयोपेतः किं करोमीति राघवम्} %||6-110-2||

\twolineshloka
{तमब्रवीन्महातेजा लक्ष्मणस्योपशृण्वतः}
{विमृश्य राघवो वाक्यमिदं स्नेहपुरस्कृतम्} %||6-110-3||

\twolineshloka
{कृतप्रयत्नकर्माणो विभीषण वनौकसः}
{रत्नैरर्थैश्च विविभैर्भूषणैश्चाभिपूजय} %||6-110-4||

\twolineshloka
{सहैभिरर्दिता लङ्का निर्जिता राक्षसेश्वर}
{हृष्टैः प्राणभयं त्यक्त्वा सङ्ग्रामेष्वनिवर्तिभिः} %||6-110-5||

\twolineshloka
{एवं सम्मानिताश्चेमे मानार्हा मानद त्वया}
{भविष्यन्ति कृतज्ञेन निर्वृता हरियूथपाः} %||6-110-6||

\twolineshloka
{त्यागिनं सङ्ग्रहीतारं सानुक्रोशं यशस्विनम्}
{यतस्त्वामवगच्छन्ति ततः सम्बोधयामि ते} %||6-110-7||

\twolineshloka
{एवमुक्तस्तु रामेण वानरांस्तान्विभीषणः}
{रत्नार्थैः संविभागेन सर्वानेवान्वपूजयत्} %||6-110-8||

\twolineshloka
{ततस्तान्पूजितान्दृष्ट्वा रत्नैरर्थैश्च यूथपान्}
{आरुरोह ततो रामस्तद्विमानमनुत्तमम्} %||6-110-9||

\twolineshloka
{अङ्केनादाय वैदेहीं लज्जमानां यशस्विनीम्}
{लक्ष्मणेन सह भ्रात्रा विक्रान्तेन धनुष्मता} %||6-110-10||

\twolineshloka
{अब्रवीच्च विमानस्थः काकुत्स्थः सर्ववानरान्}
{सुग्रीवं च महावीर्यं राक्षसं च विभीषणम्} %||6-110-11||

\twolineshloka
{मित्रकार्यं कृतमिदं भवद्भिर्वानरोत्तमाः}
{अनुज्ञाता मया सर्वे यथेष्टं प्रतिगच्छत} %||6-110-12||

\threelineshloka
{यत्तु कार्यं वयस्येन सुहृदा वा परन्तप}
{कृतं सुग्रीव तत्सर्वं भवता धर्मभीरुणा}
{किष्किन्धां प्रतियाह्याशु स्वसैन्येनाभिसंवृतः} %||6-110-13||

\twolineshloka
{स्वराज्ये वस लङ्कायां मया दत्ते विभीषण}
{न त्वां धर्षयितुं शक्ताः सेन्द्रा अपि दिवौकसः} %||6-110-14||

\twolineshloka
{अयोध्यां प्रतियास्यामि राजधानीं पितुर्मम}
{अभ्यनुज्ञातुमिच्छामि सर्वानामन्त्रयामि वः} %||6-110-15||

\threelineshloka
{एवमुक्तास्तु रामेण वानरास्ते महाबलाः}
{ऊचुः प्राञ्जलयो रामं राक्षसश्च विभीषणः}
{अयोध्यां गन्तुमिच्छामः सर्वान्नयतु नो भवान्} %||6-110-16||

\twolineshloka
{दृष्ट्वा त्वामभिषेकार्द्रं कौसल्यामभिवाद्य च}
{अचिरेणागमिष्यामः स्वान्गृहान्नृपतेः सुत} %||6-110-17||

\twolineshloka
{एवमुक्तस्तु धर्मात्मा वानरैः सविभीषणैः}
{अब्रवीद्राघवः श्रीमान्ससुग्रीवविभीषणान्} %||6-110-18||

\twolineshloka
{प्रियात्प्रियतरं लब्धं यदहं ससुहृज्जनः}
{सर्वैर्भवद्भिः सहितः प्रीतिं लप्स्ये पुरीं गतः} %||6-110-19||

\twolineshloka
{क्षिप्रमारोह सुग्रीव विमानं वानरैः सह}
{त्वमध्यारोह सामात्यो राक्षसेन्द्रविभीषण} %||6-110-20||

\twolineshloka
{ततस्तत्पुष्पकं दिव्यं सुग्रीवः सह सेनया}
{अध्यारोहत्त्वरञ्शीघ्रं सामात्यश्च विभीषणः} %||6-110-21||

\twolineshloka
{तेष्वारूढेषु सर्वेषु कौबेरं परमासनम्}
{राघवेणाभ्यनुज्ञातमुत्पपात विहायसम्} %||6-110-22||

\twolineshloka
{ययौ तेन विमानेन हंसयुक्तेन भास्वता}
{प्रहृष्टश्च प्रतीतश्च बभौ रामः कुबेरवत्} %||6-111-23||


{॥इत्यार्षे श्रीमद्रामायणे वाल्मीकीये आदिकाव्ये युद्धकाण्डे दशमाधिकशततमः सर्गः॥}

\sect{एकादशाधिकशततमः सर्गः}

\twolineshloka
{अनुज्ञातं तु रामेण तद्विमानमनुत्तमम्}
{उत्पपात महामेघः श्वसनेनोद्धतो यथा} %||6-111-1||

\twolineshloka
{पातयित्वा ततश्चक्षुः सर्वतो रघुनन्दनः}
{अब्रवीन्मैथिलीं सीतां रामः शशिनिभाननाम्} %||6-111-2||

\twolineshloka
{कैलासशिखराकारे त्रिकूटशिखरे स्थिताम्}
{लङ्कामीक्षस्व वैदेहि निर्मितां विश्वकर्मणा} %||6-111-3||

\twolineshloka
{एतदायोधनं पश्य मांसशोणितकर्दमम्}
{हरीणां राक्षसानां च सीते विशसनं महत्} %||6-111-4||

\twolineshloka
{तवहेतोर्विशालाक्षि रावणो निहतो मया}
{कुम्भकर्णोऽत्र निहतः प्रहस्तश्च निशाचरः} %||6-111-5||

\twolineshloka
{लक्ष्मणेनेन्द्रजिच्चात्र रावणिर्निहतो रणे}
{विरूपाक्षश्च दुष्प्रेक्ष्यो महापार्श्वमहोदरौ} %||6-111-6||

\twolineshloka
{अकम्पनश्च निहतो बलिनोऽन्ये च राक्षसाः}
{त्रिशिराश्चातिकायश्च देवान्तकनरान्तकौ} %||6-111-7||

\twolineshloka
{अत्र मन्दोदरी नाम भार्या तं पर्यदेवयत्}
{सपत्नीनां सहस्रेण सास्रेण परिवारिता} %||6-111-8||

\twolineshloka
{एतत्तु दृश्यते तीर्थं समुद्रस्य वरानने}
{यत्र सागरमुत्तीर्य तां रात्रिमुषिता वयम्} %||6-111-9||

\twolineshloka
{एष सेतुर्मया बद्धः सागरे सलिलार्णवे}
{तवहेतोर्विशालाक्षि नलसेतुः सुदुष्करः} %||6-111-10||

\twolineshloka
{पश्य सागरमक्षोभ्यं वैदेहि वरुणालयम्}
{अपारमभिगर्जन्तं शङ्खशुक्तिनिषेवितम्} %||6-111-11||

\twolineshloka
{हिरण्यनाभं शैलेन्द्रं काञ्चनं पश्य मैथिलि}
{विश्रमार्थं हनुमतो भित्त्वा सागरमुत्थितम्} %||6-111-12||

{अत्र राक्षसराजोऽयमाजगाम विभीषणः॥\devanumber{13}॥} %||6-111-13|| (Check)

\addtocounter{shlokacount}{1}

\twolineshloka
{एषा सा दृश्यते सीते किष्किन्धा चित्रकानना}
{सुग्रीवस्य पुरी रम्या यत्र वाली मया हतः} %||6-111-14||

\twolineshloka
{दृश्यतेऽसौ महान्सीते सविद्युदिव तोयदः}
{ऋश्यमूको गिरिश्रेष्ठः काञ्चनैर्धातुभिर्वृतः} %||6-111-15||

\twolineshloka
{अत्राहं वानरेन्द्रेण सुग्रीवेण समागतः}
{समयश्च कृतः सीते वधार्थं वालिनो मया} %||6-111-16||

\twolineshloka
{एषा सा दृश्यते पम्पा नलिनी चित्रकानना}
{त्वया विहीनो यत्राहं विललाप सुदुःखितः} %||6-111-17||

\twolineshloka
{अस्यास्तीरे मया दृष्टा शबरी धर्मचारिणी}
{अत्र योजनबाहुश्च कबन्धो निहतो मया} %||6-111-18||

\threelineshloka
{दृश्यतेऽसौ जनस्थाने सीते श्रीमान्वनस्पतिः}
{यत्र युद्धं महद्वृत्तं तवहेतोर्विलासिनि}
{रावणस्य नृशंसस्य जटायोश्च महात्मनः} %||6-111-19||

\twolineshloka
{खरश्च निहतश्सङ्ख्ये दूषणश्च निपातितः}
{त्रिशिराश्च महावीर्यो मया बाणैरजिह्मगैः} %||6-111-20||

\twolineshloka
{पर्णशाला तथा चित्रा दृश्यते शुभदर्शना}
{यत्र त्वं राक्षसेन्द्रेण रावणेन हृता बलात्} %||6-111-21||

\twolineshloka
{एषा गोदावरी रम्या प्रसन्नसलिला शिवा}
{अगस्त्यस्याश्रमो ह्येष दृश्यते पश्य मैथिलि} %||6-111-22||

\twolineshloka
{वैदेहि दृश्यते चात्र शरभङ्गाश्रमो महान्}
{उपयातः सहस्राक्षो यत्र शक्रः पुरन्दरः} %||6-111-23||

\threelineshloka
{एते ते तापसावासा दृश्यन्ते तनुमध्यमे}
{अत्रिः कुलपतिर्यत्र सूर्यवैश्वानरप्रभः}
{अत्र सीते त्वया दृष्टा तापसी धर्मचारिणी} %||6-111-24||

{अस्मिन्देशे महाकायो विराधो निहतो मया॥\devanumber{25}॥} %||6-111-25|| (Check)

\addtocounter{shlokacount}{1}

\twolineshloka
{असौ सुतनुशैलेन्द्रश्चित्रकूटः प्रकाशते}
{यत्र मां कैकयीपुत्रः प्रसादयितुमागतः} %||6-111-26||

\twolineshloka
{एषा सा यमुना दूराद्दृश्यते चित्रकानना}
{भरद्वाजाश्रमो यत्र श्रीमानेष प्रकाशते} %||6-111-27||

\twolineshloka
{एषा त्रिपथगा गङ्गा दृश्यते वरवर्णिनि}
{शृङ्गवेरपुरं चैतद्गुहो यत्र समागतः} %||6-111-28||

\twolineshloka
{एषा सा दृश्यतेऽयोध्या राजधानी पितुर्मम}
{अयोध्यां कुरु वैदेहि प्रणामं पुनरागता} %||6-111-29||

\twolineshloka
{ततस्ते वानराः सर्वे राक्षसश्च विभीषणः}
{उत्पत्योत्पत्य ददृशुस्तां पुरीं शुभदर्शनाम्} %||6-111-30||

\fourlineindentedshloka
{ततस्तु तां पाण्डुरहर्म्यमालिनीम्}
{विशालकक्ष्यां गजवाजिसङ्कुलाम्}
{पुरीमयोध्यां ददृशुः प्लवङ्गमाः}
{पुरीं महेन्द्रस्य यथामरावतीम्} %||6-112-31||


{॥इत्यार्षे श्रीमद्रामायणे वाल्मीकीये आदिकाव्ये युद्धकाण्डे एकादशाधिकशततमः सर्गः॥}

\sect{द्वादशाधिकशततमः सर्गः}

\twolineshloka
{पूर्णे चतुर्दशे वर्षे पञ्चभ्यां लक्ष्मणाग्रजः}
{भरद्वाजाश्रमं प्राप्य ववन्दे नियतो मुनिम्} %||6-112-1||

\threelineshloka
{सोऽपृच्छदभिवाद्यैनं भरद्वाजं तपोधनम्}
{शृणोषि कचिद्भगवन्सुभिक्षानामयं पुरे}
{कच्चिच्च युक्तो भरतो जीवन्त्यपि च मातरः} %||6-112-2||

\twolineshloka
{एवमुक्तस्तु रामेण भरद्वाजो महामुनिः}
{प्रत्युवाच रघुश्रेष्ठं स्मितपूर्वं प्रहृष्टवत्} %||6-112-3||

\twolineshloka
{पङ्कदिग्धस्तु भरतो जटिलस्त्वां प्रतीक्षते}
{पादुके ते पुरस्कृत्य सर्वं च कुशलं गृहे} %||6-112-4||

\twolineshloka
{त्वां पुरा चीरवसनं प्रविशन्तं महावनम्}
{स्त्रीतृतीयं च्युतं राज्याद्धर्मकामं च केवलम्} %||6-112-5||

\twolineshloka
{पदातिं त्यक्तसर्वस्वं पितुर्वचनकारिणम्}
{स्वर्गभोगैः परित्यक्तं स्वर्गच्युतमिवामरम्} %||6-112-6||

\twolineshloka
{दृष्ट्वा तु करुणा पूर्वं ममासीत्समितिंजय}
{कैकेयीवचने युक्तं वन्यमूलफलाशनम्} %||6-112-7||

\twolineshloka
{साम्प्रतं सुसमृद्धार्थं समित्रगणबान्धवम्}
{समीक्ष्य विजितारिं त्वां मम प्रीतिरनुत्तमा} %||6-112-8||

\twolineshloka
{सर्वं च सुखदुःखं ते विदितं मम राघव}
{यत्त्वया विपुलं प्राप्तं जनस्थानवधादिकम्} %||6-112-9||

\twolineshloka
{ब्राह्मणार्थे नियुक्तस्य रक्षतः सर्वतापसान्}
{मारीचदर्शनं चैव सीतोन्मथनमेव च} %||6-112-10||

\twolineshloka
{कबन्धदर्शनं चैव पम्पाभिगमनं तथा}
{सुग्रीवेण च ते सख्यं यच्च वाली हतस्त्वया} %||6-112-11||

\threelineshloka
{मार्गणं चैव वैदेह्याः कर्म वातात्मजस्य च}
{विदितायां च वैदेह्यां नलसेतुर्यथा कृतः}
{यथा च दीपिता लङ्का प्रहृष्टैर्हरियूथपैः} %||6-112-12||

\twolineshloka
{सपुत्रबान्धवामात्यः सबलः सह वाहनः}
{यथा च निहतः सङ्ख्ये रावणो देवकण्टकः} %||6-112-13||

\twolineshloka
{समागमश्च त्रिदशैर्यथादत्तश्च ते वरः}
{सर्वं ममैतद्विदितं तपसा धर्मवत्सल} %||6-112-14||

\twolineshloka
{अहमप्यत्र ते दद्मि वरं शस्त्रभृतां वर}
{अर्घ्यं प्रतिगृहाणेदमयोध्यां श्वो गमिष्यसि} %||6-112-15||

\twolineshloka
{तस्य तच्छिरसा वाक्यं प्रतिगृह्य नृपात्मजः}
{बाढमित्येव संहृष्टः श्रीमान्वरमयाचत} %||6-112-16||

\twolineshloka
{अकालफलिनो वृक्षाः सर्वे चापि मधुस्रवाः}
{भवन्तु मार्गे भगवन्नयोध्यां प्रति गच्छतः} %||6-112-17||

\twolineshloka
{निष्फलाः फलिनश्चासन्विपुष्पाः पुष्पशालिनः}
{शुष्काः समग्रपत्रास्ते नगाश्चैव मधुस्रवाः} %||6-113-18||


{॥इत्यार्षे श्रीमद्रामायणे वाल्मीकीये आदिकाव्ये युद्धकाण्डे द्वादशाधिकशततमः सर्गः॥}

\sect{त्रयोदशाधिकशततमः सर्गः}

\twolineshloka
{अयोध्यां तु समालोक्य चिन्तयामास राघवः}
{चिन्तयित्वा ततो दृष्टिं वानरेषु न्यपातयत्} %||6-113-1||

\twolineshloka
{प्रियकामः प्रियं रामस्ततस्त्वरितविक्रमम्}
{उवाच धीमांस्तेजस्वी हनूमन्तं प्लवङ्गमम्} %||6-113-2||

\twolineshloka
{अयोध्यां त्वरितो गच्छ क्षिप्रं त्वं प्लवगोत्तम}
{जानीहि कच्चित्कुशली जनो नृपतिमन्दिरे} %||6-113-3||

\twolineshloka
{शृङ्गवेरपुरं प्राप्य गुहं गहनगोचरम्}
{निषादाधिपतिं ब्रूहि कुशलं वचनान्मम} %||6-113-4||

\twolineshloka
{श्रुत्वा तु मां कुशलिनमरोगं विगतज्वरम्}
{भविष्यति गुहः प्रीतः स ममात्मसमः सखा} %||6-113-5||

\twolineshloka
{अयोध्यायाश्च ते मार्गं प्रवृत्तिं भरतस्य च}
{निवेदयिष्यति प्रीतो निषादाधिपतिर्गुहः} %||6-113-6||

\twolineshloka
{भरतस्तु त्वया वाच्यः कुशलं वचनान्मम}
{सिद्धार्थं शंस मां तस्मै सभार्यं सहलक्ष्मणम्} %||6-113-7||

\twolineshloka
{हरणं चापि वैदेह्या रावणेन बलीयसा}
{सुग्रीवेण च संवादं वालिनश्च वधं रणे} %||6-113-8||

\twolineshloka
{मैथिल्यन्वेषणं चैव यथा चाधिगता त्वया}
{लङ्घयित्वा महातोयमापगापतिमव्ययम्} %||6-113-9||

\twolineshloka
{उपयानं समुद्रस्य सागरस्य च दर्शनम्}
{यथा च कारितः सेतू रावणश्च यथा हतः} %||6-113-10||

\twolineshloka
{वरदानं महेन्द्रेण ब्रह्मणा वरुणेन च}
{महादेवप्रसादाच्च पित्रा मम समागमम्} %||6-113-11||

\twolineshloka
{जित्वा शत्रुगणान्रामः प्राप्य चानुत्तमं यशः}
{उपयाति समृद्धार्थः सह मित्रैर्महाबलः} %||6-113-12||

\twolineshloka
{एतच्छ्रुत्वा यमाकारं भजते भरतस्ततः}
{स च ते वेदितव्यः स्यात्सर्वं यच्चापि मां प्रति} %||6-113-13||

\twolineshloka
{ज्ञेयाः सर्वे च वृत्तान्ता भरतस्येङ्गितानि च}
{तत्त्वेन मुखवर्णेन दृष्ट्या व्याभाषणेन च} %||6-113-14||

\twolineshloka
{सर्वकामसमृद्धं हि हस्त्यश्वरथसङ्कुलम्}
{पितृपैतामहं राज्यं कस्य नावर्तयेन्मनः} %||6-113-15||

\twolineshloka
{सङ्गत्या भरतः श्रीमान्राज्येनार्थी स्वयं भवेत्}
{प्रशास्तु वसुधां सर्वामखिलां रघुनन्दनः} %||6-113-16||

\twolineshloka
{तस्य बुद्धिं च विज्ञाय व्यवसायं च वानर}
{यावन्न दूरं याताः स्मः क्षिप्रमागन्तुमर्हसि} %||6-113-17||

\twolineshloka
{इति प्रतिसमादिष्टो हनूमान्मारुतात्मजः}
{मानुषं धारयन्रूपमयोध्यां त्वरितो ययौ} %||6-113-18||

\twolineshloka
{लङ्घयित्वा पितृपथं भुजगेन्द्रालयं शुभम्}
{गङ्गायमुनयोर्भीमं संनिपातमतीत्य च} %||6-113-19||

\twolineshloka
{शृङ्गवेरपुरं प्राप्य गुहमासाद्य वीर्यवान्}
{स वाचा शुभया हृष्टो हनूमानिदमब्रवीत्} %||6-113-20||

\twolineshloka
{सखा तु तव काकुत्स्थो रामः सत्यपराक्रमः}
{ससीतः सह सौमित्रिः स त्वां कुशलमब्रवीत्} %||6-113-21||

\twolineshloka
{पञ्चमीमद्य रजनीमुषित्वा वचनान्मुनेः}
{भरद्वाजाभ्यनुज्ञातं द्रक्ष्यस्यद्यैव राघवम्} %||6-113-22||

\twolineshloka
{एवमुक्त्वा महातेजाः सम्प्रहृष्टतनूरुहः}
{उत्पपात महावेगो वेगवानविचारयन्} %||6-113-23||

\twolineshloka
{सोऽपश्यद्रामतीर्थं च नदीं वालुकिनीं तथा}
{गोमतीं तां च सोऽपश्यद्भीमं सालवनं तथा} %||6-113-24||

\twolineshloka
{स गत्वा दूरमध्वानं त्वरितः कपिकुञ्जरः}
{आससाद द्रुमान्फुल्लान्नन्दिग्रामसमीपजान्} %||6-113-25||

\twolineshloka
{क्रोशमात्रे त्वयोध्यायाश्चीरकृष्णाजिनाम्बरम्}
{ददर्श भरतं दीनं कृशमाश्रमवासिनम्} %||6-113-26||

\twolineshloka
{जटिलं मलदिग्धाङ्गं भ्रातृव्यसनकर्शितम्}
{फलमूलाशिनं दान्तं तापसं धर्मचारिणम्} %||6-113-27||

\twolineshloka
{समुन्नतजटाभारं वल्कलाजिनवाससम्}
{नियतं भावितात्मानं ब्रह्मर्षिसमतेजसम्} %||6-113-28||

\twolineshloka
{पादुके ते पुरस्कृत्य शासन्तं वै वसुन्धराम्}
{चतुर्वर्ण्यस्य लोकस्य त्रातारं सर्वतो भयात्} %||6-113-29||

\twolineshloka
{उपस्थितममात्यैश्च शुचिभिश्च पुरोहितैः}
{बलमुख्यैश्च युक्तैश्च काषायाम्बरधारिभिः} %||6-113-30||

\twolineshloka
{न हि ते राजपुत्रं तं चीरकृष्णाजिनाम्बरम्}
{परिमोक्तुं व्यवस्यन्ति पौरा वै धर्मवत्सलाः} %||6-113-31||

\twolineshloka
{तं धर्ममिव धर्मज्ञं देववन्तमिवापरम्}
{उवाच प्राञ्जलिर्वाक्यं हनूमान्मारुतात्मजः} %||6-113-32||

\twolineshloka
{वसन्तं दण्डकारण्ये यं त्वं चीरजटाधरम्}
{अनुशोचसि काकुत्स्थं स त्वा कुशलमब्रवीत्} %||6-113-33||

\twolineshloka
{प्रियमाख्यामि ते देव शोकं त्यक्ष्यसि दारुणम्}
{अस्मिन्मुहूर्ते भ्रात्रा त्वं रामेण सह सङ्गतः} %||6-113-34||

\twolineshloka
{निहत्य रावणं रामः प्रतिलभ्य च मैथिलीम्}
{उपयाति समृद्धार्थः सह मित्रैर्महाबलैः} %||6-113-35||

\twolineshloka
{लक्ष्मणश्च महातेजा वैदेही च यशस्विनी}
{सीता समग्रा रामेण महेन्द्रेण शची यथा} %||6-113-36||

\twolineshloka
{एवमुक्तो हनुमता भरतः कैकयीसुतः}
{पपात सहसा हृष्टो हर्षान्मोहं जगाम ह} %||6-113-37||

\twolineshloka
{ततो मुहूर्तादुत्थाय प्रत्याश्वस्य च राघवः}
{हनूमन्तमुवाचेदं भरतः प्रियवादिनम्} %||6-113-38||

\twolineshloka
{अशोकजैः प्रीतिमयैः कपिमालिङ्ग्य सम्भ्रमात्}
{सिषेच भरतः श्रीमान्विपुलैरश्रुबिन्दुभिः} %||6-113-39||

\twolineshloka
{देवो वा मानुषो वा त्वमनुक्रोशादिहागतः}
{प्रियाख्यानस्य ते सौम्य ददामि ब्रुवतः प्रियम्} %||6-113-40||

\twolineshloka
{गवां शतसहस्रं च ग्रामाणां च शतं परम्}
{सकुण्डलाः शुभाचारा भार्याः कन्याश्च षोडश} %||6-113-41||

\twolineshloka
{हेमवर्णाः सुनासोरूः शशिसौम्याननाः स्त्रियः}
{सर्वाभरणसम्पन्ना सम्पन्नाः कुलजातिभिः} %||6-113-42||

\fourlineindentedshloka
{निशम्य रामागमनं नृपात्मजः}
{कपिप्रवीरस्य तदाद्भुतोपमम्}
{प्रहर्षितो रामदिदृक्षयाभवत्}
{पुनश्च हर्षादिदमब्रवीद्वचः} %||6-114-43||


{॥इत्यार्षे श्रीमद्रामायणे वाल्मीकीये आदिकाव्ये युद्धकाण्डे त्रयोदशाधिकशततमः सर्गः॥}

\sect{चतुर्दशाधिकशततमः सर्गः}

\twolineshloka
{बहूनि नाम वर्षाणि गतस्य सुमहद्वनम्}
{शृणोम्यहं प्रीतिकरं मम नाथस्य कीर्तनम्} %||6-114-1||

\twolineshloka
{कल्याणी बत गाथेयं लौकिकी प्रतिभाति मे}
{एति जीवन्तमानन्दो नरं वर्षशतादपि} %||6-114-2||

\twolineshloka
{राघवस्य हरीणां च कथमासीत्समागमः}
{कस्मिन्देशे किमाश्रित्य तत्त्वमाख्याहि पृच्छतः} %||6-114-3||

\twolineshloka
{स पृष्टो राजपुत्रेण बृस्यां समुपवेशितः}
{आचचक्षे ततः सर्वं रामस्य चरितं वने} %||6-114-4||

\twolineshloka
{यथा प्रव्रजितो रामो मातुर्दत्ते वरे तव}
{यथा च पुत्रशोकेन राजा दशरथो मृतः} %||6-114-5||

\twolineshloka
{यथा दूतैस्त्वमानीतस्तूर्णं राजगृहात्प्रभो}
{त्वयायोध्यां प्रविष्टेन यथा राज्यं न चेप्सितम्} %||6-114-6||

\twolineshloka
{चित्रकूटं गिरिं गत्वा राज्येनामित्रकर्शनः}
{निमन्त्रितस्त्वया भ्राता धर्ममाचरिता सताम्} %||6-114-7||

\twolineshloka
{स्थितेन राज्ञो वचने यथा राज्यं विसर्जितम्}
{आर्यस्य पादुके गृह्य यथासि पुनरागतः} %||6-114-8||

\twolineshloka
{सर्वमेतन्महाबाहो यथावद्विदितं तव}
{त्वयि प्रतिप्रयाते तु यद्वृत्तं तन्निबोध मे} %||6-114-9||

\twolineshloka
{अपयाते त्वयि तदा समुद्भ्रान्तमृगद्विजम्}
{प्रविवेशाथ विजनं सुमहद्दण्डकावनम्} %||6-114-10||

\twolineshloka
{तेषां पुरस्ताद्बलवान्गच्छतां गहने वने}
{विनदन्सुमहानादं विराधः प्रत्यदृश्यत} %||6-114-11||

\twolineshloka
{तमुत्क्षिप्य महानादमूर्ध्वबाहुमधोमुखम्}
{निखाते प्रक्षिपन्ति स्म नदन्तमिव कुञ्जरम्} %||6-114-12||

\twolineshloka
{तत्कृत्वा दुष्करं कर्म भ्रातरौ रामलक्ष्मणौ}
{सायाह्ने शरभङ्गस्य रम्यमाश्रममीयतुः} %||6-114-13||

\twolineshloka
{शरभङ्गे दिवं प्राप्ते रामः सत्यपराक्रमः}
{अभिवाद्य मुनीन्सर्वाञ्जनस्थानमुपागमत्} %||6-114-14||

\twolineshloka
{चतुर्दशसहस्राणि रक्षसां भीमकर्मणाम्}
{हतानि वसता तत्र राघवेण महात्मना} %||6-114-15||

\twolineshloka
{ततः पश्चाच्छूर्पणखा रामपार्श्वमुपागता}
{ततो रामेण सन्दिष्टो लक्ष्मणः सहसोत्थितः} %||6-114-16||

\twolineshloka
{प्रगृह्य खड्गं चिच्छेद कर्णनासे महाबलः}
{ततस्तेनार्दिता बाला रावणं समुपागता} %||6-114-17||

\twolineshloka
{रावणानुचरो घोरो मारीचो नाम राक्षसः}
{लोभयामास वैदेहीं भूत्वा रत्नमयो मृगः} %||6-114-18||

\twolineshloka
{सा राममब्रवीद्दृष्ट्वा वैदेही गृह्यतामिति}
{अहो मनोहरः कान्त आश्रमे नो भविष्यति} %||6-114-19||

\twolineshloka
{ततो रामो धनुष्पाणिर्धावन्तमनुधावति}
{स तं जघान धावन्तं शरेणानतपर्वणा} %||6-114-20||

\threelineshloka
{अथ सौम्या दशग्रीवो मृगं याते तु राघवे}
{लक्ष्मणे चापि निष्क्रान्ते प्रविवेशाश्रमं तदा}
{जग्राह तरसा सीतां ग्रहः खे रोहिणीमिव} %||6-114-21||

\twolineshloka
{त्रातुकामं ततो युद्धे हत्वा गृध्रं जटायुषम्}
{प्रगृह्य सीतां सहसा जगामाशु स रावणः} %||6-114-22||

\threelineshloka
{ततस्त्वद्भुतसङ्काशाः स्थिताः पर्वतमूर्धनि}
{सीतां गृहीत्वा गच्छन्तं वानराः पर्वतोपमाः}
{ददृशुर्विस्मितास्तत्र रावणं राक्षसाधिपम्} %||6-114-23||

{प्रविवेर्श तदा लङ्कां रावणो लोकरावणः॥\devanumber{24}॥} %||6-114-24|| (Check)

\addtocounter{shlokacount}{1}

\twolineshloka
{तां सुवर्णपरिक्रान्ते शुभे महति वेश्मनि}
{प्रवेश्य मैथिलीं वाक्यैः सान्त्वयामास रावणः} %||6-114-25||

{निवर्तमानः काकुत्स्थो दृष्ट्वा गृध्रं प्रविव्यथे॥\devanumber{26}॥} %||6-114-26|| (Check)

\addtocounter{shlokacount}{1}

\threelineshloka
{गृध्रं हतं तदा दग्ध्वा रामः प्रियसखं पितुः}
{गोदावरीमनुचरन्वनोद्देशांश्च पुष्पितान्}
{आसेदतुर्महारण्ये कबन्धं नाम राक्षसम्} %||6-114-27||

\twolineshloka
{ततः कबन्धवचनाद्रामः सत्यपराक्रमः}
{ऋश्यमूकं गिरिं गत्वा सुग्रीवेण समागतः} %||6-114-28||

\twolineshloka
{तयोः समागमः पूर्वं प्रीत्या हार्दो व्यजायत}
{इतरेतर संवादात्प्रगाढः प्रणयस्तयोः} %||6-114-29||

\twolineshloka
{रामः स्वबाहुवीर्येण स्वराज्यं प्रत्यपादयत्}
{वालिनं समरे हत्वा महाकायं महाबलम्} %||6-114-30||

\twolineshloka
{सुग्रीवः स्थापितो राज्ये सहितः सर्ववानरैः}
{रामाय प्रतिजानीते राजपुत्र्यास्तु मार्गणम्} %||6-114-31||

\twolineshloka
{आदिष्टा वानरेन्द्रेण सुग्रीवेण महात्मना}
{दशकोट्यः प्लवङ्गानां सर्वाः प्रस्थापिता दिशः} %||6-114-32||

\twolineshloka
{तेषां नो विप्रनष्टानां विन्ध्ये पर्वतसत्तमे}
{भृशं शोकाभितप्तानां महान्कालोऽत्यवर्तत} %||6-114-33||

\twolineshloka
{भ्राता तु गृध्रराजस्य सम्पातिर्नाम वीर्यवान्}
{समाख्याति स्म वसतिं सीताया रावणालये} %||6-114-34||

\twolineshloka
{सोऽहं दुःखपरीतानां दुःखं तज्ज्ञातिनां नुदन्}
{आत्मवीर्यं समास्थाय योजनानां शतं प्लुतः} %||6-114-35||

\twolineshloka
{तत्राहमेकामद्राक्षमशोकवनिकां गताम्}
{कौशेयवस्त्रां मलिनां निरानन्दां दृढव्रताम्} %||6-114-36||

\twolineshloka
{तया समेत्य विधिवत्पृष्ट्वा सर्वमनिन्दिताम्}
{अभिज्ञानं मणिं लब्ध्वा चरितार्थोऽहमागतः} %||6-114-37||

\twolineshloka
{मया च पुनरागम्य रामस्याक्लिष्टकर्मणः}
{अभिज्ञानं मया दत्तमर्चिष्मान्स महामणिः} %||6-114-38||

\twolineshloka
{श्रुत्वा तां मैथिलीं हृष्टस्त्वाशशंसे स जीवितम्}
{जीवितान्तमनुप्राप्तः पीत्वामृतमिवातुरः} %||6-114-39||

\twolineshloka
{उद्योजयिष्यन्नुद्योगं दध्रे लङ्कावधे मनः}
{जिघांसुरिव लोकांस्ते सर्वाँल्लोकान्विभावसुः} %||6-114-40||

\twolineshloka
{ततः समुद्रमासाद्य नलं सेतुमकारयत्}
{अतरत्कपिवीराणां वाहिनी तेन सेतुना} %||6-114-41||

\twolineshloka
{प्रहस्तमवधीन्नीलः कुम्भकर्णं तु राघवः}
{लक्ष्मणो रावणसुतं स्वयं रामस्तु रावणम्} %||6-114-42||

\twolineshloka
{स शक्रेण समागम्य यमेन वरुणेन च}
{सुरर्षिभिश्च काकुत्स्थो वराँल्लेभे परन्तपः} %||6-114-43||

\twolineshloka
{स तु दत्तवरः प्रीत्या वानरैश्च समागतः}
{पुष्पकेण विमानेन किष्किन्धामभ्युपागमत्} %||6-114-44||

\twolineshloka
{तं गङ्गां पुनरासाद्य वसन्तं मुनिसंनिधौ}
{अविघ्नं पुष्ययोगेन श्वो रामं द्रष्टुमर्हसि} %||6-114-45||

\fourlineindentedshloka
{ततः स सत्यं हनुमद्वचो महत्}
{निशम्य हृष्टो भरतः कृताञ्जलिः}
{उवाच वाणीं मनसः प्रहर्षिणी}
{चिरस्य पूर्णः खलु मे मनोरथः} %||6-115-46||


{॥इत्यार्षे श्रीमद्रामायणे वाल्मीकीये आदिकाव्ये युद्धकाण्डे चतुर्दशाधिकशततमः सर्गः॥}

\sect{पञ्चदशाधिकशततमः सर्गः}

\twolineshloka
{श्रुत्वा तु परमानन्दं भरतः सत्यविक्रमः}
{हृष्टमाज्ञापयामास शत्रुघ्नं परवीरहा} %||6-115-1||

\twolineshloka
{दैवतानि च सर्वाणि चैत्यानि नगरस्य च}
{सुगन्धमाल्यैर्वादित्रैरर्चन्तु शुचयो नराः} %||6-115-2||

\twolineshloka
{राजदारास्तथामात्याः सैन्याः सेनागणाङ्गनाः}
{अभिनिर्यान्तु रामस्य द्रष्टुं शशिनिभं मुखम्} %||6-115-3||

\twolineshloka
{भरतस्य वचः श्रुत्वा शत्रुघ्नः परवीरहा}
{विष्टीरनेकसाहस्रीश्चोदयामास वीर्यवान्} %||6-115-4||

\twolineshloka
{समीकुरुत निम्नानि विषमाणि समानि च}
{स्थानानि च निरस्यन्तां नन्दिग्रामादितः परम्} %||6-115-5||

\twolineshloka
{सिञ्चन्तु पृथिवीं कृत्स्नां हिमशीतेन वारिणा}
{ततोऽभ्यवकिरंस्त्वन्ये लाजैः पुष्पैश्च सर्वतः} %||6-115-6||

\twolineshloka
{समुच्छ्रितपताकास्तु रथ्याः पुरवरोत्तमे}
{शोभयन्तु च वेश्मानि सूर्यस्योदयनं प्रति} %||6-115-7||

\twolineshloka
{स्रग्दाममुक्तपुष्पैश्च सुगन्धैः पञ्चवर्णकैः}
{राजमार्गमसम्बाधं किरन्तु शतशो नराः} %||6-115-8||

\threelineshloka
{मत्तैर्नागसहस्रैश्च शातकुम्भविभूषितः}
{अपरे हेमकक्ष्याभिः सगजाभिः करेणुभिः}
{निर्ययुस्त्वरया युक्ता रथैश्च सुमहारथाः} %||6-115-9||

\twolineshloka
{ततो यानान्युपारूढाः सर्वा दशरथस्त्रियः}
{कौसल्यां प्रमुखे कृत्वा सुमित्रां चापि निर्ययुः} %||6-115-10||

\twolineshloka
{अश्वानां खुरशब्देन रथनेमिस्वनेन च}
{शङ्खदुन्दुभिनादेन सञ्चचालेव मेदिनी} %||6-115-11||

\twolineshloka
{कृत्स्नं च नगरं तत्तु नन्दिग्राममुपागमत्}
{द्विजातिमुख्यैर्धर्मात्मा श्रेणीमुख्यैः सनैगमैः} %||6-115-12||

\twolineshloka
{माल्यमोदक हस्तैश्च मन्त्रिभिर्भरतो वृतः}
{शङ्खभेरीनिनादैश्च बन्दिभिश्चाभिवन्दितः} %||6-115-13||

\twolineshloka
{आर्यपादौ गृहीत्वा तु शिरसा धर्मकोविदः}
{पाण्डुरं छत्रमादाय शुक्लमाल्योपशोभितम्} %||6-115-14||

\twolineshloka
{शुक्ले च वालव्यजने राजार्हे हेमभूषिते}
{उपवासकृशो दीनश्चीरकृष्णाजिनाम्बरः} %||6-115-15||

\twolineshloka
{भ्रातुरागमनं श्रुत्वा तत्पूर्वं हर्षमागतः}
{प्रत्युद्ययौ तदा रामं महात्मा सचिवैः सह} %||6-115-16||

\threelineshloka
{समीक्ष्य भरतो वाक्यमुवाच पवनात्मजम्}
{कच्चिन्न खलु कापेयी सेव्यते चलचित्तता}
{न हि पश्यामि काकुत्स्थं राममार्यं परन्तपम्} %||6-115-17||

\twolineshloka
{अथैवमुक्ते वचने हनूमानिदमब्रवीत्}
{अर्थं विज्ञापयन्नेव भरतं सत्यविक्रमम्} %||6-115-18||

\twolineshloka
{सदा फलान्कुसुमितान्वृक्षान्प्राप्य मधुस्रवान्}
{भरद्वाजप्रसादेन मत्तभ्रमरनादितान्} %||6-115-19||

\twolineshloka
{तस्य चैष वरो दत्तो वासवेन परन्तप}
{ससैन्यस्य तदातिथ्यं कृतं सर्वगुणान्वितम्} %||6-115-20||

\twolineshloka
{निस्वनः श्रूयते भीमः प्रहृष्टानां वनौकसाम्}
{मन्ये वानरसेना सा नदीं तरति गोमतीम्} %||6-115-21||

\twolineshloka
{रजोवर्षं समुद्भूतं पश्य वालुकिनीं प्रति}
{मन्ये सालवनं रम्यं लोलयन्ति प्लवङ्गमाः} %||6-115-22||

\twolineshloka
{तदेतद्दृश्यते दूराद्विमलं चन्द्रसंनिभम्}
{विमानं पुष्पकं दिव्यं मनसा ब्रह्मनिर्मितम्} %||6-115-23||

\twolineshloka
{रावणं बान्धवैः सार्धं हत्वा लब्धं महात्मना}
{धनदस्य प्रसादेन दिव्यमेतन्मनोजवम्} %||6-115-24||

\twolineshloka
{एतस्मिन्भ्रातरौ वीरौ वैदेह्या सह राघवौ}
{सुग्रीवश्च महातेजा राक्षसेन्द्रो विभीषणः} %||6-115-25||

\twolineshloka
{ततो हर्षसमुद्भूतो निस्वनो दिवमस्पृशत्}
{स्त्रीबालयुववृद्धानां रामोऽयमिति कीर्तितः} %||6-115-26||

\twolineshloka
{रथकुञ्जरवाजिभ्यस्तेऽवतीर्य महीं गताः}
{ददृशुस्तं विमानस्थं नराः सोममिवाम्बरे} %||6-115-27||

\twolineshloka
{प्राञ्जलिर्भरतो भूत्वा प्रहृष्टो राघवोन्मुखः}
{स्वागतेन यथार्थेन ततो राममपूजयत्} %||6-115-28||

\twolineshloka
{मनसा ब्रह्मणा सृष्टे विमाने लक्ष्मणाग्रजः}
{रराज पृथुदीर्घाक्षो वज्रपाणिरिवापरः} %||6-115-29||

\twolineshloka
{ततो विमानाग्रगतं भरतो भ्रातरं तदा}
{ववन्दे प्रणतो रामं मेरुस्थमिव भास्करम्} %||6-115-30||

\twolineshloka
{आरोपितो विमानं तद्भरतः सत्यविक्रमः}
{राममासाद्य मुदितः पुनरेवाभ्यवादयत्} %||6-115-31||

\twolineshloka
{तं समुत्थाप्य काकुत्स्थश्चिरस्याक्षिपथं गतम्}
{अङ्के भरतमारोप्य मुदितः परिषष्वजे} %||6-115-32||

\twolineshloka
{ततो लक्ष्मणमासाद्य वैदेहीं च परन्तपः}
{अभ्यवादयत प्रीतो भरतो नाम चाब्रवीत्} %||6-115-33||

\twolineshloka
{सुग्रीवं कैकयी पुत्रो जाम्बवन्तं तथाङ्गदम्}
{मैन्दं च द्विविदं नीलमृषभं चैव सस्वजे} %||6-115-34||

\twolineshloka
{ते कृत्वा मानुषं रूपं वानराः कामरूपिणः}
{कुशलं पर्यपृष्हन्त प्रहृष्टा भरतं तदा} %||6-115-35||

\twolineshloka
{विभीषणं च भरतः सान्त्वयन्वाक्यमब्रवीत्}
{दिष्ट्या त्वया सहायेन कृतं कर्म सुदुष्करम्} %||6-115-36||

\twolineshloka
{शत्रुघ्नश्च तदा राममभिवाद्य सलक्ष्मणम्}
{सीतायाश्चरणौ पश्चाद्ववन्दे विनयान्वितः} %||6-115-37||

\twolineshloka
{रामो मातरमासाद्य विषण्णं शोककर्शिताम्}
{जग्राह प्रणतः पादौ मनो मातुः प्रसादयन्} %||6-115-38||

\twolineshloka
{अभिवाद्य सुमित्रां च कैकेयीं च यशस्विनीम्}
{स मातॄश्च तदा सर्वाः पुरोहितमुपागमत्} %||6-115-39||

\twolineshloka
{स्वागतं ते महाबाहो कौसल्यानन्दवर्धन}
{इति प्राञ्जलयः सर्वे नागरा राममब्रुवन्} %||6-115-40||

\twolineshloka
{तन्यञ्जलिसहस्राणि प्रगृहीतानि नागरैः}
{आकोशानीव पद्मानि ददर्श भरताग्रजः} %||6-115-41||

\twolineshloka
{पादुके ते तु रामस्य गृहीत्वा भरतः स्वयम्}
{चरणाभ्यां नरेन्द्रस्य योजयामास धर्मवित्} %||6-115-42||

\twolineshloka
{अब्रवीच्च तदा रामं भरतः स कृताञ्जलिः}
{एतत्ते रक्षितं राजन्राज्यं निर्यातितं मया} %||6-115-43||

\twolineshloka
{अद्य जन्म कृतार्थं मे संवृत्तश्च मनोरथः}
{यस्त्वां पश्यामि राजानमयोध्यां पुनरागतम्} %||6-115-44||

\twolineshloka
{अवेक्षतां भवान्कोशं कोष्ठागारं पुरं बलम्}
{भवतस्तेजसा सर्वं कृतं दशगुणं मया} %||6-115-45||

\twolineshloka
{तथा ब्रुवाणं भरतं दृष्ट्वा तं भ्रातृवत्सलम्}
{मुमुचुर्वानरा बाष्पं राक्षसश्च विभीषणः} %||6-115-46||

\twolineshloka
{ततः प्रहर्षाद्भरतमङ्कमारोप्य राघवः}
{ययौ तेन विमानेन ससैन्यो भरताश्रमम्} %||6-115-47||

\twolineshloka
{भरताश्रममासाद्य ससैन्यो राघवस्तदा}
{अवतीर्य विमानाग्रादवतस्थे महीतले} %||6-115-48||

\twolineshloka
{अब्रवीच्च तदा रामस्तद्विमानमनुत्तमम्}
{वह वैश्रवणं देवमनुजानामि गम्यताम्} %||6-115-49||

\twolineshloka
{ततो रामाभ्यनुज्ञातं तद्विमानमनुत्तमम्}
{उत्तरां दिशमुद्दिश्य जगाम धनदालयम्} %||6-115-50||

\fourlineindentedshloka
{पुरोहितस्यात्मसमस्य राघवो}
{बृहस्पतेः शक्र इवामराधीअफ्}
{निपीड्य पादौ पृथगासने शुभे}
{सहैव तेनोपविवेश वीर्यवान्} %||6-116-51||


{॥इत्यार्षे श्रीमद्रामायणे वाल्मीकीये आदिकाव्ये युद्धकाण्डे पञ्चदशाधिकशततमः सर्गः॥}

\sect{षोडशाधिकशततमः सर्गः}

\twolineshloka
{शिरस्यञ्जलिमादाय कैकेयीनन्दिवर्धनः}
{बभाषे भरतो ज्येष्ठं रामं सत्यपराक्रमम्} %||6-116-1||

\twolineshloka
{पूजिता मामिका माता दत्तं राज्यमिदं मम}
{तद्ददामि पुनस्तुभ्यं यथा त्वमददा मम} %||6-116-2||

\twolineshloka
{धुरमेकाकिना न्यस्तामृषभेण बलीयसा}
{किशोरवद्गुरुं भारं न वोढुमहमुत्सहे} %||6-116-3||

\twolineshloka
{वारिवेगेन महता भिन्नः सेतुरिव क्षरन्}
{दुर्बन्धनमिदं मन्ये राज्यच्छिद्रमसंवृतम्} %||6-116-4||

\twolineshloka
{गतिं खर इवाश्वस्य हंसस्येव च वायसः}
{नान्वेतुमुत्सहे देव तव मार्गमरिन्दम} %||6-116-5||

\twolineshloka
{यथा च रोपितो वृक्षो जातश्चान्तर्निवेशने}
{महांश्च सुदुरारोहो महास्कन्धः प्रशाखवान्} %||6-116-6||

\twolineshloka
{शीर्येत पुष्पितो भूत्वा न फलानि प्रदर्शयेत्}
{तस्य नानुभवेदर्थं यस्य हेतोः स रोप्यते} %||6-116-7||

\twolineshloka
{एषोपमा महाबाहो त्वमर्थं वेत्तुमर्हसि}
{यद्यस्मान्मनुजेन्द्र त्वं भक्तान्भृत्यान्न शाधि हि} %||6-116-8||

\twolineshloka
{जगदद्याभिषिक्तं त्वामनुपश्यतु सर्वतः}
{प्रतपन्तमिवादित्यं मध्याह्ने दीप्ततेजसम्} %||6-116-9||

\twolineshloka
{तूर्यसङ्घातनिर्घोषैः काञ्चीनूपुरनिस्वनैः}
{मधुरैर्गीतशब्दैश्च प्रतिबुध्यस्व शेष्व च} %||6-116-10||

\twolineshloka
{यावदावर्तते चक्रं यावती च वसुन्धरा}
{तावत्त्वमिह सर्वस्य स्वामित्वमभिवर्तय} %||6-116-11||

\twolineshloka
{भरतस्य वचः श्रुत्वा रामः परपुरंजयः}
{तथेति प्रतिजग्राह निषसादासने शुभे} %||6-116-12||

\twolineshloka
{ततः शत्रुघ्नवचनान्निपुणाः श्मश्रुवर्धकाः}
{सुखहस्ताः सुशीघ्राश्च राघवं पर्युपासत} %||6-116-13||

\twolineshloka
{पूर्वं तु भरते स्नाते लक्ष्मणे च महाबले}
{सुग्रीवे वानरेन्द्रे च राक्षसेन्द्रे विभीषणे} %||6-116-14||

\twolineshloka
{विशोधितजटः स्नातश्चित्रमाल्यानुलेपनः}
{महार्हवसनोपेतस्तस्थौ तत्र श्रिया ज्वलन्} %||6-116-15||

\twolineshloka
{प्रतिकर्म च रामस्य कारयामास वीर्यवान्}
{लक्ष्मणस्य च लक्ष्मीवानिक्ष्वाकुकुलवर्धनः} %||6-116-16||

\twolineshloka
{प्रतिकर्म च सीतायाः सर्वा दशरथस्त्रियः}
{आत्मनैव तदा चक्रुर्मनस्विन्यो मनोहरम्} %||6-116-17||

\twolineshloka
{ततो राघवपत्नीनां सर्वासामेव शोभनम्}
{चकार यत्नात्कौसल्या प्रहृष्टा पुत्रवत्सला} %||6-116-18||

\twolineshloka
{ततः शत्रुघ्नवचनात्सुमन्त्रो नाम सारथिः}
{योजयित्वाभिचक्राम रथं सर्वाङ्गशोभनम्} %||6-116-19||

\twolineshloka
{अर्कमण्डलसङ्काशं दिव्यं दृष्ट्वा रथं स्थितम्}
{आरुरोह महाबाहू रामः सत्यपराक्रमः} %||6-116-20||

\twolineshloka
{अयोध्यायां तु सचिवा राज्ञो दशरथस्य ये}
{पुरोहितं पुरस्कृत्य मन्त्रयामासुरर्थवत्} %||6-116-21||

\threelineshloka
{मन्त्रयन्रामवृद्ध्यर्थं वृत्त्यर्थं नगरस्य च}
{सर्वमेवाभिषेकार्थं जयार्हस्य महात्मनः}
{कर्तुमर्हथ रामस्य यद्यन्मङ्गलपूर्वकम्} %||6-116-22||

\twolineshloka
{इति ते मन्त्रिणः सर्वे सन्दिश्य तु पुरोहितम्}
{नगरान्निर्ययुस्तूर्णं रामदर्शनबुद्धयः} %||6-116-23||

\twolineshloka
{हरियुक्तं सहस्राक्षो रथमिन्द्र इवानघः}
{प्रययौ रथमास्थाय रामो नगरमुत्तमम्} %||6-116-24||

\twolineshloka
{जग्राह भरतो रश्मीञ्शत्रुघ्नश्छत्रमाददे}
{लक्ष्मणो व्यजनं तस्य मूर्ध्नि सम्पर्यवीजयत्} %||6-116-25||

\twolineshloka
{श्वेतं च वालव्यजनं सुग्रीवो वानरेश्वरः}
{अपरं चन्द्रसङ्काशं राक्षसेन्द्रो विभीषणः} %||6-116-26||

\twolineshloka
{ऋषिसङ्घैर्तदाकाशे देवैश्च समरुद्गणैः}
{स्तूयमानस्य रामस्य शुश्रुवे मधुरध्वनिः} %||6-116-27||

\twolineshloka
{ततः शत्रुंजयं नाम कुञ्जरं पर्वतोपमम्}
{आरुरोह महातेजाः सुग्रीवो वानरेश्वरः} %||6-116-28||

\twolineshloka
{नवनागसहस्राणि ययुरास्थाय वानराः}
{मानुषं विग्रहं कृत्वा सर्वाभरणभूषिताः} %||6-116-29||

\twolineshloka
{शङ्खशब्दप्रणादैश्च दुन्दुभीनां च निस्वनैः}
{प्रययू पुरुषव्याघ्रस्तां पुरीं हर्म्यमालिनीम्} %||6-116-30||

\twolineshloka
{ददृशुस्ते समायान्तं राघवं सपुरःसरम्}
{विराजमानं वपुषा रथेनातिरथं तदा} %||6-116-31||

\twolineshloka
{ते वर्धयित्वा काकुत्स्थं रामेण प्रतिनन्दिताः}
{अनुजग्मुर्महात्मानं भ्रातृभिः परिवारितम्} %||6-116-32||

\twolineshloka
{अमात्यैर्ब्राह्मणैश्चैव तथा प्रकृतिभिर्वृतः}
{श्रिया विरुरुचे रामो नक्षत्रैरिव चन्द्रमाः} %||6-116-33||

\twolineshloka
{स पुरोगामिभिस्तूर्यैस्तालस्वस्तिकपाणिभिः}
{प्रव्याहरद्भिर्मुदितैर्मङ्गलानि ययौ वृतः} %||6-116-34||

\twolineshloka
{अक्षतं जातरूपं च गावः कन्यास्तथा द्विजाः}
{नरा मोदकहस्ताश्च रामस्य पुरतो ययुः} %||6-116-35||

\threelineshloka
{सख्यं च रामः सुग्रीवे प्रभावं चानिलात्मजे}
{वानराणां च तत्कर्म व्याचचक्षेऽथ मन्त्रिणाम्}
{श्रुत्वा च विस्मयं जग्मुरयोध्यापुरवासिनः} %||6-116-36||

\twolineshloka
{द्युतिमानेतदाख्याय रामो वानरसंवृतः}
{हृष्टपुष्टजनाकीर्णामयोध्यां प्रविवेश ह} %||6-116-37||

\twolineshloka
{ततो ह्यभ्युच्छ्रयन्पौराः पताकास्ते गृहे गृहे}
{ऐक्ष्वाकाध्युषितं रम्यमाससाद पितुर्गृहम्} %||6-116-38||

\twolineshloka
{पितुर्भवनमासाद्य प्रविश्य च महात्मनः}
{कौसल्यां च सुमित्रां च कैकेयीं चाभ्यवादयत्} %||6-116-39||

\twolineshloka
{अथाब्रवीद्राजपुत्रो भरतं धर्मिणां वरम्}
{अथोपहितया वाचा मधुरं रघुनन्दनः} %||6-116-40||

\twolineshloka
{यच्च मद्भवनं श्रेष्ठं साशोकवनिकं महत्}
{मुक्तावैदूर्यसङ्कीर्णं सुग्रीवस्य निवेदय} %||6-116-41||

\twolineshloka
{तस्य तद्वचनं श्रुत्वा भरतः सत्यविक्रमः}
{पाणौ गृहीत्वा सुग्रीवं प्रविवेश तमालयम्} %||6-116-42||

\twolineshloka
{ततस्तैलप्रदीपांश्च पर्यङ्कास्तरणानि च}
{गृहीत्वा विविशुः क्षिप्रं शत्रुघ्नेन प्रचोदिताः} %||6-116-43||

\twolineshloka
{उवाच च महातेजाः सुग्रीवं राघवानुजः}
{अभिषेकाय रामस्य दूतानाज्ञापय प्रभो} %||6-116-44||

\twolineshloka
{सौवर्णान्वानरेन्द्राणां चतुर्णां चतुरो घटान्}
{ददौ क्षिप्रं स सुग्रीवः सर्वरत्नविभूषितान्} %||6-116-45||

\twolineshloka
{यथा प्रत्यूषसमये चतुर्णां सागराम्भसाम्}
{पूर्णैर्घटैः प्रतीक्षध्वं तथा कुरुत वानराः} %||6-116-46||

\twolineshloka
{एवमुक्ता महात्मानो वानरा वारणोपमाः}
{उत्पेतुर्गगनं शीघ्रं गरुडा इव शीघ्रगाः} %||6-116-47||

\threelineshloka
{जाम्बवांश्च हनूमांश्च वेगदर्शी च वानरः}
{ऋषभश्चैव कलशाञ्जलपूर्णानथानयन्}
{नदीशतानां पञ्चानां जले कुम्भैरुपाहरन्} %||6-116-48||

\twolineshloka
{पूर्वात्समुद्रात्कलशं जलपूर्णमथानयत्}
{सुषेणः सत्त्वसम्पन्नः सर्वरत्नविभूषितम्} %||6-116-49||

{ऋषभो दक्षिणात्तूर्णं समुद्राज्जलमाहरत्॥\devanumber{50}॥} %||6-116-50|| (Check)

\addtocounter{shlokacount}{1}

\twolineshloka
{रक्तचन्दनकर्पूरैः संवृतं काञ्चनं घटम्}
{गवयः पश्चिमात्तोयमाजहार महार्णवात्} %||6-116-51||

\twolineshloka
{रत्नकुम्भेन महता शीतं मारुतविक्रमः}
{उत्तराच्च जलं शीघ्रं गरुडानिलविक्रमः} %||6-116-52||

\twolineshloka
{अभिषेकाय रामस्य शत्रुघ्नः सचिवैः सह}
{पुरोहिताय श्रेष्ठाय सुहृद्भ्यश्च न्यवेदयत्} %||6-116-53||

\twolineshloka
{ततः स प्रयतो वृद्धो वसिष्ठो ब्राह्मणैः सह}
{रामं रत्नमयो पीठे सहसीतं न्यवेशयत्} %||6-116-54||

\twolineshloka
{वसिष्ठो वामदेवश्च जाबालिरथ काश्यपः}
{कात्यायनः सुयज्ञश्च गौतमो विजयस्तथा} %||6-116-55||

\twolineshloka
{अभ्यषिञ्चन्नरव्याघ्रं प्रसन्नेन सुगन्धिना}
{सलिलेन सहस्राक्षं वसवो वासवं यथा} %||6-116-56||

\twolineshloka
{ऋत्विग्भिर्ब्राह्मणैः पूर्वं कन्याभिर्मन्त्रिभिस्तथा}
{योधैश्चैवाभ्यषिञ्चंस्ते सम्प्रहृष्टाः सनैगमैः} %||6-116-57||

\twolineshloka
{सर्वौषधिरसैश्चापि दैवतैर्नभसि स्थितैः}
{चतुर्हिर्लोकपालैश्च सर्वैर्देवैश्च सङ्गतैः} %||6-116-58||

\threelineshloka
{छत्रं तस्य च जग्राह शत्रुघ्नः पाण्डुरं शुभम्}
{श्वेतं च वालव्यजनं सुग्रीवो वानरेश्वरः}
{अपरं चन्द्रसङ्काशं राक्षसेन्द्रो विभीषणः} %||6-116-59||

\twolineshloka
{मालां ज्वलन्तीं वपुषा काञ्चनीं शतपुष्कराम्}
{राघवाय ददौ वायुर्वासवेन प्रचोदितः} %||6-116-60||

\twolineshloka
{सर्वरत्नसमायुक्तं मणिरत्नविभूषितम्}
{मुक्ताहारं नरेन्द्राय ददौ शक्रप्रचोदितः} %||6-116-61||

\twolineshloka
{प्रजगुर्देवगन्धर्वा ननृतुश्चाप्सरो गणाः}
{अभिषेके तदर्हस्य तदा रामस्य धीमतः} %||6-116-62||

\twolineshloka
{भूमिः सस्यवती चैव फलवन्तश्च पादपाः}
{गन्धवन्ति च पुष्पाणि बभूवू राघवोत्सवे} %||6-116-63||

\twolineshloka
{सहस्रशतमश्वानां धेनूनां च गवां तथा}
{ददौ शतं वृषान्पूर्वं द्विजेभ्यो मनुजर्षभः} %||6-116-64||

\twolineshloka
{त्रिंशत्कोटीर्हिरण्यस्य ब्राह्मणेभ्यो ददौ पुनः}
{नानाभरणवस्त्राणि महार्हाणि च राघवः} %||6-116-65||

\twolineshloka
{अर्करश्मिप्रतीकाशां काञ्चनीं मणिविग्रहाम्}
{सुग्रीवाय स्रजं दिव्यां प्रायच्छन्मनुजर्षभः} %||6-116-66||

\twolineshloka
{वैदूर्यमणिचित्रे च वज्ररत्नविभूषिते}
{वालिपुत्राय धृतिमानङ्गदायाङ्गदे ददौ} %||6-116-67||

\twolineshloka
{मणिप्रवरजुष्टं च मुक्ताहारमनुत्तमम्}
{सीतायै प्रददौ रामश्चन्द्ररश्मिसमप्रभम्} %||6-116-68||

\twolineshloka
{अरजे वाससी दिव्ये शुभान्याभरणानि च}
{अवेक्षमाणा वैदेही प्रददौ वायुसूनवे} %||6-116-69||

\twolineshloka
{अवमुच्यात्मनः कण्ठाद्धारं जनकनन्दिनी}
{अवैक्षत हरीन्सर्वान्भर्तारं च मुहुर्मुहुः} %||6-116-70||

\twolineshloka
{तामिङ्गितज्ञः सम्प्रेक्ष्य बभाषे जनकात्मजाम्}
{प्रदेहि सुभगे हारं यस्य तुष्टासि भामिनि} %||6-116-71||

\twolineshloka
{पौरुषं विक्रमो बुद्धिर्यस्मिन्नेतानि नित्यदा}
{ददौ सा वायुपुत्राय तं हारमसितेक्षणा} %||6-116-72||

\twolineshloka
{हनूमांस्तेन हारेण शुशुभे वानरर्षभः}
{चन्द्रांशुचयगौरेण श्वेताभ्रेण यथाचलः} %||6-116-73||

\twolineshloka
{ततो द्विविद मैन्दाभ्यां नीलाय च परन्तपः}
{सर्वान्कामगुणान्वीक्ष्य प्रददौ वसुधाधिपः} %||6-116-74||

\twolineshloka
{सर्ववानरवृद्धाश्च ये चान्ये वानरेश्वराः}
{वासोभिर्भूषणैश्चैव यथार्हं प्रतिपूजिताः} %||6-116-75||

\twolineshloka
{यथार्हं पूजिताः सर्वे कामै रत्नैश्च पुष्कलैर्}
{प्रहृष्टमनसः सर्वे जग्मुरेव यथागतम्} %||6-116-76||

\twolineshloka
{राघवः परमोदारः शशास परया मुदा}
{उवाच लक्ष्मणं रामो धर्मज्ञं धर्मवत्सलः} %||6-116-77||

\fourlineindentedshloka
{आतिष्ठ धर्मज्ञ मया सहेमाम्}
{गां पूर्वराजाध्युषितां बलेन}
{तुल्यं मया त्वं पितृभिर्धृता या}
{तां यौवराज्ये धुरमुद्वहस्व} %||6-116-78||

\fourlineindentedshloka
{सर्वात्मना पर्यनुनीयमानो}
{यदा न सौमित्रिरुपैति योगम्}
{नियुज्यमानो भुवि यौवराज्ये}
{ततोऽभ्यषिञ्चद्भरतं महात्मा} %||6-116-79||

\twolineshloka
{राघवश्चापि धर्मात्मा प्राप्य राज्यमनुत्तमम्}
{ईजे बहुविधैर्यज्ञैः ससुहृद्भ्रातृबान्धवः} %||6-116-80||

\twolineshloka
{पौण्डरीकाश्वमेधाभ्यां वाजपेयेन चासकृत्}
{अन्यैश्च विविधैर्यज्ञैरयजत्पार्थिवर्षभः} %||6-116-81||

\twolineshloka
{राज्यं दशसहस्राणि प्राप्य वर्षाणि राघवः}
{शताश्वमेधानाजह्रे सदश्वान्भूरिदक्षिणान्} %||6-116-82||

\twolineshloka
{आजानुलम्बिबाहुश्च महास्कन्धः प्रतापवान्}
{लक्ष्मणानुचरो रामः पृथिवीमन्वपालयत्} %||6-116-83||

\twolineshloka
{न पर्यदेवन्विधवा न च व्यालकृतं भयम्}
{न व्याधिजं भयं वापि रामे राज्यं प्रशासति} %||6-116-84||

\twolineshloka
{निर्दस्युरभवल्लोको नानर्थः कञ्चिदस्पृशत्}
{न च स्म वृद्धा बालानां प्रेतकार्याणि कुर्वते} %||6-116-85||

\twolineshloka
{सर्वं मुदितमेवासीत्सर्वो धर्मपरोऽभवत्}
{राममेवानुपश्यन्तो नाभ्यहिंसन्परस्परम्} %||6-116-86||

\twolineshloka
{आसन्वर्षसहस्राणि तथा पुत्रसहस्रिणः}
{निरामया विशोकाश्च रामे राज्यं प्रशासति} %||6-116-87||

\twolineshloka
{नित्यपुष्पा नित्यफलास्तरवः स्कन्धविस्तृताः}
{कालवर्षी च पर्जन्यः सुखस्पर्शश्च मारुतः} %||6-116-88||

\twolineshloka
{स्वकर्मसु प्रवर्तन्ते तुष्ठाः स्वैरेव कर्मभिः}
{आसन्प्रजा धर्मपरा रामे शासति नानृताः} %||6-116-89||

\twolineshloka
{सर्वे लक्षणसम्पन्नाः सर्वे धर्मपरायणाः}
{दशवर्षसहस्राणि रामो राज्यमकारयत्} %||6-117-90||


{॥इत्यार्षे श्रीमद्रामायणे वाल्मीकीये आदिकाव्ये युद्धकाण्डे षोडशाधिकशततमः सर्गः॥}

